\documentclass[11pt, a4paper, oneside]{book}

% =====================================================================
% بسته‌های ضروری
% =====================================================================
% توجه: \usepackage[utf8]{inputenc} و \usepackage[T1]{fontenc} حذف شدند
% چون با xelatex/xepersian ناسازگارند و خطای "inputenc package ignored" می‌دهند
\usepackage{amsmath, amssymb, amsfonts}
\usepackage{graphicx}
\usepackage{geometry}
\geometry{margin=2.5cm}
\usepackage{hyperref}
\hypersetup{
	colorlinks=true,
	linkcolor=blue,
	filecolor=magenta,      
	urlcolor=blue,
	pdftitle={QKD Simulation Framework Documentation},
}
\usepackage{listings}
\usepackage{subcaption}
\usepackage{float}
\usepackage{xcolor}
\usepackage{longtable}
\usepackage{booktabs}
\usepackage{titlesec}
\usepackage{tocloft}
\usepackage{fancyhdr}

% =====================================================================
% xepersian باید همیشه آخرین بسته باشد
% =====================================================================
\usepackage{xepersian}
\settextfont{Vazirmatn}

% =====================================================================
% رفع خطای: Font shape `TU/Vazirmatn(0)/m/it' undefined
% علت: فونت Vazirmatn حالت italic ندارد
% این تعریف، italic را به regular هدایت می‌کند
% =====================================================================
\DeclareFontShape{TU}{Vazirmatn(0)}{m}{it}{<-> ssub * Vazirmatn(0)/m/n}{}
\DeclareFontShape{TU}{Vazirmatn(0)}{bx}{it}{<-> ssub * Vazirmatn(0)/bx/n}{}

% =====================================================================
% رفع خطای: Make it at least 13.59999pt (headheight)
% علت: fancyhdr با هدر italic نیاز به ارتفاع بیشتر دارد
% =====================================================================
\setlength{\headheight}{14pt}
\setlength{\headsep}{0.5cm}

% =====================================================================
% تنظیمات هدر و فوتر
% =====================================================================
\pagestyle{fancy}
\fancyhf{}
\fancyhead[L]{\textit{QKD Simulation Framework}}
\fancyhead[R]{\textit{\leftmark}}
\fancyfoot[C]{\thepage}
\renewcommand{\headrulewidth}{0.4pt}

% =====================================================================
% تنظیمات کدنویسی پایتون
% =====================================================================
\definecolor{codegreen}{rgb}{0,0.6,0}
\definecolor{codegray}{rgb}{0.5,0.5,0.5}
\definecolor{codepurple}{rgb}{0.58,0,0.82}
\definecolor{backcolour}{rgb}{0.95,0.95,0.92}

\lstdefinestyle{mystyle}{
	backgroundcolor=\color{backcolour},   
	commentstyle=\color{codegreen},
	keywordstyle=\color{magenta},
	numberstyle=\tiny\color{codegray},
	stringstyle=\color{codepurple},
	basicstyle=\ttfamily\footnotesize\setLTR,   % \setLTR حیاتی برای RTL
	breakatwhitespace=false,         
	breaklines=true,                 
	captionpos=b,                    
	keepspaces=true,                 
	numbers=left,                    
	numbersep=5pt,                  
	showspaces=false,                
	showstringspaces=false,
	showtabs=false,                  
	tabsize=2,
	mathescape=false,                % جلوگیری از تفسیر $ به‌عنوان ریاضی
	extendedchars=true,
	literate={→}{$\rightarrow$}{1}
	{≤}{$\leq$}{1}
	{≥}{$\geq$}{1},
}
\lstset{style=mystyle, language=Python}

% =====================================================================
% رفع Overfull \hbox
% =====================================================================
\setlength{\emergencystretch}{3em}
\hbadness=10000
\hfuzz=2pt

% =====================================================================
% اطلاعات سند
% =====================================================================
\title{\lr{\textbf{Comprehensive Documentation\\ of QKD Simulation Framework}}}

\author{\lr{\textbf{Reza Aramjou}}}
\date{\today}

\begin{document}
	
	% --------------------------- صفحه عنوان ---------------------------
	\maketitle
	\frontmatter
	
	% --------------------------- پیشگفتار ---------------------------
	\addcontentsline{toc}{chapter}{Preface}
	\chapter*{پیشگفتار}
\addcontentsline{toc}{chapter}{پیشگفتار}

\section*{پیدایش یک فریم‌ورک امن کوانتومی}

\begin{sloppypar}
تحول از ایده‌های نظری در توزیع کلید کوانتومی (\lr{Quantum Key Distribution, QKD}) به سامانه‌هایی که بتوان آن‌ها را در سطح پیاده‌سازی، اعتبارسنجی، بهینه‌سازی و نهایتاً استقرار شبکه‌ای بررسی کرد، مستلزم عبور از یک شکاف بنیادی میان دو جهان متفاوت است: از یک سو، جهان فیزیک کوانتومی و نظریهٔ اطلاعات که با صورت‌بندی‌های انتزاعی، کران‌های امنیتی، فروض ایده‌آل و مدل‌های تحلیلی سروکار دارد؛ و از سوی دیگر، جهان مهندسی سامانه و نرم‌افزار که در آن، جزئیات اجرایی، قیود محاسباتی، ناپایداری‌های سخت‌افزاری، ساختار داده، همزمانی، تکرارپذیری و اعتبارسنجی، نقش تعیین‌کننده در قابل‌اعتماد بودن یک نتیجه ایفا می‌کنند. در عمل، بخش عمده‌ای از دشواری QKD نه در فهم اصول اولیهٔ آن، بلکه در پر کردن دقیق، روشمند و قابل دفاع همین فاصله نهفته است.

مستندی که پیش رو قرار دارد، حاصل تلاشی برای پاسخ دادن به این نیاز است: طراحی، توسعه، سازمان‌دهی و اعتبارسنجی یک \textbf{فریم‌ورک شبیه‌سازی QKD} که بتواند به‌عنوان یک بستر پژوهشی جامع، دقیق، توسعه‌پذیر و مهندسی‌شده، رفتار سامانه‌های \lr{Discrete-Variable QKD} را از سطح تولید پالس نوری تا سطح پردازش امنیتی و حتی مسیریابی شبکه‌ای مدل‌سازی کند. این فریم‌ورک صرفاً یک ابزار محاسبه‌گر نرخ کلید نیست، بلکه تلاشی است برای بازسازی یک خط لولهٔ کامل از رخدادهای فیزیکی، آماری، الگوریتمی و معماری که در یک سامانهٔ واقعی QKD حضور دارند. در این معنا، پروژهٔ حاضر نه‌تنها به مدل‌سازی یک پروتکل، بلکه به ساخت یک زیرساخت تحقیقاتی برای مطالعهٔ خانواده‌ای از سناریوهای QKD می‌پردازد.

انگیزهٔ اصلی شکل‌گیری این فریم‌ورک از یک مشاهدهٔ نسبتاً روشن اما عمیق آغاز شد: بسیاری از ابزارها و مدل‌های موجود، اگرچه برای محاسبهٔ تقریبی نرخ کلید، بازتولید چند نمودار مرجع، یا اجرای سناریوهای آموزشی مفیدند، اما برای پژوهش‌های جدی‌تر با محدودیت‌های بنیادی مواجه می‌شوند. بخشی از این محدودیت‌ها از ساده‌سازی بیش از حد لایهٔ فیزیکی ناشی می‌شود؛ برای مثال، مدل‌سازی ایده‌آل یا ناقص آشکارسازها، نادیده‌گرفتن آبشارهای زمان مرده، صرف‌نظر کردن از \lr{afterpulsing}، یا فروکاستن نویز به یک ثابت ساده. بخش دیگر از محدودیت‌ها به لایهٔ امنیتی بازمی‌گردد؛ جایی که برخی ابزارها تنها یک اثبات امنیتی خاص را به‌صورت سخت‌کدگذاری‌شده پیاده‌سازی می‌کنند و در نتیجه، امکان مقایسهٔ دقیق و منصفانهٔ چند اثبات بر روی یک موتور فیزیکی واحد را از بین می‌برند. افزون بر این، در بسیاری از شبیه‌سازهای موجود، تحلیل‌ها در رژیم مجانبی انجام می‌شوند؛ در حالی که در سامانه‌های واقعی، کلید محدود، نوسانات آماری و بودجهٔ امنیتی، بخش جدایی‌ناپذیر مسئله هستند.

فریم‌ورک حاضر برای رفع این مصالحه‌ها توسعه یافت. هدف آن ایجاد محیطی یکپارچه بود که در آن، پژوهشگر بتواند بدون قربانی کردن دقت فیزیکی، انعطاف‌پذیری رمزنگاری، شفافیت معماری و قابلیت توسعه، سناریوهای متنوع QKD را مطالعه کند. از همین رو، طراحی این سامانه از ابتدا بر اساس چند اصل هدایت‌کننده شکل گرفت: تفکیک روشن لایه‌ها، ماژولار بودن، قابلیت اعتبارسنجی تحلیلی، پشتیبانی از سناریوهای \lr{Finite-key}، توانایی مقایسهٔ چند اثبات امنیتی، و امکان توسعه به سطح شبکه‌های چندگرهی. این اصول نه صرفاً ترجیحات مهندسی، بلکه پاسخ‌هایی ضروری به ماهیت میان‌رشته‌ای و حساس QKD بودند.

این پروژه در عین آنکه از نظر مفهومی بر مبانی استاندارد QKD تکیه دارد، به‌صورت خاص برای پاسخ به نیازهای پژوهش تجربی و شبیه‌سازی دقیق طراحی شده است. به همین دلیل، تمرکز آن تنها بر «درست بودن» فرمول‌ها نیست، بلکه بر «درست بودن رفتار سامانه در تعامل اجزا» نیز هست. در عمل، یک اثبات امنیتی دقیق نیز اگر بر داده‌هایی سوار شود که به‌دلیل مدل‌سازی ناقص آشکارساز یا مدیریت نادرست رویدادهای همزمان مخدوش شده‌اند، خروجی معناداری نخواهد داشت. به همین ترتیب، یک موتور فیزیکی دقیق نیز اگر نتواند داده‌های خود را به‌درستی در اختیار لایهٔ پساپردازش قرار دهد، ارزش پژوهشی محدودی خواهد داشت. از این رو، این فریم‌ورک بر پیوستگی دقیق و مهندسی‌شدهٔ کل زنجیره تأکید دارد.

به بیان دیگر، این فریم‌ورک تلاشی است برای آنکه شبیه‌سازی QKD از سطح یک اسکریپت محاسباتی مجزا فراتر رود و به یک محیط پژوهشی قابل اتکا تبدیل شود؛ محیطی که بتواند برای طراحی آزمایش، مقایسهٔ اثبات‌ها، تحلیل حساسیت، بهینه‌سازی پارامتر، بررسی سناریوهای چندگرهی و تولید نتایج بازتولیدپذیر به کار رود. این مستند نیز با همین نگاه تدوین شده است: نه فقط برای توضیح آنچه پیاده‌سازی شده، بلکه برای روشن‌سازی منطق علمی، معماری نرم‌افزاری و مسیر اعتبارسنجی‌ای که این فریم‌ورک بر آن بنا شده است.

\section*{فلسفه معماری}

در طراحی این فریم‌ورک، معماری نرم‌افزار تنها یک مسئلهٔ سازمان‌دهی کد نبوده، بلکه بخشی از منطق علمی پروژه به‌شمار آمده است. هنگامی که یک سامانه قرار است به‌طور همزمان فیزیک منبع، تلفات کانال، دینامیک آشکارساز، آمارداده‌های کلید محدود، اثبات‌های امنیتی متنوع، بهینه‌سازی عددی و حتی رفتار شبکه‌ای را در خود جای دهد، هرگونه آمیختگی نامنظم میان این اجزا به‌سرعت به پیچیدگی کنترل‌ناپذیر، دشواری در اعتبارسنجی و کاهش قابلیت اطمینان منجر خواهد شد. به همین دلیل، فلسفهٔ معماری این پروژه بر پایهٔ \textbf{تفکیک مسئولیت‌ها، ماژولار بودن، صراحت در قراردادهای داده، و استقلال نسبی لایه‌ها} شکل گرفته است.

در هستهٔ این طراحی، تمایز روشنی میان دو ناحیهٔ اصلی برقرار شده است: نخست، \textbf{لایهٔ فیزیکی} که مسئول مدل‌سازی آن چیزی است که در سطح منبع، کانال و آشکارساز رخ می‌دهد؛ و دوم، \textbf{لایهٔ پساپردازش و امنیت} که وظیفهٔ تفسیر داده‌های مشاهده‌شده، غربال‌گری پایه، تخمین پارامتر، اعمال اثبات‌های امنیتی و محاسبهٔ نرخ کلید امن را بر عهده دارد. این جداسازی دو مزیت مهم ایجاد می‌کند. از یک سو، امکان آن را فراهم می‌سازد که مدل‌های فیزیکی بدون دست‌کاری مستقیم در لایهٔ اثبات‌ها توسعه یابند؛ و از سوی دیگر، به پژوهشگر اجازه می‌دهد چندین اثبات امنیتی را بر روی دقیقاً یک داده‌مولد فیزیکی مشترک اجرا کند. به این ترتیب، مقایسه‌ها نه تحت تأثیر تفاوت پیاده‌سازی‌های پراکنده، بلکه بر پایهٔ یک بستر همگن و کنترل‌شده انجام می‌شوند.

یکی از عناصر شاخص این معماری، \textbf{موتور توزیع چندکارگری (\lr{Multi-Worker Dispatch Engine})} است. در شبیه‌سازی‌های رویدادمحور QKD، مقیاس مسئله به‌سرعت بزرگ می‌شود. زمانی که تحلیل در رژیم \lr{Finite-key}، روی تعداد قابل توجهی پالس، با مدل‌های واقع‌گرایانه و مجموعه‌ای از پارامترها انجام می‌شود، هزینهٔ محاسباتی اجرای ترتیبی به‌شدت افزایش می‌یابد. در چنین شرایطی، استفاده از پردازش موازی دیگر یک بهینه‌سازی لوکس نیست، بلکه یک الزام عملی است. با این حال، موازی‌سازی در این حوزه با چالش‌هایی همراه است که در بسیاری از سامانه‌های ساده نادیده گرفته می‌شوند. مهم‌ترین این چالش‌ها، حفظ \textbf{یکپارچگی هویت پالس‌ها و سازگاری رویدادها} در محیطی است که اقلام کاری روی کارگرهای مختلف، با زمان‌بندی‌های متفاوت و احتمالاً با تأخیرهای ناهمگون پردازش می‌شوند.

برای پاسخ به این مسئله، معماری فریم‌ورک از مکانیزمی مبتنی بر \lr{\texttt{work item}}های مستقل و خودبسنده بهره می‌برد. هر قلم کاری، اطلاعات کافی دربارهٔ پالس، پایه، شدت، حالت و زمینهٔ آماری خود را به همراه دارد و به‌گونه‌ای تعریف می‌شود که وابستگی پنهان یا شکننده به وضعیت اجرایی دیگر بخش‌ها نداشته باشد. این تصمیم معماری باعث می‌شود که نتایج به ترتیب زمان‌بندی کارگرها وابسته نشوند و پردازش موازی به تغییر خاموش در معنای داده‌ها منجر نگردد. در یک سامانهٔ رمزنگاری یا شبه‌رمزنگاری، این موضوع اهمیتی فراتر از کارایی دارد؛ زیرا هرگونه ناسازگاری خاموش در هویت پالس‌ها یا تعلق رویدادها می‌تواند مستقیماً به تفسیر نادرست آمار امنیتی بینجامد.

از منظر طراحی داده، این فریم‌ورک تلاش می‌کند تا حد ممکن از قراردادهای صریح، ساختارهای دادهٔ مشخص و مسیرهای پردازش قابل رهگیری استفاده کند. این امر صرفاً برای زیبایی کد یا سهولت توسعه نیست؛ بلکه برای آن است که یک نتیجهٔ عددی نهایی، در صورت نیاز، قابل ردیابی به مراحل پیشین خود باشد. پژوهشگر باید بتواند بفهمد یک نرخ کلید یا یک نرخ خطا چگونه از میان لایه‌های مختلف شکل گرفته است، چه عواملی بر آن اثر گذاشته‌اند، و در صورت بروز انحراف، منشأ آن در کدام بخش قرار دارد. در اینجا، معماری خوب عملاً نقش ابزار تبیین علمی را نیز ایفا می‌کند.

فلسفهٔ معماری پروژه همچنین بر \textbf{توسعه‌پذیری کنترل‌شده} تأکید دارد. در سامانه‌های پژوهشی، نیازها ایستا نیستند. ممکن است در نسخه‌های بعدی، مدل جدیدی از منبع نوری، یک تقریب بهتر برای \lr{afterpulsing}، یک اثبات امنیتی متفاوت، یا یک سیاست تازه برای مدیریت شبکه مورد نیاز باشد. اگر افزودن هر قابلیت جدید مستلزم بازنویسی چندین ماژول و شکستن فرضیات قبلی باشد، سامانه خیلی زود به نقطهٔ فرسودگی می‌رسد. از این رو، طراحی کلاس‌های پایه، لایه‌های انتزاعی و رابط‌های داخلی به‌گونه‌ای انجام شده است که افزودن جزء جدید، عمدتاً مستلزم پیاده‌سازی همان جزء باشد، نه بازآرایی کل سیستم.

در نهایت، فلسفهٔ این معماری را می‌توان چنین خلاصه کرد: \textit{دقت علمی تنها زمانی ارزشمند است که در یک ساختار مهندسی‌شده، شفاف، تکرارپذیر و قابل توسعه استقرار یابد.} این فریم‌ورک تلاش کرده است دقیقاً چنین بستری را فراهم کند؛ بستری که در آن، شبیه‌سازی QKD از یک محاسبهٔ موردی و محدود فراتر رود و به یک دارایی پژوهشی و نرم‌افزاری پایدار تبدیل شود.

\section*{فرآیند جامع اعتبارسنجی}

هیچ فریم‌ورک شبیه‌سازی، صرف‌نظر از پیچیدگی ظاهری، تنوع قابلیت‌ها یا ظرافت معماری، بدون یک فرایند اعتبارسنجی دقیق و چندلایه، شایستهٔ اعتماد علمی نخواهد بود. این گزاره در مورد سامانه‌های QKD اهمیتی دوچندان دارد؛ زیرا خطا در چنین محیطی تنها یک اختلاف عددی ساده نیست، بلکه می‌تواند به تفسیر نادرست امنیت، خوش‌بینی کاذب نسبت به نرخ کلید، یا مقایسه‌های گمراه‌کننده میان مدل‌ها و اثبات‌ها منجر شود. از همین رو، اعتبارسنجی در این پروژه نه یک مرحلهٔ پسینی و تکمیلی، بلکه بخشی مرکزی از فرایند توسعه بوده است.

رویکرد اعتبارسنجی این فریم‌ورک بر چند اصل بنا شده است. نخست آنکه اعتبارسنجی نباید به سطح تست واحد توابع محدود شود، هرچند این تست‌ها لازم‌اند. دوم آنکه هر لایه باید تا حد امکان در برابر یک مرجع مستقل سنجیده شود: لایهٔ فیزیک در برابر روابط تحلیلی و پیش‌بینی‌های فرم‌بسته، لایهٔ امنیتی در برابر خواص مورد انتظار اثبات‌ها و رفتار مرزی آن‌ها، و لایهٔ شبکه در برابر منطق ساختاری سناریوهای مسیریابی و تأمین کلید. سوم آنکه اعتبارسنجی باید هم در سطح محلی و هم در سطح برهم‌کنش میان لایه‌ها انجام شود؛ زیرا بسیاری از خطاهای مهم نه در خود اجزا، بلکه در مرزهای میان آن‌ها پدیدار می‌شوند.

در این پروژه، اعتبارسنجی با اتکا به \textbf{بنچمارک‌های تحلیلی حقیقتِ پایه (\lr{Ground-Truth Analytical Benchmarks})}، آزمون‌های آماری، تست‌های واحد، تست‌های یکپارچه‌سازی، و سناریوهای کامل اجرایی صورت گرفته است. هدف از این ترکیب آن بوده است که از یک سو، اجزای منفرد به‌صورت ایزوله کنترل شوند، و از سوی دیگر، عملکرد کل خط لوله نیز در سناریوهای واقع‌گرایانه سنجیده شود.

\subsection*{۱. یکپارچگی لایه فیزیکی}

نخستین و بنیادی‌ترین لایه‌ای که باید اعتبارسنجی شود، لایهٔ فیزیکی است. در این لایه، منبع نوری، کانال انتقال و آشکارسازها مجموعه‌ای از رخدادها را تولید می‌کنند که تمام مراحل بعدی بر آن‌ها استوار است. اگر این لایه دچار انحراف مفهومی یا عددی باشد، حتی دقیق‌ترین اثبات امنیتی نیز بر ورودی نادرست سوار خواهد شد. به همین دلیل، بخش مهمی از فرایند توسعه به اطمینان از صحت این لایه اختصاص یافت.

در سطح منبع، آمار نشر پالس‌ها برای مدل‌های مورد استفاده—از جمله توزیع پواسونی و در صورت نیاز مدل‌های حرارتی—در برابر روابط تحلیلی شناخته‌شده بررسی شد. در سطح کانال، مدل تضعیف فیبر نوری با تکیه بر قانون بیر–لمبرت و پارامترهای استاندارد انتقال بررسی گردید. در سطح آشکارسازی، رفتار نویز شمارش تاریک، احتمال کلیک، زمان مردهٔ غیرفلج‌شونده (\lr{Non-paralyzable Dead Time}) و آثار وابسته به آن با مدل‌های مرجع تطبیق داده شد.

برای آنکه این تطبیق تنها به مقایسهٔ بصری یا تقریبی محدود نماند، از چارچوب‌های آماری صریح استفاده شد. به‌طور خاص، با به‌کارگیری یک آزمون آماری \(Z\) دوجمله‌ای و اتخاذ سطح معنی‌داری سخت‌گیرانهٔ \(\alpha = 10^{-4}\)، بررسی شد که کمیت‌هایی مانند بهرهٔ کل یا \(Q_\mu\)، و نرخ خطای بیت کوانتومی (\lr{QBER})، در فواصل و سناریوهای مختلف، در تلورانس تعریف‌شده با پیش‌بینی‌های نظری هم‌خوان باشند. نتیجهٔ این فرایند نشان داد که موتور فیزیکی، در محدوده‌های عملی مورد مطالعه، رفتار مورد انتظار را با دقت قابل قبول بازتولید می‌کند. این موضوع اهمیت زیادی دارد، زیرا نشان می‌دهد انحرافات مشاهده‌شده در مراحل بعد، در صورت وقوع، احتمالاً ناشی از بخش‌های دیگر زنجیره‌اند، نه از لایهٔ پایهٔ فیزیکی.

\subsection*{۲. سازگارسازی چند‌اثباتی رمزنگاری}

یکی از چالش‌های مهم در توسعهٔ فریم‌ورک‌های QKD آن است که بسیاری از پیاده‌سازی‌ها از همان ابتدا حول یک اثبات امنیتی خاص طراحی می‌شوند. در چنین وضعیتی، مسیر داده، ساختار آماری، توابع کمکی و حتی نام‌گذاری متغیرها به‌شکلی ناخواسته به منطق آن اثبات خاص وابسته می‌شود. نتیجه آن است که افزودن اثباتی دیگر، به‌جای آنکه صرفاً توسعهٔ یک ماژول جدید باشد، به بازنویسی بخش‌هایی از کل خط لوله نیاز پیدا می‌کند. این وابستگی پنهان، یکی از موانع اصلی برای مقایسهٔ دقیق اثبات‌های امنیتی در محیط‌های فیزیکی مشترک است.

در نسخه‌های اولیه‌تر این پروژه نیز خط لولهٔ اجرایی تا حدی به اثبات کلید محدود \texttt{Lim2014} گره خورده بود. یکی از مهم‌ترین نقاط عطف در بلوغ معماری، معرفی یک \textbf{لایهٔ سازگارساز (\lr{Adapter Layer})} در اسکریپت اجرایی اصلی بود. این لایه نقش مترجم ساختاری را ایفا می‌کند: داده‌ها و آمارهای خام یا نیمه‌پردازش‌شدهٔ حاصل از شبیه‌سازی را دریافت می‌کند و آن‌ها را به امضاها، ساختارها و قراردادهای مورد انتظار هر اثبات امنیتی تبدیل می‌نماید. به این ترتیب، منطق فیزیکی و منطق اثباتی از یکدیگر جدا می‌شوند، در حالی که همچنان از یک زبان دادهٔ مشترک بهره می‌برند.

بر پایهٔ این بازآرایی معماری، فریم‌ورک توانست چند خانوادهٔ اثباتی را با موفقیت پشتیبانی و اجرا کند:
\begin{itemize}
	\item \textbf{2014 Lim}: به‌عنوان یکی از چارچوب‌های مهم \texttt{Finite-key}، این اثبات توانست در مقیاس‌هایی نظیر \(N = 10^7\)، نرخ‌های کلید محدود بهینه و معناداری تولید کند و به‌عنوان مرجع اصلی در بسیاری از مقایسه‌ها به‌کار رود.
	\item \textbf{2005 Wang}: این اثبات نرخ‌های کلیدی متمایزتر و در مواردی محافظه‌کارانه‌تر تولید کرد و نشان داد که تفاوت در روش کران‌بندی آماری و مفروضات تحلیلی چگونه مستقیماً بر خروجی نهایی اثر می‌گذارد.
	\item \textbf{2005 Ma}: اجرای این اثبات، به‌ویژه در نواحی‌ای که مفروضات گاوسی یا کفایت آماری به‌خوبی برقرار نیستند، به‌درستی رفتار محافظه‌کارانه نشان داد و در برخی رژیم‌ها منجر به نرخ کلید صفر شد. همین رفتار، خود نشانه‌ای از صحت سازگارسازی است؛ زیرا سامانه در این حالت تلاش نکرده است به‌صورت مصنوعی نرخ مثبت تولید کند.
	\item \textbf{Proof Tight}: این شاخه با بهره‌گیری از ساختار Lim2014 و در عین حال اعمال تخمین‌های بهینه‌تر برای برخی کمیت‌ها، رفتار متفاوتی را در تحلیل امنیتی ارائه داد و نشان داد که فریم‌ورک آمادهٔ مقایسهٔ نسخه‌های سخت‌گیرانه‌تر یا دقیق‌تر از اثبات‌های مرجع است.
\end{itemize}

اهمیت این دستاورد تنها در «اجرا شدن» چند اثبات خلاصه نمی‌شود. ارزش اصلی آن در این است که پژوهشگر اکنون می‌تواند دقیقاً یک سناریوی فیزیکی واحد را تحت چند لنز امنیتی مختلف مشاهده و تحلیل کند. این قابلیت، برای بررسی اثر انتخاب اثبات بر نرخ کلید، مقایسهٔ محافظه‌کاری یا تهاجمی بودن کران‌ها، و مطالعهٔ حساسیت خروجی به فروض آماری، اهمیت پژوهشی بالایی دارد.

\subsection*{۳. بهینه‌سازی پارامترها با برنامه‌ریزی خطی (LP)}

سامانه‌های QKD ذاتاً به پارامترهای عملیاتی خود حساس‌اند. شدت سیگنال (\(\mu\))، شدت‌های فریب (\(\nu\))، احتمال انتخاب هر شدت، نسبت پایه‌ها، و گاه حتی تنظیمات اجرایی آشکارساز یا پنجره‌های گیتینگ، همگی می‌توانند به‌طور مستقیم و غیرمستقیم بر نرخ کلید امن، QBER، بازده و برد نهایی سامانه اثر بگذارند. در بسیاری از مطالعات ساده، این پارامترها یا به‌صورت دستی تنظیم می‌شوند یا از مقادیر مرجع در مقاله‌های پیشین اقتباس می‌گردند. با این حال، چنین رویکردی معمولاً برای دستیابی به بهترین عملکرد در یک سناریوی مشخص کافی نیست.

در این فریم‌ورک، یک لایهٔ بهینه‌سازی برای جست‌وجوی نظام‌مند در فضای پارامترها توسعه یافته است. این لایه با استفاده از حلگرهای عددی مانند \lr{Nelder-Mead} و همچنین ابزارهای مبتنی بر \textbf{برنامه‌ریزی خطی (\lr{LP})}، می‌تواند فضای شدت‌ها و احتمال‌ها را پیمایش کرده و پیکربندی‌هایی را بیابد که نرخ کلید امن را در شرایط معین بیشینه می‌کنند. اعتبارسنجی این بخش صرفاً به بررسی همگرایی عددی محدود نشد، بلکه خروجی‌های بهینه‌سازی با سناریوهای مرجع و مقادیر غیر‌بهینه مقایسه شدند تا اثر واقعی بهینه‌سازی بر عملکرد نهایی سنجیده شود.

نتایج این بررسی‌ها نشان دادند که حلگر LP و لایهٔ بهینه‌سازی نه‌تنها به‌درستی در فضای پارامتر حرکت می‌کنند، بلکه در بسیاری از فواصل عملی—به‌ویژه در بازهٔ ۰ تا ۱۰۰ کیلومتر—توانسته‌اند پیکربندی‌هایی بیابند که به \textbf{افزایش بسیار معنادار در نرخ کلید امن} نسبت به مقادیر ثابت و سخت‌کدگذاری‌شده منجر می‌شوند. این نتیجه از منظر مهندسی اهمیت زیادی دارد، زیرا نشان می‌دهد بخشی از شکاف عملکردی میان یک سامانهٔ «قابل اجرا» و یک سامانهٔ «خوب تنظیم‌شده» را می‌توان به‌صورت محاسباتی و سیستماتیک کاهش داد. از منظر پژوهشی نیز، وجود این لایه باعث می‌شود که مقایسهٔ اثبات‌ها یا مدل‌ها بر اساس پارامترهای منصفانه‌تر و نزدیک‌تر به بهینه انجام شود، نه صرفاً بر اساس انتخاب‌های دستی.

\subsection*{۴. توپولوژی شبکه و رله‌های قابل اعتماد}

بخش بزرگی از ادبیات کلاسیک QKD بر لینک‌های نقطه‌به‌نقطه متمرکز است؛ حال آنکه در کاربردهای عملی، به‌ویژه زمانی که هدف فراتر رفتن از اتصال مستقیم دو نقطه باشد، مسئلهٔ شبکه‌سازی کوانتومی مطرح می‌شود. در چنین محیطی، دیگر تنها نرخ کلید یک لینک خاص مهم نیست، بلکه نحوهٔ توزیع و مصرف کلیدها در سراسر توپولوژی، مدیریت لینک‌های میانی، تاب‌آوری در برابر خرابی، و منطق مسیریابی نیز به مسئله افزوده می‌شود. به همین دلیل، یکی از ابعاد مهم اعتبارسنجی این پروژه، توسعه و بررسی یک لایهٔ شبکه بر بستر موتور فیزیکی و امنیتی بود.

اوج این مرحله از اعتبارسنجی در شبیه‌سازی یک \textbf{توپولوژی ۷ گرهی} حاصل شد؛ توپولوژی‌ای که شامل فرستنده‌ها، گیرنده‌ها و رله‌های قابل اعتماد (\lr{Trusted Relays}) بوده و امکان بررسی سناریوهای چندمرحله‌ای را فراهم می‌کرد. در این سناریو، موتور شبکه باید علاوه بر مدیریت تولید کلید روی لینک‌های منفرد، بتواند از این کلیدها در سطح بین‌گرهی نیز استفاده کند. این مسئله شامل نگهداری و به‌روزرسانی \textit{استخرهای کلید} (\lr{Key Pools})، ارزیابی هزینهٔ مسیرها، تصمیم‌گیری برای انتخاب مسیر مناسب، و در صورت لزوم، رله‌گذاری مرحله‌به‌مرحلهٔ کلید میان گره‌هاست.

اعتبارسنجی این لایه نشان داد که سامانه توانایی انجام موارد زیر را دارد:
\begin{itemize}
	\item اجرای دقیق منطق \textbf{کوتاه‌ترین مسیر اول (\lr{SPF+})} بر اساس هزینهٔ لینک و وضعیت استخر کلیدها؛
	\item انجام \textbf{رله‌گذاری مرحله‌به‌مرحله (\lr{Hop-by-Hop Relay})} و مدیریت ترکیب کلید در مسیرهای چندگامی؛
	\item نمایش \textbf{Failover پویا}، به این معنا که در صورت حذف یا قطع عمدی یک لینک اصلی، سامانه بتواند مسیر جایگزین مناسبی را شناسایی و سرویس توزیع کلید را از مسیر جدید ادامه دهد؛
	\item تغذیه و بازیابی خودکار استخرهای کلید برای جلوگیری از فروپاشی عملکرد شبکه در بارهای متغیر.
\end{itemize}

این سطح از اعتبارسنجی نشان می‌دهد که فریم‌ورک حاضر صرفاً ابزاری برای تولید چند شاخص فیزیکی نیست، بلکه می‌تواند در مطالعهٔ معماری‌های گسترده‌تر QKD و رفتار شبکه‌ای نیز نقش ایفا کند. این ویژگی به‌ویژه برای پژوهش‌هایی که به استقرار شبکه‌های کوانتومی کاربردی، مدیریت کلید در مقیاس چندگرهی، یا ارزیابی تاب‌آوری مسیرها علاقه‌مندند، اهمیت ویژه دارد.

\section*{تاب‌آوری مهندسی}

هر پروژهٔ نرم‌افزاری که به‌طور جدی در معرض توسعه، بازطراحی، افزودن قابلیت و اعتبارسنجی مستمر قرار گیرد، دیر یا زود با موارد لبه، تضادهای معماری، یا ناهماهنگی‌های پنهانی روبه‌رو می‌شود. در سامانه‌ای مانند این فریم‌ورک، که در آن لایه‌های فیزیک، آمار، امنیت، همزمانی و شبکه در تعامل‌اند، چنین مواردی نه استثنا، بلکه بخشی طبیعی از مسیر بلوغ مهندسی هستند. تفاوت اصلی میان یک نمونهٔ آزمایشگاهی و یک ابزار پژوهشی قابل اتکا، تا حد زیادی در نحوهٔ مواجهه با همین جزئیات نهفته است.

در طول توسعه و اعتبارسنجی این فریم‌ورک، چندین مورد لبه شناسایی و اصلاح شدند که رفع آن‌ها نقشی مهم در افزایش تاب‌آوری سیستم ایفا کرد. بخشی از این اصلاحات به \textbf{تضادهای ارث‌بری} در امضاهای متدهای برخی اثبات‌ها، از جمله شاخهٔ \texttt{Tight}، مربوط می‌شد. اگرچه این دسته از مسائل ممکن است در نگاه اول صرفاً به سازگاری کدنویسی مربوط به نظر برسند، اما در عمل می‌توانند منجر به اجرای نادرست متدها، تفسیر غلط داده‌های ورودی، یا ناسازگاری خاموش در مسیرهای اجرایی شوند. رفع این تعارض‌ها باعث شد قراردادهای میان کلاس‌های پایه و مشتق‌شده شفاف‌تر و پایدارتر شوند.

بخش دیگری از اصلاحات به \textbf{سریال‌سازی و بازیابی وضعیت}، به‌ویژه در سطح آشکارساز، مربوط بود. هنگامی که یک سامانهٔ شبیه‌سازی قرار است حالت اجزای خود را ذخیره و بازیابی کند، عدم کنترل دقیق نسخهٔ طرح‌واره (\lr{Schema Version}) می‌تواند به برش خاموش داده‌ها، جابه‌جایی فیلدها، یا بازیابی ناقص وضعیت منجر شود. به همین دلیل، بررسی‌های سخت‌گیرانه‌تری برای سازگاری نسخهٔ طرح‌واره و درستی بازخوانی داده‌ها اعمال شد. این موضوع برای سناریوهایی که در آن‌ها تداوم وضعیت آشکارساز، carry-over رویدادها یا ادامهٔ اجرای دسته‌ای اهمیت دارد، حیاتی است.

همچنین، برخی اصلاحات به \textbf{چینش فیلدهای تغییرناپذیر در ساختارهای داده} و نحوهٔ پذیرش پارامترهای امنیتی چندگانه بازمی‌گشت. در ساختارهای مبتنی بر \textit{dataclass}، کوچک‌ترین ناسازگاری در ترتیب فیلدها یا قرارداد سازنده می‌تواند در نسخه‌های مختلف یا هنگام گسترش مدل، به رفتارهای شکننده منتهی شود. بازتنظیم این بخش‌ها، انعطاف‌پذیری سیستم در افزودن پارامترهای جدید را افزایش داد و احتمال بروز خطاهای پنهان در ساخت اشیا را کاهش داد.

اهمیت این اصلاحات در آن است که نشان می‌دهند تاب‌آوری یک فریم‌ورک، تنها از طریق داشتن مدل‌های نظری خوب حاصل نمی‌شود. یک ابزار پژوهشی باید در برابر ورودی‌های مرزی، سناریوهای غیرمعمول، توسعهٔ تدریجی و تنوع مسیرهای اجرا نیز پایدار باشد. از این منظر، تاب‌آوری مهندسی بخشی از اعتبار علمی پروژه است؛ زیرا نتایجی که بر بستر نرم‌افزار شکننده تولید شوند، حتی اگر از نظر فرمول‌بندی صحیح باشند، در سطح اعتمادپذیری پژوهشی با محدودیت مواجه خواهند بود.

\section*{راهنمای مطالعه این مستند}

این مستند با این فرض تدوین شده است که خواننده ممکن است از یکی از چند مسیر متفاوت وارد آن شود: از مسیر نظریهٔ \texttt{QKD}، از مسیر پیاده‌سازی نرم‌افزار، از مسیر اعتبارسنجی پژوهشی، یا از مسیر کاربردهای شبکه‌ای. به همین دلیل، ساختار سند به‌گونه‌ای طراحی شده که هم بتوان آن را به‌صورت ترتیبی از ابتدا تا انتها مطالعه کرد و هم بتوان از آن به‌عنوان یک مرجع موضوع‌محور استفاده نمود.

به‌طور کلی، این سند در چند بخش اصلی سازمان‌دهی شده است:

\begin{itemize}
	\item \textbf{بخش اول: معماری فریم‌ورک و فیزیک پایه} \\
	در این بخش، مبانی مفهومی و محاسباتی لایهٔ فیزیکی معرفی می‌شوند. مدل‌های منبع نوری، کانال، تلفات، نویز، و آشکارسازها در این قسمت تشریح می‌گردند و روشن می‌شود که هر یک چگونه در قالب ساختار نرم‌افزاری پروژه پیاده‌سازی شده‌اند.
	
	\item \textbf{بخش دوم: اثبات‌های امنیتی رمزنگاری} \\
	این بخش به منطق امنیتی فریم‌ورک اختصاص دارد. در آن، چارچوب‌های پشتیبانی‌شده برای تخمین پارامتر، کران‌بندی آماری و محاسبهٔ نرخ کلید امن ارائه می‌شوند. افزون بر فرمول‌بندی نظری، ارتباط این اثبات‌ها با رابط‌های نرم‌افزاری و داده‌های ورودی نیز توضیح داده می‌شود.
	
	\item \textbf{بخش سوم: خط لوله اجرایی و لایه شبکه} \\
	در این بخش، مسیر کامل پردازش داده از مرحلهٔ اجرای شبیه‌سازی تا تولید خروجی و همچنین منطق موتور توزیع، کارگرها، هماهنگی موازی و قابلیت‌های شبکه‌ای تشریح می‌شود. این بخش برای فهم رفتار سامانه در مقیاس بزرگ و سناریوهای چندگرهی اهمیت ویژه دارد.
	
	\item \textbf{بخش چهارم: اعتبارسنجی و تست} \\
	در این قسمت، بنچمارک‌های تحلیلی، روش‌های آزمون، نتایج اعتبارسنجی، و سناریوهای تستی که برای راستی‌آزمایی فریم‌ورک به‌کار رفته‌اند، ارائه می‌شود. این بخش برای خوانندگانی که به اعتمادپذیری نتایج و بازتولیدپذیری علاقه‌مندند، اهمیت محوری دارد.
\end{itemize}

در مجموع، هدف از این ساختاربندی آن است که مستند حاضر صرفاً توصیفی از یک کدبیس نباشد، بلکه پلی میان \textbf{مبانی نظری QKD}، \textbf{مهندسی دقیق نرم‌افزار} و \textbf{اعتبارسنجی پژوهشی} ایجاد کند. به بیان دیگر، این سند باید هم برای کسی که می‌خواهد منطق فیزیکی و امنیتی فریم‌ورک را درک کند مفید باشد، هم برای توسعه‌دهنده‌ای که می‌خواهد آن را گسترش دهد، و هم برای پژوهشگری که می‌خواهد به نتایج آن استناد کند یا آن‌ها را بازتولید نماید.

امید آن است که این فریم‌ورک بتواند به‌عنوان ابزاری \textbf{قابل‌اعتماد، شفاف، توسعه‌پذیر و بازتولیدپذیر} در اختیار پژوهشگران حوزهٔ رمزنگاری کوانتومی قرار گیرد؛ ابزاری که نه‌تنها برای تحلیل یک سناریوی منفرد، بلکه برای ساختن و ارزیابی نسل بعدی مطالعات در حوزهٔ استقرار سامانه‌ها و شبکه‌های کوانتومی کاربردی مفید باشد. اگر این پروژه بتواند حتی بخشی از مسیر میان صورت‌بندی‌های نظری QKD و سامانه‌های عملی و مهندسی‌شدهٔ آن را روشن‌تر کند، به هدف اصلی خود نزدیک شده است.
\end{sloppypar}
\vspace{2cm}
\begin{flushleft}
	\textit{رضا آرام جو}\\
	\textit{\today}
\end{flushleft}

	
	% --------------------------- فهرست مطالب ---------------------------
	\tableofcontents
	\listoftables
	\lstlistoflistings
	\newpage
	
	% =========================== بخش اول: معماری ===========================
	\mainmatter
	\part{معماری چارچوب و فیزیک بنیادی}
	
	\chapter{مقدمه و ساختار پروژه}
\label{ch:intro}

\section{زمینه و انگیزه}

امنیت اطلاعات در دهه‌های اخیر به یکی از ارکان اصلی زیرساخت‌های دیجیتال، مخابراتی، مالی، صنعتی و حیاتی تبدیل شده است. بسیاری از سامانه‌های ارتباطی مدرن، از تبادلات بانکی و شبکه‌های داده گرفته تا زیرساخت‌های ابری، سامانه‌های کنترلی و ارتباطات حساس، بر این فرض تکیه دارند که داده‌ها می‌توانند با استفاده از کلیدهای رمزنگاری امن میان دو طرف محافظت و مبادله شوند.
 در رمزنگاری کلاسیک، به‌ویژه در شاخهٔ رمزنگاری کلید عمومی، امنیت بسیاری از سازوکارهای رایج نه از قوانین بنیادین فیزیک، بلکه از دشواری محاسباتی برخی مسائل ریاضی حاصل می‌شود. برای مثال، امنیت الگوریتم‌هایی نظیر \lr{RSA} به دشواری تجزیهٔ اعداد صحیح بزرگ به عوامل اول و امنیت روش‌هایی مانند \lr{ECC} به دشواری حل مسئلهٔ لگاریتم گسسته روی منحنی‌های بیضوی وابسته است.

با وجود موفقیت گستردهٔ این خانواده از الگوریتم‌ها در کاربردهای عملی، این نوع امنیت ذاتاً \textit{امنیت محاسباتی} است، نه \textit{امنیت اطلاعات‌نظری}. به بیان دیگر، این سامانه‌ها تا زمانی امن تلقی می‌شوند که هیچ مهاجم عملی، در بازهٔ زمانی قابل‌قبول، توان محاسباتی کافی برای حل مسائل زیربنایی آن‌ها را در اختیار نداشته باشد. این فرض برای دهه‌ها مبنای طراحی زیرساخت‌های امنیتی بوده است، اما ظهور محاسبات کوانتومی، به‌ویژه الگوریتم شور یا \lr{Shor's algorithm}، نشان داد که یک رایانهٔ کوانتومی مقیاس‌پذیر و دارای تصحیح خطای کافی، از نظر نظری قادر است مسائل فاکتورگیری و لگاریتم گسسته را در زمان چندجمله‌ای حل کند.

تهدید محاسبات کوانتومی تنها به زمان دسترسی عملی به چنین رایانه‌هایی محدود نیست. یکی از نگرانی‌های مهم در امنیت بلندمدت، سناریوی \lr{Harvest Now, Decrypt Later} است؛ به این معنا که مهاجم می‌تواند امروز داده‌های رمزشده را ذخیره کند و در آینده، پس از دستیابی به توان پردازشی لازم، آن‌ها را رمزگشایی کند. بنابراین، حتی اگر رایانه‌های کوانتومی قدرتمند هنوز به‌صورت عملی و گسترده در دسترس نباشند، داده‌هایی که امروز با سازوکارهای سنتی منتقل می‌شوند، ممکن است در آینده محرمانگی بلندمدت خود را از دست بدهند. همین واقعیت، نیاز به روش‌هایی را برجسته می‌کند که امنیت آن‌ها تا حد امکان به توان محاسباتی مهاجم وابسته نباشد.

البته توزیع کلید کوانتومی تنها مسیر پاسخ به تهدید محاسبات کوانتومی نیست. رمزنگاری پساکوانتومی یا \lr{Post-Quantum Cryptography (PQC)} نیز با هدف طراحی الگوریتم‌های کلاسیکی مقاوم در برابر حملات کوانتومی توسعه یافته است. تفاوت اصلی این رویکرد با توزیع کلید کوانتومی در آن است که \lr{PQC} همچنان بر دشواری محاسباتی مسائل ریاضی جدید تکیه دارد، در حالی که توزیع کلید کوانتومی مسئلهٔ توزیع کلید را با استفاده از خواص فیزیکی سامانه‌های کوانتومی بررسی می‌کند. در عمل، این دو رویکرد می‌توانند مکمل یکدیگر باشند؛ به‌ویژه از آن جهت که سامانه‌های \lr{QKD} نیز برای کانال کلاسیک خود به احراز هویت نیاز دارند.

در این بستر، توزیع کلید کوانتومی یا \lr{Quantum Key Distribution (QKD)} به‌عنوان یکی از مهم‌ترین دستاوردهای نظری و عملی در تقاطع فیزیک کوانتومی، نظریهٔ اطلاعات و امنیت ارتباطات مطرح می‌شود. ایدهٔ اصلی \lr{QKD} آن است که دو طرف ارتباطی، که به‌طور قراردادی آلیس (\lr{Alice}) و باب (\lr{Bob}) نامیده می‌شوند، با ارسال حالت‌های کوانتومی و انجام پردازش کلاسیک احراز هویت‌شده، یک کلید محرمانهٔ مشترک تولید کنند. امنیت \lr{QKD} در چارچوب یک مدل امنیتی مشخص و تحت فروض پیاده‌سازی معین تحلیل می‌شود. در این چارچوب، تلاش مهاجم برای کسب اطلاعات از حالت‌های کوانتومی می‌تواند میان اطلاعات مهاجم و اختلال قابل‌مشاهده در داده‌های آلیس و باب مصالحه ایجاد کند. این اختلال از طریق آزمون آماری داده‌های نمونه‌برداری‌شده ارزیابی می‌شود.

این ویژگی، \lr{QKD} را به‌ویژه از سازوکارهای رایج توزیع کلید مبتنی بر رمزنگاری کلید عمومی متمایز می‌کند؛ زیرا امنیت آن، در صورت برقراری فروض مدل امنیتی و وجود یک کانال کلاسیک احراز هویت‌شده، نه بر فرضیات پیچیدگی محاسباتی، بلکه بر اصولی مانند اثر اندازه‌گیری بر حالت کوانتومی، اصل عدم قطعیت و قضیهٔ عدم کلونینگ استوار است. با این حال، این بیان نباید به معنای امنیت مطلق و مستقل از پیاده‌سازی تلقی شود. در سامانه‌های واقعی، اعتبار مدل منبع، آشکارساز، کانال، تصادفی‌سازی، نبود کانال‌های جانبی و نحوهٔ تخصیص پارامترهای امنیتی همگی در نتیجهٔ نهایی اثرگذارند.

در تحلیل‌های مدرن، امنیت کلید نهایی معمولاً با پارامتر شکست \(\epsilon\) بیان می‌شود. به‌طور شهودی، کلید \(\epsilon\)-امن کلیدی است که میزان تمایزپذیری آن از یک کلید ایده‌آلِ یکنواخت و مستقل از اطلاعات مهاجم، با پارامتر \(\epsilon\) کران‌بندی شود. این پارامتر معمولاً احتمال شکست یا عدم تحقق کامل شرایط امنیتی مورد نظر را به‌صورت ترکیبی نمایش می‌دهد. این نگاه به‌ویژه در تحلیل‌های کلید محدود اهمیت دارد، زیرا در سامانه‌های عملی تعداد نمونه‌ها محدود است و نوسانات آماری نمی‌توانند نادیده گرفته شوند.

\section{چالش گذار از \lr{QKD} نظری به \lr{QKD} عملی}

پروتکل \lr{BB84} که در سال ۱۹۸۴ توسط \lr{Bennett} و \lr{Brassard} معرفی شد، نخستین و شناخته‌شده‌ترین پروتکل \lr{QKD} است و همچنان یکی از چارچوب‌های مرجع در تحلیل‌های نظری و پیاده‌سازی‌های عملی به‌شمار می‌رود. در این پروتکل، آلیس بیت‌های اطلاعاتی را در یکی از دو پایهٔ ناسازگار، که در صورت‌بندی ایده‌آل می‌توان آن‌ها را دو پایهٔ متقابلاً نااریب دانست، کدگذاری کرده و برای باب ارسال می‌کند. باب نیز بدون اطلاع از انتخاب پایهٔ آلیس، به‌صورت تصادفی یکی از پایه‌ها را برای اندازه‌گیری انتخاب می‌کند. پس از تبادل کوانتومی، دو طرف از طریق یک کانال کلاسیک احراز هویت‌شده، پایه‌های مورد استفاده را مقایسه کرده و تنها نمونه‌هایی را که در آن‌ها پایه‌ها منطبق بوده‌اند حفظ می‌کنند. این فرایند معمولاً \textit{غربال‌سازی پایه‌ها} یا \lr{basis sifting} نامیده می‌شود. سپس با برآورد نرخ خطا و اجرای مراحل پساپردازش کلاسیک، کلید نهایی استخراج می‌شود.

اهمیت \lr{BB84} از آنجا ناشی می‌شود که اندازه‌گیری حالت کوانتومی در پایه‌ای ناسازگار می‌تواند اختلالی آماری در داده‌ها ایجاد کند. بنابراین، حضور شنودگر نه صرفاً از طریق حدس یا شواهد غیرمستقیم، بلکه از طریق تحلیل آماری نرخ خطا و سایر کمیت‌های مشاهده‌شده ارزیابی می‌شود. در تحلیل امنیت، نرخ خطای مشاهده‌شده تنها یک شاخص عملکردی نیست، بلکه برای کران‌بندی اطلاعات بالقوهٔ مهاجم و، در بسیاری از صورت‌بندی‌ها، برای تخمین کمیت‌هایی مانند خطای فازی نیز به‌کار می‌رود.

با وجود این مبنای نظری، فاصلهٔ میان مدل ایده‌آل و پیاده‌سازی واقعی در عمل تعیین‌کننده است. بسیاری از تحلیل‌های ساده‌شده بر فرض‌هایی مانند وجود منبع تک‌فوتونی کامل، آشکارسازهای بی‌نقص، نبود نویز محیطی و تعداد نامحدود نمونه‌ها تکیه دارند؛ در حالی که هیچ‌یک از این فروض در سامانه‌های واقعی به‌طور کامل برقرار نیست. از همین رو، امنیت و عملکرد یک سامانهٔ \lr{QKD} عملی تنها با بررسی پروتکل ایده‌آل قابل ارزیابی نیست و باید اثر اجزای فیزیکی، آماری و نرم‌افزاری نیز در تحلیل وارد شود.

یکی از محدودیت‌های مهم در پیاده‌سازی عملی \lr{BB84} آن است که ساخت منابع تک‌فوتونی ایده‌آل از نظر فنی دشوار و پرهزینه است. در بسیاری از پیاده‌سازی‌ها، از پالس‌های همدوس تضعیف‌شده یا \lr{Weak Coherent Pulses (WCP)} همراه با تصادفی‌سازی فاز استفاده می‌شود. در این حالت، پالس‌ها را می‌توان به‌صورت مخلوطی آماری از مؤلفه‌های \(n\)-فوتونی با توزیع پواسونی مدل کرد. بنابراین، احتمال غیرصفری برای ارسال پالس‌های چندفوتونی وجود دارد.

وجود پالس‌های چندفوتونی زمینهٔ حملاتی را فراهم می‌کند که در مدل منبع تک‌فوتونی مطرح نیستند. مشهورترین نمونهٔ این تهدید، حملهٔ \lr{Photon-Number Splitting (PNS)} است. در مدل ایده‌آل این حمله، مهاجم می‌تواند با تشخیص تعداد فوتون‌ها، پالس‌های چندفوتونی را شناسایی کند، بخشی از فوتون‌ها را نگه دارد و باقی پالس را به باب برساند. پس از اعلام پایه‌ها، مهاجم می‌تواند از فوتون ذخیره‌شده برای کسب اطلاعات دربارهٔ بخشی از کلید استفاده کند. اگر این رفتار در تحلیل امنیتی لحاظ نشود، ممکن است مهاجم اطلاعاتی دربارهٔ کلید به‌دست آورد، بی‌آنکه نرخ خطای مشاهده‌شده به‌تنهایی برای آشکارسازی حمله کافی باشد. بنابراین، کنترل اثر پالس‌های چندفوتونی بخش مهمی از تحلیل امنیت عملی سامانه‌های مبتنی بر \lr{WCP} است.

برای مقابله با این ضعف ساختاری، روش حالت‌های فریب یا \lr{Decoy-State Method} معرفی شد که از مهم‌ترین پیشرفت‌ها در عملی‌سازی \lr{QKD} محسوب می‌شود. در این رویکرد، آلیس به‌جای ارسال تمام پالس‌ها با یک شدت ثابت، به‌صورت تصادفی از چند شدت متفاوت استفاده می‌کند؛ معمولاً یک شدت سیگنال، یک یا چند شدت فریب ضعیف، و در بسیاری از موارد یک حالت خلأ یا \lr{vacuum decoy}. هدف این روش آن است که رفتار کانال و آشکارسازی برای مؤلفه‌های فوتونی مختلف به‌صورت آماری تخمین زده شود.

در روش حالت‌های فریب، با فرض تصادفی‌سازی کامل فاز و نبود کانال جانبی وابسته به شدت، حالت‌های \(n\)-فوتونی حاصل از شدت‌های مختلف از دید مهاجم غیرقابل‌تمایز در نظر گرفته می‌شوند. در نتیجه، بازده و نرخ خطای متناظر با هر مؤلفهٔ \(n\)-فوتونی باید میان حالت‌های سیگنال و فریب مشترک باشد. آلیس و باب با مقایسهٔ آمار شدت‌های مختلف، می‌توانند بر کمیت‌هایی مانند بازده تک‌فوتونی \(Y_1\) و نرخ خطای تک‌فوتونی \(e_1\) کران بگذارند. این تخمین‌ها زیربنای استخراج نرخ کلید امن در حضور منابع غیرایده‌آل را تشکیل می‌دهند.

به بیان دیگر، روش حالت‌های فریب امکان می‌دهد سهم پالس‌های تک‌فوتونی، که برای امنیت پروتکل نقش اصلی دارند، از سهم پالس‌های چندفوتونی، که بالقوه آسیب‌پذیرند، به‌صورت آماری تفکیک شود. در غیاب این روش، عملکرد بسیاری از سامانه‌های مبتنی بر \lr{WCP} در فواصل متوسط و بلند یا از نظر امنیتی ناکافی می‌بود یا نرخ کلید آن‌ها به‌شدت افت می‌کرد. با بهره‌گیری از این تکنیک، می‌توان در شرایط واقع‌گرایانه و با در نظر گرفتن محدودیت‌های عملی، به نرخ کلید امن قابل‌استفاده دست یافت.

افزون بر مدل منبع، تحلیل سامانه‌های \lr{QKD} باید میان دو رژیم مهم تمایز قائل شود: رژیم \textit{مجانبی} یا \lr{asymptotic} و رژیم \textit{کلید محدود} یا \lr{finite-key}. در مدل مجانبی فرض می‌شود که تعداد پالس‌های مبادله‌شده آن‌قدر زیاد است که نوسانات آماری ناچیز باشند و پارامترها تقریباً دقیق تخمین زده شوند. اما در عمل، هر آزمایش و هر سامانهٔ عملی با تعداد محدودی از پالس‌ها و زمان اجرایی محدود روبه‌رو است. در این وضعیت، نوسانات آماری و خطاهای نمونه‌برداری می‌توانند اثر مهمی بر برآورد کمیت‌های امنیتی بگذارند.

از این رو، تحلیل‌های کلید محدود برای ارزیابی امنیت عملی سامانه‌هایی با اندازهٔ بلوک محدود ضروری‌اند. تحلیل مجانبی همچنان برای فهم رفتار پایه، مقایسه‌های اولیه و اعتبارسنجی تحلیلی ارزشمند است، اما اثر نوسانات آماری و بودجهٔ امنیتی را به‌طور کامل بازنمایی نمی‌کند. در بسیاری از صورت‌بندی‌های کلید محدود، طول کلید نهایی به کمیت‌هایی مانند کران پایین تعداد رویدادهای تک‌فوتونی، کران بالای خطای متناظر، نشت اطلاعات در تصحیح خطا و جمله‌های وابسته به پارامترهای امنیتی بستگی دارد. به‌صورت مفهومی، بیشینهٔ طول کلید قابل استخراج را می‌توان با رابطه‌ای از شکل زیر کران‌بندی کرد:
\[
\ell \lesssim s_{1}^{L}\bigl[1-h_2(e_{1}^{U})\bigr]
-\lambda_{\mathrm{EC}}
-\Delta(\epsilon),
\]
در این رابطه، \(h_2(\cdot)\) تابع آنتروپی دودویی، \(s_{1}^{L}\) کران پایین سهم تک‌فوتونی‌ها، \(e_{1}^{U}\) کران بالای نرخ خطای متناظر، \(\lambda_{\mathrm{EC}}\) نشت اطلاعات در مرحلهٔ تصحیح خطا، و \(\Delta(\epsilon)\) جمله‌ای وابسته به بودجهٔ امنیتی است. شکل دقیق این رابطه به اثبات امنیتی و فروض مدل وابسته است.

بنابراین، مطالعه و توسعهٔ سامانه‌های \lr{QKD} تنها به شناخت اصول بنیادی فیزیک کوانتومی محدود نمی‌شود. چنین مطالعه‌ای مستلزم درک تعامل میان منبع نوری، کانال انتقال، نویزها، آشکارسازها، فرایندهای آماری، اثبات‌های امنیتی و مراحل پساپردازش کلاسیک است. پروژه‌ای که در این گزارش معرفی می‌شود، در همین نقطه قرار دارد: طراحی و پیاده‌سازی یک بستر شبیه‌سازی که بتواند فاصلهٔ میان مدل‌های نظری و رفتار عملی سامانه‌های \lr{QKD} را به‌صورت کنترل‌شده و قابل‌اعتبارسنجی مطالعه کند و امکان تحلیل، اعتبارسنجی، مقایسه و توسعهٔ سناریوهای پیشرفته را فراهم کند.

\section{ضرورت شبیه‌سازی در سامانه‌های \lr{QKD}}


هرچند اصول بنیادی \lr{QKD} در چارچوب نظری به‌خوبی شناخته شده‌اند، گذار از اثبات امنیت روی کاغذ به یک سامانهٔ عملی و قابل‌اتکا، فرایندی پیچیده و چندلایه است. در یک پیاده‌سازی واقعی، هیچ‌یک از اجزای فیزیکی رفتار کاملاً ایده‌آل ندارند و همین انحراف از ایده‌آل، هم بر عملکرد و هم بر امنیت نهایی اثر می‌گذارد. در نتیجه، طراحی، تحلیل و بهینه‌سازی سامانه‌های \lr{QKD} بدون اتکا به ابزارهای شبیه‌سازی دقیق، در بسیاری از موارد ناقص خواهد بود.

در لایهٔ منبع، پالس‌های نوری می‌توانند با نوساناتی در شدت، فاز، زمان‌بندی و حتی توزیع آماری تعداد فوتون همراه باشند. در عمل، منبع ممکن است دچار جیتر شدت یا \lr{intensity jitter} شود، پالس‌ها از نظر شکل زمانی یکنواخت نباشند، یا در برخی شرایط از مدل ایده‌آل فاصله بگیرند. این انحرافات می‌توانند مستقیماً بر تخمین پارامترهای امنیتی و همچنین بر نرخ کلید مؤثر اثر بگذارند. از سوی دیگر، کانال انتقال، به‌ویژه در بستر فیبر نوری، با تضعیف، پراکندگی، وابستگی طول‌موجی، ناپایداری‌های محیطی و نویز پس‌زمینه همراه است. در فواصل بلند، کاهش بازده انتقال می‌تواند بر نرخ آشکارسازی و در نتیجه بر امکان‌پذیری تولید کلید امن اثر قابل‌ملاحظه‌ای بگذارد.

در سمت گیرنده نیز وضعیت به همین اندازه پیچیده است. آشکارسازهای تک‌فوتونی، به‌خصوص آشکارسازهای مبتنی بر \lr{SPAD}، دارای رفتارهایی هستند که در تحلیل ایده‌آل معمولاً نادیده گرفته می‌شوند اما در عمل نقش تعیین‌کننده دارند. از جملهٔ این رفتارها می‌توان به شمارش تاریک یا \lr{dark count}، زمان مرده یا \lr{dead time}، پس‌تپش یا \lr{afterpulsing}، محدودیت‌های گیتینگ زمانی، اشباع‌شدگی در نرخ‌های بالا، و سیاست‌های متفاوت در برخورد با رخدادهای هم‌زمان آشکارسازها یا \lr{double clicks} اشاره کرد.

هر یک از این عوامل می‌تواند هم بر نرخ خطای کوانتومی یا \lr{QBER} و هم بر نرخ کلید نهایی اثر بگذارد. حتی در برخی موارد، عدم مدل‌سازی دقیق همین جزئیات موجب می‌شود نتیجهٔ شبیه‌سازی به‌ظاهر امیدوارکننده باشد، در حالی که سامانهٔ واقعی آن عملکرد را محقق نکند. بنابراین، شبیه‌سازی زمانی ارزشمند است که بتواند اثر این جزئیات را با دقت مناسب و در چارچوب فروض روشن وارد تحلیل کند.

اهمیت موضوع زمانی بیشتر روشن می‌شود که بدانیم بسیاری از تصمیمات طراحی در \lr{QKD} به‌شدت پارامترمحور هستند. انتخاب شدت‌های سیگنال و فریب، احتمال استفاده از هر شدت، نرخ تکرار پالس، بازهٔ زمانی گیت، مدل تلفات کانال، بازده آشکارسازی، و سیاست پردازش داده در لایهٔ کلاسیک، همگی بر خروجی نهایی اثر می‌گذارند. از آنجا که این پارامترها معمولاً به‌صورت کوپل‌شده عمل می‌کنند، فهم اثر مستقل هر یک بدون وجود یک محیط شبیه‌سازی کنترل‌شده دشوار است. برای مثال، افزایش شدت سیگنال ممکن است نرخ آشکارسازی را بالا ببرد، اما هم‌زمان احتمال رخداد پالس‌های چندفوتونی را نیز افزایش دهد. از سوی دیگر، کاهش بیش از حد شدت، اگرچه آسیب‌پذیری در برابر حملات مرتبط با چندفوتونی‌ها را کاهش می‌دهد، ممکن است به افت شدید نرخ آشکارسازی و غلبهٔ نویز تاریک منجر شود. چنین مصالحه‌هایی تنها از طریق تحلیل کمی و تکرارپذیر قابل ارزیابی‌اند.

از منظر پژوهشی و مهندسی، آزمایش‌های سخت‌افزاری \lr{QKD} معمولاً پرهزینه، زمان‌بر، حساس به شرایط محیطی، و محدود از نظر قابلیت پیمایش فضای پارامتر هستند. ساخت یا پیکربندی مجدد یک آزمایش برای هر مجموعهٔ جدید از پارامترها، نه‌تنها هزینهٔ زمانی و مالی دارد، بلکه در بسیاری از موارد امکان اجرای گستردهٔ سناریوهای مقایسه‌ای را محدود می‌کند. در مقابل، یک شبیه‌ساز می‌تواند تعداد زیادی سناریو را در یک محیط کنترل‌شده بررسی کند، نقاط بحرانی عملکرد را بیابد، فرضیه‌ها را پیش از پیاده‌سازی سخت‌افزاری آزمایش کند، و مسیر طراحی آزمایش واقعی را هدایت نماید. بنابراین، در بسیاری از مطالعات عملی \lr{QKD}، شبیه‌سازی نه‌تنها ابزاری کمکی، بلکه بخشی از فرایند طراحی، اعتبارسنجی و بهینه‌سازی سامانه محسوب می‌شود.

یک شبیه‌ساز مفید برای \lr{QKD} نباید تنها به محاسبهٔ روابط بسته و فرمول‌های سادهٔ مجانبی محدود شود. بسیاری از تحلیل‌های نظری در ساده‌ترین حالت، پارامترهایی مانند بهرهٔ آشکارسازی، نرخ خطا یا نرخ کلید را بر اساس روابط تحلیلی در شرایط ایده‌آل یا نیمه‌واقعی برآورد می‌کنند. این روابط برای فهم شهودی و اعتبارسنجی اولیه ارزشمندند، اما برای مدل‌سازی رفتار واقعی سامانه، به‌خصوص در مقیاس کلید محدود، کافی نیستند. در عمل، رویدادها ماهیتاً گسسته و آماری‌اند: فوتون یا به آشکارساز می‌رسد یا نمی‌رسد، کلیک رخ می‌دهد یا رخ نمی‌دهد، پالس ممکن است به‌دلیل زمان مرده از دست برود، و نوسانات تصادفی می‌توانند بر تخمین پارامترها اثر بگذارند.

بنابراین، شبیه‌سازی رویدادمحور و مبتنی بر مونت‌کارلو ابزاری مهم برای بازنمایی رخدادهای گسسته، وابستگی‌های زمانی و نوسانات آماری است. با این حال، نتایج چنین شبیه‌سازی‌ای باید در سناریوهای مناسب با روابط تحلیلی یا بنچمارک‌های مستقل مقایسه شوند تا هزینهٔ محاسباتی بیشتر آن به افزایش قابل‌سنجش در دقت مدل منجر شود.

علاوه بر این، در یک سامانهٔ پژوهشی، صرف شبیه‌سازی لایهٔ کوانتومی کافی نیست. پس از مرحلهٔ آشکارسازی، مجموعه‌ای از پردازش‌های کلاسیک شامل غربال‌سازی پایه‌ها، نمونه‌برداری برای برآورد خطا، تخمین پارامترهای حالت فریب، مدل‌سازی هزینهٔ تصحیح خطا، تخصیص بودجه‌های امنیتی، و محاسبهٔ طول یا نرخ کلید امن انجام می‌شود. اگر این مراحل با دقت آماری کافی مدل نشوند، نتیجهٔ شبیه‌سازی از منظر امنیتی قابل‌استناد نخواهد بود. به بیان دیگر، شبیه‌ساز مطلوب باید میان لایهٔ فیزیکی سامانه و لایهٔ نظریهٔ اطلاعات ارتباط برقرار کند.

در نهایت، شبیه‌سازی دقیق امکان می‌دهد مدل‌های فیزیکی مختلف، اثبات‌های امنیتی متفاوت، یا معماری‌های اجرایی متنوع در یک بستر واحد با یکدیگر مقایسه شوند. چنین قابلیتی برای پژوهش‌هایی که هدف آن‌ها توسعهٔ مدل جدید، اعتبارسنجی یک تقریب فیزیکی، سنجش اثر یک پارامتر، یا مقایسهٔ چند روش تحلیل امنیت است، اهمیت زیادی دارد. از همین رو، توسعهٔ یک بستر شبیه‌سازی توسعه‌پذیر و اعتبارسنجی‌پذیر، نه‌تنها برای پاسخ‌گویی به نیازهای این پروژه، بلکه به‌عنوان زیرساختی برای مطالعات بعدی نیز ارزشمند است.

\section{بیان مسئله و هدف پروژه}

مسئلهٔ اصلی این پروژه، طراحی و پیاده‌سازی یک فریم‌ورک شبیه‌سازی رویدادمحور برای تحلیل سامانه‌های \lr{DV-QKD} است که بتواند مدل‌سازی فیزیکی، رفتار آماری، محدودیت‌های سخت‌افزاری، تحلیل‌های امنیتی و بخشی از نیازهای معماری شبکه‌ای را در یک بستر نرم‌افزاری واحد ترکیب کند. تمرکز نسخهٔ فعلی فریم‌ورک بر خانوادهٔ سامانه‌های مبتنی بر \lr{BB84}، منابع مبتنی بر پالس‌های همدوس تضعیف‌شده، روش حالت‌های فریب و تحلیل‌های مجانبی و کلید محدود است. با این حال، معماری آن به‌گونه‌ای طراحی شده است که افزودن مدل‌ها و سناریوهای جدید در آینده با کمترین وابستگی غیرضروری به بخش‌های موجود انجام شود.

در این گزارش، منظور از فریم‌ورک، مجموعه‌ای از ماژول‌های نرم‌افزاری، رابط‌های برنامه‌نویسی، ابزارهای اجرا، و اسکریپت‌های اعتبارسنجی است که به‌صورت یکپارچه برای تعریف، اجرا و تحلیل سناریوهای \lr{QKD} به‌کار می‌روند. هدف این فریم‌ورک تنها محاسبهٔ چند خروجی نهایی نیست، بلکه بازنمایی مرحله‌به‌مرحلهٔ زنجیرهٔ تولید پالس، عبور از کانال، اثر نویزها، آشکارسازی، تجمیع آمار، تخمین پارامترها و محاسبهٔ کمیت‌های امنیتی است.

هدف‌های اصلی پروژه را می‌توان به‌صورت زیر خلاصه کرد:

\begin{enumerate}
	\item فراهم‌کردن محیطی برای مدل‌سازی پدیده‌های فیزیکی مؤثر بر عملکرد سامانه‌های \lr{QKD}، از جمله منبع، کانال، نویز و آشکارساز؛
	
	\item ایجاد امکان مقایسهٔ چند خانواده از تحلیل‌های امنیتی تحت شرایط فیزیکی و آماری مشترک؛
	
	\item پشتیبانی از تحلیل‌های کلید محدود، در کنار تحلیل‌های مجانبی، برای بررسی اثر نوسانات آماری و بودجه‌های امنیتی؛
	
	\item فراهم‌کردن زیرساخت لازم برای اجرای شبیه‌سازی‌های رویدادمحور در مقیاس بزرگ و با قابلیت پردازش دسته‌ای یا موازی؛
	
	\item فراهم‌سازی معماری نرم‌افزاری ماژولار برای افزودن مدل‌های جدید منبع، کانال، آشکارساز، حمله یا اثبات امنیت؛
	
	\item ایجاد قابلیت اولیه برای مطالعهٔ سناریوهای چندگرهی و شبکه‌ای مبتنی بر لینک‌های \lr{QKD}.
\end{enumerate}

در نسخهٔ فعلی، تصحیح خطا و فشرده‌سازی محرمانگی در سطح اثر آن‌ها بر طول یا نرخ کلید نهایی مدل می‌شوند. به بیان دیگر، تمرکز اصلی فریم‌ورک بر محاسبهٔ کمیت‌های امنیتی و هزینه‌های اطلاعاتی متناظر با مراحل پساپردازش است، نه الزاماً اجرای کامل یک پروتکل عملی تصحیح خطا یا اعمال فیزیکی یک تابع هش جهانی برای تولید بیت‌های نهایی کلید. جزئیات دقیق این مدل‌سازی در فصل‌های مربوط به پساپردازش و تحلیل امنیت ارائه می‌شود.

همچنین، اگرچه تمرکز اصلی فریم‌ورک بر تحلیل لینک‌های \lr{DV-QKD} است، معماری آن به‌گونه‌ای توسعه یافته که بتوان چندین لینک را در قالب سناریوهای شبکه‌ای مبتنی بر رلهٔ مورد اعتماد ترکیب کرد. در چنین معماری‌ای، امنیت انتهابه‌انتها به اعتماد به گره‌های میانی وابسته است؛ زیرا این گره‌ها در فرایند انتقال یا بازتوزیع کلید به اطلاعات کلیدی دسترسی دارند.

\section{نمای کلی فریم‌ورک شبیه‌سازی}

پروژهٔ حاضر یک فریم‌ورک رویدادمحور برای شبیه‌سازی سامانه‌های \lr{Discrete-Variable Quantum Key Distribution (DV-QKD)} ارائه می‌کند. این فریم‌ورک با هدف ترکیب مدل‌سازی فیزیکی، تحلیل آماری، محاسبهٔ کمیت‌های امنیتی و معماری نرم‌افزاری توسعه‌پذیر طراحی شده است. رویکرد محاسباتی آن بر پایهٔ شبیه‌سازی مونت‌کارلو و پردازش رویدادهاست؛ به این معنا که پالس‌ها، آشکارسازی‌ها و رخدادهای مرتبط با آن‌ها به‌صورت نمونه‌های گسسته تولید و تحلیل می‌شوند.

در این معماری، هر پالس نوری به‌صورت یک رویداد مستقل با شناسه و ویژگی‌های مشخص مدل می‌شود. ویژگی‌های پالس می‌توانند در لایه‌های مختلف سامانه، از لحظهٔ تولید در منبع نوری تا عبور از کانال، مواجهه با تلفات و نویز، رسیدن یا نرسیدن به آشکارساز، ثبت کلیک، و ورود به مراحل پساپردازش کلاسیک، تغییر کنند. این رویکرد امکان می‌دهد نه‌تنها خروجی نهایی، بلکه مسیر ایجاد آن و سهم هر بخش از سامانه نیز تحلیل شود.

زنجیرهٔ شبیه‌سازی را می‌توان به‌صورت مفهومی در چند لایه توصیف کرد:

\begin{enumerate}
	\item \textbf{لایهٔ منبع:} منبع نوری پالس‌ها را بر اساس مدل آماری انتخاب‌شده تولید می‌کند. در سناریوهای مبتنی بر منابع همدوس تضعیف‌شده، تعداد فوتون‌ها معمولاً با توزیع پواسونی و شدت‌های مختلف حالت فریب مدل می‌شود.
	
	\item \textbf{لایهٔ کانال:} کانال انتقال اثر تلفات، طول مسیر، بازده انتقال و در صورت نیاز نویزهای پس‌زمینه را بر پالس‌ها اعمال می‌کند.
	
	\item \textbf{لایهٔ مدل تهدید:} در صورت تعریف سناریوی حمله یا مدل تهدید، اثر آن با فروض مشخص به فرایند شبیه‌سازی افزوده می‌شود.
	
	\item \textbf{لایهٔ آشکارسازی:} زیرسامانهٔ آشکارسازی رفتارهایی مانند احتمال کلیک، شمارش تاریک، زمان مرده، گیتینگ زمانی، رخدادهای هم‌زمان و سیاست رسیدگی به آن‌ها را مدل می‌کند.
	
	\item \textbf{لایهٔ پساپردازش کلاسیک:} داده‌های خام وارد مراحل غربال‌سازی، تجمیع شمارش‌ها، تخمین پارامترهای امنیتی، اعمال مدل تصحیح خطا و محاسبهٔ نرخ یا طول کلید امن می‌شوند.
	
	\item \textbf{لایهٔ شبکه:} در صورت استفاده از سناریوهای چندگرهی، خروجی لینک‌ها در سطح شبکه ترکیب می‌شود و مفاهیمی مانند استخر کلید، مسیرهای جایگزین و رمزنگاری \lr{hop-by-hop} قابل بررسی می‌شوند.
\end{enumerate}

طراحی نرم‌افزاری فریم‌ورک بر چند اصل استوار است. نخست، منطق مربوط به پارامترها، منبع، کانال، آشکارساز، تحلیل امنیت، بهینه‌سازی و شبکه تا حد امکان به‌صورت ماژول‌های مستقل سازمان‌دهی شده است. این تفکیک باعث می‌شود تغییر یا جایگزینی یک مدل، کمترین اثر را بر سایر بخش‌ها داشته باشد. دوم، هویت و ویژگی‌های رویدادها در طول پردازش حفظ می‌شود تا تحلیل خطا، اشکال‌زدایی و اعتبارسنجی ممکن باشد. سوم، رابط‌های نرم‌افزاری به‌گونه‌ای طراحی شده‌اند که افزودن مدل‌های جدید بدون بازنویسی کل سیستم انجام شود. چهارم، خروجی‌ها باید بتوانند در سناریوهای ساده‌تر با روابط تحلیلی شناخته‌شده یا بنچمارک‌های مستقل مقایسه شوند.

با توجه به حجم بالای رویدادها در شبیه‌سازی‌های بزرگ، فریم‌ورک از پردازش دسته‌ای و سازوکارهای توزیع بار محاسباتی نیز پشتیبانی می‌کند. در چنین حالتی، حفظ یکپارچگی شناسهٔ پالس‌ها، همگام‌سازی نتایج دسته‌ای و مدیریت رویدادهایی که اثرشان از یک دستهٔ پردازشی به دستهٔ بعد منتقل می‌شود، اهمیت زیادی دارد. ارزیابی کمی کارایی، شامل زمان اجرا، مصرف حافظه و رفتار مقیاس‌پذیری، در فصل مربوط به نتایج و اعتبارسنجی ارائه می‌شود.

در بخش تحلیل امنیت، فریم‌ورک امکان مقایسهٔ چند خانواده از روش‌های محاسبهٔ نرخ کلید را فراهم می‌کند؛ از جمله تحلیل‌های مجانبی کلاسیک برای \lr{Decoy-State BB84}، تحلیل‌های کلید محدود مبتنی بر کران‌های آماری، و نسخه‌هایی که با هدف ارائهٔ کران‌های عددی تنگ‌تر توسعه یافته‌اند. هدف از قرار دادن این روش‌ها در یک بستر مشترک آن است که اثر تفاوت اثبات امنیتی از اثر تفاوت مدل فیزیکی یا فروض نرم‌افزاری تا حد امکان جدا شود.

بهینه‌سازی پارامترها نیز یکی از قابلیت‌های مهم فریم‌ورک است. انتخاب شدت سیگنال \(\mu\)، شدت فریب \(\nu\)، احتمال استفاده از هر شدت، و برخی تنظیمات کانال یا آشکارساز، اثر چشمگیری بر عملکرد سامانه دارد. از این رو، فریم‌ورک می‌تواند از روش‌های عددی و حلگرهای بهینه‌سازی برای پیمایش فضای پارامتر و شناسایی پیکربندی‌های مناسب‌تر استفاده کند. در مسائل خطی یا مراحل کران‌بندی مرتبط با تخمین پارامترها نیز استفاده از ابزارهای برنامه‌ریزی خطی قابل بررسی است.

در مجموع، این فریم‌ورک را می‌توان یک زیرساخت پژوهشی برای مطالعهٔ سامانه‌های \lr{QKD} دانست که با هدف توسعهٔ تدریجی، مقایسهٔ روش‌ها، تولید نتایج تکرارپذیر و پشتیبانی از مطالعات پیچیده‌تر در حوزهٔ \lr{QKD} و شبکه‌های کوانتومی طراحی شده است.

\section{مشارکت‌ها و قابلیت‌های اصلی فریم‌ورک}

مشارکت‌های اصلی این پروژه را می‌توان در چند محور خلاصه کرد:

\begin{enumerate}
	\item \textbf{شبیه‌سازی رویدادمحور لینک QKD:} توسعهٔ یک زنجیرهٔ شبیه‌سازی که پالس‌ها، کانال، نویز، آشکارسازی و پساپردازش را به‌صورت مرحله‌به‌مرحله مدل می‌کند.
	
	\item \textbf{پشتیبانی از منابع همدوس تضعیف‌شده و حالت‌های فریب:} مدل‌سازی منابع مبتنی بر \lr{WCP} و استفاده از شدت‌های مختلف برای تخمین آماری کمیت‌های تک‌فوتونی.
	
	\item \textbf{در نظر گرفتن محدودیت‌های سخت‌افزاری:} وارد کردن عواملی مانند تلفات کانال، شمارش تاریک، زمان مرده، گیتینگ و رخدادهای هم‌زمان آشکارسازها در تحلیل عملکرد.
	
	\item \textbf{ترکیب تحلیل مجانبی و کلید محدود:} فراهم‌کردن امکان بررسی سامانه در دو رژیم مجانبی و کلید محدود، با توجه به اثر نوسانات آماری و بودجهٔ امنیتی.
	
	\item \textbf{معماری ماژولار و توسعه‌پذیر:} تفکیک لایه‌های منبع، کانال، نویز، آشکارساز، اثبات امنیت، بهینه‌سازی و شبکه برای تسهیل نگهداری و افزودن مدل‌های جدید.
	
	\item \textbf{پشتیبانی از اعتبارسنجی داخلی:} استفاده از آزمون‌ها و سناریوهای کنترلی برای بررسی سازگاری روابط فیزیکی، پایداری خروجی‌ها و اتصال صحیح لایه‌های نرم‌افزاری.
	
	\item \textbf{زیرساخت اولیه برای سناریوهای شبکه‌ای:} فراهم‌کردن امکان مطالعهٔ چند لینک و چند گره در قالب معماری‌های مبتنی بر رلهٔ مورد اعتماد.
\end{enumerate}

این موارد دامنهٔ اصلی پروژه را مشخص می‌کنند و مرز میان هدف پژوهشی، مدل‌سازی فیزیکی و پیاده‌سازی نرم‌افزاری را روشن می‌سازند. جزئیات هر یک از این محورها در فصل‌های بعدی گزارش ارائه خواهد شد.

\section{ساختار نرم‌افزاری و اجرای پروژه}

ساختار نرم‌افزاری پروژه بر پایهٔ تفکیک مسئولیت‌ها طراحی شده است. در سطح مفهومی، ماژول‌های اصلی فریم‌ورک شامل مدیریت پارامترها، تعریف ساختارهای داده، مدل‌سازی منبع، کانال و نویز، مدل‌سازی آشکارساز، تحلیل امنیت، بهینه‌سازی، اجرای سناریوها و آزمون‌های اعتبارسنجی هستند. این تفکیک باعث می‌شود که هر بخش از سامانه بتواند مستقل‌تر توسعه یابد و اثر تغییرات آن بر سایر بخش‌ها کنترل شود.

در ریشهٔ پروژه، فایل‌های اجرایی و اسکریپت‌های هماهنگ‌کننده قرار دارند. این فایل‌ها نقش نقطهٔ ورود برای اجرای سناریوهای اصلی، اجرای پردازش‌های دسته‌ای یا موازی، اعتبارسنجی روابط فیزیکی و اجرای نمونه‌سناریوهای شبکه‌ای را بر عهده دارند. جزئیات نام فایل‌ها، وابستگی‌های دقیق آن‌ها و نحوهٔ اجرای هر سناریو در فصل پیاده‌سازی یا پیوست راهنمای اجرا ارائه می‌شود.

هستهٔ فریم‌ورک در قالب مجموعه‌ای از ماژول‌ها سازمان‌دهی شده است. ماژول مدیریت پارامترها، تنظیمات فیزیکی، امنیتی و اجرایی سناریو را نگهداری می‌کند. ماژول‌های داده‌ای، ساختارهای لازم برای انتقال شمارش‌ها، رویدادها و نتایج میانی میان لایه‌ها را تعریف می‌کنند. ماژول‌های فیزیکی، روابط مربوط به تلفات، تولید فوتون، نویز و آشکارسازی را پیاده‌سازی می‌کنند. لایهٔ اثبات‌های امنیتی نیز از داده‌های خروجی لایهٔ فیزیکی استفاده می‌کند تا کمیت‌هایی مانند نرخ کلید، طول کلید، کران‌های آماری و اثر هزینه‌های پساپردازش را محاسبه کند.

برای کنترل صحت توسعه، مجموعه‌ای از آزمون‌های واحد، آزمون‌های یکپارچه‌سازی و در صورت نیاز آزمون‌های رگرسیون در پروژه در نظر گرفته شده است. آزمون‌های واحد برای سنجش رفتار توابع یا کلاس‌های کوچک و مستقل به‌کار می‌روند؛ برای مثال، بررسی صحت محاسبهٔ تلفات یا رفتار یک ساختار داده در شرایط مرزی. آزمون‌های یکپارچه‌سازی تعامل میان ماژول‌ها را بررسی می‌کنند؛ برای نمونه، سازگاری خروجی آشکارساز با ورودی تحلیل امنیت، یا حفظ هویت رویدادها در پردازش دسته‌ای. آزمون‌های رگرسیون نیز برای جلوگیری از بازگشت خطاهای قدیمی پس از اعمال تغییرات جدید مفیدند.

از نظر اجرایی، پیاده‌سازی اصلی پروژه با زبان پایتون انجام شده است. انتخاب پایتون به دلیل اکوسیستم علمی غنی آن، سرعت بالای توسعه، سهولت اعتبارسنجی پژوهشی و دسترسی به کتابخانه‌های عددی و آماری انجام شده است. کتابخانه‌هایی مانند \texttt{NumPy}، \texttt{SciPy}، \texttt{pandas} و \texttt{pytest} در بخش‌های مختلف پروژه برای محاسبات عددی، بهینه‌سازی، مدیریت داده‌های خروجی و اجرای آزمون‌ها مورد استفاده قرار می‌گیرند. بسته به تنظیمات پروژه و ویژگی‌های فعال‌شده، ممکن است وابستگی‌های کمکی دیگری نیز برای ثبت گزارش، ترسیم نمودار یا مدیریت داده‌ها استفاده شوند.

برای اجرای یک شبیه‌سازی پایه، کاربر معمولاً پارامترهای سناریو را در پیکربندی مربوط تنظیم کرده و سپس نقطهٔ ورود مناسب را اجرا می‌کند. در طی اجرا، سیستم پارامترها را بارگذاری کرده، اجزای لازم را نمونه‌سازی می‌کند، رویدادها را طبق سناریوی مشخص تولید و پردازش می‌نماید، و در نهایت خروجی‌ها را در قالب فایل‌های ساخت‌یافته ذخیره می‌کند. در طراحی خروجی‌ها، پایداری داده‌ها و کاهش خطر خرابی فایل در اجراهای طولانی‌مدت یا دسته‌ای مورد توجه قرار گرفته است.

جزئیات کامل ساختار فایل‌ها، نام اسکریپت‌های اجرایی، وابستگی‌های نرم‌افزاری، نحوهٔ نصب، روش اجرای سناریوهای پایه و پیشرفته، و شیوهٔ بازتولید نتایج در فصل پیاده‌سازی و پیوست راهنمای اجرا ارائه خواهد شد. به این ترتیب، فصل حاضر از ورود به جزئیات پایین‌سطح نرم‌افزاری پرهیز می‌کند و تمرکز خود را بر معرفی مسئله، دامنهٔ پروژه و معماری مفهومی فریم‌ورک نگه می‌دارد.

\section{ساختار گزارش}

ادامهٔ این گزارش به‌گونه‌ای سازمان‌دهی شده است که خواننده ابتدا با مبانی نظری و سپس با مدل‌سازی، پیاده‌سازی، اعتبارسنجی و نتایج پروژه آشنا شود. به‌طور کلی، ساختار فصل‌های بعدی به شرح زیر است:

\begin{itemize}
	\item در فصل مبانی نظری، اصول \lr{QKD}، پروتکل \lr{BB84}، منابع همدوس تضعیف‌شده، روش حالت‌های فریب، تحلیل مجانبی و تحلیل کلید محدود با جزئیات بیشتری بررسی می‌شوند.
	
	\item در فصل مدل فیزیکی و آماری، مدل‌های مورد استفاده برای منبع، کانال، نویز، آشکارساز و رخدادهای رویدادمحور معرفی می‌شوند.
	
	\item در فصل تحلیل امنیت، روش‌های محاسبهٔ نرخ یا طول کلید امن، کران‌بندی پارامترها، اثر نوسانات آماری و نقش بودجهٔ امنیتی توضیح داده می‌شوند.
	
	\item در فصل معماری و پیاده‌سازی نرم‌افزار، ساختار ماژول‌ها، ارتباط میان لایه‌ها، نحوهٔ مدیریت داده‌ها، پردازش دسته‌ای و جزئیات اجرایی فریم‌ورک ارائه می‌شود.
	
	\item در فصل اعتبارسنجی و نتایج، خروجی‌های شبیه‌سازی با روابط تحلیلی یا سناریوهای مرجع مقایسه شده و اثر پارامترهای مهم بر عملکرد سامانه بررسی می‌شود.
	
	\item در صورت ارائهٔ سناریوهای شبکه‌ای، فصل مربوط به شبکه نحوهٔ ترکیب لینک‌های \lr{QKD}، مدیریت استخر کلید، مسیرهای جایگزین و معماری مبتنی بر رلهٔ مورد اعتماد را بررسی می‌کند.
	
	\item در پیوست‌ها، جزئیات نصب، اجرای کد، ساختار فایل‌ها، پیکربندی‌ها و اطلاعات تکمیلی لازم برای بازتولید نتایج آورده می‌شود.
\end{itemize}

در نتیجه، این فصل نقش نقشهٔ راه گزارش را ایفا می‌کند. ابتدا زمینهٔ نظری و انگیزهٔ پژوهش را توضیح می‌دهد، سپس چالش‌های گذار از مدل ایده‌آل به سامانهٔ عملی را روشن می‌سازد، مسئله و دامنهٔ پروژه را مشخص می‌کند، نمایی کلی از فریم‌ورک ارائه می‌دهد، و در نهایت خواننده را برای ورود به فصل‌های تخصصی‌تر آماده می‌کند.
        % این فایل باید \chapter و \label خودش را داشته باشد
	
	\chapter{پیکربندی هسته و پارامترها}
	\label{ch:params}
	% =====================================================================
% فصل جامع مستندسازی فریم‌ورک شبیه‌سازی QKD
% بخش‌های مربوط به params.py, constants_helpers.py,
% constants_definitions.py و constants.py
% =====================================================================

% ---------------------------------------------------------------------
% ===== فایل اول: params.py =====
% ---------------------------------------------------------------------

\section{نمای کلی و هدف ماژول \texttt{params.py}}

ماژول \texttt{params.py} به‌عنوان لایه‌ی مرکزی مدیریت و اعتبارسنجی پیکربندی در فریم‌ورک شبیه‌سازیِ توزیع کلید کوانتومی (\lr{Quantum Key Distribution}) عمل می‌کند. این ماژول یک لایه‌ی انتزاعیِ یکپارچه با محدودیت تغییرپذیری برای نگهداری پارامترهای ورودی یک اجرایِ شبیه‌سازی فراهم می‌سازد؛ از مشخصاتِ منبعِ نوری و مدلِ کانالِ فیبرِ نوری تا ساختارِ آشکارساز، نوع پروتکل، پارامترهای امنیتیِ ترکیب‌پذیر (\lr{composable security parameters}) و مدل‌های فیزیکیِ اختلال‌های خطی و غیرخطی در سیستم‌های چندکاناله. هدفِ بنیادینِ این ماژول آن است که پارامترهای اصلی پیکربندی در یک شیءِ مشخص و اعتبارسنجی‌شده متمرکز شوند و سازگاری آن‌ها با فرض‌های مدل بررسی شود تا از پخش‌شدنِ خطاهای پیکربندی در سراسرِ پایگاهِ کد جلوگیری به‌عمل آید.

در معماریِ کلیِ فریم‌ورک، خروجیِ این ماژول یک نمونه از کلاس \texttt{QKDParams} است که به‌صورت یک قراردادِ داده‌ایِ نهایی (\lr{data contract}) برای موتورِ شبیه‌سازی، ماژول‌های محاسبه‌ی نرخ کلید (\lr{key rate computation}) و زیرسیستم‌های گزارش‌گیری عمل می‌کند. هر تغییری در هر یک از اجزای فریم‌ورک — اعم از \lr{OpticalSource}، \lr{FiberChannel}، \lr{ThresholdDetector} و \lr{Protocol} — باید از طریق \texttt{QKDParams} یا متدِ ایمنِ \texttt{with\_channel\_rebuilt()} عبور کند تا سازگاریِ درونیِ پارامترها حفظ شود. این الگو که گاهی «نقطه‌ی واحدِ حقیقت» (\lr{Single Source of Truth}) نامیده می‌شود، به کاهش ناهماهنگی بین پارامترهای مرتبط کمک می‌کند و شناسایی زودهنگام ناسازگاری‌ها را ساده‌تر می‌سازد.

این ماژول افزون بر تعریفِ ساختارِ داده، چند مسئولیت تکمیلی را نیز بر عهده دارد. نخست، اعتبارسنجیِ جامع (\lr{validation}) که در متد \texttt{\_validate()} متمرکز شده و بازه‌ی مجازِ بیش از پنجاه پارامتر را بررسی می‌کند. دوم، سریال‌سازی و بازسازیِ بازگشتی (\lr{serialization / deserialization}) از طریق \texttt{to\_serializable\_dict()} و \texttt{from\_dict()} که امکانِ ذخیره‌سازی، بارگذاری و بازتولیدِ نتایج را فراهم می‌سازد. سوم، مدیریتِ نسخه‌سازیِ شِما (\lr{schema versioning}) با متغیر \texttt{CURRENT\_SCHEMA\_VERSION = "1.9"} که سازگاریِ نسلیِ پیکربندی‌ها را تضمین می‌کند. چهارم، ارائه‌ی توابعِ کارخانِ از پیش‌تنظیم‌شده (\lr{factory presets}) نظیر \texttt{load\_lim2014\_dedicated\_params()} و \texttt{load\_lim2014\_dwdm\_params()} که پیکربندیِ مقالاتِ مرجع را با دقت بالا بازتولید می‌کنند و به تکرارپذیریِ نتایجِ علمی کمک می‌کنند.

به‌عنوان یک ماژولِ کلیدی، \texttt{params.py} با بیش از ۱۷۰۰ خطِ کد و طراحیِ عمیقاً تایپ‌محور (\lr{type-driven design})، مرزِ واسطِ میانِ کاربرِ نهایی و موتورِ فیزیکیِ فریم‌ورک را تشکیل می‌دهد. هر خطای پیکربندی که از این سد عبور کند، می‌تواند در محاسباتِ پایین‌دستی بر نتایج عددی و تحلیل امنیتی اثر بگذارد؛ از این رو، طراحیِ این ماژول بر پایه‌ی اصلِ «سریع‌شکستن» (\lr{fail-fast}) استوار است و هر پیکربندیِ ناسازگار پیش از ورود به حلقه‌ی شبیه‌سازی رد می‌شود.

\section{مبانی تئوری و فیزیکی در \texttt{params.py}}

تئوریِ پایه‌ی این ماژول بر سه ستونِ اصلی استوار است: تئوریِ امنیتیِ ترکیب‌پذیرِ کوانتومی، فیزیکِ اتلاف در کانال‌های نوری، و مدل‌های اختلالِ غیرخطی در سیستم‌های \lr{WDM}. در ادامه‌ی این بخش هر یک از این ستون‌ها را با جزئیات بسط می‌دهیم.

\subsection{امنیتِ ترکیب‌پذیر و پارامترهای اِپسیلون}

در چارچوبِ امنیتیِ ترکیب‌پذیر که توسط \lr{Renner} و \lr{König} توسعه یافته، امنیتِ یک پروتکلِ \lr{QKD} با یک پارامترِ کلی $\varepsilon$ مشخص می‌شود که احتمالِ انحرافِ خروجیِ پروتکل از توزیعِ ایده‌آل را محدود می‌کند. این پارامتر کل به چند مؤلفه‌ی مستقل تجزیه می‌شود:
\begin{equation}
	\varepsilon_{\text{tot}} = \varepsilon_{\text{sec}} + \varepsilon_{\text{cor}} + \varepsilon_{\text{PE}} + \varepsilon_{\text{smooth}} + \varepsilon_{\text{PA}} + \varepsilon_{\text{SIF}}
\end{equation}
که در آن $\varepsilon_{\text{sec}}$ امنیتِ کلی، $\varepsilon_{\text{cor}}$ صحت، $\varepsilon_{\text{PE}}$ خطای تخمینِ پارامتر، $\varepsilon_{\text{smooth}}$ هموارسازی، $\varepsilon_{\text{PA}}$ تقویتِ حریمِ خصوصی و $\varepsilon_{\text{SIF}}$ فیلترِ اطلاعاتِ جانبی را نشان می‌دهد. ماژول \texttt{params.py} اعتبارِ هر یک از این شش پارامتر را به‌صورت جداگانه بررسی می‌کند و سپس شرطِ ترکیبی را به‌صورت:
\begin{equation}
	\sum_i \varepsilon_i < 1.0
\end{equation}
اعمال می‌نماید. این شرط، شرطِ لازم برای آن است که مفهومِ «امنیتِ محاسباتیِ کوانتومی» معنای فیزیکی داشته باشد؛ چرا که $\varepsilon \geq 1$ به معنایِ نبودِ هرگونه تضمینِ امنیتی است.

\subsection{\texorpdfstring{کاراییِ تصحیحِ خطا و ضریب $f_{\text{EC}}$}{کارایی تصحیح خطا و ضریب f-EC}}

در مرحله‌ی تصحیحِ خطا، یک پروتکلِ کلاسیک مانند \lr{Cascade} یا \lr{LDPC} برای هم‌سازسازیِ کلیدهای آلیس و باب به‌کار می‌رود. کاراییِ این مرحله با ضریبِ $f_{\text{EC}}$ مشخص می‌شود که نسبتِ اطلاعاتِ نشت‌کرده به حدِ نظریِ \lr{Shannon} را نشان می‌دهد:
\begin{equation}
	\text{leak}_{\text{EC}} = f_{\text{EC}} \cdot n \cdot h_2(Q_{\mu})
\end{equation}
که در آن $h_2(x) = -x \log_2 x - (1-x) \log_2(1-x)$ آنتروپیِ باینری و $Q_{\mu}$ نرخِ خطای کوانتومی (\lr{QBER}) است. در فریم‌ورک، $f_{\text{EC}}$ باید در بازه‌ی $[1.0, 2.0]$ باشد و مقدارِ پیش‌فرضِ $1.16$ به‌عنوان نمونه‌ای از کاراییِ متداول در پیاده‌سازی‌های عملی در نظر گرفته می‌شود.

\subsection{اتلافِ کانال و مدل‌های چند‌جزئی}

اتلافِ کلِ کانال در فریم‌ورک از فرمولِ تجمعیِ زیر پیروی می‌کند:
\begin{equation}
	\alpha_{\text{tot}} = \alpha_{\text{fiber}} \cdot L + \alpha_{\text{filter}} + \alpha_{\text{mod}} + \alpha_{\text{circ}}
\end{equation}
که در آن $\alpha_{\text{fiber}}$ ضریبِ اتلافِ فیبر بر حسب \lr{dB/km}، $L$ طولِ فیبر، $\alpha_{\text{filter}}$ اتلافِ فیلتر، $\alpha_{\text{mod}}$ اتلافِ ماژولاتور و $\alpha_{\text{circ}}$ اتلافِ سیرکولاتور است. اتلافِ فیلتر خود به سه مدل تقسیم می‌شود: \lr{ideal} (بدون اتلاف)، \lr{awg} (اتلافِ ثابتِ شبکه‌ی آرایه‌ی موج‌بر، معمولاً ۳ \lr{dB}) و \lr{fbg\_serial} (اتلافِ تجمعیِ \lr{FBG} که با تعدادِ کانال‌ها رشد می‌کند):
\begin{equation}
	\alpha_{\text{filter}}^{\text{fbg}} = N_{\text{ch}} \cdot \alpha_{\text{FBG, per-ch}}
\end{equation}

\subsection{اختلال‌های غیرخطیِ \lr{WDM}: پراکندگیِ رامان}

در حالتِ هم‌زیستیِ کانال‌های کوانتومی و کلاسیکِ \lr{WDM}، پدیده‌ی پراکندگیِ رامانِ خودبه‌خودی (\lr{Spontaneous Raman Scattering}) منبعِ اصلیِ نویزِ پس‌زمینه در کانالِ کوانتومی است. مدلِ فریم‌ورک این نویز را به‌صورت زیر محاسبه می‌کند:
\begin{equation}
	P_{\text{Raman}} = P_{\text{launch}} \cdot (N_{\text{WDM}} - 1) \cdot \gamma_{\text{Raman}} \cdot L_{\text{eff}}
\end{equation}
که در آن طولِ مؤثرِ غیرخطی $L_{\text{eff}}$ برابر است با:
\begin{equation}
	L_{\text{eff}} = \frac{1 - e^{-\alpha L}}{\alpha}, \qquad \alpha = 0.2 \times 0.23026 \, \text{Np/km}
\end{equation}
در تابعِ کارخانه‌ی \texttt{load\_lim2014\_dwdm\_params()}، توانِ لانچ به‌گونه‌ای محاسبه می‌شود که توانِ دریافتی در آشکارساز برابر با $-34$ dBm (حداقل حساسیتِ گیرنده‌های تجاریِ \lr{DWDM SFP}) باشد:
\begin{equation}
	P_{\text{launch}}^{\text{dBm}} = -34 + \alpha_{\text{fiber}} \cdot L + \alpha_{\text{AWG}}
\end{equation}

\subsection{فیزیکِ ماژولاسیونِ زیرحامل (\lr{SCM})}

در پروتکل‌های \lr{SCM-QKD}، سیگنالِ کوانتومی روی یک زیرحاملِ فرکانسی ماژوله می‌شود. پارامتر $\beta$ (\lr{modulation\_index}) عمقِ ماژولاسیون را کنترل می‌کند و نویزِ محصولاتِ تقاطعیِ مرتبه‌ی دوم (\lr{CSO}) و مرتبه‌ی سوم (\lr{CTB}) را تحت تأثیر قرار می‌دهد. مدلِ \lr{cso\_model} بین سه حالت \lr{approximate}، \lr{discrete\_exact} و \lr{combinatorial\_uniform} انتخاب می‌شود که هر یک مرتبه‌ی دقتِ محاسبات را تعیین می‌کند.

\subsection{پراکندگی‌جبرانیِ فعال}

در کانال‌های طولانی، پراکندگیِ کروماتیک پالس‌ها را پهن می‌کند و نرخِ کلیکِ تک‌فوتونی را کاهش می‌دهد. پارامتر $D$ (\lr{dispersion\_parameter\_ps\_nm\_km}) این پراکندگی را با واحدِ \lr{ps/(nm·km)} مشخص می‌کند و در فیبرِ استانداردِ \lr{SMF-28} حدود $17$ ps/(nm·km) در طول‌موجِ $1550$ nm است. جبرانِ فعال با پرچمِ \texttt{active\_dispersion\_compensation} فعال می‌شود.

\section{تحلیل معماری و ساختار داده‌ها در \texttt{params.py}}

\subsection{کلاس \texttt{QKDParams}: یک دیتاکلاسِ ایستا}

مرکزِ معماریِ این ماژول، کلاسِ \texttt{QKDParams} است که با دکوراتورهای زیر تعریف می‌شود:

\begin{lstlisting}[language=Python]
	@dataclass(frozen=True, slots=True)
	class QKDParams:
	protocol: ProtocolFromModule
	source: OpticalSource
	channel: FiberChannel
	detector: ThresholdDetector
	num_bits: int
	photon_number_cap: int
	...
\end{lstlisting}



پارامتر \texttt{frozen=True} از انتسابِ مستقیم به فیلدهای نمونه پس از ساخت جلوگیری می‌کند و هر اصلاح باید از طریقِ \texttt{dataclasses.replace()} یا \texttt{with\_channel\_rebuilt()} انجام شود. این محدودیت، ضمنِ تحمیلِ انضباطِ برنامه‌نویسی، اشتراک‌گذاریِ خواندنیِ نمونه‌ها را ساده‌تر می‌کند؛ اما به‌تنهایی تضمین‌کننده‌ی ایمنی کامل در محیط‌های همزمان نیست. دکوراتورِ \texttt{slots=True} که در پایتون ۳.۱۰ معرفی شده، می‌تواند با حذفِ \texttt{\_\_dict\_\_} سربارِ حافظه را کاهش دهد؛ میزانِ اثر آن بر سرعت و مصرف حافظه به نسخه‌ی Python و الگوی استفاده وابسته است.

\subsection{دسته‌بندیِ فیلدها}

فیلدهای \texttt{QKDParams} به سه گروهِ اصلی تقسیم می‌شوند:

\textbf{۱. مؤلفه‌های فیزیکی (بدون مقدارِ پیش‌فرض):} \texttt{protocol}، \texttt{source}، \texttt{channel}، \texttt{detector}. این چهار مؤلفه، چهار رکنِ یک سیستمِ \lr{QKD} را تشکیل می‌دهند و باید توسط کاربر به‌صراحت ارائه شوند.

\textbf{۲. پارامترهای امنیتی و محاسباتی:} \texttt{num\_bits}، \texttt{photon\_number\_cap}، \texttt{batch\_size}، \texttt{num\_workers}، \texttt{eps\_sec}، \texttt{eps\_cor}، \texttt{eps\_pe}، \texttt{eps\_smooth}، \texttt{security\_proof}، \texttt{ci\_method}. این فیلدها بسته به نقششان در مدل فیزیکی و تحلیل امنیتی، می‌توانند بر خروجی نهایی و نرخ کلید اثر بگذارند.
\textbf{۳. پارامترهای فیزیکیِ پیشرفته (با مقادیرِ پیش‌فرض):} این گروه شامل بیش از سی فیلد است و پارامترهای \lr{SCM}، \lr{WDM}، \lr{Chapman 2018}، \lr{Lim 2014} و \lr{Ma 2005} را پوشش می‌دهد. استفاده از مقادیرِ پیش‌فرض به کاربرِ مبتدی اجازه می‌دهد بدون درگیری با جزئیاتِ فیزیکی، یک شبیه‌سازیِ پایه را اجرا کند.

\subsection{پروتکل‌های تایپی و \lr{Protocol}}

فریم‌ورک از دو پروتکلِ تایپی برای تعریفِ رابط‌ها استفاده می‌کند:

\begin{lstlisting}[language=Python]
	class SerializableComponent(TypingProtocol):
	def to_config_dict(self) -> Dict[str, Any]: ...
	
	class QKDProtocol(SerializableComponent, TypingProtocol):
	protocol_name: str
\end{lstlisting}

این پروتکل‌ها قراردادهایی هستند که هر مؤلفه‌ی فیزیکی باید آن‌ها را رعایت کند: متدِ \texttt{to\_config\_dict()} برای سریال‌سازی و صفتِ \texttt{protocol\_name} برای شناسایی. این طراحی امکانِ بررسیِ ایستا (\lr{static checking}) با \lr{mypy} را فراهم می‌سازد و خطاهای نوعی را پیش از اجرا کشف می‌کند.

\subsection{شمارش‌ها (\lr{Enums})}

ماژول از چندین شمارش برای محدود کردن مقادیرِ مجاز استفاده می‌کند: \texttt{SecurityProof} با مقادیر \lr{LIM\_2014}، \lr{MDI\_QKD}، \lr{PAPER\_2009}؛ \texttt{ConfidenceBoundMethod} با روش‌های \lr{CLOPPER\_PEARSON} و \lr{NORMAL\_APPROXIMATION}؛ \texttt{DecoderArchitecture} با حالت‌های \lr{LOCAL}، \lr{ENTANGLING} و... . این شمارش‌ها، انضباطِ دامنه‌ای را اعمال می‌کنند و از خطاهای املایی جلوگیری می‌کنند.

\subsection{سیستمِ نسخه‌سازیِ شِما}

ثابت \texttt{CURRENT\_SCHEMA\_VERSION = "1.9"} و تابعِ کمکی \texttt{\_parse\_schema\_version()} نسخه‌ی پیکربندی را به‌صورت یک چندتاییِ قابل‌مقایسه نمایش می‌دهند:

\begin{lstlisting}[language=Python]
	def _parse_schema_version(v: Any) -> tuple:
	parts = [int(p) for p in v.split(".")]
	while len(parts) > 1 and parts[-1] == 0:
	parts.pop()
	return tuple(parts) if parts else (0,)
\end{lstlisting}

این رویکرد تضمین می‌کند که نسخه‌های \lr{"1.9"}، \lr{"1.9.0"} و \lr{"1.9.0.0"} هم‌ارز در نظر گرفته شوند و در عین حال \lr{"1.10"} بزرگ‌تر از \lr{"1.9"} باشد (ترتیبِ عددی، نه لغوی).

\subsection{محدودیت‌های سریال‌سازی}

سه ثابت \texttt{MAX\_SERIALIZABLE\_LIST\_LEN = 1000}، \texttt{MAX\_SERIALIZABLE\_DICT\_LEN = 1000} و \texttt{MAX\_SERIALIZABLE\_RECURSION\_DEPTH = 20} مرزهای امنیتیِ سریال‌سازی را تعیین می‌کنند. این محدودیت‌ها از حملاتِ انکارِ سرویسِ مبتنی بر ساختارهای بازگشتیِ عمیق یا لیست‌های بسیار بزرگ جلوگیری می‌کنند و سازگار با توصیه‌های \lr{OWASP} برای رمزگشاییِ \lr{JSON} هستند.

\section{پیاده‌سازی ریاضی و الگوریتمی در \texttt{params.py}}

\subsection{فرمولِ واحدِ اتلافِ کل}

نقطه‌ی حقیقتِ واحد برای محاسبه‌ی اتلاف، تابعِ مستقلِ زیر است:

\begin{lstlisting}[language=Python]
	def _compute_total_loss_db_from_scalars(
	*, distance_km, fiber_loss_db_km, filter_model,
	filter_awg_loss_db, N_channels,
	filter_fbg_loss_per_channel_db,
	modulator_insertion_loss_db,
	circulator_insertion_loss_db,
	) -> float:
	if filter_model == "awg":
	filter_loss = filter_awg_loss_db
	elif filter_model == "fbg_serial":
	filter_loss = N_channels * filter_fbg_loss_per_channel_db
	elif filter_model == "ideal":
	filter_loss = 0.0
	else:
	raise ParameterValidationError(...)
	return (distance_km * fiber_loss_db_km
	+ filter_loss
	+ modulator_insertion_loss_db
	+ circulator_insertion_loss_db)
\end{lstlisting}

این تابع به‌صورت معناداری \lr{keyword-only} است (نشانه‌ی \lr{\texttt{*}}) تا از اشتباهِ جایگاهیِ آرگومان‌ها جلوگیری شود. این تابع توسط سه مسیرِ مجزا فراخوانی می‌شود: \texttt{\_compute\_total\_loss\_db()} (متدِ کلاس)، \texttt{from\_dict()} و توابعِ کارخانه‌ی \lr{\texttt{load\_lim2014\_*}}. این الگو از یکپارچگیِ فرمول اطمینان می‌دهد و خطرِ انحرافِ تدریجیِ (\lr{drift}) بین مسیرهای مختلف را از بین می‌برد.

\subsection{اعتبارسنجیِ امنیتِ ترکیب‌پذیر}

حلقه‌ی اعتبارسنجیِ پارامترهای اِپسیلون در متدِ \texttt{\_validate()} به‌صورت زیر پیاده‌سازی شده است:

\begin{lstlisting}[language=Python]
	for eps_name, eps_val in (("eps_sec", self.eps_sec),
	("eps_cor", self.eps_cor),
	("eps_pe", self.eps_pe),
	("eps_smooth", self.eps_smooth),
	("eps_pa", self.eps_pa),
	("eps_sif", self.eps_sif)):
	if not isinstance(eps_val, (int, float)) or isinstance(eps_val, bool):
	raise ParameterValidationError(...)
	if not math.isfinite(eps_val):
	raise ParameterValidationError(...)
	if eps_val <= 0.0:
	raise ParameterValidationError(...)
	if eps_val >= 1.0:
	raise ParameterValidationError(...)
	
	total_epsilon = (self.eps_sec + self.eps_cor + self.eps_pe
	+ self.eps_smooth + self.eps_pa + self.eps_sif)
	if not math.isfinite(total_epsilon) or total_epsilon >= 1.0:
	raise ParameterValidationError(...)
\end{lstlisting}

نکته‌ی ظریفِ این پیاده‌سازی، بررسیِ جداگانه‌ی هر مؤلفه پیش از جمع زدن است. این رویکرد از یک خطای نامحسوس جلوگیری می‌کند: اگر یکی از مؤلفه‌ها \lr{NaN} باشد، جمعِ کلی نیز \lr{NaN} خواهد شد و مقایسه‌ی $\text{NaN} \geq 1.0$ همواره \lr{False} برمی‌گرداند و در نتیجه یک پارامترِ نامعتبر به‌صورت خاموش عبور می‌کند. بررسیِ \texttt{math.isfinite()} در سطحِ مؤلفه این حفره را می‌بندد.

\subsection{الگوریتمِ تبدیلِ ایمنِ شمارش‌ها}

تابعِ \texttt{\_coerce\_enum()} یک الگوریتمِ چندمرحله‌ای برای تبدیلِ ورودی‌های متنوع به اعضایِ شمارش ارائه می‌دهد:

\begin{lstlisting}[language=Python]
	def _coerce_enum(enum_cls, value):
	if isinstance(value, enum_cls):
	return value
	if isinstance(value, str):
	if enum_cls not in _ENUM_LOOKUP_CACHE:
	_ENUM_LOOKUP_CACHE[enum_cls] = _build_enum_lookup(enum_cls)
	table = _ENUM_LOOKUP_CACHE[enum_cls]
	s = value.strip()
	hit = table.get(s) or table.get(s.upper())
	if hit is not None:
	return hit
	if isinstance(value, bool):
	raise ParameterValidationError(...)
	try:
	return enum_cls(value)
	except (ValueError, TypeError) as e:
	candidates = [m.name for m in enum_cls]
	if isinstance(value, str):
	matches = difflib.get_close_matches(
	value.strip().upper(),
	[c.upper() for c in candidates], n=1)
	if matches:
	suggestion = f" Did you mean '{matches[0]}'?"
	raise ParameterValidationError(...) from e
\end{lstlisting}

این تابع از یک کشِ ایستا (\lr{\_ENUM\_LOOKUP\_CACHE}) استفاده می‌کند تا جدولِ جستجویِ نادیده‌بگیرِ بزرگی و کوچکیِ حروف تنها یک‌بار برای هر شمارش ساخته شود. علاوه بر این، با استفاده از \texttt{difflib.get\_close\_matches()} پیشنهادِ هوشمندانه‌ای به کاربر ارائه می‌دهد؛ مثلاً اگر کاربر \lr{"lim2014"} را تایپ کند، پیام \lr{Did you mean 'LIM\_2014'?} دریافت می‌کند. این تجربه‌ی کاربریِ غنی، در پیکربندی‌های پیچیده‌ی علمی زمانِ قابل‌توجهی صرفه‌جویی می‌کند.

\subsection{سریال‌سازیِ بازگشتی و حفاظت در برابر نشتِ seed}

تابعِ \texttt{\_to\_serializable()} یک الگوریتمِ بازگشتی برای تبدیلِ ساختارهای داده‌ی پیچیده‌ی پایتون به فرمتِ \lr{JSON-compatible} است:

\begin{lstlisting}[language=Python]
	def _to_serializable(o, _depth=0):
	if _depth > MAX_SERIALIZABLE_RECURSION_DEPTH:
	raise ParameterValidationError(...)
	if (not isinstance(o, type)
	and hasattr(o, "to_config_dict")
	and callable(o.to_config_dict)):
	return _to_serializable(o.to_config_dict(), _depth + 1)
	if isinstance(o, np.generic):
	return o.item()
	if isinstance(o, np.ndarray):
	if o.size > MAX_SERIALIZABLE_LIST_LEN:
	raise ParameterValidationError(...)
	return o.tolist()
	if isinstance(o, Enum):
	return o.value
	if dataclasses.is_dataclass(o) and not isinstance(o, type):
	return {f.name: _to_serializable(getattr(o, f.name), _depth + 1)
		for f in dataclasses.fields(o)}
	if isinstance(o, (set, frozenset)):
	_ordered = sorted(o, key=repr)
	return [_to_serializable(i, _depth + 1) for i in _ordered]
	...
\end{lstlisting}

این تابع چند نکته‌ی مهم دارد: نخست، محدودیتِ عمقِ بازگشت (\lr{20}) از سرریزِ پشته جلوگیری می‌کند. دوم، آرایه‌های \lr{NumPy} به‌صورت خودکار به لیست تبدیل می‌شوند با بررسیِ اندازه. سوم، مجموعه‌ها (\lr{set}) به‌صورت مرتب سریال‌سازی می‌شوند تا هشِ خروجی تکرارپذیر باشد و قابل استفاده در سیستم‌های کنترلِ نسخه باشد.

تابعِ \texttt{\_redact\_seed\_recursive()} لایه‌ی دفاعیِ مضاعفی است که اطمینان می‌دهد هیچ فیلدِ شبیهِ seed (مانند \texttt{master\_seed}، \texttt{seed}، \texttt{random\_seed}) نمی‌تواند به خروجیِ سریال‌سازی راه یابد. این محافظت از نشتِ تصادفیِ seed در گزارش‌های آزمایشی — که ممکن است در مخزنِ عمومیِ \lr{Git} منتشر شوند — جلوگیری می‌کند و امنیتِ عملیاتیِ شبیه‌سازی‌های مبتنی بر اعدادِ تصادفیِ شبه‌تصادفی را حفظ می‌نماید.

\subsection{متد \texttt{from\_dict}: بازسازیِ سازگار}

متدِ کلاسی \texttt{from\_dict()} یک پیاده‌سازیِ طولانی و دقیق از الگوی «مصرف‌کنـد-سپس‌شناسایی‌ناشناخته» (\lr{consume-then-detect-unknowns}) است. این متد با استفاده از \texttt{\_coerce\_type()} که هر کلید را از دیکشنری \texttt{pop} می‌کند، در پایان می‌تواند هر کلیدِ باقی‌مانده را به‌عنوان «ناشناخته» شناسایی کند و خطای دقیقی ارائه دهد:

\begin{lstlisting}[language=Python]
	d_copy = {k: v for k, v in d_copy.items() if not k.startswith("_")}
	if d_copy:
	raise ConfigurationError(
	f"Unknown parameter keys provided: {sorted(d_copy.keys())}")
\end{lstlisting}

این رویکرد در تجربه‌ی کاربریِ پیکربندی‌های علمی بسیار ارزشمند است؛ چرا که اشتباهاتِ املایی نظیر \lr{eps\_secc} به‌جای \lr{eps\_sec} بلافاصله شناسایی می‌شوند و از نتایجِ خاموشِ نادرست جلوگیری می‌کنند.

\subsection{متد \texttt{with\_channel\_rebuilt} و حفاظتِ یکپارچگیِ اتلاف}

هنگامی که کاربر می‌خواهد طولِ فیبر یا ضرایبِ اتلاف را تغییر دهد، نمی‌تواند به‌سادگی از \texttt{dataclasses.replace()} استفاده کند؛ چرا که فیلدِ \texttt{channel.total\_loss\_db} باید با فیلدهای مبدأ همگام بماند. متدِ \texttt{with\_channel\_rebuilt()} این مشکل را با محاسبه‌ی پیشینِ اتلافِ جدید و بازسازیِ کانال حل می‌کند:

\begin{lstlisting}[language=Python]
	def with_channel_rebuilt(self, **changes):
	if "distance_km" in changes:
	new_distance_km = float(changes.pop("distance_km"))
	else:
	new_distance_km = self.channel.distance_km
	_new = {f.name: getattr(self, f.name) for f in dataclasses.fields(self)}
	_new.update(changes)
	new_total_loss = _compute_total_loss_db_from_scalars(
	distance_km=new_distance_km,
	fiber_loss_db_km=_new["fiber_loss_db_km"],
	...
	)
	new_channel = FiberChannel.from_total_loss(new_distance_km, new_total_loss)
	return dataclasses.replace(self, channel=new_channel, **changes)
\end{lstlisting}

علاوه بر این، متدِ \texttt{\_validate()} دارای گاردِ سازگاریِ اتلاف است:

\begin{lstlisting}[language=Python]
	_expected_total_loss = self._compute_total_loss_db()
	if abs(self.channel.total_loss_db - _expected_total_loss) > 1e-9 * max(1.0, abs(_expected_total_loss)):
	raise ParameterValidationError(
	f"channel.total_loss_db ({self.channel.total_loss_db}) does not match "
	f"recomputed loss ({_expected_total_loss}). Use "
	f"QKDParams.with_channel_rebuilt() or from_dict() when changing "
	f"loss-affecting fields.")
\end{lstlisting}

این گارد از ناسازگاریِ پنهان بین فیلدها جلوگیری می‌کند: اگر کاربر از طریق \texttt{replace()} فیلدِ \texttt{filter\_model} را بدون بازسازیِ کانال تغییر دهد، اتلافِ ذخیره‌شده در کانال دیگر با فیلدهای مبدأ تطابق نخواهد داشت و از نادرست شدن تدریجیِ نتایج جلوگیری می‌کند.
\subsection{توابعِ کارخانه‌ی \lr{Lim 2014}}

دو تابعِ \texttt{load\_lim2014\_dedicated\_params()} و \texttt{load\_lim2014\_dwdm\_params()} پیکربندی‌های مقاله‌ی مرجع \lr{Lim et al. (2014)} را بازتولید می‌کنند. در تابعِ دوم، محاسبه‌ی توانِ لانچ بر اساسِ فرمولِ زیر است:

\begin{lstlisting}[language=Python]
	if wdm_channel_power_dbm is None:
	_fiber_loss_db = distance_km * params.fiber_loss_db_km
	wdm_channel_power_dbm = -34.0 + _fiber_loss_db + _awg_filter_loss_db
	if wdm_channel_power_dbm > 30.0:
	raise ParameterValidationError(...)
\end{lstlisting}

این فرمول طوری تنظیم شده است که توانِ دریافتی در گیرنده‌ی کلاسیک نزدیک به $-34$ dBm (حساسیتِ گیرنده‌های تجاری) باقی بماند و در نتیجه‌ی مقایسه‌ی عادلانه‌ای بینِ سناریوهای فاصله‌ی مختلف میسر شود.

\section{ارسال و دریافت با سایر ماژول‌ها در \texttt{params.py}}

\subsection{وابستگی‌های ورودی}

ماژول \texttt{params.py} به‌صورت گسترده‌ای از سایر بخش‌های فریم‌ورک واردات (\lr{import}) انجام می‌دهد. این وابستگی‌ها را می‌توان در پنج دسته‌ی منطقی قرار داد:

\textbf{۱. مؤلفه‌های فیزیکی:} \texttt{OpticalSource} از \texttt{.sources}، \texttt{FiberChannel} از \texttt{.channel}، \texttt{ThresholdDetector} از \texttt{.detectors}. این سه کلاس ساختمان‌های داده‌ایِ اجزای فیزیکی هستند که \texttt{QKDParams} در خود نگه می‌دارد.

\textbf{۲. پروتکل‌ها:} \texttt{Protocol} از \texttt{.protocols} به‌همراه کلاس‌های \texttt{BB84DecoyProtocol}، \texttt{MDIQKDProtocol}، \texttt{B92Protocol} و \texttt{RedundantTransmissionProtocol} که در متدِ \texttt{from\_dict()} به‌صورت داخلی وارد می‌شوند.

\textbf{۳. نوع‌های داده و شمارش‌ها:} از \texttt{.datatypes} شامل \texttt{DoubleClickPolicy}، \texttt{SecurityProof}، \texttt{ConfidenceBoundMethod}، \texttt{PulseTypeConfig}، \texttt{SourceStatisticsType}، \texttt{DecoderArchitecture}، \texttt{SourceErrorModel} و \texttt{PulseEnsembleConfig}.

\textbf{۴. استثناها:} \texttt{ParameterValidationError} و \texttt{ConfigurationError} از \texttt{.exceptions} که در تمام مسیرهای خطا استفاده می‌شوند.

\textbf{۵. ثابت‌ها و ابزارها:} \texttt{LP\_SOLVER\_METHODS} و \texttt{DEFAULT\_POISSON\_TAIL\_THRESHOLD} از \texttt{.constants}، \texttt{parse\_bool} از \texttt{.utils.validation}، \texttt{AfterpulseModel} از \texttt{.detectors}.

\subsection{خروجی‌ها و نقاطِ مصرف}

خروجیِ اصلیِ ماژول، نمونه‌های \texttt{QKDParams} است که توسط چندین مصرف‌کننده‌ی پایین‌دستی استفاده می‌شود:

\begin{itemize}
	\item \textbf{موتورِ شبیه‌سازی:} از فیلدهای \texttt{protocol}، \texttt{source}، \texttt{channel}، \texttt{detector} و \texttt{num\_bits} برای اجرای حلقه‌ی اصلیِ تولید و اندازه‌گیری استفاده می‌کند.
	\item \textbf{محاسبه‌ی نرخِ کلید:} از \texttt{security\_proof}، \texttt{ci\_method}، \texttt{f\_error\_correction}، تمامی پارامترهای \lr{\texttt{eps\_*}} و \texttt{photon\_number\_cap} برای محاسبه‌ی نرخِ کلیدِ نهایی استفاده می‌کند.
	\item \textbf{حلقه‌ی بهینه‌سازی:} از \texttt{auto\_optimize\_lim2014} و \texttt{auto\_optimize\_ma2005} برای تصمیم‌گیری در موردِ اجرایِ بهینه‌سازیِ پارامترهای منبع استفاده می‌کند.
	\item \textbf{سیستمِ گزارش‌گیری:} از \texttt{to\_serializable\_dict()} برای تولیدِ \lr{JSON} قابلِ بازتولید استفاده می‌کند.
\end{itemize}

\subsection{قراردادِ سریال‌سازی}

متد \texttt{to\_serializable\_dict()} یک دیکشنریِ تودرتو با ساختارِ زیر تولید می‌کند:

\begin{lstlisting}[language=Python]
	{
		"schema_version": "1.9",
		"protocol_name": "bb84-decoy",
		"protocol_config": {...},
		"source_config": {...},
		"channel_config": {...},
		"detector_config": {...},
		"num_bits": ..., "eps_sec": ..., ...
	}
\end{lstlisting}

این ساختار با ورودیِ \texttt{from\_dict()} همخوانی دارد و هدف آن حفظ هم‌ارزیِ رفت‌وبرگشت (\lr{round-trip identity}) در شرایطِ متداول است؛ به‌طوری که بازسازیِ نمونه از روی دیکشنریِ سریال‌سازی‌شده، شیئی با ویژگی‌های عملکردی یکسان ارائه می‌دهد. اگر خروجیِ \texttt{to\_serializable\_dict()} یک نمونه را به \texttt{from\_dict()} بدهیم، نمونه‌ای هم‌ارزِ اصلی به‌دست می‌آید. این ویژگی در تست‌های خودکار (\lr{test\_round\_trip\_identity}) به‌صورت سیستماتیک بررسی می‌شود.

% ---------------------------------------------------------------------
% ===== فایل دوم: constants_helpers.py =====
% ---------------------------------------------------------------------

\section{نمای کلی و هدف ماژول \texttt{constants\_helpers.py}}

ماژول \texttt{constants\_helpers.py} مجموعه‌ای از توابعِ خالص (\lr{pure functions}) برای اعتبارسنجیِ عددی، نرمال‌سازیِ احتمالات و تبدیل‌های یکانی (\lr{unit conversions}) را در فریم‌ورکِ \lr{QKD} ارائه می‌کند. اگرچه این ماژول از نظرِ حجمی کوچک به‌نظر می‌رسد (تنها حدود ۲۱۵ خط)، و به‌عنوان ابزار کمکی برای بهبود یکنواختیِ محاسباتِ عددی و کاهش خطاهای ناشی از مقایسه‌های اعشاری در فریم‌ورک به‌کار می‌رود. تمامی مقایسه‌های ممیزِ شناور، تمامی اعتبارسنجی‌های بازه‌ی احتمال و تمامی تبدیل‌های \lr{dB-to-linear} در این فریم‌ورک از طریق این ماژول عبور می‌کنند و این تمرکز، اطمینان می‌دهد که خطاهای عددی به‌صورت یکسان در سراسرِ پایگاهِ کد شناسایی می‌شوند.

هدفِ بنیادینِ این ماژول رفعِ یک مشکلِ کلاسیک در محاسباتِ علمیِ پایتون است: مقایسه‌های نامعتبرِ اعدادِ اعشاری. در ریاضیاتِ ممیزِ شناور، تساویِ دقیقِ \texttt{a == b} تقریباً همواره نادرست است؛ چرا که خطاهای گردسازی در طولِ محاسبات تجمع می‌کنند. به‌عنوان مثال، \texttt{0.1 + 0.2 == 0.3} در پایتون \texttt{False} برمی‌گرداند. این ماژول با ارائه‌ی تابعِ \texttt{is\_close()} که از تOLERانس‌های مطلق و نسبی استفاده می‌کند، این مشکل را به‌صورت سیستماتیک حل می‌کند.

علاوه بر این، ماژول مسئولیتِ حفاظت از یکپارچگیِ محاسباتِ آنتروپی را بر عهده دارد. آنتروپیِ شانون، $H(p) = -\sum p_i \log p_i$، در مرزهای $p = 0$ و $p = 1$ دارایِ رفتارِ نامطلوب است: در $p = 0$ عبارتِ $0 \log 0$ از نظرِ ریاضی به $0$ میل می‌کند، اما در محاسباتِ عددی به \lr{NaN} یا $-\infty$ منجر می‌شود. تابعِ \texttt{clamp\_probability()} با محدود کردنِ احتمالات به بازه‌ی $[\varepsilon, 1 - \varepsilon]$ این مشکل را برطرف می‌سازد و امکانِ ارزیابیِ پایدارِ آنتروپی را در طولِ شبیه‌سازی فراهم می‌کند.

\section{مبانی تئوری و فیزیکی در \texttt{constants\_helpers.py}}

\subsection{تحلیلِ عددیِ ممیزِ شناور}

استاندارد \lr{IEEE 754} که در پایتون از طریقِ نوع \texttt{float} پیاده‌سازی شده، اعدادِ اعشاری را با ۵۳ بیتِ مانتیس (\lr{mantissa}) ارائه می‌دهد که تقریباً معادلِ ۱۵ رقمِ اعشارِ معنادار است. با این حال، خطای ماشین (\lr{machine epsilon}) که به‌صورت $\varepsilon_{\text{mach}} = 2^{-52} \approx 2.22 \times 10^{-16}$ تعریف می‌شود، حدِ پایینِ خطاهای نسبی را مشخص می‌کند. در فریم‌ورکِ \lr{QKD}، جایی که احتمالاتِ بسیار کوچک (مانند $10^{-10}$) با احتمالاتِ بزرگ (نزدیک به ۱) در یک محاسبه ترکیب می‌شوند، درکِ این محدودیت حیاتی است.

برای مقایسه‌ی دو عدد $a$ و $b$ در حضورِ خطای گردسازی، فرمولِ استانداردِ زیر به‌کار می‌رود:
\begin{equation}
	|a - b| \leq \max(\text{rel\_tol} \cdot \max(|a|, |b|), \text{abs\_tol})
\end{equation}
که در آن $\text{rel\_tol}$ تلورانسِ نسبی و $\text{abs\_tol}$ تلورانسِ مطلق است. تابع \texttt{math.isclose()} پایتون این فرمول را پیاده‌سازی می‌کند و تابعِ \texttt{is\_close()} فریم‌ورک با استفاده از تلورانس‌های متمرکز در \texttt{constants\_definitions.py} این رفتار را بومی‌سازی می‌کند.

\subsection{آنتروپیِ شانون و ناپایداریِ عددی}

در محاسباتِ نرخِ کلیدِ \lr{QKD}، آنتروپیِ شانونِ نرخِ خطای کوانتومی به‌کرات محاسبه می‌شود:
\begin{equation}
	H_2(Q) = -Q \log_2 Q - (1-Q) \log_2(1-Q)
\end{equation}
این تابع در $Q = 0$ و $Q = 1$ به‌صورت ریاضی به‌صورت $\lim_{x \to 0^+} x \log x = 0$ تعریف می‌شود، اما در محاسباتِ عددی $\log(0) = -\infty$ و در نتیجه $0 \cdot (-\infty) = $ NaN تولید می‌شود. راه‌حلِ استاندارد، محدود کردنِ (\lr{clamping}) احتمالات به بازه‌ی $[\delta, 1 - \delta]$ با $\delta$ کوچک (مانند $10^{-12}$) است:
\begin{equation}
	p_{\text{clamped}} = \min(\max(p, \delta), 1 - \delta)
\end{equation}
این رویکرد، خطای ناشی از تغییرِ $\delta$ را در حدِ $O(\delta \log \delta)$ نگه می‌دارد که برای $\delta = 10^{-12}$ خطی در حدِ $10^{-11}$ بیت است — بسیار کوچک‌تر از دقتِ امنیتیِ موردِ نیاز.

\subsection{نرمال‌سازیِ \lr{L1} و توزیعِ احتمال}

یک توزیعِ احتمالِ گسسته باید شرطِ $\sum_i p_i = 1$ را برآورده کند. در محاسباتِ عددی، خطاهای گردسازی می‌توانند باعث شوند که مجموع از ۱ منحرف شود. نرمال‌سازیِ \lr{L1} این مشکل را با تقسیمِ هر درایه بر مجموع حل می‌کند:
\begin{equation}
	p_i^{\text{norm}} = \frac{p_i}{\sum_j p_j}
\end{equation}
با این حال، اگر $\sum_j p_j = 0$، این فرمول به‌تقسیم بر صفر منجر می‌شود. در این حالت، یک \lr{fallback} قطعی به یک بردارِ یک‌داغ (\lr{one-hot}) که تمامِ جرم را در اندیسِ صفر قرار می‌دهد، استفاده می‌شود. این انتخاب، گزینه‌ای قراردادی برای جلوگیری از توقف شبیه‌سازی در مواجهه با بردارهای صفر است؛ هرچند از نظر تئوری اطلاعات، قرار دادن کل جرم روی یک اندیس با مفهوم توزیع حداکثر آنتروپی (یکنواخت) متفاوت است.

\subsection{تبدیلِ دسی‌بل به خطی}

در مخابراتِ نوری، اتلاف و تقویت به‌صورت سنتی با واحدِ دسی‌بل (\lr{dB}) بیان می‌شوند:
\begin{equation}
	x_{\text{dB}} = 10 \log_{10}(x_{\text{linear}})
\end{equation}
تبدیلِ معکوس، که در فریم‌ورک پیاده‌سازی شده، به‌صورت زیر است:
\begin{equation}
	x_{\text{linear}} = 10^{x_{\text{dB}}/10}
\end{equation}
این تبدیل به‌دلیلِ ماهیتِ نماییِ آن، دارایِ خطرِ سرریز (\lr{overflow}) است: برای $x_{\text{dB}} = 3080$، خروجی برابر با $10^{308}$ است که در نزدیکیِ حدِ بالایِ \texttt{float64} ($\approx 1.8 \times 10^{308}$) قرار دارد. برای $x_{\text{dB}} = 3100$، سرریز به $+\infty$ رخ می‌دهد. تابع \texttt{db\_to\_linear()} با بررسیِ صریحِ این حد، از تولیدِ \lr{inf} جلوگیری می‌کند.

\subsection{بردارِ احتمالِ معتبر}

یک بردارِ احتمالِ معتبر باید چهار شرط را برآورده کند: یک‌بعدی بودن، مقدارِ حقیقیِ متناهی، نامنفی بودنِ هر درایه (با تلورانس)، و مجموع برابر با ۱ (با تلورانس). این شروط چهارگانه پایه‌ی ریاضیِ تمامِ توزیع‌های احتمالی است که در فریم‌ورک استفاده می‌شود — از توزیعِ پواسونِ فوتون‌ها تا توزیعِ پسینِ (\lr{posterior}) خطا.

\section{تحلیل معماری و ساختار داده‌ها در \texttt{constants\_helpers.py}}

\subsection{فهرستِ توابعِ عمومی}

ماژول با تعریفِ یک چندتاییِ \lr{Final} از نامِ توابعِ عمومی آغاز می‌شود:

\begin{lstlisting}[language=Python]
	PUBLIC_FUNCTION_NAMES: Final[tuple[str, ...]] = (
	"clamp_probabilities",
	"clamp_probability",
	"db_to_linear",
	"is_close",
	"is_finite_non_negative",
	"is_prob_vector",
	"is_valid_probability",
	"renormalize_probabilities",
	)
	__all__ = list(PUBLIC_FUNCTION_NAMES)
\end{lstlisting}

این الگو دو مزیت دارد: نخست، فهرستِ \lr{API} عمومی به‌صورت یک منبعِ واحد نگهداری می‌شود و امکانِ بازبینیِ آسان را فراهم می‌سازد. دوم، از آنجا که این چندتایی با \texttt{Final} علامت‌گذاری شده، بررسی‌گرهای نوع ایستا هشدار می‌دهند اگر به‌اشتباه مقدارِ آن تغییر کند. ماژولِ \texttt{constants.py} از همین فهرست برای بازصادراتِ (\lr{re-export}) توابع استفاده می‌کند.

\subsection{توابعِ کمکیِ خصوصی}

دو تابعِ خصوصیِ \texttt{\_is\_real\_scalar()} و \texttt{\_to\_real\_scalar()} لایه‌ی اولیه‌ی فیلترینگِ نوع را تشکیل می‌دهند:

\begin{lstlisting}[language=Python]
	def _is_real_scalar(value: Any) -> bool:
	return isinstance(value, (Real, np.integer, np.floating)) and \
	not isinstance(value, (bool, np.bool_))
\end{lstlisting}

نکته‌ی ظریفِ این پیاده‌سازی، حذفِ صریحِ \texttt{bool} از مجموعه‌ی اعدادِ حقیقی است. در پایتون، \texttt{bool} یک زیرکلاسِ \texttt{int} است و در نتیجه \texttt{isinstance(True, int)} مقدارِ \texttt{True} برمی‌گرداند. این رفتار می‌تواند منجر به خطاهای خاموش شود: مثلاً \texttt{True + True} برابر با \texttt{2} است. در فریم‌ورکِ \lr{QKD}، جایی که احتمالات باید اعدادِ حقیقیِ کوچک باشند، یک \texttt{True} یا \texttt{False} به‌اشتباه به‌عنوان احتمال می‌تواند منجر به نرخِ کلیدِ نادرست شود. حذفِ صریحِ \texttt{bool} این حفره را می‌بندد.

\subsection{تابعِ \texttt{\_to\_real\_1d\_array()}: اعتبارسنجیِ آرایه‌ای}

این تابع کمکی، ورودی‌های متنوع را به آرایه‌ی یک‌بعدیِ \texttt{float64} تبدیل می‌کند و چندین اعتبارسنجیِ حیاتی انجام می‌دهد:

\begin{lstlisting}[language=Python]
	def _to_real_1d_array(values, *, name, allow_empty):
	arr = np.asarray(values)
	if arr.dtype == np.bool_:
	raise TypeError(f"{name} must contain only real numeric values, not booleans.")
	if arr.ndim == 0:
	raise ValueError(f"{name} must be a 1D array, but a scalar was provided.")
	if arr.ndim != 1:
	raise ValueError(f"{name} must be a 1D array, but got {arr.ndim}D input.")
	if not allow_empty and arr.size == 0:
	raise ValueError(f"{name} must not be empty.")
	if np.iscomplexobj(arr):
	raise TypeError(f"{name} must contain only real values.")
	arr = arr.astype(np.float64, copy=False)
	if not np.all(np.isfinite(arr)):
	raise ValueError(f"{name} must contain only finite values.")
	return arr
\end{lstlisting}

این تابع چندین فیلترِ دفاعی دارد: رد کردنِ \texttt{bool}، رد کردنِ آرایه‌های چندبعدی، رد کردنِ مقادیرِ مختلط، و رد کردنِ \lr{NaN} و \lr{Inf}. این رویکردِ چندلایه، تطبیق‌پذیریِ نوع (\lr{type strictness}) را با پیام‌های خطای دقیق ترکیب می‌کند.

\subsection{تایپ‌هینت‌های \lr{ArrayLike} و \lr{NDArray}}

ماژول از تایپ‌هینت‌های \lr{NumPy} برای بیانِ دقیقِ قراردادِ توابع استفاده می‌کند:

\begin{lstlisting}[language=Python]
	from numpy.typing import ArrayLike, NDArray
	
	def clamp_probabilities(probs: ArrayLike) -> NDArray[np.float64]: ...
\end{lstlisting}

\texttt{ArrayLike} یک نوعِ «پذیرنده» است که شاملِ لیست، چندتایی، و آرایه‌های \lr{NumPy} می‌شود — هر چیزی که \texttt{np.asarray()} بتواند آن را تبدیل کند. \texttt{NDArray[np.float64]} نوعِ خروجیِ دقیق را مشخص می‌کند. این ترکیب، هم انعطاف‌پذیریِ پایتونِ اردوگانی (\lr{duck-typing}) را حفظ می‌کند و هم دقتِ تایپ‌هینت را ارائه می‌دهد.

\section{پیاده‌سازی ریاضی و الگوریتمی در \texttt{constants\_helpers.py}}

\subsection{تابع \texttt{is\_close()}: مقایسه‌ی ایمن}

پیاده‌سازیِ این تابع ساده اما حیاتی است:

\begin{lstlisting}[language=Python]
	def is_close(a, b, *, rel_tol=NUMERIC_REL_TOL, abs_tol=NUMERIC_ABS_TOL):
	a_val = _to_real_scalar(a, name="a")
	b_val = _to_real_scalar(b, name="b")
	return math.isclose(a_val, b_val, rel_tol=rel_tol, abs_tol=abs_tol)
\end{lstlisting}

با استفاده از \texttt{math.isclose()} پایتون (نسخه‌ی ۳.۵ به بعد)، این تابع از فرمولِ استاندارد استفاده می‌کند. تلورانس‌های پیش‌فرض از \texttt{constants\_definitions.py} وارد می‌شوند (\texttt{NUMERIC\_REL\_TOL = 1e-9} و \texttt{NUMERIC\_ABS\_TOL = 1e-12}) که به‌دقت تنظیم شده‌اند تا هم حساسیتِ کافی برای محاسباتِ امنیتی داشته باشند و هم در برابرِ خطاهای گردسازی مقاوم باشند.

\subsection{تابع \texttt{is\_valid\_probability()}: اعتبارسنجیِ بازه‌ی احتمال}

\begin{lstlisting}[language=Python]
	def is_valid_probability(p: Any) -> bool:
	if not _is_real_scalar(p):
	return False
	p_val = float(p)
	return math.isfinite(p_val) and -NUMERIC_ABS_TOL <= p_val <= 1.0 + NUMERIC_ABS_TOL
\end{lstlisting}

این تابع یک تلورانسِ کوچک $\varepsilon = 10^{-12}$ در هر دو مرز اضافه می‌کند. این تلورانس به مقادیرِ نظیرِ $-10^{-15}$ (که از خطای گردسازی ناشی می‌شوند) اجازه می‌دهد به‌عنوان $0$ معتبر شناخته شوند. این انعطاف در محاسباتِ بازگشتی، جایی که خطاهای کوچک تجمع می‌کنند، حیاتی است.

\subsection{تابع \texttt{clamp\_probability()}: برشِ احتمال}

\begin{lstlisting}[language=Python]
	def clamp_probability(p: Real) -> float:
	p_val = _to_real_scalar(p, name="p")
	return min(max(p_val, ENTROPY_PROB_CLAMP), 1.0 - ENTROPY_PROB_CLAMP)
\end{lstlisting}

این تابع با فرمولِ $\min(\max(p, \delta), 1 - \delta)$ عمل می‌کند که در آن $\delta = $ ENTROPY\_PROB\_CLAMP. این فرمول، \lr{p} را در بازه‌ی $[\delta, 1 - \delta]$ نگه می‌دارد. همان‌طور که پیش‌تر ذکر شد، این برش برای پایداریِ محاسباتِ آنتروپی حیاتی است.

\subsection{تابع \texttt{clamp\_probabilities()}: نسخه‌ی برداری}

نسخه‌ی برداری از تابعِ قبلی، با استفاده از \texttt{np.clip()}:

\begin{lstlisting}[language=Python]
	def clamp_probabilities(probs: ArrayLike) -> NDArray[np.float64]:
	arr = _to_real_1d_array(probs, name="probs", allow_empty=True)
	return np.clip(arr, ENTROPY_PROB_CLAMP, 1.0 - ENTROPY_PROB_CLAMP)
\end{lstlisting}

پارامترِ \texttt{allow\_empty=True} به این معناست که یک بردارِ خالی نیز مجاز است و خروجیِ آن نیز خالی خواهد بود. این رفتار برای توابعِ پایین‌دستی که ممکن است با لیست‌های خالی سروکار داشته باشند، مفید است.

\subsection{تابع \texttt{renormalize\_normalize\_probabilities()}: نرمال‌سازی با \lr{fallback}}

این تابع، پیچیده‌ترین تابعِ ماژول است:

\begin{lstlisting}[language=Python]
	def renormalize_probabilities(probs: ArrayLike) -> NDArray[np.float64]:
	arr = _to_real_1d_array(probs, name="probs", allow_empty=False)
	p_clamped = clamp_probabilities(arr)
	total_p = float(np.sum(p_clamped))
	
	if total_p < NUMERIC_ABS_TOL:
	warnings.warn(
	"renormalize_probabilities: total mass ≈ 0 after clamping; "
	"falling back to deterministic one-hot at index 0.",
	RuntimeWarning, stacklevel=2)
	p_new = np.zeros_like(p_clamped)
	p_new[0] = 1.0
	return p_new
	
	return p_clamped / total_p
\end{lstlisting}

این پیاده‌سازی چند نکته‌ی مهم دارد. نخست، ورودی‌های خالی رد می‌شوند (\texttt{allow\_empty=False}). دوم، پیش از نرمال‌سازی، احتمالات بریده می‌شوند تا از تقسیم بر صفرِ ناشی از $0 \log 0$ جلوگیری شود. سوم، اگر کلِ جرم پس از برش بسیار کوچک باشد، یک هشدارِ \texttt{RuntimeWarning} صادر می‌شود و یک \lr{one-hot} برگردانده می‌شود. این \lr{fallback} قطعی، رفتارِ تابع را در حضورِ ورودی‌های نامعتبر قابل‌پیش‌بینی می‌کند — یک ویژگیِ حیاتی برای سیستم‌های تولیدی.

\subsection{تابع \texttt{db\_to\_linear()}: تبدیلِ ایمن}

\begin{lstlisting}[language=Python]
	def db_to_linear(val_db: Real) -> float:
	val_db_f = _to_real_scalar(val_db, name="val_db")
	exponent = val_db_f / 10.0
	if exponent > 308:
	raise OverflowError(f"db_to_linear({val_db_f}) would overflow to infinity.")
	result = math.pow(10.0, exponent)
	if result == 0.0 and val_db_f != 0.0:
	raise OverflowError(f"db_to_linear({val_db_f}) underflows to zero.")
	return result
\end{lstlisting}

این تابع دو محافظتِ حیاتی دارد. نخست، اگر $x_{\text{dB}}/10 > 308$، سرریزِ $10^{308}$ به $+\infty$ پیش‌بینی شده و یک \texttt{OverflowError} صادر می‌شود. دوم، در صورتی که ورودی غیرصفر باشد اما خروجی به دلیل محدودیت‌های نمایشِ ممیز شناور به صفر تقریب داده شود (\lr{underflow})، سیستم با صدور خطا یا هشدار مانع از جریان یافتن مقادیر صفرِ ناخواسته در محاسبات حساسِ بعدی می‌شود؛ مگر آنکه این رفتار به‌وضوح در مدل فیزیکی تعریف شده باشد. این محافظت‌ها از تولیدِ خاموشِ $\infty$ یا $0$ در محاسباتِ پایین‌دستی جلوگیری می‌کنند.

\subsection{تابع \texttt{is\_prob\_vector()}: اعتبارسنجیِ بردارِ احتمال}

\begin{lstlisting}[language=Python]
	def is_prob_vector(p: ArrayLike) -> bool:
	try:
	arr = _to_real_1d_array(p, name="p", allow_empty=False)
	except (TypeError, ValueError) as exc:
	msg = str(exc)
	if ("must be a 1D array" in msg
	or "must contain only" in msg
	or "must not be empty" in msg):
	return False
	raise
	
	if np.any(arr < -NUMERIC_ABS_TOL):
	return False
	return bool(np.isclose(np.sum(arr), 1.0, atol=PROB_SUM_TOL, rtol=0.0))
\end{lstlisting}

این تابع به‌جای صدور خطا، \texttt{False} برمی‌گرداند — یک الگوی متفاوت از سایر توابع. این رفتار برای استفاده در شرط‌های \texttt{if} طراحی شده است: کاربر می‌تواند بپرسد «آیا این یک بردارِ احتمالِ معتبر است؟» بدون آنکه استثنا رخ دهد. با این حال، خطاهای ساختاری (مانندِ مختلط بودن) همچنان به‌صورت استثنا ظاهر می‌شوند، چرا که آن‌ها خطاهای برنامه‌نویسی هستند نه خطاهای داده.

\section{ارسال و دریافت با سایر ماژول‌ها در \texttt{constants\_helpers.py}}

\subsection{وابستگی‌های ورودی}

این ماژول تنها از دو منبع واردات انجام می‌دهد:

\begin{lstlisting}[language=Python]
	import math
	import warnings
	from numbers import Real
	from typing import Any, Final
	
	import numpy as np
	from numpy.typing import ArrayLike, NDArray
	
	from .constants_definitions import (
	ENTROPY_PROB_CLAMP,
	NUMERIC_ABS_TOL,
	NUMERIC_REL_TOL,
	PROB_SUM_TOL,
	)
\end{lstlisting}

همچنین یک \texttt{assert} در زمانِ واردات اجرا می‌شود:

\begin{lstlisting}[language=Python]
	assert 0.0 < ENTROPY_PROB_CLAMP < 0.5, (
	f"ENTROPY_PROB_CLAMP must be in (0, 0.5), got {ENTROPY_PROB_CLAMP}"
	)
\end{lstlisting}

این \texttt{assert} یک \lr{contract check} است که در زمانِ بارگذاریِ ماژول اجرا می‌شود و اطمینان می‌دهد که مقدارِ \texttt{ENTROPY\_PROB\_CLAMP} در بازه‌ی منطقی $(0, 0.5)$ قرار دارد. اگر به‌اشتباه مقدارِ $0.6$ تنظیم شود، کلِ فریم‌ورک از بارگذاری باز می‌ایستد و خطای صریح صادر می‌شود.

\subsection{نقاطِ مصرف}

توابعِ این ماژول توسط چندین ماژولِ دیگر فریم‌ورک استفاده می‌شوند:

\begin{itemize}
	\item \textbf{ماژولِ منابع نوری (\texttt{sources.py}):} برای اعتبارسنجیِ توزیعِ احتمالِ تعدادِ فوتون و محاسبه‌ی احتمالاتِ نهایی.
	\item \textbf{ماژولِ آشکارسازها (\texttt{detectors.py}):} برای محاسبه‌ی احتمالاتِ کلیک و نرمال‌سازیِ پس ازِ تطبیقِ با \lr{dark count}.
	\item \textbf{ماژولِ کانال (\texttt{channel.py}):} برای تبدیلِ اتلافِ \lr{dB} به خطی و محاسبه‌ی نرخِ انتقال.
	\item \textbf{ماژولِ محاسبه‌ی نرخِ کلید (\texttt{keyrate.py}):} برای محاسبه‌ی آنتروپی و مقایسه‌های عددی.
	\item \textbf{ماژولِ ثابت‌ها (\texttt{constants.py}):} برای بازصادراتِ عمومی از طریقِ فسادِ (\lr{facade}).
\end{itemize}

\subsection{قراردادِ رابط}

این ماژول چند قراردادِ رابطِ مهم را رعایت می‌کند:

\textbf{۱. خالص‌بودن (\lr{purity}):} هیچ‌یک از توابع دارای اثرِ جانبی نیستند. ورودی‌ها تغییر نمی‌کنند و وضعیتِ سراسری خوانده نمی‌شود (به‌جز ثابت‌های \texttt{Final} که تنها خوانده می‌شوند).

\textbf{۲. قطعیت (\lr{determinism}):} برای ورودیِ یکسان، خروجی همواره یکسان است. تنها استثنا صدورِ هشدار (\texttt{warnings.warn}) در حالتِ \lr{fallback} است که یک اثرِ جانبیِ قابل‌قبول است.

\textbf{۳. اعتبارسنجیِ پیش از اجرا:} تمام توابع ورودی را پیش از پردازشِ اصلی اعتبارسنجی می‌کنند. این الگو، \lr{fail-fast} نامیده می‌شود و از اجرایِ محاسباتِ نادرست جلوگیری می‌کند.

\textbf{۴. پیام‌های خطای دقیق:} تمام خطاها نامِ پارامترِ متناظر را در بر دارند، که اشکال‌زدایی را ساده می‌کند.

% ---------------------------------------------------------------------
% ===== فایل سوم: constants_definitions.py =====
% ---------------------------------------------------------------------

\section{نمای کلی و هدف ماژول \texttt{constants\_definitions.py}}

ماژول \texttt{constants\_definitions.py} منبعِ مرکزیِ ثابت‌ها و نوع‌های تایپیِ فریم‌ورکِ \lr{QKD} است. این ماژول کوچک (تنها ۶۳ خط) نقشِ یک \lr{axiomatic foundation} را ایفا می‌کند: تمامی مقادیرِ عددیِ موردِ استفاده در فریم‌ورک — از تلورانس‌های عددی تا ثابت‌های فیزیکیِ بنیادین — در اینجا تعریف می‌شوند و در هیچ جایِ دیگری از پایگاهِ کد بازتعریف نمی‌گردند. این تمرکز، امکانِ تنظیمِ دقیق و یکنواختِ رفتارِ عددیِ کلِ فریم‌ورک را از یک نقطه‌ی واحد فراهم می‌سازد.

هدفِ بنیادینِ این ماژول، رفعِ یک مشکلِ رایج در پروژه‌های علمیِ بزرگ است: پخش‌شدنِ مقادیرِ جادویی (\lr{magic numbers}) در سراسرِ کد. در غیابِ این ماژول، ممکن است مقدارِ $10^{-9}$ به‌عنوانِ تلورانسِ نسبی در یک فایل، مقدارِ $10^{-8}$ در فایلی دیگر و $1e-10$ در فایلی سوم به‌کار رود. این ناهماهنگی منجر به رفتارِ ناسازگار و اشکال‌زداییِ دشوار می‌شود. با متمرکز کردنِ تمام ثابت‌ها در این ماژول، فریم‌ورک رفتاری یکنواخت و قابلِ توضیح ارائه می‌دهد.

علاوه بر ثابت‌های عددی، این ماژول مسئولیتِ تعریفِ نوع‌های تایپیِ مشترک را نیز بر عهده دارد. \texttt{LPSolverMethod} به‌عنوان یک \lr{Literal Type}، گزینه‌های مجازِ حلقه‌ی \lr{LP} را محدود می‌کند و امکانِ بررسیِ ایستا را فراهم می‌سازد. فهرستِ \texttt{PUBLIC\_CONSTANT\_NAMES} به‌عنوان یک مرجعِ مرتب، ترتیبِ سریال‌سازی و خروجیِ گزارش‌گیری را پایدار نگه می‌دارد و در نتیجه گزارش‌های تولیدشده توسط نسخه‌های مختلفِ فریم‌ورک قابلِ مقایسه باقی می‌مانند.

\section{مبانی تئوری و فیزیکی در \texttt{constants\_definitions.py}}

\subsection{سلسله‌مراتبِ تلورانس}

یکی از مفاهیمِ کلیدیِ این ماژول، سلسله‌مراتبِ تلورانس است که در داک‌استرینگِ ماژول به‌صورت زیر توضیح داده شده است:

\begin{equation}
	\text{ENTROPY\_PROB\_CLAMP} \leq \text{NUMERIC\_ABS\_TOL} \ll \text{PROB\_SUM\_TOL}
\end{equation}

این سلسله‌مراتب بر پایه‌ی تحلیلِ دقیقِ انتشارِ خطا (\lr{error propagation}) طراحی شده است. فرض کنید یک بردارِ احتمالِ به طولِ $N$ داریم که هر درایه با خطایِ مطلقِ حداکثر $\varepsilon_{\text{abs}}$ ذخیره شده است. خطای تجمعی در مجموعِ بردار برابر است با:
\begin{equation}
	\left| \sum_i p_i - 1 \right| \leq N \cdot \varepsilon_{\text{abs}}
\end{equation}

برای آنکه این خطا از تلورانسِ مجموع $\varepsilon_{\text{sum}}$ فراتر نرود، باید:
\begin{equation}
	\varepsilon_{\text{abs}} \leq \frac{\varepsilon_{\text{sum}}}{N}
\end{equation}

با مقادیرِ فریم‌ورک ($\varepsilon_{\text{abs}} = 10^{-12}$ و $\varepsilon_{\text{sum}} = 10^{-8}$، این شرط برای $N < 10^4$ برقرار است. برای بردارهای بسیار بلند، ممکن است نیاز به بازنگری در \lr{PROB\_SUM\_TOL} باشد.

\subsection{خطای ماشین و \lr{Y1\_SAFE\_THRESHOLD}}

ثابتِ \lr{Y1\_SAFE\_THRESHOLD} برای محافظت از محاسباتِ تک‌فوتونی (\lr{single-photon yield}) طراحی شده است. در پروتکل‌های \lr{Decoy-State}، تخمینِ $Y_1$ (بازدهیِ تک‌فوتونی) یک مرحله‌ی بحرانی است که در آن، نسبت‌های عددیِ بسیار کوچک محاسبه می‌شوند. اگر مخرجِ این نسبت‌ها به خطای ماشین نزدیک شود، نتیجه ممکن است ناپایدار یا غیرقابل‌اعتماد شود.

فرمولِ محاسبه‌ی این آستانه به‌صورت زیر است:
\begin{equation}
	\text{Y1\_SAFE\_THRESHOLD} = \max(10^{-12}, 10^4 \cdot \varepsilon_{\text{mach}})
\end{equation}
که در آن $\varepsilon_{\text{mach}} \approx 2.22 \times 10^{-16}$ خطای ماشینِ \lr{IEEE 754 double} است. ضریبِ $10^4$ یک حاشیه‌ی امنیتی است که با تجربه‌ی عددی به‌دست آمده و کمک می‌کند که آستانه در برابر خطاهای گردکردن و شرایط مرزی محافظت شود.

\subsection{ثابت‌های فیزیکیِ \lr{SI 2019}}

سه ثابتِ فیزیکیِ این ماژول بر اساسِ بازتعریفِ \lr{SI 2019} مقداردهی شده‌اند که در آن این مقادیر با تعریف‌های مرجع \lr{SI 2019} سازگار هستند:

\begin{equation}
	h = 6.62607015 \times 10^{-34} \, \text{J}\cdot\text{s}
\end{equation}
\begin{equation}
	k_B = 1.380649 \times 10^{-23} \, \text{J/K}
\end{equation}
\begin{equation}
	e = 1.602176634 \times 10^{-19} \, \text{C}
\end{equation}

ثابتِ پلانک $h$ در محاسبه‌ی انرژیِ فوتون $E = h\nu$، ثابتِ بولتزمن $k_B$ در محاسبه‌ی نویزِ حرارتی، و بارِ الکترون $e$ در محاسبه‌ی نویزِ نوارِ عبور \lr{shot noise} استفاده می‌شوند. این دقت با تعریف‌های مرجع و محاسبات مبتنی بر \lr{SI} سازگار است.

\subsection{بذرِ تولیدنده‌ی عددِ تصادفیِ \lr{PCG64}}

ثابتِ \texttt{MAX\_SEED\_INT\_PCG64} حدِ بالایِ بذرِ مجاز برای تولیدنده‌ی عددِ تصادفیِ \lr{PCG64} را مشخص می‌کند:

\begin{equation}
	\text{MAX\_SEED\_INT\_PCG64} = 2^{63} - 1 = 9223372036854775807
\end{equation}

این مقدار برابر با $2^{63} - 1$ است که حدِ بالایِ \texttt{int64} با علامت است. تولیدکننده‌ی \lr{PCG64} بذرِ ۱۲۸ بیتی می‌پذیرد، اما در فریم‌ورک تنها ۶۳ بیت برای سازگاری با \lr{JSON} و سیستم‌های سریال‌سازی استفاده می‌شود. این محدودیت، سازگاری با \lr{JavaScript} (که اعدادِ صحیحِ بزرگ‌تر از $2^{53}$ را به‌دقت نمایش نمی‌دهد) را قربانی می‌کند، اما در عوض تکرارپذیریِ ساده‌تر را فراهم می‌سازد.

\subsection{آستانه‌ی دمِ پواسون}

ثابت \texttt{DEFAULT\_POISSON\_TAIL\_THRESHOLD = 1e-9} برای کنترلِ دقتِ تقریبِ دمِ توزیعِ پواسون استفاده می‌شود. در محاسباتِ \lr{Decoy-State}، توزیعِ پواسونِ تعدادِ فوتون دارای دم‌های طولانی است که در محاسباتِ عددی باید بریده شوند. این آستانه تعیین می‌کند که در کجا دم بریده شود:
\begin{equation}
	P(n > n_{\text{max}}) < \varepsilon_{\text{tail}}
\end{equation}
با $\varepsilon_{\text{tail}} = 10^{-9}$، می‌توان $n_{\text{max}}$ را با دقتِ بالا محاسبه کرد و در عین حال زمانِ محاسبه را معقول نگه داشت.

\subsection{روش‌های \lr{LP Solver}}

نوعِ \texttt{LPSolverMethod} سه روشِ حلِ \lr{LP} را مشخص می‌کند:
\begin{itemize}
	\item \texttt{"highs"}: روشِ پیش‌فرض که \lr{HiGHS} را با تشخیصِ خودکارِ روش (\lr{dual simplex} یا \lr{interior point}) فراخوانی می‌کند.
	\item \texttt{"highs-ds"}: \lr{Dual Simplex} — مناسب برای مسائلِ کوچک و متوسط.
	\item \texttt{"highs-ipm"}: \lr{Interior Point Method} — مناسب برای مسائلِ بزرگِ تنک.
\end{itemize}

این سه گزینه از کتابخانه‌ی \lr{SciPy} (نسخه‌ی ۱.۹ به بعد) استفاده می‌کنند که \lr{HiGHS} را به‌عنوان حلقه‌ی اصلی به‌کار می‌گیرد. \lr{HiGHS} یکی از سریع‌ترین و پایدارترین حلقه‌های \lr{LP} متن‌باز است.

\section{تحلیل معماری و ساختار داده‌ها در \texttt{constants\_definitions.py}}

\subsection{استفاده‌ی گسترده از \texttt{Final}}

یکی از ویژگی‌های برجسته‌ی این ماژول، استفاده‌ی سیستماتیک از نوعِ \texttt{Final} است:

\begin{lstlisting}[language=Python]
	from typing import Final, Literal, TypeAlias
	
	MAX_SEED_INT_PCG64: Final[int] = (1 << 63) - 1
	NUMERIC_ABS_TOL: Final[float] = 1e-12
	NUMERIC_REL_TOL: Final[float] = 1e-9
	EPS: Final[float] = float(np.finfo(float).eps)
\end{lstlisting}

نوعِ \texttt{Final} که در \lr{PEP 591} معرفی شده، به بررسی‌گرهای نوع ایستا می‌گوید که این متغیرها نباید دوباره مقداردهی شوند. این محدودیت، امنیتِ ثابت‌ها را در سطحِ تحلیلِ ایستا تضمین می‌کند و از تغییرِ تصادفیِ آن‌ها — که می‌تواند منجر به رفتارِ نادرستِ فریم‌ورک شود — جلوگیری می‌کند.

\subsection{نوعِ \texttt{Literal} و \lr{TypeAlias}}

نوعِ \texttt{LPSolverMethod} یک نمونه‌ی نمادین از کاربردِ \texttt{Literal} است:

\begin{lstlisting}[language=Python]
	LPSolverMethod: TypeAlias = Literal["highs", "highs-ds", "highs-ipm"]
	LP_SOLVER_METHODS: Final[tuple[LPSolverMethod, ...]] = ("highs", "highs-ds", "highs-ipm")
\end{lstlisting}

در اینجا \texttt{Literal} یک نوعِ تایپی ایجاد می‌کند که تنها سه رشته‌ی مشخص را می‌پذیرد. این نوع به‌عنوان \lr{TypeAlias} نام‌گذاری شده و سپس در یک چندتاییِ \lr{Final} به‌عنوان مرجعِ اجرایی استفاده می‌شود. این الگوی دوگانه — نوع برای بررسیِ ایستا و چندتایی برای بررسیِ پویا — امکانِ اعتبارسنجیِ هر دو در نوع و زمانِ اجرا را فراهم می‌سازد.

\subsection{فهرستِ \texttt{PUBLIC\_CONSTANT\_NAMES}}

این فهرستِ مرتب، ترتیبِ سریال‌سازی و گزارش‌گیری را پایدار نگه می‌دارد:

\begin{lstlisting}[language=Python]
	PUBLIC_CONSTANT_NAMES: Final[tuple[str, ...]] = (
	"MAX_SEED_INT_PCG64",
	"MAX_SEED_INT_UINT32",
	"LP_SOLVER_METHODS",
	"MIN_SUCCESSFUL_Z1_BASIS_EVENTS_FOR_PHASE_EST",
	"EPS",
	"NUMERIC_ABS_TOL",
	"NUMERIC_REL_TOL",
	"Y1_SAFE_THRESHOLD",
	"ENTROPY_PROB_CLAMP",
	"PROB_SUM_TOL",
	"DEFAULT_POISSON_TAIL_THRESHOLD",
	"LP_CONSTRAINT_VIOLATION_TOL",
	"CONST_PLANCK",
	"CONST_BOLTZMANN",
	"CONST_ELECTRON_CHARGE",
	)
\end{lstlisting}

این ترتیبِ خاص تصادفی نیست: ثابت‌های عددیِ صحیح ابتدا، سپس تلورانس‌ها، و در پایان ثابت‌های فیزیکی. این ترتیب در خروجیِ \texttt{as\_dict()} ماژولِ \texttt{constants.py} منعکس می‌شود و گزارش‌های تولیدشده در نسخه‌های مختلفِ فریم‌ورک قابلِ مقایسه باقی می‌مانند.

\subsection{ثابتِ \texttt{MIN\_SUCCESSFUL\_Z1\_BASIS\_EVENTS\_FOR\_PHASE\_EST}}

این ثابتِ خاص، حداقلِ تعدادِ رویدادِ موفقِ \lr{Z1} (تک‌فوتون در پایه‌ی \lr{Z}) برای تخمینِ نرخِ خطای فاز را مشخص می‌کند:
\begin{equation}
	n_{Z1}^{\min} = 10
\end{equation}
اگر تعدادِ رویدادها کمتر از این مقدار باشد، تخمینِ خطای فاز به‌شدت ناپایدار می‌شود و می‌تواند به نتایجِ غیرفیزیکی منجر گردد. این حداقل، یک تعهدِ عملی است که با تجربه‌ی عددی به‌دست آمده و در پروتکل‌های \lr{MDI-QKD} و \lr{Decoy-State BB84} کاربرد دارد.

\section{پیاده‌سازی ریاضی و الگوریتمی در \texttt{constants\_definitions.py}}

\subsection{محاسبه‌ی \lr{EPS} از \lr{NumPy}}

ثابتِ \texttt{EPS} به‌صورت پویا از \lr{NumPy} محاسبه می‌شود:

\begin{lstlisting}[language=Python]
	EPS: Final[float] = float(np.finfo(float).eps)
\end{lstlisting}

تابع \texttt{np.finfo(float).eps} مقدارِ خطای ماشین را برای نوعِ \texttt{float} برمی‌گرداند، که برابر با $2^{-52} \approx 2.22044605 \times 10^{-16}$ است. این مقدار به \texttt{float} ساده تبدیل می‌شود تا با توابعِ ماژولِ \texttt{math} سازگار باشد. این رویکرد، خودسازگاری با پلتفرم را تضمین می‌کند: اگر فریم‌ورک در آینده به \texttt{float128} مهاجرت کند، \texttt{EPS} به‌صورت خودکار به‌روز می‌شود.

\subsection{محاسبه‌ی \lr{Y1\_SAFE\_THRESHOLD} با ماکزیمم}

این ثابت با فرمولی دو-بخشی محاسبه می‌شود:

\begin{lstlisting}[language=Python]
	_Y1_EPS_MULTIPLIER: Final[float] = 1e4
	Y1_SAFE_THRESHOLD: Final[float] = max(1e-12, _Y1_EPS_MULTIPLIER * EPS)
\end{lstlisting}

فرمولِ ریاضی:
\begin{equation}
	\text{Y1\_SAFE\_THRESHOLD} = \max(10^{-12}, 10^4 \cdot \varepsilon_{\text{mach}})
\end{equation}

استفاده از \texttt{max} تضمین می‌کند که حتی اگر \lr{EPS} به‌طور غیرعادی کوچک شود (مثلاً در یک پلتفرمِ با دقتِ چهارگانه)، آستانه هرگز از $10^{-12}$ پایین‌تر نیاید. این حدِ پایین، پایداریِ عددی را در شرایطِ غیرعادی تضمین می‌کند.

\subsection{تنظیمِ \texttt{ENTROPY\_PROB\_CLAMP} برابر با \texttt{NUMERIC\_ABS\_TOL}}

\begin{lstlisting}[language=Python]
	ENTROPY_PROB_CLAMP: Final[float] = NUMERIC_ABS_TOL
\end{lstlisting}

این انتخابِ آگاهانه به این معناست که برشِ احتمال برای محاسباتِ آنتروپی در همان سطحِ تلورانسِ مطلقِ کلی انجام می‌شود. این یکپارچگی، رفتارِ یکنواخت را تضمین می‌کند: اگر یک احتمال به $-10^{-15}$ برسد (که با تلورانسِ مطلق به‌عنوان $0$ پذیرفته می‌شود)، در محاسباتِ آنتروپی نیز به $10^{-12}$ بریده می‌شود و در نتیجه محاسباتِ متوالی سازگار باقی می‌مانند.

\subsection{محاسبه‌ی \texttt{LP\_CONSTRAINT\_VIOLATION\_TOL}}

\begin{lstlisting}[language=Python]
	LP_CONSTRAINT_VIOLATION_TOL: Final[float] = NUMERIC_REL_TOL
\end{lstlisting}

این تلورانس، حدِ مجازِ نقضِ محدودیت‌ها در حلقه‌ی \lr{LP} را مشخص می‌کند. استفاده از \texttt{NUMERIC\_REL\_TOL = 1e-9} تضمین می‌کند که حلقه‌ی \lr{LP} محدودیت‌ها را با دقتی در حدِ تلورانسِ نسبیِ کلی فریم‌ورک رعایت کند. این تنظیم، تعادلِ مناسبی بینِ سخت‌گیریِ بیش‌ازحد (که باعثِ شکستِ حلقه می‌شود) و سهل‌گیریِ بیش‌ازحد (که منجر به نتایجِ نادرست می‌شود) برقرار می‌سازد.

\subsection{ثابت‌های \lr{SI 2019}}

ثابت‌های فیزیکی به‌صورت مستقیم و با دقتِ کامل تعریف می‌شوند:

\begin{lstlisting}[language=Python]
	CONST_PLANCK: Final[float] = 6.62607015e-34       # J·s (exact, 2019 SI)
	CONST_BOLTZMANN: Final[float] = 1.380649e-23       # J/K (exact, 2019 SI)
	CONST_ELECTRON_CHARGE: Final[float] = 1.602176634e-19  # C (exact, 2019 SI)
\end{lstlisting}

این مقادیر از بازتعریفِ \lr{SI 2019} گرفته شده‌اند که در آن این سه ثابت به‌صورت دقیق (\lr{exact}) تعریف می‌شوند — نه اندازه‌گیری شده. این بدان معناست که در نظامِ \lr{SI}، این اعداد بدونِ خطا هستند و در نتیجه فریم‌ورک از بالاترین دقتِ ممکن برخوردار است.

\subsection{ماکروی بذرِ \lr{PCG64}}

\begin{lstlisting}[language=Python]
	MAX_SEED_INT_PCG64: Final[int] = (1 << 63) - 1
	MAX_SEED_INT_UINT32: Final[int] = (1 << 32) - 1
\end{lstlisting}

استفاده از عملگرِ \lr{\texttt{<<}} (شیفتِ چپ) به‌جای \lr{\texttt{**}} (توان)، یک بهینه‌سازیِ جزئی است که در سطحِ بایت‌کد پایتون، سریع‌تر اجرا می‌شود. فرمولِ ریاضی:
\begin{equation}
	\text{MAX\_SEED\_INT\_PCG64} = 2^{63} - 1
\end{equation}
\begin{equation}
	\text{MAX\_SEED\_INT\_UINT32} = 2^{32} - 1
\end{equation}

این دو مقدار، حدِ بالایِ بذرِ مجاز برای دو تولیدکننده‌ی متفاوت را مشخص می‌کنند: \lr{PCG64} (۶۳ بیت) و تولیدکننده‌های مبتنی بر \lr{uint32} (مانند \lr{MT19937}).

\section{ارسال و دریافت با سایر ماژول‌ها در \texttt{constants\_definitions.py}}

\subsection{عدمِ وابستگیِ ورودی}

این ماژول تقریباً هیچ وابستگیِ ورودی ندارد:

\begin{lstlisting}[language=Python]
	from __future__ import annotations
	from typing import Final, Literal, TypeAlias
	import numpy as np
\end{lstlisting}

تنها وابستگیِ خارجی، \lr{NumPy} است که برای محاسبه‌ی \texttt{EPS} استفاده می‌شود. این استقلال، ماژول را به یک پایه‌ی صلبِ فریم‌ورک تبدیل می‌کند: هیچ ماژولِ دیگری نمی‌تواند بر آن تأثیر بگذارد و در نتیجه رفتارِ آن کاملاً قابل‌پیش‌بینی است.

\subsection{نقاطِ مصرف}

این ماژول توسط دو مصرف‌کننده‌ی اصلی استفاده می‌شود:

\begin{itemize}
	\item \textbf{ماژولِ \texttt{constants\_helpers.py}:} تمامی تلورانس‌های عددی و آستانه‌های احتمال از این ماژول وارد می‌شوند.
	\item \textbf{ماژولِ \texttt{constants.py}:} به‌عنوان فسادِ (\lr{facade}) عمومی، تمامی نام‌های ثابت عمومی و توابعِ کمکی را بازصادرات می‌کند.
\end{itemize}

علاوه بر این، چندین ماژولِ دیگر نیز به‌صورت مستقیم از این ماژول واردات انجام می‌دهند — به‌ویژه \texttt{params.py} که \texttt{LP\_SOLVER\_METHODS} و \texttt{DEFAULT\_POISSON\_TAIL\_THRESHOLD} را استفاده می‌کند.

\subsection{قراردادِ \texttt{PUBLIC\_CONSTANT\_NAMES}}

این فهرست، قراردادی برای ماژول‌های پایین‌دستی است:

\begin{lstlisting}[language=Python]
	__all__ = list(PUBLIC_CONSTANT_NAMES)
\end{lstlisting}

این بدان معناست که \lr{\texttt{from .constants\_definitions import *}} تنها نام‌های فهرست‌شده را وارد می‌کند. این قرارداد، مرزِ \lr{API} عمومی را به‌روشنی مشخص می‌کند و از وابستگیِ تصادفی به ثابت‌های خصوصی جلوگیری می‌کند.

% ---------------------------------------------------------------------
% ===== فایل چهارم: constants.py =====
% ---------------------------------------------------------------------

\section{نمای کلی و هدف ماژول \texttt{constants.py}}

ماژول \texttt{constants.py} فسادِ (\lr{facade}) عمومیِ لایه‌ی ثابت‌ها و توابعِ کمکیِ فریم‌ورکِ \lr{QKD} است. این ماژول کوچک (تنها ۵۶ خط) نقشِ یک \lr{public API surface} را ایفا می‌کند: تمامی ماژول‌های بیرونی که نیاز به دسترسی به ثابت‌ها یا توابعِ کمکی دارند، تنها از این ماژول واردات انجام می‌دهند و هرگز به‌صورت مستقیم به \texttt{constants\_definitions.py} یا \texttt{constants\_helpers.py} ارجاع نمی‌دهند. این لایه‌ی واسط، انعطاف‌پذیریِ معماری را به‌شدت افزایش می‌دهد: اگر در آینده تصمیم به بازسازیِ داخلیِ این دو ماژول گرفته شود، تنها \texttt{constants.py} نیاز به به‌روزرسانی خواهد داشت و سایر ماژول‌ها بی‌تأثیر باقی می‌مانند.

هدفِ بنیادینِ این ماژول، پیاده‌سازیِ الگوی \lr{Facade Pattern} است که توسط \lr{Gamma et al.} در کتابِ کلاسیکِ الگوهای طراحی معرفی شده. در این الگو، یک واسطِ ساده، زیرسیستم‌های پیچیده‌تر را پنهان می‌کند و یک نقطه‌ی دسترسیِ یکپارچه فراهم می‌سازد. در فریم‌ورکِ \lr{QKD}، این الگو به‌ویژه ارزشمند است؛ چرا که تعدادِ ثابت‌ها و توابعِ کمکی در طولِ زمانِ توسعه افزایش می‌یابد و در نبودِ یک فساد، وابستگی‌های مستقیمِ پیچیده‌ای شکل می‌گیرد که بازسازی را دشوار می‌سازد.

علاوه بر نقشِ فساد، این ماژول مسئولیتِ ارائه‌ی متاداده‌ی نسخه و یک تابعِ کمکیِ \texttt{as\_dict()} را نیز بر عهده دارد. متغیر \texttt{\_\_version\_\_ = "4.3.0"} نسخه‌ی فعلیِ \lr{API} را مشخص می‌کند و امکانِ بررسیِ سازگاری در زمانِ اجرا را فراهم می‌سازد. تابع \texttt{as\_dict()} یک دیکشنریِ مرتب از تمامی ثابت‌های عمومی تولید می‌کند که برای گزارش‌گیری، لاگ‌گذاری و ردیابیِ آزمایش‌ها استفاده می‌شود.

\section{مبانی تئوری و فیزیکی در \texttt{constants.py}}

\subsection{الگوی فساد در مهندسی نرم‌افزار}

الگوی فساد بر پایه‌ی اصلِ \lr{Law of Demeter} (قانونِ دمتر) استوار است: یک ماژول نباید درباره‌ی ساختارِ داخلیِ ماژول‌هایِ دیگر بداند. در نبودِ فساد، اگر ماژول \lr{A} نیاز به ثابتِ \lr{X} داشته باشد، باید بداند که \lr{X} در کدام زیرماژول قرار دارد. این دانش، وابستگیِ غیرضروری ایجاد می‌کند و بازسازی را دشوار می‌سازد.

فساد این مشکل را با ارائه‌ی یک نقطه‌ی واحدِ دسترسی حل می‌کند. ماژول \lr{A} تنها از \texttt{constants} واردات انجام می‌دهد و در نتیجه از ساختارِ داخلیِ فریم‌ورک بی‌خبر باقی می‌ماند. این رزایلی، در پروژه‌هایِ بزرگ که تعدادِ ماژول‌ها به‌صورت نمایی رشد می‌کند، حیاتی است.

\subsection{پایداریِ رابط و نسخه‌سازی}

در مهندسیِ نرم‌افزارِ حرفه‌ای، \lr{API stability} یک ارزشِ بنیادین است. کاربرانِ یک کتابخانه انتظار دارند که نسخه‌های جدید، رفتارِ نسخه‌های قبلی را حفظ کنند یا با هشدارِ روشن تغییر دهند. ماژول \texttt{constants.py} با ارائه‌ی یک فهرستِ صریحِ \texttt{\_\_all\_\_}، \lr{API} عمومیِ خود را به‌روشنی تعریف می‌کند:
\begin{equation}
	\text{\_\_all\_\_} = \text{PUBLIC\_CONSTANT\_NAMES} \cup \text{PUBLIC\_FUNCTION\_NAMES} \cup \{\text{as\_dict}\}
\end{equation}

این تعریفِ صریح، ابزاری برای \lr{semantic versioning} فراهم می‌سازد: هر تغییری در \texttt{\_\_all\_\_} به‌عنوان یک تغییرِ ناسازگار (\lr{breaking change}) تلقی می‌شود و نیازمندِ افزایشِ نسخه‌ی اصلی است.

\subsection{ثابت‌های فیزیکیِ تکمیلی}

علاوه بر ثابت‌های تعریف‌شده در \texttt{constants\_definitions.py}، این ماژول دو ثابتِ فیزیکیِ تکمیلی را نیز تعریف می‌کند:

\begin{lstlisting}[language=Python]
	CONST_PLANCK = 6.62607015e-34          # J*s
	CONST_SPEED_OF_LIGHT = 2.99792458e8    # m/s
\end{lstlisting}

سرعتِ نور در خلأ، $c = 2.99792458 \times 10^8$ m/s، یکی از ثابت‌های دقیقِ \lr{SI 2019} است که در محاسبه‌ی طول‌موجِ به‌فرکانس ($\lambda = c/\nu$) و محاسبه‌ی زمانِ پروازِ فوتون در فیبر استفاده می‌شود. جالب توجه است که \texttt{CONST\_PLANCK} در این ماژول نیز تعریف شده — یک تعریفِ موازی با \texttt{constants\_definitions.py}. این تعریفِ موازی، اگرچه از نظرِ اصلِ \lr{DRY} (\lr{Don't Repeat Yourself}) ناقص به‌نظر می‌رسد، اما سازگاریِ با przeszłe را تضمین می‌کند: کدهای قدیمی که از \texttt{constants.CONST\_PLANCK} استفاده می‌کردند، همچنان کار خواهند کرد.

\section{تحلیل معماری و ساختار داده‌ها در \texttt{constants.py}}

\subsection{ساختارِ واردات}

ساختارِ وارداتِ این ماژول، الگوی فساد را به‌روشنی نشان می‌دهد:

\begin{lstlisting}[language=Python]
	from .constants_definitions import (
	CONST_BOLTZMANN, CONST_PLANCK, CONST_ELECTRON_CHARGE,
	DEFAULT_POISSON_TAIL_THRESHOLD, ENTROPY_PROB_CLAMP, EPS,
	LP_CONSTRAINT_VIOLATION_TOL, LP_SOLVER_METHODS, LPSolverMethod,
	MAX_SEED_INT_PCG64, MAX_SEED_INT_UINT32,
	MIN_SUCCESSFUL_Z1_BASIS_EVENTS_FOR_PHASE_EST,
	NUMERIC_ABS_TOL, NUMERIC_REL_TOL, PROB_SUM_TOL,
	PUBLIC_CONSTANT_NAMES, Y1_SAFE_THRESHOLD,
	)
	from .constants_helpers import (
	PUBLIC_FUNCTION_NAMES, clamp_probabilities, clamp_probability,
	db_to_linear, is_close, is_finite_non_negative, is_prob_vector,
	is_valid_probability, renormalize_probabilities,
	)
\end{lstlisting}

دو بلوکِ واردات، دو منبعِ مجزا را نشان می‌دهند: ثابت‌ها از \texttt{constants\_definitions} و توابع از \texttt{constants\_helpers}. هر دو بلوک تمامی نام‌های عمومی را به‌صورت صریح فهرست می‌کنند — نه از \lr{\texttt{*}} استفاده می‌کنند. این صراحت، اگرچه کد را طولانی‌تر می‌کند، اما وابستگی‌ها را به‌روشنی مستند می‌سازد و از وارداتِ تصادفیِ نام‌های خصوصی جلوگیری می‌کند.

\subsection{تعریفِ \texttt{\_\_all\_\_} پویا}

\begin{lstlisting}[language=Python]
	__version__ = "4.3.0"
	
	__all__ = [*PUBLIC_CONSTANT_NAMES, *PUBLIC_FUNCTION_NAMES, "as_dict"]
\end{lstlisting}

این تعریفِ پویا، \texttt{\_\_all\_\_} را از دو منبعِ مرجع می‌سازد. مزیتِ این رویکرد آن است که اگر در آینده یک ثابت جدید به \texttt{PUBLIC\_CONSTANT\_NAMES} اضافه شود، به‌صورت خودکار به \lr{API} عمومی نیز اضافه می‌شود — بدون آنکه نیاز به ویرایشِ دستیِ \texttt{constants.py} باشد. این \lr{Single Source of Truth}، بارِ نگهداری را به‌شدت کاهش می‌دهد.

\subsection{تابع \texttt{as\_dict()}}

تنها تابعِ تعریف‌شده در این ماژول، \texttt{as\_dict()} است:

\begin{lstlisting}[language=Python]
	def as_dict() -> dict[str, Any]:
	"""
	Return public constants from this module as an ordered dictionary copy.
	"""
	module_globals = globals()
	return {name: module_globals[name] for name in PUBLIC_CONSTANT_NAMES}
\end{lstlisting}

این تابع با استفاده از \texttt{globals()} تمامی ثابت‌های عمومی را به‌صورت یک دیکشنریِ مرتب برمی‌گرداند. ترتیبِ کلیدها به ترتیبِ \texttt{PUBLIC\_CONSTANT\_NAMES} است، که این تضمین را فراهم می‌سازد که خروجی در نسخه‌های مختلفِ فریم‌ورک پایدار بماند. این پایداری برای گزارش‌گیری و مقایسه‌ی آزمایش‌ها حیاتی است.

\subsection{عدمِ استفاده از \texttt{Final}}

جالب توجه است که ثابت‌های تعریف‌شده در این ماژول (\texttt{CONST\_PLANCK} و \texttt{CONST\_SPEED\_OF\_LIGHT}) با \texttt{Final} علامت‌گذاری نشده‌اند. این یک ناسازگاریِ سبکی با \texttt{constants\_definitions.py} است که در آینده باید اصلاح شود. در نسخه‌ی فعلی، این دو ثابت از نظرِ فنی قابل‌تغییر هستند، اما در عمل هرگز تغییر نمی‌کنند.

\section{پیاده‌سازی ریاضی و الگوریتمی در \texttt{constants.py}}

\subsection{پیاده‌سازیِ \texttt{as\_dict()}}

پیاده‌سازیِ این تابع، اگرچه کوتاه است، اما چند نکته‌ی ظریف دارد:

\begin{lstlisting}[language=Python]
	def as_dict() -> dict[str, Any]:
	module_globals = globals()
	return {name: module_globals[name] for name in PUBLIC_CONSTANT_NAMES}
\end{lstlisting}

نخست، استفاده از \texttt{globals()} به‌جای \texttt{getattr(module, name)} یک بهینه‌سازیِ کوچک است: \texttt{globals()} مستقیماً به دیکشنریِ ماژول دسترسی دارد، در حالی که \texttt{getattr} نیاز به جستجویِ صفت دارد.

دوم، استفاده از \lr{dictionary comprehension} به‌جای حلقه‌ی \texttt{for}، هم خوانایی را افزایش می‌دهد و هم از نظرِ بایت‌کد پایتون، کارآمدتر است.

سوم، ترتیبِ کلیدها در خروجی دقیقاً برابر با ترتیبِ \texttt{PUBLIC\_CONSTANT\_NAMES} است، چرا که دیکشنریِ پایتون (از نسخه‌ی ۳.۷ به بعد) ترتیبِ درج را حفظ می‌کند. این تضمینِ ترتیب، برای تولیدِ \lr{hash} پایدار و مقایسه‌ی \lr{JSON} حیاتی است.

\subsection{مقدارِ \lr{SI 2019} سرعتِ نور}

\begin{lstlisting}[language=Python]
	CONST_SPEED_OF_LIGHT = 2.99792458e8    # m/s
\end{lstlisting}

این مقدار، طبقِ \lr{SI 2019}، دقیق و غیرقابل‌تغییر است. فرمولِ ریاضی:
\begin{equation}
	c = 299\,792\,458 \, \text{m/s}
\end{equation}
این ثابت در محاسباتِ زیر استفاده می‌شود:
\begin{itemize}
	\item تبدیلِ طول‌موج به فرکانس: $\nu = c/\lambda$
	\item محاسبه‌ی زمانِ پرواز در فیبر: $t = L \cdot n / c$ که $n$ ضریبِ شکستِ فیبر است.
	\item محاسبه‌ی انرژیِ فوتون: $E = h\nu = hc/\lambda$
\end{itemize}

\subsection{تعریفِ موازیِ \texttt{CONST\_PLANCK}}

همان‌طور که پیش‌تر ذکر شد، \texttt{CONST\_PLANCK} هم در این ماژول و هم در \texttt{constants\_definitions.py} تعریف شده است. مقدارِ هر دو یکسان است:
\begin{equation}
	h = 6.62607015 \times 10^{-34} \, \text{J}\cdot\text{s}
\end{equation}

این تعریفِ موازی، در نگاهِ اول نقضِ اصلِ \lr{DRY} به‌نظر می‌رسد، اما در عمل یک \lr{compatibility shim} است. در گذشته، \texttt{CONST\_PLANCK} تنها در \texttt{constants.py} تعریف می‌شد. با ظهورِ \texttt{constants\_definitions.py}، تعریفِ اصلی به آن ماژول منتقل شد، اما برای حفظِ سازگاری با کدهای قدیمی، تعریفِ موازی در اینجا نگه داشته شد. در آینده، با حذفِ تدریجیِ کدهای قدیمی، این تعریفِ موازی نیز باید حذف شود.

\subsection{ترکیبِ \lr{API} از دو منبع}

ساختارِ \texttt{\_\_all\_\_} نشان می‌دهد که چگونه دو منبع در یک \lr{API} یکپارچه ترکیب می‌شوند:

\begin{lstlisting}[language=Python]
	__all__ = [*PUBLIC_CONSTANT_NAMES, *PUBLIC_FUNCTION_NAMES, "as_dict"]
\end{lstlisting}

این ترکیب، سه دسته از نام‌ها را در بر می‌گیرد:
\begin{enumerate}
	\item ثابت‌های عمومی (از \texttt{constants\_definitions})
	\item توابعِ کمکیِ عمومی (از \texttt{constants\_helpers})
	\item تابعِ \texttt{as\_dict} (تعریف‌شده در این ماژول)
\end{enumerate}

این الگو، یک \lr{API} یکپارچه را ارائه می‌دهد که در آن تمامی منابعِ موردِ نیاز کاربر در یک نقطه در دسترس است.

\section{ارسال و دریافت با سایر ماژول‌ها در \texttt{constants.py}}

\subsection{وابستگی‌های ورودی}

این ماژول تنها به دو ماژولِ داخلی وابسته است:

\begin{itemize}
	\item \texttt{constants\_definitions.py}: برای تمامی ثابت‌ها و نوع‌های تایپی.
	\item \texttt{constants\_helpers.py}: برای تمامی توابعِ کمکی.
\end{itemize}

وابستگیِ خارجیِ تنها، کتابخانه‌ی استانداردِ \texttt{typing} (برای نوعِ \texttt{Any}) است. این استقلالِ کم، این ماژول را به یک نقطه‌ی سبک و سریع تبدیل می‌کند.

\subsection{نقاطِ مصرف}

این ماژول توسط اکثرِ ماژول‌های فریم‌ورک استفاده می‌شود:

\begin{itemize}
	\item \textbf{ماژولِ \texttt{params.py}:} از \texttt{LP\_SOLVER\_METHODS} و \texttt{DEFAULT\_POISSON\_TAIL\_THRESHOLD} استفاده می‌کند.
	\item \textbf{ماژولِ \texttt{sources.py}:} از \texttt{clamp\_probabilities}، \texttt{is\_valid\_probability} و ثابت‌های تلورانس استفاده می‌کند.
	\item \textbf{ماژولِ \texttt{detectors.py}:} از \texttt{db\_to\_linear} و \texttt{is\_close} استفاده می‌کند.
	\item \textbf{ماژولِ \texttt{channel.py}:} از \texttt{db\_to\_linear} و ثابت‌های فیزیکی استفاده می‌کند.
	\item \textbf{اسکریپت‌های آزمایشی:} از \texttt{as\_dict()} برای گزارش‌گیری و ثبتِ پیکربندی استفاده می‌کنند.
\end{itemize}

\subsection{قراردادِ نسخه‌گذاری}

متغیر \texttt{\_\_version\_\_ = "4.3.0"} یک قراردادِ مهم را ارائه می‌دهد. این نسخه‌بندی از \lr{Semantic Versioning} پیروی می‌کند:
\begin{itemize}
	\item \lr{Major (4)}: تغییراتِ ناسازگارِ \lr{API}.
	\item \lr{Minor (3)}: افزودنِ قابلیت‌های جدید به‌صورت سازگار.
	\item \lr{Patch (0)}: اصلاحِ خطاها و بهبودِ عملکرد.
\end{itemize}

ماژول‌های پایین‌دستی می‌توانند با بررسیِ این نسخه، سازگاریِ خود را تضمین کنند:
\begin{equation}
	\text{compatible} \iff \text{major}_{\text{required}} = \text{major}_{\text{actual}}
\end{equation}

این قرارداد، در آینده‌ی فریم‌ورک که ممکن است \lr{API} تغییراتِ اساسی تجربه کند، حیاتی خواهد بود.

\subsection{جمع‌بندیِ معماریِ چهار ماژول}

چهار ماژولِ بررسی‌شده در این فصل، یک معماریِ لایه‌ایِ تمیز را تشکیل می‌دهند. در پایینِ لایه‌بندی، \texttt{constants\_definitions.py} به‌عنوان منبعِ نهاییِ ثابت‌ها قرار دارد. در لایه‌ی بعدی، \texttt{constants\_helpers.py} توابعِ کمکیِ عددی را بر پایه‌ی این ثابت‌ها ارائه می‌دهد. در لایه‌ی فساد، \texttt{constants.py} یک \lr{API} یکپارچه از این دو ماژول فراهم می‌سازد. در بالاترین لایه، \texttt{params.py} از این \lr{API} برای ساختِ شیءِ پیکربندیِ مرکزیِ \texttt{QKDParams} استفاده می‌کند.

این جداسازیِ لایه‌ها، چند مزیتِ کلیدی دارد. نخست، \lr{Single Source of Truth} برای هر مفهوم: ثابت‌ها تنها در یک نقطه تعریف می‌شوند، توابعِ کمکی تنها در یک نقطه پیاده‌سازی می‌شوند، و پیکربندیِ مرکزی تنها از یک کلاس نشأت می‌گیرد. دوم، قابلیتِ آزمونپذیریِ بالا: هر لایه می‌تواند به‌صورت مستقل آزمون شود. سوم، انعطاف‌پذیریِ بازسازی: تغییراتِ داخلی در یک لایه، بر لایه‌های بالاتر تأثیر نمی‌گذارد. این معماری، نمونه‌ای از بهترین شیوه‌های مهندسیِ نرم‌افزار در پروژه‌هایِ علمیِ بزرگ است.

	
	\chapter{فیزیک و شبیه‌سازی کانال}
	\label{ch:physics}
	% ============================================================
%  QKD Simulation Framework — Chapter: Channel, Physics, Noise
%  Authors: Documentation generated for academic thesis
%  Encoding: UTF-8, XePersian-compatible
% ============================================================

% --- تنظیمات رفع سرریز جعبه‌افقی (Overfull \hbox) ---
\sloppy
\emergencystretch=3em
\hbadness=10000
\vbadness=10000

% --- پیکربندی بسته‌ی listings برای نمایش صحیح کد پایتون در متن RTL ---
\lstset{
	basicstyle=\ttfamily\small,
	keywordstyle=\color{blue}\bfseries,
	stringstyle=\color{teal},
	commentstyle=\color{gray!60!black}\itshape,
	numberstyle=\tiny\color{gray},
	showstringspaces=false,
	breaklines=true,
	breakatwhitespace=true,
	frame=single,
	rulecolor=\color{gray!40},
	backgroundcolor=\color{gray!5},
	language=Python,
	extendedchars=true,
	inputencoding=utf8,
	tabsize=4,
	keepspaces=true,
	columns=fullflexible,
	basewidth=0.55em,
	aboveskip=1em,
	belowskip=1em,
	xleftmargin=1em,
	xrightmargin=1em,
	escapeinside={(*@}{@*)},
	literate=
	{→}{{$\rightarrow$}}1
	{∂}{{$\partial$}}1
	{α}{{$\alpha$}}1
	{β}{{$\beta$}}1
	{γ}{{$\gamma$}}1
	{δ}{{$\delta$}}1
	{ε}{{$\varepsilon$}}1
	{θ}{{$\theta$}}1
	{λ}{{$\lambda$}}1
	{μ}{{$\mu$}}1
	{ν}{{$\nu$}}1
	{π}{{$\pi$}}1
	{σ}{{$\sigma$}}1
	{η}{{$\eta$}}1
	{≤}{{$\leq$}}1
	{≥}{{$\geq$}}1
	{≠}{{$\neq$}}1
	{∈}{{$\in$}}1
	{∞}{{$\infty$}}1
	{Σ}{{$\Sigma$}}1
	{Δ}{{$\Delta$}}1
	{Ω}{{$\Omega$}}1
}

\providecommand{\code}[1]{\texttt{\LR{#1}}}

\section{نمای کلی و هدف ماژول}
\label{sec:channel-overview}

این فصل به بررسی سه لایه‌ی به‌هم‌پیوسته در مدل‌سازی کانال کوانتومی در فریم‌ورک شبیه‌سازی \lr{QKD} می‌پردازد: (الف) فایل \lr{\texttt{channel.py}} که حاوی کلاس اصلی \lr{\texttt{FiberChannel}} و کلاس‌های کمکیِ آن است؛ (ب) فایل \lr{\texttt{\_physics.py}} که توابع محاسبات فیزیکی مشتق‌شده (مانند محاسبه‌ی \lr{Transmittance}، کسر تلفات، و ترکیب آبشاری کانال‌ها) را از بدنه‌ی کلاس جدا می‌سازد؛ و (ج) فایل \lr{\texttt{noise\_models.py}} که مدل‌های نویز الکتریکی سمت منبع (مانند نویز حرارتی \lr{Johnson-Nyquist} و نویز شات \lr{Schottky}) را به‌صورت یک \lr{dataclass} مستقل پیاده‌سازی می‌کند. این سه ماژول در کنار یکدیگر، چرخه‌ی کاملِ تبدیلِ پارامترهای فیزیکیِ فیبر (طول $L$ و ضریب تلفات $\alpha$) به کمیت‌های قابل‌استفاده در محاسباتِ نرخ کلید \lr{QKD} را فراهم می‌سازند.

هدف بنیادین این مجموعه، ارائه‌ی یک مدلِ تغییریافته‌ناپذیر \lr{(immutable)}، قابل سریال‌سازی \lr{(serializable)} و از نظر عددی پایدار از کانال فیبر نوری است که بر پایه‌ی قانون \lr{Beer-Lambert} بنا شده است. در این قانون، ضریب انتقال کل کانال به‌صورت تابعی نمایی از حاصل‌ضربِ طول فیبر در ضریب تلفات تعریف می‌شود. این رابطه که در ظاهر ساده به‌نظر می‌رسد، در عمل با چالش‌های متعددی روبه‌رو است: سرریز عددی \lr{(underflow)} برای فواصل طولانی، ظهور مقادیر زیرنرمال \lr{(subnormal)} که دقت محاسبات آنتروپی را تخریب می‌کنند، عدم تقارنِ \lr{round-trip} میان $\alpha \cdot L$ و $L_{\text{total}}$ به‌دلیل خطای گردکردنِ \lr{IEEE-754}، و نیاز به سازوکارهایی برای تشخیص پیکربندی‌های فیزیکی مشکوک \lr{(physical suspicion)} بدون توقف کاملِ شبیه‌سازی. این فصل به‌تفصیل نشان می‌دهد که چگونه هر یک از این چالش‌ها در سطح کدِ پایتون آدرس‌دهی شده‌اند.

طراحیِ این ماژول بر چهار اصل بنیادین استوار است. نخست، \lr{``fail-fast''}: هرگونه ناسازگاری در زمان بارگذاری پیکربندی باید بلافاصله با یک استثنای صریح افشا گردد تا از انتشار خطا به لایه‌های پایین‌تر جلوگیری شود. دوم، \lr{``strict mode'}': کاربر می‌تواند با فعال‌سازی این حالت، رفتار ماژول را از هشدار دادن \lr{(warning)} به پرتاب استثنا \lr{(raising)} تغییر دهد تا در خطوط لوله‌ی تولید \lr{(production pipelines)} هیچ پیکربندیِ مشکوکی پذیرفته نشود. سوم، \lr{``reproducibility'}': فراداده‌ی زمان ساخت \lr{(construction-time metadata)}، حالت \lr{strict} در زمان ساخت، و مسیرِ ساخت \lr{(construction path)} در هر نمونه ثبت می‌شوند تا بازتولید آزمایش‌ها در محیط‌های مختلف تضمین گردد. چهارم، \lr{``type-safety'}': استفاده‌ی گسترده از \lr{\texttt{TypeVar}}، \lr{\texttt{Generic}}، \lr{\texttt{Protocol}} و \lr{\texttt{Final}} اطمینان می‌دهد که تحلیل‌گرهای ایستا مانند \lr{mypy} و \lr{pyright} می‌توانند بسیاری از خطاهای نوعی را پیش از اجرای کد شناسایی کنند.

در سطح کلان، فایل \lr{\texttt{channel.py}} با بیش از ۴۵۰۰ خط کد، سنگین‌ترین و پیچیده‌ترین جزء این مجموعه است. این فایل نه‌تنها کلاس \lr{\texttt{FiberChannel}} را تعریف می‌کند، بلکه زیرکلاسِ \lr{\texttt{AttenuationOnlyFiberChannel}} را به‌عنوان نشانگر نوعی برای کانال‌های صرفاً تلفاتی، توابع برداریِ \lr{\texttt{transmittance\_array}} و \lr{\texttt{transmittance\_scalar}} را برای حلقه‌های داغِ \lr{Monte Carlo}، کلاس \lr{\texttt{ChannelSimParams}} را به‌عنوان بسته‌ی یخ‌زده‌ی کمیت‌های مشتق‌شده، و سازوکار \lr{strict mode} را به‌صورت یک \lr{\texttt{ContextVar}} پیاده‌سازی می‌کند. فایل \lr{\texttt{\_physics.py}} با حدود ۶۲۰ خط، توابعی را که در نسخه‌های پیشین به‌عنوان متدهای کلاس تعریف می‌شدند، به توابعی مستقل تبدیل کرده است؛ این تفکیکِ نگرانی‌ها \lr{(separation of concerns)} امکان تست واحدِ آسان‌تر و استفاده‌ی مجددِ منطق را فراهم می‌سازد. فایل \lr{\texttt{noise\_models.py}} با تنها ۸۰ خط، یک ماژول فشرده اما مهم برای مدل‌سازی نویز الکتریکیِ فوتودیودِ نظارتی در سمت فرستنده است که در سیستم‌های \lr{SCM-QKD} و \lr{WDM} برای تخمین نویز اضافی به‌کار می‌رود.

\section{مبانی تئوری و فیزیکی}
\label{sec:channel-theory}

\subsection{قانون Beer-Lambert و تلفات فیبر نوری}

در فیبر نوری سیلیکا، تضعیفِ سیگنال نوری ناشی از سه مکانیزم فیزیکی اصلی است: پراکندگی \lr{Rayleigh} ناشی از نوسانات چگالی در مقیاسِ کوچک‌تر از طول موج، جذب ناشی از ناخالصی‌های یونی (به‌ویژه یون‌های OH$^-$ در طول موج $1383\,\text{nm}$)، و پراکندگی \lr{Mie} ناشی از ناخالصی‌های با مقیاسِ بزرگ‌تر. در طول موج‌های پنجره‌ی مخابراتی ($1310\,\text{nm}$ و $1550\,\text{nm}$)، پراکندگی \lr{Rayleigh} مکانیزم غالب است و منجر به ضریب تلفاتِ تقریباً ثابت به ازای واحد طول می‌گردد. این رفتار خطی با قانون \lr{Beer-Lambert} مدل می‌شود:
\begin{equation}
	T(L) = \frac{P(L)}{P(0)} = 10^{-\alpha\,L/10},
	\label{eq:beer-lambert}
\end{equation}
که در آن $T$ ضریب انتقال \lr{(transmittance)}، $P(L)$ توان نوری خروجی در طول $L$، $P(0)$ توان ورودی، $\alpha$ ضریب تلفات به واحد \lr{dB/km} و $L$ طول فیبر به کیلومتر است. تلفات کل به دسی‌بل از رابطه‌ی خطی زیر به دست می‌آید:
\begin{equation}
	L_{\text{total}}\,[\text{dB}] = \alpha\,[\text{dB/km}] \times L\,[\text{km}],
	\label{eq:total-loss-linear}
\end{equation}
و در نتیجه رابطه‌ی نماییِ معادل برای ضریب انتقال به‌صورت $T = 10^{-L_{\text{total}}/10}$ نوشته می‌شود. در طول موج $1550\,\text{nm}$ مقدار معمولِ $\alpha \approx 0.2\,\text{dB/km}$ است که به فاصله‌ی تقریبی $50\,\text{km}$ نرخ کلید غیرصفری در \lr{BB84-Decoy} با آشکارسازهای \lr{SNSPD} اجازه می‌دهد.

\subsection{تفکیک $\eta_{\text{channel}}$ از بهره‌ی کلی آشکارساز}

در یک حلقه‌ی بودجه‌ی پیوندِ \lr{QKD}، احتمال کلی آشکارسازیِ یک فوتونِ منتشرشده از سمت آلیس به‌صورت زیر تجزیه می‌شود:
\begin{equation}
	\eta = \eta_{\text{channel}} \cdot \eta_{\text{detector}} \cdot \eta_{\text{coupling}},
	\label{eq:link-budget}
\end{equation}
که در آن $\eta_{\text{channel}} = T$ همان ضریب انتقالِ فیبر است، $\eta_{\text{detector}}$ بازدهی تک‌آشکارساز (معمولاً $0.1$ تا $0.3$ برای \lr{SNSPD} و $0.05$ تا $0.15$ برای \lr{SPAD})، و $\eta_{\text{coupling}}$ تلفات اتصالات، اتصال‌دهنده‌ها و جوش‌های فیبر (معمولاً $0.5$ تا $0.9$) است. یکی از خطاهای متداول در ادبیات شبیه‌سازی، استفاده‌ی مستقیم $T$ به‌جای $\eta$ در فرمول بهره‌ی \lr{Decoy-State}:
\begin{equation}
	Q_\mu = \sum_{n=0}^{\infty} Y_n \cdot \frac{\mu^n\,e^{-\mu}}{n!}, \qquad Y_n = 1 - (1-\eta)^n + Y_0,
	\label{eq:decoy-gain}
\end{equation}
است که می‌تواند منجر به جابهجاییِ منحنی نرخ کلید تا $50$ تا $100\,\text{km}$ گردد. به همین دلیل، کلاس \lr{\texttt{FiberChannel}} صراحتاً تنها $\eta_{\text{channel}}$ را ارائه می‌دهد و تابع کمکیِ \lr{\texttt{link\_efficiency}} ترکیب سه‌گانه‌ی رابطه‌ی \eqref{eq:link-budget} را در یک نقطه‌ی متمرکز انجام می‌دهد تا اطمینان حاصل شود که تمام محاسبات (تحلیل متقاطع \lr{Rogers et al. (2007)} و فرمول بهره) از یک ترکیب‌بندیِ واحد استفاده می‌کنند.

\subsection{پدیده‌ی زیرنرمال بودن و سرریز عددی}

در فواصل بسیار طولانی (مثلاً $L = 5000\,\text{km}$ با $\alpha = 0.2$)، مقدار $L_{\text{total}} = 1000\,\text{dB}$ و $T = 10^{-100}$ است که به‌راحتی در محدوده‌ی نرمالِ \lr{double} نمایش داده می‌شود. اما برای $L = 50000\,\text{km}$ با همان $\alpha$، مقدار $T = 10^{-1000}$ است که بسیار کوچک‌تر از کوچک‌ترین عدد نرمالِ مثبت در \lr{IEEE-754} ($\approx 2.225 \times 10^{-308}$) می‌باشد. در این حالت دو سناریو رخ می‌دهد: (الف) اگر مقدار به‌صورتِ تدریجی‌کاهش‌یافته \lr{(gradual underflow)} محاسبه شود، یک عدد زیرنرمال با تنها $1$ تا $2$ بیت معتبرِ مانتیس تولید می‌شود؛ (ب) اگر زیرنرمال‌ها در محیطِ اجرا غیرفعال باشند \lr{(FTZ/DAZ modes)}، مقدار به صفر دقیق تبدیل می‌شود.

هر دو حالت برای محاسبات \lr{QKD} فاجعه‌بارند: یک عدد زیرنرمال در فرمولِ بهره‌ی $Y_n = 1 - (1-\eta)^n$ به دلیل از دست رفتن بیت‌های معتبر، منجر به نتایجِ کاملاً اشتباه می‌شود؛ و یک مقدار صفر دقیق، در لگاریتمِ آنتروپیِ شانون $H_2(p) = -p\log_2 p - (1-p)\log_2(1-p)$ منجر به $\log(0) = -\infty$ و در نتیجه \lr{NaN} می‌گردد. به همین دلیل، ماژول \lr{\texttt{channel.py}} سه سطح تشخیص را پیاده‌سازی می‌کند: آستانه‌ی زیرنرمال \lr{\texttt{TRANSMITTANCE\_SUBNORMAL\_THRESHOLD}} (برابرِ \lr{\texttt{sys.float\_info.min}})، آستانه‌ی سرریزِ مطلق \lr{\texttt{\_TRANSMITTANCE\_UNDERFLOW\_ABS\_TOL}} (همان مقدار)، و آستانه‌ی قراردادِ $L,T$ به مقدار $10^{-12}$. در حالتِ \lr{strict mode}، هر یک از این شرایط به پرتابِ استثنای \lr{\texttt{ParameterValidationError}} منجر می‌شوند.

\subsection{نویز حرارتی Johnson-Nyquist}

مقاومتِ اهمی در دمای غیرصفر، به‌دلیلِ حرکتِ حرارتیِ تصادفیِ الکترون‌ها، یک ولتاژِ نویزِ تصادفی تولید می‌کند که به افتخارِ کاشفان آن، نویز \lr{Johnson-Nyquist} نامیده می‌شود. توانِ طیفیِ این نویز در محدوده‌ی فرکانسیِ پایین (نسبت به دمای پلاسمایی) مسطح است و از رابطه‌ی زیر پیروی می‌کند:
\begin{equation}
	\langle v_n^2 \rangle = 4\,k_B\,T\,R\,\Delta f,
	\label{eq:johnson-noise}
\end{equation}
که در آن $k_B = 1.380649 \times 10^{-23}\,\text{J/K}$ ثابت بولتزمن، $T$ دمای مطلق، $R$ مقاومت بار \lr{(load resistance)} و $\Delta f$ پهنای‌باندِ الکتریکی است. انحراف معیارِ ولتاژ نویزِ حرارتی از ریشه‌گیریِ رابطه‌ی فوق به دست می‌آید:
\begin{equation}
	\sigma_v^{\text{(thermal)}} = \sqrt{4\,k_B\,T\,R\,\Delta f}.
	\label{eq:thermal-std}
\end{equation}
در دمای اتاق $T = 298.15\,\text{K}$، با $R = 50\,\Omega$ و $\Delta f = 1\,\text{GHz}$، این مقدار تقریباً $28.9\,\mu\text{V}$ است.

\subsection{نویز شات Schottky}

عبورِ جریان الکتریکی از یک فراسو، به‌دلیلِ گسسته‌بودنِ بارِ الکترون، با نویزِ شات همراه است. توانِ طیفیِ جریانِ نویزِ شات از رابطه‌ی \lr{Schottky} به دست می‌آید:
\begin{equation}
	\langle i_n^2 \rangle = 2\,e\,I_{\text{DC}}\,\Delta f,
	\label{eq:shot-noise}
\end{equation}
که در آن $e = 1.60217663 \times 10^{-19}\,\text{C}$ بار الکترون، $I_{\text{DC}}$ جریانِ مستمرِ فوتودیودِ نظارتی، و $\Delta f$ پهنای‌باند است. با عبور از مقاومتِ بار، این نویز به ولتاژ تبدیل می‌شود:
\begin{equation}
	\sigma_v^{\text{(shot)}} = R \cdot \sqrt{2\,e\,I_{\text{DC}}\,\Delta f}.
	\label{eq:shot-std}
\end{equation}
چون نویزِ حرارتی و نویزِ شات از فرایندهای فیزیکیِ مستقل هستند، ترکیبِ آن‌ها به‌صورتِ جمعِ توان‌ها انجام می‌شود:
\begin{equation}
	\sigma_v^{\text{(total)}} = \sqrt{(\sigma_v^{\text{(thermal)}})^2 + (\sigma_v^{\text{(shot)}})^2}.
	\label{eq:total-noise}
\end{equation}
این رابطه در متدِ \lr{\texttt{total\_voltage\_std}} کلاس \lr{\texttt{ElectricalNoiseConfig}} پیاده‌سازی شده است.

\subsection{کسر تلفات تک‌کیلومتری و تقریب خطی}

کسرِ توانِ نوریِ از دست‌رفته در دقیقاً یک کیلومتر فیبر از رابطه‌ی زیر به دست می‌آید:
\begin{equation}
	\ell = 1 - T(1\,\text{km}) = 1 - 10^{-\alpha/10}.
	\label{eq:loss-fraction}
\end{equation}
برای $\alpha = 0.2\,\text{dB/km}$، این مقدار $\ell \approx 0.0451$ است. یک اشتباهِ متداول در ادبیاتِ مهندسی، ضربِ این کسر در طول فیبر برای تخمینِ تلفاتِ کل است؛ اما این کار نادرست است زیرا کسر تلفات به‌صورتِ نمایی با فاصله ترکیب می‌شود:
\begin{equation}
	\ell_{\text{total}}(L) = 1 - (1-\ell)^L \neq \ell \cdot L.
	\label{eq:nonlinear-composition}
\end{equation}
برای $\alpha = 0.2$ و $L = 100\,\text{km}$، تقریب خطیِ $\ell \cdot L \approx 4.51$ است که کاملاً غیرفیزیکی (بزرگ‌تر از ۱) می‌باشد، در حالی که مقدار دقیق $\ell_{\text{total}} \approx 0.990$ است. به همین دلیل، ماژول \lr{\texttt{channel.py}} متدِ قدیمیِ \lr{\texttt{loss\_fraction\_per\_km}} را منسوخ کرده و نامِ جدیدِ \lr{\texttt{single\_km\_power\_loss\_fraction}} را برای آن انتخاب کرده است تا ماهیتِ غیرخطیِ این کمیت به‌طور صریح در نام بازتاب یابد.

برای $\alpha$‌های بسیار کوچک (نظیر $\alpha = 10^{-6}\,\text{dB/km}$ در شبیه‌سازی‌های فرضی)، فرمول مستقیمِ $1 - 10^{-\alpha/10}$ دچار لغوِ فاجعه‌بار \lr{(catastrophic cancellation)} می‌شود: مقدار $10^{-10^{-7}} \approx 0.99999976$ تنها $6$ رقم معتبر دارد و تفریق از $1.0$ منجر به از دست رفتنِ بیشترِ بیت‌ها می‌شود. برای جلوگیری از این ناپایداری، از تابعِ \lr{\texttt{math.expm1}} که به‌صورتِ تخصصی برای محاسبه‌ی دقیقِ $e^x - 1$ برای $x$‌های کوچک طراحی شده، استفاده می‌گردد:
\begin{equation}
	\ell = -\mathrm{expm1}\!\left(-\frac{\alpha \cdot \ln 10}{10}\right).
	\label{eq:expm1-stable}
\end{equation}
این فرمول برای هر مقدار $\alpha$، تمامی بیت‌های معتبر را حفظ می‌کند.

تقریبِ خطیِ مرتبه‌ی اولِ تیلور از رابطه‌ی \eqref{eq:loss-fraction} به‌صورت زیر است:
\begin{equation}
	\ell_{\text{linear}} \approx \frac{\alpha \cdot \ln 10}{10} \approx 0.2303 \cdot \alpha,
	\label{eq:linear-approx}
\end{equation}
که برای $\alpha = 0.2$ مقدار $\ell_{\text{linear}} \approx 0.0461$ را می‌دهد. این تقریب واقعاً خطی است (به $L$ وابسته نیست) و با فاصله به‌صورتِ خطی ترکیب می‌شود، اما تنها برای $\alpha \cdot L$‌های کوچک معتبر است. برای $\alpha > 4.34\,\text{dB/km}$ (که در آن $10/\ln 10 \approx 4.34$) این تقریب از $1$ فراتر می‌رود و فیزیکی نمی‌باشد.

\subsection{ترکیب آبشاری کانال‌ها}

در یک زنجیره‌ی چندبُعدی از فیبرها با پارامترهای متفاوت، ضریب انتقالِ کل از حاصل‌ضربِ ضرایبِ انتقالِ هر بخش به دست می‌آید:
\begin{equation}
	T_{\text{total}} = \prod_{i=1}^{N} T_i = \prod_{i=1}^{N} 10^{-\alpha_i L_i / 10} = 10^{-\sum_i \alpha_i L_i / 10},
	\label{eq:cascade-product}
\end{equation}
که معادل با جمعِ تلفات در واحد دسی‌بل است: $L_{\text{total}} = \sum_i \alpha_i L_i$. اگر بخواهیم یک کانالِ معادلِ تک‌بُعدی با همان تلفاتِ کل و طولِ کل تعریف کنیم، ضریب تلفاتِ میانگینِ آن به‌صورت زیر محاسبه می‌شود:
\begin{equation}
	\alpha_{\text{avg}} = \frac{L_{\text{total}}}{L_{\text{total\_distance}}} = \frac{\sum_i \alpha_i L_i}{\sum_i L_i}.
	\label{eq:cascade-alpha-avg}
\end{equation}
این یک میانگینِ حسابیِ وزن‌دار است که به هیچ نوع فیبرِ استانداردی در هیچ طول موجی اشاره ندارد. به‌عنوان مثال، ترکیبِ یک فیبر $100\,\text{km}$ با $\alpha = 0.2$ و یک فیبر $10\,\text{km}$ با $\alpha = 3.0$ به $\alpha_{\text{avg}} \approx 0.47\,\text{dB/km}$ می‌انجامد که در هیچ طول موجی فیزیکی نیست. به همین دلیل، متد \lr{\texttt{cascade}} در کد به‌صراحت هشدار می‌دهد که کانالِ معادل، یک \lr{``composite loss-equivalent''} است و نباید در مقایسه‌های ادبی به‌عنوان فیبر واقعی به‌کار رود.

\subsection{پاشش کروماتیک و محدودیت‌های مدل}

پاشش کروماتیک \lr{(Chromatic Dispersion)} باعث پهن‌شدنِ پالس نوری در طول فیبر می‌شود و در صورتی که پهنای زمانیِ پالس از پنجره‌ی زمانیِ آشکارساز فراتر رود، به از بین رفتنِ اطلاعات می‌انجامد. ضریب پاشش $D(\lambda)$ به واحد \lr{ps/(nm· km)} بیان می‌شود و پهن‌شدنِ زمانی از رابطه‌ی زیر تقریب زده می‌شود:
\begin{equation}
	\Delta\tau \approx D(\lambda) \cdot L \cdot \Delta\lambda,
	\label{eq:chromatic-dispersion}
\end{equation}
که در آن $\Delta\lambda$ پهنای طیفیِ منبع است. در فیبر استانداردِ تک‌حالته در $1550\,\text{nm}$، $D \approx 17\,\text{ps/(nm· km)}$ است که برای منبعِ لیزری با $\Delta\lambda = 0.1\,\text{nm}$ در $100\,\text{km}$، پهن‌شدنِ $170\,\text{ps}$ تولید می‌کند. این پدیده در مدل \lr{\texttt{FiberChannel}} پیاده‌سازی نشده است (زیرا بر ضریب انتقال اثر مستقیم ندارد) اما در \lr{\texttt{ChannelSimParams}} به‌عنوان فراداده‌ی عددی ثبت می‌گردد تا توسط موتورِ شبیه‌سازیِ آشکارساز برای مدل‌سازیِ پنجره‌ی زمانیِ قابل‌استفاده باشد.

\section{تحلیل معماری و ساختار داده‌ها}
\label{sec:channel-architecture}

\subsection{کلاس \lr{\texttt{FiberChannel}} به‌عنوان \lr{Frozen Dataclass}}

کلاسِ \lr{\texttt{FiberChannel}} با دو decorator اصلی تعریف می‌شود:
\begin{LTR}
	\begin{lstlisting}[language=Python]
		@dataclass(frozen=True, slots=True)
		class FiberChannel:
		"""Models a simple optical fiber channel defined by attenuation."""
		attenuation: AttenuationConfig
		_construction_path: str = field(default="direct", repr=False)
		_total_loss_db_override: Optional[float] = field(
		init=False, default=None, repr=False
		)
		_transmittance: Any = field(init=False, repr=False, default=_UNSET)
		_strict_mode_at_construction: bool = field(init=False, repr=False, default=False)
		wavelength_nm: Optional[float] = field(default=None, repr=False)
		_metadata_extra: Mapping[str, Any] = field(
		init=False, repr=False, default_factory=lambda: _EMPTY_METADATA_EXTRA
		)
		__hash__ = None  # type: ignore[assignment,method-assign]
	\end{lstlisting}
\end{LTR}
پرچمِ \lr{\texttt{frozen=True}} کلاس را تغییرناپذیر می‌سازد؛ پس از ساختِ نمونه، هیچ فیلدی نمی‌تواند با تخصیصِ معمولی تغییر کند. این ویژگی برای تضمینِ \lr{reproducibility} حیاتی است. هرگاه تغییر ضروری باشد (مانند تنظیمِ \lr{\texttt{\_total\_loss\_db\_override}} در \lr{\texttt{from\_total\_loss}})، باید از \lr{\texttt{object.\_\_setattr\_\_(...)}} استفاده شود که تنها راه قانونیِ دور زدنِ محدودیتِ \lr{frozen} است. پرچمِ \lr{\texttt{slots=True}} از تخصیصِ \lr{\_\_dict\_\_} جلوگیری کرده و به‌جای آن فضا را از طریقِ \lr{\_\_slots\_\_} تخصیص می‌دهد؛ این بهینه‌سازی مصرفِ حافظه را تا $40\%$ کاهش داده و سرعتِ دسترسی به فیلدها را افزایش می‌دهد.

\subsection{فیلدهای کلاس و نقش آن‌ها}

\textbf{فیلدِ \lr{\texttt{attenuation}} (نوع: \lr{\texttt{AttenuationConfig}}):} تنها فیلدِ اجباریِ کاربر. این فیلد، یک \lr{dataclass} دیگر از ماژولِ \lr{\texttt{qkd.datatypes}} است که دو فیلدِ \lr{\texttt{fiber\_length}} (به \lr{km}) و \lr{\texttt{attenuation\_coefficient}} (به \lr{dB/km}) را در بر می‌گیرد. قراردادِ \lr{\texttt{AttenuationConfig}} آن است که \lr{\texttt{total\_loss\_db = fiber\_length * attenuation\_coefficient}} باید برقرار باشد.

\textbf{فیلدِ \lr{\texttt{\_construction\_path}} (نوع: \lr{\texttt{str}}، پیش‌فرض: \lr{\texttt{"direct"}}):} این فیلد مسیرِ ساختِ نمونه را ثبت می‌کند تا در پیام‌های هشدار و فراداده‌ی بازتولید، ریشه‌ی خطا قابل ردیابی باشد. مقادیرِ ممکن عبارت‌اند از: \lr{\texttt{"direct"}} (ساخت مستقیم)، \lr{\texttt{"from\_total\_loss"}}، \lr{\texttt{"from\_tuple"}}، \lr{\texttt{"from\_config\_dict"}}، \lr{\texttt{"with\_distance"}}، \lr{\texttt{"with\_loss"}}، \lr{\texttt{"fast\_construct"}} و \lr{\texttt{"cascade"}}.

\textbf{فیلدِ \lr{\texttt{\_total\_loss\_db\_override}} (نوع: \lr{\texttt{Optional[float]}}، \lr{\texttt{init=False}}):} این فیلد برای حفظِ دقتِ \lr{round-trip} در کانال‌های ساخته‌شده با \lr{\texttt{from\_total\_loss}} طراحی شده است. هنگامی که کاربر $L_{\text{total}}$ را به‌طور صریح ارائه می‌دهد، $\alpha = L_{\text{total}} / L$ محاسبه می‌شود که به‌دلیلِ گردکردنِ \lr{IEEE-754} ممکن است $\alpha \cdot L$ با $L_{\text{total}}$ اصلی در سطحِ \lr{ULP} متفاوت باشد. برای جلوگیری از این ناسازگاری، مقدارِ اصلیِ $L_{\text{total}}$ در این فیلد ذخیره شده و در محاسباتِ \lr{\texttt{total\_loss\_db}} و \lr{\texttt{transmittance}} به‌جای $\alpha \cdot L$ به‌کار می‌رود.

\textbf{فیلدِ \lr{\texttt{\_transmittance}} (نوع: \lr{\texttt{Any}}، \lr{\texttt{init=False}}، پیش‌فرض: \lr{\texttt{\_UNSET}}):} این فیلد مقدارِ از پیش محاسبه‌شده‌ی ضریب انتقال را نگه می‌دارد. استفاده از نوعِ \lr{\texttt{Any}} و پیش‌فرضِ \lr{\texttt{\_UNSET}} (یک \lr{singleton} از کلاسِ \lr{\texttt{\_Unset}}) به ما اجازه می‌دهد بینِ حالتِ ``هنوز محاسبه نشده'' و ``مقدارِ عددیِ معتبر'' تمایز قائل شویم؛ این امر در \lr{\texttt{\_\_repr\_\_}} و \lr{\texttt{\_\_setstate\_\_}} حیاتی است.

\textbf{فیلدِ \lr{\texttt{\_strict\_mode\_at\_construction}} (نوع: \lr{\texttt{bool}}، \lr{\texttt{init=False}}):} این فیلد مقدارِ \lr{strict mode} را در زمانِ ساختِ نمونه ثبت می‌کند (نه در زمانِ سریال‌سازی). این تمایز مهم است زیرا اگر کاربر بینِ ساخت و سریال‌سازی، حالت \lr{strict} را تغییر دهد، فراداده‌ی بازتولید نباید جعلی شود.

\textbf{فیلدِ \lr{\texttt{wavelength\_nm}} (نوع: \lr{\texttt{Optional[float]}}):} طول موجِ عملیاتی به نانومتر. هنگامی که تنظیم شود، یک بررسیِ سازگاریِ فیزیکی بر اساسِ جدولِ \lr{\texttt{TYPICAL\_WAVELENGTH\_ALPHA\_RANGES}} انجام می‌شود تا مقادیرِ $\alpha$ نامتعارف شناسایی شوند.

\textbf{فیلدِ \lr{\texttt{\_metadata\_extra}} (نوع: \lr{\texttt{Mapping[str, Any]}}):} این فیلد کلیدهای ناشناخته‌ی \lr{\texttt{\_metadata}} را که از \lr{\texttt{from\_config\_dict}} عبور کرده‌اند، برای \lr{round-trip} حفظ می‌کند. استفاده از \lr{\texttt{MappingProxyType}} به‌جای \lr{\texttt{dict}}، تغییرناپذیریِ این فیلد را در یک \lr{dataclass} یخ‌زده تضمین می‌کند.

\subsection{\lr{\texttt{\_\_hash\_\_ = None}: ناسازگاریِ تساویِ مبتنی بر تلورانس}}

یکی از تصمیماتِ طراحیِ ظریف، تنظیمِ \lr{\texttt{\_\_hash\_\_ = None}} پس از تعریفِ کلاس است:
\begin{LTR}
	\begin{lstlisting}[language=Python]
		FiberChannel.__hash__ = None  # type: ignore[assignment,method-assign]
	\end{lstlisting}
\end{LTR}
این کار به این دلیل ضروری است که تساویِ مبتنی بر \lr{\texttt{is\_close}} غیرانتقالی است: اگر $a \approx b$ و $b \approx c$، نمی‌توان نتیجه گرفت که $a \approx c$. این ویژگی با قراردادِ هش $a == b \Rightarrow \text{hash}(a) = \text{hash}(b)$ در تضاد است. به همین دلیل، کلاس به‌صورتِ صریح غیرقابل‌هش \lr{(unhashable)} می‌شود و تلاش برای استفاده از نمونه‌ها به‌عنوان کلیدِ \lr{dict} یا عضوِ \lr{set} با \lr{\texttt{TypeError}} مواجه می‌گردد.

\subsection{زیرکلاسِ \lr{\texttt{AttenuationOnlyFiberChannel}}}

این زیرکلاس به‌عنوان یک نشانگرِ نوعی \lr{(type marker)} عمل می‌کند:
\begin{LTR}
	\begin{lstlisting}[language=Python]
		class AttenuationOnlyFiberChannel(FiberChannel):
		__slots__ = ()
		def __eq__(self, other: Any) -> bool:
		if not isinstance(other, AttenuationOnlyFiberChannel):
		return NotImplemented
		return super().__eq__(other)
	\end{lstlisting}
\end{LTR}
بدنه‌ی کلاس خالی است (\lr{\texttt{\_\_slots\_\_ = ()}} برای جلوگیری از تداخلِ \lr{slot} در زیرکلاس‌سازی آینده) زیرا رفتارِ آن با کلاسِ پایه یکسان است. تنها تفاوت در \lr{\texttt{\_\_eq\_\_}} است که \lr{\texttt{type(self) == type(other)}} را الزامی می‌کند؛ به این ترتیب، یک \lr{\texttt{FiberChannel}} و یک \lr{\texttt{AttenuationOnlyFiberChannel}} با همان اسکالرها دیگر برابر در نظر گرفته نمی‌شوند. این امر به کدهای پایین‌دستی اجازه می‌دهد با \lr{\texttt{isinstance}} بررسی کنند که آیا کانالِ ارائه‌شده واقعاً صرفاً تلفاتی است یا خیر.

\subsection{ثابت‌های آستانه‌ای}

ماژول چندین ثابتِ آستانه‌ای در سطحِ ماژول تعریف می‌کند:
\begin{LTR}
	\begin{lstlisting}[language=Python]
		TRANSMITTANCE_SUBNORMAL_THRESHOLD: float = sys.float_info.min
		TRANSMITTANCE_UNDERFLOW_THRESHOLD: float = TRANSMITTANCE_SUBNORMAL_THRESHOLD
		MAX_DISTANCE_KM: float = 10_000.0
		WARNING_DISTANCE_KM: float = 1e-6
		MAX_FIBER_LOSS_DB_KM: float = 100.0
		_CONTRACT_CHECK_REL_TOL: float = 1e-12
		_CONTRACT_CHECK_ABS_TOL: float = 1e-12
		_ZERO_LOSS_ABS_TOL: float = 1e-12
		_TRANSMITTANCE_UNDERFLOW_ABS_TOL: float = sys.float_info.min
	\end{lstlisting}
\end{LTR}
این ثابت‌ها هر یک نقشِ مشخصی دارند. \lr{\texttt{TRANSMITTANCE\_SUBNORMAL\_THRESHOLD}} کوچک‌ترین عدد نرمالِ مثبت در \lr{IEEE-754} را نشان می‌دهد و آستانه‌ی تشخیصِ مقادیر زیرنرمال است. \lr{\texttt{MAX\_DISTANCE\_KM}} حداکثر طولِ فیبرِ قابل پذیرش برای یک قطعه‌ی منفرد را تعیین می‌کند (پیش از این $1000\,\text{km}$ بود که در نسخه‌های جدید به $10000\,\text{km}$ افزایش یافت تا لینک‌های اولترالانگ با تقویت‌کننده پوشش داده شوند). \lr{\texttt{MAX\_FIBER\_LOSS\_DB\_KM}} حداکثر ضریب تلفاتِ مجاز را تعیین می‌کند (افزایش از $10$ به $100\,\text{dB/km}$ برای پوششِ فیبرهای پلاستیکی و فیبرهای تخصصی). \lr{\texttt{\_CONTRACT\_CHECK\_*}} تلورانس‌های بررسیِ قراردادِ $L = \alpha \cdot L$ هستند. \lr{\texttt{\_ZERO\_LOSS\_ABS\_TOL}} آستانه‌ی یکپارچه‌ی ``تلفاتِ صفر'' است.

\subsection{مدیریت \lr{strict mode} با \lr{ContextVar}}

یکی از نقاطِ قوتِ معماری، استفاده از \lr{\texttt{contextvars.ContextVar}} برای مدیریتِ \lr{strict mode} است:
\begin{LTR}
	\begin{lstlisting}[language=Python]
		_STRICT_MODE: "contextvars.ContextVar[bool]" = contextvars.ContextVar(
		"QKD_STRICT",
		default=os.environ.get("QKD_STRICT", "").lower()
		in ("1", "true", "yes", "on"),
		)
		
		def get_strict_mode() -> bool:
		return _STRICT_MODE.get()
		
		def set_strict_mode(enabled: bool) -> contextvars.Token:
		return _STRICT_MODE.set(bool(enabled))
		
		def reset_strict_mode(token: contextvars.Token) -> None:
		_STRICT_MODE.reset(token)
	\end{lstlisting}
\end{LTR}
این رویکرد سه مزیت دارد: (الف) \lr{thread-safe} و \lr{async-safe} است (هر \lr{thread} یا \lr{task} مقدارِ مستقلِ خود را دارد)؛ (ب) امکان \lr{scoped strictness} را فراهم می‌کند (یک \lr{token} بازگردانده می‌شود که می‌توان با آن به مقدارِ قبلی بازگشت)؛ (ج) مقدارِ پیش‌فرض از متغیرِ محیطی \lr{\texttt{QKD\_STRICT}} خوانده می‌شود تا کاربر بتواند بدون تغییرِ کد، حالت را فعال کند.

\subsection{ثبتِ هشدارهای فیزیکی مشکوک}

کلاسِ \lr{\texttt{PhysicalSuspicionWarning}} به‌عنوان زیرکلاسِ \lr{\texttt{UserWarning}} تعریف می‌شود:
\begin{LTR}
	\begin{lstlisting}[language=Python]
		class PhysicalSuspicionWarning(UserWarning):
		"""Warning emitted when a channel configuration is physically suspicious."""
		
		def _warn_physical_suspicion(msg: str, *, stacklevel: int = 2) -> None:
		warnings.warn(msg, PhysicalSuspicionWarning, stacklevel=stacklevel)
	\end{lstlisting}
\end{LTR}
استفاده از \lr{\texttt{warnings.warn}} به‌جای \lr{\texttt{logger.warning}} تضمین می‌کند که هشدارها به‌صورت پیش‌فرض به کاربر نمایش داده می‌شوند (برخلاف \lr{\texttt{logger}} که نیازمند پیکربندی است). این هشدارها در مواردی مانند فاصله‌های زیرمیلی‌متری، $\alpha = 0$ با فاصله‌ی غیرصفر، و ناسازگاری طول‌موج با $\alpha$ صادر می‌شوند.


\subsection{کلاسِ \lr{\texttt{ElectricalNoiseConfig}}}

این کلاس در \lr{\texttt{noise\_models.py}} به‌صورت یک \lr{dataclass} با \lr{\texttt{slots=True}} (اما نه \lr{\texttt{frozen=True}}، زیرا مقادیرِ آن ممکن است در طولِ شبیه‌سازی به‌روز شوند) تعریف می‌شود:
\begin{LTR}
	\begin{lstlisting}[language=Python]
		@dataclass(slots=True)
		class ElectricalNoiseConfig:
		temperature_k: float = 298.15
		bandwidth_hz: float = 1e9
		load_resistance_ohm: float = 50.0
		monitor_photocurrent_a: Optional[float] = None
		
		def validate(self) -> None:
		if self.temperature_k <= 0:
		raise ParameterValidationError("temperature_k must be positive.")
		if self.bandwidth_hz < 0:
		raise ParameterValidationError("bandwidth_hz must be non-negative.")
		if self.load_resistance_ohm <= 0:
		raise ParameterValidationError("load_resistance_ohm must be positive.")
		if self.monitor_photocurrent_a is not None and self.monitor_photocurrent_a < 0:
		raise ParameterValidationError("monitor_photocurrent_a must be non-negative.")
	\end{lstlisting}
\end{LTR}
چهار فیلدِ آن پارامترهای فیزیکیِ نویز الکتریکی را در بر می‌گیرند: دما، پهنای‌باند، مقاومتِ بار، و جریانِ فوتودیودِ نظارتی. متدِ \lr{\texttt{validate}} به‌صورت صریح محدوده‌های فیزیکی را بررسی می‌کند؛ به‌عنوان مثال، دمای صفر یا منفی به‌دلیلِ نقضِ ترمودینامیک رد می‌شود.

\subsection{کلاسِ \lr{\texttt{ChannelSimParams}}}

این کلاس به‌عنوان یک بسته‌ی یخ‌زده‌ی کمیت‌های مشتق‌شده عمل می‌کند:
\begin{LTR}
	\begin{lstlisting}[language=Python]
		@dataclasses.dataclass(frozen=True)
		class ChannelSimParams:
		transmittance: float
		total_loss_db: float
		fiber_loss_db_km: float
		distance_km: float
		dispersion_parameter_ps_nm_km: float = 0.0
		wavelength_nm: Optional[float] = None
	\end{lstlisting}
\end{LTR}
این کلاس با ثابت‌کردنِ مقادیر در زمانِ ساخت، اطمینان حاصل می‌کند که کدِ شبیه‌سازیِ پایین‌دستی همیشه از یک تصوِیر لحظه‌ای \lr{(snapshot)} یکپارچه می‌خواند و خطرِ خواندنِ مقادیرِ قدیمی از دیکشنریِ پیکربندی را از بین می‌برد. تابعِ کارخانه‌ی \lr{\texttt{build\_channel\_sim\_params}} تنها مسیرِ تأییدشده‌ی ساختِ این کلاس است.

\section{پیاده‌سازی ریاضی و الگوریتمی}
\label{sec:channel-implementation}

\subsection{محاسبه‌ی ضریب انتقال در \lr{\texttt{\_compute\_transmittance}}}

این متد که در \lr{\texttt{channel.py}} و نسخه‌ی مستقلِ آن در \lr{\texttt{\_physics.py}} پیاده‌سازی شده، قلبِ محاسباتِ فیزیکی ماژول است:
\begin{LTR}
	\begin{lstlisting}[language=Python]
		def _compute_transmittance(self) -> float:
		total_db = self._effective_total_loss_db()
		return self._transmittance_from_loss_db(total_db)
		
		@staticmethod
		def _transmittance_from_loss_db(total_db: float) -> float:
		if total_db == 0.0:
		if math.copysign(1.0, total_db) < 0:
		_warn_physical_suspicion(
		f"_transmittance_from_loss_db: total_db is negative zero (-0.0)...",
		stacklevel=2,
		)
		return 1.0
		return math.pow(10.0, -total_db / 10.0)
	\end{lstlisting}
\end{LTR}
این کد سه نکته‌ی ظریف را آدرس‌دهی می‌کند. نخست، استفاده از \lr{\texttt{self.\_effective\_total\_loss\_db()}} به‌جایِ \lr{\texttt{self.attenuation.total\_loss\_db}} امکانِ استفاده از \lr{\texttt{\_total\_loss\_db\_override}} را فراهم می‌کند. دوم، شاخه‌ی \lr{\texttt{total\_db == 0.0}} برای جلوگیری از محاسبه‌ی \lr{\texttt{math.pow(10.0, 0.0)}} است که همیشه $1.0$ را برمی‌گرداند؛ اما این شاخه‌ی صریح به ما اجازه می‌دهد نویز صفرِ منفی \lr{(-0.0)} را که می‌تواند نشانه‌ی خطای بالادستی باشد، تشخیص دهیم. سوم، استفاده از \lr{\texttt{math.pow}} به‌جای عملگرِ \lr{\texttt{**}} به این دلیل است که \lr{\texttt{math.pow}} در برابرِ \lr{\texttt{OverflowError}} در پلتفرم‌های مختلف، رفتارِ یکدست‌تری دارد.

\subsection{بررسیِ سلامتِ ضریب انتقال}

متدِ \lr{\texttt{\_check\_transmittance\_health}} یک بررسیِ چهارمرحله‌ای انجام می‌دهد:
\begin{LTR}
	\begin{lstlisting}[language=Python]
		def _check_transmittance_health(self, transmittance: float) -> None:
		# NaN: reject unconditionally.
		if math.isnan(transmittance):
		raise ParameterValidationError(f"Transmittance computed as NaN...")
		
		# Negative zero: suspicious, may mask upstream sign error.
		if transmittance == 0.0 and math.copysign(1.0, transmittance) < 0:
		neg_zero_msg = (...)
		if get_strict_mode():
		raise ParameterValidationError(neg_zero_msg, ...)
		_warn_physical_suspicion(neg_zero_msg, stacklevel=4)
		
		total_db = self._effective_total_loss_db()
		
		# Underflow: exact-zero transmittance with non-zero total_db.
		if transmittance == 0.0 and total_db > 0.0:
		underflow_msg = (...)
		if get_strict_mode():
		raise ParameterValidationError(underflow_msg, ...)
		_warn_physical_suspicion(underflow_msg, stacklevel=4)
		# Subnormal: between 0 and smallest normal float.
		elif 0.0 < transmittance < TRANSMITTANCE_SUBNORMAL_THRESHOLD:
		subnormal_msg = (...)
		if get_strict_mode():
		raise ParameterValidationError(...)
		_warn_physical_suspicion(subnormal_msg, stacklevel=4)
		
		# Range validation.
		if not (0.0 <= transmittance <= 1.0):
		raise ParameterValidationError(f"Transmittance out of range [0, 1]...")
	\end{lstlisting}
\end{LTR}
هر یک از این مراحل به یک ناپایداریِ عددیِ متفاوت اشاره دارد. \lr{NaN} معمولاً ناشی از ورودی‌های \lr{NaN} است (که نباید از \lr{\texttt{AttenuationConfig}} معتبر عبور کنند، اما ممکن است از طریقِ \lr{pickle} یا \lr{monkey-patching} وارد شوند). زیرنرمال‌بودن نشانه‌ی فواصلِ بسیار طولانی است که دقت محاسبات را به‌شدت کاهش می‌دهد. سرریز به صفر نیز مشکلِ مشابهی دارد اما شدیدتر است زیرا در لگاریتم‌ها به $-\infty$ می‌انجامد.

\subsection{محاسبه‌ی کسر تلفاتِ تک‌کیلومتری}

این متد در \lr{\texttt{\_physics.py}} به‌صورت یک تابعِ مستقل پیاده‌سازی شده است:
\begin{LTR}
	\begin{lstlisting}[language=Python]
		def compute_single_km_power_loss_fraction(channel: Any) -> float:
		alpha = get_fiber_loss_db_km(channel)
		if not is_finite_non_negative(alpha):
		raise ParameterValidationError(
		f"single_km_power_loss_fraction requires finite non-negative "
		f"alpha, got {alpha!r}",
		param_name="fiber_loss_db_km",
		param_value=alpha,
		)
		if alpha == 0.0:
		return 0.0
		return -math.expm1(-alpha * math.log(10.0) / 10.0)
	\end{lstlisting}
\end{LTR}
محاسبه‌ی ریاضی آن بر پایه‌ی رابطه‌ی \eqref{eq:expm1-stable} است. شاخه‌ی \lr{\texttt{alpha == 0.0}} برای جلوگیری از محاسبه‌ی \lr{\texttt{expm1(0)}} است که همیشه $0$ را برمی‌گرداند اما صرفه‌جویی در زمانِ محاسبه می‌کند. مهم‌تر از آن، این شاخه‌ی صریح، رفتار را در حالتِ $\alpha = 0$ (کانال بدون تلفات) تضمین می‌کند.

\subsection{تقریب خطیِ تلفات و هشدارِ منسوخ‌سازی}

متدِ \lr{\texttt{compute\_linear\_loss\_per\_km}} در \lr{\texttt{\_physics.py}} به‌صورت زیر پیاده‌سازی شده:
\begin{LTR}
	\begin{lstlisting}[language=Python]
		def compute_linear_loss_per_km(channel: Any) -> float:
		global _LINEAR_LOSS_WARNED
		with _LINEAR_LOSS_LOCK:
		if not _LINEAR_LOSS_WARNED:
		warnings.warn(
		"FiberChannel.linear_loss_per_km() is deprecated; use "
		"loss_fraction_per_km() for the exact per-km power-loss "
		"fraction...",
		DeprecationWarning,
		stacklevel=2,
		)
		_LINEAR_LOSS_WARNED = True
		
		alpha = get_fiber_loss_db_km(channel)
		val = alpha * math.log(10.0) / 10.0
		if val > 1.0:
		msg = (f"linear_loss_per_km() returned {val:.6g} > 1.0 for alpha={alpha:.6g}...")
		if get_strict_mode():
		raise ParameterValidationError(msg, ...)
		_warn_physical_suspicion(msg, stacklevel=2)
		val = 1.0  # clamp to physical range [0, 1]
		return val
	\end{lstlisting}
\end{LTR}
این متد سه ویژگی دارد. نخست، هشدارِ منسوخ‌سازی تنها یک‌بار در هر پردازش \lr{(process)} با استفاده از \lr{\texttt{threading.Lock}} برای \lr{thread-safety} emit می‌شود. دوم، مقدار بازگردانده‌شده بر اساسِ رابطه‌ی \eqref{eq:linear-approx} محاسبه می‌شود که واقعاً خطی است. سوم، اگر مقدار از $1$ فراتر رود (که برای $\alpha > 4.34\,\text{dB/km}$ رخ می‌دهد)، در حالت \lr{strict} استثنا پرتاب می‌شود و در غیر این صورت، مقدار به $1$ بریده \lr{(clamp)} می‌شود.

\subsection{ترکیب آبشاری کانال‌ها}

این متد در \lr{\texttt{\_physics.py}} به‌صورت زیر پیاده‌سازی شده:
\begin{LTR}
	\begin{lstlisting}[language=Python]
		def compute_cascade(channel: Any, other: Any) -> Any:
		if not isinstance(other, type(channel)):
		raise TypeError(f"cascade requires a {type(channel).__name__} argument, ...")
		total_loss_db = get_total_loss_db(channel) + get_total_loss_db(other)
		total_distance = get_distance_km(channel) + get_distance_km(other)
		if total_distance == 0.0:
		instance = type(channel)(
		AttenuationConfig(fiber_length=0.0, attenuation_coefficient=0.0),
		_construction_path="cascade",
		)
		else:
		instance = type(channel).from_total_loss(
		total_distance, total_loss_db, _construction_path="cascade",
		)
		merged_metadata: Dict[str, Any] = dict(channel._metadata_extra)
		merged_metadata.update(other._metadata_extra)
		if merged_metadata:
		object.__setattr__(instance, "_metadata_extra",
		types.MappingProxyType(merged_metadata))
		if total_distance > 0.0:
		avg_alpha = total_loss_db / total_distance
		if avg_alpha < 0.05 or avg_alpha > 5.0:
		msg = (f"cascade: The cascaded average alpha={avg_alpha:.6g} dB/km "
		f"is outside typical fiber ranges (0.05--5 dB/km)...")
		if get_strict_mode():
		raise ParameterValidationError(msg, ...)
		_warn_physical_suspicion(msg, stacklevel=2)
		return instance
	\end{lstlisting}
\end{LTR}
این کد چندین نکته‌ی طراحی را بازتاب می‌دهد. استفاده از \lr{\texttt{type(channel)}} به‌جای \lr{\texttt{FiberChannel}} به زیرکلاس‌ها اجازه می‌دهد نمونه‌ای از همان نوعِ زیرکلاس تولید کنند. شاخه‌ی \lr{\texttt{total\_distance == 0.0}} برای حالتی است که هر دو کانال طولِ صفر دارند؛ در این حالت، کانالِ معادلِ هویت \lr{(identity)} با $\alpha = 0$ ساخته می‌شود. ترکیبِ فراداده‌ها با تقدمِ کانالِ دوم \lr{(\texttt{other})} در صورت تداخلِ کلید انجام می‌شود. هشدارِ $\alpha_{\text{avg}}$ خارج از محدوده، کاربر را از اشتباهِ گرفتنِ کانالِ معادل با فیبرِ واقعی بازمی‌دارد.

\subsection{ترکیبِ آبشاریِ بدون اعتبارسنجی}

متدِ \lr{\texttt{compute\_cascade\_unchecked}} برای حلقه‌های \lr{Monte Carlo} با کاراییِ بالا طراحی شده:
\begin{LTR}
	\begin{lstlisting}[language=Python]
		def compute_cascade_unchecked(cls: Type[Any], channels: Sequence[Any]) -> Any:
		if not channels:
		raise ParameterValidationError("_cascade_unchecked requires at least one channel", ...)
		for ch in channels:
		if not isinstance(ch, cls):
		raise TypeError(f"Expected {cls.__name__}, got {type(ch).__name__}")
		total_loss = math.fsum(get_total_loss_db(ch) for ch in channels)
		total_distance = math.fsum(get_distance_km(ch) for ch in channels)
		if not is_finite_non_negative(total_loss):
		raise ParameterValidationError(...)
		if not is_finite_non_negative(total_distance):
		raise ParameterValidationError(...)
		metadata_extra: Dict[str, Any] = {}
		for ch in channels:
		metadata_extra.update(ch._metadata_extra)
		if total_distance > MAX_DISTANCE_KM:
		raise ParameterValidationError(...)
		if total_distance == 0.0:
		instance = fast_construct(cls, AttenuationConfig(0.0, 0.0),
		_construction_path="cascade")
		else:
		alpha = total_loss / total_distance
		if alpha > MAX_FIBER_LOSS_DB_KM:
		raise ParameterValidationError(...)
		if alpha < 0.05 or alpha > 5.0:
		_warn_physical_suspicion(...)
		instance = fast_construct(cls, AttenuationConfig(total_distance, alpha),
		_construction_path="cascade",
		_total_loss_db_override=total_loss)
		_eff_loss = get_total_loss_db(instance)
		_expected_T = transmittance_from_loss_db(_eff_loss)
		assert math.isclose(instance._transmittance, _expected_T,
		rel_tol=0.0, abs_tol=_CONTRACT_CHECK_ABS_TOL), \
		f"(L, T) contract violation after cascade: ..."
		if metadata_extra:
		object.__setattr__(instance, "_metadata_extra",
		types.MappingProxyType(metadata_extra))
		return instance
	\end{lstlisting}
\end{LTR}
نکاتِ کلیدی این پیاده‌سازی عبارت‌اند از: (الف) بررسیِ نوع پیش از دسترسی به ویژگی‌ها برای جلوگیری از \lr{\texttt{AttributeError}}؛ (ب) استفاده از \lr{\texttt{math.fsum}} برای جمعِ پایدارِ عددی در هزاران کانال؛ (ج) بررسیِ \lr{NaN} و \lr{Inf} در جمعِ تجمعی؛ (د) استفاده از \lr{\texttt{fast\_construct}} برای دور زدنِ اعتبارسنجیِ هر کانال؛ (هـ) \lr{assert} برای بررسیِ قراردادِ $(L,T)$ پس از ساخت (که در حالتِ \lr{-O} حذف می‌شود).

\subsection{محاسبه‌ی نویز حرارتی}

این متد در \lr{\texttt{noise\_models.py}} پیاده‌سازی شده:
\begin{LTR}
	\begin{lstlisting}[language=Python]
		def thermal_voltage_std(self) -> float:
		self.validate()
		return float(
		np.sqrt(
		4.0
		* CONST_BOLTZMANN
		* self.temperature_k
		* self.bandwidth_hz
		* self.load_resistance_ohm
		)
		)
	\end{lstlisting}
\end{LTR}
این کد مستقیماً رابطه‌ی \eqref{eq:thermal-std} را پیاده‌سازی می‌کند. فراخوانیِ \lr{\texttt{self.validate()}} پیش از محاسبه، اطمینان حاصل می‌کند که پارامترها در محدوده‌ی فیزیکی هستند. استفاده از \lr{\texttt{numpy.sqrt}} به‌جای \lr{\texttt{math.sqrt}} به این دلیل است که در صورت عبورِ یک آرایه (که ممکن است در حالت‌های پردازشِ برداری رخ دهد)، کد همچنان کار کند.

\subsection{محاسبه‌ی نویز شات}

\begin{LTR}
	\begin{lstlisting}[language=Python]
		def shot_voltage_std(self) -> float:
		self.validate()
		if self.monitor_photocurrent_a is None or self.monitor_photocurrent_a == 0.0:
		return 0.0
		i_std = np.sqrt(
		2.0
		* CONST_ELECTRON_CHARGE
		* self.monitor_photocurrent_a
		* self.bandwidth_hz
		)
		return float(i_std * self.load_resistance_ohm)
	\end{lstlisting}
\end{LTR}
این کد رابطه‌ی \eqref{eq:shot-std} را پیاده‌سازی می‌کند. شاخه‌ی بازگشتِ $0$ در صورتِ نبودِ جریان، یک بهینه‌سازیِ کوچک است اما مهم‌تر از آن، از محاسبه‌ی $\sqrt{0}$ و تولیدِ $0.0$ جلوگیری می‌کند که در برخی محیط‌ها ممکن است منجر به هشدارهای \lr{runtime} شود.

\subsection{ترکیبِ نویزهای مستقل}

\begin{LTR}
	\begin{lstlisting}[language=Python]
		def total_voltage_std(self) -> float:
		vt = self.thermal_voltage_std()
		vs = self.shot_voltage_std()
		return float(np.sqrt(vt * vt + vs * vs))
	\end{lstlisting}
\end{LTR}
این متد رابطه‌ی \eqref{eq:total-noise} را پیاده‌سازی می‌کند. استفاده از $v_t^2 + v_s^2$ به‌جای $v_t \cdot v_t + v_s \cdot v_s$ (که ریاضیاً معادل است) ممکن است به‌دلیلِ بهینه‌سازیِ پردازنده، اندکی سریع‌تر باشد. این الگوی جمعِ توان‌ها در همه‌ی مدل‌های نویزِ مستقلِ فریم‌ورک استفاده می‌شود.

\subsection{متدِ \lr{\texttt{from\_total\_loss}} و ساختِ دو مرحله‌ای}

این متد که برای حفظِ دقتِ \lr{round-trip} طراحی شده، الگوی ساختِ دو مرحله‌ای را به‌کار می‌برد:
\begin{LTR}
	\begin{lstlisting}[language=Python]
		@classmethod
		def from_total_loss(cls, distance_km, total_loss_db, *,
		_construction_path="from_total_loss",
		wavelength_nm=None):
		coerced_distance = _coerce_numeric(distance_km, param_name="distance_km")
		coerced_total_loss = _coerce_numeric(total_loss_db, param_name="total_loss_db")
		
		if not is_finite_non_negative(coerced_distance):
		raise ParameterValidationError(...)
		if not is_finite_non_negative(coerced_total_loss):
		raise ParameterValidationError(...)
		
		if coerced_distance == 0.0:
		if not is_close(coerced_total_loss, 0.0, abs_tol=_ZERO_LOSS_ABS_TOL):
		raise ParameterValidationError(f"distance_km=0 requires total_loss_db=0...")
		attenuation = AttenuationConfig(fiber_length=0.0, attenuation_coefficient=0.0)
		return cls(attenuation=attenuation, _construction_path=_construction_path,
		wavelength_nm=wavelength_nm)
		
		if coerced_distance > MAX_DISTANCE_KM:
		raise ParameterValidationError(...)
		
		if is_close(coerced_total_loss, 0.0, abs_tol=_ZERO_LOSS_ABS_TOL):
		msg = (f"from_total_loss: distance_km={coerced_distance:.6g} > 0 "
		f"but total_loss_db={coerced_total_loss:.6g} is zero...")
		if get_strict_mode():
		raise ParameterValidationError(msg, ...)
		_warn_physical_suspicion(msg, stacklevel=2)
		
		fiber_loss_db_km = coerced_total_loss / coerced_distance
		if not math.isfinite(fiber_loss_db_km):
		raise ParameterValidationError(...)
		if fiber_loss_db_km > MAX_FIBER_LOSS_DB_KM:
		raise ParameterValidationError(...)
		
		attenuation = AttenuationConfig(fiber_length=coerced_distance,
		attenuation_coefficient=fiber_loss_db_km)
		
		# Two-phase construction: skip health check during __post_init__
		_skip_health_token = _SKIP_TRANSMITTANCE_HEALTH.set(True)
		try:
		instance = cls(attenuation=attenuation,
		_construction_path=_construction_path,
		wavelength_nm=wavelength_nm)
		finally:
		_SKIP_TRANSMITTANCE_HEALTH.reset(_skip_health_token)
		object.__setattr__(instance, "_total_loss_db_override", coerced_total_loss)
		new_t = cls._transmittance_from_loss_db(coerced_total_loss)
		object.__setattr__(instance, "_transmittance", new_t)
		instance._check_transmittance_health(new_t)
		return instance
	\end{lstlisting}
\end{LTR}
این متد سه ویژگیِ مهم دارد. نخست، الگوی \lr{ContextVar} به نام \lr{\texttt{\_SKIP\_TRANSMITTANCE\_HEALTH}} برای دور زدنِ بررسیِ سلامت در \lr{\texttt{\_\_post\_init\_\_}} استفاده می‌شود؛ این کار از هشدارهای متناقض جلوگیری می‌کند (مقدارِ موقتاً محاسبه‌شده از $\alpha \cdot L$ ممکن است متفاوت از مقدارِ نهاییِ محاسبه‌شده از \lr{override} باشد). دوم، \lr{override} پس از ساختِ نمونه با \lr{\texttt{object.\_\_setattr\_\_}} تنظیم می‌شود. سوم، بررسیِ نهاییِ سلامت روی مقدارِ نهایی انجام می‌شود.

\subsection{متدِ \lr{\texttt{\_fast\_construct}}}

این متد برای حلقه‌های \lr{Monte Carlo} طراحی شده و تمام اعتبارسنجی‌ها را دور می‌زند:
\begin{LTR}
	\begin{lstlisting}[language=Python]
		@classmethod
		def _fast_construct(cls, attenuation, *,
		_construction_path="fast_construct",
		_total_loss_db_override=None,
		_wavelength_nm=None):
		if not isinstance(attenuation, AttenuationConfig):
		raise TypeError(f"_fast_construct requires an AttenuationConfig, got ...")
		instance = object.__new__(cls)
		object.__setattr__(instance, "attenuation", attenuation)
		object.__setattr__(instance, "_construction_path", _construction_path)
		object.__setattr__(instance, "_total_loss_db_override", _total_loss_db_override)
		object.__setattr__(instance, "_strict_mode_at_construction", get_strict_mode())
		object.__setattr__(instance, "_metadata_extra", _EMPTY_METADATA_EXTRA)
		object.__setattr__(instance, "wavelength_nm", _wavelength_nm)
		
		total_db = (_total_loss_db_override if _total_loss_db_override is not None
		else attenuation.total_loss_db)
		t = cls._transmittance_from_loss_db(total_db)
		if not math.isfinite(t):
		raise ParameterValidationError(...)
		if not (0.0 <= t <= 1.0):
		raise ParameterValidationError(...)
		if 0.0 < t < TRANSMITTANCE_SUBNORMAL_THRESHOLD:
		logger.debug("_fast_construct: computed transmittance %r is subnormal...")
		object.__setattr__(instance, "_transmittance", t)
		return instance
	\end{lstlisting}
\end{LTR}
این متد از \lr{\texttt{object.\_\_new\_\_(cls)}} به‌جای فراخوانیِ سازنده‌ی کلاس استفاده می‌کند تا \lr{\texttt{\_\_post\_init\_\_}} اجرا نشود. سپس تمام فیلدها به‌صورت دستی با \lr{\texttt{object.\_\_setattr\_\_}} تنظیم می‌شوند. با وجودِ دور زدنِ اعتبارسنجی، یک بررسیِ حداقلیِ \lr{\texttt{isfinite}} و بازه‌ی $[0,1]$ روی ضریب انتقال انجام می‌شود تا از تولیدِ نمونه‌های کاملاً نادرست جلوگیری شود.

\subsection{متدِ \lr{\texttt{validate}}}

این متد برای بررسیِ سازگاریِ حالتِ مشتق‌شده پس از عملیاتی که ممکن است آن را فاسد کند، طراحی شده:
\begin{LTR}
	\begin{lstlisting}[language=Python]
		def validate(self) -> None:
		self.validate_parameter("distance_km", self.distance_km, MAX_DISTANCE_KM)
		self.validate_parameter("fiber_loss_db_km", self.fiber_loss_db_km,
		MAX_FIBER_LOSS_DB_KM)
		
		if self._total_loss_db_override is None:
		expected_total_db = (self.attenuation.fiber_length
		* self.attenuation.attenuation_coefficient)
		if not is_close(self.total_loss_db, expected_total_db,
		rel_tol=_CONTRACT_CHECK_REL_TOL,
		abs_tol=_CONTRACT_CHECK_ABS_TOL):
		raise ParameterValidationError(f"Derived state inconsistency: ...")
		
		expected_transmittance = self._transmittance_from_loss_db(self.total_loss_db)
		
		if expected_transmittance == 0.0 and self.transmittance == 0.0:
		logger.warning("validate(): both expected and stored transmittance are 0.0...")
		elif not is_close(self.transmittance, expected_transmittance,
		rel_tol=1e-9,
		abs_tol=_TRANSMITTANCE_UNDERFLOW_ABS_TOL):
		raise ParameterValidationError(f"Derived state inconsistency: transmittance...")
		
		self._check_transmittance_health(self._transmittance)
		
		if 0.0 < self.transmittance < TRANSMITTANCE_SUBNORMAL_THRESHOLD:
		subnormal_warning = (f"validate(): Transmittance {self.transmittance:.3e} "
		f"is subnormal...")
		if get_strict_mode():
		raise ParameterValidationError(subnormal_warning, ...)
		_warn_physical_suspicion(subnormal_warning, stacklevel=2)
	\end{lstlisting}
\end{LTR}
این متد سه سطحِ بررسی دارد: بازه‌ی پارامترها، قراردادِ $L = \alpha \cdot L$ (تنها در صورت نبودِ \lr{override})، و سازگاریِ ضریبِ انتقال. تلورانسِ \lr{\texttt{\_TRANSMITTANCE\_UNDERFLOW\_ABS\_TOL}} به‌عنوان تلورانسِ مطلق برای همه‌ی مقایسه‌ها (حتی برای مقادیرِ نرمال) استفاده می‌شود تا مقادیرِ زیرنرمال که تنها $1$ یا $2$ بیتِ معتبر دارند، به‌اشتباهِ ناسازگار تشخیص داده نشوند.

\subsection{متدِ \lr{\texttt{link\_efficiency}}}

این تابع کمکی ترکیبِ سه‌گانه‌ی رابطه‌ی \eqref{eq:link-budget} را متمرکز می‌کند:
\begin{LTR}
	\begin{lstlisting}[language=Python]
		def link_efficiency(channel_transmittance, detector_efficiency,
		coupling_efficiency=1.0):
		result = float(channel_transmittance) * float(detector_efficiency) \
		* float(coupling_efficiency)
		if result < 0.0:
		return 0.0
		if result > 1.0:
		return 1.0
		return result
	\end{lstlisting}
\end{LTR}
برش به بازه‌ی $[0,1]$ در اینجا عمدی است: اگر کاربر به‌اشتباه یک بازدهیِ آشکارسازِ بزرگ‌تر از $1$ ارسال کند، نتیجه به‌جای پرتابِ استثنا به $1$ بریده می‌شود. این رفتار برای حلقه‌های \lr{Monte Carlo} که در آنجا سرعت مهم‌تر از تشخیصِ خطاست، مناسب است.

\subsection{توابع برداریِ \lr{\texttt{transmittance\_array}} و \lr{\texttt{transmittance\_scalar}}}

این دو تابع برای حلقه‌های \lr{Monte Carlo} با کاراییِ بالا طراحی شده‌اند:
\begin{LTR}
	\begin{lstlisting}[language=Python]
		def transmittance_array(distances_km, fiber_loss_db_km):
		distances_km = np.asarray(distances_km, dtype=np.float64)
		if fiber_loss_db_km == 0.0:
		return np.ones_like(distances_km)
		total_loss_db = fiber_loss_db_km * distances_km
		return np.float_power(10.0, -total_loss_db / 10.0)
		
		def transmittance_scalar(distance_km, fiber_loss_db_km):
		if fiber_loss_db_km == 0.0 or distance_km == 0.0:
		return 1.0
		total_db = fiber_loss_db_km * distance_km
		return math.pow(10.0, -total_db / 10.0)
	\end{lstlisting}
\end{LTR}
تابعِ \lr{\texttt{transmittance\_array}} از قابلیتِ \lr{vectorization} \lr{NumPy} استفاده می‌کند و برای محاسبه‌ی همزمانِ ضرایبِ انتقال برای هزاران فاصله طراحی شده. تابعِ \lr{\texttt{transmittance\_scalar}} تنها از \lr{\texttt{math.pow}} و حسابِ پایه استفاده می‌کند تا با \lr{Numba @njit} سازگار باشد. هیچ‌کدام از این توابع هشدارِ اعتبارسنجی emit نمی‌کنند زیرا در حلقه‌های داغ، هزینه‌ی هشدارهای نقطه‌به‌نقطه بیش از حد است.

\section{ارسال و دریافت با سایر ماژول‌ها}
\label{sec:channel-interface}

این سه ماژول با طیفِ گسترده‌ای از سایر اجزای فریم‌ورک در ارتباط هستند. در این بخش، این ارتباطات را به‌تفصیل بررسی می‌کنیم.

\subsection{وابستگی‌های ورودی به \lr{\texttt{channel.py}}}

\textbf{ماژولِ \lr{\texttt{qkd.exceptions}}:} دو استثنای \lr{\texttt{ParameterValidationError}} و \lr{\texttt{ConfigurationError}} از این ماژول وارد می‌شوند. این دو استثنا نقش‌های متفاوتی دارند: \lr{\texttt{ParameterValidationError}} برای خطاهای مربوط به مقدارِ یک پارامتر خاص (مانند \lr{\texttt{distance\_km = -1}}) استفاده می‌شود، در حالی که \lr{\texttt{ConfigurationError}} برای خطاهای ساختاریِ پیکربندی (مانند بخشِ گم‌شده‌ی \lr{channel\_config}) به‌کار می‌رود. علاوه بر این، کلاسِ \lr{\texttt{ChannelError}} که در \lr{\texttt{channel.py}} به‌صورت محلی تعریف شده، به‌عنوان کلاسِ پایه برای تمام استثناهای مرتبط با کانال عمل می‌کند (هرچند که در نسخه‌ی فعلی، \lr{\texttt{ConfigurationError}} و \lr{\texttt{ParameterValidationError}} هنوز از آن ارث نمی‌برند). تا این ارث‌بری پیاده‌سازی شود، تاپلِ \lr{\texttt{CHANNEL\_ERRORS = (ConfigurationError, ParameterValidationError)}} به‌عنوان سازوکارِ ایمنِ \lr{catch-all} توصیه می‌شود.

\textbf{ماژولِ \lr{\texttt{qkd.constants}}:} دو تابعِ \lr{\texttt{is\_close}} و \lr{\texttt{is\_finite\_non\_negative}} از این ماژول وارد می‌شوند. \lr{\texttt{is\_close}} برای مقایسه‌های مبتنی بر تلورانس در تساوی و بررسیِ قراردادها استفاده می‌شود. \lr{\texttt{is\_finite\_non\_negative}} برای اعتبارسنجیِ پارامترها به‌کار می‌رود و اطمینان حاصل می‌کند که مقادیر \lr{NaN}، \lr{Inf} و منفی رد می‌شوند.

\textbf{ماژولِ \lr{\texttt{qkd.datatypes}}:} کلاسِ \lr{\texttt{AttenuationConfig}} از این ماژول وارد می‌شود. این کلاس یک \lr{dataclass} ساده با دو فیلدِ \lr{\texttt{fiber\_length}} و \lr{\texttt{attenuation\_coefficient}} است. قراردادِ آن این است که \lr{\texttt{total\_loss\_db = fiber\_length * attenuation\_coefficient}} باید برقرار باشد.

\textbf{ماژولِ \lr{\texttt{\_schema}} (داخلی):} این ماژول ثابت‌ها، تایپ‌ها و توابعِ پایه‌ای را که در نسخه‌ی تجزیه‌شده‌ی \lr{\texttt{channel.py}} به اشتراک گذاشته می‌شوند، فراهم می‌سازد. از جمله می‌توان به \lr{\texttt{TRANSMITTANCE\_SUBNORMAL\_THRESHOLD}}، \lr{\texttt{MAX\_DISTANCE\_KM}}، \lr{\texttt{MAX\_FIBER\_LOSS\_DB\_KM}}، \lr{\texttt{\_CONTRACT\_CHECK\_ABS\_TOL}}، تایپ‌های \lr{\texttt{Transmittance}}، \lr{\texttt{Distance}}، \lr{\texttt{Loss}} و \lr{\texttt{Numeric}}، و توابعِ \lr{\texttt{get\_strict\_mode}} و \lr{\texttt{\_warn\_physical\_suspicion}} اشاره کرد.

\textbf{ماژولِ \lr{\texttt{\_gating}} (داخلی):} توابعِ \lr{\texttt{transmittance\_from\_loss\_db}}، \lr{\texttt{check\_transmittance\_health}}، \lr{\texttt{validate\_parameter}} و \lr{\texttt{effective\_total\_loss\_db}} از این ماژول وارد می‌شوند. این توابع لایه‌ی ``گیتینگ'' \lr{(gating)} را تشکیل می‌دهند که اعتبارسنجی و محاسباتِ پایه‌ای را فراهم می‌سازد.

\textbf{ماژولِ \lr{\texttt{\_casting}} (داخلی):} تابعِ \lr{\texttt{fast\_construct}} از این ماژول وارد می‌شود. این تابع همانند متدِ \lr{\texttt{\_fast\_construct}} در کلاس، اما در سطحِ ماژول عمل می‌کند.

\textbf{بسته‌های استاندارد:} \lr{\texttt{contextlib}} برای مدیریتِ زمینه، \lr{\texttt{contextvars}} برای \lr{ContextVar}، \lr{\texttt{dataclasses}} برای \lr{dataclass}، \lr{\texttt{logging}} برای لاگ‌گذاری، \lr{\texttt{math}} برای توابعِ ریاضی، \lr{\texttt{types}} برای \lr{\texttt{MappingProxyType}}، \lr{\texttt{warnings}} برای هشدارها، و \lr{\texttt{numpy}} برای محاسباتِ برداری.

\subsection{وابستگی‌های ورودی به \lr{\texttt{noise\_models.py}}}

\textbf{ماژولِ \lr{\texttt{qkd.exceptions}}:} استثنای \lr{\texttt{ParameterValidationError}} وارد می‌شود.

\textbf{ماژولِ \lr{\texttt{qkd.constants}}:} دو ثابتِ فیزیکیِ \lr{\texttt{CONST\_BOLTZMANN}} و \lr{\texttt{CONST\_ELECTRON\_CHARGE}} وارد می‌شوند. این ثابت‌ها در \lr{\texttt{constants\_definitions.py}} با مقادیر دقیقِ \lr{SI} تعریف شده‌اند.

\textbf{بسته‌های استاندارد:} \lr{\texttt{dataclasses}} برای \lr{dataclass}، \lr{\texttt{typing}} برای \lr{\texttt{Optional}}، و \lr{\texttt{numpy}} برای \lr{\texttt{sqrt}}.

\subsection{خروجی‌ها و مصرف‌کنندگان}

نمونه‌ی \lr{\texttt{FiberChannel}} پس از ساخت، در سراسر فریم‌ورک به‌عنوان مرجعِ مرکزیِ مدلِ کانال استفاده می‌شود. مصرف‌کنندگانِ اصلی عبارت‌اند از:

\textbf{موتور شبیه‌سازی اصلی (\lr{\texttt{main\_optimized.py}}):} این موتور از \lr{\texttt{FiberChannel}} برای استخراجِ \lr{\texttt{transmittance}}، \lr{\texttt{total\_loss\_db}}، \lr{\texttt{fiber\_loss\_db\_km}} و \lr{\texttt{distance\_km}} استفاده می‌کند. تابعِ \lr{\texttt{build\_channel\_sim\_params}} تمام این کمیت‌ها را در یک \lr{\texttt{ChannelSimParams}} یخ‌زده بسته‌بندی می‌کند تا موتور شبیه‌سازی همیشه از یک تصویرِ لحظه‌ایِ یکپارچه بخواند. این طراحی از خطرِ خواندنِ مقادیرِ قدیمی از دیکشنریِ پیکربندی در طولِ شبیه‌سازی جلوگیری می‌کند.

\textbf{ماژول‌های اثباتِ امنیتی (\lr{\texttt{proofs/}}):} هر یک از ماژول‌های اثبات (مانند \lr{\texttt{Ma 2005}}، \lr{\texttt{Lim 2014}} و \lr{\texttt{MDI-QKD}}) از \lr{\texttt{transmittance}} برای محاسبه‌ی بهره‌های $Y_n$ و نرخ‌های خطای $e_n$ استفاده می‌کنند. به‌عنوان مثال، در اثباتِ \lr{Ma 2005}، نرخ کلیدِ مجانبی از رابطه‌ی \eqref{eq:decoy-gain} محاسبه می‌شود که در آن $\eta$ به‌طور مستقیم از \lr{\texttt{link\_efficiency(channel.transmittance, detector\_efficiency, coupling\_efficiency)}} به دست می‌آید. استفاده از تابعِ متمرکزِ \lr{\texttt{link\_efficiency}} تضمین می‌کند که تمام اثبات‌ها از یک ترکیب‌بندیِ واحد استفاده می‌کنند و از ناسازگاریِ پارامترها جلوگیری می‌شود.

\textbf{ماژول آشکارساز (\lr{\texttt{qkd.detectors}}):} این ماژول از \lr{\texttt{ChannelSimParams}} برای استخراجِ پارامترهای مرتبط با کانال (مانند \lr{\texttt{dispersion\_parameter\_ps\_nm\_km}}) که در مدل‌سازیِ پنجره‌ی زمانیِ آشکارساز استفاده می‌شوند، بهره می‌برد. مقادیرِ \lr{\texttt{transmittance}} و \lr{\texttt{total\_loss\_db}} نیز در محاسبه‌ی نرخِ شمارشِ تاریکِ مؤثر \lr{(effective dark count rate)} که به‌دلیلِ نشتِ فوتون از کانالِ کلاسیک  مجاور در سیستم‌های \lr{DWDM} افزایش می‌یابد، استفاده می‌شوند.

\textbf{ماژول \lr{Privacy Amplification}:} این ماژول از \lr{\texttt{total\_loss\_db}} برای تخمینِ $\varepsilon_{\text{PE}}$ (خطای تخمین پارامتر) استفاده می‌کند؛ هرچه تلفاتِ کل بیشتر باشد، تعدادِ رویدادهای قابل مشاهده کمتر شده و $\varepsilon_{\text{PE}}$ باید بزرگ‌تر در نظر گرفته شود.

\textbf{ماژول‌های لاگ‌گذاری و بازتولید:} متدِ \lr{\texttt{to\_config\_dict}} برای ذخیره‌ی پیکربندی در فایل‌های \lr{JSON} به‌کار می‌رود. این فایل‌ها امکان بازتولیدِ دقیقِ آزمایش‌ها را فراهم می‌سازند. متدِ \lr{\texttt{\_\_repr\_\_}} و \lr{\texttt{\_\_str\_\_}} نیز برای لاگ‌های خوانا استفاده می‌شوند؛ این متدها برای مقادیرِ بسیار کوچکِ $T$ به‌طور خودکار به نمای علمی تغییر می‌کنند.

\textbf{ماژول \lr{ElectricalNoiseConfig}:} این کلاس به‌طور مستقل توسط ماژولِ منبع نوری \lr{(\texttt{qkd.sources})} برای مدل‌سازیِ نویز الکتریکیِ فوتودیودِ نظارتی استفاده می‌شود. خروجیِ \lr{\texttt{total\_voltage\_std()}} به موتورِ شبیه‌سازی ارسال می‌شود تا به‌عنوان نویزِ افزودنی در مدلِ آشکارساز اعمال گردد.

\subsection{جریان داده از ورودی تا خروجی}

برای روشن‌تر شدنِ ارتباطات، جریانِ داده را می‌توان به‌صورت زیر خلاصه کرد:

\begin{enumerate}
	\item \textbf{ورودی:} کاربر یک دیکشنریِ پیکربندی (یا یک فایلِ \lr{JSON}) ارائه می‌دهد که شامل کلیدهای \lr{\texttt{distance\_km}}، \lr{\texttt{fiber\_loss\_db\_km}} (یا \lr{\texttt{total\_loss\_db}}) و به‌اختیار \lr{\texttt{wavelength\_nm}} است.
	\item \textbf{بارگذاری:} متدِ \lr{\texttt{from\_config\_dict}} دیکشنری را دریافت کرده و با استفاده از \lr{\texttt{\_coerce\_numeric}} مقادیر را به \lr{float} تبدیل می‌کند. در صورت وجودِ \lr{\texttt{total\_loss\_db}}، از \lr{\texttt{from\_total\_loss}} استفاده می‌کند تا دقتِ \lr{round-trip} حفظ شود.
	\item \textbf{اعتبارسنجی:} سازنده‌ی \lr{\texttt{FiberChannel}} متدِ \lr{\texttt{\_\_post\_init\_\_}} را فراخوانی می‌کند که به‌نوبه‌ی خود هفت مرحله‌ی اعتبارسنجی را اجرا می‌کند: بررسیِ نوع، بازه‌ی پارامترها، معناداریِ $\alpha$، معناداریِ فاصله، محدودبودنِ تلفاتِ کل، قراردادِ $L = \alpha \cdot L$، و سازگاریِ طول موج.
	\item \textbf{محاسبه‌ی ضریب انتقال:} متدِ \lr{\texttt{\_compute\_transmittance}} مقدارِ $T$ را محاسبه کرده و در \lr{\texttt{\_transmittance}} ذخیره می‌کند. سپس متدِ \lr{\texttt{\_check\_transmittance\_health}} بررسیِ چهارمرحله‌ایِ \lr{NaN}، زیرنرمال، سرریز و بازه را انجام می‌دهد.
	\item \textbf{استفاده:} نمونه‌ی تأییدشده به موتور شبیه‌سازی، ماژول‌های اثباتِ امنیتی و ماژول آشکارساز ارسال می‌شود. تابعِ \lr{\texttt{build\_channel\_sim\_params}} تمام کمیت‌های لازم را در یک \lr{\texttt{ChannelSimParams}} یخ‌زده بسته‌بندی می‌کند.
	\item \textbf{ذخیره‌سازی:} متدِ \lr{\texttt{to\_config\_dict}} پیکربندیِ نهایی را برای بازتولیدِ آینده ذخیره می‌کند. این خروجی شامل \lr{\texttt{schema\_version}}، فراداده‌ی زمانِ ساخت و تمام پارامترهای فیزیکی است.
	\item \textbf{لاگ‌گذاری:} متدِ \lr{\texttt{\_\_repr\_\_}} یک نمایشِ خوانا از کانال تولید می‌کند که در لاگ‌های شبیه‌سازی استفاده می‌شود.
\end{enumerate}

\subsection{مدیریتِ نسخه‌گذاری و سازگاریِپسین}

فیلدِ \lr{\texttt{\_metadata.version}} در خروجیِ \lr{\texttt{to\_config\_dict}} مدیریت می‌شود. در هنگامِ بارگذاری، اگر نسخه‌ی ذخیره‌شده با نسخه‌ی فعلی (\lr{\texttt{\_\_version\_\_ = "4.2.1"}}) متفاوت باشد، یک هشدار ثبت می‌شود. علاوه بر این، چندین سازوکارِ سازگاری‌پسین پیاده‌سازی شده است: نامِ قدیمیِ \lr{\texttt{STRICT\_MODE}} (که اکنون یک \lr{ContextVar} است) از طریقِ \lr{\texttt{\_\_getattr\_\_}} در سطحِ ماژول \lr{(PEP 562)} پشتیبانی می‌شود و هشدارِ منسوخ‌سازی emit می‌کند. نامِ قدیمیِ \lr{\texttt{loss\_fraction\_per\_km}} همچنان به‌عنوان نامِ مستعارِ منسوخ‌شده‌ی \lr{\texttt{single\_km\_power\_loss\_fraction}} کار می‌کند. متدِ \lr{\texttt{from\_config}} نیز با هشدارِ منسوخ‌سازی پشتیبانی می‌شود و کاربران به استفاده‌ی مستقیم از سازنده تشویق می‌شوند.

\subsection{رفتار در برابر خطا و پیام‌های تشخیصی}

یکی از ویژگی‌های برجسته‌ی این ماژول، کیفیتِ پیام‌های خطا است. تمامی استثناها با \lr{context} دقیق تولید می‌شوند: نامِ پارامترِ نامعتبر، مقدارِ دریافتی و نوعِ آن، نوعِ مورد انتظار، مقادیرِ مجاز (در صورتِ وجود)، و پیشنهادهای هوشمند در صورتِ وجودِ خطایِ تایپیِ نزدیک. این رویکرد، تجربه‌ی کاربر را به‌طور چشمگیری بهبود می‌بخشد. علاوه بر این، استفاده‌ی متمایز از \lr{\texttt{logger.warning}} (برای مواردِ فنیِ قابل‌توجه اما غیربحرانی)، \lr{\texttt{\_warn\_physical\_suspicion}} (برای مواردِ فیزیکیِ مشکوک) و \lr{\texttt{raise}} (برای مواردِ غیرقابل‌قبول) به کاربر اجازه می‌دهد تا به‌سرعت تشخیص دهد که آیا شبیه‌سازی می‌تواند ادامه یابد یا خیر. در حالتِ \lr{strict mode}، تمام هشدارهای \lr{\texttt{\_warn\_physical\_suspicion}} به استثنا تبدیل می‌شوند تا در خطوط لوله‌ی تولید، هیچ پیکربندیِ مشکوکی پذیرفته نشود.

\subsection{جمع‌بندیِ معماری}

به‌عنوان جمع‌بندی، این سه ماژول با طراحیِ چندلایه، تغییرناپذیر و نوع-امن خود، زیربنایِ پایدارِ فریم‌ورک شبیه‌سازیِ \lr{QKD} را در لایه‌ی کانال فراهم می‌سازند. تفکیکِ سه‌گانه‌ی \lr{\texttt{channel.py}} (کلاس اصلی و توابعِ کمکی)، \lr{\texttt{\_physics.py}} (توابعِ فیزیکیِ مستقل) و \lr{\texttt{noise\_models.py}} (مدل‌های نویزِ الکتریکی) الگوی تمیزی برای سازماندهیِ منطقِ پیچیده در پروژه‌های بزرگ فراهم می‌کند. کلاسِ \lr{\texttt{FiberChannel}} با بیش از پنجاه متد و فیلدِ دقیقاً نوع‌دار، یک ظرفِ پیکربندیِ کامل و قابل اتکا فراهم می‌سازد که در عینِ سادگیِ استفاده، قدرتِ بیانِ کاملِ پیچیدگی‌های فیزیکی و عددیِ مدل‌سازیِ کانالِ کوانتومی را دارد. کلاسِ \lr{\texttt{ElectricalNoiseConfig}} به‌عنوان یک ماژولِ مستقل اما مرتبط، امکانِ مدل‌سازیِ نویزِ الکتریکیِ سمتِ منبع را بدون افزوده‌کردنِ پیچیدگی به کلاسِ اصلیِ کانال فراهم می‌سازد. در مجموع، این معماری تعادلیِ ظریف میانِ دقتِ فیزیکی، پایداریِ عددی، کاراییِ محاسباتی و قابلیتِ استفاده فراهم می‌سازد که برای شبیه‌سازی‌های \lr{QKD} در سطحِ تولید ضروری است.

	
	\chapter{منابع نوری}
	\label{ch:sources}
	% ============================================================
%  QKD Simulation Framework — Chapter: Optical Sources & Modulators
%  Authors: Documentation generated for academic thesis
%  Encoding: UTF-8, XePersian-compatible
% ============================================================

% --- تنظیمات رفع سرریز جعبه‌افقی (Overfull \hbox) ---
\sloppy
\emergencystretch=3em
\hbadness=10000
\vbadness=10000

% --- پیکربندی بسته‌ی listings برای نمایش صحیح کد پایتون در متن RTL ---
\lstset{
	basicstyle=\ttfamily\small,
	keywordstyle=\color{blue}\bfseries,
	stringstyle=\color{teal},
	commentstyle=\color{gray!60!black}\itshape,
	numberstyle=\tiny\color{gray},
	showstringspaces=false,
	breaklines=true,
	breakatwhitespace=true,
	frame=single,
	rulecolor=\color{gray!40},
	backgroundcolor=\color{gray!5},
	language=Python,
	extendedchars=true,
	inputencoding=utf8,
	tabsize=4,
	keepspaces=true,
	columns=fullflexible,
	basewidth=0.55em,
	aboveskip=1em,
	belowskip=1em,
	xleftmargin=1em,
	xrightmargin=1em,
	escapeinside={(*@}{@*)},
	literate=
	{→}{{$\rightarrow$}}1
	{∂}{{$\partial$}}1
	{α}{{$\alpha$}}1
	{β}{{$\beta$}}1
	{γ}{{$\gamma$}}1
	{δ}{{$\delta$}}1
	{ε}{{$\varepsilon$}}1
	{θ}{{$\theta$}}1
	{λ}{{$\lambda$}}1
	{μ}{{$\mu$}}1
	{ν}{{$\nu$}}1
	{π}{{$\pi$}}1
	{ρ}{{$\rho$}}1
	{σ}{{$\sigma$}}1
	{τ}{{$\tau$}}1
	{φ}{{$\varphi$}}1
	{ψ}{{$\psi$}}1
	{ω}{{$\omega$}}1
	{Γ}{{$\Gamma$}}1
	{Δ}{{$\Delta$}}1
	{Θ}{{$\Theta$}}1
	{Λ}{{$\Lambda$}}1
	{Σ}{{$\Sigma$}}1
	{Ω}{{$\Omega$}}1
	{≤}{{$\leq$}}1
	{≥}{{$\geq$}}1
	{≠}{{$\neq$}}1
	{≈}{{$\approx$}}1
	{·}{{$\cdot$}}1
	{×}{{$\times$}}1
	{∞}{{$\infty$}}1
	{η}{{$\eta$}}1
	{ℏ}{{$\hbar$}}1
}

\section{نمای کلی و هدف ماژول}
\label{sec:sources-overview}

ماژولِ \lr{\texttt{sources.py}} به‌همراه ماژولِ کمکیِ \lr{\texttt{modulators.py}}، لایه‌ی پایه‌ای مدل‌سازی منابع نوری را در فریم‌ورک شبیه‌سازیِ توزیع کلید کوانتومی \lr{(QKD)} تشکیل می‌دهند. این دو فایل با هم، چرخه‌ی کاملِ تولید پالس نوری، اعمال ناکامل‌های کلاسیکِ شدت، تبدیل به تعداد فوتون و نمونه‌برداری آماری را پوشش می‌دهند. خروجیِ این ماژول، ورودیِ اصلیِ ماژولِ کانال \lr{\texttt{channel.py}} است که در آن ضریب انتقالِ فیبر بر تعداد میانگین فوتون اعمال می‌شود و سپس به آشکارسازها ارسال می‌گردد.

هدف بنیادینِ این ماژول، ارائه‌ی یک مدلِ واحد، تغییرناپذیر \lr{(immutable)} و قابل سریال‌سازی از منبع نوری است که هم برای منابعِ لیزریِ تضعیف‌شده با آمارِ پواسون و هم برای منابعِ نورِ گرمایی با آمارِ بوز-آینشتین و هم برای منابعِ حالت-کوانتومی-کلی با ماتریس چگالی در پایه‌ی فوک کار کند. این عموم‌بخشی از طریق یک سلسله‌مراتبِ وراثتی با کلاسِ پایه‌ی \lr{\texttt{OpticalSource}} و دو زیرکلاسِ \lr{\texttt{DensityMatrixSource}} و \lr{\texttt{PoissonSource}} محقق می‌شود. ماژول \lr{\texttt{modulators.py}} نیز کلاسِ \lr{\texttt{MZMConfig}} را برای مدل‌سازی مودولاتورِ ماخ-زندر \lr{(Mach-Zehnder Modulator)} فراهم می‌سازد که یکی از اجزای کلیدیِ لایه‌ی ناکامل‌های کلاسیک است.

حوزه‌ی responsibility این ماژول به‌دقت محدود شده است: این ماژول صراحتاً مدل‌سازیِ تصادفی‌سازی فاز، کانال‌های جانبی طیفی/زمانی/قطبش، همدوسی فاز میان پالس‌ها و پدیده‌ی \lr{afterpulsing} (که یک اثر آشکارسازی است) را پوشش نمی‌دهد. این تصمیم طراحی بر اساس اصلِ تک‌مسئولیتی \lr{(Single Responsibility Principle)} اتخاذ شده است تا هر ماژول تنها یک لایه‌ی فیزیکی را پوشش دهد و وابستگی‌های متقابلِ نامطلوب ایجاد نکند. به‌خصوص، تصادفی‌سازی فاز که در اثبات‌های امنیتی \lr{Decoy-State BB84} فرضِ بنیادین است، در این ماژول صرفاً به‌عنوان یک پرچمِ فراداده‌ای \lr{(\texttt{phase\_randomization\_assumed})} ثبت می‌شود؛ این پرچم یک \lr{claim} از سوی کاربر است که تجهیزات خارجیِ وی تصادفی‌سازی فاز هر پالس را تضمین می‌کند، نه یک ویژگیِ پیاده‌سازی‌شده در کد.

معماریِ داخلیِ ماژول بر دو لایه‌ی مجزا استوار است که تا حد امکان درونِ یک انتزاعِ واحد نگه داشته می‌شوند. لایه‌ی نخست، \lr{State Preparation} یا انتخاب نوع پالس است که مشخص می‌کند کدام پالس (سیگنال یا یکی از دکوی‌ها) با چه احتمال و چه میانگین فوتونِ پایه تولید شود. لایه‌ی دوم، \lr{Photon-Number Sampling} است که بر اساس آمارِ پواسون، گرمایی/هندسی یا قطرِ ماتریس چگالی در پایه‌ی فوک، تعداد فوتون را نمونه‌برداری می‌کند و سپس به‌صورت اختیاری، ناکامل‌های کلاسیکِ شدت (تضعیف \lr{MZM}، تقسیم کانال، ضریب باند جانبی \lr{SCM}، جیتر شدت و شکست گسیل) را اعمال می‌کند. این جداسازی به کاربر اجازه می‌دهد بدون تغییر در لایه‌ی فیزیکی، پیکربندیِ پالس‌ها را به‌صورت مستقل تنظیم کند.

یک نکته‌ی بحرانی درباره‌ی \lr{phase\_randomization\_assumed}: این پرچم در اصلاحیه‌ی \lr{F-04} از \lr{phase\_randomized} تغییر نام یافت تا ماهیتِ صرفاً فراداده‌ایِ آن به‌روشنی در نام بازتاب یابد. وقتی این پرچم \lr{True} باشد، کاربر اطمینان می‌دهد که نمونه‌برداریِ قطریِ پواسون/گرمایی با اثبات امنیتی سازگار است زیرا تصادفی‌سازیِ فاز، هر پالس را به یک مخلوطِ قطری از حالت‌های فوک تبدیل می‌کند. وقتی \lr{False} باشد، نمونه‌برداریِ قطری برای تحلیل امنیتیِ \lr{Decoy-State BB84} مناسب نیست و کاربر باید از \lr{\texttt{DensityMatrixSource}} با ماتریس‌های چگالیِ کامل استفاده کند. این تمایز در تحلیل‌های امنیتی حیاتی است و نادیده‌گرفتنِ آن می‌تواند به تخمینِ نادرستِ نرخ کلید منجر شود.

\section{مبانی تئوری و فیزیکی}
\label{sec:sources-theory}

\subsection{آمار پواسون و منابع لیزری تضعیف‌شده}

منبعِ استاندارد در پروتکل‌های \lr{QKD}، لیزری است که به‌شدت تضعیف می‌شود تا میانگین تعداد فوتون به‌ازای پالس به حدود $\mu < 1$ برسد. چنین منبعی به‌دلیل ماهیتِ آماریِ فرآیندِ گسیلِ خودبه‌خودیِ \lr{spontaneous emission} در محیطِ بهره‌ی لیزر، از آمارِ پواسون پیروی می‌کند. تابع توزیع احتمالِ تعداد فوتون در یک پالس با میانگین $\mu$ به‌صورت زیر بیان می‌شود:
\begin{equation}
	P(n;\mu) = \frac{\mu^n e^{-\mu}}{n!}, \quad n = 0, 1, 2, \dots
	\label{eq:poisson-pmf}
\end{equation}
که در آن $\mu$ میانگین تعداد فوتون، $n$ عدد فوتون و $e$ پایه‌ی لگاریتم طبیعی است. این توزیع دارای ویژگی‌های مهمی است: میانگین و واریانس هر دو برابرِ $\mu$ هستند ($\mathbb{E}[n] = \text{Var}(n) = \mu$)، که به معنای پراکندگیِ نسبتاً پایین است. برای $\mu = 0.5$ (یک مقدار متداول در \lr{BB84-Decoy})، احتمالِ پالسِ خالی $P(0) = e^{-0.5} \approx 0.607$ است، احتمالِ تک‌فوتونی $P(1) = 0.5 \cdot e^{-0.5} \approx 0.303$ و احتمالِ چندفوتونی $P(n \geq 2) \approx 0.090$. این کسرِ غیرقابل‌چشم‌پوشیِ پالس‌های چندفوتونی، پایه‌ی حمله‌ی \lr{Photon-Number Splitting (PNS)} است که \lr{Decoy-State} برای مقابله با آن معرفی شده است.

محاسبه‌ی عددیِ پایدارِ توزیعِ پواسون برای $\mu$‌های بزرگ یا $n$‌های بالا نیاز به احتیاط دارد. محاسبه‌ی مستقیمِ $\mu^n / n!$ به سرریز \lr{(overflow)} می‌انجامد زیرا برای $\mu = 10$ و $n = 20$، مقدارِ $10^{20}/20! \approx 4.1 \times 10^{-4}$ است اما فاکتوریلِ ۲۰ به‌خودی‌خود $\approx 2.4 \times 10^{18}$ است. پیاده‌سازیِ این فریم‌ورک از روش بازگشتیِ پایدارِ عددی استفاده می‌کند:
\begin{align}
	P(0;\mu) &= e^{-\mu}, \label{eq:poisson-recurrence-init} \\
	P(n+1;\mu) &= P(n;\mu) \cdot \frac{\mu}{n+1}, \label{eq:poisson-recurrence}
\end{align}
که در آن ابتدا $P(0)$ از طریقِ \lr{\texttt{math.exp(-mu)}} محاسبه می‌شود (که برای $\mu$‌های معمول در \lr{QKD}، یعنی $\mu < 10$، هیچ خطر سرریز ندارد) و سپس هر جمله از جمله‌ی قبلی با ضرب در $\mu/(n+1)$ به‌دست می‌آید. این روش علاوه بر پایداریِ عددی، از مزیتِ سرعت نیز برخوردار است زیرا نیازی به محاسبه‌ی فاکتوریل یا توان نیست و هر مرحله تنها یک ضرب و یک تقسیم دارد.

\subsection{آمار گرمایی (بوز-آینشتین) و منابع نور گرمایی}

برخلاف لیزر که گسیلِ تحریکیِ همدوس دارد، منابعِ نورِ گرمایی (مانند \lr{blackbody} یا \lr{chaotic light}) از آمارِ بوز-آینشتین پیروی می‌کنند. توزیعِ تعداد فوتون در یک مُدِ گرماییِ تک‌فرکانس با میانگین $\mu$ به‌صورت زیر است:
\begin{equation}
	P(n;\mu) = \frac{\mu^n}{(1+\mu)^{n+1}}, \quad n = 0, 1, 2, \dots
	\label{eq:thermal-pmf}
\end{equation}
که این توزیع با توزیع هندسی با پارامترِ $p = 1/(1+\mu)$ معادل است. ویژگیِ بارزِ این توزیع، واریانسِ بزرگ‌تر از میانگین است: $\text{Var}(n) = \mu + \mu^2 = \mu(\mu+1)$ که برای $\mu = 0.5$ برابر با $0.75$ است (در مقایسه با $0.5$ در پواسون). این پراکندگیِ بیشتر، منابعِ گرمایی را برای \lr{QKD} نامطلوب می‌سازد زیرا کسرِ پالس‌های چندفوتونی به‌طور قابل توجهی بالاتر است: برای $\mu = 0.5$، $P(n \geq 2) \approx 0.222$ در مقایسه با $0.090$ در پواسون.

با این حال، مدل‌سازی منابع گرمایی در فریم‌ورک برای دو هدف ضروری است: نخست، تحلیلِ امنیتی در برابر حمله‌ای که مهاجم منبعِ سیگنال را با منبعِ گرمایی جایگزین کند؛ و دوم، مدل‌سازیِ پس‌زمینه‌ی تابشیِ محیطی که در آشکارسازها به‌عنوان نویز ظاهر می‌شود. مانند پواسون، محاسبه‌ی پایدارِ عددی از طریق بازگشت انجام می‌شود:
\begin{align}
	P(0;\mu) &= \frac{1}{1+\mu}, \label{eq:thermal-recurrence-init} \\
	P(n+1;\mu) &= P(n;\mu) \cdot \frac{\mu}{1+\mu}, \label{eq:thermal-recurrence}
\end{align}
که در آن نسبتِ $\mu/(1+\mu)$ یک بار پیش‌محاسبه می‌شود و سپس هر جمله از جمله‌ی قبلی با یک ضرب به‌دست می‌آید. این روش برای $\mu$‌های بزرگ نیز پایدار است زیرا نسبت همواره کوچک‌تر از ۱ است.

\subsection{ماتریس چگالی در پایه‌ی فوک و مدل‌سازی حالت‌های غیرکلاسیک}

برای منابعی که حالت‌های غیرکلاسیک مانند حالت‌های تک‌فوتونیِ واقعی \lr{(Fock states)}، حالت‌های فشرده \lr{(squeezed states)} یا حالت‌های گربه‌ی شرودینگر \lr{(cat states)} تولید می‌کنند، آمار پواسون یا گرمایی کافی نیست. این فریم‌ورک با کلاسِ \lr{\texttt{DensityMatrixSource}} این حالت‌ها را پشتیبانی می‌کند که در آن هر پالس با یک ماتریس چگالی $\rho$ در پایه‌ی فوک $\{|0\rangle, |1\rangle, |2\rangle, \dots\}$ مشخص می‌شود. قطرِ این ماتریس، توزیعِ احتمالِ تعداد فوتون را تشکیل می‌دهد:
\begin{equation}
	p_n = \langle n | \rho | n \rangle = \rho_{nn}, \quad \sum_n p_n = \text{Tr}(\rho) = 1.
	\label{eq:dm-diagonal-pn}
\end{equation}
اعتبارِ فیزیکیِ ماتریس چگالی نیازمند سه شرط است: هرمیتی بودن ($\rho = \rho^\dagger$)، مثبت‌معینِ نیمه ($\rho \succeq 0$) و یک بودن رد ($\text{Tr}(\rho) = 1$). این سه شرط در تابع کمکیِ \lr{\texttt{\_validate\_density\_matrix}} بررسی می‌شوند. هرمیتی بودن از طریق مقایسه‌ی $\rho$ با $\rho^\dagger$ در تلورانسِ عددی تأیید می‌شود. مثبت‌معین بودن از طریق محاسبه‌ی مقدار ویژه‌ی کوچک‌تر و بررسیِ غیرمنفی بودن آن تأیید می‌شود. یک بودن رد نیز با مجموعِ قطر بررسی می‌گردد.

عنصرهای خارج از قطر ($\rho_{mn}$ برای $m \neq n$) که نشانگر همدوسیِ کوانتومی میان حالت‌های فوک هستند، در فرآیند نمونه‌برداری \lr{(photon-number sampling)} نادیده گرفته می‌شوند زیرا اندازه‌گیریِ تعداد فوتون یک اندازه‌گیریِ قطری در پایه‌ی فوک است و همدوسی‌ها را از بین می‌برد. با این حال، این عنصرها برای تحلیل‌های امنیتیِ پیشرفته‌تر (مانند محاسبه‌ی آنتروپیِ فون‌نویمان) معتبر هستند و در ساختار داده ذخیره می‌شوند. در ادامه‌ی این بخش، مدل‌سازیِ فرآیندهای اتلاف روی ماتریس چگالی را بررسی می‌کنیم.

\subsection{فرآیندهای اتلاف روی ماتریس چگالی: تقسیم کانال و شکست گسیل}

وقتی منبعِ ماتریس-چگالی با $N_{\text{channels}} > 1$ پیکربندی شود، یک شبکه‌ی تقسیم‌گر پرتو \lr{(beam-splitter network)} $N$-کاناله فرض می‌شود که انرژی به‌طور یکنواخت میان کانال‌ها تقسیم می‌شود. ضریب گذارندگی هر کانال $\eta = 1/N$ است. تأثیر این تقسیم بر قطرِ ماتریس چگالی، یک تبدیلِ اتلافِ برنولی-دوجمله‌ای \lr{(binomial loss)} است:
\begin{equation}
	p'_k = \sum_{n=k}^{\infty} p_n \binom{n}{k} \eta^k (1-\eta)^{n-k},
	\label{eq:binomial-loss-channel-split}
\end{equation}
که در آن $p'_k$ احتمالِ مشاهده‌ی $k$ فوتون در یک کانالِ خروجی است، به‌شرطی که $n$ فوتون در ورودی تقسیم‌گر بوده است. این تبدیل، فیزیکِ یک تقسیم‌گر پرتو ایده‌آل را به‌دقت بازتولید می‌کند: هر فوتون با احتمال $\eta$ به کانال مورد نظر می‌رود و با احتمال $1-\eta$ به یکی از کانال‌های دیگر. این فرمول در زمان ساختِ \lr{\texttt{DensityMatrixSource}} روی جدولِ احتمالاتِ عدد-فوتون \lr{(\texttt{\_number\_probs\_table})} اعمال می‌شود تا قطرِ ماتریس از ابتدا روی وضعیتِ پس از تقسیم تنظیم شود.

شکست گسیل \lr{(emission failure)} نیز یک فرآیند اتلاف است که با احتمال $1 - p_{\text{emit}}$ منبع اصلاً پالس گسیل نمی‌کند. مدلِ فیزیکیِ این پدیده، مخلوط‌کردنِ حالتِ گسیل‌شده با حالتِ خلأ است:
\begin{equation}
	\rho' = p_{\text{emit}} \cdot \rho + (1 - p_{\text{emit}}) \cdot |0\rangle\langle 0|,
	\label{eq:emission-failure-mixture}
\end{equation}
که بر قطر به‌صورت زیر تأثیر می‌گذارد:
\begin{align}
	p'_0 &= p_{\text{emit}} \cdot p_0 + (1 - p_{\text{emit}}), \label{eq:emission-failure-p0} \\
	p'_k &= p_{\text{emit}} \cdot p_k, \quad k > 0. \label{eq:emission-failure-pk}
\end{align}
این تبدیل نیز در زمان ساخت اعمال می‌شود تا جدول احتمالات از ابتدا روی وضعیتِ پس از شکست گسیل تنظیم گردد. توجه شود که این دو فرآیندِ اتلاف، بر خلاف \lr{MZM}، \lr{SCM} و جیتر شدت، به‌صورت تبدیل‌های قطری روی ماتریس چگالی قابل بیان هستند و به همین دلیل در \lr{\texttt{DensityMatrixSource}} به‌صورت خودکار اعمال می‌شوند. سایر ناکامل‌های کلاسیک که حالت را به‌صورت غیرقطری تغییر می‌دهند، در این کلاس نادیده گرفته می‌شوند و کاربر باید آن‌ها را پیش از ساخت به ماتریس چگالی اعمال کند.

\subsection{مودولاتور ماخ-زندر و انتقال آرام‌شیب}

مودولاتورِ ماخ-زندر \lr{(Mach-Zehnder Modulator, MZM)} یکی از اجزای کلیدی در سیستم‌های \lr{QKD} با مدولاسیون شدت است. این مودولاتور از دو شاخه‌ی تداخلی تشکیل شده که با اعمال ولتاژ روی یکی از شاخه‌ها، اختلاف فاز $\phi$ میان دو مسیر ایجاد می‌شود. خروجیِ مودولاتور به‌صورت زیر مدل می‌شود:
\begin{equation}
	\mu_{\text{out}} = \mu_{\min} + (\mu_{\max} - \mu_{\min}) \cos^2(\phi),
	\label{eq:mzm-transfer}
\end{equation}
که در آن $\mu_{\min}$ میانگین فوتونِ نشتی \lr{(leakage)} (مرتبط با نسبتِ خاموشی)، $\mu_{\max}$ حداکثر شدتِ خروجی و $\phi$ فازِ کل است که از جمعِ فازِ سیگنال و فاز بایاس به‌دست می‌آید:
\begin{equation}
	\phi = \frac{\pi V}{2 V_\pi} + \phi_{\text{bias}},
	\label{eq:mzm-phase}
\end{equation}
که در آن $V$ ولتاژِ اعمال‌شده، $V_\pi$ ولتاژِ نیم‌موج \lr{(half-wave voltage)} (ولتاژی که برای تغییرِ فاز $\pi$ نیاز است) و $\phi_{\text{bias}}$ فاز بایاسِ ثابت است. برای $\phi_{\text{bias}} = \pi/2$، نقطه‌ی کار روی شیبِ نیمه-حداکثر تنظیم می‌شود که بیشترین حساسیت خطی را فراهم می‌کند.

نسبتِ خاموشی \lr{(extinction ratio)} که با واحد دسی‌بل بیان می‌شود، نسبتی است میان حداکثر توانِ خروجی و حداقل آن (نشتی):
\begin{equation}
	\text{ER}_{\text{dB}} = 10 \log_{10}\left(\frac{\mu_{\max}}{\mu_{\min}}\right).
	\label{eq:extinction-ratio}
\end{equation}
از این رابطه، میانگین فوتون نشتی به‌صورت زیر محاسبه می‌شود:
\begin{equation}
	\mu_{\min} = \mu_{\max} \cdot 10^{-\text{ER}_{\text{dB}}/10}.
	\label{eq:mzm-leakage}
\end{equation}
این نشتی یک کانال جانبی امنیتی ایجاد می‌کند: حتی وقتی مودولاتور در حالت ``خاموش'' است، مقدار کمی نور گسیل می‌شود که می‌تواند اطلاعاتِ بیت را به مهاجم نشت دهد. در اثبات‌های امنیتیِ \lr{QKD}، این نشتی باید به‌صورت صریح در نظر گرفته شود و $\mu_{\min}$ در محاسبه‌ی نرخ کلید لحاظ گردد. در این فریم‌ورک، $\mu_{\min}$ به‌عنوان کفِ سختی در عملیاتِ \lr{\texttt{clip}} اعمال می‌شود تا خروجیِ مودولاتور هرگز کمتر از این مقدار نشود.

نویز ولتاژِ اعمالی که از راهِ درایور الکتریکی وارد می‌شود، باعث fluctuation تصادفیِ فاز می‌شود. اگر $V_{\text{noise}} \sim \mathcal{N}(0, \sigma_V^2)$ باشد، فازِ نویزی برابر است با:
\begin{equation}
	\phi_{\text{noisy}} = \frac{\pi (V_{\text{signal}} + V_{\text{noise}})}{2 V_\pi} + \phi_{\text{bias}},
	\label{eq:mzm-noisy-phase}
\end{equation}
و خروجی نهایی به‌صورت $\mu_{\min} + (\mu_{\max} - \mu_{\min}) \cos^2(\phi_{\text{noisy}})$ محاسبه می‌شود. این نویز، یک fluctuations غیرخطی روی شدت ایجاد می‌کند که در ولتاژهای نویز بزرگ، می‌تواند به توزیعِ غیر-گاوسی روی $\mu_{\text{out}}$ منجر شود.

\subsection{مدولاسیون زیرحامل (SCM) و اتحاد ژاکوبی-آنگر}

مدولاسیونِ زیرحامل \lr{(Subcarrier Modulation, SCM)} تکنیکی است که در آن فازِ حاملِ نوری با یک فرکانسِ رادیویی مدوله می‌شود تا چندین کانالِ فرعی روی یک حاملِ مشترک قرار گیرند. وقتی فاز با شاخصِ مدولاسیون $m$ مدوله می‌شود، میدانِ خروجی از طریق اتحادِ ژاکوبی-آنگر بسط می‌یابد:
\begin{equation}
	e^{i m \sin(\Omega t)} = \sum_{n=-\infty}^{\infty} J_n(m) e^{i n \Omega t},
	\label{eq:jacobi-anger}
\end{equation}
که در آن $J_n(m)$ تابع بسل از نوع اول و مرتبه‌ی $n$ است و $\Omega$ فرکانسِ زیرحامل. این بسط، باند‌های جانبیِ متقارن در $n = \pm 1, \pm 2, \dots$ تولید می‌کند. توانِ هر باند جانبی از مرتبه‌ی $n$، متناسب با $J_n(m)^2$ است. از آنجا که $J_{-n}(m) = (-1)^n J_n(m)$، توانِ باندهای $+n$ و $-n$ یکسان است.

در این فریم‌ورک، تنها باندهای جانبی مرتبه‌ی اول ($n = \pm 1$) در نظر گرفته می‌شوند. ضریب مقیاسِ \lr{SCM} که به‌عنوان یک ضرب‌کننده‌ی سراسری روی میانگین فوتون اعمال می‌شود، به‌صورت زیر تعریف می‌شود:
\begin{equation}
	S_{\text{SCM}}(m) = N_{\text{sb}} \cdot J_1(m)^2,
	\label{eq:scm-exact}
\end{equation}
که در آن $N_{\text{sb}}$ تعداد باندهای جانبیِ مورد استفاده است: $N_{\text{sb}} = 1$ برای فیلترینگ تک‌باند جانبی \lr{(SSB)} و $N_{\text{sb}} = 2$ برای آشکارسازی دو‌باندی \lr{(DSB)}. برای مقادیر کوچک $m$، از تقریبِ زاویه کوچک استفاده می‌شود:
\begin{equation}
	J_1(m) \approx \frac{m}{2} - \frac{m^3}{16} + \mathcal{O}(m^5) \implies J_1(m)^2 \approx \frac{m^2}{4},
	\label{eq:bessel-small-angle}
\end{equation}
که در نتیجه:
\begin{equation}
	S_{\text{SCM}}(m) \approx N_{\text{sb}} \cdot \frac{m^2}{4}, \quad m \ll 1.
	\label{eq:scm-small-angle}
\end{equation}
این تقریب در آستانه‌ی $m \leq 10^{-3}$ با خطای نسبی کمتر از $2.5 \times 10^{-7}$ معتبر است (محاسبه‌شده از بسطِ $J_1(m) \approx m/2 - m^3/16$). این آستانه در ثابتِ \lr{\texttt{SCM\_SMALL\_ANGLE\_LIMIT}} تعریف شده است و در مقایسه با آستانه‌ی قدیمیِ $0.1$، خطای مرزی را ۱۰۰۰ برابر کاهش می‌دهد که برای بهینه‌سازیِ مبتنی بر گرادیان حیاتی است.

نکته‌ی ظریف درباره‌ی ناپیوستگیِ معنایی در $m = 0$: وقتی $m = 0$ باشد، \lr{SCM} غیرفعال است و $\mu$ به‌عنوان میانگین فوتونِ پس از \lr{SCM} تفسیر می‌شود، در حالی که برای هر $m > 0$، $\mu$ به‌عنوان انرژیِ کلِ گسیل‌شده تفسیر می‌شود و ضریب $S_{\text{SCM}}(m)$ کسرِ باند جانبی را می‌دهد. این تفسیرِ دوگانه در پرچمِ مستقلِ \lr{\texttt{scm\_enabled}} جدا شده است تا ناپیوستگیِ عددی به ناپیوستگیِ فیزیکی تبدیل شود. بهینه‌سازیِ مبتنی بر گرادیان در سراسرِ $m = 0$ نامعتبر است زیرا تفسیر فیزیکی به‌صورت ناپیوسته تغییر می‌کند.

\subsection{جیتر شدت: مدل گوسی با تصحیح بایاس}

جیتر شدت \lr{(intensity jitter)} fluctuations پالس-به-پالسِ میانگین فوتون را مدل می‌کند که از ناپایداریِ لیزرِ نیمه‌رسانا نشأت می‌گیرد. این فریم‌ورک دو مدل را پشتیبانی می‌کند: \lr{\texttt{RANDOM\_GAUSSIAN}} برای نویز گوسی افزودنی و \lr{\texttt{ADVERSARIAL\_BLOCK}} برای fluctuations همبسته‌ی بلوکی.

در مدلِ گوسی، نویز به‌صورت زیر به میانگین فوتون اعمال می‌شود:
\begin{equation}
	\mu_{\text{new}} = \mu \cdot |1 + \sigma z|, \quad z \sim \mathcal{N}(0, 1),
	\label{eq:gaussian-jitter}
\end{equation}
که در آن $\sigma$ جیتر نسبی \lr{(std/mean)} است. استفاده از قدرمطلق $|\cdot|$ به‌جای clip کردن ($\max(0, \cdot)$) به این دلیل است که بایاسِ قدرمطلق از مرتبه‌ی دومِ $P(X < 0)$ است، در حالی که بایاسِ clip از مرتبه‌ی اول است. با این حال، مدلِ قدرمطلق دارای بایاسِ مثبتِ سیستماتیک است:
\begin{equation}
	\mathbb{E}[|1 + \sigma z|] = \sigma \sqrt{\frac{2}{\pi}} e^{-1/(2\sigma^2)} + \text{erf}\left(\frac{1}{\sigma\sqrt{2}}\right) > 1,
	\label{eq:reflection-bias}
\end{equation}
که در آن $\text{erf}$ تابع خطا است. این بایاس در $\sigma = 0.5$ حدود $2\%$ و در $\sigma = 1.0$ حدود $16\%$ است. برای رفع این بایاس، فریم‌ورک از تصحیحِ تحلیلی استفاده می‌کند: مقدارِ $\mu_{\text{new}}$ بر $\text{bias\_factor}$ تقسیم می‌شود تا مدلِ میانگین-حفظ‌کننده شود ($\mathbb{E}[\mu_{\text{new}}] = \mu$):

\begin{LTR}
	\begin{lstlisting}[language=Python]
		if jitter > NUMERIC_ABS_TOL:
		inv_sigma = 1.0 / jitter
		bias_factor = (
		jitter * np.sqrt(2.0 / np.pi) * np.exp(-0.5 * inv_sigma * inv_sigma)
		+ scipy_erf(inv_sigma / np.sqrt(2.0))
		)
		mus = mus / bias_factor  # mean-preserving correction
	\end{lstlisting}
\end{LTR}

در مدلِ \lr{\texttt{ADVERSARIAL\_BLOCK}}، پالس‌ها به بلوک‌های متوالیِ اندازه‌ی $B$ تقسیم می‌شوند و تمامِ پالس‌های یک بلوک، یک ضریبِ fluctuation مشترک دارند که از توزیع لگاریتم-نرمال نمونه‌برداری می‌شود:
\begin{equation}
	\mu_{\text{new}} = \mu \cdot F_b, \quad F_b \sim \text{LogNormal}\left(-\frac{\sigma^2_{\ln}}{2}, \sigma^2_{\ln}\right),
	\label{eq:adversarial-block}
\end{equation}
که در آن $b$ شناسه‌ی بلوک، $F_b$ ضریبِ مشترکِ بلوک و $\sigma^2_{\ln} = \ln(1 + \sigma^2)$ پارامترِ لگاریتم-نرمال است. این پارامتری‌سازی به‌گونه‌ای انتخاب شده که $\mathbb{E}[F_b] = 1$ (میانگین-حفظ‌کننده) و $\text{Var}(F_b) = \sigma^2$ (واریانس برابر با جیتر نسبی). این مدل، سناریویی را شبیه‌سازی می‌کند که مهاجم بتواند شدتِ منبع را به‌صورت همبسته در بلوک‌هایی از پالس‌ها دستکاری کند، که در تحلیل‌های امنیتیِ پیشرفته حائز اهمیت است.

\section{تحلیل معماری و ساختار داده‌ها}
\label{sec:sources-architecture}

\subsection{سلسله‌مراتب وراثتی کلاس‌های منبع}

معماریِ کلاس‌های منبع بر یک سلسله‌مراتبِ وراثتی با کلاسِ پایه‌ی \lr{\texttt{OpticalSource}} و دو زیرکلاسِ \lr{\texttt{DensityMatrixSource}} و \lr{\texttt{PoissonSource}} استوار است. کلاسِ پایه با دو decorator تعریف می‌شود:

\begin{LTR}
	\begin{lstlisting}[language=Python]
		@dataclass(slots=True)
		class OpticalSource:
		"""Generic optical source backed by OpticalSourceConfig."""
		config: OpticalSourceConfig
		security_metadata: Optional[SourceSecurityMetadata] = None
		# ... per-instance fields like _base_mus_cache, _probabilities_cache, ...
	\end{lstlisting}
\end{LTR}

استفاده از \lr{\texttt{@dataclass(slots=True)}} به‌جای \lr{\texttt{frozen=True}} در این کلاس، یک تصمیمِ طراحیِ مهم است. برخلافِ \lr{\texttt{FiberChannel}} که کاملاً یخ‌زده است، \lr{\texttt{OpticalSource}} نیاز به state داخلیِ قابل تغییر دارد: کش‌های \lr{\texttt{\_base\_mus\_cache}} و \lr{\texttt{\_probabilities\_cache}} در زمانِ ساخت محاسبه می‌شوند، شمارنده‌های بلوکِ حریصانه \lr{(\texttt{\_adversarial\_block\_counter})} در طولِ شبیه‌سازی به‌روز می‌شوند، و \lr{\texttt{security\_metadata}} در هر فراخوانیِ \lr{\texttt{generate\_photons}} به‌روز می‌گردد. برای جلوگیری از تغییرِ ناخواسته‌ی پیکربندی، \lr{\texttt{\_\_setattr\_\_}} override شده است تا فقط فیلدهای مجاز را قابل تغییر سازد.

زیرکلاسِ \lr{\texttt{DensityMatrixSource}} دو فیلد اضافی دارد:
\begin{itemize}
	\item \lr{\texttt{density\_matrices: Optional[Dict[str, np.ndarray]]}}: نگاشت نام پالس به ماتریس چگالی.
	\item \lr{\texttt{missing\_dm\_policy: str}}: سیاستِ برخورد با پالس‌های بدون ماتریس (\lr{"error"} یا \lr{"vacuum"}).
	\item \lr{\texttt{heterogeneous\_dim\_policy: str}}: سیاستِ برخورد با ابعادِ ناهمگن (\lr{"warn"} یا \lr{"error"}).
\end{itemize}
این زیرکلاس الگوی طراحیِ \lr{Template Method} را با استفاده از \lr{\texttt{\_pre\_init\_hook}} و \lr{\texttt{\_post\_init\_hook}} پیاده می‌کند. این \texttt{hook}‌ها به زیرکلاس‌ها اجازه می‌دهند قبل و بعد از اعتبارسنجیِ پایه، کد اختصاصی خود را اجرا کنند. \lr{\texttt{DensityMatrixSource.\_pre\_init\_hook}} ماتریس‌های چگالی را coerce می‌کند و سیاست‌ها را اعتبارسنجی می‌نماید، در حالی که \lr{\texttt{\_post\_init\_hook}} تانسورِ \lr{\texttt{\_rho\_tensor}} و جدولِ \lr{\texttt{\_number\_probs\_table}} را می‌سازد.

زیرکلاسِ \lr{\texttt{PoissonSource}} یک convenience subclass است که صرفاً تایید می‌کند \lr{\texttt{statistics\_type == POISSON}} باشد. این کلاس برای خواناییِ کد و تمایز صریح از \lr{\texttt{DensityMatrixSource}} معرفی شده است.

\subsection{کلاس \lr{\texttt{SourceSimParams}}: بسته‌ی یخ‌زده‌ی پارامترهای شبیه‌سازی}

این کلاس با \lr{\texttt{@dataclasses.dataclass(frozen=True)}} تعریف می‌شود و شامل ده فیلد است که در زمان اجرا توسط \lr{\texttt{main\_optimized.py}} استفاده می‌شوند:

\begin{LTR}
	\begin{lstlisting}[language=Python]
		@dataclasses.dataclass(frozen=True)
		class SourceSimParams:
		pulse_period_ns: float
		source_rate: float
		statistics_type: str  # .value of SourceStatisticsType enum
		error_model: str      # .value of SourceErrorModel enum
		modulation_index: float
		intensity_jitter: float
		N_channels: int
		ideal_emission_probability: float
		n_pulse_types: int
		signal_pulse_index: int
	\end{lstlisting}
\end{LTR}

این بسته از طریقِ تابعِ \lr{\texttt{build\_source\_params(source)}} از یک \lr{\texttt{OpticalSource}} معتبر ساخته می‌شود. یخ‌زده بودنِ این کلاس تضمین می‌کند که پارامترهای منبع در طولِ شبیه‌سازی ثابت بمانند و از رقابتِ \lr{race condition} در شبیه‌سازی‌های چندرشته‌ای جلوگیری شود. این الگو، مشابهِ \lr{\texttt{ChannelSimParams}} و \lr{\texttt{DetectorSimParams}} است که در ماژول‌های دیگر فریم‌ورک پیاده‌سازی شده‌اند.

\subsection{کلاس \lr{\texttt{MZMConfig}}: مودولاتور ماخ-زندر}

این کلاس با \lr{\texttt{@dataclass(slots=True)}} در ماژولِ جداگانه‌ی \lr{\texttt{modulators.py}} تعریف می‌شود و چهار فیلد دارد:
\begin{itemize}
	\item \lr{\texttt{v\_pi: float = 3.5}}: ولتاژ نیم‌موج به ولت.
	\item \lr{\texttt{phi\_bias: float = np.pi / 2.0}}: فاز بایاس.
	\item \lr{\texttt{max\_intensity\_mu: float = 1.0}}: حداکثر میانگین فوتون.
	\item \lr{\texttt{extinction\_ratio\_db: float | None = None}}: نسبت خاموشی به دسی‌بل.
\end{itemize}

متدِ \lr{\texttt{validate()}} سه شرط را بررسی می‌کند: $V_\pi > 0$، $\mu_{\max} \geq 0$ و $\text{ER}_{\text{dB}} \geq 0$. متدِ \lr{\texttt{leakage\_mu()}} میانگین فوتونِ نشتی را از رابطه‌ی \eqref{eq:mzm-leakage} محاسبه می‌کند. متدِ \lr{\texttt{apply(target\_mus, voltage\_noise\_std, rng)}} عملیاتِ اصلیِ مودولاتور را انجام می‌دهد که در بخش بعد به‌تفصیل بررسی می‌شود.

\subsection{کلاس \lr{\texttt{\_BoundedLRUCache}}: کشِ LRU با اندازه‌ی محدود}

این کلاس از \lr{\texttt{OrderedDict}} ارث می‌برد و یک کشِ \lr{LRU} \lr{(Least Recently Used)} با اندازه‌ی حداکثر \lr{\texttt{ADVERSARIAL\_CACHE\_MAX\_SIZE = 10000}} پیاده‌سازی می‌کند. این کش برای نگهداریِ ضرایبِ fluctuationِ بلوک‌های حریصانه استفاده می‌شود تا همان بلوک فیزیکی همواره همان ضریب را دریافت کند، مستقل از نحوه‌ی chunkingِ شبیه‌سازی:

\begin{LTR}
	\begin{lstlisting}[language=Python]
		class _BoundedLRUCache(OrderedDict):
		def __init__(self, max_size: int = ADVERSARIAL_CACHE_MAX_SIZE):
		super().__init__()
		self._max_size = max_size
		self._lock = threading.Lock()  # thread-safety
		
		def __setitem__(self, key, value):
		with self._lock:
		if key in self:
		self.move_to_end(key)
		super().__setitem__(key, value)
		if len(self) > self._max_size:
		self.popitem(last=False)  # evict oldest
	\end{lstlisting}
\end{LTR}

استفاده از \lr{\texttt{threading.Lock}} برای thread-safety در شبیه‌سازی‌های چندرشته‌ای ضروری است. بدون این قفل، دسترسیِ همزمان به کش می‌تواند منجر به فساد داده شود. الگوریتم \lr{LRU}: وقتی یک کلید جدید اضافه می‌شود و اندازه از حد تجاوز کند، قدیمی‌ترین ورودی \lr{(\texttt{last=False})} حذف می‌گردد.

\subsection{ثابت‌های عددی ماژول-محلی}

این ماژول چندین ثابتِ عددی ماژول-محلی تعریف می‌کند که به‌دقت انتخاب شده‌اند و نباید از لایه‌های دیگر فریم‌ورک بازاستفاده شوند. جدول زیر مهم‌ترین این ثابت‌ها را فهرست می‌کند:

\begin{itemize}
	\item \lr{\texttt{SCM\_SMALL\_ANGLE\_LIMIT = 1e-3}}: آستانه‌ی تقریب زاویه کوچک برای $J_1(m)$. خطای نسبی در این آستانه $\sim 2.5 \times 10^{-7}$ است.
	\item \lr{\texttt{SCM\_POWER\_DIVISOR = 4.0}}: مقسوم‌علیه‌ی تقریبِ خطی $m^2/4$ برای توانِ باند جانبی.
	\item \lr{\texttt{SCM\_NUM\_SIDEBANDS = 1}}: تعداد باندهای جانبیِ اول (۱ برای \lr{SSB}، ۲ برای \lr{DSB}).
	\item \lr{\texttt{MAX\_GAMMA\_SHAPE = 1e12}}: سقفِ ایمن برای پارامتر شکل گاما.
	\item \lr{\texttt{GAUSSIAN\_LARGE\_JITTER\_THRESHOLD = 0.1}}: آستانه‌ی جیتر بزرگ.
	\item \lr{\texttt{SOURCE\_NORMALIZATION\_TOL = 1e-10}}: تلورانسِ نرمال‌سازیِ جمعِ ردیفِ احتمالات.
	\item \lr{\texttt{SOURCE\_DM\_MEAN\_RTOL = 1e-3}}: تلورانسِ نسبیِ بررسیِ سازگاریِ $\mu$.
	\item \lr{\texttt{PN\_TRUNCATION\_DIM = 20}}: بعدِ پیش‌فرضِ برشِ تعداد فوتون.
	\item \lr{\texttt{MAX\_MU = 1e6}}: سقفِ ایمن برای میانگین فوتون.
	\item \lr{\texttt{ADVERSARIAL\_CACHE\_MAX\_SIZE = 10000}}: حداکثر اندازه‌ی کشِ بلوک حریصانه.
	\item \lr{\texttt{PHASE\_RANDOMIZATION\_ASSUMED\_DEFAULT = True}}: پیش‌فرضِ پرچمِ فرضِ تصادفی‌سازی فاز.
	\item \lr{\texttt{SCM\_ENABLED\_DEFAULT = False}}: پیش‌فرضِ پرچمِ فعال‌بودن \lr{SCM}.
\end{itemize}

این ثابت‌ها عمدتاً از اصلاحیه‌های بازبینی‌های چندباره (مثل \lr{F-04}, \lr{F-08}, \lr{F-25}, \lr{F-27}) نشأت می‌گیرند و هر یک پاسخی به یک خطای عددی یا فیزیکی مشخص هستند. مثلاً \lr{\texttt{MAX\_MU}} برای جلوگیری از \lr{\texttt{rng.poisson(np.inf)}} معرفی شده که در غیر این صورت \lr{\texttt{ValueError}} پرتاب می‌کند.

\section{پیاده‌سازی ریاضی و الگوریتمی}
\label{sec:sources-implementation}

\subsection{توابع تحلیلی توزیع تعداد فوتون}

دو تابعِ \lr{\texttt{poisson\_pn\_array}} و \lr{\texttt{thermal\_pn\_array}} توابعِ خالصِ \lr{NumPy} هستند که برای کار در حلقه‌های \lr{Numba} طراحی شده‌اند و هیچ شیء پایتونی نمی‌سازند. تابعِ \lr{\texttt{photon\_number\_distribution}} یک wrapper است که بر اساس نوع آمار، به یکی از این دو تابع dispatch می‌کند:

\begin{LTR}
	\begin{lstlisting}[language=Python]
		def poisson_pn_array(mu: float, max_n: int = PN_TRUNCATION_DIM) -> np.ndarray:
		if mu < 0.0:
		raise ParameterValidationError("mu must be >= 0 ...")
		if mu == 0.0:
		result = np.zeros(max_n, dtype=np.float64)
		result[0] = 1.0
		return result
		result = np.empty(max_n, dtype=np.float64)
		# Numerically stable recurrence: P(0)=exp(-mu), P(n+1)=P(n)*mu/(n+1)
		result[0] = math.exp(-mu)
		for n in range(1, max_n):
		result[n] = result[n - 1] * mu / n
		return result
	\end{lstlisting}
\end{LTR}

خط‌به‌خط:
\begin{itemize}
	\item خطوط ۱-۲: اعتبارسنجی ورودی. $\mu < 0$ نامعتبر است زیرا میانگین فوتون نمی‌تواند منفی باشد.
	\item خطوط ۳-۶: حالت خاص $\mu = 0$. در این حالت $P(0) = 1$ و $P(n > 0) = 0$ که معادل حالت خلأ است.
	\item خط ۷: تخصیص آرایه‌ی \lr{float64} با طولِ \lr{\texttt{max\_n}}.
	\item خط ۹: محاسبه‌ی $P(0) = e^{-\mu}$ با \lr{\texttt{math.exp}} که برای $\mu < 10$ پایدار است.
	\item خطوط ۱۰-۱۱: اعمال بازگشتِ \eqref{eq:poisson-recurrence}. هر جمله از جمله‌ی قبلی با یک ضرب و یک تقسیم به‌دست می‌آید که از سرریز فاکتوریل جلوگیری می‌کند.
\end{itemize}

تابعِ \lr{\texttt{thermal\_pn\_array}} الگوی مشابهی دارد اما با نسبتِ $\mu/(1+\mu)$ که همواره کوچک‌تر از ۱ است:

\begin{LTR}
	\begin{lstlisting}[language=Python]
		def thermal_pn_array(mu: float, max_n: int = PN_TRUNCATION_DIM) -> np.ndarray:
		if mu == 0.0:
		result = np.zeros(max_n, dtype=np.float64)
		result[0] = 1.0
		return result
		result = np.empty(max_n, dtype=np.float64)
		ratio = mu / (1.0 + mu)  # P(n+1) / P(n), always < 1
		result[0] = 1.0 / (1.0 + mu)
		for n in range(1, max_n):
		result[n] = result[n - 1] * ratio
		return result
	\end{lstlisting}
\end{LTR}

استفاده از نسبتِ پیش‌محاسبه‌شده‌ی $\mu/(1+\mu)$ به‌جای محاسبه‌ی مستقیمِ $\mu^n/(1+\mu)^{n+1}$ هم از نظر سرعت (یک ضرب به‌جای توان و فاکتوریل) و هم از نظر پایداریِ عددی (نسبت همواره کوچک‌تر از ۱) برتری دارد. این الگوی بازگشتی در تمامِ توابع تحلیلیِ این فریم‌ورک برای توزیع‌های آماری استفاده می‌شود.

\subsection{متد \lr{\texttt{\_scm\_scaling\_factor}}: ضریب مقیاس \lr{SCM}}

این متد، ضریبِ مقیاسِ باند جانبیِ \lr{SCM} را محاسبه می‌کند. پیش‌اولویتِ اعمال به این ترتیب است: (۱) اگر \lr{\texttt{scm\_enabled}} خاموش باشد، ضریب $1.0$؛ (۲) اگر \lr{\texttt{use\_linear\_modulation\_approximation}} روشن باشد، تقریبِ خطیِ اجباری؛ (۳) اگر \lr{\texttt{use\_small\_angle\_approximation}} روشن و $m \leq 10^{-3}$، تقریبِ زاویه کوچک؛ (۴) در غیر این صورت، فرمولِ دقیقِ بسل.

\begin{LTR}
	\begin{lstlisting}[language=Python]
		def _scm_scaling_factor(self) -> float:
		if not getattr(self, 'scm_enabled', SCM_ENABLED_DEFAULT):
		return 1.0
		m = self.modulation_index
		nsb = self.scm_num_sidebands
		if self.use_linear_modulation_approximation:
		return float(nsb * (m ** 2) / SCM_POWER_DIVISOR)
		if self.use_small_angle_approximation and m <= SCM_SMALL_ANGLE_LIMIT:
		return float(nsb * (m ** 2) / SCM_POWER_DIVISOR)
		return float(nsb * (jv(1, m)) ** 2)
	\end{lstlisting}
\end{LTR}

خط‌به‌خط:
\begin{itemize}
	\item خطوط ۱-۳: بررسی \lr{\texttt{scm\_enabled}}. استفاده از \lr{\texttt{getattr}} با پیش‌فرض، سازگاریِ پسین را با نسخه‌های قدیمی که این پرچم را نداشتند تضمین می‌کند.
	\item خط ۴: خواندن شاخصِ مدولاسیون $m$.
	\item خط ۵: خواندن \lr{\texttt{scm\_num\_sidebands}} که در هر نمونه به‌صورت مستقل ذخیره می‌شود (اصلاحیه‌ی \lr{I.4} از بازبینیِ ۷ام برای thread-safety).
	\item خطوط ۶-۷: تقریبِ خطیِ اجباری. این شاخه برای بهینه‌سازی‌های عددی که نیاز به گرادیانِ تحلیلی دارند، معرفی شده است.
	\item خطوط ۸-۹: تقریبِ زاویه کوچک. خطای نسبی در آستانه‌ی $m = 10^{-3}$ حدود $2.5 \times 10^{-7}$ است که برای بهینه‌سازیِ مبتنی بر گرادیان کافی است.
	\item خط ۱۰: فرمولِ دقیق با تابع بسل $J_1(m)$ از \lr{\texttt{scipy.special.jv}}.
\end{itemize}

نکته‌ی مهم درباره‌ی ناپیوستگیِ معنایی در $m = 0$: در نسخه‌های قدیمی، آستانه‌ی $m \leq 10^{-9}$ برای ``غیرفعال‌سازی \lr{SCM}'' استفاده می‌شد که یک ناپیوستگیِ عددیِ ۳۰-مرتبه‌ای ایجاد می‌کرد. اصلاحیه‌ی \lr{F-03} این ناپیوستگی را با جدا کردنِ پرچمِ فیزیکی \lr{\texttt{scm\_enabled}} از پارامترِ پیوسته‌ی $m$ برطرف کرد.

\subsection{متد \lr{\texttt{\_calculate\_effective\_mus}}: اعمال ناکامل‌های کلاسیک}

این متد، تمامِ ناکامل‌های کلاسیکِ شدت را به ترتیب فیزیکیِ درست روی میانگین فوتونِ پایه اعمال می‌کند. ترتیبِ عملیات به این صورت است:

\begin{LTR}
	\begin{lstlisting}[language=Python]
		def _calculate_effective_mus(self, base_mus, rng, pulse_positions=None):
		# Step 1: Intensity jitter (laser fluctuation, before modulation).
		effective_mus = self._apply_intensity_jitter(base_mus, rng, pulse_positions)
		# Step 2: MZM attenuation (with electrical-driver voltage noise).
		effective_mus = self.mzm.apply(effective_mus, self._driver_voltage_noise_std, rng)
		# Step 3: SCM sideband factor.
		effective_mus = effective_mus * self._scm_scaling_factor()
		# Step 4: Channel split.
		if self.N_channels > 1:
		effective_mus = effective_mus / float(self.N_channels)
		return effective_mus
	\end{lstlisting}
\end{LTR}

ترتیبِ فیزیکیِ این مراحل بر اساسِ فیزیکِ منبع است: جیترِ شدت از لیزرِ نیمه‌رسانا نشأت می‌گیرد و قبل از مودولاسیونِ خارجی اعمال می‌شود. سپس \lr{MZM} با ولتاژِ درایور، شدت را مدوله می‌کند. \lr{SCM} در حوزه‌ی فرکانس بعد از مدولاسیونِ شدت اعمال می‌شود. در نهایت، تقسیمِ کانال در مرحله‌ی چندگانه‌سازیِ نوری انجام می‌گیرد. برای مدل‌های کاملاً ضربی (مثل \lr{MZM} خطی و \lr{SCM} ضربی)، این ترتیب از نظر ریاضی معادلِ هر ترتیب دیگری است، اما برای مدل‌های غیرخطی (مثل \lr{MZM} با نشتی یا نویز افزودنی)، ترتیب مهم است.

\subsection{متد \lr{\texttt{\_apply\_intensity\_jitter}}: مدل‌های جیتر}

این متد، مدل‌های مختلفِ جیترِ شدت را پیاده‌سازی می‌کند. شاخه‌ی \lr{\texttt{RANDOM\_GAUSSIAN}} با تصحیحِ بایاس:

\begin{LTR}
	\begin{lstlisting}[language=Python]
		if self.error_model == SourceErrorModel.RANDOM_GAUSSIAN:
		jitter = float(self.intensity_jitter)
		z = rng.standard_normal(size=mus.shape)  # ALWAYS N draws
		factor = 1.0 + jitter * z
		mus = mus * np.abs(factor)  # 0 * |factor| = 0 for vacuum
		if jitter > NUMERIC_ABS_TOL:
		inv_sigma = 1.0 / jitter
		bias_factor = (
		jitter * np.sqrt(2.0 / np.pi) * np.exp(-0.5 * inv_sigma * inv_sigma)
		+ scipy_erf(inv_sigma / np.sqrt(2.0))
		)
		mus = mus / bias_factor  # mean-preserving correction
		return mus
	\end{lstlisting}
\end{LTR}

نکته‌ی کلیدی: طبقِ اصلاحیه‌ی \lr{F-10} از بازبینیِ ۸ام، همواره \lr{N} نمونه از \lr{RNG} مصرف می‌شود حتی برای پالس‌های خلأ ($\mu = 0$). این کار یکنواختیِ جریانِ \lr{RNG} را در سراسرِ اسکن‌های پارامتری حفظ می‌کند: اگر تعداد نمونه‌های مصرفی به مقدارِ $\mu$ وابسته باشد، دو شبیه‌سازی با پیکربندی‌های مختلفِ دکوی، جریان‌های \lr{RNG} متفاوتی خواهند داشت و قابلیت بازتولید از بین می‌رود.

شاخه‌ی \lr{\texttt{ADVERSARIAL\_BLOCK}} پیچیده‌تر است و از یک کشِ \lr{LRU} پایدار برای حفظ همبستگیِ بلوک در سراسرِ فراخوانی‌ها استفاده می‌کند. پارامترِ $\sigma^2_{\ln} = \ln(1 + \sigma^2)$ به‌گونه‌ای انتخاب شده که توزیع لگاریتم-نرمال میانگین ۱ و واریانس $\sigma^2$ داشته باشد:

\begin{LTR}
	\begin{lstlisting}[language=Python]
		if self.error_model == SourceErrorModel.ADVERSARIAL_BLOCK:
		block_size = int(self.adversarial_block_size)
		n = len(mus)
		sigma_sq = np.log(1.0 + self.intensity_jitter ** 2)
		if pulse_positions is not None:
		pulse_positions = np.asarray(pulse_positions, dtype=np.int64)
		block_ids = pulse_positions // block_size
		# ... lookup in _adversarial_block_factor_cache ...
		fresh = rng.lognormal(mean=-0.5 * sigma_sq,
		sigma=np.sqrt(sigma_sq),
		size=len(uncached_block_ids))
		# ... cache and apply ...
		return mus * factors
	\end{lstlisting}
\end{LTR}

وقتی \lr{\texttt{pulse\_positions}} ارائه شود (روش توصیه‌شده)، شناسه‌ی بلوک مستقیماً از $\text{position} // B$ محاسبه می‌شود و کش پایدار، همان بلوک فیزیکی همواره همان ضریب را دریافت می‌کند. این اصلاحیه (بازبینیِ ۷ام، مسئله‌ی \lr{I.1}) یک خطای همبستگی را برطرف کرد که در آن، فراخوانی‌های غیرمتوالی به بلوک‌های یکسان، ضرایبِ مستقل دریافت می‌کردند.

\subsection{متد \lr{\texttt{generate\_photons}}: تولید نهایی فوتون}

این متد، نقطه‌ی ورودیِ اصلیِ شبیه‌سازی است که از منبع، تعداد فوتون تولید می‌کند. الگوریتم:

\begin{LTR}
	\begin{lstlisting}[language=Python]
		def generate_photons(self, alice_pulse_indices, rng, num_samples=1, *,
		pulse_positions=None):
		if alice_pulse_indices is None:
		indices = self.sample_pulse_indices(rng, num_samples)
		is_scalar = False
		else:
		indices, is_scalar = self._validate_pulse_indices(alice_pulse_indices)
		if len(indices) == 0:
		return np.array([], dtype=np.int64)
		pp = None
		if pulse_positions is not None:
		pp = np.asarray(pulse_positions, dtype=np.int64)
		selected_mus = self._calculate_effective_mus(
		self._base_mus_cache[indices], rng, pulse_positions=pp
		)
		photons = self._sample_poisson_or_thermal(selected_mus, rng)
		# Emission-failure thinning
		p_emit = clamp_probability(self.ideal_emission_probability)
		n_attempted = len(indices)
		emit_mask = rng.random(size=n_attempted) < p_emit
		photons = np.where(emit_mask, photons, 0).astype(np.int64, copy=False)
		n_emitted = int(np.count_nonzero(emit_mask))
		self.update_internal_tallies(
		num_pulses_generated=n_emitted,
		num_pulses_attempted=n_attempted,
		)
		return int(photons[0]) if is_scalar else photons
	\end{lstlisting}
\end{LTR}

مراحل کلیدی:
\begin{itemize}
	\item خطوط ۲-۷: اگر شاخص پالس \lr{None} باشد، \lr{\texttt{num\_samples}} شاخص به‌صورت تصادفی نمونه‌برداری می‌شود. در غیر این صورت، شاخص‌های ارائه‌شده اعتبارسنجی می‌شوند.
	\item خطوط ۸-۹: مدیریتِ حالتِ خالی (\lr{num\_samples = 0}) که آرایه‌ی خالی برمی‌گرداند.
	\item خطوط ۱۰-۱۲: تبدیل \lr{\texttt{pulse\_positions}} به آرایه برای مدل بلوک حریصانه.
	\item خطوط ۱۳-۱۵: محاسبه‌ی $\mu_{\text{eff}}$ از $\mu_{\text{base}}$ با اعمالِ تمامِ ناکامل‌های کلاسیک.
	\item خط ۱۶: نمونه‌برداریِ تعداد فوتون از پواسون یا گرمایی.
	\item خطوط ۱۷-۲۰: اعمالِ شکست گسیل به‌صورت thinning. طبقِ اصلاحیه‌ی \lr{F-41} و \lr{II.3}، فراخوانیِ \lr{\texttt{rng.random}} همواره انجام می‌شود تا جریان \lr{RNG} در سراسرِ تغییراتِ $p_{\text{emit}}$ یکنواخت بماند.
	\item خطوط ۲۱-۲۵: به‌روزرسانیِ شمارنده‌های \lr{security\_metadata} برای تحلیل‌های پسین.
\end{itemize}

\subsection{متد \lr{\texttt{MZMConfig.apply}}: عملیات مودولاتور}

این متد در فایل \lr{\texttt{modulators.py}} پیاده‌سازی شده و به‌صورت زیر کار می‌کند:

\begin{LTR}
	\begin{lstlisting}[language=Python]
		def apply(self, target_mus, voltage_noise_std, rng):
		self.validate()
		i_min = self.leakage_mu()
		i_max = self.max_intensity_mu
		delta_i = i_max - i_min
		target = np.clip(np.asarray(target_mus, dtype=np.float64), i_min, i_max)
		if delta_i < NUMERIC_ABS_TOL:
		normalized = np.zeros_like(target)
		else:
		normalized = np.clip((target - i_min) / delta_i, 0.0, 1.0)
		theta_target = np.arccos(np.sqrt(normalized))
		v_signal = (2.0 * self.v_pi / np.pi) * (theta_target - self.phi_bias)
		if voltage_noise_std > 0.0:
		v_noise = rng.normal(loc=0.0, scale=voltage_noise_std, size=target.shape)
		else:
		v_noise = 0.0
		v_total = v_signal + v_noise
		noisy_phase = (np.pi * v_total) / (2.0 * self.v_pi) + self.phi_bias
		return i_min + delta_i * (np.cos(noisy_phase) ** 2)
	\end{lstlisting}
\end{LTR}

این پیاده‌سازی، الگوریتمِ معکوس‌سازیِ تابع انتقالِ \lr{MZM} را برای یافتنِ ولتاژِ لازم، و سپس اعمالِ نویز روی آن ولتاژ، و در نهایت اعمالِ تابع انتقال به ولتاژِ نویزی را انجام می‌دهد. این روشِ سه‌مرحله‌ای، فیزیکِ واقعیِ مودولاتور را بازتولید می‌کند: ابتدا ولتاژِ سیگنال از روی شدتِ هدف محاسبه می‌شود (معکوسِ تابع انتقال از رابطه‌ی \eqref{eq:mzm-transfer})، سپس نویزِ الکتریکی به آن افزوده می‌شود، و در نهایت خروجیِ واقعی از تابع انتقالِ $\cos^2(\phi)$ به‌دست می‌آید. این رویکرد، برخلافِ افزودنِ مستقیمِ نویز به شدت، اثرِ غیرخطیِ \lr{MZM} را به‌درستی مدل می‌کند.

خط‌به‌خط:
\begin{itemize}
	\item خط ۲: اعتبارسنجی پارامترهای مودولاتور.
	\item خطوط ۳-۵: محاسبه‌ی $\mu_{\min}$ (نشتی)، $\mu_{\max}$ و $\Delta = \mu_{\max} - \mu_{\min}$.
	\item خط ۶: clip کردنِ شدتِ هدف به بازه‌ی $[\mu_{\min}, \mu_{\max}]$. شدت‌های خارج از این بازه غیرفیزیکی هستند.
	\item خطوط ۷-۱۰: نرمال‌سازیِ شدت به بازه‌ی $[0, 1]$. اگر $\Delta$ بسیار کوچک باشد (مثلاً وقتی $\mu_{\max} = \mu_{\min}$)، همه‌ی شدت‌ها به صفر نرمال می‌شوند.
	\item خط ۱۱: معکوسِ تابع انتقال. از $\cos^2(\theta) = \text{normalized}$ نتیجه می‌گیریم $\theta = \arccos(\sqrt{\text{normalized}})$.
	\item خط ۱۲: محاسبه‌ی ولتاژِ سیگنال از معکوسِ رابطه‌ی \eqref{eq:mzm-phase}: $V = (2 V_\pi / \pi)(\theta - \phi_{\text{bias}})$.
	\item خطوط ۱۳-۱۶: نمونه‌برداریِ نویز گوسی با انحراف معیار $\sigma_V$.
	\item خط ۱۷: جمعِ ولتاژِ سیگنال و نویز.
	\item خط ۱۸: محاسبه‌ی فازِ نویزی از رابطه‌ی \eqref{eq:mzm-phase}.
	\item خط ۱۹: اعمالِ تابع انتقالِ $\cos^2(\phi)$ و افزودنِ کفِ نشتی.
\end{itemize}

\subsection{متد \lr{\texttt{\_build\_rho\_tensor}}: ساخت تانسور ماتریس چگالی}

این متد در \lr{\texttt{DensityMatrixSource}}، ماتریس‌های چگالیِ با ابعادِ ناهمگن را در یک تانسورِ سه‌بعدی با بعدِ یکسان embed می‌کند. هر ماتریس در گوشه‌ی بالا-چپِ یک ماتریس صفرِ $\text{max\_dim} \times \text{max\_dim}$ قرار می‌گیرد. این zero-padding فیزیکی معنادار است: حالت‌های فوک با احتمال صفر را نشان می‌دهد. پس از \texttt{padding}، قطر نرمال می‌شود تا جمع برابر ۱ باشد.

\subsection{متد \lr{\texttt{\_apply\_dm\_loss\_processes}}: اعمال اتلاف روی ماتریس چگالی}

این متد، دو فرآیندِ اتلافِ قابل‌بیان به‌صورت قطری (تقسیم کانال و شکست گسیل) را در زمانِ ساخت اعمال می‌کند:

\begin{itemize}
	\item \textbf{تقسیم کانال}: تبدیلِ برنولی-دوجمله‌ای از رابطه‌ی \eqref{eq:binomial-loss-channel-split} با $\eta = 1/N$.
	\item \textbf{شکست گسیل}: مخلوط با خلأ از رابطه‌های \eqref{eq:emission-failure-p0} و \eqref{eq:emission-failure-pk}.
\end{itemize}

این دو فرآیند در زمانِ ساخت روی \lr{\texttt{\_number\_probs\_table}} اعمال می‌شوند تا جدولِ از ابتدا روی وضعیتِ پس از اتلاف تنظیم شود. سایر ناکامل‌های کلاسیک (مثل \lr{MZM}، \lr{SCM} و جیتر) که به‌صورت غیرقطری عمل می‌کنند، در \lr{\texttt{DensityMatrixSource}} نادیده گرفته می‌شوند و در صورت پیکربندی غیرپیش‌فرض، هشدار ثبت می‌شود.

\section{ارسال و دریافت با سایر ماژول‌ها}
\label{sec:sources-interfacing}

\subsection{ورودی‌های ماژول}

ماژولِ \lr{\texttt{sources.py}} از طریق کلاسِ \lr{\texttt{OpticalSourceConfig}} (که در \lr{\texttt{datatypes.py}} تعریف شده) پیکربندی می‌شود. این پیکربندی شامل موارد زیر است:
\begin{itemize}
	\item \lr{\texttt{source\_rate}}: نرخ تکرار منبع به هرتز.
	\item \lr{\texttt{pulse\_configs}}: تاپلِ پیکربندیِ پالس‌ها (سیگنال و دکوی‌ها).
	\item \lr{\texttt{statistics\_type}}: نوع آمار (\lr{POISSON} یا \lr{THERMAL}).
	\item \lr{\texttt{error\_model}}: مدل جیتر شدت.
	\item \lr{\texttt{intensity\_jitter}}: جیتر نسبی.
	\item \lr{\texttt{modulation\_index}}: شاخص مدولاسیون \lr{SCM}.
	\item \lr{\texttt{N\_channels}}: تعداد کانال‌ها.
	\item \lr{\texttt{ideal\_emission\_probability}}: احتمال گسیل.
	\item \lr{\texttt{mzm}}: شیء مودولاتور \lr{MZM} از نوع \lr{\texttt{MZMConfig}}.
	\item \lr{\texttt{electrical\_noise}}: شیء نویز الکتریکی از \lr{\texttt{noise\_models.py}}.
\end{itemize}

منابع متعددی برای ساختِ این پیکربندی وجود دارد: \lr{\texttt{from\_dict}} (از دیکشنری سریال‌شده)، \lr{\texttt{from\_protocol\_parameters}} (از \lr{\texttt{ProtocolParameters}})، \lr{\texttt{from\_pulse\_ensemble}} (از \lr{\texttt{PulseEnsembleConfig}})، و \lr{\texttt{from\_intensity\_config}} (از \lr{\texttt{IntensityConfig}}). هر یک از این \lr{\texttt{factory method‌}}ها، coerce و اعتبارسنجیِ مختصر به خود را انجام می‌دهند.

برای \lr{\texttt{DensityMatrixSource}}، ورودیِ اضافی \lr{\texttt{density\_matrices}} (نگاشت نام پالس به ماتریس چگالی) نیز لازم است. این ماتریس‌ها باید در پایه‌ی فوک باشند و شروطِ هرمیتی، مثبت‌معین و یک‌بودنِ رد را برآورده کنند.

\subsection{خروجی‌های ماژول}

خروجیِ اصلیِ این ماژول، تعداد فوتون تولیدشده به‌ازای پالس است که از طریقِ متدِ \lr{\texttt{generate\_photons}} در دسترس قرار می‌گیرد. این خروجی به‌صورت یک آرایه‌ی \lr{int64} (یا یک اسکالر در حالتِ تک‌پالس) برگردانده می‌شود و مستقیماً به ماژولِ کانال \lr{\texttt{channel.py}} ارسال می‌شود تا ضریبِ انتقالِ فیبر روی آن اعمال شود. به‌علاوه، متدِ \lr{\texttt{calculate\_effective\_mus}} میانگین فوتون مؤثر را (بدون نمونه‌برداری) برمی‌گرداند که برای تحلیل‌های تحلیلیِ نرخ کلید استفاده می‌شود.

خروجیِ تحلیلیِ دیگر، توزیعِ احتمالِ تعداد فوتون $P(n)$ است که از طریقِ متدِ \lr{\texttt{number\_probabilities\_for\_pulse}} برای هر پالس قابل دسترسی است. این خروجی در ماژول‌های تحلیل امنیتی (مانند تخمینِ $Y_1$ و $e_1$ در \lr{Decoy-State}) استفاده می‌شود. توابعِ ماژول-سطحی \lr{\texttt{poisson\_pn\_array}} و \lr{\texttt{thermal\_pn\_array}} نیز برای محاسباتِ تحلیلیِ مستقل (بدون نیاز به شیء منبع) در دسترس هستند.

بسته‌ی \lr{\texttt{SourceSimParams}} نیز از طریقِ \lr{\texttt{build\_source\_sim\_params}} به ماژول \lr{\texttt{main\_optimized.py}} ارسال می‌شود تا پارامترهای اجراییِ شبیه‌سازی (مثل دوره‌ی پالس، نرخ منبع، نوع آمار) در اختیارِ موتورِ شبیه‌سازی قرار گیرد. این بسته‌ی یخ‌زده تضمین می‌کند که پارامترها در طولِ شبیه‌سازی ثابت بمانند.

\subsection{وابستگی‌های داخلی}

این ماژول به چندین ماژولِ دیگر وابسته است:
\begin{itemize}
	\item \lr{\texttt{constants.py}}: برای ثابت‌های عددی مثل \lr{\texttt{NUMERIC\_ABS\_TOL}} و \lr{\texttt{MAX\_SEED\_INT\_PCG64}} و توابعِ کمکی مثل \lr{\texttt{clamp\_probability}} و \lr{\texttt{is\_close}}.
	\item \lr{\texttt{datatypes.py}}: برای تعاریفِ \lr{\texttt{OpticalSourceConfig}}، \lr{\texttt{PulseTypeConfig}}، \lr{\texttt{SourceErrorModel}} و \lr{\texttt{SourceStatisticsType}}.
	\item \lr{\texttt{exceptions.py}}: برای \lr{\texttt{ParameterValidationError}}.
	\item \lr{\texttt{security\_metadata.py}}: برای \lr{\texttt{SourceSecurityMetadata}} که شمارنده‌های امنیتی را نگه می‌دارد.
	\item \lr{\texttt{type\_defs.py}}: برای \lr{\texttt{RNGType}} (نوعِ \lr{NumPy Generator}).
	\item \lr{\texttt{utils/utils.py}}: برای \lr{\texttt{sanitize\_for\_serialization}}.
	\item \lr{\texttt{scipy.special}}: برای توابع بسل \lr{\texttt{jv}} و تابع خطا \lr{\texttt{erf}}.
\end{itemize}

این وابستگی‌ها، یک گرافِ وابستگیِ خطی از لایه‌ی پایین (ثابت‌ها و تایپ‌ها) به لایه‌ی بالای فیزیکی (منابع) تشکیل می‌دهند. ماژول \lr{\texttt{modulators.py}} به‌صورت مستقل عمل می‌کند و تنها به \lr{\texttt{constants.py}} و \lr{\texttt{exceptions.py}} وابسته است، که به \lr{\texttt{OpticalSource}} اجازه می‌دهد آن را به‌عنوان یک کامپوننتِ قابل‌تعویض استفاده کند.

\subsection{جریانِ داده در شبیه‌سازیِ کامل}

جریانِ داده از منبع تا آشکارساز به‌صورت زیر است:
\begin{enumerate}
	\item کاربر، \lr{\texttt{OpticalSource}} را از طریقِ \lr{\texttt{from\_dict}} یا \lr{\texttt{from\_protocol\_parameters}} می‌سازد.
	\item \lr{\texttt{build\_source\_sim\_params}} بسته‌ی اجرایی را برای موتور شبیه‌سازی می‌سازد.
	\item در هر تکرارِ شبیه‌سازی، \lr{\texttt{generate\_photons}} فراخوانی می‌شود و آرایه‌ای از تعداد فوتون تولید می‌کند.
	\item این آرایه به \lr{\texttt{FiberChannel.transmittance}} ارسال می‌شود تا ضریبِ انتقال روی هر فوتون اعمال شود (به‌صورت ضرب در $T$ و سپس گردکردن به عدد صحیح).
	\item خروجیِ کانال به ماژولِ آشکارساز \lr{\texttt{detectors.py}} ارسال می‌شود تا کاراییِ آشکارساز، نرخ شمارش تاریک و پدیده‌ی afterpulsing اعمال شود.
	\item در نهایت، رویدادهای آشکارسازی به ماژولِ \lr{post-processing} (تصحیح خطا و تقویت حریم خصوصی) ارسال می‌شوند.
\end{enumerate}

این جریانِ خطی، مدولار بودنِ فریم‌ورک را تضمین می‌کند: هر ماژول مسئولیتِ یک لایه‌ی فیزیکی مشخص را دارد و خروجیِ آن، ورودیِ ماژولِ بعدی است. این معماری، امکانِ جایگزینیِ آسانِ هر ماژول با یک پیاده‌سازیِ متفاوت (مثلاً مدلِ آشکارسازِ پیشرفته‌تر) را بدون تغییر در سایر ماژول‌ها فراهم می‌سازد.

	
	\chapter{مدل‌های آشکارساز و ساختارهای داده}
	\label{ch:detectors}
	\sloppy
\emergencystretch=3em
\hbadness=10000
\vbadness=10000
\sloppy
\emergencystretch=3em
\hbadness=10000
\vbadness=10000

\lstset{
	basicstyle=\ttfamily\footnotesize,
	breakatwhitespace=false,
	breaklines=true,
	captionpos=b,
	keepspaces=true,
	numbers=left,
	numbersep=5pt,
	numberstyle=\tiny\color{gray},
	showspaces=false,
	showstringspaces=false,
	showtabs=false,
	tabsize=2,
	language=Python,
	frame=single,
	framesep=2pt,
	rulecolor=\color{gray!50},
	backgroundcolor=\color{gray!5},
	keywordstyle=\color{blue!70!black}\bfseries,
	commentstyle=\color{green!50!black}\itshape,
	stringstyle=\color{red!60!black},
	emph={np, math, self, rng, DetectorType, DoubleClickPolicy, DeadTimeModel, AfterpulseModel, DetectorTopology, APRescheduleMode, ParameterValidationError, ConfigurationError, DetectionDiagnostics, DetectionResult, SinglePhotonDetector, EntanglingDecoder, DetectorConstants, DetectorSimParams, _DetectorState, CarryOverEvent, PhotonResolvedResult, np.ndarray},
	emphstyle=\color{purple!70!black}\bfseries,
}

\section{نمای کلی و هدف ماژول}

ماژول \lr{\texttt{qkd.detectors}} یکی از بزرگ‌ترین، پیچیده‌ترین و از منظر فیزیکی حساس‌ترین اجزای فریم‌ورک شبیه‌سازی \lr{QKD} است. این ماژول مدل کاملی از آشکارسازهای تک‌فوتونی را برای شبیه‌سازی پروتکل \lr{BB84-Decoy} به‌صورت رویداد-محور \lr{(event-driven)} پیاده‌سازی می‌کند و سه نوع آشکارساز را به‌صورت یکپارچه پشتیبانی می‌نماید: آشکارسازهای آوانش نوری \lr{(SPAD)} در حالت آستانه‌ای، آشکارسازهای ابررسانا نانوسیمی \lr{(SNSPD)} با مدل‌سازی اختصاصی از وابستگی دمایی نرخ شمارش تاریک، و آشکارسازهای تفکیک‌پذیر عدد فوتون \lr{(PNRD)} با خروجی شمارشی. هدف اصلی این ماژول، فراهم‌آوردن یک موتور شبیه‌سازی واقع‌گرایانه است که اثرات پدیده‌های مخرب فیزیکی نظیر زمان مرده \lr{(dead time)}، پالس‌های پسین \lr{(afterpulsing)}، شمارش تاریک \lr{(dark counts)}، لرزش زمانی \lr{(timing jitter)} و پهن‌شدگی کروماتیک \lr{(chromatic dispersion)} را با دقت ریاضی و با وفاداری به ادبیات علمی، به‌ویژه مرجع سنتی \lr{Rogers et al. (2007)}، بازتولید کند.

این ماژول نقش گره اصلی در گراف وابستگی‌های فریم‌ورک را ایفا می‌کند؛ به این معنا که خروجی کانال فیبر نوری (آرایه‌های تعداد فوتون و پیامد ایده‌آل مسیریابی) به‌عنوان ورودی این ماژول عرضه می‌شود و خروجی آن، یعنی آرایه‌های کلیک بولی \lr{D0} و \lr{D1} به‌همراه اشیای تشخیصی \lr{DetectionDiagnostics}، پایه‌ی محاسبات نرخ کلید امن در پروتکل‌های \lr{decoy-state} و \lr{MDI-QKD} را شکل می‌دهد. پیچیدگی این ماژول نه‌تنها از تنوع پدیده‌های فیزیکی مدل‌سازی‌شده برمی‌آید، بلکه از لزوم حفظ سازگاری نتایج میان مسیرهای مختلف شبیه‌سازی (آستانه‌ای در برابر \lr{PNRD}) و لزوم تضمین بازتولیدپذیری آماری \lr{(reproducibility)} برای کارهای وابسته به انتشار علمی ناشی می‌شود.

یک ویژگی برجسته این ماژول، سیستم کنترل نسخه‌بندی حالت داخلی است که با ثابت \texttt{STATE\_VERSION = 19} و فهرست نسخه‌های قدیمی \texttt{\_LEGACY\_STATE\_VERSIONS} مدیریت می‌شود. این سیستم امکان مهاجرت رو به جلوی حالت‌های ذخیره‌شده از نسخه‌های قدیمی‌تر را فراهم می‌سازد و در عین حال با مکانیزم کنترل یکپارچگی هش پیکربندی \lr{(SHA-256 config hash)} و کنترل نسخه کد فیزیک \lr{(\_PHYSICS\_CODE\_VERSION = 10)} از بارگذاری ناسازگار حالت‌های ذخیره‌شده با طرح متفاوت \lr{const\_params} یا معناشناسی متفاوت \lr{AP-reschedule} جلوگیری می‌کند. این طراحی محافظه‌کارانه برای کاربردهای انتشار علمی که در آن‌ها لازم است نتایج یک شبیه‌سازی با همان بذر تصادفی دقیقاً بازتولید شوند، حیاتی است.

ماژول همچنین از یک مکانیزم حالت سخت‌گیرانه \lr{(strict\_mode)} که به‌صورت پیش‌فرض فعال است، بهره می‌برد. این مکانیزم تقریباً هر ساده‌سازی فیزیکی مسئله‌ساز را یا به یک استثنای صریح تبدیل می‌کند یا هشدار می‌دهد، تا از انتشار نتایج نادرست فیزیکی جلوگیری شود. برای مثال، حالت سخت‌گیرانه استفاده از مدل زمان مرده پارالایز \lr{(paralyzable)} در سطح یک \lr{SPAD} منفرد را که بر اساس \lr{Rogers et al.} نادرست است، رد می‌کند و در عوض پارالایزپذیری ظهور-سطح-پایه را تنها در صورت فعال‌سازی سیاست \lr{ROGERS\_2007} شبیه‌سازی می‌نماید. این رویکرد دفاعی، ماژول را برای استفاده در پروژه‌های پایان‌نامه و مقالات ژورنالی که در آن‌ها دقت فیزیکی فراتر از سرعت اجرا اهمیت دارد، مناسب می‌سازد.

\section{مبانی تئوری و فیزیکی}

مدل‌سازی آشکارسازهای تک‌فوتونی در این ماژول بر چند ستون فیزیکی استوار است که در ادامه به تفصیل بسط داده می‌شوند. نخست، فیزیک \lr{SPAD} بر پدیده شکست بهرهاوری \lr{(avalanche breakdown)} در دیودهای معکوس‌بایاس استوار است. وقتی یک فوتون در ناحیه تخلیه جذب می‌شود و یک حامل بار تولید می‌کند، میدان الکتریکی قوی حامل را شتاب می‌دهد و شکست بهرهاوری آغاز می‌شود؛ در نتیجه یک پالس جریان قابل‌اندازه‌گیری ایجاد می‌گردد. در این حالت که به حالت گایگر \lr{(Geiger mode)} معروف است، آشکارساز به هر فوتون با احتمال $\eta$ (بازده آشکارساز) پاسخ می‌دهد، اما پس از هر کلیک نیازمند یک پنجره زمانی برای خاموش‌سازی \lr{(quenching)} و بازیابی به حالت آماده‌به‌کار است که به آن زمان مرده $\tau_d$ می‌گوئیم. این زمان مرده در مدل حاضر از نوع غیر-پارالایز \lr{(non-paralyzable)} است؛ یعنی فوتون‌هایی که در طول پنجره مرده می‌رسند هیچ اثری بر آشکارساز ندارند و پنجره مرده را تمدید نمی‌کنند. این انتخاب بر اساس \lr{Rogers et al. (2007) §3 و §6} صورت گرفته که صراحتاً بیان می‌کنند \lr{SPAD}های منفرد غیر-پارالایز هستند و پارالایزپذیری صرفاً یک ویژگی نوپدید \lr{(emergent)} در سطح جفت آشکارسازهای پایه به‌همراه قانون سیفتنگ \lr{(sifting rule)} است.

دومین ستون فیزیکی، پدیده پالس‌های پسین \lr{(afterpulsing)} است. در \lr{SPAD}های \lr{InGaAs/InP} که در بازه مخابراتی \lr{1550 nm} استفاده می‌شوند، حامل‌های بار در طول شکست بهرهاوری می‌توانند در تله‌های سطحی \lr{(surface trap levels)} به‌دام بیفتند و سپس با یک تأخیر تصادفی آزاد شوند. اگر آزادسازی در زمانی رخ دهد که آشکارساز در حالت آماده‌به‌کار باشد، یک کلیک کاذب \lr{(afterpulse)} تولید می‌شود. مدل زمان تأخیر در این ماژول از دو نوع توزیع پشتیبانی می‌کند: توزیع نمایی پیوسته با میانگین $\tau_{ap}$ (پارامتر \texttt{afterpulse\_lifetime\_ns}) و توزیع هندسی گسسته روی شبکه دوره پالس. در حالت هندسی، پارامتر $p$ به‌صورت زیر از طول عمر به‌دست می‌آید:

\begin{equation}
	p = 1 - \exp\left(-\frac{T}{\tau_{ap}}\right)
\end{equation}

که در آن $T$ دوره پالس \lr{(pulse\_period\_ns)} است. این انتخاب تضمین می‌کند که $\tau_{ap}$ در هر دو حالت به‌عنوان زمان مشخصه توزین تأخیر تفسیر شود. احتمال تولید پالس پسین پس از هر کلیک با $p_{ap}$ (پارامتر \texttt{afterpulse\_prob}) کنترل می‌شود و عمق آبشاری \lr{(cascade depth)} با کران سخت \lr{\texttt{MAX\_AFTERPULSE\_CASCADE = 32}} محدود می‌گردد تا از زنجیره‌های بازگشتی بی‌نهایت جلوگیری شود. وقتی یک پالس پسین در طول پنجره زمان مرده فعال می‌شود، هسته پردازشی آن را به زمانی پس از پایان پنجره مرده برنامه‌ریزی مجدد می‌کند؛ این برنامه‌ریزی مجدد می‌تواند با سیاست \lr{END} (پایان پنجره مرده) یا سیاست \lr{UNIFORM} (توزیع یکنواخت در بازه \lr{[last\_fire + dt, last\_fire + dt + ap\_lifetime]}) انجام شود.

سومین ستون، شمارش تاریک \lr{(dark count)} است. شمارش تاریک در \lr{SPAD}ها به‌دو مکانیسم اصلی برمی‌گردد: نشر تونلی ناشی از تله \lr{(trap-assisted tunneling)} و تولید حرارتی-نمایی \lr{(thermal generation)}. مدل حاضر از تقریب آرنیوس \lr{(Arrhenius)} با انرژی فعال‌سازی $E_a = 0.15$ eV برای بخش دمایی و یک وابستگی درجه دوم به بایاس مازاد \lr{(excess bias)} برای بخش ولتاژی استفاده می‌کند. بایاس مازاد به‌صورت $V_{eb} = V_{bias} - V_{bd}$ تعریف می‌شود، که در آن $V_{bd}$ ولتاژ شکست است. اصلاح \lr{Review-v19 F-01} وابستگی دمایی ولتاژ شکست را با ضریب $50$ mV/K برای \lr{InGaAs} مدل می‌کند:

\begin{equation}
	V_{bd}^{op} = V_{bd}^{ref} + \alpha_{bd}\,(T_{op} - T_{ref})
\end{equation}

که در آن $\alpha_{bd}$ پارامتر \texttt{breakdown\_temp\_coeff\_mv\_per\_k} است. این اصلاح از یک باگ خاموش می‌کشد که پیش از آن همان $V_{bd}$ برای هر دو نقطه عملیاتی و مرجع استفاده می‌شد و برای اختلاف دمایی $50$ K مقادیر اشتباه $\text{voltage\_factor} = 1.0$ به‌جای $0.25$ تولید می‌کرد. مقیاس‌کننده نهایی نرخ تاریک \lr{SPAD} به‌صورت زیر به‌دست می‌آید:

\begin{equation}
	\mathcal{S}_{SPAD}(T, V) = \exp\!\left[\frac{E_a}{k_B\,T_{ref}} - \frac{E_a}{k_B\,T}\right] \cdot \left(\frac{V_{eb}^{op}}{V_{eb}^{ref}}\right)^{2}
\end{equation}

که در آن $k_B$ ثابت بولتزمن است. باید تأکید کرد که این مدل صرفاً یک تقریب کیفی است و برای کارهای انتشار علمی توصیه می‌شود نرخ تاریک اندازه‌گیری‌شده در نقطه عملیاتی مستقیماً تأمین شود (با قرار دادن $T_{ref}=T_{op}$ تا مقیاس‌کننده به $1.0$ فروکد کند).

برای \lr{SNSPD}ها، مکانیسم شمارش تاریک متفاوت است: آشکارساز در دماهای بسیار پایین (نظیر $2$--$4$ \lr{K}) کار می‌کند و منبع غالب نویز، تابش جسم سیاه \lr{(black-body radiation)} از محیط اطراف است که با طیف پلانک توصیف می‌شود. مدل حاضر شار فوتون جسم سیاه را در فرکانس‌های بالاتر از فرکانس برش $\nu_{cut} = c/\lambda_{cut}$ (با $\lambda_{cut}\approx 1550$ nm) به‌صورت زیر محاسبه می‌کند:

\begin{equation}
	\Phi(T) = \frac{2(k_B T)^3}{c^2 h^3} \sum_{n=1}^{\infty} e^{-nx}\left(\frac{x^2}{n} + \frac{2x}{n^2} + \frac{2}{n^3}\right)
\end{equation}

که در آن $x = h\nu_{cut}/(k_B T)$ است. این سری با حد به $n \to \infty$ همگراست و در عمل با کران \lr{\texttt{SNSPD\_PLANCK\_MAX\_TERMS = 200}} و تلورانس نسبی $10^{-20}$ متوقف می‌شود. مقیاس‌کننده دمای \lr{SNSPD} نسبت $\Phi(T)/\Phi(T_{ref})$ است؛ به‌علاوه یک فاکتور تجربی وابسته به جریان بایاس $f(x) = \exp(-\alpha/(1-x))$ (با $x = I_{bias}/I_{sw}$) اعمال می‌شود که به‌صورت کیفی رشد نرخ تاریک را هنگام نزدیک شدن جریان بایاس به جریان سوئیچینگ بازتولید می‌کند. این مقیاس‌کننده صریحاً به‌عنوان کیفی-تنها \lr{(QUALITATIVE ONLY)} مستند شده و در حالت سخت‌گیرانه اگر انحراف آن از $1.0$ بیش از $5\%$ باشد، استثنا صادر می‌کند.

چهارمین ستون، لرزش زمانی \lr{(timing jitter)} و پهن‌شدگی کروماتیک \lr{(chromatic dispersion)} است. لرزش زمانی به‌منظور بازتولید عدم قطعیت در زمان‌بندی الکترونیک آشکارساز و پهن‌شدگی پالس نوری در فیبر، در مرحله بن کردن \lr{(binning)} به هر رویداد کلیک اعمال می‌شود. مدل از توزیع گوسی با عرض نیم-ماکزیمم $\text{FWHM} = \text{jitter\_fwhm\_ns}$ استفاده می‌کند که به انحراف معیار از رابطه آشنای زیر تبدیل می‌شود:

\begin{equation}
	\sigma = \frac{\text{FWHM}}{2\sqrt{2\ln 2}}
\end{equation}

نکته فیزیکی مهم این است که پهن‌شدگی کروماتیک صرفاً به رویدادهای کلیک با منشأ سیگنال اعمال می‌شود و نه به شمارش‌های تاریک و پالس‌های پسین، زیرا این دو پدیده درون آشکارساز تولید می‌شوند و از فیبر عبور نمی‌کنند. این تفکیک با پرچم \texttt{is\_sig} در هسته پردازشی پیگیری می‌شود (اصلاح \lr{Review-v12 F3}).

پنجمین ستون، سیاست سيفتنگ دو-کلیک \lr{(double-click sifting policy)} است که با \lr{Enum} به نام \lr{DoubleClickPolicy} پیاده‌سازی شده و چهار حالت دارد: \lr{DISCARD} (رد هر دو کلیک هم‌زمان)، \lr{RANDOM} (انتخاب تصادفی یکی از دو کلیک)، \lr{KEEP\_BOTH} (نگه‌داشتن هر دو) و \lr{ROGERS\_2007}. حالت \lr{ROGERS\_2007} مطابق مقاله \lr{Rogers et al. (2007)} یک دنباله تشخیص \lr{(detection sequence)} را به‌صورت یک اجرای ماکزیمال از رویدادهای تشخیص روی D0 $\cup$ D1 تعریف می‌کند که در آن رویدادهای متوالی با فاصله کمتر از $\tau_d$ جدا شده‌اند. هر دنباله نهایتاً یک بیت سیفتی \lr{(sifted bit)} (اولین رویداد) تولید می‌کند و سایر رویدادها در همان دنباله دور ریخته می‌شوند و در شمارنده \texttt{rogers\_discarded\_events} ثبت می‌گردند.

\section{تحلیل معماری و ساختار داده‌ها}

معماری این ماژول چندلایه است و از یک مجموعه غنی از \lr{dataclass}ها، \lr{Enum}ها و \lr{NamedTuple}ها تشکیل شده که هر کدام نقش مشخصی در جداسازی نگرانی‌ها \lr{(separation of concerns)} ایفا می‌کنند. در رأس این معماری، کلاس اصلی \texttt{SinglePhotonDetector} قرار دارد که یک \lr{dataclass} با \texttt{slots=True} (اما بدون \texttt{frozen=True}) است تا امکان تغییر حالت درون‌رونده \lr{(in-process state)} را حفظ کند و در عین حال مصرف حافظه را بهینه نماید. این کلاس به‌صورت هدفمند غیر-ایستا \lr{(non-frozen)} طراحی شده زیرا باید حالت درون‌رونده‌ی آشکارساز (مانند زمان آخرین کلیک، زمان \lr{AP} معلق، عمق آبشار، زمان \lr{DC} بعدی و رویدادهای \lr{carry-over}) را در طول شبیه‌سازی به‌روزرسانی کند. با این‌حال، \textit{پیکربندی} آشکارساز (یعنی تمام فیلدهای فیزیکی مانند بازده، نرخ تاریک، زمان مرده و غیره) پس از ساخت از طریق اعتبارسنجی در \texttt{\_\_post\_init\_\_} به‌صورت عملی تغییرناپذیر باقی می‌مانند؛ زیرا هیچ متد عمومی برای تغییر آن‌ها وجود ندارد.

فیلدهای پیکربندی به سه گروه تقسیم می‌شوند. گروه نخست، فیلدهای پایه عملکرد آشکارساز هستند: \texttt{det\_eff\_d0} و \texttt{det\_eff\_d1} (بازده‌های نامتقارن دو آشکارساز)، \texttt{dark\_rate} و \texttt{dark\_rate\_d1} (نرخ‌های تاریک نامتقارن)، \texttt{qber\_intrinsic} و \texttt{misalignment} که به‌ترتیب نرخ خطای ذاتی کوانتومی و ناهم‌راستایی پایه را کنترل می‌کنند، و \texttt{detector\_type} که بین سه مسیر \lr{SPD}، \lr{SNSPD} و \lr{PNRD} انتخاب می‌کند. گروه دوم، فیلدهای زمان‌بندی و نویز هستند: \texttt{dead\_time\_ns}، \texttt{jitter\_fwhm\_ns}، \texttt{afterpulse\_prob}، \texttt{afterpulse\_lifetime\_ns}، و حالت‌های نامتقارن لرزش \texttt{jitter\_fwhm\_ns\_d0/d1}. این فیلدها با \lr{metadata=\{"noise": True\}} علامت‌گذاری شده‌اند تا در عملیات \lr{noiseless\_copy()} به‌صورت خودکار صفر شوند و امکان شبیه‌سازی حالت ایده‌آل (بدون نویز) فراهم شود. گروه سوم، فیلدهای سخت‌افزاری/محیطی هستند: \texttt{temperature\_k}، \texttt{bias\_voltage}، \texttt{breakdown\_voltage}، \texttt{ref\_temperature\_k}، \texttt{ref\_bias\_voltage}، \texttt{bias\_current}، \texttt{switching\_current} که برای محاسبه مقیاس‌کننده‌های دمایی/ولتاژی نرخ تاریک به کار می‌روند.

حالت درون‌رونده در یک \lr{dataclass} جداگانه به نام \texttt{\_DetectorState} نگهداری می‌شود که به‌صورت یک فیلد \texttt{init=False} درون آشکارساز قرار گرفته و در فرآیند سریال‌سازی \texttt{to\_config\_dict()} از آن عبور نمی‌کند. این جداسازی صریح میان \lr{پیکربندی} و \lr{حالت} طراحی امنی برای جلوگیری از نشت حالت درون‌رونده به خروجی پیکربندی است و به کاربر اجازه می‌دهد یک آشکارساز را با همان پیکربندی اما با حالت بازنشانی‌شده بسازد. فیلدهای \texttt{\_DetectorState} شامل: \texttt{last\_abs\_fire\_time\_ns\_d0/d1} (با سنتینل $-\infty$ به‌معنای «هرگز کلیک نکرده»)، \texttt{pending\_ap\_time\_d0/d1} (با سنتینل $<0$ به‌معنای «بدون \lr{AP} معلق»)، \texttt{ap\_depth\_d0/d1} (عمق آبشار \lr{AP} معلق)، \texttt{next\_dc\_time\_d0/d1} (زمان برنامه‌ریزی‌شده شمارش تاریک بعدی)، \texttt{carry\_over\_events} (تاپل از \lr{CarryOverEvent}ها که رویدادهای ورودی پس از پایان بلاک را برای بلاک بعدی نگه می‌دارند)، \texttt{total\_time\_processed\_ns} (ساعت زمان مطلق پردازش‌شده)، \texttt{afterpulse\_total\_count} (شمارنده تجمعی \lr{AP}) و \texttt{last\_path} (\lr{"threshold"} یا \lr{"pnrd"} که برای کشف تغییر مسیر و هشدار درباره از دست رفتن حالت پیوسته به کار می‌رود) است.

\lr{NamedTuple} به نام \texttt{CarryOverEvent} برای ساختاردهی به رویدادهای \lr{carry-over} معرفی شده است. این ساختار با نام‌گذاری صریح فیلدها (\texttt{time\_abs\_ns}، \texttt{det\_id}، \texttt{src\_pulse\_idx}، \texttt{is\_signal\_origin}) از جابه‌جایی تصادفی مؤلفه‌ها هنگام تغییر شکل تاپل جلوگیری می‌کند و کد را خود-مستندکننده می‌سازد. خروجی شبیه‌سازی در دو \lr{NamedTuple} جداگانه بازگردانده می‌شود: \texttt{DetectionResult} برای خروجی استاندارد (با آرایه‌های کلیک، تشخیصی، و \lr{state\_snapshot} به‌صورت رشته \lr{JSON}) و \texttt{PhotonResolvedResult} برای تحلیل‌های \lr{decoy-state yield} که در آن فیلدهای اضافی \texttt{arrivals}، \texttt{target\_d0} و \texttt{ideal\_target\_d0} نیز در دسترس هستند.

اشیای تشخیصی در \lr{dataclass} ایستای \texttt{DetectionDiagnostics} جمع‌آوری می‌شوند که بیش از $25$ فیلد شمارشی دارد: از \texttt{double\_clicks\_total} و \texttt{resolved\_to\_d0/d1} (برای سیاست \lr{RANDOM}) گرفته تا \texttt{rogers\_discarded\_events}، \texttt{deadtime\_dropped}، \texttt{combined\_flips\_with\_arrivals}، \texttt{qber\_flips\_only}، \texttt{isi\_events} (تداخل بین-سمبلی)، \texttt{afterpulse\_events}، \texttt{dark\_count\_events}، \texttt{imd\_click\_events} و \texttt{multiple\_fires\_per\_slot}. این کلاس دارای \lr{property}های مفیدی مانند \texttt{tossed\_breakdown} (که تجزیه دسته‌بندی‌شده رویدادهای دورریخته را برمی‌گرداند) و \texttt{total\_discarded} است. شاخص \texttt{kernel\_tossed} که پیش از \lr{Review-v19 F-11} با تفریق محاسبه می‌شد و ممکن بود منفی شود، اکنون مستقیماً از diagnose هسته \lr{(DIAG\_TOSSED\_CORE)} پر می‌شود تا از پنهان‌شدن باگ‌های حسابداری تشخیصی جلوگیری شود.

یک \lr{dataclass} ایستای دیگر به نام \texttt{DetectorConstants} تمام ثابت‌های تنظیمی را متمرکز می‌کند تا تغییر آن‌ها بدون اعتبارسنجی مجدد در برابر داده‌های مشخصه‌سازی امکان‌پذیر نباشد. این کلاس شامل مقادیری مانند \texttt{MAX\_AFTERPULSE\_CASCADE = 32}، \texttt{SPAD\_DARK\_ACTIVATION\_EV = 0.15}، \texttt{SNSPD\_PLANCK\_MAX\_TERMS = 200}، \texttt{DEFAULT\_RNG\_BUFFER\_MULT = 10.0}، \texttt{LOG\_SAFE\_EPS = 1e-300} و چندین ثابت مرتبط با تحلیل \lr{Rogers} است. \lr{Enum}های پشتیبانی‌شده شامل \texttt{DetectorType} (\lr{SPD/SNSPD/PNRD})، \texttt{DeadTimeModel} (\lr{NON\_PARALYZABLE/PARALYZABLE} که دومی منسوخ‌شده است)، \texttt{AfterpulseModel} (\lr{EXPONENTIAL/GEOMETRIC})، \texttt{DetectorTopology} (\lr{INDEPENDENT\_SPADS/SHARED\_SPAD})، \texttt{DoubleClickPolicy}، \texttt{DecoderArchitecture} و \texttt{APRescheduleMode} (\lr{END/UNIFORM}) هستند.

کلاس کمکی \texttt{EntanglingDecoder} برای رمزگشای درهم‌تنیدگی در پروتکل‌های \lr{MDI-QKD} و بازگشت متقارن \lr{(redundant return)} طراحی شده است. این \lr{dataclass} ایستا با \texttt{slots=True} از عملگرهای کوانتومی \lr{CPTP} برای اعمال دیفیزینگ \lr{(dephasing)} آگاه از دیدپذیری \lr{(visibility-aware)} استفاده می‌کند: درایه‌های خارج‌قطری ماتریس چگالی مشترک $\rho_{joint}$ با ضریب $V$ (دیدپذیری) مقیاس می‌شوند در حالی که قطر (جمعیت‌ها) دست‌نخورده باقی می‌ماند. این کلاس از \lr{redundancy\_M} در \lr{\{2, 3\}} پشتیبانی می‌کند و برای $M=2$ از پروژکتورهای بل-حالت $|\Psi^-\rangle$ و $|\Psi^+\rangle$ استفاده می‌کند:

\begin{equation}
	|\Psi^\pm\rangle = \frac{1}{\sqrt{2}}\left(|01\rangle \pm |10\rangle\right)
\end{equation}

و برای $M=3$ از کلمات کد کیوبیتی-۳ برای کد فلیپ-بیت که حالت‌های $|000\rangle$، $|011\rangle$، $|101\rangle$ و $|110\rangle$ را به‌عنوان پایه‌های سندروم دو-بیتی در نظر می‌گیرد.

\section{پیاده‌سازی ریاضی و الگوریتمی}

مسیر آستانه‌ای \lr{(threshold path)} که برای \lr{SPD/SNSPD} فعال می‌شود، یک موتور رویداد-محور کامل را پیاده‌سازی می‌کند. در گام نخست، نازک‌سازی کانال \lr{(channel thinning)} با توزیع دو-جمله‌ای انجام می‌شود: برای هر پالس، تعداد فوتون‌های رسیده به آشکارساز با $\text{arrivals} \sim \text{Binomial}(n, \eta_{ch})$ نمونه‌برداری می‌شود که در آن $n$ تعداد فوتون منبع و $\eta_{ch} = \eta_{eff} \cdot F_{int}$ شامل شدت بازده کانال و فاکتور ناهم‌خوانی شدت است. سپس احتمال کلیک در هر پالس برای $m$ فوتون وارد شده با فرمول زیر محاسبه می‌شود:

\begin{equation}
	P_{click}(m, \eta) = 1 - (1-\eta)^m
\end{equation}

این محاسبه با $\texttt{log1p}$ و $\texttt{expm1}$ برای پایداری عددی پیاده‌سازی می‌شود تا برای مقادیر کوچک $\eta$ خطای گردکردن \lr{IEEE-754} به حداقل برسد. در ادامه، مسیریابی به \lr{D0/D1} با احتمال ترکیبی فلیپ $p_{combined}$ انجام می‌شود که پیش‌فرض آن از فرمول \lr{XOR} زیر به‌دست می‌آید:

\begin{equation}
	p_{combined} = m(1-p_{qber}) + (1-m)\,p_{qber}
\end{equation}

که در آن $m$ ناهم‌راستایی \lr{(misalignment)} و $p_{qber} = q_{intrinsic} + (1-V)/2$ است که $V$ دیدپذیری \lr{(visibility)} تداخل است. این فرمول \lr{XOR} فرض استقلال آماری میان ناهم‌راستایی و خطای \lr{QBER} را دارد که در \lr{BB84} واقعی از نظر فیزیکی نادرست است (هر دو از همان نقص‌های دیدپذیری ناشی می‌شوند)؛ به همین دلیل حالت سخت‌گیرانه استفاده از این فرمول را نیازمند ارائه صریح \texttt{flip\_prob\_override} یا فعال‌سازی \texttt{allow\_xor\_flip\_formula=True} می‌کند.

محاسبه نرخ تاریک پویا در متد \texttt{\_calculate\_dynamic\_rates()} انجام می‌شود که برای \lr{SPAD} از مقیاس‌کننده آرنیوس-درجه دوم و برای \lr{SNSPD} از مقیاس‌کننده طیف پلانک استفاده می‌کند. این متد برای جلوگیری از محاسبه مکرر انتگرال پلانک (که تا $200$ جمله سری می‌تواند داشته باشد) در هر فراخوانی شبیه‌سازی، کش می‌شود. کد زیر پیاده‌سازی مقیاس‌کننده \lr{SPAD} را نشان می‌دهد:

\begin{LTR}
	\begin{lstlisting}[language=Python]
		def _spad_dark_scaler(self) -> float:
		# Review-v19 F-01: temperature-dependent breakdown voltage
		coeff_V_per_K = self.breakdown_temp_coeff_mv_per_k * 1e-3
		delta_T = self.temperature_k - self.ref_temperature_k
		bdv_op = self.breakdown_voltage + coeff_V_per_K * delta_T
		excess_bias = self.bias_voltage - bdv_op
		ref_excess_bias = self.ref_bias_voltage - self.breakdown_voltage
		if ref_excess_bias <= EPS:
		ref_excess_bias = _REF_EXCESS_BIAS_FALLBACK_V
		voltage_factor = (excess_bias / ref_excess_bias) ** 2
		
		kT = CONST_BOLTZMANN * self.temperature_k
		ref_kT = CONST_BOLTZMANN * self.ref_temperature_k
		e_act = _SPAD_DARK_ACTIVATION_EV * _EV_TO_JOULE
		
		exponent = (e_act / ref_kT) - (e_act / kT)
		if exponent < -700.0 or exponent > 700.0:
		# guard against numerical overflow in math.exp
		if self.strict_mode:
		raise ParameterValidationError(...)
		exponent = max(-700.0, min(700.0, exponent))
		temp_factor = math.exp(exponent)
		return temp_factor * voltage_factor
	\end{lstlisting}
\end{LTR}

در خط نخست، ضریب $\text{mV/K}$ به $\text{V/K}$ تبدیل می‌شود تا در محاسبه جبری با $V$ سازگار باشد. سپس اختلاف دما $\Delta T$ و ولتاژ شکست عملیاتی $V_{bd}^{op}$ با اصلاح دمایی محاسبه می‌شوند. بایاس مازاد عملیاتی و مرجع به‌ترتیب $V_{eb}^{op} = V_{bias} - V_{bd}^{op}$ و $V_{eb}^{ref} = V_{bias}^{ref} - V_{bd}^{ref}$ هستند. اگر بایاس مازاد مرجع نزدیک به صفر باشد، از مقدار جایگزین \texttt{\_REF\_EXCESS\_BIAS\_FALLBACK\_V} استفاده می‌شود تا از تقسیم بر صفر جلوگیری شود. فاکتور ولتاژ به‌صورت درجه دوم و فاکتور دما به‌صورت نمایی \lr{Arrhenius} محاسبه می‌شوند. محافظت عددی برای بازه $[-700, 700]$ اعمال می‌شود تا از سرریز $\exp$ جلوگیری شود؛ این محدوده برای اکثر پیکربندی‌های واقع‌گرایانه (دماهای $200$--$400$ \lr{K}) کافی است.

هسته پردازش رویداد که در ماژول جداگانه \texttt{detector\_kernel.py} پیاده‌سازی شده (که فصل جداگانه‌ای به آن اختصاص دارد)، در این ماژول توسط متد \texttt{\_simulate\_event\_driven()} فراخوانی می‌شود. این متد ابتدا آرایه‌های رویداد ساختار-یافته را می‌سازد، سپس \lr{const\_params} (یک آرایه \lr{float64} با $15$ مؤلفه که شامل زمان مرده، احتمال \lr{AP}، طول عمر \lr{AP}، دوره پالس، زمان شروع بلاک، مدت بلاک، نرخ‌های تاریک \lr{D0/D1} و... است) و بافر \lr{RNG} پیش‌کشیده را آماده می‌کند، و سپس هسته را فرامی‌خواند. خروجی هسته شامل زمان‌های کلیک، شناسه آشکارساز، پرچم منشأ سیگنال، آرایه تشخیصی، و تاپل حالت-خروج است که برای به‌روزرسانی \texttt{\_DetectorState} استفاده می‌شود.

مسیر \lr{PNRD} در متد \lr{\texttt{\_simulate\_pnrd()}} به‌صورت کاملاً متفاوتی پیاده‌سازی شده است. در این مسیر، اثر زمان مرده، پس‌پالس و فرآیند زمانی پیوستهٔ شمارش‌های تاریک لحاظ نمی‌شود؛ در عوض، برای هر پالس از یک مدل پواسونیِ مستقل جهت شبیه‌سازی شمارش‌های تاریک استفاده می‌شود. این انتخاب از نظر سازگاری درونی مسیر قابل‌قبول است، زیرا آشکارسازهای واقعی \lr{PNRD} معمولاً در رژیمی کار می‌کنند که در آن زمان مرده ناچیز است. کد زیر، نازک‌سازی کانال و مسیریابی را در مسیر \lr{PNRD} نشان می‌دهد:


\begin{LTR}
	\begin{lstlisting}[language=Python]
		# Channel thinning (binomial)
		photon_numbers_i = np.asarray(photon_numbers, dtype=np.int64)
		arrivals = rng.binomial(photon_numbers_i, eff_trans)
		
		if per_photon_routing_pnrd and has_multiphoton:
		# Per-photon routing (the default, physically correct)
		n_total = int(np.sum(arrivals))
		if n_total > 0:
		photon_flips = rng.random(n_total) < p_combined_flip
		ideal_per_pulse = np.asarray(ideal_outcomes_d0, dtype=np.bool_)
		photon_ideal = np.repeat(ideal_per_pulse, arrivals)
		photon_actual_to_d0 = photon_ideal ^ photon_flips
		pulse_idx = np.repeat(np.arange(num_pulses), arrivals)
		arrivals_d0 = np.zeros(num_pulses, dtype=np.int64)
		np.add.at(arrivals_d0, pulse_idx,
		photon_actual_to_d0.astype(np.int64))
		arrivals_d1 = arrivals - arrivals_d0
		# ...
		det_counts_d0 = rng.binomial(arrivals_d0, self.det_eff_d0)
		det_counts_d1 = rng.binomial(arrivals_d1, self.det_eff_d1)
		mu_dc_d0 = max(0.0, dr0) * float(pulse_period_ns) * 1e-9
		dark_counts_d0 = rng.poisson(mu_dc_d0, size=num_pulses)
		final_d0 = det_counts_d0 + dark_counts_d0 + imd_counts_d0
	\end{lstlisting}
\end{LTR}

این کد نشان می‌دهد که مسیر \lr{PNRD} از \texttt{rng.binomial} برای نازک‌سازی کانال و بازده آشکارساز استفاده می‌کند و \texttt{rng.poisson} را برای شمارش تاریک هر-پالس به کار می‌گیرد. حالت پیش‌فرض \texttt{per\_photon\_routing=True} مسیریابی هر-فوتون را (که فیزیکی صحیح است) انتخاب می‌کند؛ اگر ورودی چند-فوتونی در حالت سخت‌گیرانه با \texttt{per\_photon\_routing=False} ارائه شود، استثنا صادر می‌شود زیرا مسیریابی هر-پالس بازده \lr{decoy-state} $Y_n$ برای $n \geq 2$ را به‌صورت خاموش تحریف می‌کند.

تحلیل‌های تحلیلی \lr{Rogers et al. (2007)} به‌صورت چند تابع سطح-ماژول پیاده‌سازی شده‌اند. تابع \lr{\texttt{rogers\_P\_00(p, k, eps)}}
 حالت پایدار احتمال اینکه هر دو آشکارساز در یک پایه زنده باشند را با تکرار ارزش \lr{(value iteration)} روی فضای حالت دوبعدی $\{(i,j): 0 \leq i, j \leq k\}$ محاسبه می‌کند. این تابع برای $k$ بزرگ با کران تکرار تطبیقی $\min(50000, \max(5000, k^2 \cdot 100))$ و تلورانس $10^{-12}$ همگرا می‌شود. تابع\lr{\texttt{rogers\_T\_N(N, p, k)}} توزیع احتمال طول دنباله را با فرمول هندسی بازتولید می‌کند:

\begin{equation}
	T_N = (1 - s)^{N-1} \cdot s, \quad s = (1 - 2p)^k
\end{equation}

که در آن $p = L/8$ (پارامتر از-دست‌رفته پیوند) و $k = \rho_{TX} \cdot \tau$ (تعداد دوره‌های ارسال در هر زمان مرده) است. تابع \lr{\texttt{rogers\_S(p, k)}}
احتمال سيفتنگ موفق یک بیت از یک دنباله را به‌صورت زیر محاسبه می‌کند:

\begin{equation}
	S(p, k) = \sum_{N=1}^{N_{max}} T_N(p, k) \cdot \left(1 - \frac{1}{2^N}\right)
\end{equation}

و نرخ بیت سيفتی \lr{(SBR)} از معادله \lr{(15)} مقاله به دست می‌آید:

\begin{equation}
	\text{SBR} = \rho_{TX} \cdot 8p \cdot P_{00}(p, k) \cdot S(p, k)
\end{equation}

همچنین تقریب‌های مجانبی مقاله در توابع \texttt{rogers\_sbr\_max(\textbackslash tau) = 1.433/(2\textbackslash tau)} و \texttt{rogers\_rho\_tx\_max(L, \textbackslash tau) = 5.92/(L\textbackslash tau)} پیاده‌سازی شده‌اند که به ترتیب نرخ بیت سيفتی بیشینه و نرخ ارسال بیشینه را در شرط بهینه‌سازی ارائه می‌کنند.

پیاده‌سازی سیاست سيفتنگ \lr{ROGERS\_2007} در متد ایستا \texttt{\_apply\_rogers\_2007\_sifting()} قرار دارد. این متد یک اجرای ماکزیمال از رویدادهای کلیک را با بررسی فاصله زمانی متوالی شناسایی می‌کند: یک رویداد به‌عنوان شروع دنباله در نظر گرفته می‌شود اگر اولین رویداد باشد یا شکاف زمانی $\geq \tau_d$ از رویداد قبلی داشته باشد. سپس تنها رویدادهای شروع دنباله (که در پنجره گیت نیز قرار دارند) در آرایه‌های کلیک نهایی نگه داشته می‌شوند و سایر رویدادها به‌عنوان \texttt{rogers\_discarded\_events} شمارش می‌گردند. این متد از زمان‌های واقعی \lr{(true times)} برای مقایسه شکاف با زمان مرده استفاده می‌کند (نه زمان‌های لرزش‌دیده) که از نظر فیزیکی صحیح است: ساعت زمان مرده آشکارساز در زمان واقعی رسیدن فوتون شروع به کار می‌کند، نه در زمان لرزش‌دیده‌ی اندازه‌گیری.

\section{ارسال و دریافت با سایر ماژول‌ها}

این ماژول به‌عنوان گره میانی میان کانال فیبر نوری و لایه تحلیل کلید امن عمل می‌کند و رابط‌های ورودی و خروجی متعددی دارد. از منظر ورودی، متد عمومی \texttt{simulate\_detection()} یک سیگنال چندوجهی دریافت می‌کند که از ماژول‌های بالا‌دست می‌آید: پارامتر \texttt{channel\_transmittance} (یک عدد \lr{float} در بازه $[0, 1]$) از ماژول کانال فیبر نوری \lr{qkd.channel} می‌آید و حاصل قانون بیر-لامبرت $T = 10^{-L_{total}/10}$ است؛ آرایه \texttt{photon\_numbers} (از نوع \lr{NDArray[np.int64]}) از ماژول منبع نوری \lr{qkd.source} می‌آید و تعداد فوتون‌های ارسالی آلیس در هر پالس را با توزیع پواسون یا توزیع چندجمله‌ای \lr{decoy} ارائه می‌کند؛ آرایه \texttt{ideal\_outcomes\_d0} (از نوع \lr{NDArray[np.bool\_]}) از ماژول رمزگذاری پایه \lr{qkd.encoder} می‌آید و نشان می‌دهد کدام پالس‌ها در حالت ایده‌آل به \lr{D0} مسیریابی شده‌اند؛ آرایه \texttt{basis\_selector} اختیاری نیز از همان ماژول می‌آید و برای فعال‌سازی بازده پایه-وابسته \lr{(basis-dependent efficiency)} با فیلدهای \texttt{det\_eff\_d0\_z/x} و \texttt{det\_eff\_d1\_z/x} به کار می‌رود.

شیء \texttt{rng} از نوع \lr{numpy.random.Generator} یا از ماژول بالادستی \lr{qkd.simulation\_driver} می‌آید یا توسط کاربر به‌صورت مستقیم ساخته می‌شود و بذر آن برای بازتولیدپذیری شبیه‌سازی حیاتی است. ماژول از مکانیزم \lr{RNG state snapshot/restore} برای اطمینان از اینکه تلاش‌های مجدد بافر \lr{(buffer retries)} نتایج یکسانی با همان بذر تولید کنند، استفاده می‌کند. کلمات کلیدی \texttt{kwargs} شامل پارامترهای فیزیکی \lr{SCM} مانند \texttt{visibility\_mismatch\_dm}، \texttt{visibility\_bias\_drift\_psi1/psi2}، \texttt{modulation\_index}، \texttt{mu\_signal} و \texttt{N\_channels} است که از ماژول تحلیل \lr{SCM} \lr{(qkd.utils.scm\_analysis)} می‌آیند و برای محاسبه شدت \lr{IMD/FWM} به کار می‌روند. پارامترهای \lr{WDM} مانند \texttt{include\_wdm\_fwm\_noise} نیز از ماژول تحلیل \lr{WDM} عرضه می‌شوند.

از منظر خروجی، ماژول یک \texttt{DetectionResult} (یا \texttt{PhotonResolvedResult}) برمی‌گرداند که شامل سه بخش اصلی است. آرایه‌های \texttt{click0} و \texttt{click1} (از نوع \lr{NDArray[np.bool\_]} برای آستانه‌ای و \lr{NDArray[np.int64]} برای \lr{PNRD}) به طول \texttt{num\_pulses} هستند و به ماژول سيفتنگ پایه \lr{(qkd.sifting)} ارسال می‌شوند تا بیت‌های سيفتی استخراج شوند. شیء \texttt{DetectionDiagnostics} به ماژول گزارش‌گیری و تجزیه‌وتحلیل آماری \lr{(qkd.analysis)} ارسال می‌شود و در آن معیارهایی مانند \lr{QBER} مشاهده‌شده، بازده \lr{decoy-state} $Y_n$، نرخ \lr{IMD} و نرخ \lr{ISI} محاسبه می‌شوند. رشته \texttt{state\_snapshot} (یک \lr{JSON} سریال‌شده با \texttt{allow\_nan=False}) به ماژول مدیریت حالت \lr{(qkd.state\_manager)} ارسال می‌شود تا در صورت نیاز به ادامه شبیه‌سازی در بلاک‌های بعدی، حالت آشکارساز به‌صورت دقیق قابل بازیابی باشد.

از منظر وابستگی‌های درون‌پروژه‌ای، این ماژول از چندین ماژول زیرین واردات دارد: \texttt{qkd.datatypes} برای \lr{Enum}ها و \lr{dataclass}های مشترک مانند \lr{Double\-Click\-Policy}، \lr{Detector\-Type}، \lr{Detection\-Config}، \lr{Decoder\-Architecture} و \lr{Tally\-Counts}؛ \texttt{qkd.exceptions} برای \lr{Configuration\-Error} و \lr{Parameter\-Validation\-Error} که با زمینه‌های ساختاریافته \lr{(param\_name, param\_value, context)} هستند؛ \texttt{qkd.constants} برای توابع کمری مانند \lr{is\_valid\_probability}، \lr{is\_finite\_non\_negative}، \lr{clamp\_probability} و ثابت‌های فیزیکی \lr{CONST\_BOLTZMANN}، \lr{CONST\_ELECTRON\_CHARGE}، \lr{CONST\_PLANCK} و \lr{CONST\_SPEED\_OF\_LIGHT}؛ \texttt{qkd.utils.utils} برای \lr{sanitize\_for\_serialization}؛ و \texttt{qkd.utils.scm\_analysis} برای توابع \lr{calculate\_visibility}، \lr{calculate\_intensity\_mismatch\_factor} و \lr{calculate\_wdm\_fwm\_noise}.

این ماژول همچنین یک لایه سازگار-کاهش \lr{(graceful fallback)} برای زمانی که بسته \lr{qkd} کامل نصب نباشد (مانند زمان اجرای تست‌های مستقل) دارد که در آن نسخه‌های ساده‌شده‌ی تمام \lr{Enum}ها، ثابت‌ها و توابع کمری به‌صورت درون‌ماژولی تعریف می‌شوند. این لایه با پرچم \lr{\_HAS\_QKD\_PACKAGE} کنترل می‌شود و امکان تست‌گذاری جداگانه ماژول را بدون نیاز به کل زیرساخت پروژه فراهم می‌سازد. با این‌حال، در کاربردهای عملی، همیشه نسخه کامل بسته \lr{qkd} باید نصب باشد تا تمام اعتبارسنجی‌ها و رفتارهای سخت‌گیرانه فعال شوند.

از منظر مصرف‌کنندگان اصلی، ماژول \lr{qkd.main\_optimized} (موتور اصلی شبیه‌سازی) از \texttt{build\_detector\_sim\_params()} برای ساخت یک بسته ایستای \lr{DetectorSimParams} از آشکارساز اعتبارسنجی‌شده استفاده می‌کند؛ این بسته شامل یازده پارامتر مشتق‌شده از آشکارساز است و برای جلوگیری از خواندن مکرر از دیکشنری پیکربندی در حلقه اصلی شبیه‌سازی به کار می‌رود. ماژول‌های تحلیل کلید امن مانند \lr{qkd.key\_rate} نیز از خروجی‌های \texttt{DetectionDiagnostics} به‌ویژه \texttt{combined\_flips\_with\_arrivals}، \texttt{dark\_count\_events} و \texttt{afterpulse\_events} برای تغذیه فرمول‌های نرخ کلید \lr{GLLP} و \lr{decoy-state} استفاده می‌کنند. در نهایت، تابع \texttt{is\_threshold\_detector()} به‌عنوان یک تابع کمکی برای سایر ماژول‌هایی که نیاز به بررسی نوع آشکارساز دارند (مانند انتخاب مسیر محاسبه بازده) فراهم شده است.


\lstset{
	basicstyle=\ttfamily\footnotesize,
	breakatwhitespace=false,
	breaklines=true,
	captionpos=b,
	keepspaces=true,
	numbers=left,
	numbersep=5pt,
	numberstyle=\tiny\color{gray},
	showspaces=false,
	showstringspaces=false,
	showtabs=false,
	tabsize=2,
	language=Python,
	frame=single,
	framesep=2pt,
	rulecolor=\color{gray!50},
	backgroundcolor=\color{gray!5},
	keywordstyle=\color{blue!70!black}\bfseries,
	commentstyle=\color{green!50!black}\itshape,
	stringstyle=\color{red!60!black},
	emph={np, math, njit, NUMBA_AVAILABLE, _process_events_logic, _process_events_numba, event_np_dtype, _SENTINEL_T, STATUS_OK, STATUS_RNG_EXHAUSTED, STATUS_OUTPUT_BUFFER_EXHAUSTED, STATUS_CARRY_OVER_BUFFER_EXHAUSTED, STATUS_ITERATION_CAP_EXCEEDED, DIAG_DEADTIME, DIAG_AFTERPULSE, DIAG_DARKCOUNT, DIAG_TOSSED_CORE, DIAG_IMD_CLICKS, DIAG_CARRY_OVER, DIAG_DEADTIME_DROPPED, DIAG_MULTIPLE_FIRES_PER_SLOT, DIAG_COUNT, CP_DEAD_TIME, CP_AP_PROB, CP_AP_LIFETIME, CP_PULSE_PERIOD, CP_TIME_START, CP_BATCH_DURATION, CP_DR0, CP_DR1, CP_AP_GEOMETRIC, CP_MAX_CASCADE, CP_GATED, CP_GATE_WIDTH, CP_GATE_OFFSET, CP_SHARED_SPAD, CP_AP_RESCHEDULE_UNIFORM, CP_LEN},
	emphstyle=\color{purple!70!black}\bfseries,
}

\section{نمای کلی و هدف ماژول}

ماژول \lr{\texttt{qkd.detector\_kernel}} هسته پردازش رویداد \lr{(event-processing kernel)} و پوشش کامپایل \lr{Numba JIT} را برای موتور شبیه‌سازی آشکارساز فریم‌ورک \lr{QKD} در خود جای داده است. این ماژول نقش حلقه داغ \lr{(hot loop)} پردازش متوالی رویدادها را ایفا می‌کند و عملکرد کل فریم‌ورک شبیه‌سازی به‌طور مستقیم به بهینه‌سازی این ماژول وابسته است. هدف اصلی، پیاده‌سازی یک هسته پردازشی است که رویدادهای کلیک آشکارساز را با وفاداری کامل به فیزیک پدیده‌های موقت \lr{(time-domain phenomena)} - شامل زمان مرده غیر-پارالایز، تولید پالس پسین با کران آبشاری، زمان‌بندی پیوسته شمارش تاریک و \lr{carry-over} رویدادهای ورودی پس از پایان بلاک - پردازش کند و در عین حال با مکانیزم \lr{Numba JIT} به سرعت اجرای نزدیک به \lr{C} دست یابد. این ماژول از ماژول مادر \texttt{qkd.detectors} استخراج شده تا جداسازی صریح میان منطق پردازش رویداد (پُر از بحت‌های عددی و فیزیکی) و منطق سطح بالای آشکارساز (ساخت پیکربندی، اعتبارسنجی، سریال‌سازی) را فراهم سازد.

پیچیدگی اصلی این ماژول از چند منبع نشئت می‌گیرد. نخست، هسته باید با ترتیب زمانی سختگیرانه رویدادها را پردازش کند و در زمان‌های تساوی، اولویت صریح \lr{input > afterpulse > dark count} را اعمال نماید. این اولویت‌بندی برای تضمین بازتولیدپذیری \lr{(reproducibility)} نتایج در پلتفرم‌های مختلف و در حضور نویز تصادفی حیاتی است. دوم، هسته باید حالت آشکارساز را در طول چندین بلاک شبیه‌سازی به‌صورت پیوسته نگه دارد؛ به این معنا که رویدادهای \lr{afterpulse} و شمارش تاریک که در پایان یک بلاک برنامه‌ریزی شده‌اند، در بلاک بعدی به‌صورت دقیق در زمان‌های مطلق صحیح خود پردازش شوند. این نیازمند یک \lr{state-out tuple} غنی است که تمام حالت موقت را در خود نگه دارد. سوم، هسته باید با چندین سناریوی فرسایشی \lr{(exhaustion scenarios)} مقابله کند: فرسایش بافر \lr{RNG}، فرسایش بافر خروجی، فرسایش بافر \lr{carry-over}، و فراتر رفتن از کران تکرار. در هر یک از این موارد، هسته باید به‌صورت تشخیصی \lr{(diagnostic)} رفتار کند و خروجی معتبری تولید نماید که به کاربر امکان دیباگ و تلاش مجدد با بافر بزرگ‌تر را بدهد.

این ماژول به‌صورت آگاهانه بر روی مکانیزم \lr{Numba JIT} تکیه دارد اما در صورت عدم دسترسی یا شکست کامپایل، به‌صورت شفاف به پیاده‌سازی \lr{pure-Python} بازمی‌گردد. این طراحی \lr{graceful fallback} با پرچم \texttt{NUMBA\_AVAILABLE} کنترل می‌شود و در صورت وقوع، تنها یک هشدار لاگ صادر می‌کند. نکته مهم این است که کامپایل \lr{Numba} با \texttt{fastmath=False} انجام می‌شود تا از بازترتیب عملیات \lr{IEEE-754-violating} توسط کامپایلر جلوگیری شود و نتایج \lr{bit-identical} میان مسیر \lr{Numba-compiled} و \lr{pure-Python} (در محدوده عملیاتی هسته) تضمین گردد. با این‌حال، مستندات ماژول به‌صراحت اشاره می‌کنند که این تضمین به‌صورت رسمی \lr{bit-identical} نیست و نقاط واگرایی شناخته‌شده‌ای (مانند \texttt{math.log} روی ورودی‌های زیر-نرمال و عرض \lr{int()} در پایتون در برابر \lr{Numba}) وجود دارند که در محدوده عملیاتی هسته (زمان‌مهرهای کمتر از $10^{16}$ و \texttt{max\_steps} کمتر از $10^{7}$) غیرقابل‌دسترس هستند.

ماژول از ماژول خواهری \texttt{qkd.detector\_datastructures} دو نهاد حیاتی را وارد می‌کند: ثابت \texttt{\_SENTINEL\_T} (یک مقدار شناور بسیار بزرگ که به‌عنوان سنتینل «هیچ رویداد» استفاده می‌شود) و \texttt{event\_np\_dtype} (یک \lr{NumPy structured dtype} که شکل هر رویداد را در آرایه ورودی تعریف می‌کند). این وابستگی به‌صورت سختگیرانه اعمال می‌شود و در صورت عدم دسترسی، به‌جای \lr{fallback} خاموش، استثنای \lr{ImportError} صادر می‌کند؛ زیرا یک \lr{fallback} با \lr{dtype} متفاوت می‌تواند به واگرایی پنهان میان مسیرهای \lr{Numba-compiled} و \lr{pure-Python} منجر شود که دیباگ آن بسیار دشوار است. این تصمیم طراحی بر اساس ممیزی \lr{F9} گرفته شده و سیاست صریح ماژول را به \lr{«no silent fallback for dtype divergence»} تبدیل می‌کند.

\section{مبانی تئوری و فیزیکی}

هسته پردازش رویداد بر پایه چند مفهوم بنیادین از فیزیک آشکارسازهای تک‌فوتونی و تئوری صف‌های رویداد استوار است. نخست، مدل زمان مرده غیر-پارالایز \lr{(non-paralyzable dead-time model)} که طبق \lr{Rogers et al. (2007)} رفتار صحیح یک \lr{SPAD} منفرد را توصیف می‌کند. در این مدل، وقتی یک کلیک در زمان $t_f$ رخ می‌دهد، آشکارساز تا زمان $t_f + \tau_d$ به هیچ ورودی حساس نیست و فوتون‌های رسیده در طول این پنجره هیچ تأثیری بر آشکارساز ندارند - نه کلیک تولید می‌کنند و نه پنجره مرده را تمدید می‌نمایند. این رفتار با بررسی ساده زیر در هسته پیاده‌سازی می‌شود:

\begin{equation}
	\text{if } (t_{current} - t_{last\_fire}) < \tau_d \Rightarrow \text{suppress}
\end{equation}

نکته فیزیکی مهم این است که این بررسی با \textit{زمان واقعی} \lr{(true time)} رویداد انجام می‌شود، نه با زمان لرزش‌دیده \lr{(jittered time)}. این تمایز حیاتی است زیرا ساعت زمان مرده آشکارساز در زمان واقعی رسیدن فوتون شروع به کار می‌کند، در حالی که زمان ثبت‌شده در خروجی (که برای بن کردن به پالس‌ها استفاده می‌شود) تحت تأثیر لرزش الکترونیک است.

دومین مفهوم، مدل تولید پالس پسین \lr{(afterpulse generation)} با توزیع تأخیر است. وقتی یک کلیک رخ می‌دهد، با احتمال $p_{ap}$ یک پالس پسین تولید می‌شود که زمان آغاز آن از توزیع تأخیر نمایی یا هندسی نمونه‌برداری می‌شود. در توزیع نمایی پیوسته، تأخیر با $\text{decay} = -\ln(u) \cdot \tau_{ap}$ نمونه‌برداری می‌شود که در آن $u \sim \text{Uniform}(0, 1)$ و $\tau_{ap}$ طول عمر پالس پسین است:

\begin{equation}
	P(\text{delay} > t) = \exp\left(-\frac{t}{\tau_{ap}}\right)
\end{equation}

در توزیع هندسی گسسته روی شبکه دوره پالس $T$، پارامتر هندسی $p = 1 - \exp(-T/\tau_{ap})$ به‌گونه‌ای انتخاب می‌شود که میانگین تأخیر در هر دو توزیع برابر $\tau_{ap}$ باشد. تعداد گام‌های تأخیر با $K = \lceil \ln(u) / \ln(1-p) \rceil$ نمونه‌برداری می‌شود و $\text{decay} = K \cdot T$ تنظیم می‌گردد. یک کران سخت \texttt{max\_steps} (که تطبیقی با طول عمر و مدت بلاک است) از نمونه‌برداری تأخیرهای بسیار بزرگ جلوگیری می‌کند تا از سرریز بافر \lr{RNG} جلوگیری شود.

سومین مفهوم، عمق آبشار پالس پسین \lr{(afterpulse cascade depth)} است. وقتی خود یک پالس پسین کلیک می‌کند، می‌تواند پالس پسین دیگری تولید کند که این چرخه می‌تواند به‌صورت بی‌نهایت ادامه یابد. در عمل، احتمال تولید در هر مرحله به‌صورت نمایی کاهش می‌یابد اما برای جلوگیری از زنجیره‌های طولانی، یک کران سخت \texttt{MAX\_AFTERPULSE\_CASCADE = 32} اعمال می‌شود. وقتی این کران رسید، عمق آبشار به مقدار کران قفل می‌شود تا شروع فوری یک آبشار جدید در کلیک بعدی همان آشکارساز جلوگیری شود. یک آبشار جدید تنها می‌تواند از یک کلیک غیر-\lr{AP} (سیگنال، شمارش تاریک یا \lr{IMD}) آغاز شود که عمق را به $1$ بازنشانی می‌کند.

چهارمین مفهوم، زمان‌بندی پیوسته شمارش تاریک \lr{(continuous-time dark-count scheduling)} است. شمارش تاریک به‌صورت یک فرآیند پواسون همگن با نرخ $\lambda$ مدل می‌شود که در آن زمان میان-رسیدها از توزیع نمایی با میانگین $1/\lambda$ پیروی می‌کند:

\begin{equation}
	P(\text{inter-arrival} > t) = \exp(-\lambda t)
\end{equation}

نمونه‌برداری با تبدیل معکوس \lr{(inverse transform sampling)} انجام می‌شود: $\text{wait} = -\ln(u) / \lambda$ که در آن $u \sim \text{Uniform}(0, 1)$. برای جلوگیری از $\ln(0) = -\infty$، یک کران پایین \texttt{log\_safe\_eps = 1e-300} اعمال می‌شود که به حداکثر تأخیر $\approx 690/\lambda$ منجر می‌شود (که برای هر مدت بلاک واقع‌گرایانه عملاً نامحدود است). نکته مهم این است که وقتی نرخ شاری تاریک میان بلاک‌ها به صفر تغییر می‌کند، یک زمان شاری تاریک معلق از بلاک قبلی باید به‌صورت صریح باطل شود؛ در غیر این صورت یک شاری تاریک در بلاکی که نرخ صفر است اجرا می‌شود که از نظر فیزیکی محال است (اصلاح \lr{F-H9}).

پنجمین مفهوم، \lr{carry-over} رویدادهای ورودی است. وقتی بلاک شبیه‌سازی پایان می‌یابد، ممکن است رویدادهای ورودی وجود داشته باشند که زمان مطلق آن‌ها فراتر از \lr{time\_end\_abs} است. این رویدادها نباید در بلاک فعلی پردازش شوند؛ به‌جای آن، در بافر \lr{carry-over} ذخیره می‌شوند تا در بلاک بعدی به‌عنوان رویدادهای ورودی اولیه فراخوانی شوند. نکته ظریف این است که رویدادهای ورودی که در طول پنجره زمان مرده رخ داده‌اند (و در نتیجه در بلاک فعلی سرکوب شده‌اند) \textbf{نباید} به بافر \lr{carry-over} منتقل شوند؛ زیرا پنجره مرده در بلاک بعدی نیز بر اساس \texttt{last\_fire} قبلی اعمال می‌شود و انتقال چنین رویدادی منجر به سرکوب مضاعف یا حتی رفتار نامشخص می‌شود. این نکته در اصلاح \lr{C4} به‌صورت صریح پیاده‌سازی شده است.

ششمین مفهوم، برنامه‌ریزی مجدد پالس پسین در زمان مرده \lr{(AP reschedule in dead time)} است. وقتی یک پالس پسین در طول پنجره زمان مرده فعال می‌شود، نمی‌تواند کلیک تولید کند و باید به زمانی پس از پایان پنجره مرده برنامه‌ریزی مجدد شود. این ماژول دو سیاست برنامه‌ریزی مجدد را پشتیبانی می‌کند: سیاست \lr{END} که پالس پسین را در انتهای پنجره مرده ($t_{last\_fire} + \tau_d$) قرار می‌دهد، و سیاست \lr{UNIFORM} که آن را به‌صورت یکنواخت در بازه $[t_{last\_fire} + \tau_d, t_{last\_fire} + \tau_d + \tau_{ap}]$ نمونه‌برداری می‌کند. سیاست \lr{UNIFORM} فیزیکاً صحیح‌تر است زیرا زمان آزادسازی تله‌ها به‌صورت یکنواخت در طول طول عمر تله توزیع شده است. اصلاح \lr{Review-v12 F2} یک باگ حیاتی را برطرف کرد که در آن بازه نمونه‌برداری به‌اشتباه $[t_{last\_fire}, t_{last\_fire} + \tau_d]$ (درون پنجره مرده) بود که منجر به حلقه‌های بی‌نهایت برنامه‌ریزی مجدد می‌شد.

\section{تحلیل معماری و ساختار داده‌ها}

معماری این ماژول حول دو نهاد اصلی سازمان‌دهی شده: تابع \texttt{\_process\_events\_logic()} که پیاده‌سازی \lr{pure-Python} هسته است و متغیر \texttt{\_process\_events\_numba} که یا به همین تابع (در صورت عدم دسترسی \lr{Numba}) یا به نسخه \lr{JIT-compiled} آن اشاره می‌کند. این طراحی به مصرف‌کنندگان اجازه می‌دهد همیشه از \texttt{\_process\_events\_numba} استفاده کنند و نگران دسترس‌پذیری \lr{Numba} نباشند. در زمان \lr{import} ماژول، در صورت دسترس بودن \lr{Numba}، تلاش کامپایل با \texttt{njit(cache=True, fastmath=False)} انجام می‌شود و یک کامپایل اجباری \lr{(one-shot)} با ورودی‌های کوچک برای اشکال‌زدایی زودهنگام خطاهای کامپایل و پر کردن کش دیسک اجرا می‌گردد. اگر کامپایل شکست بخورد، یک هشدار لاگ صادر شده و \texttt{\_process\_events\_numba} به \texttt{\_process\_events\_logic} بازمی‌گردد.

امضای تابع \texttt{\_process\_events\_logic()} به‌صورت زیر است:

\begin{LTR}
	\begin{lstlisting}[language=Python]
		def _process_events_logic(
		sorted_events: NDArray[Any],          # structured array
		last_abs_fire_d0: float,              # state-in: last fire times
		last_abs_fire_d1: float,
		pending_ap_d0: float,                 # state-in: pending AP times
		pending_ap_d1: float,
		next_dc_d0_in: float,                 # state-in: next DC times
		next_dc_d1_in: float,
		const_params: NDArray[np.float64],    # 15-float bundle
		rng_floats: NDArray[np.float64],      # pre-drawn RNG buffer
		co_buffer: NDArray[Any],              # carry-over output buffer
		out_times: NDArray[np.float64],       # fire-time output buffer
		out_dets: NDArray[np.int64],          # detector-id output buffer
		out_sigs: NDArray[np.bool_],          # signal-origin flag buffer
		cascade_depths_in: NDArray[np.int64], # 2-element AP depth state
		ap_overflow_times: NDArray[np.float64],   # 4-element multi-trap overflow
		ap_overflow_depths: NDArray[np.int64],    # 4-element overflow depths
		):
	\end{lstlisting}
\end{LTR}

این امضای طولانی به‌دلیل نیاز هسته به دسترسی مستقیم به تمام حالت و بافرهای ورودی/خروجی بدون عبور از لایه‌های شیء-گرا است؛ این طراحی برای کارایی \lr{Numba} حیاتی است زیرا \lr{Numba} نمی‌تواند به‌صورت کارآمد با اشیاء پایتون تعامل داشته باشد و آرایه‌های \lr{NumPy} بهترین نوع داده برای کامپایل \lr{JIT} هستند. پارامتر \texttt{sorted\_events} یک آرایه ساختار-یافته با \lr{dtype} \texttt{event\_np\_dtype} است که شامل چهار فیلد \texttt{time}، \texttt{det}، \texttt{idx} و \texttt{sig} برای هر رویداد است؛ این فیلدها به‌ترتیب زمان مطلق، شناسه آشکارساز، شاخص پالس منبع و پرچم منشأ سیگنال را نگه می‌دارند.

پارامتر \texttt{const\_params} یک آرایه \lr{float64} با طول ثابت \texttt{CP\_LEN = 15} است که تمام پارامترهای فیزیکی و پیکربندی را در خود جای داده. این پارامترها با شاخص‌های نام‌دار \texttt{CP\_DEAD\_TIME}، \texttt{CP\_AP\_PROB}، \texttt{CP\_AP\_LIFETIME}، \texttt{CP\_PULSE\_PERIOD}، \texttt{CP\_TIME\_START}، \texttt{CP\_BATCH\_DURATION}، \texttt{CP\_DR0}، \texttt{CP\_DR1}، \texttt{CP\_AP\_GEOMETRIC}، \texttt{CP\_MAX\_CASCADE}، \texttt{CP\_GATED}، \texttt{CP\_GATE\_WIDTH}، \texttt{CP\_GATE\_OFFSET}، \texttt{CP\_SHARED\_SPAD} و \texttt{CP\_AP\_RESCHEDULE\_UNIFORM} دسترسی می‌شوند. این رویکرد با شاخص‌های نام‌دار (به‌جای دیکشنری یا کلاس) برای کارایی \lr{Numba} ضروری است زیرا دسترسی به آرایه با شاخص عددی سریع‌ترین روش دسترسی در کد کامپایل‌شده است. مقادیر بولی به‌صورت \lr{float64} ذخیره می‌شوند و با مقایسه $\geq 0.5$ به بولی تبدیل می‌شوند (نه $> 0.5$ که در نقطه میانی $0.5$ رفتار شکننده‌ای دارد).

پارامترهای \texttt{ap\_overflow\_times} و \texttt{ap\_overflow\_depths} آرایه‌هایی با طول $4$ هستند که \lr{multi-trap overflow} را برای هر آشکارساز پیگیری می‌کنند. دو شکاف نخست ($0, 1$) برای \lr{D0} و دو شکاف دوم ($2, 3$) برای \lr{D1} اختصاص داده شده‌اند. این طراحی به هسته اجازه می‌دهد تا دو \lr{AP} معلق اضافی به ازای هر آشکارساز را پیگیری کند، که تقریبی بهتر از مدل تک-\lr{AP} معلق قبلی است. وقتی یک \lr{AP} جدید باید برنامه‌ریزی شود و هر سه شکاف (یک اصلی + دو سرریز) پر هستند، جدیدترین \lr{AP} در شکاف سرریز با جدیدترین زمان بازنویسی می‌شود (تصمیمی که کمترین تأثیر را دارد زیرا \lr{AP} جدید احتمالاً زودتر از \lr{AP} سرریز قدیمی فعال می‌شود). این بازنویسی در شمارنده تشخیصی \texttt{DIAG\_AP\_OVERWRITE} ثبت می‌گردد.

کدهای وضعیت بازگشتی توسط ثابت‌های ماژولی \texttt{STATUS\_OK = 0}، \texttt{STATUS\_RNG\_EXHAUSTED = 1}، \texttt{STATUS\_OUTPUT\_BUFFER\_EXHAUSTED = 2}، \texttt{STATUS\_CARRY\_OVER\_BUFFER\_EXHAUSTED = 3} و \texttt{STATUS\_ITERATION\_CAP\_EXCEEDED = 4} تعریف شده‌اند. این کدها به مصرف‌کننده اجازه می‌دهند به‌صورت برنامه‌نویسی رفتار متناسب با وضعیت را اتخاذ کنند؛ مثلاً در صورت \texttt{STATUS\_RNG\_EXHAUSTED}، مصرف‌کننده می‌تواند بافر \lr{RNG} بزرگ‌تری تخصیص دهد و تلاش مجدد کند. شاخص‌های تشخیصی \texttt{DIAG\_DEADTIME}، \texttt{DIAG\_AFTERPULSE}، \texttt{DIAG\_DARKCOUNT}، \texttt{DIAG\_TOSSED\_CORE}، \texttt{DIAG\_IMD\_CLICKS}، \texttt{DIAG\_CARRY\_OVER}، \texttt{DIAG\_DEADTIME\_DROPPED} و \texttt{DIAG\_MULTIPLE\_FIRES\_PER\_SLOT} شاخص‌هایی در آرایه خروجی \texttt{diags} هستند که آمار تفکیکی از انواع رویدادها را در اختیار مصرف‌کننده قرار می‌دهند.

خروجی تابع، یک تاپل هشت‌عضوی به فرمت \lr{\texttt{(final\_times, final\_dets, final\_sigs, diags, state\_out, co\_count, status\_code)}} است. متغیر \lr{\texttt{state\_out}} خود یک ساختار داده‌ای جامع است که شامل $14$ مؤلفه برای به‌روزرسانی وضعیت آشکارساز می‌باشد: زمان‌های آخرین کلیک برای \lr{D0/D1}، زمان‌های \lr{AP} معلق برای \lr{D0/D1}، زمان‌های \lr{DC} بعدی برای \lr{D0/D1}، عمق‌های آبشار برای \lr{D0/D1}، اشاره‌گر \lr{RNG}، شمارنده‌های بازنویسی \lr{AP} و کپی‌هایی از آرایه‌های سرریز \lr{AP}. وجود این کپی‌ها ضروری است؛ زیرا آرایه‌های محلی ممکن است در بلاک زمانی بعدی تغییر کنند و وضعیت باید مستقل از آن‌ها حفظ شود.

\section{پیاده‌سازی ریاضی و الگوریتمی}

هسته پردازش از یک حلقه اصلی تشکیل شده که در هر تکرار، رویداد بعدی را با توجه به اولویت \lr{input > afterpulse > dark count} انتخاب و پردازش می‌کند. در ادامه، گزیده‌ای از کد انتخاب رویداد ارائه شده است:

\begin{LTR}
	\begin{lstlisting}[language=Python]
		while True:
		iteration += 1
		if iteration > max_iterations:
		status_code = STATUS_ITERATION_CAP_EXCEEDED
		break
		
		t_input = _SENTINEL_T
		if input_ptr < input_len:
		t_input = float(sorted_events[input_ptr]["time"])
		
		# Find the earliest in-batch AP across both detectors
		t_next_ap = _SENTINEL_T
		ap_det = -1
		if t_ap_0 >= 0.0 and t_ap_0 < time_end_abs and t_ap_0 < t_next_ap:
		t_next_ap = t_ap_0; ap_det = 0
		if t_ap_1 >= 0.0 and t_ap_1 < time_end_abs and t_ap_1 < t_next_ap:
		t_next_ap = t_ap_1; ap_det = 1
		
		# Filter past-batch-end DCs; track which detector wins
		_dc0_c = t_dc_0 if (t_dc_0 < _SENTINEL_T and t_dc_0 < time_end_abs) else _SENTINEL_T
		_dc1_c = t_dc_1 if (t_dc_1 < _SENTINEL_T and t_dc_1 < time_end_abs) else _SENTINEL_T
		if _dc0_c <= _dc1_c and _dc0_c < _SENTINEL_T:
		t_next_dc = _dc0_c; dc_det = 0
		elif _dc1_c < _SENTINEL_T:
		t_next_dc = _dc1_c; dc_det = 1
		else:
		t_next_dc = _SENTINEL_T; dc_det = -1
		
		# Early exit: no more events to process
		if (input_ptr >= input_len
		and (t_ap_0 < 0.0 or t_ap_0 >= time_end_abs)
		and (t_ap_1 < 0.0 or t_ap_1 >= time_end_abs)
		and (t_dc_0 >= _SENTINEL_T or t_dc_0 >= time_end_abs)
		and (t_dc_1 >= _SENTINEL_T or t_dc_1 >= time_end_abs)):
		break
		
		# Priority selection: input > afterpulse > dark count
		if t_input <= t_next_ap and t_input <= t_next_dc:
		# Input wins (priority 1)
		current_time = t_input
		current_det = int(sorted_events[input_ptr]["det"])
		src_idx = int(sorted_events[input_ptr]["idx"])
		is_sig = bool(sorted_events[input_ptr]["sig"])
		input_ptr += 1
		elif t_next_ap <= t_next_dc:
		# Afterpulse wins (priority 2 over dark count)
		current_time = t_next_ap
		current_det = ap_det
		is_ap = True
		# ...
		else:
		# Dark count wins
		current_time = t_next_dc
		current_det = dc_det
		is_dark = True
	\end{lstlisting}
\end{LTR}

این کد چندین ویژگی طراحی مهم را نشان می‌دهد. نخست، بررسی «خروج زودهنگام» \lr{(early exit)} پیش از بررسی فرسایش بافر انجام می‌شود؛ این ترتیب بر اساس اصلاح \lr{Review-v12 F13} تغییر کرده تا از خروج اشتباه با کد فرسایش وقتی حلقه به‌صورت تمیز پایان می‌یافت جلوگیری شود. دوم، مقایسه‌ها با $\leq$ دقیق انجام می‌شوند (نه با تلورانس \lr{EPS}) که این مقایسه‌ها را انتقالی \lr{(transitive)} و قابل بازتولید در پلتفرم‌های مختلف می‌سازد. سوم، فیلتر کردن \lr{AP}ها و \lr{DC}های پس از پایان بلاک به‌صورت صریح انجام می‌شود تا فقط رویدادهای درون-بلاک به‌عنوان کاندیدای «رویداد بعدی» انتخاب شوند.

پس از انتخاب رویداد، چندین بررسی متوالی انجام می‌شود. ابتدا بررسی زمان مرده:

\begin{LTR}
	\begin{lstlisting}[language=Python]
		if shared_spad:
		# Shared dead-time clock: most recent fire on EITHER detector
		last_fire = t_fire_0 if t_fire_0 > t_fire_1 else t_fire_1
		else:
		last_fire = t_fire_0 if current_det == 0 else t_fire_1
		
		# Non-paralyzable dead-time check
		if (current_time - last_fire) < dt_ns:
		# Distinguish batch-end input events (DIAG_DEADTIME_DROPPED)
		# from in-batch events (DIAG_DEADTIME)
		if current_time >= time_end_abs and not is_dark and not is_ap:
		diags[DIAG_DEADTIME_DROPPED] += 1
		continue
		diags[DIAG_DEADTIME] += 1
		
		if is_ap:
		# Reschedule the AP after the dead-time window
		if ap_reschedule_uniform:
		u_release = rng_floats[rng_ptr]; rng_ptr += 1
		release_time = last_fire + dt_ns + u_release * ap_lifetime
		else:
		release_time = last_fire + dt_ns
		rescheduled = release_time + _KERNEL_AP_RESCHEDULE_SLACK_NS
		# ... restore to t_ap_0 or t_ap_1
		continue
		
		# In-batch input/dark events in dead time: silently dropped
		continue
	\end{lstlisting}
\end{LTR}

این بخش نشان می‌دهد که چگونه مدل زمان مرده غیر-پارالایز پیاده‌سازی شده: اگر رویداد در پنجره مرده باشد، بدون تأثیر بر ساعت زمان مرده سرکوب می‌شود. تمایز مهمی میان رویدادهای ورودی در پایان بلاک (که در \texttt{DIAG\_DEADTIME\_DROPPED} شمارش می‌شوند اما در \lr{carry-over} قرار نمی‌گیرند) و رویدادهای درون-بلاک (که در \texttt{DIAG\_DEADTIME} شمارش می‌شوند) وجود دارد. برای پالس‌های پسین که در زمان مرده فعال می‌شوند، سیاست \lr{END} یا \lr{UNIFORM} برای برنامه‌ریزی مجدد اعمال می‌گردد؛ یک \texttt{slack} کوچک (\lr{\_KERNEL\_AP\_RESCHEDULE\_SLACK\_NS = 1e-6}) برای جلوگیری از شرط مرزی دقیق $t = t_{last\_fire} + \tau_d$ اضافه می‌شود.

تولید پالس پسین پس از هر کلیک موفق با کد زیر انجام می‌شود:

\begin{LTR}
	\begin{lstlisting}[language=Python]
		if ap_prob > 0.0:
		u_fire = rng_floats[rng_ptr]; rng_ptr += 1
		if u_fire < ap_prob:
		# Determine the current cascade depth for this detector
		if is_ap:
		# AP-originated fire: depth = parent_depth + 1
		current_depth = _consumed_ap_depth + 1
		else:
		# Signal/dark/IMD-originated fire: first-generation AP
		current_depth = 1
		
		if (current_depth <= per_detector_depth_limit
		and global_ap_count < global_cascade_limit):
		u_decay = rng_floats[rng_ptr]; rng_ptr += 1
		u_decay = max(log_safe_eps, u_decay)
		
		if ap_model_geometric:
		if pulse_period_ns > EPS and p_geometric > 0.0:
		# Standard geometric sampling:
		# K = ceil(log(U)/log(1-p))
		raw_steps = int(math.ceil(math.log(u_decay) / log_1mp))
		steps = max(1, min(raw_steps, max_steps))
		decay = steps * pulse_period_ns
		else:
		decay = 0.0
		else:
		# Exponential sampling: decay = -ln(u) * tau
		decay = -math.log(u_decay) * ap_lifetime
		
		sched_time = current_time + max(0.0, decay)
		
		# Schedule the AP, handling overflow slots
		if current_det == 0:
		if t_ap_0 < 0.0:
		t_ap_0 = sched_time
		ap_depth_0 = current_depth
		else:
		# Try overflow slots 0, 1 for D0
		# ... (overflow management logic)
		# ...
		global_ap_count += 1
	\end{lstlisting}
\end{LTR}

این کد چندین نکته مهم را نشان می‌دهد. نخست، عمق آبشار به‌صورت صریح پیگیری می‌شود: اگر خود کلیک یک \lr{AP} باشد، عمق پالس پسین جدید $\text{depth}_{parent} + 1$ است؛ در غیر این صورت (سیگنال، شمارش تاریک یا \lr{IMD})، عمق از $1$ آغاز می‌شود. این تمایز برای جلوگیری از شروع فوری یک آبشار جدید پس از رسیدن به کران ضروری است. دوم، نمونه‌برداری هندسی از فرمول استاندارد $K = \lceil \ln(u)/\ln(1-p) \rceil$ استفاده می‌کند که $P(K = k) = (1-p)^{k-1} \cdot p$ را برای $k = 1, 2, \ldots$ تولید می‌کند. سوم، کران \texttt{max\_steps} (که تطبیقی با طول عمر و مدت بلاک است) از نمونه‌برداری تأخیرهای بسیار بزرگ جلوگیری می‌کند. چهارم، مدیریت \lr{overflow} برای \lr{AP}های معلق اضافی پیاده‌سازی شده که تقریبی بهتر از مدل تک-\lr{AP} معلق قبلی ارائه می‌دهد.

محاسبه کران تکرار تطبیقی \lr{(adaptive iteration cap)} بر اساس تعداد تخمینی کل رویدادها انجام می‌شود:

\begin{equation}
	\text{max\_iterations} = \max\left(10^7,\; 100 \cdot \text{est\_total\_events}\right)
\end{equation}

که در آن $\text{est\_total\_events} = \text{input\_len} + (\text{dr}_0 + \text{dr}_1) \cdot 10^{-9} \cdot \text{batch\_duration} + 100$ است. این کران تطبیقی (که بر اساس اصلاح \lr{Review-v15 F-17} معرفی شده) از رسیدن زودرس به کران ثابت قبلی $10^7$ برای بلاک‌های بزرگ با $10^6$ پالس و شمارش تاریک بالا جلوگیری می‌کند. یک کران سخت بالا $10^8$ نیز اعمال می‌شود تا از مسیرهای پاتولوژیک جلوگیری شود.

محاسبه پارامتر هندسی $p$ با کد زیر انجام می‌شود:

\begin{LTR}
	\begin{lstlisting}[language=Python]
		if ap_model_geometric and ap_lifetime > EPS and pulse_period_ns > EPS:
		p_geometric = 1.0 - math.exp(-pulse_period_ns / ap_lifetime)
		# Clamp to [0, 1 - EPS] to avoid log1p(-1) = -inf
		if p_geometric < 0.0:
		p_geometric = 0.0
		elif p_geometric >= 1.0 - EPS:
		p_geometric = 1.0 - EPS
		if p_geometric <= 0.0:
		log_1mp = 0.0  # degenerate; effectively no afterpulse
		else:
		log_1mp = math.log1p(-p_geometric)
		else:
		p_geometric = 0.0
		log_1mp = 0.0
	\end{lstlisting}
\end{LTR}

محاسبهٔ \lr{\texttt{log\_1mp}} با استفاده از \lr{\texttt{math.log1p(-p)}} به‌جای \lr{\texttt{math.log(1 - p)}}، به‌منظور بهبود دقت عددی در نزدیکی $p = 0$ انجام می‌شود. اعمال کران $1 - \varepsilon$ از بروز خطای $\log(0) = -\infty$ جلوگیری می‌کند. این پارامترها یک‌بار در ابتدای هر فراخوانی هسته محاسبه شده و در طول حلقهٔ اصلی استفاده می‌شوند.


پوشش کامپایل \lr{Numba} با کد زیر پیاده‌سازی شده:

\begin{LTR}
	\begin{lstlisting}[language=Python]
		# Default to the pure-Python implementation
		_process_events_numba = _process_events_logic
		
		if NUMBA_AVAILABLE:
		try:
		_process_events_numba = njit(cache=True, fastmath=False)(
		_process_events_logic
		)
		# Force a one-shot compile to surface compilation errors eagerly
		_seed_events = np.zeros(2, dtype=event_np_dtype)
		_seed_events["time"] = (1.0, 2.0)
		_seed_events["det"] = (0, 1)
		_seed_events["idx"] = (0, 1)
		_seed_events["sig"] = (True, False)
		_seed_cp = np.zeros(CP_LEN, dtype=np.float64)
		_seed_cp[CP_PULSE_PERIOD] = 1.0
		_seed_cp[CP_DEAD_TIME] = 50.0
		_seed_cp[CP_BATCH_DURATION] = 1000.0
		_seed_cp[CP_AP_LIFETIME] = 50.0
		_seed_cp[CP_MAX_CASCADE] = 5.0
		_ = _process_events_numba(
		_seed_events,
		-1.0e9, -1.0e9, -1.0, -1.0, -1.0, -1.0,
		_seed_cp,
		np.zeros(32, dtype=np.float64),
		np.zeros(8, dtype=event_np_dtype),
		np.zeros(16, dtype=np.float64),
		np.zeros(16, dtype=np.int64),
		np.zeros(16, dtype=np.bool_),
		np.zeros(2, dtype=np.int64),
		np.full(4, -1.0, dtype=np.float64),
		np.zeros(4, dtype=np.int64),
		)
		except Exception as e:
		logger.warning(
		"Failed to compile Numba logic, falling back to pure Python: %s",
		e,
		)
		_process_events_numba = _process_events_logic
	\end{lstlisting}
\end{LTR}

این پیاده‌سازی چندین ویژگی دارد. نخست، \texttt{njit(cache=True)} کش دیسک را فعال می‌کند تا در اجراهای بعدی، زمان کامپایل اولیه حذف شود. دوم، \texttt{fastmath=False} از بازترتیب عملیات \lr{IEEE-754-violating} توسط کامپایلر جلوگیری می‌کند و در نتیجه نتایج \lr{bit-identical} (در محدوده عملیاتی) میان مسیر \lr{Numba} و \lr{pure-Python} را تضمین می‌نماید. سوم، کامپایل اجباری با ورودی‌های کوچک در زمان \lr{import} خطاهای کامپایل را زودهنگام آشکار می‌کند و کش دیسک را پر می‌کند. چهارم، در صورت شکست کامپایل، با هشدار لاگ به \lr{pure-Python} بازمی‌گردد تا ماژول همواره قابل استفاده باشد. این طراحی \lr{graceful} برای محیط‌هایی که \lr{Numba} در آن‌ها نصب نیست یا نسخه ناسازگار دارد، حیاتی است.

\section{ارسال و دریافت با سایر ماژول‌ها}

این ماژول به‌عنوان هستهٔ پردازشی درونی، در نقش یک سرویس‌دهندهٔ سطح پایین برای ماژول سطح بالای \lr{\texttt{qkd.detectors}} عمل کرده و رابط‌های مشخصی با ماژول‌های هم‌راستا دارد. مصرف‌کنندهٔ اصلی این ماژول، متد \lr{\texttt{\_simulate\_event\_driven()}} از کلاس \lr{SinglePhotonDetector} است که هسته را به این صورت فراخوانی می‌کند: نخست، آرایه‌های رویداد ساختاریافته را با استفاده از متد کمکی \lr{\texttt{\_build\_event\_array()}} می‌سازد؛ سپس آرایهٔ \lr{\texttt{const\_params}} را با $15$ مؤلفه از فیلدهای اعتبارسنجی‌شدهٔ آشکارساز پر می‌کند؛ سپس بافر \lr{RNG} را با \lr{\texttt{rng.random(est\_events * rng\_buffer\_mult)}} پیش‌بارگذاری می‌کند؛ و در نهایت هسته را فراخوانی کرده و خروجی را به‌صورت ساختاریافته پردازش می‌نماید. این الگوی فراخوانی، جداسازی صریحی میان لایهٔ بالا (ساخت پارامترها و پردازش پسین) و لایهٔ پایین (پردازش رویداد) ایجاد می‌کند.


از منظر ورودی، هسته چندین نوع داده دریافت می‌کند. آرایه \texttt{sorted\_events} که از \texttt{qkd.detector\_datastructures} با \lr{dtype} \texttt{event\_np\_dtype} وارد می‌شود، شامل رویدادهای کلیک سیگنال (که از خروجی کانال فیبر و آشکارساز منبع مشتق شده‌اند) و رویدادهای کلیک \lr{IMD/FWM} (که از تحلیل \lr{SCM/WDM} می‌آیند) است. این رویدادها پیش از فراخوانی هسته بر اساس زمان مطلق مرتب می‌شوند تا هسته بتواند به‌صورت کارآمد با یک عبور خطی آن‌ها را پردازش کند. آرایه \texttt{cascade\_depths\_in} که عمق آبشار فعلی \lr{AP} معلق برای هر آشکارساز را در بر دارد، از \texttt{\_DetectorState} آشکارساز استخراج می‌شود.

از منظر خروجی، هسته چندین ساختار داده تولید می‌کند. آرایه‌های \texttt{out\_times}، \texttt{out\_dets} و \texttt{out\_sigs} (که به‌عنوان بافر از سوی فراخواننده تخصیص داده می‌شوند و هسته تنها آن‌ها را پر می‌کند) حاوی زمان‌های کلیک، شناسه آشکارساز و پرچم منشأ سیگنال برای رویدادهای کلیک موفق هستند. این آرایه‌ها در متد \texttt{\_simulate\_event\_driven()} به مرحله بن کردن \lr{(binning)} ارسال می‌شوند که در آن زمان‌های کلیک با لرزش گوسی به پالس‌های متناظر نگاشته می‌شوند و در نهایت آرایه‌های بولی \texttt{click0} و \texttt{click1} تولید می‌گردند. آرایه \texttt{diags} شامل $8$ شمارنده تشخیصی است که در شیء \texttt{DetectionDiagnostics} کپی می‌شوند و به مصرف‌کنندگان نهایی (مانند تحلیل کلید امن) ارائه می‌گردند.

تاپل \texttt{state\_out} که حاوی $14$ مؤلفه است، برای به‌روزرسانی \texttt{\_DetectorState} آشکارساز استفاده می‌شود. این به‌روزرسانی به‌صورت زیر انجام می‌شود: زمان‌های آخرین کلیک \lr{D0/D1} در فیلدهای \texttt{last\_abs\_fire\_time\_ns\_d0/d1} ذخیره می‌شوند؛ زمان‌های \lr{AP} معلق \lr{D0/D1} در \texttt{pending\_ap\_time\_d0/d1}؛ عمق‌های آبشار در \texttt{ap\_depth\_d0/d1}؛ زمان‌های \lr{DC} بعدی در \texttt{next\_dc\_time\_d0/d1}؛ و آرایه‌های سرریز \lr{AP} در فیلدهای مخصوص به خود. این به‌روزرسانی پیوسته تضمین می‌کند که بلاک بعدی شبیه‌سازی می‌تواند حالت دقیق را از بلاک قبلی به ارث ببرد و پیوستگی فیزیکی زمان مرده، پالس پسین و شمارش تاریک را حفظ کند.

از منظر وابستگی‌های درون‌پروژه‌ای، این ماژول به‌صورت سختگیرانه به ماژول \texttt{qkd.detector\_datastructures} وابسته است و در صورت عدم دسترسی، به‌جای \lr{fallback} خاموش، استثنای \lr{ImportError} صادر می‌کند. این وابستگی شامل دو نهاد است: ثابت \texttt{\_SENTINEL\_T} (یک مقدار شناور بسیار بزرگ که به‌عنوان سنتینل «هیچ رویداد» استفاده می‌شود) و \texttt{event\_np\_dtype} (یک \lr{NumPy structured dtype} که شکل هر رویداد را تعریف می‌کند). ثابت‌های شاخص \texttt{CP\_*} و \texttt{DIAG\_*} و \texttt{STATUS\_*} نیز از ماژول \texttt{qkd.detector\_datastructures} (یا ماژول خواهری آن) واردات می‌شوند تا یکپارچگی شاخص‌ها در سراسر فریم‌ورک تضمین شود. این طراحی یک-منبع \lr{(single source of truth)} برای ثابت‌های شاخص از واگرایی احتمالی جلوگیری می‌کند.

تعامل با کتابخانه \lr{Numba} نیز از طریق \texttt{try/except ImportError} انجام می‌شود. در صورت دسترس بودن \lr{Numba}، دکوراتور \texttt{njit} به‌صورت پویا واردات می‌شود و در غیر این صورت، یک دکوراتور \lr{no-op} جایگزین تعریف می‌گردد. این رویکرد امکان استفاده از ماژول را در محیط‌هایی که \lr{Numba} نصب نیست (مانند برخی محیط‌های \lr{CI} سبک) فراهم می‌سازد. تصمیم آگاهانه برای عدم اعلام صریح امضا روی تابع هسته \lr{(no explicit signature on the hot-loop kernel)} برای جلوگیری از وابستگی‌های شکننده \lr{ABI} به \texttt{numba.types.Tuple} و \texttt{numba.types.Array} که در نسخه‌های مختلف \lr{Numba} تغییر کرده‌اند، اتخاذ شده است. این طراحی به ماژول اجازه می‌دهد با طیف وسیع‌تری از نسخه‌های \lr{Numba} سازگار بماند.

لاگر ماژول (\texttt{logging.getLogger(\_\_name\_\_)}) در سطح \lr{WARNING} برای گزارش شکست کامپایل \lr{Numba} و در سطح \lr{INFO} برای گزارش رویدادهای مهم (مانند بازگشت به \lr{pure-Python}) به کار می‌رود. این سطح‌بندی به کاربران اجازه می‌دهد با تنظیم سطح لاگ، میزان جزئیات گزارش‌ها را کنترل کنند. لاگ سطح \lr{DEBUG} در این ماژول استفاده نمی‌شود زیرا هسته در حلقه داغ اجرا می‌شود و هرگونه لاگ اضافی به‌طور قابل‌توجهی عملکرد را کاهش می‌دهد؛ به‌جای آن، آمار تفکیکی در آرایه \texttt{diags} جمع‌آوری می‌شود که پس از پایان اجرای هسته توسط مصرف‌کننده قابل بررسی است.

در نهایت، این ماژول به‌صورت مستقیم با مکانیزم \lr{RNG state snapshot/restore} ماژول \texttt{qkd.detectors} تعامل دارد. مصرف‌کننده پیش از فراخوانی هسته، حالت \lr{RNG} را با \texttt{\_copy\_rng\_state(rng.bit\_generator.state)} عکس‌برداری می‌کند و در صورت نیاز به تلاش مجدد با بافر بزرگ‌تر، حالت را بازیابی می‌نماید. این مکانیزم تضمین می‌کند که تلاش‌های مجدد با همان بذر \lr{RNG} نتایج یکسانی تولید کنند (تا جایی که بافر \lr{RNG} بزرگ‌تر اجازه می‌دهد). نکته مهم این است که پس از اصلاح \lr{Review-v12 F1}، حالت \lr{RNG} دیگر پس از \lr{pre-draw} بازگردانده نمی‌شود؛ زیرا این بازگردانی به‌صورت قطعی کش‌های \lr{RNG} هسته و پس-هسته را همبسته می‌کرد و آمار \lr{ISI}، \lr{gate-loss} و \lr{double-click} را تحریف می‌نمود. این تصمیم طراحی برای بازتولیدپذیری صحیح آماری حیاتی است.


\sloppy
\emergencystretch=5em
\hbadness=10000
\vbadness=10000
\tolerance=2000
\providecolor{codebg}{RGB}{248,248,248}
\providecolor{codeframe}{RGB}{200,200,200}
\providecolor{kwblue}{RGB}{0,0,140}
\providecolor{strred}{RGB}{163,21,21}
\providecolor{cmtgreen}{RGB}{0,110,0}
\providecolor{emphpurple}{RGB}{110,0,140}
\lstset{
	basicstyle=\ttfamily\footnotesize,
	breakatwhitespace=false,
	breaklines=true,
	captionpos=b,
	keepspaces=true,
	numbers=left,
	numbersep=5pt,
	numberstyle=\tiny\color{gray},
	showspaces=false,
	showstringspaces=false,
	showtabs=false,
	tabsize=2,
	language=Python,
	frame=single,
	framesep=2pt,
	rulecolor=\color{gray!50},
	backgroundcolor=\color{gray!5},
	keywordstyle=\color{kwblue}\bfseries,
	commentstyle=\color{cmtgreen}\itshape,
	stringstyle=\color{strred},
	emph={math, channel, get_strict_mode, ParameterValidationError, PhysicalSuspicionWarning, _warn_physical_suspicion, is_close, is_finite_non_negative, AttenuationConfig, validate_parameter, validate_channel, check_type, validate_parameter_ranges, check_meaningful_alpha, check_meaningful_distance, check_total_loss_finite, check_attenuation_config_contract, check_wavelength_consistency, validate_wavelength_range, compute_transmittance, effective_total_loss_db, transmittance_from_loss_db, check_transmittance_health},
	emphstyle=\color{emphpurple}\bfseries,
}

\section{نمای کلی و هدف ماژول}

ماژول \lr{\texttt{qkd.\_gating}} یک لایه اعتبارسنجی و گارد \lr{(validation / gating layer)} است که به‌صورت هدفمند از کلاس اصلی \lr{\texttt{FiberChannel}} استخراج شده تا نگرانی‌های مرتبط با کنترل سلامت فیزیکی و سازگاری پیکربندی کانال را در یک مجموعه منسجم از توابع مستقل متمرکز کند. این ماژول در راستای اصل مهندسی نرم‌افزارِ تفکیک نگرانی‌ها \lr{(separation of concerns)} طراحی شده و هر تابع آن به‌جای ارجاع به \texttt{self} از پارامتر صریح \texttt{channel} به‌عنوان آرگومان اول استفاده می‌کند؛ این طراحی امکان آزمون‌پذیری واحد \lr{(unit testability)}، استفاده مجدد در زمینه‌های متفاوت (مانند اعتبارسنجی پس از \lr{deserialization}) و امکان جایگزینی آسان در \lr{mock}های تست را فراهم می‌سازد. هدف نهایی، فراهم‌آوردن یک لایه دفاعی چندگانه \lr{(defense in depth)} است که هرگونه پیکربندی غیرفیزیکی کانال را پیش از ورود به حلقه اصلی شبیه‌سازی شناسایی و رد کند.

اهمیت این ماژول از آنجا ناشی می‌شود که فریم‌ورک شبیه‌سازی \lr{QKD} به‌عنوان یک ابزار تحقیقاتی-آکادمیک، همواره با پیکربندی‌های متنوع و گاهی غیرعادی از سوی کاربران مواجه است؛ مقادیر منفی برای طول فیبر، ضرایب تضعیف غیرمنطقی (مانند $\alpha=0$ با طول غیرصفر که فیبر بی‌تلفاتی را بر روی فاصله متناهی ارائه می‌دهد و از نظر فیزیکی محال است)، طول‌موج‌های خارج از بازه مخابراتی، و تطبیق نادرست میان \texttt{total\_loss\_db} و حاصلضرب $\alpha \cdot L$ تنها چند نمونه از خطاهای متداول هستند. اگر این خطاها به‌صورت خاموش از لایه اعتبارسنجی عبور کنند، در محاسبات پایین‌دستی مانند فرمول‌های بازده \lr{decoy-state}، نرخ کلید امن و \lr{QBER} به نتایج ظاهراً معتبر اما فیزیکی غلط منجر می‌شوند که می‌توانند ماه‌ها تلاش تحقیقاتی را بی‌اثر کنند. این ماژول با سیاست دو-سطحی \lr{strict/non-strict} (که توسط \lr{ContextVar} جهانی \texttt{get\_strict\_mode()} کنترل می‌شود) به کاربر اجازه می‌دهد در حالت پیش‌فرض سخت‌گیرانه هرگونه تخلف را به‌صورت استثنا مشاهده کند و در حالت اکتشافی \lr{(exploratory)} صرفاً هشدار دریافت نماید.

ماژول \lr{\_gating} شامل چهارده تابع عمومی است که در \texttt{\_\_all\_\_} فهرست شده‌اند: \texttt{check\_type}، \texttt{validate\_parameter\_ranges}، \texttt{check\_meaningful\_alpha}، \texttt{check\_meaningful\_distance}، \texttt{check\_total\_loss\_finite}، \texttt{check\_attenuation\_config\_contract}، \texttt{check\_wavelength\_consistency}، \texttt{validate\_wavelength\_range}، \texttt{compute\_transmittance}، \texttt{effective\_total\_loss\_db}، \texttt{transmittance\_from\_loss\_db}، \texttt{check\_transmittance\_health}، \texttt{validate\_parameter} و \texttt{validate\_channel}. هر یک از این توابع به یک جنبه خاص از سلامت کانال می‌پردازند و در مجموع یک شبکه اعتبارسنجی کامل را تشکیل می‌دهند که از نوع شیء \texttt{attenuation} آغاز شده و به انسجام حالت مشتق‌شده \lr{(derived state)} مانند \texttt{transmittance} ختم می‌شود.

این ماژول به‌عنوان زیرسیستمی فرعی از معماری بزرگ‌تر \lr{qkd.channel} عمل می‌کند و در خط لوله ساخت کانال به‌صورت متوالی فراخوانی می‌شود. در معماری کلی، پس از ساخت اولیه شیء \lr{FiberChannel} (که خود یک \lr{dataclass} ایستا با فیلدهای \texttt{attenuation}، \texttt{wavelength\_nm} و \texttt{\_total\_loss\_db\_override} است)، متد \texttt{\_\_post\_init\_\_} این کلاس توابع گیتینگ را به‌ترتیب زیر فرامی‌خواند: نخست \texttt{check\_type}، سپس \texttt{validate\_parameter\_ranges}، سپس \texttt{check\_meaningful\_alpha} و \texttt{check\_meaningful\_distance}، سپس \texttt{check\_total\_loss\_finite} و \texttt{check\_attenuation\_config\_contract}، و در نهایت \texttt{check\_wavelength\_consistency}، \texttt{validate\_wavelength\_range} و \texttt{check\_transmittance\_health}. این ترتیب به‌گونه‌ای انتخاب شده که هر گیت بر اساس خروجی‌های گیت‌های قبلی عمل کند و در صورت بروز خطا در هر مرحله، مراحل بعدی اجرا نشوند.

\section{مبانی تئوری و فیزیکی}

بنای نظری این ماژول بر چند اصل بنیادین فیزیک فیبر نوری و نظریه اعتبارسنجی عددی استوار است. نخست، قانون بیر-لامبرت \lr{(Beer-Lambert law)} که پایه محاسبه انتقال \lr{(transmittance)} کانال است. این قانون بیان می‌کند که شدت نور پس از عبور از یک محیط با ضریب تضعیف $\alpha$ به طول $L$ به‌صورت نمایی کاهش می‌یابد:

\begin{equation}
	T = \frac{I(L)}{I(0)} = 10^{-\alpha L / 10}
\end{equation}

که در آن $\alpha$ بر حسب \lr{dB/km}، $L$ بر حسب \lr{km} و $T$ یک عدد بی‌بعد در بازه $[0, 1]$ است. در این ماژول، تابع \texttt{transmittance\_from\_loss\_db()} این محاسبه را با استفاده از \texttt{math.pow} (به‌جای عملگر \texttt{**}) پیاده‌سازی می‌کند تا در برابر \lr{OverflowError} در پلتفرم‌های مختلف مقاوم باشد. یک نکته ظریف اما مهم این است که برای $L_{total} = 0$ (یعنی فیبر بی‌تلفات با طول صفر)، انتقال برابر $1$ بازگردانده می‌شود، اما این شاخه با استفاده از برابری دقیق $L_{total} == 0.0$ (به‌جای \texttt{is\_close} با تلورانس $10^{-12}$) ارزیابی می‌شود. این انتخاب عمدی است: تلورانس $10^{-12}$ قبلی برای $\alpha = 0.2$ و $L = 5 \times 10^{-12}$ km، مقدار غیرصفر $L_{total}$ را به‌صورت خاموش صفر تلقی می‌کرد و انتقال نادرست $1.0$ بازمی‌گرداند.

دومین اصل، رفتار اعداد ممیز-شناور \lr{IEEE-754} و پدیده‌های پاتولوژیک آن، مانند صفر منفی \lr{-0.0} و اعداد زیرنرمال \lr{subnormal} است. در استاندارد \lr{IEEE-754}، رابطهٔ $-0.0 == 0.0$ برقرار است، اما این دو مقدار از نظر فیزیکی متفاوت تلقی می‌شوند؛ زیرا \lr{-0.0} معمولاً نشان‌دهندهٔ یک خطای علامت در بالادست \lr{upstream sign error} یا حاصلِ سرریزِ یک مقدار منفیِ بسیار کوچک (مانند $-1 \times 10^{-324}$ که به \lr{-0.0} تبدیل می‌شود) است. این ماژول با استفاده از \lr{\texttt{math.copysign(1.0, value) < 0}} به‌صورت صریح \lr{-0.0} را تشخیص داده و در حالت سخت‌گیرانه، استثنا صادر می‌کند. به‌طور مشابه، اعداد زیرنرمال \lr{subnormal floats} که کوچک‌تر از \lr{\texttt{sys.float\_info.min}} هستند، نهایتاً یک بیت دقت مانتیس دارند و در فرمول‌های \lr{QKD} (مانند محاسبهٔ $Y_n$ یا $e_1$) به نتایج کاملاً نادرست منجر می‌شوند؛ لذا این ماژول هرگونه عدد زیرنرمال را به‌صورت غیرقابل‌بازگشت رد می‌کند.


سومین اصل، مدل‌سازی فیزیکی ضرایب تضعیف بر اساس طول‌موج است. فیبرهای سیلیکایی در طول‌موج‌های استاندارد مخابراتی دارای ضرایب تضعیف مشخصی هستند: در $850$ \lr{nm} (پنجره اول) نزدیک به $2.5$ \lr{dB/km}، در $1310$ \lr{nm} (پنجره دوم) نزدیک به $0.35$ \lr{dB/km}، و در $1550$ \lr{nm} (پنجره سوم) نزدیک به $0.20$ \lr{dB/km}. این ماژول با استفاده از دیکشنری \texttt{TYPICAL\_WAVELENGTH\_ALPHA\_RANGES} بازه‌های مجاز را تعریف می‌کند و تابع \texttt{check\_wavelength\_consistency()} نزدیک‌ترین طول‌موج شناخته‌شده را با \texttt{WAVELENGTH\_CHECK\_TOLERANCE\_NM} پیدا کرده و اعتبارسنجی می‌کند. این بررسی از ارسال پیکربندی‌هایی مانند $\alpha = 0.05$ \lr{dB/km} در $1550$ \lr{nm} (که از نظر فیزیکی محال است زیرا کم‌تلفات‌ترین فیبر تجاری در این طول‌موج نزدیک به $0.15$ \lr{dB/km} است) جلوگیری می‌کند.

چهارمین اصل، مفهوم قرارداد ساختاری \lr{(structural contract)} میان \lr{dataclass}های مرتبط است. در این فریم‌ورک، \texttt{AttenuationConfig} یک \lr{dataclass} با سه فیلد اصلی \texttt{fiber\_length}، \texttt{attenuation\_coefficient} و \texttt{total\_loss\_db} است و قرارداد آن این است که \texttt{total\_loss\_db} باید برابر با \texttt{fiber\_length * attenuation\_coefficient} باشد. این قرارداد به‌صورت ضمنی توسط تمامی کدهای پایین‌دستی فرض می‌شود؛ به‌ویژه، محاسبه انتقال $T = 10^{-L_{total}/10}$ به این قرارداد وابسته است. تابع \texttt{check\_attenuation\_config\_contract()} این قرارداد را با تلورانس نسبی \lr{\_CONTRACT\_CHECK\_REL\_TOL} و تلورانس مطلق \lr{\_CONTRACT\_CHECK\_ABS\_TOL} (که هر دو در \texttt{\_schema} تعریف شده‌اند) اعتبارسنجی می‌کند. یک جنبه ظریف این است که وقتی \texttt{\_total\_loss\_db\_override} تنظیم شده باشد (یعنی کاربر به‌صورت دستی مقدار \texttt{total\_loss\_db} را تنظیم کرده باشد)، قرارداد همچنان بررسی می‌شود اما با تلورانس نسبی سخت‌گیرانه $10^{-9}$؛ زیرا فرض می‌شود \texttt{override} باید نهایتاً چند \lr{ULP} از $\alpha \cdot L$ فاصله داشته باشد و انحراف بزرگ‌تر نشانگر یک پیکربندی ناسازگار است.

پنجمین اصل، مدیریت هشدارهای فیزیکی-مشکوک \lr{(physically suspicious)} است. این ماژول دو سطح پاسخ را تفکیک می‌کند: استثنای سخت \lr{(ParameterValidationError)} که در حالت سخت‌گیرانه برای خطاهای قطعی (مانند $T = $ NaN یا $\alpha = 0$ با $L > 0$) صادر می‌شود، و هشدار نرم \lr{(PhysicalSuspicionWarning)} که از طریق تابع کمکی \texttt{\_warn\_physical\_suspicion()} برای موارد مشکوک اما نه قطعاً غلط (مانند $\alpha < 0.05$ بدون تنظیم طول‌موج) منتشر می‌شود. این هشدارها به‌صورت پیش‌فرض توسط پایتون نمایش داده می‌شوند (زیرا از \lr{warnings.warn} با کلاس هشدار غیر-\lr{DeprecationWarning} استفاده می‌کنند) و در عین حال توسط کاربر قابل فیلتر هستند.

\section{تحلیل معماری و ساختار داده‌ها}

معماری این ماژول تابع-محور \lr{(function-oriented)} است و از الگوی تزریق وابستگی \lr{(dependency injection)} برای جداسازی منطق اعتبارسنجی از شیء کانال استفاده می‌کند. هر تابع به‌جای دسترسی به فیلدهای \texttt{self}، یک شیء \texttt{channel} (از نوع \lr{Any}) دریافت می‌کند و از طریق ویژگی‌های آن شیء (\texttt{channel.attenuation}، \texttt{channel.wavelength\_nm}، \texttt{channel.\_total\_loss\_db\_override} و غیره) به داده‌ها دسترسی پیدا می‌کند. این طراحی مزایای متعددی دارد: نخست، توابع قابل استفاده مجدد برای هر شیء‌ای که پروتکل \lr{duck-typed} مشابهی داشته باشد (نه فقط \lr{FiberChannel} اصلی) هستند؛ دوم، توابع به‌راحتی با اشیاء \lr{mock} قابل آزمون هستند؛ سوم، وابستگی‌های ماژول حداقل است و تنها به ماژول‌های \texttt{qkd.exceptions}، \texttt{qkd.constants}، \texttt{qkd.datatypes} و \texttt{qkd.\_schema} نیاز دارد.

ورودی‌های ماژول از طریق \texttt{import} از ماژول \lr{\_schema} فراهم می‌شوند و شامل چندین دسته از ثابت‌ها و توابع هستند. ثابت‌های فیزیکی و سخت‌گیرانه شامل: \texttt{TRANSMITTANCE\_SUBNORMAL\_THRESHOLD} (آستانه تشخیص اعداد زیر-نرمال، معمولاً $\approx 2.2 \times 10^{-308}$)، \texttt{MAX\_DISTANCE\_KM} و \texttt{WARNING\_DISTANCE\_KM} (کران‌های فیزیکی برای طول فیبر)، \texttt{MAX\_FIBER\_LOSS\_DB\_KM} (کران فیزیکی برای ضریب تضعیف)، \texttt{TYPICAL\_WAVELENGTH\_ALPHA\_RANGES} (دیکشنری بازه‌های مجاز $\alpha$ بر اساس طول‌موج)، و \texttt{WAVELENGTH\_CHECK\_TOLERANCE\_NM} (تلورانس جست‌وجوی طول‌موج) هستند. ثابت‌های تلورانس عددی شامل: \texttt{\_TRANSMITTANCE\_UNDERFLOW\_ABS\_TOL} (که برابر با \lr{\texttt{sys.float\_info.min}} است و برای مقایسه انتقال‌های زیر-نرمال به کار می‌رود)، \texttt{\_CONTRACT\_CHECK\_REL\_TOL} و \texttt{\_CONTRACT\_CHECK\_ABS\_TOL} (برای اعتبارسنجی قرارداد \lr{AttenuationConfig})، و \texttt{\_ZERO\_LOSS\_ABS\_TOL} هستند.

ماژول از مکانیزم \lr{ContextVar} پایتون برای کنترل حالت سخت‌گیرانه استفاده می‌کند. تابع \texttt{get\_strict\_mode()} مقدار بولی این حالت را برمی‌گرداند و در تمامی توابع گیتینگ برای تصمیم‌گیری میان صدور استثنا یا هشدار به کار می‌رود. این مکانیزم به‌جای \lr{ClassVar} از \lr{ContextVar} استفاده می‌کند تا امکان تنظیم محلی \lr{(thread-local / context-local)} حالت سخت‌گیرانه را در تست‌ها و شبیه‌سازی‌های موازی فراهم کند؛ به این معنا که یک تست می‌تواند بدون تأثیر بر سایر تست‌ها، حالت سخت‌گیرانه را غیرفعال کند. علاوه بر این، \texttt{\_SKIP\_TRANSMITTANCE\_HEALTH} یک \lr{ContextVar} بولی دیگر است که اجازه می‌دهد در مسیرهای سریع ساخت \lr{(fast-construction paths)} از بررسی سلامت انتقال صرف‌نظر شود.

دیکشنری \texttt{VALIDATED\_UNITS} نقشه‌ای از نام پارامتر به واحد آن (مانند \texttt{\{"distance\_km": "km", "fiber\_loss\_db\_km": "dB/km", ...\}}) است که در پیام‌های خطا برای ارائه توضیحات کاربر-پسند استفاده می‌شود. این رویکرد به ماژول امکان می‌دهد بدون سخت‌کدنویسی واحد در هر پیام خطا، پیام‌های یکپارچه و قابل‌توسعه تولید کند. تابع کمکی \texttt{validate\_parameter(name, value, upper\_bound)} از این دیکشنری برای استخراج واحد استفاده می‌کند و در عین حال با بررسی ویژگی‌های \texttt{shape} و \texttt{dtype} ورودی، آرایه‌های \lr{NumPy/JAX/Torch} را به‌صورت زودهنگام رد می‌کند؛ زیرا \texttt{math.isfinite} بر روی آرایه‌ها با \lr{TypeError} مبهم مواجه می‌شود. این لایه دفاعی جلوی عبور آرایه‌ها به لایه‌های پایین‌تر اعتبارسنجی که انتظار اسکالر دارند را می‌گیرد.

استثنای اصلی ماژول، \lr{\texttt{ParameterValidationError}} است که از ماژول \lr{\texttt{qkd.exceptions}} فراخوانی شده و دارای امضای جامعی به صورت \lr{\texttt{(msg, *, param\_name=None, param\_value=None, context=None)}} می‌باشد. این امضای ساختاریافته به کاربر اجازه می‌دهد خطاها را به‌صورت برنامه‌نویسی پردازش کند (مانند فیلتر بر اساس \lr{\texttt{param\_name}}) و در عین حال، پیام‌های خطای خوانایی برای انسان تولید نماید. فیلد \lr{\texttt{context}} یک دیکشنری است که اطلاعات تکمیلی نظیر مقادیر مرتبط، مسیر ساخت و تلورانس‌های استفاده‌شده را در خود نگه می‌دارد؛ این داده‌ها برای دیباگ و گزارش‌گیری خودکار بسیار ارزشمند هستند.


لاگر ماژول از طریق \texttt{logger = logging.getLogger(\_\_name\_\_)} ساخته شده و نام آن به \texttt{"qkd.\_gating"} تنظیم می‌شود. این لاگر در سطح \lr{DEBUG} برای گزارش وضعیت اعتبارسنجی موفق و در سطح \lr{WARNING} برای هشدارهای غیرسخت‌گیرانه به کار می‌رود. این دوگانگی به کاربر اجازه می‌دهد با تنظیم سطح لاگ، میزان جزئیات گزارش‌ها را کنترل کند: سطح \lr{WARNING} (پیش‌فرض) تنها هشدارهای مهم را نشان می‌دهد، در حالی که سطح \lr{DEBUG} جریان کامل اعتبارسنجی را برای دیباگ نمایش می‌دهد.

\section{پیاده‌سازی ریاضی و الگوریتمی}

الگوریتم محاسبه انتقال در تابع \texttt{transmittance\_from\_loss\_db()} با چندین محافظت عددی پیاده‌سازی شده است. کد زیر نسخه کامل این تابع را نشان می‌دهد:

\begin{LTR}
	\begin{lstlisting}[language=Python]
		def transmittance_from_loss_db(total_db: float) -> float:
		if total_db == 0.0:
		# F-31 / P-06 / N-1: check for negative zero.
		# ``-0.0 == 0.0`` is True in IEEE 754, so this branch
		# activates for -0.0.  Negative zero can only arise from
		# (a) an upstream sign error, or (b) a very small
		# negative value underflowing (e.g. -1e-324 -> -0.0),
		# which would represent gain.  In both cases returning
		# transmittance=1.0 is physically unsafe, so raise
		# unconditionally.
		if math.copysign(1.0, total_db) < 0:
		raise ParameterValidationError(
		f"total_loss_db is negative zero (-0.0). ...",
		param_name="total_loss_db",
		param_value=total_db,
		)
		return 1.0
		if total_db < 0.0:
		raise ParameterValidationError(
		f"total_loss_db is negative ({total_db!r}), implying "
		f"gain rather than attenuation. This is physically "
		f"impossible for a passive fiber channel. ...",
		param_name="total_loss_db",
		param_value=total_db,
		)
		return math.pow(10.0, -total_db / 10.0)
	\end{lstlisting}
\end{LTR}

این پیاده‌سازی سه شاخه دارد. شاخه نخست ($L_{total} = 0$) با برابری دقیق ارزیابی می‌شود و در آن ابتدا علامت صفر با \texttt{copysign} بررسی می‌شود؛ اگر منفی باشد، استثنا صادر می‌شود زیرا $-0.0$ معمولاً نشانه یک خطای بالا-دست است. شاخه دوم ($L_{total} < 0$) به‌صورت غیرقابل-بازگشت رد می‌شود زیرا $L_{total} < 0$ به‌معنای تقویت \lr{(gain)} است که برای یک فیبر پسیو فیزیکاً محال است. شاخه سوم از \texttt{math.pow} به‌جای عملگر \texttt{**} استفاده می‌کند زیرا \texttt{math.pow} در برابر \lr{OverflowError} در پلتفرم‌های مختلف رفتار یکنواخت‌تری دارد. باید توجه داشت که برای مقادیر بسیار بزرگ $L_{total}$ (مانند $L_{total} > 308 \times 10 / \ln 10 \approx 1338$ dB)، \texttt{math.pow} به $0.0$ سرریز می‌شود که خود در تابع \texttt{check\_transmittance\_health()} به‌عنوان خطای \lr{underflow} شناسایی می‌شود.

تابع \texttt{check\_transmittance\_health()} چهار سطح کنترل سلامت را به‌صورت متوالی اعمال می‌کند. این تابع پس از محاسبه انتقال در مسیر ساخت کانال فراخوانی می‌شود و آن را برای پاتولوژی‌های مختلف بررسی می‌کند:

\begin{LTR}
	\begin{lstlisting}[language=Python]
		def check_transmittance_health(channel, transmittance):
		# NaN: reject unconditionally.
		if math.isnan(transmittance):
		raise ParameterValidationError(
		f"Transmittance computed as NaN for distance_km="
		f"{channel.attenuation.fiber_length!r}, ...",
		param_name="transmittance",
		param_value=transmittance,
		context={...},
		)
		
		# P-06: Detect negative-zero transmittance.
		if transmittance == 0.0 and math.copysign(1.0, transmittance) < 0:
		# ... issue PhysicalSuspicionWarning or raise in strict mode
		
		total_db = effective_total_loss_db(channel)
		
		# Underflow: exact-zero transmittance with non-zero total_db.
		if transmittance == 0.0 and total_db > 0.0:
		raise ParameterValidationError(
		f"Transmittance underflowed to exactly 0.0 for ...",
		param_name="transmittance", ...
		)
		
		# Subnormal transmittance.
		elif 0.0 < transmittance < TRANSMITTANCE_SUBNORMAL_THRESHOLD:
		raise ParameterValidationError(
		f"Transmittance {transmittance!r} is subnormal ...",
		param_name="transmittance", ...
		)
		
		# Range validation.
		if not (0.0 <= transmittance <= 1.0):
		raise ParameterValidationError(
		f"Transmittance {transmittance!r} is out of the physical "
		f"range [0.0, 1.0]. ...",
		param_name="transmittance", ...
		)
	\end{lstlisting}
\end{LTR}

این تابع به‌ترتیب چهار بررسی را انجام می‌دهد: \lr{NaN} (رد غیرقابل-بازگشت)، صفر منفی (هشدار یا رد بر اساس حالت سخت‌گیرانه)، \lr{underflow} (یعنی $T = 0$ با $L_{total} > 0$ که رد غیرقابل-بازگشت است زیرا در فرمول‌های \lr{QKD} منجر به تقسیم بر صفر می‌شود)، \lr{subnormal} (رد غیرقابل-بازگشت زیرا اعداد زیر-نرمال نهایتاً یک بیت دقت دارند و در محاسبات به نتایج تصادفی منجر می‌شوند) و در نهایت بررسی بازه $[0, 1]$. توجه شود که هیچ‌کدام از این بررسی‌ها به‌صورت خاموش مقدار را تصحیح نمی‌کنند (\lr{clamp} نمی‌کنند)؛ زیرا چنین تصحیحی می‌تواند خطای بالا-دست را پنهان کند و به نتایج ظاهراً معتبر اما فیزیکی غلط منجر شود. این سیاست صریح در داک‌استرینگ تابع به‌صورت \textit{We explicitly do NOT clamp silently} مستند شده است.

تابع \texttt{check\_attenuation\_config\_contract()} دارای دو مسیر است: مسیر \lr{override} (وقتی \texttt{\_total\_loss\_db\_override} تنظیم شده) و مسیر قرارداد استاندارد. در مسیر \lr{override}، تابع بررسی می‌کند که آیا \lr{override} به‌طور وحشتناکی با $\alpha \cdot L$ ناسازگار است یا خیر. این بررسی با تلورانس نسبی $10^{-9}$ انجام می‌شود که عمداً سخت‌گیرانه است زیرا فرض می‌شود \lr{override} نهایتاً چند \lr{ULP} از $\alpha \cdot L$ فاصله داشته باشد:

\begin{equation}
	\left| L_{override} - \alpha \cdot L \right| > 10^{-9} \cdot \max\left( \left| \alpha \cdot L \right|, 1 \right) \Rightarrow \text{خطا}
\end{equation}

در مسیر قرارداد استاندارد، یعنی در حالتی که \lr{\texttt{override}} فعال نیست، تابع با استفاده از تلورانس‌های \lr{\texttt{\_CONTRACT\_CHECK\_REL\_TOL}} و \lr{\texttt{\_CONTRACT\_CHECK\_ABS\_TOL}} بررسی می‌کند که مقدار \lr{\texttt{total\_loss\_db}} با مقدار $\alpha L$ سازگار باشد. نکتهٔ مهم در مسیر \lr{\texttt{override}} آن است که اگر $L = 0$، $\alpha = 0$ و مقدار \lr{\texttt{override}} نیز برابر با صفر باشد، تابع بدون ایجاد خطا بازمی‌گردد؛ زیرا این حالت متناظر با یک کانال بدیهی است که نیازی به بررسی ندارد. در مقابل، اگر $L > 0$ و $\alpha = 0$ باشد، یعنی با فیبری بی‌تلفات و دارای طول ناصفر مواجه باشیم، بررسی قرارداد همچنان انجام می‌شود؛ زیرا چنین حالتی از نظر فیزیکی نامعتبر است.

تابع \lr{\texttt{check\_wavelength\_consistency()}} با جست‌وجو در \lr{\texttt{TYPICAL\_WAVELENGTH\_ALPHA\_RANGES}}، نزدیک‌ترین طول‌موج شناخته‌شده را تعیین می‌کند. این جست‌وجو به‌صورت ترتیبی روی طول‌موج‌های مرتب‌شده انجام می‌شود تا در صورت تساوی فاصله، نزدیک‌ترین طول‌موج کوتاه‌تر به‌طور قطعی برگزیده شود. تلورانس جست‌وجو با استفاده از \lr{\texttt{math.nextafter(WAVELENGTH\_CHECK\_TOLERANCE\_NM, float("inf"))}} تنظیم می‌گردد؛ استفاده از \lr{\texttt{nextafter}} در اینجا آگاهانه است و از ردِ طول‌موج‌هایی که در عمل روی مرز تلورانس هستند اما به‌دلیل خطاهای گردکردن \lr{IEEE-754} اندکی فراتر می‌روند، جلوگیری می‌کند. در ادامه، بررسی می‌شود که آیا $\alpha$ در بازهٔ مجاز $[\alpha_{\min}, \alpha_{\max}]$ برای آن طول‌موج قرار دارد یا خیر؛ در صورت عدم تطابق، در حالت سخت‌گیرانه استثنا صادر می‌شود و در غیر این صورت، هشدار \lr{\texttt{PhysicalSuspicionWarning}} نمایش داده می‌شود.


تتابع \lr{\texttt{validate\_parameter(name, value, upper\_bound)}} شامل چندین سازوکار حفاظتی است. ابتدا، خودِ پارامتر \lr{\texttt{upper\_bound}} بررسی می‌شود تا نامنفی بودن آن تضمین شود؛ زیرا مقدار منفی برای این پارامتر می‌تواند همهٔ مقادیر نامنفی را به‌صورت خاموش معتبر تلقی کرده و عملاً سازوکار اعتبارسنجی را غیرفعال کند. سپس، ورودی پیش از ادامهٔ پردازش بررسی می‌شود تا از ایجاد خطاهای مبهم از نوع \lr{\texttt{TypeError}} جلوگیری شود. اگر ورودی دارای ویژگی‌های \lr{\texttt{shape}} و \lr{\texttt{dtype}} باشد، به این معنا که یک آرایه از نوع \lr{\texttt{NumPy/JAX/Torch}} است، تابع استثنا صادر می‌کند.

در مرحلهٔ بعد، تابع \lr{\texttt{is\_finite\_non\_negative(value)}} درون یک بلوک \lr{\texttt{try/except}} فراخوانی می‌شود تا \lr{\texttt{TypeError}} ناشی از ورودی‌های مختلط \lr{\texttt{complex}} به \lr{\texttt{ParameterValidationError}} تبدیل شود. در نهایت، اگر شرط $value > upper\_bound$ برقرار باشد، استثنای \lr{\texttt{ParameterValidationError}} همراه با واحد متناظرِ استخراج‌شده از \lr{\texttt{VALIDATED\_UNITS}} صادر می‌شود.


تابع مرکزی \texttt{validate\_channel()} تمامی بررسی‌های بالا را به‌صورت یکپارچه فراخوانی می‌کند و علاوه بر آن‌ها، انسجام حالت مشتق‌شده را نیز اعتبارسنجی می‌کند. این تابع برای کشف فساد حالت ناشی از \lr{pickling}، \lr{monkey-patching} یا ارتقای ماژول \lr{qkd.datatypes} طراحی شده و در \lr{API} عمومی به‌عنوان نقطه ورودی برای اعتبارسنجی مجدد در نظر گرفته شده است. تابع به‌صورت صریح در \texttt{\_\_post\_init\_\_} فراخوانی نمی‌شود زیرا آن بررسی‌ها از قبل انجام شده‌اند؛ اما کاربر می‌تواند آن را به‌صورت دستی پس از عملیاتی که ممکن است حالت را فاسد کند فراخوانی نماید. مکانیزم اختیاری \texttt{set\_validate\_on\_construction(True)} اجازه می‌دهد این تابع در انتهای هر \texttt{\_\_post\_init\_\_} به‌صورت خودکار فراخوانی شود تا اطمینان اضافی فراهم شود.

\section{ارسال و دریافت با سایر ماژول‌ها}

این ماژول به‌عنوان یک لایه میانی در معماری فریم‌ورک عمل می‌کند و رابط‌های مشخصی با ماژول‌های بالا-دست و پایین-دست دارد. از منظر بالا-دست، مصرف‌کننده اصلی این ماژول کلاس \lr{FiberChannel} در \texttt{qkd.channel} است که در متد \texttt{\_\_post\_init\_\_} خود توابع گیتینگ را به‌ترتیب زیر فرامی‌خواند: نخست \texttt{check\_type()} برای اطمینان از اینکه \texttt{channel.attenuation} یک \lr{AttenuationConfig} واقعی است (نه مثلاً یک دیکشنری خام که بعداً منجر به \lr{AttributeError} می‌شود)، سپس \texttt{validate\_parameter\_ranges()} برای بررسی $L \leq$ MAX\_DISTANCE\_KM و $\alpha \leq$ MAX\_FIBER\_LOSS\_DB\_KM، سپس چهار بررسی معناداری فیزیکی \lr{(\_meaningful\_alpha/distance)}، سپس \texttt{check\_total\_loss\_finite()} و \texttt{check\_attenuation\_config\_contract()}، و در نهایت بررسی‌های طول‌موج و سلامت انتقال. این ترتیب به‌گونه‌ای است که هر گیت بر اساس خروجی گیت‌های قبلی عمل کند.

ماژول‌های دیگر که از این لایه استفاده می‌کنند شامل: \texttt{qkd.simulation\_driver} که در زمان ساخت کانال‌های متعدد برای شبیه‌سازی \lr{Monte Carlo}، توابع گیتینگ را به‌صورت ضمنی از طریق سازنده \lr{FiberChannel} فراخوانی می‌کند؛ \texttt{qkd.state\_manager} که پس از بارگذاری حالت کانال از دیسک، \texttt{validate\_channel()} را فراخوانی می‌کند تا از انسجام حالت بارگذاری‌شده اطمینان حاصل کند؛ و \texttt{qkd.cli} که در زمان پردازش آرگومان‌های خط فرمان، \texttt{validate\_parameter()} را برای اعتبارسنجی اولیه پارامترها پیش از ساخت کانال به کار می‌برد.

از منظر پایین-دست، خروجی اصلی این ماژول مقادیر اعتبارسنجی‌شده انتقال $T$ و $L_{total}$ هستند که توسط کانال فیبر در اختیار ماژول آشکارساز \lr{(qkd.detectors)} قرار می‌گیرند. به‌طور خاص، تابع \texttt{compute\_transmittance()} انتقال نهایی را به‌عنوان یک \lr{float} در بازه $[0, 1]$ بازمی‌گرداند که سپس به‌عنوان آرگومان \texttt{channel\_transmittance} به متد \texttt{simulate\_detection()} آشکارساز ارسال می‌شود. اعتبارسنجی سخت‌گیرانه این ماژول تضمین می‌کند که این مقدار هرگز \lr{NaN}، زیر-نرمال، یا خارج از بازه فیزیکی نباشد، که در نتیجه از خطاهای خاموش در محاسبات بازده \lr{decoy-state} و نرخ کلید جلوگیری می‌کند.

از منظر وابستگی‌های درون‌پروژه‌ای، این ماژول به‌صورت مستقیم به چهار ماژول زیرین وابسته است: \texttt{qkd.exceptions} برای \texttt{ParameterValidationError}؛ \texttt{qkd.constants} برای توابع کمری \texttt{is\_close} و \texttt{is\_finite\_non\_negative}؛ \texttt{qkd.datatypes} برای \texttt{AttenuationConfig}؛ و \texttt{qkd.\_schema} برای ثابت‌ها، \lr{ContextVar}ها، تایپ‌ها و سنتینل‌ها. این وابستگی‌ها عمدتاً به‌صورت \lr{static} (در زمان \lr{import}) حل می‌شوند و این ماژول بر خلاف ماژول‌هایی مانند \texttt{qkd.detectors} لایه \lr{fallback} برای حالت standalone ندارد؛ زیرا فرض بر این است که این ماژول همواره در بستر کامل بسته \lr{qkd} اجرا می‌شود.

لاگر ماژول در سطح \lr{WARNING} برای هشدارهای غیرسخت‌گیرانه و در سطح \lr{DEBUG} برای گزارش جزئیات اعتبارسنجی موفق به کار می‌رود. این دو سطح به کاربران اجازه می‌دهد با تنظیم سطح لاگ، میزان جزئیات گزارش‌ها را کنترل کنند: تنظیم \texttt{logging.getLogger("qkd.\_gating").setLevel(logging.DEBUG)} جریان کامل اعتبارسنجی را برای دیباگ نمایش می‌دهد، در حالی که سطح پیش‌فرض \lr{WARNING} تنها هشدارهای مهم را نشان می‌دهد. این طراحی به‌خصوص برای دیباگ پیکربندی‌های پیچیده که در آن‌ها چندین کانال به‌صورت موازی ساخته می‌شوند (مانند شبیه‌سازی \lr{sweep} روی طول‌موج یا فاصله) مفید است.

در نهایت، باید به این نکته اشاره کرد که این ماژول به‌صورت مستقیم با مکانیزم \lr{ContextVar} پایتون برای کنترل حالت سخت‌گیرانه تعامل دارد. تابع \texttt{get\_strict\_mode()} در هر فراخوانی مقدار فعلی \lr{ContextVar} را می‌خواند، که به این معنا است که تغییر حالت سخت‌گیرانه در طول اجرا (مانند \texttt{with strict\_mode\_context(False):}) به‌صورت پویا توسط تمامی توابع گیتینگ مشاهده می‌شود. این مکانیزم برای تست‌هایی که نیاز به غیرفعال‌کردن موقت حالت سخت‌گیرانه دارند (مانند تست‌های \lr{boundary case} که عمداً پیکربندی غیرفیزیکی می‌سازند) حیاتی است و در عین حال تضمین می‌کند که حالت پیش‌فرض (سخت‌گیرانه) برای کاربردهای انتشار علمی همواره فعال باشد.

	
	% =========================== بخش دوم: لایه کریپتوگرافی ===========================
	\part{اثبات‌های امنیتی رمزنگاری}
	
	\chapter{ساختار بنیادین اثبات امنیت}
	\label{ch:base_proof}
	\sloppy
\emergencystretch=5em
\hbadness=10000
\vbadness=10000
\tolerance=2000
\providecolor{codebg}{RGB}{248,248,248}
\providecolor{codeframe}{RGB}{200,200,200}
\providecolor{keyword}{RGB}{0,0,180}
\providecolor{string}{RGB}{163,21,21}
\providecolor{comment}{RGB}{0,128,0}
\lstset{
	basicstyle=\ttfamily\scriptsize,
	breaklines=true,
	breakatwhitespace=false,
	columns=fullflexible,
	keepspaces=true,
	frame=single,
	framesep=2pt,
	showstringspaces=false,
	backgroundcolor=\color{codebg},
	rulecolor=\color{codeframe},
	keywordstyle=\color{keyword}\bfseries,
	stringstyle=\color{string},
	commentstyle=\color{comment}\itshape,
	numberstyle=\tiny\color{gray},
	numbers=left,
	numbersep=6pt,
	xleftmargin=14pt,
	framexleftmargin=14pt,
	literate=
	{ą}{{\v{a}}}1 {ę}{{\k{e}}}1 {ż}{{\.z}}1 {ś}{{\'s}}1
	{ł}{{\l}}1 {ć}{{\'c}}1 {ń}{{\'n}}1 {ó}{{\'o}}1 {ź}{{\'z}}1
	{→}{{$\rightarrow$}}1 {←}{{$\leftarrow$}}1 {≠}{{$\neq$}}1
}

\section{نمای‌ کلی و هدف ماژول}

فایل \lr{\texttt{qkd/proofs/base.py}} ستون فقرات لایه اثبات‌های امنیتی در فریم‌ورک شبیه‌سازی توزیع کلید کوانتومی است. این ماژول یک \lr{Abstract Base Class} به نام \lr{\texttt{FiniteKeyProof}} معرفی می‌کند که قرارداد رابطی همه اثبات‌های امنیتی کلید متناهی را در سراسر سیستم استاندارد می‌سازد. اثبات‌های امنیتی کلید متناهی به دسته‌ای از الگوریتم‌ها گفته می‌شود که به جای اتکای صرف به حد asymptotic توزیع کلید، تأمین امنیت برای طول کلید متناهی را با پارامتر امنیتی $\varepsilon_{\text{sec}}$ قابل اثبات تضمین می‌کنند. در دنیای واقعی، تعداد پالس‌های ارسالی همواره محدود است و الگوهای asymptotic مانند \lr{Gottesman-Lo-Lütkenhaus-Preskill} سنتی به تنهایی پاسخگو نیستند؛ از این رو لایه‌بندی صریحی از تقسیم $\varepsilon$-بودجه، تخمین پارامتر، محاسبه بازده تک‌فوتونی و کسر فاز خطا ضروری می‌شود.

هدف اصلی این ماژول، فراهم آوردن یک رابط یکپارچه برای همه زیرکلاس‌های اثبات است؛ به گونه‌ای که هر زیرکلاس فقط به پیاده‌سازی چهار متد انتزاعی \lr{\texttt{notation\_map}}، \lr{\texttt{allocate\_epsilons}}، \lr{\texttt{get\_epsilon\_policy}} و \lr{\texttt{estimate\_yields\_and\_errors}} بپردازد و مابقی خدمات مشترک شامل محاسبه کران‌های آماری، آنتروپی باینری، تقسیم ایمن، تشخیص خطاهای عددی و گزارش‌گیری \lr{audit} را از کلاس پایه به ارث ببرد. این طراحی از الگوی \lr{Template Method} پیروی می‌کند و اطمینان می‌دهد که هر پیاده‌سازی اثبات، بدون توجه به اینکه مبتنی بر چارچوب \lr{Ma 2005} باشد یا \lr{Lim 2014} یا \lr{GLLP}، خروجی‌ها و رفتار سازگار و قابل مقایسه تولید می‌کند. به علاوه، کلاس پایه لایه‌ای از دفاع عددی شامل \lr{clamp} احتمالات، \lr{safe divide}، \lr{safe log} و تشخیص مقادیر نامتناهی را فراهم می‌سازد که در محاسبات آنتروپی و \lr{log-likelihood} اهمیت حیاتی دارد.

این فایل علاوه بر کلاس پایه، چندین \lr{dataclass} مورد نیاز برای انتقال داده میان لایه‌ها را تعریف می‌کند. این \lr{dataclass}ها شامل \lr{\texttt{TallyCounts}} برای ثبت شمارش‌های آماری، \lr{\texttt{EpsilonAllocation}} برای ذخیره تقسیم بودجه $\varepsilon$، \lr{\texttt{SolverDiagnostics}} برای گزارش وضعیت \lr{solver} برنامه‌ریزی خطی، \lr{\texttt{DecoyEstimates}} برای نگهداری کران‌های پایین بازده تک‌فوتونی و کران بالای خطای تک‌فوتونی و نهایتاً \lr{\texttt{KeyCalculationResult}} برای ارائه نتیجه نهایی محاسبه طول کلید امن است. همه این ساختارها با \lr{\texttt{@dataclass(frozen=True)}} تعریف شده‌اند تا از تغییر ناخواسته پس از ساخت جلوگیری شود و قابلیت \lr{hashing} و استفاده به عنوان کلید در \lr{cache} را داشته باشند. این تصمیم طراحی در اثبات‌های امنیتی بحرانی است، زیرا هرگونه جهش ناگهانی در مقادیر میانی می‌تواند به نتیجه نهایی رنگ امنیتی بزند.

وظیفه نهایی این ماژول، تضمین یکپارچگی و تکرارپذیری نتایج امنیتی است. به همین دلیل، متغیرهای محیطی مانند \lr{run\_id}، \lr{lib\_version}، نسخه \lr{numpy} و \lr{scipy} در گزارش‌های \lr{audit} ثبت می‌شوند تا هر اجرا به طور کامل قابل بازتولید باشد. این ویژگی برای گزارش‌های امنیتی حساس که نیاز به بازبینی توسط شخص ثالث دارند، الزامی است.

\section{مبانی تئوری و فیزیکی}

\subsection{امنیت قابل ترکیب و تقسیم بودجه \texorpdfstring{$\varepsilon$}{epsilon}}

تعریف مدرن امنیت در QKD بر پایه مفهوم \lr{composable security} استوار است که نخستین بار توسط \lr{Ben-Or} و همکاران و سپس توسط \lr{Canetti} بسط داده شد. در این چارچوب، امنیت یک پروتکل با یک پارامتر $\varepsilon_{\text{sec}} \in (0,1)$ کمیت‌سازی می‌شود که نشان می‌دهد فاصله آماری میان خروجی پروتکل و یک حالت ایده‌آل (توزیع کلید یکنواخت و مستقل از اطلاعات دشمن) چقدر است. اگر $\rho_{\text{real}}$ حالت نهایی مشترک آلیس، باب و ایو و $\rho_{\text{ideal}}$ حالت ایده‌آل باشد، شرط امنیت به صورت زیر نوشته می‌شود:
\begin{equation}
	\frac{1}{2}\,\|\rho_{\text{real}} - \rho_{\text{ideal}}\|_1 \;\le\; \varepsilon_{\text{sec}},
\end{equation}
که در آن $\|\cdot\|_1$ نرم رد ماتریس چگالی است. این تعریف به ما اجازه می‌دهد پروتکل QKD را مانند یک جعبه سیاه با سایر پروتکل‌های رمزنگاری ترکیب کنیم و جمع کل $\varepsilon$ها از قانون زیر پیروی کند:
\begin{equation}
	\varepsilon_{\text{total}} \;\le\; \sum_i \varepsilon_i.
\end{equation}
در فایل \lr{\texttt{base.py}}، این بودجه به چهار مؤلفه شکسته می‌شود:
\begin{equation}
	\varepsilon_{\text{sec}} \;=\; \varepsilon_{\text{PE}} + \varepsilon_{\text{EC}} + \varepsilon_{\text{PA}} + \varepsilon_{\text{sm}},
\end{equation}
که در آن $\varepsilon_{\text{PE}}$ احتمال شکست تخمین پارامتر، $\varepsilon_{\text{EC}}$ احتمال شکست تصحیح خطا، $\varepsilon_{\text{PA}}$ احتمال شکست تقویت حریم خصوصی و $\varepsilon_{\text{sm}}$ پارامتر هموارسازی در \lr{smooth min-entropy} است. این تقسیم در \lr{\texttt{EpsilonAllocation}} ذخیره شده و در متد \lr{\texttt{validate}} با یک \lr{policy function} بررسی می‌شود که مجموع $\varepsilon$ها از کل مجاز تجاوز نکند.

\subsection{تخمین پارامتر و کران‌های آماری}

در پروتکل‌های \lr{Decoy-State}، آلیس به طور تصادفی شدت پالس را میان چندین حالت (سیگنال، \lr{decoy} کم‌شدت و خلأ) تغییر می‌دهد و باب تعداد آشکارسازی‌ها و خطاها را در هر حالت می‌شمارد. هدف، تخمین کران پایین بازده تک‌فوتونی $Y_1^L$ و کران بالای خطای تک‌فوتونی $e_1^U$ است. به دلیل ماهیت تصادفی این شمارش‌ها، تخمین قطعی ممکن نیست و باید از کران‌های آماری استفاده کرد. این ماژول سه روش را پشتیبانی می‌کند: \lr{Clopper-Pearson}، \lr{Hoeffding} و \lr{Gaussian}.

\lr{Clopper-Pearson} دقیق‌ترین روش و بر پایه توزیع \lr{Beta} است. اگر $k$ آشکارسازی از $n$ پالس مشاهده شود، کران پایین و بالا برای احتمال واقعی $p$ به این شکل به دست می‌آیند:
\begin{align}
	p_{\text{lower}} &= B\!\left(\frac{\alpha}{2};\, k,\, n-k+1\right), \\
	p_{\text{upper}} &= B\!\left(1-\frac{\alpha}{2};\, k+1,\, n-k\right),
\end{align}
که در آن $B(q; a, b)$ تابع \lr{quantile} توزیع \lr{Beta} و $\alpha$ احتمال شکست است. این کران‌ها محافظه‌کارانه‌ترین کران ممکن هستند و در عین حال نیازمند محاسبه \lr{\texttt{beta.ppf}} از \lr{SciPy} با نسخه حداقل 1.8 می‌باشند؛ به همین دلیل فایل در ابتدا بررسی نسخه انجام می‌دهد. \lr{Hoeffding} تقریبی‌تر اما ساده‌تر است:
\begin{equation}
	p \;\in\; \left[\hat{p} - \sqrt{\frac{\ln(1/\alpha)}{2n}},\; \hat{p} + \sqrt{\frac{\ln(1/\alpha)}{2n}}\right],
\end{equation}
که در آن $\hat{p} = k/n$. در نهایت \lr{Gaussian} از تقریب نرمال استفاده می‌کند:
\begin{equation}
	p \;\in\; \left[\hat{p} - z_{1-\alpha/2}\sqrt{\frac{\hat{p}(1-\hat{p})}{n}},\; \hat{p} + z_{1-\alpha/2}\sqrt{\frac{\hat{p}(1-\hat{p})}{n}}\right].
\end{equation}
این سه روش در متد \lr{\texttt{get\_bounds}} با توجه به \lr{\texttt{ci\_method}} در \lr{\texttt{QKDParams}} انتخاب می‌شوند.

\subsection{آنتروپی باینری و تقویت حریم خصوصی}

آنتروپی باینری \lr{Shannon} یک عنصر محوری در محاسبه طول کلید امن است:
\begin{equation}
	h(p) \;=\; -p\,\log_2(p) - (1-p)\,\log_2(1-p).
\end{equation}
در اثبات‌های مبتنی بر \lr{GLLP} و \lr{Ma 2005}، طول کلید امن به صورت زیر به دست می‌آید:
\begin{equation}
	\ell \;=\; s_{z,1}^L\left[1 - h\!\left(e_{\text{ph}}^U\right)\right] - \lambda_{\text{EC}} - 6\,\log_2\!\left(\frac{19}{\varepsilon_{\text{sec}}}\right) - \log_2\!\left(\frac{2}{\varepsilon_{\text{cor}}}\right),
\end{equation}
که در آن $s_{z,1}^{L}$ کران پایین تعداد آشکارسازی‌های تک‌فوتونی در پایهٔ $Z$، $e_{\mathrm{ph}}^{U}$ کران بالای خطای فاز، $\lambda_{\mathrm{EC}}$ میزان نشت اطلاعات ناشی از تصحیح خطا و $\varepsilon_{\mathrm{cor}}$ پارامتر صحت است. محاسبهٔ $h(p)$ در نزدیکی کران‌های $p = 0$ و $p = 1$، به‌دلیل حضور عبارت‌های $p\log p$ و $(1-p)\log(1-p)$، از نظر عددی ناپایدار می‌شود؛ هرچند طبق قرارداد ریاضی، $0\log 0 = 0$ در نظر گرفته می‌شود. ازاین‌رو، متد \lr{\texttt{binary\_entropy}} با استفاده از \lr{\texttt{np.clip}} مقدار $p$ را به بازهٔ $[\delta, 1-\delta]$ محدود می‌کند، که در آن
\[
\delta = \texttt{ENTROPY\_PROB\_CLAMP}.
\]
این محدودسازی، همراه با استفاده از \lr{\texttt{np.nan\_to\_num}}، تضمین می‌کند که خروجی همواره مقداری متناهی و قابل‌استفاده باشد.


\subsection{برنامه‌ریزی خطی در تخمین \lr{Decoy}}

تخمین $Y_1^L$ معمولاً به حل یک مسأله برنامه‌ریزی خطی (\lr{LP}) نیاز دارد. در چارچوب \lr{Ma 2005}، این مسأله با حداقل‌سازی $Y_1$ تحت قیود فیزیکی که از معادلات تعادل بازده به دست می‌آیند، صورت می‌گیرد. اگر \lr{solver} نتواند جواب بهینه پیدا کند یا باقی‌مانده \lr{LP} بزرگ باشد، \lr{\texttt{SolverDiagnostics}} این وضعیت را ثبت می‌کند و متد \lr{\texttt{\_conservative\_decoy\_estimate}} یک تخمین محافظه‌کارانه با $Y_1^L = 0$ و $e_1^U = 0.5$ برمی‌گرداند. این مقدار به معنای آن است که در بدترین حالت، هیچ تک‌فوتونی امن قابل استخراج نیست و طول کلید صفر می‌شود.

\section{تحلیل معماری و ساختار داده‌ها}

\subsection{واردسازی وابستگی‌ها}

فایل با یک مجموعه \lr{try/except} ساختارمند برای واردسازی وابستگی‌های داخلی و خارجی آغاز می‌شود. این الگو دو هدف را دنبال می‌کند: نخست، در صورت نبود یک وابستگی، پیام خطای روشنی با \lr{\texttt{ConfigurationError}} تولید شود و دوم، بررسی نسخه \lr{SciPy} حداقل 1.8 تضمین گردد. این بررسی به دلیل رفتار متفاوت \lr{\texttt{beta.ppf}} در نسخه‌های قدیمی ضروری است. کد مربوطه به این شکل است:

\begin{LTR}
	\begin{lstlisting}[language=Python]
		try:
		import scipy
		from scipy import special
		from scipy.stats import beta
		if tuple(map(int, scipy.__version__.split('.')[:2])) < (1, 8):
		raise ConfigurationError(
		"SciPy >= 1.8 is required for reliable beta.ppf behavior.",
		code=QKDErrorCode.CONFIG,
		context={"installed_version": scipy.__version__,
			"required_version": ">=1.8"}
		)
		except ImportError as exc:
		raise ConfigurationError(
		"SciPy is required for Clopper-Pearson confidence intervals.",
		code=QKDErrorCode.CONFIG,
		cause=exc
		) from exc
	\end{lstlisting}
\end{LTR}

در خط اول، \lr{\texttt{scipy}} به عنوان ماژول اصلی وارد می‌شود. در خط دوم، زیرماژول \lr{\texttt{special}} برای توابع خاص ریاضی مانند \lr{\texttt{betainc}} و در خط سوم توزیع \lr{beta} برای محاسبه \lr{quantile} در کران‌های \lr{Clopper-Pearson} وارد می‌شوند. خط چهارم نسخه را تجزیه و مقایسه می‌کند و در صورت عدم تطابق، استثنای سفارشی با \lr{\texttt{QKDErrorCode.CONFIG}} پرتاب می‌کند. این کد طراحی شده تا در محیط‌های عملیاتی که نسخه‌های قدیمی \lr{SciPy} ممکن است نصب باشند، پیش از وقوع خطای عددی غیرقابل ردیابی، متوقف شود.

علاوه بر این، یک واردسازی امن برای \lr{\texttt{mpmath}} وجود دارد که در صورت موجود بودن، امکان محاسبات با دقت دلخواه را فراهم می‌کند:

\begin{LTR}
	\begin{lstlisting}[language=Python]
		try:
		import mpmath
		except ImportError:
		mpmath = None
	\end{lstlisting}
\end{LTR}

این الگو از نوع \lr{optional dependency} است؛ یعنی اگر \lr{\texttt{mpmath}} نصب نباشد، برنامه متوقف نمی‌شود بلکه متغیر \lr{\texttt{mpmath}} روی \lr{\texttt{None}} تنظیم می‌شود و کد می‌تواند با استفاده از شرط \lr{\texttt{use\_mpmath}} تصمیم بگیرد که از مسیر دقت بالا استفاده کند یا خیر.

\subsection{ثابت‌های ماژول و \texorpdfstring{\lr{TypeAlias}}{TypeAlias}}

دو ثابت نمادین در سطح ماژول برای کلیدهای آماری تعریف می‌شوند:

\begin{LTR}
	\begin{lstlisting}[language=Python]
		TALLY_KEY_Z_SIGNAL = "Z_signal"
		TALLY_KEY_X_SIGNAL = "X_signal"
	\end{lstlisting}
\end{LTR}

این دو ثابت نشان‌دهنده پایه‌های اندازه‌گیری در پروتکل \lr{BB84} هستند. پایه $Z$ پایه محاسباتی استاندارد و پایه $X$ پایه \lr{Hadamard} است. در \lr{Decoy-State}، آلیس و باب پس از تبادل، آمار را بر اساس پایه تفکیک می‌کنند تا بتوانند خطای فاز $e_{\text{ph}}$ را به صورت جداگانه از خطای بیت $e_{\text{bit}}$ تخمین بزنند. استفاده از ثابت به جای رشته خام، خطاهای تایپی را حذف می‌کند و در بازبینی کد به وضوح مشخص می‌کند که چه کلیدی مورد نظر است.

همچنین \lr{\texttt{T = TypeVar("T")}} برای استفاده در توابع عمومی تعریف شده و \lr{\texttt{\_\_version\_\_ = "4.2.0"}} نسخه فعلی ماژول را مشخص می‌کند. این نسخه در گزارش‌های \lr{audit} درج می‌شود تا در صورت بروز مشکل، نسخه دقیق کد قابل ردیابی باشد.

\subsection{شمارش‌ها و \lr{Enum}های خطا}

\lr{\texttt{ProofMode}} یک \lr{Enum} سه‌حالتی است که رفتار سیستم در مواجهه با خطاهای عددی و هشدارها را تعیین می‌کند:

\begin{LTR}
	\begin{lstlisting}[language=Python]
		class ProofMode(Enum):
		PRODUCTION = auto()
		DEBUG = auto()
		AUDIT = auto()
	\end{lstlisting}
\end{LTR}

حالت \lr{PRODUCTION} حالت پیش‌فرض است که در آن خطاهای غیربحرانی با هشدار \lr{log} می‌شوند و اجرا ادامه می‌یابد. حالت \lr{DEBUG} برای محیط توسعه است و اطلاعات بیشتری در \lr{log} ثبت می‌کند. حالت \lr{AUDIT} سخت‌گیرانه‌ترین حالت است که در آن هر گونه انحراف از شرط‌های فیزیکی یا عددی به استثنای کامل منجر می‌شود. این حالت برای کاربردهای حساس مانند \lr{certification} یا بازبینی توسط شخص ثالث طراحی شده است. متدهای \lr{\texttt{is\_audit\_mode}} و \lr{\texttt{is\_debug\_mode}} در کلاس پایه به این \lr{Enum} متکی‌اند تا تصمیم‌گیری در مورد رفتار خطا را ساده کنند.

\lr{\texttt{ErrorCode}} شش کد خطای متمایز را تعریف می‌کند که در گزارش نتیجه محاسبه کلید استفاده می‌شوند تا علت دقیق تولید نشدن کلید قابل ردیابی باشد:

\begin{LTR}
	\begin{lstlisting}[language=Python]
		class ErrorCode(Enum):
		INSUFFICIENT_STATISTICS = auto()
		LP_INFEASIBLE = auto()
		NUMERICAL_CLAMP_APPLIED = auto()
		LEAKAGE_EXCEEDS_ENTROPY = auto()
		HIGH_SOLVER_RESIDUAL = auto()
		SAFE_DIVIDE_FALLBACK = auto()
	\end{lstlisting}
\end{LTR}

هر یک از این کدها به یک شرایط مشخص اشاره دارد. \lr{INSUFFICIENT\_STATISTICS} زمانی فعال می‌شود که $Y_1^L$ از آستانه \lr{\texttt{min\_y1\_threshold}} کمتر باشد. \lr{LP\_INFEASIBLE} نشان‌دهنده ناپذیرفتگی مسأله برنامه‌ریزی خطی است. \lr{NUMERICAL\_CLAMP\_APPLIED} هشدار می‌دهد که یک مقدار منفی به صفر \lr{clamp} شده است. \lr{LEAKAGE\_EXCEEDS\_ENTROPY} بحرانی‌ترین کد است و نشان می‌دهد که نشت تصحیح خطا و تقویت حریم خصوصی از آنتروپی قابل استخراج بیشتر شده است. \lr{HIGH\_SOLVER\_RESIDUAL} باقی‌مانده بالای \lr{solver} را علامت می‌زند و \lr{SAFE\_DIVIDE\_FALLBACK} نشان می‌دهد که تقسیم ایمن به مقدار پیش‌فرض بازگشته است.

\subsection{\lr{Dataclass}های داده}

\subsubsection{\lr{TallyCounts}}

این \lr{dataclass} ساده اما حیاتی است و دو فیلد دارد: \lr{detections} و \lr{errors}. هر دو از نوع \lr{\texttt{int}} و با مقدار پیش‌فرض صفر تعریف شده‌اند. با استفاده از \lr{\texttt{@dataclass(frozen=True)}} غیرقابل تغییر شده و در متد \lr{\texttt{\_\_post\_init\_\_}} دو شرط اعتبارسنجی دارد:

\begin{LTR}
	\begin{lstlisting}[language=Python]
		@dataclass(frozen=True)
		class TallyCounts:
		detections: int = 0
		errors: int = 0
		
		def __post_init__(self):
		if self.detections < 0 or self.errors < 0:
		raise QKDParamError(
		"Tally counts cannot be negative.",
		code=QKDErrorCode.PARAM_VALIDATION,
		context={"detections": self.detections, "errors": self.errors}
		)
		if self.errors > self.detections:
		raise QKDParamError(
		f"Error count ({self.errors}) cannot exceed detection count ({self.detections}).",
		code=QKDErrorCode.PARAM_VALIDATION,
		context={"detections": self.detections, "errors": self.errors}
		)
	\end{lstlisting}
\end{LTR}

شرط نخست عدم منفی بودن شمارش‌ها را تضمین می‌کند که یک ضرورت فیزیکی است. شرط دوم منطقی‌تر است: تعداد خطاها نمی‌تواند از تعداد آشکارسازی‌ها بیشتر باشد، زیرا هر خطا نیازمند یک آشکارسازی است. استفاده از \lr{\texttt{frozen=True}} اطمینان می‌دهد که پس از ساخت، شمارش‌ها قابل تغییر نیستند و در نتیجه در سراسر محاسبات ثابت می‌مانند. این ویژگی برای صحت‌سنجی نتایج و جلوگیری از تغییرات ناخواسته در طول محاسبات بسیار مهم است.

\subsubsection{\lr{EpsilonAllocation}}

این ساختار تقسیم بودجه $\varepsilon$ را در چهار مؤلفه نگه می‌دارد:

\begin{LTR}
	\begin{lstlisting}[language=Python]
		@dataclass
		class EpsilonAllocation:
		eps_sec: float
		eps_pe: float
		eps_smooth: float
		eps_pa: float
		
		def validate(self, policy: Callable[["EpsilonAllocation"], float]) -> None:
		consumed = policy(self)
		if consumed > self.eps_sec:
		raise QKDParamError(
		f"Epsilon allocation exceeds budget: {consumed:.2e} > {self.eps_sec:.2e}",
		code=QKDErrorCode.PARAM_VALIDATION,
		context={"consumed": consumed, "budget": self.eps_sec,
			"policy": policy.__name__}
		)
	\end{lstlisting}
\end{LTR}

برخلاف \lr{TallyCounts}، این کلاس \lr{frozen} نیست زیرا ممکن است زیرکلاس‌ها در طول زمان تخصیص را تنظیم کنند. متد \lr{\texttt{validate}} یک \lr{policy function} می‌گیرد که جمع موزون $\varepsilon$ها را محاسبه می‌کند و اگر از کل بودجه $\varepsilon_{\text{sec}}$ تجاوز کند، استثنا پرتاب می‌کند. این الگو به زیرکلاس‌ها اجازه می‌دهد سیاست‌های متفاوتی برای تخصیص $\varepsilon$ پیاده‌سازی کنند (مثلاً تخصیص مساوی، تخصیص بر اساس اهمیت نسبی یا تخصیص بهینه) بدون آنکه کلاس پایه وابستگی به سیاست خاصی داشته باشد.

\subsubsection{\lr{SolverDiagnostics}}

این \lr{dataclass} برای ثبت اطلاعات مربوط به حل مسأله \lr{LP} در تخمین \lr{Decoy} به کار می‌رود:

\begin{LTR}
	\begin{lstlisting}[language=Python]
		@dataclass(frozen=True)
		class SolverDiagnostics:
		solver_name: str
		is_success: bool
		status_message: str
		residual_norm: Optional[float] = None
		timings: Dict[str, float] = field(default_factory=dict)
		numeric_diagnostics: Dict[str, Any] = field(default_factory=dict)
	\end{lstlisting}
\end{LTR}

فیلد \lr{\texttt{solver\_name}} نام \lr{solver} (مانند \lr{scipy.optimize.linprog} یا \lr{cvxopt}) را ذخیره می‌کند. \lr{\texttt{is\_success}} نشان می‌دهد که آیا \lr{solver} به راه‌حل بهینه دست یافته است یا خیر. \lr{\texttt{status\_message}} پیام متنی \lr{solver} را نگه می‌دارد. \lr{\texttt{residual\_norm}} نرم باقی‌مانده مسأله را نشان می‌دهد و در تشخیص کیفیت راه‌حل کاربرد دارد. دو فیلد \lr{\texttt{timings}} و \lr{\texttt{numeric\_diagnostics}} از نوع \lr{\texttt{Dict}} با \lr{\texttt{default\_factory=dict}} هستند تا در صورت عدم ارائه مقدار، دیکشنری خالی ایجاد شود. استفاده از \lr{\texttt{default\_factory}} به جای مقدار پیش‌فرض مستقیم، از مشکل معروف \lr{mutable default argument} در پایتون جلوگیری می‌کند.

\subsubsection{\lr{DecoyEstimates}}

این ساختار خروجی اصلی مرحله تخمین \lr{Decoy} است:

\begin{LTR}
	\begin{lstlisting}[language=Python]
		@dataclass(frozen=True)
		class DecoyEstimates:
		yield_1_lower_bound: float
		error_rate_1_upper_bound: float
		is_feasible: bool
		failure_prob_used: float
		diagnostics: SolverDiagnostics
		
		def __repr__(self) -> str:
		return (f"DecoyEstimates(Y1_L={self.yield_1_lower_bound:.2e}, "
		f"e1_U={self.error_rate_1_upper_bound:.3f}, feasible={self.is_feasible})")
		
		def as_serializable(self) -> Dict[str, Any]:
		data = asdict(self, dict_factory=_json_safe_dict)
		data['metadata'] = {
			'schema_version': '1.2', 'lib_version': __version__,
			'numpy_version': np.__version__, 'scipy_version': scipy.__version__,
		}
		return data
	\end{lstlisting}
\end{LTR}

دو فیلد اصلی \lr{\texttt{yield\_1\_lower\_bound}} (یعنی $Y_1^L$) و \lr{\texttt{error\_rate\_1\_upper\_bound}} (یعنی $e_1^U$) هستند که مستقیماً در فرمول طول کلید وارد می‌شوند. \lr{\texttt{is\_feasible}} نشان می‌دهد که آیا \lr{LP} یک راه‌حل قابل قبول پیدا کرده است یا خیر. \lr{\texttt{failure\_prob\_used}} احتمال شکستی است که در محاسبه کران‌ها استفاده شده (معمولاً $\varepsilon_{\text{PE}}$). متد \lr{\texttt{as\_serializable}} با استفاده از \lr{\texttt{\_json\_safe\_dict}} نوع‌های \lr{numpy} و \lr{Enum} را به نوع‌های بومی پایتون تبدیل می‌کند و متادیتای نسخه را اضافه می‌کند تا خروجی به راحتی به \lr{JSON} سریال‌سازی شود.

\subsubsection{\lr{KeyCalculationResult}}

این \lr{dataclass} مهم‌ترین خروجی کل اثبات امنیتی است:

\begin{LTR}
	\begin{lstlisting}[language=Python]
		@dataclass(frozen=True)
		class KeyCalculationResult:
		secure_key_length: int
		privacy_amplification_term: float
		error_correction_leakage: float
		phase_error_rate_upper_bound: float
		diagnostics: Any = field(default_factory=dict, repr=False)
		error_codes: list[ErrorCode] = field(default_factory=list)
		weak_gllp_rate: Optional[float] = None
		tagged_fraction_bound: Optional[float] = None
	\end{lstlisting}
\end{LTR}

چهار فیلد اصلی \lr{\texttt{secure\_key\_length}} ($\ell$)، \lr{\texttt{privacy\_amplification\_term}} (جریمه تقویت حریم خصوصی)، \lr{\texttt{error\_correction\_leakage}} ($\lambda_{\text{EC}}$) و \lr{\texttt{phase\_error\_rate\_upper\_bound}} ($e_{\text{ph}}^U$) هستند. \lr{\texttt{diagnostics}} از نوع \lr{\texttt{Any}} تعریف شده تا هم لیست رشته و هم دیکشنری ساختاریافته را بپذیرد. دو فیلد اختیاری \lr{\texttt{weak\_gllp\_rate}} و \lr{\texttt{tagged\_fraction\_bound}} برای مقایسه با چارچوب \lr{GLLP} سنتی طراحی شده‌اند. \lr{\texttt{tagged\_fraction\_bound}} کسر فوتون‌های برچسب‌خورده (تک‌فوتون‌هایی که اطلاعات فاز آن‌ها نشت کرده) را نشان می‌دهد و در فرمول \lr{GLLP} نرخ کلید را به صورت زیر محدود می‌کند:
\begin{equation}
	R_{\text{GLLP}} \;\ge\; q\,\left[Q_1\,(1 - h(e_1^U)) - Q_{\mu}\,f\,h(E_{\mu})\right],
\end{equation}
که در آن $Q_1$ بازده تک‌فوتونی، $Q_{\mu}$ بازده سیگنال، $E_{\mu}$ خطای بیت سیگنال و $f$ کارایی تصحیح خطا است.

\subsection{کلاس پایه انتزاعی \lr{FiniteKeyProof}}

این کلاس با \lr{\texttt{ABC}} به عنوان کلاس پایه انتزاعی تعریف شده و سه ویژگی کلاسی دارد:

\begin{LTR}
	\begin{lstlisting}[language=Python]
		class FiniteKeyProof(ABC):
		__implementation_version__: str = "0.0.0"
		allowed_ci_methods: frozenset[ConfidenceBoundMethod] = frozenset(ConfidenceBoundMethod)
		min_y1_threshold: float = 1e-9
	\end{lstlisting}
\end{LTR}

\lr{\texttt{\_\_implementation\_version\_\_}} به زیرکلاس‌ها اجازه می‌دهد نسخه پیاده‌سازی خود را ثبت کنند. \lr{\texttt{allowed\_ci\_methods}} یک \lr{frozenset} از روش‌های مجاز کران آماری است که به صورت پیش‌فرض همه روش‌های \lr{\texttt{ConfidenceBoundMethod}} را شامل می‌شود؛ زیرکلاس‌ها می‌توانند آن را محدودتر کنند. \lr{\texttt{min\_y1\_threshold}} آستانه $10^{-9}$ برای $Y_1^L$ است که اگر کمتر از آن باشد، تخمین به عنوان غیر قابل اتکا در نظر گرفته می‌شود. این آستانه بر اساس تجربه عملی تنظیم شده تا از تولید کلید نامعتبر در شرایط آمار ضعیف جلوگیری کند.

متد سازنده \lr{\texttt{\_\_init\_\_}} مجموعه‌ای از تنظیمات اولیه را انجام می‌دهد:

\begin{LTR}
	\begin{lstlisting}[language=Python]
		def __init__(self, params: QKDParams, mode: ProofMode = ProofMode.PRODUCTION):
		self.p = params
		self.mode = mode
		self.logger = logging.getLogger(self.__class__.__name__)
		self.run_id = getattr(params, 'run_id', 'no-run-id')
		self.min_prob_clamp = getattr(params, 'entropy_clamp', ENTROPY_PROB_CLAMP)
		self.use_mpmath = getattr(params, 'use_mpmath', False) and mpmath is not None
		
		self.logger.info("event=init run_id=%s proof=%s version=%s numpy=%s scipy=%s",
		self.run_id, self.__class__.__name__, self.__implementation_version__,
		np.__version__, scipy.__version__)
		
		self.validate_physical_params()
		self._check_notation_map()
		
		if self.p.ci_method not in self.allowed_ci_methods:
		msg = f"CI method '{self.p.ci_method.name}' is not recommended for this proof."
		if self.is_audit_mode():
		raise QKDParamError(
		msg,
		code=QKDErrorCode.PARAM_VALIDATION,
		context={"ci_method": self.p.ci_method.name,
			"allowed": [m.name for m in self.allowed_ci_methods]}
		)
		self.logger.warning("[%s] %s", self.run_id, msg)
		
		self.eps_alloc = self.allocate_epsilons()
		self.eps_alloc.validate()
	\end{lstlisting}
\end{LTR}

در خطوط نخست، پارامترهای QKD و حالت اجرا ذخیره می‌شوند. سپس یک \lr{logger} اختصاصی با نام کلاس ایجاد می‌شود تا لاگ‌ها به راحتی فیلتر شوند. \lr{\texttt{run\_id}} شناسه یکتای اجرا است که در صورت عدم تنظیم، مقدار \lr{\texttt{'no-run-id'}} می‌گیرد. \lr{\texttt{min\_prob\_clamp}} حد \lr{clamp} احتمالات در محاسبه آنتروپی است که از پارامترها خوانده می‌شود و در صورت نبود، از ثابت \lr{\texttt{ENTROPY\_PROB\_CLAMP}} استفاده می‌کند. \lr{\texttt{use\_mpmath}} پرچمی است که در صورت تنظیم در پارامترها و موجود بودن \lr{mpmath}، فعال می‌شود. سپس متدهای \lr{\texttt{validate\_physical\_params}} و \lr{\texttt{\_check\_notation\_map}} فراخوانی می‌شوند تا صحت پارامترها و پیاده‌سازی زیرکلاس بررسی شود. در ادامه، اگر روش \lr{CI} مورد استفاده در فهرست مجاز نباشد، در حالت \lr{AUDIT} استثنا پرتاب می‌شود و در غیر این صورت هشدار داده می‌شود. نهایتاً، \lr{\texttt{allocate\_epsilons}} فراخوانی شده و اعتبارسنجی می‌شود.

\section{پیاده‌سازی ریاضی و الگوریتمی}

\subsection{تبدیل نوع‌ها برای \lr{JSON}}

تابع کمکی \lr{\texttt{\_json\_safe\_dict}} یک \lr{dict\_factory} برای \lr{\texttt{asdict}} است که نوع‌های \lr{numpy} و \lr{Enum} را به نوع‌های بومی پایتون تبدیل می‌کند:

\begin{LTR}
	\begin{lstlisting}[language=Python]
		def _json_safe_dict(data: Iterable[Tuple[str, Any]]) -> Dict[str, Any]:
		"""Recursively converts numpy/enum types to native Python types for JSON."""
		d: Dict[str, Any] = {}
		for key, value in data:
		if isinstance(value, (np.integer, np.int64)):
		d[key] = int(value)
		elif isinstance(value, (np.floating, np.float64)):
		d[key] = float(value)
		elif isinstance(value, np.ndarray):
		if value.size > 1000:
		d[key] = {'data': value[:1000].tolist(), 'truncated': True}
		else:
		d[key] = value.tolist()
		elif isinstance(value, np.bool_):
		d[key] = bool(value)
		elif isinstance(value, Enum):
		d[key] = value.name
		elif isinstance(value, dict):
		d[key] = _json_safe_dict(value.items())
		else:
		d[key] = value
		return d
	\end{lstlisting}
\end{LTR}

این تابع برای هر جفت کلید-مقدار، نوع مقدار را بررسی می‌کند. اگر \lr{np.integer} یا \lr{np.int64} باشد به \lr{\texttt{int}} تبدیل می‌شود. اگر \lr{np.floating} یا \lr{np.float64} باشد به \lr{\texttt{float}} تبدیل می‌شود. برای آرایه‌های \lr{numpy}، اگر اندازه از 1000 تجاوز کند، تنها 1000 عنصر اول به همراه پرچم \lr{\texttt{truncated}} ذخیره می‌شوند تا حجم گزارش \lr{audit} کنترل‌پذیر بماند. این تصمیم طراحی بر اساس تجربه عملی اتخاذ شده، زیرا در شبیه‌سازی‌های بزرگ آرایه‌های ده‌ها هزار عنغانی می‌توانند گزارش را غیرقابل مدیریت کنند. \lr{np.bool\_} به \lr{\texttt{bool}} و \lr{Enum} به نام آن تبدیل می‌شود. اگر مقدار خود دیکشنری باشد، فراخوانی بازگشتی انجام می‌شود. این تابع در متدهای \lr{\texttt{as\_serializable}} \lr{DecoyEstimates} و \lr{KeyCalculationResult} استفاده می‌شود.

\subsection{کران‌های آماری در عمل}

متد \lr{\texttt{get\_bounds}} مهم‌ترین تابع آماری کلاس پایه است. این متد کران پایین و بالا را برای یک احتمال ناشناخته $p$ بر اساس $k$ موفقیت در $n$ آزمایش با احتمال شکست \lr{\texttt{failure\_prob}} محاسبه می‌کند:

\begin{LTR}
	\begin{lstlisting}[language=Python]
		def get_bounds(self, k: int, n: int, failure_prob: float,
		sided: Literal["two", "one_upper", "one_lower"] = "two",
		diagnostics: Dict | None = None) -> Tuple[float, float]:
		k, n = int(k), int(n)
		if not (0 <= k <= n):
		raise QKDParamError(
		f"k must be in [0, n], got k={k}, n={n}.",
		code=QKDErrorCode.PARAM_VALIDATION,
		context={"k": k, "n": n}
		)
		if not (1e-300 < failure_prob < 1):
		raise QKDParamError(
		f"Failure probability must be in (1e-300, 1), got {failure_prob}.",
		param_name="failure_prob",
		param_value=failure_prob,
		code=QKDErrorCode.PARAM_VALIDATION
		)
		if n == 0: return (0.0, 1.0)
		alpha = float(failure_prob if sided != "two" else failure_prob / 2.0)
		
		if diagnostics is not None:
		diagnostics['alpha_used'] = alpha
		
		sided_arg = "two-sided" if sided == "two" else ("upper" if sided == "one_upper" else "lower")
		
		if self.p.ci_method == ConfidenceBoundMethod.CLOPPER_PEARSON:
		interval = qkd_math.clopper_pearson_bounds(k, n, failure_prob, side=sided_arg)
		lower, upper = interval.lower, interval.upper
		elif self.p.ci_method == ConfidenceBoundMethod.HOEFFDING:
		interval = qkd_math.hoeffding_bounds(k, n, failure_prob, side=sided_arg)
		lower, upper = interval.lower, interval.upper
		elif self.p.ci_method == ConfidenceBoundMethod.GAUSSIAN:
		interval = qkd_math.gaussian_bounds(k, n, failure_prob, side=sided_arg)
		lower, upper = interval.lower, interval.upper
		else:
		raise QKDSimulationError(
		f"CI method '{self.p.ci_method.name}' not implemented.",
		code=QKDErrorCode.CONFIG,
		context={"ci_method": self.p.ci_method.name}
		)
		
		self._assert_finite("lower bound", lower)
		self._assert_finite("upper bound", upper)
		return (lower, upper)
	\end{lstlisting}
\end{LTR}

ابتدا $k$ و $n$ به \lr{\texttt{int}} تبدیل می‌شوند تا از خطاهای نوع \lr{numpy} جلوگیری شود. شرط $0 \le k \le n$ یک ضرورت منطقی است: تعداد موفقیت‌ها نمی‌تواند از تعداد آزمایش‌ها بیشتر باشد یا منفی باشد. شرط $10^{-300} < $ \lr{\texttt{failure\_prob}} $ < 1$ محدوده معتبر احتمال شکست را تضمین می‌کند. کران پایین $10^{-300}$ به جای صفر تنظیم شده تا از خطاهای عددی در محاسبات \lr{log} جلوگیری شود. اگر $n = 0$ باشد، یعنی هیچ آزمایشی انجام نشده، کران پیش‌فرض $(0, 1)$ برگردانده می‌شود که نشان‌دهنده عدم اطلاع است.

پارامتر \lr{\texttt{sided}} سه حالت دارد: \lr{two} برای کران دو طرفه، \lr{one\_upper} برای کران بالا فقط و \lr{one\_lower} برای کران پایین فقط. در حالت دو طرفه، $\alpha$ به دو نیم تقسیم می‌شود تا کران از هر دو طرف پوشش $1-\alpha$ داشته باشد. این تنظیم به صورت زیر فرمول‌بندی می‌شود:
\begin{equation}
	\alpha_{\text{eff}} = \begin{cases}
		\alpha/2 & \text{if } \text{sided} = \text{two}, \\
		\alpha   & \text{otherwise}.
	\end{cases}
\end{equation}

سپس بر اساس \lr{\texttt{ci\_method}} در پارامترها، تابع مربوطه از ماژول \lr{\texttt{qkd\_math}} فراخوانی می‌شود. هر سه تابع یک شیء \lr{interval} با ویژگی‌های \lr{\texttt{lower}} و \lr{\texttt{upper}} برمی‌گردانند. در پایان، متد \lr{\texttt{\_assert\_finite}} مقادیر را بررسی می‌کند تا اطمینان حاصل شود که متناهی هستند. در حالت \lr{AUDIT}، مقدار نامتناهی به استثنا منجر می‌شود.

\subsection{محاسبه آنتروپی باینری}

متد \lr{\texttt{binary\_entropy}} آنتروپی باینری را به صورت برداری محاسبه می‌کند:

\begin{LTR}
	\begin{lstlisting}[language=Python]
		def binary_entropy(self, p_err: float | np.ndarray | list) -> float | np.ndarray:
		p = np.asarray(p_err, dtype=np.float64)
		p = np.clip(p, self.min_prob_clamp, 1.0 - self.min_prob_clamp)
		with np.errstate(divide='ignore', invalid='ignore',
		over='raise' if self.is_audit_mode() else 'warn'):
		h = -p * np.log2(p) - (1 - p) * np.log2(1 - p)
		h = np.nan_to_num(h, nan=0.0)
		return h.item() if isinstance(p_err, (float, int)) else h
	\end{lstlisting}
\end{LTR}

ورودی می‌تواند یک عدد، یک لیست یا یک آرایه \lr{numpy} باشد. در خط دوم، ورودی به آرایه \lr{\texttt{float64}} تبدیل می‌شود. در خط سوم، مقدار $p$ به بازه $[\delta, 1-\delta]$ محدود می‌شود که در آن $\delta = $ \lr{\texttt{min\_prob\_clamp}}. این \lr{clip} از محاسبه $\log(0)$ که به $-\infty$ منجر می‌شود، جلوگیری می‌کند. محاسبه آنتروپی به صورت زیر انجام می‌شود:
\begin{equation}
	h(p) = -p\,\log_2(p) - (1-p)\,\log_2(1-p),
\end{equation}
که با توجه به \lr{clip}، همواره مقدار متناهی خواهد داشت. با این حال، به دلیل احتمال سرریز در محاسبات \lr{log} برای مقادیر بسیار کوچک، از \lr{\texttt{np.errstate}} برای کنترل رفتار خطا استفاده می‌شود. در حالت \lr{AUDIT}، سرریز به استثنا تبدیل می‌شود و در غیر این صورت هشدار داده می‌شود. در نهایت، \lr{\texttt{np.nan\_to\_num}} هر مقدار \lr{NaN} باقی‌مانده را به صفر تبدیل می‌کند. اگر ورودی اصلی یک عدد بوده، خروجی هم به صورت اسکالر با \lr{\texttt{.item()}} برگردانده می‌شود.

\subsection{متدهای ایمن عددی}

کلاس پایه چهار متد کمکی برای مدیریت شرایط عددی بحرانی دارد. متد \lr{\texttt{\_safe\_log}} برای محاسبه لگاریتم طبیعی با حداقل مقدار ایمن استفاده می‌شود:

\begin{LTR}
	\begin{lstlisting}[language=Python]
		def _safe_log(self, x: float) -> float:
		if x < 1e-300:
		if self.is_audit_mode():
		raise QKDSimulationError(
		"Log argument too small in AUDIT mode.",
		code=QKDErrorCode.SIMULATION,
		context={"x": x}
		)
		self.logger.warning("Log argument %e is smaller than safe limit, clamping.", x)
		x = 1e-300
		return math.log(x)
	\end{lstlisting}
\end{LTR}

اگر $x$ کمتر از $10^{-300}$ باشد، در حالت \lr{AUDIT} استثنا پرتاب می‌شود و در غیر این صورت $x$ به $10^{-300}$ \lr{clamp} می‌شود تا از خطای \lr{\texttt{math domain error}} جلوگیری شود. این متد در محاسبات \lr{log-likelihood} و \lr{relative entropy} که در برخی اثبات‌های امنیتی به کار می‌روند، حیاتی است.

متد \lr{\texttt{\_safe\_divide}} برای تقسیم ایمن با مخرج نزدیک به صفر است:

\begin{LTR}
	\begin{lstlisting}[language=Python]
		def _safe_divide(self, num: float, den: float, default: float = 0.0,
		name: str = '', error_codes: list | None = None) -> float:
		if abs(den) < NUMERIC_ABS_TOL:
		if error_codes is not None:
		error_codes.append(ErrorCode.SAFE_DIVIDE_FALLBACK)
		self.logger.warning("Denominator for '%s' is near zero (%e). Returning default value %f.",
		name, den, default)
		return default
		return num / den
	\end{lstlisting}
\end{LTR}

اگر قدر مطلق مخرج کمتر از \lr{\texttt{NUMERIC\_ABS\_TOL}} باشد، مقدار \lr{\texttt{default}} (پیش‌فرض صفر) برگردانده می‌شود و در صورت ارائه لیست \lr{\texttt{error\_codes}}، کد \lr{\texttt{SAFE\_DIVIDE\_FALLBACK}} به آن اضافه می‌شود. این متد در محاسباتی مانند نسبت خطا به آشکارسازی که ممکن است مخرج صفر داشته باشند، ضروری است.

متد \lr{\texttt{\_clamp\_nonneg}} مقادیر منفی کوچک را به صفر تبدیل می‌کند:

\begin{LTR}
	\begin{lstlisting}[language=Python]
		def _clamp_nonneg(self, x: float, name: str, error_codes: list | None = None) -> float:
		if x < 0:
		if error_codes is not None:
		error_codes.append(ErrorCode.NUMERICAL_CLAMP_APPLIED)
		if x < -NUMERIC_ABS_TOL and self.is_audit_mode():
		raise QKDSimulationError(
		f"Negative value for {name} ({x}) exceeded tolerance in AUDIT mode.",
		code=QKDErrorCode.SIMULATION,
		context={"name": name, "value": x, "tolerance": -NUMERIC_ABS_TOL}
		)
		self.logger.debug("Clamped negative value for %s: %e -> 0.0", name, x)
		return 0.0
		return x
	\end{lstlisting}
\end{LTR}

مقادیر منفی کوچک (در حد تلورانس) به صفر \lr{clamp} می‌شوند زیرا در محاسبات عددی با نوسانات ممیز شناور، ممکن است مقادیری که باید صفر باشند، منفی کوچک گزارش شوند. اما مقادیر منفی بزرگ در حالت \lr{AUDIT} به استثنا منجر می‌شوند، زیرا نشان‌دهنده خطای منطقی جدی است.

متد \lr{\texttt{\_clamp\_key\_length}} طول کلید نهایی را به عدد صحیح نامنفی تبدیل می‌کند:

\begin{LTR}
	\begin{lstlisting}[language=Python]
		def _clamp_key_length(self, l_float: float) -> int:
		return int(max(0.0, math.floor(l_float)))
	\end{lstlisting}
\end{LTR}

ابتدا با \lr{\texttt{max(0.0, ...)}} اطمینان حاصل می‌شود که مقدار منفی نیست، سپس با \lr{\texttt{math.floor}} به عدد صحیح پایین گرد می‌شود و نهایتاً به \lr{\texttt{int}} تبدیل می‌شود. این رفتار محافظه‌کارانه تضمین می‌کند که طول کلید هرگز بیش از مقدار مجاز ریاضی نباشد.

\subsection{مدیریت زمان‌سنجی و گزارش‌گیری \lr{Audit}}

متد \lr{\texttt{\_timed}} یک \lr{context manager} برای ثبت زمان اجرای قطعات کد است:

\begin{LTR}
	\begin{lstlisting}[language=Python]
		@contextmanager
		def _timed(self, timings: Dict, key: str) -> Generator:
		t0 = time.perf_counter()
		yield
		timings[key] = time.perf_counter() - t0
	\end{lstlisting}
\end{LTR}

با استفاده از \lr{\texttt{time.perf\_counter}} که بالاترین دقت زمان‌سنجی را ارائه می‌دهد، زمان شروع و پایان ثبت شده و اختلاف در دیکشنری \lr{\texttt{timings}} با کلید مربوطه ذخیره می‌شود. این اطلاعات در \lr{\texttt{SolverDiagnostics}} برای شناسایی گلوگاه‌های عملکرد به کار می‌رود.

متد \lr{\texttt{audit\_dump}} یک گزارش کامل از اجرا تولید می‌کند:

\begin{LTR}
	\begin{lstlisting}[language=Python]
		def audit_dump(self, result: KeyCalculationResult, decoy: DecoyEstimates,
		stats: Dict, start_time: float) -> Dict:
		dump = {
			"params": self.p.to_summary_dict(redact=True),
			"mode": self.mode.name,
			"proof_implementation": f"{self.__class__.__name__} v{self.__implementation_version__}",
			"stats_map": {k: asdict(v) for k, v in stats.items()},
			"decoy_estimates": decoy.as_serializable(),
			"key_calculation_result": result.as_serializable(),
			"timestamps": {'start': start_time, 'end': time.time()},
		}
		if self.is_audit_mode():
		y1_margin = decoy.yield_1_lower_bound - self.min_y1_threshold
		e1_margin = 0.5 - decoy.error_rate_1_upper_bound
		dump['sensitivity_summary'] = {'y1_margin': y1_margin, 'e1_margin': e1_margin}
		return dump
	\end{lstlisting}
\end{LTR}

گزارش شامل پارامترها (با \lr{\texttt{redact=True}} برای حذف اطلاعات حساس مانند کلیدها)، حالت اجرا، نسخه پیاده‌سازی، آمار کامل، تخمین‌های \lr{Decoy}، نتیجه محاسبه کلید و زمان‌بندی است. در حالت \lr{AUDIT}، یک خلاصه حساسیت اضافه می‌شود که حاشیه $Y_1^L$ نسبت به آستانه و حاشیه $e_1^U$ نسبت به حد 0.5 (حد بی‌اطلاعی) را نشان می‌دهد. این حاشیه‌ها معیار مهمی برای سنجش کیفیت اجرا هستند؛ هرچه حاشیه بیشتر، اجرا امن‌تر است.

\subsection{توضیح تصمیم کلید}

متد \lr{\texttt{explain\_key\_decision}} یک پیام متنی قابل فهم برای انسان تولید می‌کند که علت تولید نشدن کلید را توضیح می‌دهد:

\begin{LTR}
	\begin{lstlisting}[language=Python]
		def explain_key_decision(self, result: KeyCalculationResult, decoy: DecoyEstimates) -> str:
		if result.secure_key_length > 0:
		return f"Secure key generated. Final length: {result.secure_key_length} bits."
		reasons = []
		if ErrorCode.INSUFFICIENT_STATISTICS in result.error_codes:
		reasons.append(f"insufficient single-photon yield (Y1_L={decoy.yield_1_lower_bound:.2e} "
		f"< threshold={self.min_y1_threshold:.2e})")
		if ErrorCode.LP_INFEASIBLE in result.error_codes:
		reasons.append("decoy-state analysis was infeasible")
		if ErrorCode.LEAKAGE_EXCEEDS_ENTROPY in result.error_codes:
		reasons.append(f"information leakage ({result.error_correction_leakage + result.privacy_amplification_term:.2f}) "
		"exceeded available entropy")
		if not reasons:
		reasons.append("an unspecified condition led to a non-positive key length")
		return f"No secure key generated because {', and '.join(reasons)}."
	\end{lstlisting}
\end{LTR}

اگر طول کلید مثبت باشد، پیام موفقیت برگردانده می‌شود. در غیر این صورت، بر اساس \lr{\texttt{error\_codes}} در نتیجه، دلایل مختلف جمع‌آوری می‌شوند. اگر هیچ کد خطای مشخصی وجود نداشته باشد، پیام عمومی نمایش داده می‌شود. این متد برای \lr{logging} و گزارش به کاربر نهایی بسیار مفید است و تشخیص علت مشکلات را ساده می‌کند.

\subsection{مشتق کردن \lr{Seed} فرزند با \lr{HMAC}}

تابع \lr{\texttt{derive\_child\_seed}} برای تولید \lr{seed}های قطعی و قابل بازتولید از یک \lr{master seed} استفاده می‌شود:

\begin{LTR}
	\begin{lstlisting}[language=Python]
		@lru_cache(maxsize=128)
		def derive_child_seed(master_seed: int, index: int) -> int:
		"""Deterministically derives a child seed. Not for cryptographic use."""
		key = str(master_seed).encode()
		msg = str(index).encode()
		digest = hmac.new(key, msg, digestmod=hashlib.sha256).digest()
		return int.from_bytes(digest[:8], 'big') % MAX_SEED_INT_PCG64 or 1
	\end{lstlisting}
\end{LTR}

این تابع با استفاده از \lr{HMAC-SHA256} یک \lr{seed} فرزند از \lr{master seed} و یک اندیس تولید می‌کند. کلید \lr{HMAC} از \lr{master seed} و پیام از اندیس ساخته می‌شود. سپس 8 بایت اول \lr{digest} به عدد صحیح تبدیل شده و با \lr{\texttt{MAX\_SEED\_INT\_PCG64}} (حداکثر مقدار مجاز برای \lr{PCG64}) modulo می‌شود. عبارت \lr{\texttt{or 1}} اطمینان می‌کند که نتیجه هرگز صفر نباشد زیرا \lr{PCG64} به \lr{seed} غیر صفر نیاز دارد. این تابع با \lr{\texttt{@lru\_cache(maxsize=128)}} کش می‌شود تا در فراخوانی‌های مکرر با همان ورودی، محاسبه مجدد انجام نشود. این مهم است زیرا در شبیه‌سازی‌های بزرگ با هزاران \lr{seed} فرزند، صرفه‌جویی در زمان محاسبه قابل توجه است.

از نظر ریاضی، این تابع به صورت زیر عمل می‌کند:
\begin{equation}
	\text{seed}_i = \left(\text{int}_{64}\!\left(\text{HMAC}_{\text{SHA256}}(M, i)\right) \bmod N_{\max}\right) \;\vee\; 1,
\end{equation}
که در آن $M$ دانهٔ اصلی (\lr{master seed})، $i$ اندیس و $N_{\max} = \text{\lr{\texttt{MAX\_SEED\_INT\_PCG64}}}$ است. عملگر $\lor$ در اینجا به معنای \lr{OR} بیتی است تا صفر بودن احتمالیِ خروجیِ حاصل را به یک تبدیل کند. توجه شود که در مستندات کد (\lr{docstring}) به‌صراحت اشاره شده است که این تابع برای مصارف غیررمزنگاری طراحی شده و صرفاً برای تولید دانه‌های (\lr{seed}های) تکرارپذیر در شبیه‌سازی به کار می‌رود.



\section{ارسال و دریافت با سایر ماژول‌ها}

\subsection{ورودی‌های ماژول}

ماژول \lr{\texttt{base.py}} مجموعه‌ای از وابستگی‌های داخلی و خارجی دارد. وابستگی‌های خارجی شامل \lr{numpy} برای محاسبات برداری، \lr{scipy} برای توابع آماری و توزیع \lr{beta}، و \lr{mpmath} (اختیاری) برای دقت بالا هستند. وابستگی‌های داخلی به صورت زیر ساختارمند شده‌اند:

\begin{itemize}
	\item \lr{\texttt{qkd.params.QKDParams}}: کلاس پارامترهای QKD که تمام تنظیمات شبیه‌سازی را در خود نگه می‌دارد. این کلاس در سازنده \lr{\texttt{FiniteKeyProof}} وارد می‌شود و در سراسر طول عمر نمونه در دسترس است.
	\item \lr{\texttt{qkd.datatypes.ConfidenceBoundMethod}}: \lr{Enum} که روش کران آماری را مشخص می‌کند. این مقدار در \lr{\texttt{QKDParams.ci\_method}} ذخیره شده و در متد \lr{\texttt{get\_bounds}} استفاده می‌شود.
	\item \lr{\texttt{qkd.constants}}: شامل ثابت‌های عددی مانند \lr{\texttt{ENTROPY\_PROB\_CLAMP}}، \lr{\texttt{NUMERIC\_ABS\_TOL}} و \lr{\texttt{MAX\_SEED\_INT\_PCG64}}. این ثابت‌ها در سراسر کد برای تنظیم رفتار عددی استفاده می‌شوند.
	\item \lr{\texttt{qkd.exceptions}}: شامل استثناهای سفارشی مانند \lr{\texttt{ConfigurationError}}، \lr{\texttt{ParameterValidationError}}، \lr{\texttt{QKDSimulationError}} و \lr{\texttt{ErrorCode}}. این استثناها در سراسر کد برای گزارش خطاهای ساختاریافته به کار می‌روند.
	\item \lr{\texttt{qkd.utils.math}}: ماژول توابع ریاضی کمکی شامل \lr{\texttt{clopper\_pearson\_bounds}}، \lr{\texttt{hoeffding\_bounds}} و \lr{\texttt{gaussian\_bounds}}. این توابع پیاده‌سازی واقعی کران‌های آماری را ارائه می‌دهند و ماژول \lr{\texttt{base.py}} تنها آن‌ها را فراخوانی می‌کند.
\end{itemize}

این الگوی واردسازی با \lr{try/except} به ماژول اجازه می‌دهد در صورت بروز مشکل، پیام خطای روشنی تولید کند و به جای خطای مبهم \lr{\texttt{ImportError}}، اطلاعات زمینه‌ای درباره ماژول مفقود ارائه دهد.

\subsection{خروجی‌های ماژول}

خروجی اصلی این ماژول کلاس \lr{\texttt{FiniteKeyProof}} و \lr{dataclass}های مرتبط است که در \lr{\texttt{\_\_all\_\_}} فهرست شده‌اند:

\begin{LTR}
	\begin{lstlisting}[language=Python]
		__all__ = [
		'FiniteKeyProof', 'DecoyEstimates', 'KeyCalculationResult', 'TallyCounts',
		'EpsilonAllocation', 'SolverDiagnostics', 'ProofMode', 'ErrorCode',
		'ConfidenceBoundMethod', 'TALLY_KEY_Z_SIGNAL', 'TALLY_KEY_X_SIGNAL',
		'derive_child_seed'
		]
	\end{lstlisting}
\end{LTR}

این فهرست به ماژول‌های دیگر اجازه می‌دهد با \lr{\texttt{from qkd.proofs.base import *}} فقط نمادهای عمومی مورد نظر را وارد کنند. زیرکلاس‌های اثبات امنیتی مانند \lr{\texttt{Ma2005Proof}}، \lr{\texttt{Lim2014Proof}} و \lr{\texttt{GLLPProof}} (که در فایل‌های جداگانه پیاده‌سازی می‌شوند) از این کلاس پایه ارث می‌برند و متدهای انتزاعی را پیاده‌سازی می‌کنند.

خروجی عملی ماژول در زمان اجرا به دو شکل است:
\begin{enumerate}
	\item شیء \lr{\texttt{DecoyEstimates}} که خروجی متد \lr{\texttt{estimate\_yields\_and\_errors}} است و شامل $Y_1^L$، $e_1^U$، \lr{\texttt{is\_feasible}}، \lr{\texttt{failure\_prob\_used}} و \lr{\texttt{diagnostics}} می‌شود. این شیء به عنوان ورودی به متد \lr{\texttt{calculate\_key\_length}} داده می‌شود.
	\item شیء \lr{\texttt{KeyCalculationResult}} که خروجی متد \lr{\texttt{calculate\_key\_length}} است و شامل طول کلید امن، نشت‌ها، کران خطای فاز، کدهای خطا و اطلاعات تشخیصی است. این شیء توسط لایه بالایی فریم‌ورک (مانند \lr{\texttt{qkd.simulator}}) دریافت شده و در گزارش نهایی شبیه‌سازی استفاده می‌شود.
\end{enumerate}

\subsection{جریان داده در سراسر فریم‌ورک}

جریان داده کلی به این شکل است:
\begin{enumerate}
	\item لایه شبیه‌ساز پالس‌های نوری را تولید و از کانال کوانتومی عبور می‌دهد.
	\item آشکارسازها در سمت باب، پالس‌ها را ثبت کرده و شمارش‌ها را در \lr{\texttt{TallyCounts}} ذخیره می‌کنند.
	\item این شمارش‌ها به صورت یک دیکشنری \lr{\texttt{stats\_map}} با کلیدهایی مانند \lr{\texttt{TALLY\_KEY\_Z\_SIGNAL}} و \lr{\texttt{TALLY\_KEY\_X\_SIGNAL}} به اثبات امنیتی داده می‌شوند.
	\item متد \lr{\texttt{estimate\_yields\_and\_errors}} با استفاده از کران‌های آماری \lr{\texttt{get\_bounds}} و حل \lr{LP}، \lr{\texttt{DecoyEstimates}} را تولید می‌کند.
	\item متد \lr{\texttt{calculate\_key\_length}} با استفاده از \lr{\texttt{DecoyEstimates}}، \lr{\texttt{binary\_entropy}} و متدهای ایمن عددی، \lr{\texttt{KeyCalculationResult}} نهایی را محاسبه می‌کند.
	\item لایه بالایی می‌تواند با \lr{\texttt{explain\_key\_decision}} علت نتیجه را به صورت متنی دریافت کند و با \lr{\texttt{audit\_dump}} یک گزارش کامل برای \lr{audit} تولید کند.
\end{enumerate}

این جداسازی مراحل به گونه‌ای طراحی شده که هر مرحله به صورت مستقل قابل آزمایش و اعتبارسنجی باشد. به عنوان مثال، می‌توان \lr{\texttt{DecoyEstimates}} را به صورت دستی ساخت و به \lr{\texttt{calculate\_key\_length}} داد تا رفتار آن در شرایط خاص بررسی شود. این ماژول‌بندی از اصل \lr{Single Responsibility} پیروی می‌کند و قابلیت نگهداری و توسعه کد را به طور قابل توجهی افزایش می‌دهد.

	
	\chapter{اثبات امنیت کلید محدود لیم ۲۰۱۴}
	\label{ch:lim2014}
	\sloppy
\emergencystretch=5em
\hbadness=10000
\vbadness=10000
\tolerance=2000
\providecolor{codebg}{RGB}{248,248,248}
\providecolor{codeframe}{RGB}{200,200,200}
\providecolor{keyword}{RGB}{0,0,180}
\providecolor{string}{RGB}{163,21,21}
\providecolor{comment}{RGB}{0,128,0}
\lstset{
	basicstyle=\ttfamily\scriptsize,
	breaklines=true,
	breakatwhitespace=false,
	columns=fullflexible,
	keepspaces=true,
	frame=single,
	framesep=2pt,
	showstringspaces=false,
	backgroundcolor=\color{codebg},
	rulecolor=\color{codeframe},
	keywordstyle=\color{keyword}\bfseries,
	stringstyle=\color{string},
	commentstyle=\color{comment}\itshape,
	numberstyle=\tiny\color{gray},
	numbers=left,
	numbersep=6pt,
	xleftmargin=14pt,
	framexleftmargin=14pt,
	literate=
	{ą}{{\v{a}}}1 {ę}{{\k{e}}}1 {ż}{{\.z}}1 {ś}{{\'s}}1
	{ł}{{\l}}1 {ć}{{\'c}}1 {ń}{{\'n}}1 {ó}{{\'o}}1 {ź}{{\'z}}1
	{→}{{$\rightarrow$}}1 {←}{{$\leftarrow$}}1 {≠}{{$\neq$}}1
}

\section{نمای‌ کلی و هدف ماژول}

فایل \lr{\texttt{qkd/proofs/lim2014.py}} پیاده‌سازی وفادارانه چارچوب اثبات امنیتی کلید متناهی مبتنی بر \lr{Decoy-State} است که توسط \lr{Lim}، \lr{Curty}، \lr{Walenta}، \lr{Xu} و \lr{Zbinden} در سال 2014 در مقاله مرجع \lr{``Concise security bounds for practical decoy-state quantum key distribution''} (در مجله \lr{Physical Review A}، جلد 89، شماره 2، صفحه 022307) ارائه شده است. این چارچوب به دلیل سادگی، فشردگی و کاربردپذیری عملی، به یکی از پذیرفته‌شده‌ترین مراجع برای محاسبه طول کلید امن در سیستم‌های تجاری و آزمایشگاهی QKD تبدیل شده است. مزیت اصلی این چارچوب نسبت به روش‌های قبلی مانند \lr{Ma 2005} این است که فرمول‌های بسته و تحلیلی برای تخمین کران‌های پایین بازده خلأ $s_{X,0}$، کران پایین بازده تک‌فوتونی $s_{X,1}$، کران بالای خطای تک‌فوتونی $v_{Z,1}$ و کران بالای خطای فاز $\varphi_X$ ارائه می‌دهد که نیاز به حل برنامه‌ریزی خطی (\lr{LP}) را برای حالت استاندارد سه شدت (سیگنال، \lr{decoy} ضعیف و خلأ) از بین می‌برد. این فایل، خروجی نهایی محاسبه امنیتی کلید را در قالب شیء \lr{\texttt{KeyCalculationResult}} تولید می‌کند که شامل طول کلید امن، نشت تصحیح خطا، کران بالای خطای فاز، کدهای خطا و اطلاعات تشخیصی (\lr{diagnostics}) است.

هدف این ماژول دوگانه است: نخست، ارائه یک مرجع عددی قابل اعتماد و قابل بازتولید برای محاسبه طول کلید امن در چارچوب \lr{Lim 2014} که در آن تمام فرمول‌های معادلات (1) تا (5) مقاله دقیقاً به صورت کد پیاده‌سازی شده‌اند؛ دوم، فراهم آوردن یک لایه توسعه برای مواردی که کاربر بیش از سه شدت نور نیاز دارد. در حالت چهار یا چند شدت، چارچوب \lr{Lim 2014} فرمول بسته ارائه نمی‌کند و صرفاً در پیوست مقاله اشاره می‌کند که تحلیل به طور سرراست به هر تعداد سطح شدت قابل تعمیم است. این ماژول برای پوشش این حالت، یک مسیر جایگزین مبتنی بر \lr{LP} با محدودیت‌های غیرمنفی‌بودن و کران بالای صریح و با یک \lr{tail relaxation} ارائه می‌دهد که در گزارش \lr{diagnostics} به روشنی به عنوان توسعه فراتر از نتایج بسته مقاله علامت‌گذاری شده است. این شفافیت روش‌شناسی برای کاربردهای حساس مانند \lr{certification} و بازبینی امنیتی توسط شخص ثالث ضروری است.

یکی از نوآوری‌های کلیدی این نسخه از پیاده‌سازی، استفاده از یک استراتژی \lr{\textbf{هیبرید Hoeffding/Chernoff}} برای بازه‌های اطمینان است که در ضمیمه A مقاله صراحتاً مجاز شناخته شده است. مقاله اصلی برای سادگی فرمول بسته، از نامساوی \lr{Hoeffding} با یک نوسان $\delta$ یکسان برای همه شدت‌ها استفاده می‌کند، اما این انتخاب در شرایطی که شمارش \lr{decoy} یا خلأ بسیار کم است (مانند فواصل طولانی کانال که در آن‌ها $n_{X,\mu_3}$ ممکن است تنها چند رویداد باشد)، به‌طور پاتولوژیک شل می‌شود و کران‌های پایین منفی یا کران‌های بالا بیش از کل فیزیکی تولید می‌کند. پیاده‌سازی هیبرید در سه رژیم عمل می‌کند: برای شمارش‌های بزرگ ($n_k \ge 10\,\delta_{\text{paper}}$) از \lr{Hoeffding} وفادار به کاغذ استفاده می‌کند؛ برای شمارش‌های کوچک اما ناصفر ($0 < n_k < 10\,\delta_{\text{paper}}$) از \lr{Chernoff} به‌ازای هر شدت با $\delta_k = \sqrt{2\,n_k\,\ln(1/\varepsilon)}$ استفاده می‌کند که تنگ‌تر از \lr{Hoeffding} است؛ و برای $n_k = 0$ از کران بالای \lr{Chernoff} ضربی $\ln(1/\varepsilon_{\text{per-event}})$ (قانون سه) استفاده می‌کند که متناهی و بسیار کوچکتر از $\delta_{\text{paper}}$ است. این استراتژی به‌طور قابل اثبات کوچکتر یا مساوی \lr{Hoeffding} برای همه $n_k$ است، بنابراین امنیت هرگز تضعیف نمی‌شود.

این فایل از کلاس پایه \lr{\texttt{FiniteKeyProof}} (که در \lr{\texttt{qkd/proofs/base.py}} تعریف شده) ارث می‌برد و متدهای انتزاعی \lr{\texttt{allocate\_epsilons}}، \lr{\texttt{get\_epsilon\_policy}}، \lr{\texttt{estimate\_yields\_and\_errors}}، \lr{\texttt{calculate\_key\_length}} و \lr{\texttt{notation\_map}} را پیاده‌سازی می‌کند. همچنین یک \lr{dataclass} جدید به نام \lr{\texttt{DecoyEstimatesLim2014}} معرفی می‌کند که فیلدهای اضافی مانند $s_{X,0}^L$، $s_{X,1}^L$، $s_{Z,1}^L$، $v_{Z,1}^U$ و $\varphi_X^U$ را نگه می‌دارد و از \lr{\texttt{DecoyEstimates}} پایه ارث می‌برد. این طراحی به کاربر اجازه می‌دهد هم به رابط عمومی \lr{\texttt{FiniteKeyProof}} و هم به اطلاعات اختصاصی چارچوب \lr{Lim 2014} دسترسی داشته باشد. علاوه بر این، سازنده کلاس \lr{\texttt{\_\_init\_\_}} به‌صورت خودکار روش بازه اطمینان پارامتر \lr{\texttt{ci\_method}} را به \lr{\texttt{CHERNOFF}} بازنویسی می‌کند زیرا امنیت کلید متناهی در این چارچوب به‌طور صریح به کران‌های تمرکز \lr{Chernoff} وابسته است.

وظیفه نهایی این ماژول، تبدیل شمارش‌های آماری شبکه آشکارسازی به یک عدد صحیح مثبت (یا صفر) به عنوان طول کلید امن است. این تبدیل با رعایت سخت‌گیرانه تمام محافظه‌کاری‌های ریاضی و آماری انجام می‌شود تا هیچ طول کلیدی که از نظر تئوری قابل اثبات نیست، تولید نشود. یک لایه دفاعی فیزیکی جدید در متد \lr{\texttt{calculate\_key\_length}} اضافه شده است که طول کلید نهایی را به مجموع بیت‌های \lr{sifted} ($n_X + n_Z$) محدود می‌کند؛ این قید در صورت بروز هرگونه باگ داخلی (مانند حلقه \lr{kappa-solver}، \lr{fallback} در \lr{LP}، یا سرریز در تخمین \lr{decoy}) فعال می‌شود و یک هشدار گزارش می‌کند. در صورت بروز هرگونه ناسازگاری پیکربندی، آمار ناکافی یا خطای عددی، خروجی به جای استثنا، یک \lr{\texttt{KeyCalculationResult}} با طول کلید صفر و کد خطای مناسب تولید می‌کند تا خط لوله شبیه‌سازی متوقف نشود و علت شکست در گزارش تشخیصی ثبت گردد.

\section{مبانی تئوری و فیزیکی}

\subsection{پروتکل \lr{Decoy-State} با سه شدت}

در پروتکل استاندارد \lr{BB84} با یک منبع تک‌شدت، امنیت در برابر حمله \lr{Photon-Number Splitting (PNS)} قابل تضمین نیست؛ زیرا یک مهاجم با تقسیم فوتون‌های چندفوتونی می‌تواند بدون ایجاد خطا اطلاعات کلید را استخراج کند. راه‌حل \lr{Decoy-State} این است که آلیس به طور تصادفی شدت نور ارسالی را میان چندین سطح تغییر می‌دهد و باب پس از اندازه‌گیری، با مقایسه نرخ آشکارسازی در شدت‌های مختلف، بازده تک‌فوتونی را تخمین می‌زند. در چارچوب \lr{Lim 2014}، سه شدت $\mu_1 > \mu_2 > \mu_3 \ge 0$ در نظر گرفته می‌شود که $\mu_1$ شدت سیگنال، $\mu_2$ شدت \lr{decoy} ضعیف و $\mu_3$ شدت خلأ است. شرط کلیدی $\mu_1 > \mu_2 + \mu_3$ در تابع \lr{\texttt{\_validate\_three\_intensity\_order}} بررسی می‌شود و تضمین می‌کند که مخرج فرمول‌های تحلیلی مثبت باشد. این شرط در کنار $\mu_2 > \mu_3 \ge 0$ یک مجموعه کامل از قیود ترتیبی را تشکیل می‌دهد که در آن مخرج $D = \mu_1(\mu_2 - \mu_3) - \mu_2^2 + \mu_3^2$ در معادله (3) همیشه مثبت می‌ماند.

توزیع آماری فوتون‌های ارسالی از یک منبع لیز با توزیع \lr{Poisson} پیروی می‌کند. اگر میانگین فوتون $\mu$ باشد، احتمال ارسال $n$ فوتون برابر است با:
\begin{equation}
	P(n \mid \mu) = e^{-\mu}\,\frac{\mu^n}{n!}.
\end{equation}
احتمال کلی تولید $n$ فوتون بدون توجه به شدت انتخاب‌شده، $\tau_n$ نامیده می‌شود و به صورت زیر تعریف می‌شود:
\begin{equation}
	\tau_n = \sum_k p_k\,P(n \mid \mu_k) = \sum_k p_k\,e^{-\mu_k}\,\frac{\mu_k^n}{n!},
\end{equation}
که در آن $p_k$ احتمال انتخاب شدت $\mu_k$ توسط آلیس است. این کمیت در تابع \lr{\texttt{\_get\_tau\_n}} پیاده‌سازی شده و نقش مرکزی در همه فرمول‌های تحلیلی چارچوب دارد. توجه شود که $\tau_0$ و $\tau_1$ به‌طور خاص در معادلات (2) و (3) ظاهر می‌شوند و مثبت بودن آن‌ها یک پیش‌نیاز برای ادامه محاسبه است؛ در صورت صفر شدن این مقادیر (که می‌تواند در پیکربندی‌های غیرفیزیکی رخ دهد)، ماژول به یک نتیجه صفر با کد خطای \lr{\texttt{INSUFFICIENT\_STATISTICS}} بازمی‌گردد.

\subsection{تخصیص \texorpdfstring{$\varepsilon$}{epsilon} و ۲۱ رویداد شکست}

یکی از نوآوری‌های کلیدی \lr{Lim 2014}، ارائه یک طرح صریح برای تقسیم بودجه امنیتی $\varepsilon_{\text{sec}}$ میان ۲۱ رویداد شکست مستقل است. در ضمیمه B مقاله، نشان داده می‌شود که اگر هر رویداد شکست به یک مقدار مشترک $\varepsilon$ تنظیم شود، کل $\varepsilon_{\text{sec}}$ از رابطه زیر به دست می‌آید:
\begin{equation}
	\varepsilon_{\text{sec}} = 21\,\varepsilon \quad \Rightarrow \quad \varepsilon = \frac{\varepsilon_{\text{sec}}}{21}.
\end{equation}
این ۲۱ رویداد شامل کران‌های \lr{Hoeffding} برای شمارش آشکارسازی و خطا در چندین پایه و شدت، کران تخمین خطای فاز، کران تصحیح خطا و کران تقویت حریم خصوصی است. در کد، ثابت \lr{\texttt{LIM2014\_FAILURE\_EVENTS\_3\_INTENSITY = 21.0}} این مقدار را نگه می‌دارد و در تابع \lr{\texttt{allocate\_epsilons}} برای محاسبه $\varepsilon$ تک‌آزمون استفاده می‌شود. این تخصیص محافظه‌کارانه است زیرا فرض می‌کند همه ۲۱ رویداد همزمان در بدترین حالت ممکن شکست می‌خورند. در مسیر \lr{N-decoy} (بیش از سه شدت)، همین تخصیص ۲۱ رویدادی به‌صورت محافظه‌کارانه بازاستفاده می‌شود زیرا مقاله فرمول بسته‌ای برای این حالت ارائه نمی‌کند.

علاوه بر این، یک پارامتر مرتبط به نام \lr{\texttt{failure\_prob\_used}} در \lr{\texttt{DecoyEstimatesLim2014}} وجود دارد که در عمل برابر $\varepsilon_{\text{sec}}$ استفاده‌شده در محاسبه است. این فیلد به کاربر اجازه می‌دهد بداند آیا در مسیر \lr{kappa-solver} تکراری، یک $\varepsilon_{\text{sec}}$ متفاوت از مقدار پیش‌فرض \lr{\texttt{p.eps\_sec}} استفاده شده است یا خیر. این شفافیت برای بازبینی امنیتی و \lr{audit} حساس است زیرا تغییر $\varepsilon_{\text{sec}}$ مستقیماً بر طول کلید نهایی تأثیر می‌گذارد.

\subsection{کران‌های \lr{Hoeffding} و انگیزه استراتژی هیبرید}

برای تخمین بازده و خطا از روی شمارش‌های مشاهده‌شده، چارچوب \lr{Lim 2014} از نامساوی \lr{Hoeffding} استفاده می‌کند. اگر $n_{X,k}$ تعداد آشکارسازی‌های مشاهده‌شده در پایه $X$ با شدت $\mu_k$ باشد و $n_X = \sum_k n_{X,k}$ کل آشکارسازی‌های پایه $X$ باشد، کران‌های بالا و پایین با یک نوسان $\delta$ به صورت زیر تعریف می‌شوند:
\begin{equation}
	n_{X,k}^{\pm} = \frac{e^{\mu_k}}{p_k}\left(n_{X,k} \pm \delta\right), \qquad \delta = \sqrt{\frac{n_X}{2}\,\ln\!\left(\frac{21}{\varepsilon_{\text{sec}}}\right)}.
\end{equation}
ضریب پیشانی $e^{\mu_k}/p_k$ برای تبدیل شمارش مشاهده‌شده در شدت $k$ به تخمین کلان بازده آن شدت است. در پیاده‌سازی مقاله، $\delta$ برای همه شدت‌ها یکسان و تنها به $n_X$ (کل پایه) وابسته است. این انتخاب ساده و بسته است، اما یک ضعف عملی جدی دارد: وقتی شمارش یک شدت خاص $n_{X,k}$ بسیار کمتر از $n_X$ باشد (مانند شمارش خلأ در فواصل طولانی که در آن $n_{X,\mu_3}$ ممکن است تنها چند رویداد باشد در حالی که $n_X$ صدها هزار است)، نوسان $\delta$ مبتنی بر $n_X$ می‌تواند به‌طور قابل توجهی بزرگتر از خود $n_{X,k}$ شود. به‌عنوان مثال، در یک شبیه‌سازی بلندمسافت با $n_X = 10^6$ و $\varepsilon_{\text{sec}} = 10^{-9}$، مقدار $\delta \approx \sqrt{5\times10^5 \times \ln(21\times10^9)} \approx 141$ می‌شود، در حالی که شمارش مشاهده‌شده خلأ ممکن است تنها $n_{X,\mu_3} = 5$ باشد. در این حالت، کران بالای \lr{Hoeffding} $n_{X,\mu_3}^{+} \approx (e^{\mu_3}/p_3)\times 146$ می‌شود که حدود 114 برابر مشاهده است و کاملاً غیرفیزیکی است.

این مشکل در ضمیمه A مقاله به‌صورت صریح تأیید می‌شود: \lr{``the Hoeffding choice is for closed-form simplicity; tighter concentration inequalities may be used''}. پیاده‌سازی فعلی با معرفی یک استراتژی هیبرید، این مجوز را به‌کار می‌گیرد تا در رژیم‌هایی که \lr{Hoeffding} پاتولوژیک است، از کران‌های تمرکز تنگ‌تر استفاده کند بدون آنکه فرض امنیتی تغییر کند. این استراتژی در ادامه توضیح داده می‌شود.

\subsection{استراتژی هیبرید \lr{Hoeffding/Chernoff}}

استراتژی هیبرید بر اساس یک آستانه که نسبت خطای آماری نسبی را کنترل می‌کند، تعریف می‌شود. اگر $\delta_{\text{paper}} = \sqrt{(n_X/2)\,\ln(21/\varepsilon_{\text{sec}})}$ نوسان \lr{Hoeffding} مقاله باشد، آستانه هیبرید به‌صورت $n_k \ge 10\,\delta_{\text{paper}}$ تعریف می‌شود. این آستانه منطقی است زیرا در این نقطه، نسبت $\delta_{\text{paper}}/n_k \le 0.1$ است، یعنی خطای آماری نسبی کمتر از 10\% است که در آن \lr{Hoeffding} به‌طور معقولی تنگ عمل می‌کند. در سه رژیم زیر، بازه اطمینان به‌صورت متفاوت محاسبه می‌شود:

\textbf{رژیم 1 (شمارش بزرگ):} وقتی $n_k \ge 10\,\delta_{\text{paper}}$ و $n_k > 0$، از \lr{Hoeffding} وفادار به کاغذ استفاده می‌شود:
\begin{equation}
	n_{X,k}^{\pm} = \frac{e^{\mu_k}}{p_k}\left(n_{X,k} \pm \delta_{\text{paper}}\right).
\end{equation}
این رژیم در شرایطی که شمارش کافی وجود دارد، رفتار پیاده‌سازی را دقیقاً با مقاله همخوان می‌کند.

\textbf{رژیم 2 (شمارش کوچک اما ناصفر):} وقتی $0 < n_k < 10\,\delta_{\text{paper}}$، از \lr{Chernoff} به‌ازای هر شدت استفاده می‌شود:
\begin{equation}
	n_{X,k}^{\pm} = \frac{e^{\mu_k}}{p_k}\left(n_{X,k} \pm \delta_k\right), \qquad \delta_k = \sqrt{2\,n_k\,\ln\!\left(\frac{1}{\varepsilon_{\text{per-event}}}\right)},
\end{equation}
که در آن $\varepsilon_{\text{per-event}} = \varepsilon_{\text{sec}}/21$ طبق ضمیمه B است. توجه شود که $\ln(21/\varepsilon_{\text{sec}}) = \ln(1/\varepsilon_{\text{per-event}})$، بنابراین همان جمله لگاریتمی در هر دو رژیم استفاده می‌شود و تنها تفاوت در مقیاس نوسان است: \lr{Hoeffding} با $n_X/2$ مقیاس می‌گیرد در حالی که \lr{Chernoff} با $n_k$ مقیاس می‌گیرد. برای $n_k \ll n_X$، این تفاوت چشمگیر است و کران تنگ‌تر تولید می‌کند.

\textbf{رژیم 3 (شمارش صفر):} وقتی $n_k = 0$، از کران بالای \lr{Chernoff} ضربی استفاده می‌شود که از قانون سه برای مشاهدات صفر در توزیع \lr{Poisson}/\lr{Binomial} استخراج می‌شود:
\begin{equation}
	P(X = 0 \mid \mu) = e^{-\mu} \le \varepsilon \quad \Rightarrow \quad \mu \le \ln\!\left(\frac{1}{\varepsilon}\right),
\end{equation}
بنابراین کران بالای کران فیزیکی متناهی $\ln(1/\varepsilon_{\text{per-event}})$ است. این کران در مقایسه با \lr{Hoeffding} که برای $n_k = 0$ به $\delta_{\text{paper}}$ می‌رسد (که می‌تواند بزرگ باشد)، بسیار کوچکتر و فیزیکی‌تر است. برای $\varepsilon_{\text{sec}} = 10^{-9}$، این کران برابر $\ln(21\times10^9) \approx 23.6$ است که یک کران فوق‌العاده تنگ‌تر از $\delta_{\text{paper}} \approx 141$ در همان شرایط است.

اثبات اینکه این استراتژی هیبرید هرگز امنیت را تضعیف نمی‌کند، ساده است: هر کران تمرکزی که پوشش $1 - \varepsilon$ ارائه می‌کند، در چارچوب امنیتی معتبر است. از آنجا که \lr{Chernoff} برای شمارش‌های کوچک تنگ‌تر از \lr{Hoeffding} است (یعنی کران پایین بزرگتر و کران بالا کوچکتر تولید می‌کند، که هر دو محافظه‌کارانه‌تر هستند)، نتیجه هیبرید همیشه حداقل به اندازه \lr{Hoeffding} محافظه‌کارانه است. در نتیجه، طول کلید نهایی همیشه کمتر یا مساوی طول کلید \lr{Hoeffding} خالص است، که این دقیقاً همان جهت مطلوب برای امنیت است.

\subsection{فرمول تحلیلی کران پایین بازده خلأ}

معادله (2) مقاله، کران پایین تعداد رویدادهای خلأ در پایه $X$ را به صورت زیر ارائه می‌دهد:
\begin{equation}
	s_{X,0} \;\ge\; \tau_0\,\frac{\mu_2\,n_{X,\mu_3}^{-} - \mu_3\,n_{X,\mu_2}^{+}}{\mu_2 - \mu_3}.
\end{equation}
در این فرمول، $n_{X,\mu_3}^{-}$ کران پایین آشکارسازی در شدت خلأ است که با کمترین مقدار ممکن، صورت کسر را کوچک می‌کند و در نتیجه کران پایین محافظه‌کارانه می‌شود. به طور متقابل، $n_{X,\mu_2}^{+}$ کران بالای آشکارسازی در شدت \lr{decoy} ضعیف است که با بزرگ‌ترین مقدار ممکن، جمله کم‌شونده را بزرگ می‌کند و باز هم کران را سخت‌گیرانه می‌سازد. این انتخاب‌های جهت‌دار تضمین می‌کنند که تخمین نهایی، در بدترین حالت نوسان آماری، همچنان یک کران پایین معتبر باشد. در تابع \lr{\texttt{\_calc\_s0\_lower}} این فرمول پیاده‌سازی شده و در صورت منفی شدن صورت کسر، خروجی به صفر برده می‌شود. در نسخه فعلی، با معرفی استراتژی هیبرید، کران‌های $n_{X,\mu_3}^{-}$ و $n_{X,\mu_2}^{+}$ به‌صورت خودکار در رژیم مناسب (بزرگ، کوچک یا صفر) محاسبه می‌شوند، که به‌طور خاص برای $n_{X,\mu_3} = 0$ (که در فواصل طولانی رایج است) یک کران پایین معنادار به جای یک مقدار منفی تولید می‌کند.

\subsection{فرمول تحلیلی کران پایین بازده تک‌فوتونی}

معادله (3) مقاله، کران پایین تعداد رویدادهای تک‌فوتونی در پایه $X$ را به صورت زیر ارائه می‌دهد:
\begin{equation}
	s_{X,1} \;\ge\; \tau_1\,\mu_1\,\frac{T}{D},
\end{equation}
که در آن:
\begin{align}
	T &= n_{X,\mu_2}^{-} - n_{X,\mu_3}^{+} - \frac{\mu_2^2 - \mu_3^2}{\mu_1^2}\left(n_{X,\mu_1}^{+} - \frac{s_{X,0}}{\tau_0}\right), \\
	D &= \mu_1(\mu_2 - \mu_3) - \mu_2^2 + \mu_3^2.
\end{align}
مخرج $D$ با شرط $\mu_1 > \mu_2 + \mu_3$ مثبت می‌ماند. این فرمول پیچیده‌تر از معادله (2) است زیرا از تخمین بازده خلأ $s_{X,0}$ نیز استفاده می‌کند تا اثر فوتون‌های چندتایی حذف شود. در تابع \lr{\texttt{\_calc\_s1\_lower}} این فرمول با دریافت $\tau_1$، $\mu_1$، $\mu_2$، $\mu_3$، کران‌های \lr{Hoeffding}/\lr{Chernoff} برای هر سه شدت و نسبت $s_{X,0}/\tau_0$ پیاده‌سازی شده است. همانند معادله (2)، انتخاب جهت‌دار کران‌ها (کران پایین برای $\mu_2$ و کران بالا برای $\mu_3$ و $\mu_1$) تضمین می‌کند که نتیجه یک کران پایین معتبر باشد. در صورت منفی شدن صورت $T$، خروجی به صفر برده می‌شود زیرا تعداد رویدادهای تک‌فوتونی نمی‌تواند منفی باشد.

\subsection{فرمول تحلیلی کران بالای خطای تک‌فوتونی}

معادله (4) مقاله، کران بالای تعداد خطاهای تک‌فوتونی در پایه $Z$ را به صورت زیر ارائه می‌دهد:
\begin{equation}
	v_{Z,1} \;\le\; \tau_1\,\frac{m_{Z,\mu_2}^{+} - m_{Z,\mu_3}^{-}}{\mu_2 - \mu_3},
\end{equation}
که در آن $m_{Z,k}^{\pm}$ کران‌های \lr{Hoeffding}/\lr{Chernoff} برای شمارش خطاها در پایه $Z$ با شدت $\mu_k$ هستند. در اینجا از کران بالای خطاهای \lr{decoy} ضعیف و کران پایین خطاهای خلأ استفاده می‌شود تا کران بالای محافظه‌کارانه به دست آید. تابع \lr{\texttt{\_calc\_v1\_upper}} این محاسبه را پیاده‌سازی می‌کند و خروجی را به یک کران فیزیکی (مثلاً $n_Z$) محدود می‌سازد. توجه شود که این کران بالا برای خطاها است، بنابراین در جهت مخلاف کران‌های پایین عمل می‌کند: بزرگ‌ترین مقدار ممکن برای $m_{Z,\mu_2}^{+}$ و کوچک‌ترین مقدار ممکن برای $m_{Z,\mu_3}^{-}$ انتخاب می‌شوند تا صورت کسر بزرگ و در نتیجه کران بالا محافظه‌کارانه شود.

\subsection{کران بالای خطای فاز با نوسان نمونه‌گیری تصادفی}

معادله (5) مقاله، کران بالای خطای فاز $\varphi_X$ را با در نظر گرفتن نوسان نمونه‌گیری تصادفی (\lr{random-sampling fluctuation}) ارائه می‌دهد:
\begin{equation}
	\varphi_X = \frac{c_{X,1}}{s_{X,1}} \;\le\; \frac{v_{Z,1}}{s_{Z,1}} + \gamma(a, b, c, d),
\end{equation}
که در آن:
\begin{align}
	a &= \frac{\varepsilon_{\text{sec}}}{21}, \\
	b &= \frac{v_{Z,1}}{s_{Z,1}}, \\
	c &= s_{Z,1}, \\
	d &= s_{X,1}, \\
	\gamma(a, b, c, d) &= \sqrt{\frac{(c+d)(1-b)\,b}{c\,d\,\ln 2}\,\log_2\!\left(\frac{c+d}{c\,d\,(1-b)\,b\,a^2}\right)}.
\end{align}
این کران یکی از پیچیده‌ترین بخش‌های چارچوب \lr{Lim 2014} است و در تابع \lr{\texttt{\_calc\_phi\_X\_upper}} پیاده‌سازی شده است. در حالت‌های تباهی ($b \in \{0, 1\}$ یا $c \le 0$ یا $d \le 0$)، تابع به صورت ایمن به مقدار بدترین حالت $0.5$ بازمی‌گردد که با کاغذ همخوان است. ریشه دوم در فرمول اصلی به صراحت در مقاله آمده است (اگرچه چیدمان تایپی ممکن است آن را پنهان کند) و در کد با \lr{\texttt{math.sqrt}} پیاده‌سازی شده است. مدیریت دقیق این حالت‌های تباهی در ادامه در بخش پیاده‌سازی توضیح داده می‌شود.

\subsection{طول کلید امن و نشت تصحیح خطا}

معادله (1) مقاله، طول کلید امن نهایی را به صورت زیر ارائه می‌دهد:
\begin{equation}
	\ell = s_{X,0} + s_{X,1}\bigl(1 - h(\varphi_X)\bigr) - \lambda_{\text{EC}} - 6\log_2\!\left(\frac{21}{\varepsilon_{\text{sec}}}\right) - \log_2\!\left(\frac{2}{\varepsilon_{\text{cor}}}\right),
\end{equation}
که در آن $h(p) = -p\log_2 p - (1-p)\log_2(1-p)$ آنتروپی باینری است. نشت تصحیح خطا $\lambda_{\text{EC}}$ در حالت استاندارد به صورت زیر محاسبه می‌شود:
\begin{equation}
	\lambda_{\text{EC}} = f_{\text{EC}}\,n_X\,h(e_{\text{obs}}),
\end{equation}
که در آن $f_{\text{EC}} \ge 1$ کارایی تصحیح خطا و $e_{\text{obs}} = m_X / n_X$ خطای بیت مشاهده‌شده در پایه $X$ است. در مدل \lr{tomamichel} (اختیاری، بر اساس مرجع [37] مقاله یعنی \lr{Tomamichel et al., arXiv:1401.5194})، جمله $\log_2(2/\varepsilon_{\text{cor}})$ به $\lambda_{\text{EC}}$ اضافه می‌شود تا یک مدل نشت دقیق‌تر ارائه گردد و در نتیجه جمله جداگانه صحت صفر می‌شود تا از شمارش مضاعف جلوگیری شود. جمله $6\log_2(21/\varepsilon_{\text{sec}})$ شش رویداد شکست مرتبط با تقویت حریم خصوصی (\lr{privacy amplification}) را پوشش می‌دهد و در گزارش \lr{diagnostics} به‌صورت \lr{\texttt{pa\_term}} ذخیره می‌شود.

\subsection{قید فیزیکی دفاعی روی طول کلید}

یکی از additions مهم در این نسخه از پیاده‌سازی، یک لایه دفاعی فیزیکی در متد \lr{\texttt{calculate\_key\_length}} است که پس از محاسبه نهایی طول کلید اعمال می‌شود. این قید بر اساس یک اصل بدیهی فیزیکی است: \textbf{بیت‌های کلید امن زیرمجموعه‌ای از بیت‌های خام \lr{sifted} هستند}، پس طول کلید نمی‌تواند از مجموع $n_X + n_Z$ (یعنی کل بیت‌های \lr{sifted}) بیشتر شود. این قید در شرایط عادی هرگز فعال نمی‌شود زیرا فرمول‌های \lr{Lim 2014} به‌طور ذاتی محافظه‌کارانه هستند، اما به‌عنوان یک تور ایمنی در برابر باگ‌های داخلی طراحی شده است.

سه سناریویی که در آن‌ها این قید ممکن است فعال شود عبارتند از: (1) حلقه \lr{kappa-solver} تکراری که به‌دلیل عدم همگرایی، یک $\varepsilon_{\text{sec}}$ نامعتبر تولید کرده و در نتیجه $\ell$ بزرگ می‌شود؛ (2) \lr{fallback} در مسیر \lr{LP} که ممکن است مقادیر غیرصفر برای کران پایین بازده تولید کند حتی وقتی آمار کافی نیست؛ (3) سرریز در تخمین \lr{decoy} به‌دلیل خطاهای عددی در محاسبات \lr{Chernoff}. در هر یک از این موارد، قید دفاعی طول کلید را به $n_X + n_Z$ محدود می‌کند، یک هشدار در \lr{log} ثبت می‌کند و نتیجه بازنویسی‌شده را بازمی‌گرداند. این طراحی در فلسفه \lr{fail-safe} عمل می‌کند: حتی در صورت بروز باگ داخلی، خروجی هرگز از نظر فیزیکی نامعتبر نخواهد بود.

\section{تحلیل معماری و ساختار داده‌ها}

\subsection{ثابت‌های ماژول و کنترل عددی}

فایل با تعریف چند ثابت کلیدی آغاز می‌شود که رفتار عددی و امنیتی پیاده‌سازی را کنترل می‌کنند:

\begin{LTR}
	\begin{lstlisting}[language=Python]
		MAX_ERROR_RATE = 0.5
		HOEFFDING_FACTOR = 10.0
		PROB_SUM_TOL = 1e-9
		MIN_EPS = 1e-300
		LIM2014_FAILURE_EVENTS_3_INTENSITY = 21.0
		DEFAULT_KAPPA_MAX_ITERATIONS = 50
		DEFAULT_KAPPA_ABS_TOL_BITS = 100.0
		DEFAULT_KAPPA_REL_TOL = 1e-6
		MAX_SECURITY_EPS = 1.0 - 1e-15
	\end{lstlisting}
\end{LTR}

ثابت \lr{\texttt{MAX\_ERROR\_RATE = 0.5}} حد بالای فیزیکی خطای بیت است که در حالت بدترین حالت (یعنی زمانی که هیچ اطلاعاتی از کانال استخراج نمی‌شود) به کار می‌رود. این مقدار برای \lr{clamp} کردن خطای فاز و خطای مشاهده‌شده به بازه $[0, 0.5]$ استفاده می‌شود. \lr{\texttt{HOEFFDING\_FACTOR = 10.0}} ثابت جدیدی است که آستانه انتخاب رژیم در استراتژی هیبرید را تعیین می‌کند: وقتی $n_k \ge \text{HOEFFDING\_FACTOR} \times \delta_{\text{paper}}$، از \lr{Hoeffding} استفاده می‌شود. مقدار 10 به این معناست که وقتی خطای آماری نسبی $\delta_{\text{paper}}/n_k \le 10\%$ است، \lr{Hoeffding} به‌طور معقولی تنگ عمل می‌کند و نیازی به \lr{Chernoff} نیست. \lr{\texttt{PROB\_SUM\_TOL = 1e-9}} تلورانس قابل قبول برای مجموع احتمال‌های شدت نور است. \lr{\texttt{MIN\_EPS = 1e-300}} کف بسیار کوچکی برای $\varepsilon$ است که از خطاهای عددی در محاسبات \lr{log} جلوگیری می‌کند.

\lr{\texttt{LIM2014\_FAILURE\_EVENTS\_3\_INTENSITY = 21.0}} تعداد رویدادهای شکست در طرح تخصیص $\varepsilon$ است که در ضمیمه B مقاله توضیح داده شده است. سه ثابت \lr{\texttt{DEFAULT\_KAPPA\_MAX\_ITERATIONS}}، \lr{\texttt{DEFAULT\_KAPPA\_ABS\_TOL\_BITS}} و \lr{\texttt{DEFAULT\_KAPPA\_REL\_TOL}} برای حل تکراری معادله $\varepsilon_{\text{sec}} = \kappa \cdot \ell$ استفاده می‌شوند که در بخش \S IV مقاله معرفی شده است. در این روش، $\varepsilon_{\text{sec}}$ به طول کلید متناسب می‌شود و باید با تکرار حل شود زیرا $\ell$ خود به $\varepsilon_{\text{sec}}$ وابسته است. حداکثر 50 تکرار با تلورانس مطلق 100 بیت و تلورانس نسبی $10^{-6}$ به عنوان معیار همگرایی در نظر گرفته شده است. \lr{\texttt{MAX\_SECURITY\_EPS = 1.0 - 1e-15}} کران بالای $\varepsilon_{\text{sec}}$ است که از تنظیم آن به 1 (یعنی عدم امنیت کامل) جلوگیری می‌کند.

\subsection{\lr{DecoyEstimatesLim2014}}

این \lr{dataclass} یک ساختار داده تخصصی برای نگهداری تخمین‌های میانی چارچوب \lr{Lim 2014} است:

\begin{LTR}
	\begin{lstlisting}[language=Python]
		@dataclass(frozen=True)
		class DecoyEstimatesLim2014(DecoyEstimates):
		"""
		Intermediate decoy estimates for the Lim-2014 calculation.
		"""
		
		s_X_0_lower: float = 0.0
		s_X_1_lower: float = 0.0
		s_Z_1_lower: float = 0.0
		v_Z_1_upper: float = 0.0
		phi_X_upper: float = MAX_ERROR_RATE
		is_feasible: bool = False
		failure_prob_used: float = 1.0
		diagnostics: Dict[str, Any] = field(default_factory=dict)
		
		def to_dict(self) -> Dict[str, Any]:
		return _json_safe_dict(self)
	\end{lstlisting}
\end{LTR}

این کلاس از \lr{\texttt{DecoyEstimates}} (تعریف‌شده در \lr{\texttt{base.py}}) ارث می‌برد و فیلدهای اضافی اختصاصی چارچوب \lr{Lim 2014} را اضافه می‌کند. پنج فیلد اصلی $s_{X,0}^L$، $s_{X,1}^L$، $s_{Z,1}^L$، $v_{Z,1}^U$ و $\varphi_X^U$ هستند که مستقیماً از معادلات (2) تا (5) مقاله استخراج می‌شوند. \lr{\texttt{is\_feasible}} نشان می‌دهد که آیا تخمین معتبر بوده و طول کلید مثبت تولید شده است یا خیر. \lr{\texttt{failure\_prob\_used}} احتمال شکستی است که در محاسبه کران‌ها استفاده شده (معمولاً $\varepsilon_{\text{sec}}$ یا $\varepsilon_{\text{sec}}/21$). \lr{\texttt{diagnostics}} یک دیکشنری برای اطلاعات تشخیصی کامل است. استفاده از \lr{\texttt{frozen=True}} اطمینان می‌دهد که پس از ساخت، مقادیر قابل تغییر نیستند که این ویژگی برای صحت‌سنجی نتایج امنیتی حیاتی است. متد \lr{\texttt{to\_dict}} برای سریال‌سازی به \lr{JSON} با تبدیل نوع‌های \lr{numpy} و \lr{Enum} به نوع‌های بومی پایتون استفاده می‌شود.

\subsection{\lr{Lim2014Proof} و سازنده آن}

کلاس اصلی ماژول، \lr{\texttt{Lim2014Proof}} است که به صورت \lr{dataclass} و \lr{frozen} تعریف شده و از \lr{\texttt{FiniteKeyProof}} ارث می‌برد. سازنده این کلاس در نسخه فعلی دارای یک منطق جدید برای اعمال روش بازه اطمینان است:

\begin{LTR}
	\begin{lstlisting}[language=Python]
		@dataclass(frozen=True)
		class Lim2014Proof(FiniteKeyProof):
		"""
		Lim et al. 2014 finite-key proof implementation.
		
		The three-intensity analytical path implements Eqs. (1)-(5) of the paper
		exactly. For more than three intensities, a conservative LP-based
		N-decoy estimate is used (see module docstring).
		"""
		
		p: "QKDParams"  # type: ignore[assignment]
		
		def __init__(self, p: "QKDParams"):
		import logging
		from qkd.proofs.base import ENTROPY_PROB_CLAMP, ProofMode
		try:
		import mpmath
		has_mpmath = True
		except ImportError:
		has_mpmath = False
		
		object.__setattr__(self, 'p', p)
		object.__setattr__(self, 'mode', ProofMode.PRODUCTION)
		object.__setattr__(self, 'logger', logging.getLogger(self.__class__.__name__))
		object.__setattr__(self, 'run_id', getattr(p, 'run_id', 'no-run-id'))
		object.__setattr__(self, 'min_prob_clamp', getattr(p, 'entropy_clamp', ENTROPY_PROB_CLAMP))
		object.__setattr__(self, 'use_mpmath', getattr(p, 'use_mpmath', False) and has_mpmath)
		if getattr(p, 'ci_method', None) is not None:
		from ..datatypes import ConfidenceBoundMethod
		if p.ci_method != ConfidenceBoundMethod.CHERNOFF:
		import logging
		logging.getLogger(self.__class__.__name__).warning(
		f"Lim2014 proof requires CHERNOFF bounds for finite-key security. "
		f"Overriding '{p.ci_method.name}' with CHERNOFF."
		)
		object.__setattr__(p, 'ci_method', ConfidenceBoundMethod.CHERNOFF)
	\end{lstlisting}
\end{LTR}

از آنجا که کلاس \lr{\texttt{@dataclass(frozen=True)}} است، تخصیص مستقیم ویژگی‌ها در \lr{\texttt{\_\_init\_\_}} مجاز نیست. به جای آن از \lr{\texttt{object.\_\_setattr\_\_}} برای دور زدن این محدودیت استفاده می‌شود. این الگو به ما اجازه می‌دهد کلاس را به صورت \lr{immutable} نگه داریم در حالی که هنوز بتوانیم در زمان ساخت، وابستگی‌ها را تزریق کنیم. سپس \lr{\texttt{mpmath}} به صورت اختیاری وارد می‌شود و اگر موجود نباشد، پرچم \lr{\texttt{use\_mpmath}} به \lr{\texttt{False}} تنظیم می‌شود. این طراحی امکان استفاده در محیط‌هایی بدون \lr{mpmath} را فراهم می‌کند.

بخش جدید و حیاتی سازنده، بررسی \lr{\texttt{p.ci\_method}} است. امنیت کلید متناهی در چارچوب \lr{Lim 2014} به‌طور صریح به کران‌های تمرکز \lr{Chernoff} وابسته است (هم در رژیم \lr{Hoeffding} و هم در رژیم \lr{Chernoff} هیبرید، زیرا \lr{Hoeffding} یک حالت خاص از \lr{Chernoff} است). اگر کاربر روش دیگری مانند \lr{\texttt{CONFIDENCE\_INTERVAL}} یا \lr{\texttt{BOOTSTRAP}} را تنظیم کرده باشد، سازنده یک هشدار ثبت می‌کند و روش را به \lr{\texttt{CHERNOFF}} بازنویسی می‌کند. این \lr{override} تهاهماهنگ با فلسفه \lr{fail-safe} است: به جای ارائه یک نتیجه ناامن، پیاده‌سازی به‌طور خودکار به تنظیمات امن تغییر می‌یابد و کاربر را مطلع می‌سازد. توجه شود که این بازنویسی روی خود شیء \lr{\texttt{p}} انجام می‌شود (\lr{\texttt{object.\_\_setattr\_\_(p, ...)}})، که یک \lr{side-effect} است؛ این انتخاب آگاهانه است زیرا در غیر این صورت، تنظیمات پیکربندی کاربر در طول خط لوله شبیه‌سازی ناسازگار می‌ماند.

\subsection{متدهای استاتیک ایمن عددی}

کلاس چهار متد استاتیک برای محاسبات ایمن عددی دارد. متد \lr{\texttt{\_validate\_epsilon\_value}} اعتبارسنجی $\varepsilon$ را انجام می‌دهد:

\begin{LTR}
	\begin{lstlisting}[language=Python]
		@staticmethod
		def _validate_epsilon_value(value: float, name: str) -> None:
		if not math.isfinite(float(value)) or not (0.0 < float(value) < 1.0):
		raise ParameterValidationError(f"{name} must be finite and in the open interval (0, 1).")
	\end{lstlisting}
\end{LTR}

این متد بررسی می‌کند که مقدار متناهی باشد و در بازه باز $(0, 1)$ قرار داشته باشد. بازه باز به این معناست که $\varepsilon = 0$ (که منجر به تقسیم بر صفر در فرمول‌های \lr{log} می‌شود) و $\varepsilon = 1$ (که به معنای عدم امنیت است) رد می‌شوند. این اعتبارسنجی برای همه $\varepsilon$های ورودی به کار می‌رود: $\varepsilon_{\text{sec}}$، $\varepsilon_{\text{cor}}$ و $\varepsilon$ تک‌آزمون. در مسیر \lr{kappa-solver} تکراری نیز برای هر $\varepsilon_{\text{target}}$ جدید اعمال می‌شود تا از مقادیر نامعتبر جلوگیری شود.

متد \lr{\texttt{\_safe\_log}} و \lr{\texttt{\_safe\_log2}} لگاریتم‌های ایمن هستند:

\begin{LTR}
	\begin{lstlisting}[language=Python]
		@staticmethod
		def _safe_log(x: float) -> float:
		if not math.isfinite(float(x)) or x <= 0.0:
		raise ParameterValidationError(f"log argument must be finite and positive; got {x!r}.")
		return math.log(x)
		
		@staticmethod
		def _safe_log2(x: float) -> float:
		if not math.isfinite(float(x)) or x <= 0.0:
		raise ParameterValidationError(f"log2 argument must be finite and positive; got {x!r}.")
		return math.log2(x)
	\end{lstlisting}
\end{LTR}

برخلاف متد \lr{\texttt{\_safe\_log}} در کلاس پایه که مقدار بسیار کوچک را به $10^{-300}$ \lr{clamp} می‌کند، این پیاده‌سازی در صورت دریافت ورودی نامعتبر، استثنا پرتاب می‌کند. این رفتار سخت‌گیرانه‌تر به این دلیل است که در چارچوب \lr{Lim 2014}، هر ورودی غیرمثبت به لگاریتم نشانگر یک خطای منطقی جدی است که باید ریشه‌یابی شود، نه یک خطای عددی قابل چشم‌پوشی. متد \lr{\texttt{\_clamp\_probability}} احتمالات را به بازه $[0, u]$ محدود می‌کند که در آن $u$ به صورت پیش‌فرض 1 است اما می‌تواند به $0.5$ (برای خطاها) تنظیم شود:

\begin{LTR}
	\begin{lstlisting}[language=Python]
		@staticmethod
		def _clamp_probability(x: float, *, upper: float = 1.0) -> float:
		if not math.isfinite(float(x)):
		return upper
		return min(max(float(x), 0.0), upper)
	\end{lstlisting}
\end{LTR}

اگر ورودی نامتناهی باشد، به کران بالای بازه بازمی‌گردد که در حالت خطا به معنای بدترین حالت (یعنی عدم امنیت) است.

\subsection{اعتبارسنجی \lr{TallyCounts}}

متد \lr{\texttt{\_validate\_tally\_counts}} یک لایه دفاعی برای صحت‌سنجی شمارش‌های آماری ورودی است:

\begin{LTR}
	\begin{lstlisting}[language=Python]
		def _validate_tally_counts(self, stats: TallyCounts, pulse_name: str) -> None:
		required = ("sent", "sifted_x", "sifted_z", "errors_sifted_x", "errors_sifted_z")
		missing = [name for name in required if not hasattr(stats, name)]
		if missing:
		raise ParameterValidationError(
		f"TallyCounts for pulse {pulse_name!r} is missing required fields: {missing}."
		)
		
		sent = float(stats.sent)
		sifted_x = float(stats.sifted_x)
		sifted_z = float(stats.sifted_z)
		err_x = float(stats.errors_sifted_x)
		err_z = float(stats.errors_sifted_z)
		
		values = {
			"sent": sent,
			"sifted_x": sifted_x,
			"sifted_z": sifted_z,
			"errors_sifted_x": err_x,
			"errors_sifted_z": err_z,
		}
		
		for name, value in values.items():
		if not math.isfinite(value) or value < 0.0:
		raise ParameterValidationError(
		f"{name} for pulse {pulse_name!r} must be finite and non-negative."
		)
		
		if err_x > sifted_x:
		raise ParameterValidationError(
		f"errors_sifted_x exceeds sifted_x for pulse {pulse_name!r}: {err_x} > {sifted_x}."
		)
		if err_z > sifted_z:
		raise ParameterValidationError(
		f"errors_sifted_z exceeds sifted_z for pulse {pulse_name!r}: {err_z} > {sifted_z}."
		)
		if sifted_x + sifted_z > sent + 1e-9:
		raise ParameterValidationError(
		f"sifted_x + sifted_z exceeds sent for pulse {pulse_name!r}."
		)
	\end{lstlisting}
\end{LTR}

این متد چهار سطح اعتبارسنجی انجام می‌دهد: (1) حضور تمام فیلدهای مورد نیاز \lr{\texttt{sent}}، \lr{\texttt{sifted\_x}}، \lr{\texttt{sifted\_z}}، \lr{\texttt{errors\_sifted\_x}} و \lr{\texttt{errors\_sifted\_z}} در شیء \lr{\texttt{TallyCounts}}؛ (2) متناهی و نامنفی بودن تمام مقادیر؛ (3) عدم تجاوز خطاها از شمارش متناظر ($m_X \le n_X$ و $m_Z \le n_Z$) که یک قید فیزیکی بدیهی است؛ (4) عدم تجاوز مجموع \lr{sifted} از \lr{sent} با تلورانس $10^{-9}$. این اعتبارسنجی‌ها قبل از هر محاسبه انجام می‌شوند تا از خطاهای دیرهنگام که ریشه‌یابی آن‌ها دشوار است، جلوگیری شود. در صورت شکست هر یک از این اعتبارسنجی‌ها، یک \lr{\texttt{ParameterValidationError}} با پیام مشخص پرتاب می‌شود که نام پالس و فیلد مشکل‌دار را شامل می‌شود.

\section{پیاده‌سازی ریاضی و الگوریتمی}

\subsection{اعتبارسنجی پیکربندی}

متد \lr{\texttt{\_validate\_configuration}} تمام پیش‌نیازهای چارچوب \lr{Lim 2014} را بررسی می‌کند:

\begin{LTR}
	\begin{lstlisting}[language=Python]
		def _validate_configuration(self) -> None:
		self._validate_epsilon_value(self.p.eps_sec, "eps_sec")
		self._validate_epsilon_value(self.p.eps_cor, "eps_cor")
		
		if not math.isfinite(float(self.p.f_error_correction)) or self.p.f_error_correction < 1.0:
		raise ParameterValidationError("f_error_correction must be finite and >= 1.")
		
		ec_model = getattr(self.p, "error_correction_model", "standard")
		if ec_model not in {"standard", "lim", "tomamichel"}:
		raise ParameterValidationError(
		"error_correction_model must be one of: 'standard', 'lim', 'tomamichel'."
		)
		
		photon_cap = int(getattr(self.p, "photon_number_cap", 0))
		if photon_cap < 1:
		raise ParameterValidationError("photon_number_cap must be at least 1.")
		
		pulse_configs = list(getattr(self.p.source, "pulse_configs", []))
		if len(pulse_configs) < 3:
		raise ConfigurationError("Lim2014Proof requires at least three pulse configurations.")
		
		names = [pc.name for pc in pulse_configs]
		if len(set(names)) != len(names):
		raise ConfigurationError("Pulse configuration names must be unique.")
		
		prob_sum = 0.0
		for pc in pulse_configs:
		mu = float(pc.mean_photon_number)
		prob = float(pc.probability)
		
		if not math.isfinite(mu) or mu < 0.0:
		raise ParameterValidationError(f"Invalid mean photon number for pulse {pc.name!r}: {mu!r}.")
		if not math.isfinite(prob) or prob <= 0.0:
		raise ParameterValidationError(f"Pulse probability for {pc.name!r} must be finite and positive.")
		prob_sum += prob
		
		if abs(prob_sum - 1.0) > PROB_SUM_TOL:
		raise ParameterValidationError(f"Pulse probabilities must sum to 1; got {prob_sum!r}.")
		
		kappa = getattr(self.p, "security_constant_kappa", None)
		if kappa is not None:
		if not math.isfinite(float(kappa)) or float(kappa) <= 0.0:
		raise ParameterValidationError("security_constant_kappa must be finite and positive when set.")
	\end{lstlisting}
\end{LTR}

این متد به ترتیب اعتبارسنجی می‌کند: (1) مقادیر $\varepsilon_{\text{sec}}$ و $\varepsilon_{\text{cor}}$ در بازه $(0, 1)$ باشند؛ (2) ضریب کارایی تصحیح خطا $f_{\text{EC}} \ge 1$ باشد (مقدار کمتر از 1 از نظر تئوری رمزنگاری غیرممکن است زیرا تصحیح خطا نمی‌تواند بهتر از حد \lr{Shannon} باشد)؛ (3) مدل تصحیح خطا یکی از سه حالت مجاز \lr{\texttt{standard}}، \lr{\texttt{lim}} یا \lr{\texttt{tomamichel}} باشد؛ (4) کسر تعداد فوتون \lr{\texttt{photon\_number\_cap}} حداقل 1 باشد (برای مسیر \lr{LP} ضروری است)؛ (5) حداقل سه پیکربندی شدت نور وجود داشته باشد؛ (6) نام پیکربندی‌ها یکتا باشند؛ (7) مجموع احتمال‌ها برابر 1 باشد با تلورانس \lr{\texttt{PROB\_SUM\_TOL}}؛ (8) در صورت تنظیم $\kappa$، مقدار آن مثبت و متناهی باشد. این اعتبارسنجی‌ها قبل از هر محاسبه انجام می‌شوند.

\subsection{محاسبه توزیع \lr{Poisson} و \texorpdfstring{$\tau_n$}{tau\_n}}

تابع \lr{\texttt{\_poisson\_pmf\_list}} توزیع \lr{Poisson} را تا حد $n_{\text{cap}}$ به صورت بازگشتی و پایدار عددی محاسبه می‌کند:

\begin{LTR}
	\begin{lstlisting}[language=Python]
		@staticmethod
		def _poisson_pmf_list(mu: float, n_cap: int) -> List[float]:
		if mu < 0.0 or not math.isfinite(mu):
		raise ParameterValidationError(f"Poisson mean must be finite and non-negative; got {mu!r}.")
		pmfs = [0.0] * (n_cap + 1)
		pmfs[0] = math.exp(-mu)
		for n in range(1, n_cap + 1):
		pmfs[n] = pmfs[n - 1] * mu / float(n)
		return pmfs
	\end{lstlisting}
\end{LTR}

این پیاده‌سازی از رابطه بازگشتی $P(n \mid \mu) = P(n-1 \mid \mu) \cdot \mu / n$ استفاده می‌کند که از محاسبه مستقیم فاکتوریل $n!$ که برای $n$ بزرگ سرریز می‌شود، اجتناب می‌کند. این روش پایداری عددی بسیار بالایی دارد و حتی برای $n = 100$ نیز نتایج دقیق تولید می‌کند. مقدار $P(0 \mid \mu) = e^{-\mu}$ به عنوان نقطه شروع بازگشت استفاده می‌شود.

تابع \lr{\texttt{\_get\_tau\_n}} مقدار $\tau_n$ را برای $n$ مشخص محاسبه می‌کند:

\begin{LTR}
	\begin{lstlisting}[language=Python]
		def _get_tau_n(self, n: int, pulse_map: Mapping[str, float], p_k: Mapping[str, float]) -> float:
		"""
		tau_n := sum_k p_k * e^{-mu_k} * mu_k^n / n!
		(paper, definition below Eq. (2)).
		"""
		if n < 0:
		raise ParameterValidationError("n must be non-negative.")
		
		tau_n = 0.0
		for name, mu in pulse_map.items():
		prob_k = p_k[name]
		pmf = math.exp(-mu) if n == 0 else math.exp(-mu) * (mu**n) / math.factorial(n)
		tau_n += prob_k * pmf
		return max(0.0, tau_n)
	\end{lstlisting}
\end{LTR}

این تابع دقیقاً تعریف $\tau_n = \sum_k p_k\,P(n \mid \mu_k)$ را پیاده‌سازی می‌کند. در حالت $n = 0$ برای جلوگیری از محاسبه $0^0$، از شاخه خاص \lr{\texttt{math.exp(-mu)}} استفاده می‌شود. تابع \lr{\texttt{\_get\_tau\_list}} همان محاسبه را به صورت برداری برای همه $n$ از 0 تا $n_{\text{cap}}$ انجام می‌دهد که در مسیر \lr{LP} استفاده می‌شود. تابع \lr{\texttt{\_get\_tail\_probability}} احتمال دم را محاسبه می‌کند، یعنی $1 - \sum_{n=0}^{n_{\text{cap}}} P(n \mid \mu)$ که در \lr{tail relaxation} مسیر \lr{LP} به کار می‌رود. این تابع برای مدیریت فوتون‌های چندتایی فراتر از کسر $n_{\text{cap}}$ ضروری است.

\subsection{کران‌های هیبرید \lr{Hoeffding/Chernoff} با دو پارامتر}

متد \lr{\texttt{\_get\_bounds}} قلب استراتژی هیبرید است و کران‌های بازه اطمینان را در سه رژیم پیاده‌سازی می‌کند:

\begin{LTR}
	\begin{lstlisting}[language=Python]
		def _get_bounds(self, n_k: float, n_total: float, eps_ln_term: float,
		clip_total: Optional[float] = None) -> Tuple[float, float]:
		"""
		Hybrid Hoeffding / per-intensity Chernoff confidence interval.
		
		Paper Lim2014 Appendix A uses Hoeffding's inequality with a single
		delta = sqrt(n_total/2 * ln(21/eps_sec)) that is the SAME for all
		intensities k. This is a sufficient condition for security but is
		pathologically loose when n_k << n_total (e.g. vacuum counts at long
		distance: observed 5, Hoeffding upper 572, 114x the observation).
		
		Hybrid strategy (security-preserving, tighter for small counts):
		- When n_k is "large" (n_k >= 10 * delta_paper, i.e. relative
		statistical error < 10%), use the paper's Hoeffding bound.
		- When n_k is "small" but nonzero (0 < n_k < 10 * delta_paper),
		use per-intensity multiplicative Chernoff:
		delta_k = sqrt(2 * n_k * ln(1/eps_per_event))
		- When n_k = 0, the upper bound is ln(1/eps_per_event) (rule of
		three for Poisson/binomial zero observations).
		
		# HYBRID-HOEFFDING-CHERNOFF v1
		"""
		if clip_total is None:
		clip_total = n_total
		
		if n_total < 0.0:
		raise ParameterValidationError("n_total must be non-negative.")
		if clip_total < 0.0:
		raise ParameterValidationError("clip_total must be non-negative.")
		if eps_ln_term < 0.0 or not math.isfinite(eps_ln_term):
		raise ParameterValidationError("eps_ln_term must be finite and non-negative.")
		if n_k < 0.0:
		raise ParameterValidationError("n_k must be non-negative.")
		
		# --- Hoeffding delta (paper Eq. A1): same for all k, scales with n_total.
		delta_paper = math.sqrt((n_total / 2.0) * eps_ln_term)
		
		# --- Hybrid selection ---
		hoeffding_threshold = HOEFFDING_FACTOR * delta_paper
		
		if n_k >= hoeffding_threshold and n_k > 0.0:
		# Paper-faithful Hoeffding regime (large counts).
		lower = max(0.0, n_k - delta_paper)
		upper = min(n_k + delta_paper, clip_total)
		elif n_k > 0.0:
		# Per-intensity multiplicative Chernoff regime (small but nonzero).
		delta_k = math.sqrt(2.0 * n_k * eps_ln_term)
		lower = max(0.0, n_k - delta_k)
		upper = min(n_k + delta_k, clip_total)
		else:
		# n_k = 0: multiplicative Chernoff upper bound (rule of three).
		lower = 0.0
		upper = min(eps_ln_term, clip_total)
		
		return lower, upper
	\end{lstlisting}
\end{LTR}

پارامتر \lr{\texttt{n\_total}} مقیاس محاسبه $\delta_{\text{paper}}$ را تعیین می‌کند و \lr{\texttt{clip\_total}} کران بالای فیزیکی برای \lr{clip} کردن است. در پیاده‌سازی استاندارد مقاله، هر دو برابر $n_X$ (یا $n_Z$) هستند؛ اما در این پیاده‌سازی، \lr{\texttt{n\_total}} می‌تواند برابر $n_{X,k}$ (شمارش خود شدت) باشد تا $\delta$ کوچکتر و تنگ‌تر شود. این امر از یک مشکل عملی جلوگیری می‌کند: وقتی شمارش \lr{decoy} یا خلأ بسیار کم است، ضریب پیشانی $e^{\mu_k}/p_k$ می‌تواند $\delta$ مبتنی بر $n_X$ را به قدری بزرگ کند که کران‌های پایین منفی و کران‌های بالا بیش از کل فیزیکی شوند.

نکته کلیدی در این پیاده‌سازی، استفاده از \lr{\texttt{eps\_ln\_term}} به‌عنوان آرگومان ورودی است. این متغیر برابر $\ln(21/\varepsilon_{\text{sec}}) = \ln(1/\varepsilon_{\text{per-event}})$ است (زیرا $\varepsilon_{\text{per-event}} = \varepsilon_{\text{sec}}/21$). این یکپارچه‌سازی هوشمندانه امکان استفاده از همان جمله لگاریتمی در هر دو رژیم \lr{Hoeffding} و \lr{Chernoff} را فراهم می‌کند بدون آنکه نیاز به محاسبه جداگانه باشد. در رژیم صفر، کران بالای $\ln(1/\varepsilon_{\text{per-event}})$ مستقیماً برابر \lr{\texttt{eps\_ln\_term}} است که یک کران فیزیکی متناهی و معقول ارائه می‌دهد. در نتیجه، این متد هیچ‌گاه یک کران بالای نامتناهی یا غیرفیزیکی تولید نمی‌کند.

\subsection{مسیر \lr{N-decoy} مبتنی بر \lr{LP}}

وقتی تعداد شدت‌ها بیش از 3 باشد، پیاده‌سازی به مسیر \lr{LP} تغییر مسیر می‌دهد. متد \lr{\texttt{\_solve\_n\_decoy\_lp}} یک برنامه‌ریزی خطی محافظه‌کارانه با محدودیت‌های غیرمنفی‌بودن، کران بالای صریح و \lr{tail relaxation} را حل می‌کند. این متد سه نوع هدف را پشتیبانی می‌کند: \lr{\texttt{s0\_lower}} (کران پایین بازده خلأ)، \lr{\texttt{s1\_lower}} (کران پایین بازده تک‌فوتونی) و \lr{\texttt{v1\_upper}} (کران بالای خطای تک‌فوتونی). در حالت \lr{\texttt{v1\_upper}}، تابع هدف به $-v_1$ تبدیل می‌شود زیرا \lr{LP} استاندارد مسأله کمینه‌سازی را حل می‌کند و برای یافتن کران بالا باید منفی متغیر را کمینه کرد.

این مسیر با فراخوانی \lr{\texttt{solve\_lp}} از ماژول \lr{\texttt{qkd.proofs.utils\_lp}} که به \lr{scipy.optimize.linprog} یا \lr{cvxopt} متصل است، حل می‌شود. در صورت شکست \lr{LP} (مثلاً به‌دلیل عدم \lr{feasibility})، یک \lr{fallback} محافظه‌کارانه اعمال می‌شود: کران‌های پایین به صفر و کران‌های بالا به حد فیزیکی بازمی‌گردند. این رفتار \lr{fail-safe} تضمین می‌کند که حتی در صورت شکست \lr{solver}، خروجی هرگز خوش‌بینانه نخواهد بود. توجه شود که این مسیر در گزارش \lr{diagnostics} به‌روشنی به‌عنوان توسعه فراتر از نتایج بسته مقاله علامت‌گذاری شده است.

\subsection{پیاده‌سازی معادله (2): کران پایین بازده خلأ}

تابع \lr{\texttt{\_calc\_s0\_lower}} معادله (2) مقاله را پیاده‌سازی می‌کند:

\begin{LTR}
	\begin{lstlisting}[language=Python]
		def _calc_s0_lower(
		self,
		tau_0: float,
		mu2: float,
		mu3: float,
		n_mu2_plus: float,
		n_mu3_minus: float,
		) -> float:
		"""
		Eq. (2):
		s_{X,0} >= tau_0 * (mu2 * n^-_{X,mu3} - mu3 * n^+_{X,mu2}) / (mu2 - mu3)
		"""
		denominator = mu2 - mu3
		if denominator <= 0.0:
		raise ConfigurationError("Invalid s0 denominator; require mu2 > mu3.")
		
		numerator = tau_0 * (mu2 * n_mu3_minus - mu3 * n_mu2_plus)
		return max(0.0, numerator / denominator)
	\end{lstlisting}
\end{LTR}

پارامترهای ورودی به ترتیب $\tau_0$، $\mu_2$، $\mu_3$، $n_{X,\mu_2}^{+}$ (کران بالای آشکارسازی \lr{decoy}) و $n_{X,\mu_3}^{-}$ (کران پایین آشکارسازی خلأ) هستند. توجه شود که در صورت کسر، $\mu_2$ با کران پایین خلأ و $\mu_3$ با کران بالای \lr{decoy} ضرب می‌شوند تا صورت کسر به کوچکترین مقدار ممکن برسد و در نتیجه کران پایین محافظه‌کارانه به دست آید. در صورت منفی شدن صورت (که ممکن است در شرایط نوسان شدید آماری رخ دهد)، خروجی با \lr{\texttt{max(0.0, ...)}} به صفر برده می‌شود زیرا تعداد رویدادهای خلأ نمی‌تواند منفی باشد. نکته مهم این است که در نسخه فعلی با استراتژی هیبرید، $n_{X,\mu_3}^{-}$ دیگر هرگز به مقدار منفی \lr{Hoeffding} منجر نمی‌شود، زیرا در رژیم $n_k = 0$، کران پایین صفر و کران بالا $\ln(1/\varepsilon)$ است.

\subsection{پیاده‌سازی معادله (3): کران پایین بازده تک‌فوتونی}

تابع \lr{\texttt{\_calc\_s1\_lower}} معادله (3) مقاله را پیاده‌سازی می‌کند:

\begin{LTR}
	\begin{lstlisting}[language=Python]
		def _calc_s1_lower(
		self,
		tau_1: float,
		mu1: float,
		mu2: float,
		mu3: float,
		n_mu1_plus: float,
		n_mu2_minus: float,
		n_mu3_plus: float,
		s_0_lower_over_tau_0: float,
		) -> float:
		"""
		Eq. (3):
		s_{X,1} >= tau_1 * mu1 * T / D
		where
		T = n^-_{X,mu2} - n^+_{X,mu3} - ((mu2^2 - mu3^2)/mu1^2) * (n^+_{X,mu1} - s_{X,0}/tau_0)
		D = mu1*(mu2 - mu3) - mu2^2 + mu3^2
		"""
		mu1_sq = mu1 * mu1
		mu2_sq = mu2 * mu2
		mu3_sq = mu3 * mu3
		
		denominator = mu1 * (mu2 - mu3) - mu2_sq + mu3_sq
		if denominator <= 0.0:
		raise ConfigurationError("Invalid s1 denominator; source intensities do not satisfy formula conditions.")
		
		term = n_mu2_minus - n_mu3_plus - ((mu2_sq - mu3_sq) / mu1_sq) * (
		n_mu1_plus - s_0_lower_over_tau_0
		)
		numerator = tau_1 * mu1 * term
		return max(0.0, numerator / denominator)
	\end{lstlisting}
\end{LTR}

این تابع پیچیده‌ترین فرمول تحلیلی چارچوب است. مخرج $D = \mu_1(\mu_2 - \mu_3) - \mu_2^2 + \mu_3^2$ با شرط $\mu_1 > \mu_2 + \mu_3$ مثبت می‌ماند که در تابع \lr{\texttt{\_validate\_three\_intensity\_order}} بررسی می‌شود. در صورت کسر، سه جمله وجود دارد: $n_{X,\mu_2}^{-}$ با کران پایین \lr{decoy}، $n_{X,\mu_3}^{+}$ با کران بالای خلأ، و یک جمله تصحیحی که شامل $n_{X,\mu_1}^{+}$ (کران بالای سیگنال) و نسبت $s_{X,0}/\tau_0$ (تخمین بازده خلأ) است. این جمله تصحیحی اثر فوتون‌های چندتایی را حذف می‌کند. در صورت منفی شدن صورت، خروجی به صفر برده می‌شود.

\subsection{پیاده‌سازی معادله (4): کران بالای خطای تک‌فوتونی}

تابع \lr{\texttt{\_calc\_v1\_upper}} معادله (4) مقاله را پیاده‌سازی می‌کند:

\begin{LTR}
	\begin{lstlisting}[language=Python]
		def _calc_v1_upper(
		self,
		tau_1: float,
		mu2: float,
		mu3: float,
		m_mu2_plus: float,
		m_mu3_minus: float,
		physical_upper: float,
		) -> float:
		"""
		Eq. (4):
		v_{Z,1} <= tau_1 * (m^+_{Z,mu2} - m^-_{Z,mu3}) / (mu2 - mu3)
		"""
		denominator = mu2 - mu3
		if denominator <= 0.0:
		raise ConfigurationError("Invalid v1 denominator; require mu2 > mu3.")
		
		numerator = tau_1 * (m_mu2_plus - m_mu3_minus)
		return min(max(0.0, numerator / denominator), max(0.0, physical_upper))
	\end{lstlisting}
\end{LTR}

این تابع کران بالای تعداد خطاهای تک‌فوتونی در پایه $Z$ را محاسبه می‌کند. در اینجا از کران بالای خطاهای \lr{decoy} ضعیف ($m_{Z,\mu_2}^{+}$) و کران پایین خطاهای خلأ ($m_{Z,\mu_3}^{-}$) استفاده می‌شود تا کران بالا به دست آید. خروجی به یک کران فیزیکی (مثلاً $n_Z$ یا $m_Z$) محدود می‌شود زیرا تعداد خطاها نمی‌تواند از تعداد آشکارسازی‌ها بیشتر باشد. این \lr{clip} فیزیکی اطمینان می‌دهد که در شرایط نوسان شدید، کران بالای غیرمنطقی به دست نیاید.

\subsection{پیاده‌سازی معادله (5): کران بالای خطای فاز}

تابع \lr{\texttt{\_calc\_phi\_X\_upper}} پیچیده‌ترین بخش چارچوب را پیاده‌سازی می‌کند:

\begin{LTR}
	\begin{lstlisting}[language=Python]
		def _calc_phi_X_upper(
		self,
		v_Z_1_upper: float,
		s_Z_1_lower: float,
		s_X_1_lower: float,
		eps_a: float,
		) -> float:
		"""
		Eq. (5):
		phi_X = c_{X,1} / s_{X,1} <= v_{Z,1}/s_{Z,1} + gamma(a, b, c, d)
		"""
		self._validate_epsilon_value(eps_a, "eps_phase_est")
		
		if s_Z_1_lower <= 0.0 or s_X_1_lower <= 0.0:
		return MAX_ERROR_RATE
		
		e_Z_1 = self._safe_divide(v_Z_1_upper, s_Z_1_lower, default=MAX_ERROR_RATE)
		e_Z_1 = self._clamp_probability(e_Z_1, upper=MAX_ERROR_RATE)
		
		if e_Z_1 >= MAX_ERROR_RATE:
		return MAX_ERROR_RATE
		
		c = float(s_Z_1_lower)
		d = float(s_X_1_lower)
		
		if c <= 0.0 or d <= 0.0:
		return MAX_ERROR_RATE
		
		b = min(max(e_Z_1, MIN_EPS), MAX_ERROR_RATE - 1e-15)
		a = max(float(eps_a), MIN_EPS)
		
		denominator = c * d * (1.0 - b) * b * (a * a)
		if denominator <= 0.0 or not math.isfinite(denominator):
		return MAX_ERROR_RATE
		
		log_arg = (c + d) / denominator
		if log_arg <= 1.0 or not math.isfinite(log_arg):
		gamma = 0.0
		else:
		term1 = ((c + d) * (1.0 - b) * b) / (c * d * math.log(2.0))
		term2 = self._safe_log2(log_arg)
		product = term1 * term2
		gamma = math.sqrt(product) if product > 0.0 and math.isfinite(product) else MAX_ERROR_RATE
		
		return min(max(0.0, e_Z_1 + gamma), MAX_ERROR_RATE)
	\end{lstlisting}
\end{LTR}

این تابع شامل چندین لایه دفاعی برای مدیریت حالت‌های تباهی است. در ابتدا، $\varepsilon_a$ (یعنی $\varepsilon_{\text{sec}}/21$) اعتبارسنجی می‌شود. سپس اگر $s_{Z,1}$ یا $s_{X,1}$ صفر یا منفی باشند، تابع به صورت ایمن به $0.5$ بازمی‌گردد. نسبت $e_{Z,1} = v_{Z,1}/s_{Z,1}$ با \lr{\texttt{\_safe\_divide}} محاسبه می‌شود تا از تقسیم بر صفر جلوگیری شود و سپس به $[0, 0.5]$ \lr{clamp} می‌شود. اگر $e_{Z,1}$ به $0.5$ برسد، تابع بلافاصله $0.5$ برمی‌گرداند زیرا در این حالت هیچ امنیتی قابل استخراج نیست.

نکته ظریف این تابع، مدیریت حالت‌های تباهی $b \in \{0, 1\}$ است که در آن‌ها فرمول $\gamma$ تکین می‌شود. در این پیاده‌سازی، $b$ با \lr{\texttt{min(max(e\_Z\_1, MIN\_EPS), MAX\_ERROR\_RATE - 1e-15)}} به بازه $(10^{-300}, 0.5 - 10^{-15})$ محدود می‌شود تا از تکینگی جلوگیری شود. این \lr{clamp} بسیار کوچک است و در عمل نتیجه نهایی را تغییر نمی‌دهد، اما از خطاهای عددی مانند \lr{ZeroDivisionError} یا \lr{ValueError: math domain error} جلوگیری می‌کند. اگر مخرج $\gamma$ صفر یا نامتناهی شود، تابع به $0.5$ بازمی‌گردد که با رویکرد بدترین حالت مقاله همخوان است.

\subsection{محاسبه نشت تصحیح خطا}

تابع \lr{\texttt{\_calc\_error\_correction\_leakage}} دو مدل نشت را پشتیبانی می‌کند:

\begin{LTR}
	\begin{lstlisting}[language=Python]
		def _calc_error_correction_leakage(self, n_x: float, qber_x: float, eps_cor: float) -> Tuple[float, float]:
		"""
		Returns (leak_EC, separate_correctness_term).
		
		Models:
		- 'standard' / 'lim': paper Eq. (1).
		leak_EC      = f_EC * n_X * h(e_obs)
		corr_term    = log2(2 / eps_cor)
		- 'tomamichel': Ref. [37] of the paper. Includes the correctness
		logarithmic overhead in the leakage term.
		"""
		self._validate_epsilon_value(eps_cor, "eps_cor")
		
		qber_x = self._clamp_probability(qber_x, upper=MAX_ERROR_RATE)
		h_e = self.binary_entropy(qber_x)
		asymptotic_leak = float(self.p.f_error_correction * n_x * h_e)
		corr_term = self._safe_log2(2.0 / eps_cor)
		
		model = getattr(self.p, "error_correction_model", "standard")
		if model == "tomamichel":
		return asymptotic_leak + corr_term, 0.0
		
		return asymptotic_leak, corr_term
	\end{lstlisting}
\end{LTR}

در مدل استاندارد (که با مدل \lr{lim} نیز برابر است)، نشت تصحیح خطا $\lambda_{\text{EC}} = f_{\text{EC}} \cdot n_X \cdot h(e_{\text{obs}})$ است و جمله صحت $\log_2(2/\varepsilon_{\text{cor}})$ به صورت جداگانه در معادله (1) کم می‌شود. در مدل \lr{tomamichel}، جمله صحت به $\lambda_{\text{EC}}$ اضافه می‌شود و جمله جداگانه صفر می‌گردد تا از شمارش مضاعف جلوگیری شود. این طراحی به کاربر اجازه می‌دهد مدل دقیق‌تر \lr{Tomamichel} را انتخاب کند، در حالی که پیش‌فرض بر مدل استاندارد مقاله است.

\subsection{حل تکراری \texorpdfstring{$\varepsilon_{\text{sec}} = \kappa \cdot \ell$}{epsilon = kappa * l}}

در بخش \S IV مقاله، روشی پیشنهاد می‌شود که در آن $\varepsilon_{\text{sec}}$ به طول کلید متناسب می‌شود: $\varepsilon_{\text{sec}} = \kappa \cdot \ell$. این یک معادله ضمنی است زیرا $\ell$ خود به $\varepsilon_{\text{sec}}$ وابسته است. تابع \lr{\texttt{\_iterative\_kappa\_solver}} این معادله را با تکرار حل می‌کند:

\begin{LTR}
	\begin{lstlisting}[language=Python]
		def _iterative_kappa_solver(
		self,
		stats_map: Mapping[Union[str, int], TallyCounts],
		) -> KeyCalculationResult:
		"""
		Implements the paper's eps_sec = kappa * l convention (paper §IV).
		"""
		kappa = float(self.p.security_constant_kappa)
		if not math.isfinite(kappa) or kappa <= 0.0:
		raise ParameterValidationError("security_constant_kappa must be finite and positive.")
		
		n_total = sum(float(s.sent) for s in stats_map.values())
		l_est = max(1.0, min(n_total * 0.01, n_total))
		
		last_result: KeyCalculationResult | None = None
		converged = False
		
		for iteration in range(DEFAULT_KAPPA_MAX_ITERATIONS):
		eps_target = kappa * max(l_est, 1.0)
		# floor at 1e-10: smaller eps explodes Chernoff bounds (ln(eps/21) > 26)
		eps_target = min(max(eps_target, 1e-10), MAX_SECURITY_EPS)
		
		result = self._calculate_key_length_core(stats_map, eps_target)
		last_result = result
		l_new = float(result.secure_key_length)
		
		abs_delta = abs(l_new - l_est)
		rel_delta = abs_delta / max(1.0, abs(l_est))
		
		if abs_delta <= DEFAULT_KAPPA_ABS_TOL_BITS or rel_delta <= DEFAULT_KAPPA_REL_TOL:
		converged = True
		result.diagnostics.setdefault("dynamic_epsilon", {})
		result.diagnostics["dynamic_epsilon"].update(
		{
			"converged": True,
			"iterations": iteration + 1,
			"kappa": kappa,
			"eps_sec_final": eps_target,
		}
		)
		return result
		
		if l_new <= 0.0:
		l_est = max(1.0, 0.5 * l_est)
		else:
		l_est = 0.5 * (l_est + l_new)
		
		assert last_result is not None
		last_result.diagnostics.setdefault("dynamic_epsilon", {})
		last_result.diagnostics["dynamic_epsilon"].update(
		{
			"converged": converged,
			"iterations": DEFAULT_KAPPA_MAX_ITERATIONS,
			"kappa": kappa,
			"warning": "dynamic epsilon iteration did not converge",
		}
		)
		last_result.diagnostics.setdefault("warnings", []).append("Dynamic epsilon iteration did not converge.")
		return last_result
	\end{lstlisting}
\end{LTR}

الگوریتم با یک تخمین اولیه $l_{\text{est}} = \max(1, \min(0.01 \cdot n_{\text{total}}, n_{\text{total}}))$ آغاز می‌شود که 1\% از کل پالس‌ها را به عنوان طول کلید تقریبی در نظر می‌گیرد. در هر تکرار، $\varepsilon_{\text{target}} = \kappa \cdot \max(l_{\text{est}}, 1)$ محاسبه می‌شود. یک نکته حیاتی در این نسخه این است که \lr{floor} کف $\varepsilon_{\text{target}}$ از $10^{-15}$ به $10^{-10}$ افزایش یافته است. این تغییر ضروری است زیرا با معرفی استراتژی هیبرید \lr{Hoeffding/Chernoff}، کران‌های \lr{Chernoff} مستقیماً به $\ln(1/\varepsilon)$ وابسته هستند و برای $\varepsilon < 10^{-10}$، مقدار $\ln(21/\varepsilon) > 26$ می‌شود که باعث می‌شود کران‌های \lr{Chernoff} به‌طور پاتولوژیک شل شوند (کران بالا بیش از حد بزرگ، کران پایین منفی). با کف $10^{-10}$، حداکثر مقدار $\ln(21/\varepsilon) \approx 25.8$ است که در محدوده قابل قبول عددی باقی می‌ماند.

معیار همگرایی ترکیبی از تلورانس مطلق (100 بیت) و تلورانس نسبی ($10^{-6}$) است. در صورت عدم همگرایی، تخمین به صورت میانگین حسابی بین تخمین قبلی و مقدار جدید به‌روزرسانی می‌شود. اگر $l_{\text{new}} \le 0$ (یعنی کلیدی تولید نشده)، تخمین به نصف کاهش می‌یابد تا $\varepsilon$ سخت‌گیرانه‌تر شود. در صورت عدم همگرایی پس از 50 تکرار، آخرین نتیجه با هشدار بازگردانده می‌شود. این طراحی \lr{fail-soft} تضمین می‌کند که حتی در صورت عدم همگرایی، یک خروجی معتبر (با احتمالاً طول کلید صفر) تولید می‌شود و خط لوله شبیه‌سازی متوقف نمی‌شود.

\subsection{نرمال‌سازی \lr{stats\_map} و نتایج صفر}

متد \lr{\texttt{\_normalize\_stats\_map}} یک لایه انعطاف‌پذیر برای تطبیق کلیدهای \lr{stats\_map} با نام پالس‌ها فراهم می‌کند:

\begin{LTR}
	\begin{lstlisting}[language=Python]
		def _normalize_stats_map(
		self,
		stats_map: Mapping[Union[str, int], TallyCounts],
		pulse_names: List[str],
		) -> Dict[str, TallyCounts]:
		if all(name in stats_map for name in pulse_names):
		return {name: stats_map[name] for name in pulse_names}
		
		if all(str(i) in stats_map for i in range(len(pulse_names))):
		return {pulse_names[i]: stats_map[str(i)] for i in range(len(pulse_names))}
		
		if all(i in stats_map for i in range(len(pulse_names))):
		return {pulse_names[i]: stats_map[i] for i in range(len(pulse_names))}
		
		raise ParameterValidationError(
		f"stats_map keys do not match pulse names or positional indices. "
		f"Expected names {pulse_names!r}; got {list(stats_map.keys())!r}."
		)
	\end{lstlisting}
\end{LTR}

این متد سه فرمت کلید را پشتیبانی می‌کند: (1) کلیدهای نام پالس (مانند \lr{\texttt{"signal"}}، \lr{\texttt{"decoy"}}، \lr{\texttt{"vacuum"}}) که صریح‌ترین حالت است؛ (2) کلیدهای اندیس رشته‌ای (مانند \lr{\texttt{"0"}}، \lr{\texttt{"1"}}، \lr{\texttt{"2"}}) که معمولاً در خروجی \lr{JSON}/\lr{CSV} استفاده می‌شود؛ (3) کلیدهای اندیس صحیح (مانند \lr{\texttt{0}}، \lr{\texttt{1}}، \lr{\texttt{2}}) که در تعاملات درون‌حافظه‌ای رایج است. این انعطاف‌پذیری به ماژول اجازه می‌دهد با انواع فرمت‌های ورودی از لایه‌های مختلف شبیه‌سازی کار کند بدون آنکه نیاز به تبدیل صریح باشد. در صورت عدم تطابق با هیچ‌یک از این سه فرمت، یک \lr{\texttt{ParameterValidationError}} با پیام مشخص پرتاب می‌شود.

متد \lr{\texttt{\_zero\_result}} یک ابزار کمکی برای تولید نتایج صفر با اطلاعات تشخیصی غنی است:

\begin{LTR}
	\begin{lstlisting}[language=Python]
		def _zero_result(
		self,
		branch: str,
		diagnostics: Dict[str, Any],
		error_codes: List[ErrorCode] | None = None,
		) -> KeyCalculationResult:
		diagnostics = dict(diagnostics)
		diagnostics["return_branch"] = branch
		return KeyCalculationResult(
		secure_key_length=0,
		privacy_amplification_term=0.0,
		error_correction_leakage=0.0,
		phase_error_rate_upper_bound=MAX_ERROR_RATE,
		error_codes=error_codes or [ErrorCode.INSUFFICIENT_STATISTICS],
		diagnostics=_json_safe_dict(diagnostics.items()),
		)
	\end{lstlisting}
\end{LTR}

این متد در تمام مسیرهای شکست استفاده می‌شود: آمار ناکافی (\lr{\texttt{insufficient\_statistics\_nX\_or\_nZ}})، \lr{stats\_map} نامطبوع (\lr{\texttt{invalid\_or\_mismatched\_stats}})، $\tau$ نامعتبر (\lr{\texttt{invalid\_tau}})، و پیکربندی سه‌شدت نامعتبر (\lr{\texttt{invalid\_three\_intensity\_configuration}}). هر شاخه در \lr{diagnostics} به‌صورت \lr{\texttt{return\_branch}} ثبت می‌شود تا ریشه‌یابی شکست ساده شود. پیش‌فرض کد خطا \lr{\texttt{INSUFFICIENT\_STATISTICS}} است، اما می‌توان کدهای دیگری نیز ارسال کرد. استفاده از \lr{\texttt{\_json\_safe\_dict}} اطمینان می‌دهد که \lr{diagnostics} قابل سریال‌سازی به \lr{JSON} است.

\subsection{محاسبه اصلی طول کلید}

تابع \lr{\texttt{\_calculate\_key\_length\_core}} هسته اصلی محاسبه است. پس از اعتبارسنجی و آماده‌سازی، شمارش‌های مشاهده‌شده را به کران‌های \lr{Hoeffding/Chernoff} تبدیل می‌کند، تخمین‌های $s_{X,0}^L$، $s_{X,1}^L$، $s_{Z,1}^L$، $v_{Z,1}^U$ و $\varphi_X^U$ را محاسبه می‌کند و در نهایت معادله (1) را ارزیابی می‌کند. این متد با آماده‌سازی \lr{diagnostics} کامل آغاز می‌شود:

\begin{LTR}
	\begin{lstlisting}[language=Python]
		diagnostics: Dict[str, Any] = {
			"epsilon_accounting": {
				"scheme": "Lim2014 Appendix B (21-event split, eps_sec = 21*eps)",
				"eps_sec_used": float(eps_sec_override),
				"single_test_epsilon": float(epsilon),
				"failure_events": LIM2014_FAILURE_EVENTS_3_INTENSITY,
				"hoeffding_ln_term": float(eps_ln_term),
				"bound_method": "hoeffding_chernoff_hybrid",
				"note": (
				"Uses a hybrid Hoeffding / per-intensity Chernoff bound. "
				"For large counts (n_k >= 10 * delta_paper, relative "
				"statistical error <= 10%), uses the paper-faithful "
				"Hoeffding delta = sqrt(n_total/2 * ln(21/eps_sec)). "
				"For small counts (0 < n_k < threshold), uses the tighter "
				"per-intensity Chernoff delta_k = sqrt(2 * n_k * ln(1/eps)). "
				"For n_k = 0, uses the multiplicative Chernoff upper bound "
				"ln(1/eps_per_event) (rule of three). This is provably "
				"<= Hoeffding for all n_k, so security is never weakened."
				),
			},
			"mu_map": {k: float(v) for k, v in mu_map.items()},
			"p_k_map": {k: float(v) for k, v in p_k_map.items()},
			"warnings": [],
		}
	\end{lstlisting}
\end{LTR}

این ساختار \lr{diagnostics} غنی به کاربر اجازه می‌دهد تمام جزئیات محاسبه را بازبینی کند. فیلد \lr{\texttt{bound\_method}} با مقدار \lr{\texttt{"hoeffding\_chernoff\_hybrid"}} به‌طور صریح نشان می‌دهد که از استراتژی هیبرید استفاده شده است. فیلد \lr{\texttt{note}} یک توضیح کامل از سه رژیم ارائه می‌دهد که برای \lr{audit} امنیتی و بازبینی توسط شخص ثالث ضروری است. سپس شمارش‌های $n_X$، $n_Z$، $m_X$، $m_Z$ و تفکیک آن‌ها به‌ازای هر شدت در \lr{diagnostics} ذخیره می‌شوند.

در مسیر تحلیلی سه‌شدت، کران‌های \lr{Hoeffding/Chernoff} برای هر شدت در هر دو پایه محاسبه و در \lr{diagnostics} ذخیره می‌شوند:

\begin{LTR}
	\begin{lstlisting}[language=Python]
		n_X_k_bounds = {name: self._get_bounds(n_X_k[name], n_X, eps_ln_term, clip_total=n_X) for name in mu_map}
		n_Z_k_bounds = {name: self._get_bounds(n_Z_k[name], n_Z, eps_ln_term, clip_total=n_Z) for name in mu_map}
		m_Z_k_bounds = {name: self._get_bounds(m_Z_k[name], m_Z, eps_ln_term, clip_total=m_Z) for name in mu_map}
		
		diagnostics["n_X_k_bounds"] = {
			k: [float(v[0]), float(v[1])] for k, v in n_X_k_bounds.items()
		}
	\end{lstlisting}
\end{LTR}

سپس ضریب پیشانی $e^{\mu_k}/p_k$ برای تبدیل کران‌های خام به فرم استفاده‌شده در معادلات (2)-(4) اعمال می‌شود. به‌عنوان مثال، برای کران بالای بازده خلأ:

\begin{LTR}
	\begin{lstlisting}[language=Python]
		s_X_0_L = self._calc_s0_lower(
		tau_0,
		mu2,
		mu3,
		n_X_k_bounds[k2][1] / p2 * math.exp(mu2),  # decoy UPPER bound (n^+_{X,mu2})
		n_X_k_bounds[k3][0] / p3 * math.exp(mu3),  # vacuum LOWER bound (n^-_{X,mu3})
		)
	\end{lstlisting}
\end{LTR}

نکته مهم در این نسخه این است که کران بالای خلأ ($n_{X,\mu_3}^{+}$) اکنون به‌درستی توسط \lr{\texttt{\_get\_bounds}} به‌صورت $\ln(1/\varepsilon)$ محاسبه می‌شود وقتی مشاهده صفر است (کران تنگ \lr{Chernoff}، نه $\delta_{\text{paper}}$ مبتنی بر $n_X$). این یک \lr{patch} ضد محافظه‌کارانه نیست، بلکه یک اصلاح درست است: در نسخه‌های قبلی، وقتی $n_{X,\mu_3} = 0$، کران بالای \lr{Hoeffding} به $\delta_{\text{paper}}$ می‌رسید که می‌تواند بسیار بزرگتر از حد معقول باشد. با استراتژی هیبرید، این کران به $\ln(1/\varepsilon)$ محدود می‌شود که هم متناهی و هم تنگ‌تر است. در نتیجه، تخمین $s_{X,0}^L$ در فواصل طولانی (جایی که شمارش خلأ اغلب صفر است) به‌طور قابل توجهی بهبود می‌یابد و طول کلید نهایی بالاتر می‌شود بدون آنکه امنیت تضعیف شود.

پس از محاسبه $\varphi_X^U$ و $\lambda_{\text{EC}}$، معادله (1) به صورت مستقیم ارزیابی می‌شود:

\begin{LTR}
	\begin{lstlisting}[language=Python]
		phi_X_U = self._calc_phi_X_upper(v_Z_1_U, s_Z_1_L, s_X_1_L, epsilon)
		
		qber_x = self._safe_divide(m_X, n_X, default=MAX_ERROR_RATE)
		qber_x = self._clamp_probability(qber_x, upper=MAX_ERROR_RATE)
		
		leak_EC, corr_term = self._calc_error_correction_leakage(n_X, qber_x, self.p.eps_cor)
		
		# Eq. (1):
		#   l = s_{X,0} + s_{X,1} * (1 - h(phi_X))
		#       - lambda_EC - 6*log2(21/eps_sec) - log2(2/eps_cor)
		term_s0 = s_X_0_L
		term_s1 = s_X_1_L * (1.0 - self.binary_entropy(phi_X_U))
		
		pa_term = 6.0 * self._safe_log2(LIM2014_FAILURE_EVENTS_3_INTENSITY / eps_sec_override)
		
		key_len = term_s0 + term_s1 - leak_EC - pa_term - corr_term
		key_len_clamped = self._clamp_key_length(float(key_len))
	\end{lstlisting}
\end{LTR}

جمله $6\log_2(21/\varepsilon_{\text{sec}})$ که در مقاله به عنوان جمله \lr{privacy amplification} شناخته می‌شود، شش رویداد شکست مرتبط با تقویت حریم خصوصی را پوشش می‌دهد. در نهایت طول کلید با \lr{\texttt{\_clamp\_key\_length}} به عدد صحیح نامنفی تبدیل می‌شود. اگر طول خام منفی باشد (یعنی نشت‌ها از آنتروپی بیشتر شوند)، هشداری در \lr{diagnostics} ثبت می‌شود. کدهای خطای مناسب نیز به نتیجه اضافه می‌شوند تا علت شکست قابل ردیابی باشد.

\subsection{قید دفاعی فیزیکی در \lr{calculate\_key\_length}}

متد \lr{\texttt{calculate\_key\_length}} یک لایه دفاعی فیزیکی برای قید طول کلید به کلید \lr{sifted} کل اعمال می‌کند:

\begin{LTR}
	\begin{lstlisting}[language=Python]
		def calculate_key_length(
		self,
		stats_map: Mapping[Union[str, int], TallyCounts],
		) -> KeyCalculationResult:
		self._validate_configuration()
		
		if getattr(self.p, "security_constant_kappa", None) is not None:
		result = self._iterative_kappa_solver(stats_map)
		else:
		result = self._calculate_key_length_core(stats_map, self.p.eps_sec)
		
		# === Defensive physical clamp ===
		# The secure key length CANNOT exceed the total number of basis-matched
		# raw detections (n_X + n_Z = total sifted bits).
		try:
		_n_total = 0
		for _k, _v in stats_map.items():
		if hasattr(_v, 'sifted'):
		_n_total += int(_v.sifted)
		if _n_total > 0 and result.secure_key_length > _n_total:
		import logging as _logging
		_logging.getLogger('Lim2014Proof').warning(
		f'Defensive clamp triggered: secure_key_length {result.secure_key_length} '
		f'> total sifted {_n_total}; clamping to {_n_total}. '
		f'This indicates an internal bug in the kappa solver or decoy estimation.'
		)
		from qkd.datatypes import KeyCalculationResult as _KCR
		result = _KCR(
		secure_key_length=_n_total,
		privacy_amplification_term=result.privacy_amplification_term,
		error_correction_leakage=result.error_correction_leakage,
		phase_error_rate_upper_bound=result.phase_error_rate_upper_bound,
		error_codes=result.error_codes,
		diagnostics=result.diagnostics,
		)
		except Exception as _e:
		import logging as _logging
		_logging.getLogger('Lim2014Proof').debug(
		f'Defensive clamp check failed (non-fatal): {_e}'
		)
		
		return result
	\end{lstlisting}
\end{LTR}

این قید بر اساس یک اصل بدیهی فیزیکی است: \textbf{بیت‌های کلید امن زیرمجموعه‌ای از بیت‌های خام \lr{sifted} هستند}. در شرایط عادی، این قید هرگز فعال نمی‌شود زیرا فرمول‌های \lr{Lim 2014} به‌طور ذاتی محافظه‌کارانه هستند و طول کلید نهایی همیشه کمتر از $n_X$ (که خود زیرمجموعه‌ای از کل \lr{sifted} است) می‌شود. اما به‌عنوان یک تور ایمنی در برابر باگ‌های داخلی طراحی شده است. اگر در اثر یک باگ در حلقه \lr{kappa-solver}، \lr{fallback} در \lr{LP}، یا سرریز در تخمین \lr{decoy}، طول کلید نهایی از $n_X + n_Z$ بیشتر شود، این قید فعال می‌شود، طول کلید را به $n_X + n_Z$ محدود می‌کند، یک هشدار در \lr{log} ثبت می‌کند و نتیجه بازنویسی‌شده را بازمی‌گرداند.

نکته طراحی مهم این است که این قید در یک بلوک \lr{\texttt{try/except}} قرار دارد. اگر خود بررسی قید به هر دلیلی (مانند فقدان فیلد \lr{\texttt{sifted}} در یک \lr{TallyCounts} قدیمی) شکست بخورد، استثنا به‌صورت غیرکشنده (\lr{non-fatal}) در سطح \lr{debug} ثبت می‌شود و نتیجه اصلی بدون تغییر بازگردانده می‌شود. این طراحی تضمین می‌کند که قید دفاعی هرگز به‌طور خود یک منبع شکست جدید معرفی نکند. این فلسفه \lr{fail-safe} با اصل امنیتی \lr{defense in depth} همخوان است: چندین لایه دفاعی مستقل، هر کدام با \lr{fallback} امن.

\section{ارسال و دریافت با سایر ماژول‌ها}

\subsection{ورودی‌های ماژول}

ماژول \lr{\texttt{lim2014.py}} وابستگی‌های داخلی و خارجی متعددی دارد. وابستگی‌های داخلی شامل موارد زیر است:

\begin{itemize}
	\item \lr{\texttt{qkd.proofs.base}}: کلاس پایه \lr{\texttt{FiniteKeyProof}}، \lr{\texttt{DecoyEstimates}}، \lr{\texttt{KeyCalculationResult}}، \lr{\texttt{ErrorCode}} و تابع کمکی \lr{\texttt{\_json\_safe\_dict}} از این ماژول وارد می‌شوند. این کلاس پایه رابط عمومی و خدمات مشترک مانند آنتروپی باینری، \lr{clamp} عددی و گزارش‌گیری \lr{audit} را فراهم می‌سازد.
	\item \lr{\texttt{qkd.proofs.utils\_lp}}: تابع \lr{\texttt{solve\_lp}} برای حل مسأله برنامه‌ریزی خطی در مسیر \lr{N-decoy} وارد می‌شود. این تابع از \lr{scipy.optimize.linprog} یا \lr{cvxopt} به عنوان \lr{solver} زیرین استفاده می‌کند.
	\item \lr{\texttt{qkd.datatypes}}: \lr{\texttt{EpsilonAllocation}}، \lr{\texttt{TallyCounts}}، \lr{\texttt{ConfidenceBoundMethod}} و \lr{\texttt{KeyCalculationResult}} از این ماژول وارد می‌شوند. \lr{\texttt{TallyCounts}} شامل فیلدهای \lr{\texttt{sent}}، \lr{\texttt{sifted\_x}}، \lr{\texttt{sifted\_z}}، \lr{\texttt{errors\_sifted\_x}} و \lr{\texttt{errors\_sifted\_z}} است که شمارش‌های آماری شبکه آشکارسازی را نگه می‌دارد. \lr{\texttt{ConfidenceBoundMethod}} برای اعمال روش \lr{CHERNOFF} در سازنده استفاده می‌شود.
	\item \lr{\texttt{qkd.exceptions}}: استثناهای سفارشی \lr{\texttt{ConfigurationError}}، \lr{\texttt{LPFailureError}} و \lr{\texttt{ParameterValidationError}} برای گزارش خطاهای ساختاریافته.
	\item \lr{\texttt{qkd.params.QKDParams}}: کلاس پارامترهای QKD که شامل $\varepsilon_{\text{sec}}$، $\varepsilon_{\text{cor}}$، $f_{\text{EC}}$، \lr{\texttt{photon\_number\_cap}}، \lr{\texttt{security\_constant\_kappa}}، \lr{\texttt{error\_correction\_model}}، \lr{\texttt{lp\_solver\_method}}، \lr{\texttt{ci\_method}} و \lr{\texttt{source.pulse\_configs}} است. توجه شود که \lr{\texttt{ci\_method}} در سازنده به \lr{\texttt{CHERNOFF}} بازنویسی می‌شود زیرا امنیت کلید متناهی در این چارچوب به‌طور صریح به کران‌های \lr{Chernoff} وابسته است.
\end{itemize}

وابستگی‌های خارجی شامل \lr{\texttt{numpy}} برای محاسبات برداری در مسیر \lr{LP}، \lr{\texttt{math}} برای توابع ریاضی پایه، \lr{\texttt{logging}} برای ثبت رویدادها و \lr{\texttt{mpmath}} (اختیاری) برای محاسبات با دقت بالا هستند.

\subsection{خروجی‌های ماژول}

خروجی اصلی ماژول دو شیء است:

\begin{enumerate}
	\item \lr{\texttt{KeyCalculationResult}}: این شیء از کلاس پایه به ارث رسیده و شامل فیلدهای \lr{\texttt{secure\_key\_length}} ($\ell$)، \lr{\texttt{privacy\_amplification\_term}} (جریمه \lr{PA})، \lr{\texttt{error\_correction\_leakage}} ($\lambda_{\text{EC}}$)، \lr{\texttt{phase\_error\_rate\_upper\_bound}} ($\varphi_X^U$)، \lr{\texttt{error\_codes}} و \lr{\texttt{diagnostics}} است. \lr{diagnostics} شامل اطلاعات کامل مانند \lr{\texttt{epsilon\_accounting}} (با \lr{\texttt{bound\_method}}، \lr{\texttt{hoeffding\_ln\_term}} و \lr{\texttt{note}})، مقادیر میانی $s_{X,0}^L$، $s_{X,1}^L$، $s_{Z,1}^L$، $v_{Z,1}^U$، $\tau_0$، $\tau_1$، کران‌های \lr{Hoeffding/Chernoff} برای هر شدت، مدل تصحیح خطای انتخاب‌شده و در صورت استفاده، اطلاعات همگرایی \lr{kappa-solver} است.
	\item \lr{\texttt{DecoyEstimatesLim2014}}: این شیء خروجی متد \lr{\texttt{estimate\_yields\_and\_errors}} است و فیلدهای اختصاصی چارچوب \lr{Lim 2014} را نگه می‌دارد. این شیء برای کاربرانی که می‌خواهند تخمین‌های میانی را بدون محاسبه کامل طول کلید بررسی کنند، مفید است.
\end{enumerate}

\subsection{جریان داده در سراسر فریم‌ورک}

جریان داده کلی به این شکل است:
\begin{enumerate}
	\item لایه شبیه‌ساز پالس‌های نوری را با شدت‌های $\mu_1$، $\mu_2$، $\mu_3$ تولید و از کانال کوانتومی عبور می‌دهد.
	\item آشکارسازها در سمت باب، آشکارسازی‌ها و خطاها را در دو پایه $X$ و $Z$ و برای هر شدت ثبت کرده و در \lr{\texttt{TallyCounts}} ذخیره می‌کنند.
	\item این شمارش‌ها به صورت یک دیکشنری \lr{\texttt{stats\_map}} با کلیدهای نام شدت، اندیس رشته‌ای یا اندیس صحیح به اثبات امنیتی داده می‌شوند. متد \lr{\texttt{\_normalize\_stats\_map}} به‌طور خودکار فرمت کلید را تشخیص می‌دهد.
	\item متد \lr{\texttt{calculate\_key\_length}} ابتدا پیکربندی را اعتبارسنجی می‌کند. اگر \lr{\texttt{security\_constant\_kappa}} تنظیم شده باشد، مسیر \lr{kappa-solver} تکراری انتخاب می‌شود؛ در غیر این صورت مسیر مستقیم با $\varepsilon_{\text{sec}}$ ثابت.
	\item در مسیر مستقیم، $\varepsilon$ تک‌آزمون به صورت $\varepsilon_{\text{sec}}/21$ محاسبه می‌شود و جمله \lr{Hoeffding}/\lr{Chernoff} $\ln(21/\varepsilon_{\text{sec}})$ آماده می‌شود.
	\item سپس $\tau_0$ و $\tau_1$ محاسبه می‌شوند. اگر تعداد شدت‌ها بیش از 3 باشد، مسیر \lr{LP} انتخاب می‌شود؛ در غیر این صورت مسیر تحلیلی معادلات (2) تا (5).
	\item در مسیر تحلیلی، شدت‌ها به ترتیب نزولی مرتب می‌شوند تا نقش سیگنال، \lr{decoy} و خلأ مشخص شود. سپس کران‌های هیبرید \lr{Hoeffding/Chernoff} برای هر شدت در هر دو پایه محاسبه می‌شوند. این کران‌ها به‌طور خودکار در رژیم مناسب (بزرگ، کوچک یا صفر) عمل می‌کنند.
	\item با استفاده از این کران‌ها، $s_{X,0}^L$، $s_{X,1}^L$، $s_{Z,1}^L$ و $v_{Z,1}^U$ محاسبه می‌شوند.
	\item سپس $\varphi_X^U$ با تابع \lr{\texttt{\_calc\_phi\_X\_upper}} محاسبه می‌شود.
	\item در نهایت، معادله (1) ارزیابی می‌شود و طول کلید امن به دست می‌آید.
	\item قید دفاعی فیزیکی بررسی می‌کند که طول کلید از کل \lr{sifted} ($n_X + n_Z$) بیشتر نشود. در صورت فعال شدن، هشدار ثبت و مقدار بازنویسی می‌شود.
	\item لایه بالایی می‌تواند با \lr{\texttt{explain\_key\_decision}} (از کلاس پایه) علت نتیجه را به صورت متنی دریافت کند و با \lr{\texttt{audit\_dump}} یک گزارش کامل برای بازبینی امنیتی تولید کند.
\end{enumerate}

این جداسازی مراحل به گونه‌ای طراحی شده که هر مرحله قابل آزمایش و بازبینی مستقل باشد. به عنوان مثال، می‌توان \lr{\texttt{DecoyEstimatesLim2014}} را به صورت دستی ساخت و رفتار مرحله نهایی محاسبه کلید را در شرایط خاص بررسی کرد. این ماژول‌بندی از اصل \lr{Single Responsibility} پیروی می‌کند و قابلیت نگهداری و توسعه کد را به طور قابل توجهی افزایش می‌دهد. علاوه بر این، شفافیت کامل در \lr{diagnostics} به کاربر اجازه می‌دهد هر مرحله از محاسبه را ردیابی کرده و در صورت نیاز، فرضیات تئوریک را برای کاربرد خاص خود بازبینی کند. استراتژی هیبرید \lr{Hoeffding/Chernoff} به‌طور خاص برای شبیه‌سازی‌های بلندمسافت که در آن‌ها شمارش \lr{decoy} یا خلأ ممکن است صفر یا بسیار کم باشد، یک بهبود چشمگیر در کارایی عملی فراهم می‌کند بدون آنچه هیچ‌گونه تضعیف امنیتی به همراه داشته باشد.

	
	\chapter{برهان‌های \lr{(Ma, 2005) و (Wang, 2005)}}
	\label{ch:ma_wang}
	\sloppy
\emergencystretch=5em
\hbadness=10000
\vbadness=10000
\tolerance=2000
\providecolor{codebg}{RGB}{248,248,248}
\providecolor{codeframe}{RGB}{200,200,200}
\providecolor{keyword}{RGB}{0,0,180}
\providecolor{string}{RGB}{163,21,21}
\providecolor{comment}{RGB}{0,128,0}
\lstset{
	basicstyle=\ttfamily\scriptsize,
	breaklines=true,
	breakatwhitespace=false,
	columns=fullflexible,
	keepspaces=true,
	frame=single,
	framesep=2pt,
	showstringspaces=false,
	backgroundcolor=\color{codebg},
	rulecolor=\color{codeframe},
	keywordstyle=\color{keyword}\bfseries,
	stringstyle=\color{string},
	commentstyle=\color{comment}\itshape,
	numberstyle=\tiny\color{gray},
	numbers=left,
	numbersep=6pt,
	xleftmargin=14pt,
	framexleftmargin=14pt,
	literate=
	{ą}{{\v{a}}}1 {ę}{{\k{e}}}1 {ż}{{\.z}}1 {ś}{{\'s}}1
	{ł}{{\l}}1 {ć}{{\'c}}1 {ń}{{\'n}}1 {ó}{{\'o}}1 {ź}{{\'z}}1
	{→}{{$\rightarrow$}}1 {←}{{$\leftarrow$}}1 {≠}{{$\neq$}}1
}

\section{نمای‌ کلی و هدف ماژول}

فایل‌های \lr{\texttt{qkd/proofs/ma2005.py}} و \lr{\texttt{qkd/proofs/wang2005.py}} دو ستون فقرات تئوریک چارچوب شبیه‌سازی \lr{QKD} در زمینه پروتکل‌های \lr{Decoy-State} را تشکیل می‌دهند. این دو ماژول به صورت موازی عمل می‌کنند و هر کدام پاسخگوی نیاز متفاوتی در تحلیل امنیتی هستند: ماژول \lr{ma2005.py} پیاده‌سازی وفادارانه و چندگانه چارچوب مرجع \lr{Ma, Qi, Zhao, Lo} (سال 2005) است که در مقاله \lr{``Practical decoy state for quantum key distribution''} در مجله \lr{Physical Review A} (جلد 72، شماره 012326) منتشر شده، و ماژول \lr{wang2005.py} پیاده‌سازی چارچوب تکمیلی \lr{Wang} (سال 2008) است که در مقاله \lr{``General theory of decoy-state quantum cryptography with source errors''} در همان مجله (جلد 77، شماره 042311) ارائه شده است. در حالی که چارچوب \lr{Ma 2005} فرض می‌کند که منبع نوری آلیس یک منبع \lr{coherent state} ایده‌آل با شدت دقیقاً مشخص است، چارچوب \lr{Wang 2008} این فرض را کنار می‌گذارد و اجازه می‌دهد شدت منبع دارای نوسان باشد یا حتی خود حالت کوانتومی از یک \lr{coherent state} ایده‌آل منحرف شود. این دو چارچوب در کنار هم طیف کاملی از فرضیات منبع را پوشش می‌دهند و به کاربر اجازه می‌دهند با انتخاب ماژول مناسب، تحلیل امنیتی متناسب با شرایط واقعی منبع خود را انتخاب کند.

ماژول \lr{ma2005.py} شامل پنج پیاده‌سازی مجزا از چارچوب \lr{Ma 2005} است که هر کدام برای سناریوی خاصی بهینه شده‌اند. این پنج کلاس عبارتند از: \lr{\texttt{Ma2005VacuumWeakProof}} که پروتکل کلاسیک \lr{``Vacuum + Weak''} را با دو \lr{decoy} (یک \lr{decoy} ضعیف با شدت $\nu$ و یک \lr{decoy} خلأ) پیاده‌سازی می‌کند و در برابر نوسان شدت منبع به صورت \lr{worst-case} مقاوم است؛ \lr{\texttt{Ma2005GeneralTwoDecoyProof}} که حالت تعمیم‌یافته دو \lr{decoy} را با استفاده از فرمول‌های تحلیلی معادلات (18)، (21) و (25) مقاله پیاده می‌کند و بر خلاف \lr{Vacuum + Weak}، فرض خلأ بودن یکی از \lr{decoy}‌ها را ندارد؛ \lr{\texttt{Ma2005OneDecoyProof}} که پروتکل تک‌‌‌decoy را با دو حالت تحلیلی (با یا بدون فرض \lr{zero-background}) پیاده‌سازی می‌کند و برای شبکه‌هایی با محدودیت سخت‌افزاری مناسب است؛ \lr{\texttt{Ma2005GeneralMultiDecoyProof}} که برای حالت چند \lr{decoy} (بیش از دو) از برنامه‌ریزی خطی (\lr{LP}) برای یافتن کران‌های بهینه $Y_1$ و $e_1$ استفاده می‌کند و تنها پیاده‌سازی در این ماژول است که از توزیع \lr{Thermal} پشتیبانی می‌کند؛ و در نهایت \lr{\texttt{Ma2005AsymptoticProof}} که حد نظری \lr{``Infinite Decoy State''} را با معادلات (27) و (28) مقاله پیاده می‌کند و به عنوان مرجع بالای عملکرد (\lr{upper bound}) عمل می‌کند.

ماژول \lr{wang2005.py} یک پیاده‌سازی تک‌کلاسی به نام \lr{\texttt{Wang2005Proof}} دارد که از کلاس پایه \lr{\texttt{FiniteKeyProof}} ارث می‌برد و یک \lr{dataclass} تخصصی به نام \lr{\texttt{DecoyEstimatesWang2005}} را معرفی می‌کند که فیلد اضافی \lr{\texttt{intensity\_error\_bound}} را برای ثبت کران خطای منبع نگه می‌دارد. این ماژول از همان فرمول پایه \lr{Vacuum + Weak} استفاده می‌کند اما تمام شدت‌ها ($\mu$ و $\nu$) را با کران‌های بالا و پایین ($\mu^U$, $\mu^L$, $\nu^U$, $\nu^L$) جایگزین می‌کند تا در بدترین حالت نوسان منبع، امنیت حفظ شود. علاوه بر این، ماژول \lr{wang2005.py} از یک مفهوم جدید به نام \lr{source fidelity} پشتیبانی می‌کند که نشان می‌دهد حالت واقعی منبع چقدر به یک \lr{coherent state} ایده‌آل نزدیک است. اگر \lr{fidelity} کمتر از یک باشد، ماژول به صورت خودکار یک \lr{mapping} محافظه‌کارانه از $(1 - F)$ به یک \lr{delta} مؤثر انجام می‌دهد و آن را در فرمول‌های کران شدت جایگزین می‌کند.

وظیفه نهایی هر دو ماژول، تبدیل شمارش‌های آماری شبکه آشکارسازی (که در شیء \lr{\texttt{TallyCounts}} ذخیره شده‌اند) به یک تخمین امن از طول کلید نهایی است. این تبدیل با رعایت سخت‌گیرانه تمام محافظه‌کاری‌های ریاضی، آماری و فیزیکی انجام می‌شود تا هیچ طول کلیدی که از نظر تئوری قابل اثبات نیست، تولید نشود. در صورت بروز ناسازگاری پیکربندی، آمار ناکافی، یا خطای عددی، خروجی به جای \lr{raise} کردن استثنا، یک شیء \lr{\texttt{KeyCalculationResult}} با طول کلید صفر و کد خطای مناسب تولید می‌کند تا خط لوله شبیه‌سازی متوقف نشود و علت شکست در گزارش تشخیصی (\lr{diagnostics}) ثبت گردد. این طراحی \lr{fail-safe} برای کاربردهای تولیدی که در آن‌ها قطع شبیه‌سازی هزینه بر است، حیاتی است.

اهمیت تاریخی و علمی این دو ماژول در این است که چارچوب \lr{Ma 2005} نخستین پیاده‌سازی عملی و قابل استفاده از ایده \lr{Decoy-State} بود که توسط \lr{Hwang, Lo} (2003) ارائه شده بود. پیش از این، ایده \lr{Decoy-State} تنها یک مفهوم نظری بود و پیاده‌سازی عملی آن به دلیل پیچیدگی محاسباتی و فقدان فرمول‌های بسته ممکن نبود. \lr{Ma 2005} با ارائه فرمول‌های تحلیلی بسته برای کران‌های $Y_1$ و $e_1$، این ایده را به یک فناوری قابل اجرا تبدیل کرد. چارچوب \lr{Wang 2008} سپس این پیاده‌سازی را با در نظر گرفتن نقص‌های واقعی منبع (مانند نوسان شدت لیزر و انحراف از حالت \lr{coherent}) توسعه داد. این دو چارچوب در کنار هم پایه‌ای برای تمام پروتکل‌های \lr{Decoy-State QKD} تجاری مدرن (مانند شبکه‌های \lr{Toshiba, ID Quantique, Quantum CTek}) شده‌اند.

\section{مبانی تئوری و فیزیکی}

\subsection{حمله \lr{Photon-Number Splitting} و انگیزه \lr{Decoy-State}}

در پروتکل استاندارد \lr{BB84} با یک منبع تک‌شدت، امنیت در برابر حمله \lr{Photon-Number Splitting (PNS)} قابل تضمین نیست. در این حمله، یک مهاجم (اِو) با استفاده از یک \lr{beam splitter} غیرتعاملی (\lr{ND-PS}) یا تکنیک‌های \lr{QND}، یک فوتون از هر پالس چندفوتونی استخراج می‌کند و بقیه را به سمت باب ارسال می‌کند. از آنجا که این کار بدون ایجاد خطا انجام می‌شود، باب متوجه حضور اِو نمی‌شود. پس از پایه‌گذاری عمومی، اِو فوتون‌های ذخیره‌شده خود را در پایه درست اندازه می‌گیرد و کلید را به طور کامل به دست می‌آورد. در سیستم‌های واقعی با منبع \lr{coherent state}، بخش قابل توجهی از پالس‌ها چندفوتونی هستند (به ویژه وقتی $\mu \approx 0.5$ که برای دستیابی به نرخ کلید مطلوب انتخاب می‌شود)، و در نتیجه حمله \lr{PNS} یک تهدید جدی است.

راه‌حل \lr{Decoy-State} این است که آلیس به طور تصادفی شدت نور ارسالی را میان چندین سطح تغییر می‌دهد. باب پس از اندازه‌گیری، با مقایسه نرخ آشکارسازی در شدت‌های مختلف، می‌تواند تخمین بزند که مهاجم در حال انجام \lr{PNS} هست یا نه. ایده کلیدی این است که بازده تک‌فوتونی $Y_1$ و خطای تک‌فوتونی $e_1$ مستقل از انتخاب شدت آلیس است (زیرا کانال و آشکارساز فقط تعداد فوتون‌ها را می‌بینند نه شدت انتخاب‌شده را). در نتیجه، با ترکیب اطلاعات از چندین شدت، می‌توان $Y_1$ و $e_1$ را به صورت کران‌های بسته تخمین زد. چارچوب \lr{Ma 2005} نخستین کسی بود که این تخمین را به صورت فرمول‌های تحلیلی بسته ارائه داد.

\subsection{توزیع آماری فوتون‌های ارسالی}

توزیع آماری فوتون‌های ارسالی از یک منبع به نوع آماری آن بستگی دارد. دو نوع آماری متداول در شبیه‌سازی \lr{QKD} عبارتند از \lr{Poisson} (برای منبع \lr{coherent state} که از لیزر معمولی به دست می‌آید) و \lr{Thermal} (برای منبع \lr{thermal state} که از منابع حرارتی یا \lr{chaotic light} به دست می‌آید). اگر میانگین فوتون $\mu$ باشد، احتمال ارسال $n$ فوتون در توزیع \lr{Poisson} برابر است با:
\begin{equation}
	P(n \mid \mu) = e^{-\mu}\,\frac{\mu^n}{n!},
\end{equation}
و در توزیع \lr{Thermal} برابر است با:
\begin{equation}
	P(n \mid \mu) = \frac{\mu^n}{(1 + \mu)^{n+1}}.
\end{equation}
تفاوت کلیدی این دو توزیع در واریانس است: برای \lr{Poisson} واریانس برابر میانگین ($\text{Var} = \mu$) است، در حالی که برای \lr{Thermal} واریانس برابر $\mu(1 + \mu)$ است (بزرگ‌تر از میانگین). این بدان معناست که در \lr{Thermal}، بخش بیشتری از پالس‌ها چندفوتونی هستند و در نتیجه حمله \lr{PNS} تهدید جدی‌تری است. در کد، تابع \lr{\texttt{\_get\_prob\_n}} در کلاس \lr{\texttt{Ma2005BaseProof}} این دو توزیع را بر اساس پرچم \lr{\texttt{is\_thermal}} (که از پیکربندی منبع استخراج می‌شود) انتخاب می‌کند. تنها کلاس \lr{\texttt{Ma2005GeneralMultiDecoyProof}} از توزیع \lr{Thermal} پشتیبانی می‌کند زیرا تنها این کلاس از \lr{LP} استفاده می‌کند و فرمول‌های تحلیلی بسته برای \lr{Thermal} پیچیده‌تر و در دسترس نیستند.

\subsection{بازده، خطا و \lr{Gain}}

برای هر شدت $\mu$، دو کمیت قابل مشاهده تعریف می‌شوند: \lr{Gain} (نرخ آشکارسازی) $Q_\mu$ و \lr{QBER} $E_\mu$. این دو کمیت به ترتیب برابرند با:
\begin{align}
	Q_\mu &= \sum_{n=0}^{\infty} Y_n\,P(n \mid \mu), \\
	E_\mu\,Q_\mu &= \sum_{n=0}^{\infty} e_n\,Y_n\,P(n \mid \mu),
\end{align}
که در آن $Y_n$ بازده (احتمال آشکارسازی مشروط به ارسال $n$ فوتون) و $e_n$ خطای تک‌فوتونی مشروط به ارسال $n$ فوتون است. این کمیت‌ها مستقل از انتخاب شدت آلیس هستند و تنها به خصوصیات کانال و آشکارساز بستگی دارند. این مستقل‌بودن، پایه ریاضی کل ایده \lr{Decoy-State} است. در کد، این کمیت‌ها به صورت $Q_\mu = \text{sifted}/\text{sent}$ و $E_\mu = \text{errors\_sifted}/\text{sifted}$ از شیء \lr{\texttt{TallyCounts}} استخراج می‌شوند.

بازده $Y_n$ را می‌توان به صورت $Y_n = Y_0 + \eta_n$ تجزیه کرد که در آن $Y_0$ نرخ آشکارسازی \lr{dark count} (مستقل از فوتون ارسالی) و $\eta_n$ احتمال آشکارسازی حداقل یک فوتون از $n$ فوتون ارسالی است. برای یک کانال با عبور $\eta$ و آشکارساز با راندمان $\eta_{\text{det}}$، بازده تک‌فوتونی برابر $Y_1 = Y_0 + \eta_{\text{channel}}\,\eta_{\text{det}}$ است. این رابطه در کلاس \lr{\texttt{Ma2005AsymptoticProof}} برای استخراج $\eta$ از روی مشاهدات و سپس محاسبه $Y_1$ استفاده می‌شود.

\subsection{کران پایین بازده تک‌فوتونی در پروتکل \lr{Vacuum + Weak}}

معادله (34) مقاله \lr{Ma 2005}، کران پایین $Y_1$ را برای پروتکل \lr{Vacuum + Weak} (با شدت‌های سیگنال $\mu$، \lr{decoy} ضعیف $\nu$ و خلأ) ارائه می‌دهد:
\begin{equation}
	Y_1 \;\ge\; \frac{\mu}{\mu\nu - \nu^2}\left[Q_\nu\,e^{\nu} - Q_\mu\,e^{\mu}\,\frac{\nu^2}{\mu^2} - Y_0\,\frac{\mu^2 - \nu^2}{\mu^2}\right].
\end{equation}
در این فرمول، $Q_\nu$ و $Q_\mu$ به ترتیب \lr{Gain} مشاهده‌شده در شدت \lr{decoy} و سیگنال هستند. برای محافظه‌کاری، از کران‌های بالا و پایین آماری استفاده می‌شود: $Q_\nu$ با کران پایین $Q_\nu^L$ و $Q_\mu$ با کران بالا $Q_\mu^U$ جایگزین می‌شوند تا کران پایین $Y_1$ تا حد ممکن سخت‌گیرانه باشد. این انتخاب‌های جهت‌دار تضمین می‌کنند که تخمین نهایی، در بدترین حالت نوسان آماری، همچنان یک کران پایین معتبر باشد.

\subsection{کران بالای خطای تک‌فوتونی}

معادله (37) مقاله \lr{Ma 2005}، کران بالای $e_1$ را ارائه می‌دهد:
\begin{equation}
	e_1 \;\le\; \frac{E_\nu\,Q_\nu\,e^{\nu} - e_0\,Y_0}{Y_1\,\nu},
\end{equation}
که در آن $e_0 = 0.5$ خطای آشکارسازی خلأ است (زیرا \lr{dark count} کاملاً تصادفی است و احتمال خطا برابر 50\% دارد). در این فرمول، $E_\nu\,Q_\nu$ با کران بالا $E_\nu^U\,Q_\nu^U$ جایگزین می‌شود تا کران بالای $e_1$ تا حد ممکن بزرگ (سخت‌گیرانه) باشد. در کد، شرط $Y_1 \cdot \nu > 10^{-20}$ بررسی می‌شود تا از تقسیم بر صفر جلوگیری شود؛ اگر این شرط نقض شود، $e_1 = 0.5$ تنظیم می‌شود که معادل بدترین حالت ممکن (کامل تصادفی) است.

\subsection{کسر برچسب‌گذاری‌شده (\lr{Tagged Fraction})}

یک مفهوم کلیدی دیگر در چارچوب \lr{Ma 2005}، کسر برچسب‌گذاری‌شده $\Delta$ است که نشان می‌دهد چه کسری از پالس‌های سیگنال به جای تک‌فوتونی، چندفوتونی هستند. معادله (36) مقاله، کران بالای $\Delta$ را به صورت زیر ارائه می‌دهد:
\begin{equation}
	\Delta \;\le\; \frac{\nu}{\mu - \nu}\left[\frac{\nu\,e^{-\nu}\,Q_\mu}{\mu\,e^{-\mu}\,Q_\nu} - 1\right] + \frac{\nu\,e^{-\nu}\,Y_0}{\mu\,Q_\nu}.
\end{equation}
این کران در متد \lr{\texttt{\_calculate\_weak\_gllp\_rate}} برای محاسبه نرخ کلید در حالت \lr{Weak GLLP} (معادله 43) استفاده می‌شود. اگر $\Delta \ge 1$ باشد، به این معناست که تمام پالس‌های سیگنال چندفوتونی هستند و هیچ امنیتی وجود ندارد؛ در این حالت، نرخ کلید به صفر برده می‌شود.

\subsection{فرمول \lr{GLLP} برای نرخ کلید}

چارچوب \lr{Gottesman-Lo-Lütkenhaus-Preskill (GLLP)} فرمول کلی نرخ کلید امن را به صورت زیر ارائه می‌دهد:
\begin{equation}
	R = q\left\{Q_1\left[1 - H_2(e_1)\right] - Q_\mu\,f(E_\mu)\,H_2(E_\mu)\right\},
\end{equation}
که در آن $q$ فاکتور پایه‌گذاری (برای \lr{BB84} برابر 0.5 است زیرا نیمی از پالس‌ها در پایه اشتباه دور ریخته می‌شوند)، $Q_1 = Y_1\,P(1 \mid \mu) = Y_1\,\mu\,e^{-\mu}$ \lr{Gain} تک‌فوتونی، $H_2(\cdot)$ آنتروپی باینری، $f(E_\mu)$ بازدهی تصحیح خطا و جمله اول ترم تقویت حریم خصوصی (\lr{PA}) و جمله دوم ترم نشت تصحیح خطا (\lr{EC leakage}) است. این فرمول در متد \lr{\texttt{\_calculate\_gllp\_rate}} پیاده‌سازی شده است. خروجی این متد یک شیء \lr{\texttt{KeyCalculationResult}} است که شامل طول کلید امن (به صورت صحیح)، کران بالای خطای فاز ($e_1$)، نشت تصحیح خطا و یک دیکشنری \lr{diagnostics} با تمام مقادیر میانی است.

\subsection{فرمول \lr{Weak GLLP} (معادله 43)}

علاوه بر فرمول استاندارد \lr{GLLP}، مقاله \lr{Ma 2005} یک فرمول جایگزین به نام \lr{Weak GLLP} (معادله 43) ارائه می‌دهد که از کسر برچسب‌گذاری‌شده $\Delta$ به جای $Y_1$ و $e_1$ استفاده می‌کند:
\begin{equation}
	R = q\,Q_\mu\left[-H_2(E_\mu) + (1 - \Delta)\left(1 - H_2\!\left(\frac{E_\mu}{1 - \Delta}\right)\right)\right].
\end{equation}
این فرمول در حالتی که تخمین دقیق $Y_1$ دشوار است اما کران $\Delta$ قابل محاسبه است، مفید است. در کد، اگر پرچم \lr{\texttt{use\_weak\_gllp}} در پیکربندی فعال باشد، این فرمول به جای فرمول استاندارد استفاده می‌شود. یک شرط فیزیکی در کد بررسی می‌شود: اگر $E_\mu / (1 - \Delta) \ge 0.5$ باشد، نرخ کلید به صفر برده می‌شود زیرا آنتروپی باینری در 0.5 بیشینه است و ترم \lr{PA} منفی می‌شود.

\subsection{حد \lr{Asymptotic} با \lr{Infinite Decoy}}

معادلات (27) و (28) مقاله \lr{Ma 2005}، حد نظری \lr{Infinite Decoy State} را ارائه می‌دهند که در آن تعداد \lr{decoy}‌ها به بی‌نهایت میل می‌کند و تمام $Y_n$ و $e_n$ به طور کامل قابل تخمین هستند:
\begin{align}
	Y_1 &= Y_0 + \eta, \\
	e_1 &= \frac{e_0\,Y_0 + e_{\text{det}}\,\eta}{Y_1},
\end{align}
که در آن $\eta = \eta_{\text{channel}} \cdot \eta_{\text{det}}$ عبور کلی و $e_{\text{det}}$ خطای ذاتی آشکارساز است. در کلاس \lr{\texttt{Ma2005AsymptoticProof}}، $\eta$ به صورت غیرمستقیم از روی $Q_\mu$ استخراج می‌شود:
\begin{equation}
	Q_\mu = Y_0 + 1 - e^{-\eta\,\mu} \quad \Rightarrow \quad \eta = -\frac{\ln(1 + Y_0 - Q_\mu)}{\mu}.
\end{equation}
اگر $1 + Y_0 - Q_\mu \le 0$ (که یک ناهنجاری فیزیکی است و نشان می‌دهد $Q_\mu$ بیش از حد بزرگ است)، کد از مقدار تئوریک $\eta$ از پیکربندی به عنوان \lr{fallback} استفاده می‌کند. این کلاس به عنوان مرجع بالای عملکرد عمل می‌کند و اجازه می‌دهد کاربر ببیند در شرایط ایده‌آل، طول کلید چقدر می‌تواند باشد.

\subsection{خطای منبع و چارچوب \lr{Wang 2008}}

در چارچوب \lr{Wang 2008}، فرض می‌شود که شدت منبع به طور دقیق مشخص نیست بلکه در یک بازه $[\mu^L, \mu^U]$ قرار دارد. این مدل‌سازی برای منابع لیزر واقعی که دارای نوسان شدت هستند، ضروری است. کران‌های بالا و پایین شدت به صورت زیر تعریف می‌شوند:
\begin{align}
	\mu^U &= \mu\,(1 + \delta), \\
	\mu^L &= \mu\,(1 - \delta), \\
	\nu^U &= \nu\,(1 + \delta), \\
	\nu^L &= \nu\,(1 - \delta),
\end{align}
که در آن $\delta$ کران خطای شدت است. در کد، این $\delta$ از دو منبع ممکن است به دست آید: یا مستقیماً از \lr{\texttt{intensity\_jitter}} پیکربندی منبع خوانده می‌شود، یا از طریق \lr{source fidelity} به صورت غیرمستقیم محاسبه می‌شود. اگر \lr{source fidelity} $F < 1$ باشد، یک \lr{mapping} محافظه‌کارانه به صورت $\delta = \max(\delta_{\text{jitter}}, 1 - F)$ انجام می‌شود. این \lr{mapping} فرض می‌کند که انحراف از حالت \lr{coherent} ایده‌آل، حداقل به اندازه $(1 - F)$ در شدت مؤثر اثر دارد.

سپس فرمول کران پایین $Y_1$ در چارچوب \lr{Wang 2008} (معادله 56 مقاله) به صورت زیر بازنویسی می‌شود:
\begin{equation}
	Y_1 \;\ge\; \frac{\mu^L}{\mu^U \nu^L - (\nu^U)^2}\left[Q_\nu\,e^{\nu^L} - Q_\mu\,e^{\mu^U}\,\frac{(\nu^U)^2}{(\mu^L)^2} - Q_{\text{vac}}\,\frac{(\mu^U)^2 - (\nu^L)^2}{(\mu^L)^2}\right].
\end{equation}
تفاوت کلیدی با فرمول \lr{Ma 2005} این است که تمام شدت‌ها با کران‌های بدترین حالت جایگزین شده‌اند. به طور خاص، $\mu$ با $\mu^U$ در مخرج‌های جمله‌های کم‌شونده و با $\mu^L$ در ضریب پیشانی جایگزین شده است تا کران تا حد ممکن سخت‌گیرانه باشد. به طور متقابل، $\nu$ با $\nu^L$ در جملات مثبت و با $\nu^U$ در جملات منفی جایگزین شده است. این انتخاب‌های جهت‌دار تضمین می‌کنند که کران نهایی، در بدترین حالت نوسان منبع، همچنان معتبر باشد.

کران بالای $e_1$ نیز به صورت مشابه بازنویسی می‌شود:
\begin{equation}
	e_1 \;\le\; \frac{E_\nu\,Q_\nu\,e^{\nu^U} - E_{\text{vac}}\,Q_{\text{vac}}}{Y_1\,\nu^L}.
\end{equation}
در این فرمول، $e^{\nu}$ با $e^{\nu^U}$ (کران بالا) جایگزین شده تا صورت کسر بزرگ‌تر و در نتیجه کران بالای $e_1$ سخت‌گیرانه‌تر شود. همچنین $Y_1$ و $\nu$ در مخرج با کران پایین $Y_1^L$ و $\nu^L$ جایگزین شده‌اند تا کران تا حد ممکن بزرگ باشد. در کد، اگر $Y_1 \cdot \nu^L \le 0$ باشد، $e_1 = 0.5$ تنظیم می‌شود که معادل بدترین حالت است.

\subsection{نرخ کلید در چارچوب \lr{Wang 2008}}

محاسبه نرخ کلید در چارچوب \lr{Wang 2008} از همان فرمول استاندارد \lr{GLLP} استفاده می‌کند اما با کران‌های $Y_1$ و $e_1$ که از فرمول‌های بالا به دست آمده‌اند:
\begin{equation}
	R = \frac{1}{2}\left[Y_1^L\,\mu\,e^{-\mu}\left(1 - H_2(e_1^U)\right) - Q_\mu\,f_{\text{EC}}\,H_2(E_\mu)\right].
\end{equation}
در اینجا $q = 0.5$ برای \lr{BB84} ثابت فرض شده و $f_{\text{EC}}$ از پیکربندی خوانده می‌شود. در کد، این محاسبه در متد \lr{\texttt{calculate\_key\_length}} کلاس \lr{\texttt{Wang2005Proof}} پیاده‌سازی شده است. خروجی نهایی یک شیء \lr{\texttt{KeyCalculationResult}} است که شامل طول کلید امن، نشت تصحیح خطا و کران بالای خطای فاز است.

\section{تحلیل معماری و ساختار داده‌ها}

\subsection{سلسله‌مراتب کلاس در ماژول \lr{ma2005.py}}

ماژول \lr{ma2005.py} از یک طراحی سلسله‌مراتبی استفاده می‌کند که در آن یک کلاس پایه به نام \lr{\texttt{Ma2005BaseProof}} تمام منطق مشترک را نگه می‌دارد و پنج کلاس مشتق، هر کدام یک پروتکل خاص را پیاده‌سازی می‌کنند. کلاس پایه از کلاس انتزاعی \lr{\texttt{FiniteKeyProof}} (که در \lr{\texttt{qkd/proofs/base.py}} تعریف شده) ارث می‌برد و متدهای انتزاعی \lr{\texttt{allocate\_epsilons}}، \lr{\texttt{get\_epsilon\_policy}}، \lr{\texttt{estimate\_yields\_and\_errors}}، \lr{\texttt{calculate\_key\_length}} و \lr{\texttt{notation\_map}} را پیاده‌سازی می‌کند. این طراحی از اصل \lr{DRY (Don't Repeat Yourself)} پیروی می‌کند و از تکرار کد در پنج کلاس مشتق جلوگیری می‌کند.

در ادامه، ساختار کلاس پایه \lr{\texttt{Ma2005BaseProof}} را به صورت خلاصه نمایش می‌دهیم:

\begin{LTR}
	\begin{lstlisting}[language=Python]
		class Ma2005BaseProof(FiniteKeyProof):
		"""Base class for Ma 2005 proofs sharing the GLLP rate formula."""
		
		def __init__(self, params: QKDParams, *args, **kwargs):
		if params.ci_method != ConfidenceBoundMethod.CHERNOFF:
		logger.warning(
		f"Ma2005 proof requires CHERNOFF bounds for finite-key security. "
		f"Overriding '{params.ci_method.name}' with CHERNOFF."
		)
		object.__setattr__(params, 'ci_method', ConfidenceBoundMethod.CHERNOFF)
		
		super().__init__(params, *args, **kwargs)
		self.is_thermal = (self.p.source.statistics_type == SourceStatisticsType.THERMAL)
	\end{lstlisting}
\end{LTR}

در این کد، سازنده کلاس ابتدا بررسی می‌کند که روش فاصله اطمینان (\lr{CI}) در پیکرببری برابر \lr{CHERNOFF} باشد. اگر کاربر روش دیگری (مانند \lr{Gaussian} یا \lr{Clopper-Pearson}) انتخاب کرده باشد، یک هشدار ثبت می‌شود و روش به صورت صریح به \lr{CHERNOFF} تغییر می‌یابد. دلیل این الزام امنیتی آن است که فرمول کران پایین $Y_1^L$ شامل تفریق جمله‌های تقویت‌شده با $e^\mu$ است؛ اگر کران‌های فاصله اطمینان بیش از حد پهن باشند (مانند کران‌های \lr{Gaussian} که با وابستگی $\mathcal{O}(1/\sqrt{N})$ همگرا می‌شوند)، این تفریق می‌تواند $Y_1^L$ را منفی کند که فیزیکی نیست و به معنای شکست تخمین است. در مقابل، کران‌های \lr{Chernoff} که از نوع \lr{multiplicative} هستند (یعنی $\hat{Q} \in [Q \cdot e^{-\delta}, Q \cdot e^{+\delta}]$ با $\delta = \sqrt{2\ln(2/\varepsilon)/N}$)، به اندازه کافی فشرده هستند تا $Y_1^L > 0$ را در رژیم کلید متناهی حفظ کنند.

یک نکته فنی مهم در این پیاده‌سازی، استفاده از \lr{\texttt{object.\_\_setattr\_\_(params, ...)}} است. این الگو به این دلیل استفاده می‌شود که شیء \lr{\texttt{QKDParams}} در این فریم‌ورک به صورت \lr{frozen} (غیرقابل تغییر) تعریف شده است و تخصیص مستقیم \lr{\texttt{params.ci\_method = ...}} منجر به \lr{\texttt{FrozenInstanceError}} می‌شود. دور زدن این محدودیت با \lr{\texttt{object.\_\_setattr\_\_}} به ما اجازه می‌دهد در شرایط خاص (مانند الزام امنیتی به \lr{CHERNOFF}) مقدار فیلد را بازنویسی کنیم، در حالی که در شرایط عادی همچنان از مزایای \lr{immutability} مانند جلوگیری از تغییرات ناخواسته و قابلیت \lr{hashing} بهره می‌بریم. این یک الگوی شناخته‌شده در پایتون برای \lr{frozen dataclass}‌هاست که در آن نیاز به بازنویسی موقت یک فیلد وجود دارد. سپس پرچم \lr{\texttt{is\_thermal}} از پیکربندی منبع استخراج می‌شود که تعیین می‌کند کدام توزیع آماری (\lr{Poisson} یا \lr{Thermal}) در محاسبات استفاده شود.

\subsection{متد \lr{\texttt{\_get\_prob\_n}} و انتخاب توزیع آماری}

این متد کلیدی‌ترین تابع مشترک در کلاس پایه است که احتمال ارسال $n$ فوتون با میانگین $\mu$ را بر اساس نوع آماری منبع محاسبه می‌کند:

\begin{LTR}
	\begin{lstlisting}[language=Python]
		def _get_prob_n(self, mu: float, n: int) -> float:
		"""Calculates P(n|mu) handling both Poisson and Thermal cases."""
		if self.is_thermal:
		return (mu**n) / ((1.0 + mu)**(n + 1))
		else:
		return math.exp(-mu) * (mu**n) / math.factorial(n)
	\end{lstlisting}
\end{LTR}

این متد در دو مسیر اصلی استفاده می‌شود: نخست، در متد \lr{\texttt{\_calculate\_gllp\_rate}} برای محاسبه $P(1 \mid \mu) = \mu\,e^{-\mu}$ که برای محاسبه $Q_1 = Y_1\,P(1 \mid \mu)$ ضروری است؛ دوم، در کلاس \lr{\texttt{Ma2005GeneralMultiDecoyProof}} برای ساخت ماتریس \lr{LP} که در آن هر سطر معادله $Q_\mu = \sum_n Y_n\,P(n \mid \mu)$ را نشان می‌دهد. این متد به صورت \lr{polymorphic} عمل می‌کند: با تغییر پرچم \lr{\texttt{is\_thermal}}، کل رفتار سیستم بدون تغییر کد قابل تغییر است. این طراحی از اصل \lr{Open-Closed Principle} پیروی می‌کند: کد برای توسعه باز اما برای تغییر بسته است.


\subsection{نرمال‌سازی کلیدهای \lr{stats\_map} با \lr{\texttt{\_normalize\_stats\_map}}}

یکی از چالش‌های عملی در اتصال لایه شبیه‌سازی به لایه اثبات امنیتی، تنوع فرمت‌های ورودی است. لایه شبیه‌ساز ممکن است کلیدهای دیکشنری \lr{\texttt{stats\_map}} را به سه شکل مختلف ارائه دهد: نام‌های رشته‌ای پالس (مانند “\lr{signal}”، “\lr{decoy}”، “\lr{vacuum}”)، کلیدهای عددی رشته‌ای (مانند “\lr{0}”، “\lr{1}”، “\lr{2}”) یا کلیدهای عدد صحیح (مانند \lr{0}، \lr{1}، \lr{2}). متد \lr{\texttt{\_normalize\_stats\_map}} این سه فرمت را به یک فرمت استاندارد با کلیدهای نامی تبدیل می‌کند تا تمام کد داخلی ماژول \lr{Ma 2005} بتواند با اطمینان از کلیدهای نامی استفاده کند:

\begin{LTR}
	\begin{lstlisting}[language=Python]
		def _normalize_stats_map(
		self,
		stats_map: Mapping[Union[str, int], TallyCounts],
		pulse_names: List[str],
		) -> Dict[str, TallyCounts]:
		"""Normalize stats_map keys to use pulse name strings."""
		# Try #1: name keys already present
		if all(name in stats_map for name in pulse_names):
		return {name: stats_map[name] for name in pulse_names}
		
		# Try #2: string-integer keys ("0", "1", "2")
		if all(str(i) in stats_map for i in range(len(pulse_names))):
		return {pulse_names[i]: stats_map[str(i)]
			for i in range(len(pulse_names))}
		
		# Try #3: integer keys (0, 1, 2)
		if all(i in stats_map for i in range(len(pulse_names))):
		return {pulse_names[i]: stats_map[i]
			for i in range(len(pulse_names))}
		
		raise ParameterValidationError(
		f"stats_map keys do not match pulse names or positional indices. "
		f"Expected names {pulse_names!r}; got {list(stats_map.keys())!r}."
		)
		
		def _get_pulse_names(self) -> List[str]:
		"""Return the ordered list of pulse names from the source config."""
		return [pc.name for pc in self.p.source.pulse_configs]
	\end{lstlisting}
\end{LTR}

این متد با سه تلاش متوالی کار می‌کند: نخست، فرض می‌کند که کلیدها از قبل نام‌های پالس هستند و اگر این فرض درست باشد، همان دیکشنری را با فیلتراسیون به نام‌های معتبر برمی‌گرداند. دوم، اگر کلیدها رشته‌های عددی باشند، آن‌ها را به نام‌های پالس بر اساس ترتیب موقعیتی نگاشت می‌کند. سوم، اگر کلیدها اعداد صحیح باشند، همان کار را انجام می‌دهد. اگر هیچ یک از این سه حالت منطبق نباشد، استثنای \lr{\texttt{ParameterValidationError}} با پیام تشخیصی مشخص \lr{raise} می‌شود. این طراحی از اصل \lr{Postel‌s Law} (در سخاوتمندی باشید، در محافظه‌کاری سخت‌گیرانه) پیروی می‌کند: ورودی‌های متنوع پذیرفته می‌شوند، اما خروجی همیشه یک فرمت ثابت و معتبر است. تمام پنج کلاس مشتق از \lr{\texttt{Ma 2005}} در ابتدای متد \lr{\texttt{estimate\_yields\_and\_errors}} خود، این نرمال‌سازی را فراخوانی می‌کنند تا از سازگاری ورودی اطمینان حاصل شود.

\subsection{تخصیص بودجه \lr{$\varepsilon$} در چارچوب \lr{Ma 2005}}

متد \lr{\texttt{allocate\_epsilons}} در کلاس پایه، بودجه امنیتی کلی $\varepsilon_{\text{sec}}$ را میان اجزای مختلف تقسیم می‌کند. چارچوب \lr{Ma 2005} از معادله زیر برای تقسیم بودجه استفاده می‌کند:
\begin{equation}
	\varepsilon_{\text{sec}} \;\ge\; \varepsilon_{\text{PE}} + 2\,\varepsilon_{\text{smooth}} + \varepsilon_{\text{PA}} + \varepsilon_{\text{cor}} + \varepsilon_{\text{phase\_est}}.
\end{equation}
در این معادله، $\varepsilon_{\text{PE}}$ بودجه تخمین پارامتر، $\varepsilon_{\text{smooth}}$ بودجه هموارسازی، $\varepsilon_{\text{PA}}$ بودجه تقویت حریم خصوصی، $\varepsilon_{\text{cor}}$ بودجه تصحیح خطا و $\varepsilon_{\text{phase\_est}}$ بودجه تخمین خطای فاز است. در کد، $\varepsilon_{\text{phase\_est}} = \varepsilon_{\text{PE}} / 10$ تنظیم می‌شود که یک تقسیم محافظه‌کارانه است. اگر مجموع بودجه‌های استفاده‌شده از $\varepsilon_{\text{sec}}$ بیشتر یا مساوی شود (که ناسازگاری پیکرببری را نشان می‌دهد)، به جای جایگزینی صریح با $\varepsilon_{\text{sec}}/8$ (که قصد کاربر را به طور خام نادیده می‌گیرد)، کد از \lr{Gap [A] fix v2} استفاده می‌کند: مقادیر به صورت \emph{متناسب} نرمال می‌شوند تا مجموعشان دقیقاً برابر $\varepsilon_{\text{sec}}$ شود. این رویکرد وزن‌دهی نسبی که کاربر مشخص کرده را حفظ می‌کند و فقط مقیاس کلی را تعدیل می‌کند. فرمول نرمال‌سازی به صورت زیر است:

\begin{equation}
	\text{scale} = \frac{\varepsilon_{\text{sec}}}{\text{used\_budget}}, \qquad \varepsilon_k^{\text{norm}} = \varepsilon_k \cdot \text{scale}, \quad \forall k \in \{\text{PE}, \text{smooth}, \text{cor}\}.
\end{equation}

سپس $\varepsilon_{\text{phase\_est}}^{\text{norm}} = \varepsilon_{\text{PE}}^{\text{norm}} / 10$ و $\varepsilon_{\text{PA}}^{\text{norm}} = \varepsilon_{\text{sec}} - (\varepsilon_{\text{PE}}^{\text{norm}} + 2\,\varepsilon_{\text{smooth}}^{\text{norm}} + \varepsilon_{\text{cor}}^{\text{norm}} + \varepsilon_{\text{phase\_est}}^{\text{norm}})$ محاسبه می‌شوند. اگر $\varepsilon_{\text{PA}}^{\text{norm}} < 0$ شود، به صفر تنظیم می‌شود. در شرایط عادی (زمانی که بودجه کافی است)، $\varepsilon_{\text{PA}}$ به صورت باقی‌مانده $\varepsilon_{\text{sec}} - \text{used\_budget}$ محاسبه می‌شود تا کل بودجه به طور کامل استفاده شود. پیاده‌سازی کامل این متد در ادامه ارائه شده است:

\begin{LTR}
	\begin{lstlisting}[language=Python]
		def allocate_epsilons(self) -> EpsilonAllocation:
		# Full constraint: eps_pe + 2*eps_smooth + eps_pa + eps_cor + eps_phase_est <= eps_sec
		eps_phase_est = self.p.eps_pe / 10.0
		used_budget = (self.p.eps_pe + 2.0 * self.p.eps_smooth
		+ self.p.eps_cor + eps_phase_est)
		
		if used_budget >= self.p.eps_sec:
		# Gap [A] fix (v2): normalize proportionally so they sum to eps_sec
		scale = self.p.eps_sec / used_budget if used_budget > 0 else 1.0
		eps_cor_n   = self.p.eps_cor * scale
		eps_pe_n    = self.p.eps_pe * scale
		eps_smooth_n = self.p.eps_smooth * scale
		eps_phase_est_n = eps_pe_n / 10.0
		eps_pa_n = (self.p.eps_sec
		- (eps_pe_n + 2.0 * eps_smooth_n
		+ eps_cor_n + eps_phase_est_n))
		if eps_pa_n < 0.0:
		eps_pa_n = 0.0
		return EpsilonAllocation(
		eps_sec=self.p.eps_sec,
		eps_cor=eps_cor_n,
		eps_pe=eps_pe_n,
		eps_smooth=eps_smooth_n,
		eps_pa=eps_pa_n,
		eps_phase_est=eps_phase_est_n
		)
		
		eps_pa = self.p.eps_sec - used_budget
		return EpsilonAllocation(
		eps_sec=self.p.eps_sec,
		eps_cor=self.p.eps_cor,
		eps_pe=self.p.eps_pe,
		eps_smooth=self.p.eps_smooth,
		eps_pa=eps_pa,
		eps_phase_est=eps_phase_est
		)
	\end{lstlisting}
\end{LTR}

این پیاده‌سازی دو ویژگی مهم دارد: نخست، در مسیر نرمال‌سازی متناسب، تمام اجزا (به جز $\varepsilon_{\text{PA}}$) با همان ضریب $\text{scale}$ کوچک می‌شوند تا نسبت‌های نسبی حفظ شوند؛ دوم، $\varepsilon_{\text{PA}}$ به عنوان باقی‌مانده محاسبه می‌شود تا تضمین شود مجموع دقیقاً برابر $\varepsilon_{\text{sec}}$ است. این رویکرد نسبت به روش قبلی (جایگزینی همه با $\varepsilon_{\text{sec}}/8$) برتری دارد زیرا اگر کاربر مثلاً $\varepsilon_{\text{PE}}$ را سه برابر $\varepsilon_{\text{cor}}$ تنظیم کرده باشد، این نسبت پس از نرمال‌سازی نیز حفظ می‌شود و معنای فیزیکی پیکربندی حفظ می‌گردد.

\subsection{ساختار \lr{dataclass} در ماژول \lr{wang2005.py}}

ماژول \lr{wang2005.py} یک \lr{dataclass} جدید به نام \lr{\texttt{DecoyEstimatesWang2005}} معرفی می‌کند که از \lr{\texttt{DecoyEstimates}} پایه ارث می‌برد و فیلد اضافی \lr{\texttt{intensity\_error\_bound}} را اضافه می‌کند:

\begin{LTR}
	\begin{lstlisting}[language=Python]
		@dataclass(frozen=True)
		class DecoyEstimatesWang2005(DecoyEstimates):
		"""Decoy estimates including source error bounds."""
		intensity_error_bound: float = 0.0  # delta in Wang 2008
	\end{lstlisting}
\end{LTR}

استفاده از دکوراتور \lr{\texttt{@dataclass(frozen=True)}} دو مزیت دارد: نخست، شیء پس از ساخت غیرقابل تغییر (\lr{immutable}) می‌شود که از تغییرات ناخواسته در طول محاسبات جلوگیری می‌کند و امنیت نوع داده را تضمین می‌کند؛ دوم، شیء به طور خودکار \lr{hashable} می‌شود که اجازه می‌دهد در مجموعه‌ها یا به عنوان کلید دیکشنری استفاده شود. این ویژگی برای \lr{caching} و \lr{memoization} در محاسبات تکراری مفید است. مقدار پیش‌فرض 0.0 برای \lr{\texttt{intensity\_error\_bound}} به این معناست که در صورت عدم تنظیم صریح، فرض می‌شود منبع ایده‌آل است ($\delta = 0$).

\subsection{اعتبارسنجی پیکرببری در \lr{Wang2005Proof}}

کلاس \lr{\texttt{Wang2005Proof}} در سازنده خود، پیکرببری منبع را اعتبارسنجی می‌کند تا اطمینان حاصل شود که سه نوع پالس \lr{signal}، \lr{decoy} و \lr{vacuum} همگی موجود هستند:

\begin{LTR}
	\begin{lstlisting}[language=Python]
		def _verify_pulse_configs(self):
		required = {"signal", "decoy", "vacuum"}
		current = {p.name for p in self.p.source.pulse_configs}
		if not required.issubset(current):
		raise ConfigurationError(
		f"Wang2005Proof requires pulse types: {required}. Found: {current}"
		)
	\end{lstlisting}
\end{LTR}

این اعتبارسنجی در زمان ساخت شیء انجام می‌شود (نه در زمان محاسبه) تا خطا در همان مراحل اولیه شناسایی شود. اگر پیکرببری ناقص باشد، استثنای \lr{\texttt{ConfigurationError}} \lr{raise} می‌شود که یک استثنای سفارشی در ماژول \lr{\texttt{qkd.exceptions}} است. این طراحی \lr{fail-fast} از اجرای طولانی شبیه‌سازی و شکست در مراحل بعدی جلوگیری می‌کند. استفاده از \lr{set} برای مقایسه مورد نیاز و موجود، ترتیب پالس‌ها را نادیده می‌گیرد و فقط حضور آن‌ها را بررسی می‌کند.

\subsection{مدیریت خطای تطبیقی (\lr{Adaptive Error Correction})}

یکی از ویژگی‌های پیشرفته کلاس پایه \lr{\texttt{Ma2005BaseProof}}، پشتیبانی از مدل خطای تطبیقی است. متد \lr{\texttt{\_calculate\_f\_ec}} بر اساس نرخ خطای مشاهده‌شده، بازدهی تصحیح خطا را به صورت پویا تنظیم می‌کند:

\begin{LTR}
	\begin{lstlisting}[language=Python]
		def _calculate_f_ec(self, error_rate: float) -> float:
		"""Calculates adaptive error correction efficiency if configured."""
		config = self.p.f_ec_dynamic_config
		if config and config.get("model") == "linear":
		base = config.get("base", 1.1)
		slope = config.get("slope", 0.0)
		return base + slope * error_rate
		return self.p.f_error_correction
	\end{lstlisting}
\end{LTR}

در این کد، اگر پیکرببری \lr{\texttt{f\_ec\_dynamic\_config}} تنظیم شده باشد و مدل آن \lr{linear} باشد، بازدهی به صورت $f_{\text{EC}} = \text{base} + \text{slope} \cdot E_\mu$ محاسبه می‌شود. این مدل‌سازی برای شبیه‌سازی واقع‌بینانه‌تر از سیستم‌های واقعی ضروری است زیرا در عمل، کدهای تصحیح خطا در نرخ‌های خطای بالا کارایی کمتری دارند. مقادیر پیش‌فرض \lr{base = 1.1} و \lr{slope = 0.0} معادل یک مدل ثابت است که در صورت عدم تنظیم پیکرببری پویا، استفاده می‌شود. این طراحی انعطاف‌پذیر اجازه می‌دهد کاربر بدون تغییر کد، مدل‌های مختلف تصحیح خطا را آزمایش کند.

\subsection{نقشه‌ی نمادها (\lr{Notation Map})}

هر دو ماژول متد \lr{\texttt{notation\_map}} را پیاده‌سازی می‌کنند که یک دیکشنری از نمادهای ریاضی به نام متغیرهای کد ارائه می‌دهد. این نقشه برای ترجمه بین فرمول‌های مقاله و کد، و برای تولید گزارش‌های قابل فهم توسط انسان استفاده می‌شود. به عنوان مثال، نقشه \lr{Ma 2005} شامل موارد زیر است:

\begin{LTR}
	\begin{lstlisting}[language=Python]
		def notation_map(self) -> Dict[str, str]:
		return {
			"Y_1": "yield_1_lower_bound",
			"e_1": "error_rate_1_upper_bound",
			"Q_mu": "Gain of signal state",
			"E_mu": "QBER of signal state",
			"\\Delta": "tagged_fraction (Delta)",
			"f(e)": "f_error_correction",
			"Y_1^L": "yield_1_lower_bound",
			"e_ph": "error_rate_1_upper_bound",
			"s_z_1^L": "Calculated as Q_1 * N_sent (single-photon count)"
		}
	\end{lstlisting}
\end{LTR}

این نقشه به ابزارهای بصری‌سازی و گزارش‌گیری اجازه می‌دهد به طور خودکار نمادهای ریاضی را در خروجی‌ها نمایش دهند و در عین حال، در کد از نام‌های توصیفی استفاده کنند. این جداسازی لایه نمایش از لایه محاسبه، از اصول کلیدی مهندسی نرم‌افزار است و قابلیت نگهداری کد را به طور قابل توجهی افزایش می‌دهد.

\section{پیاده‌سازی ریاضی و الگوریتمی}

\subsection{محاسبه نرخ \lr{GLLP} استاندارد}

متد \lr{\texttt{\_calculate\_gllp\_rate}} در کلاس پایه، هسته اصلی محاسبه نرخ کلید در چارچوب \lr{Ma 2005} است. این متد دو حالت دارد: حالت استاندارد \lr{GLLP} و حالت \lr{Weak GLLP}. در ادامه، نسخه کامل این متد را با توضیح خط به خط ارائه می‌دهیم:

\begin{LTR}
	\begin{lstlisting}[language=Python]
		def _calculate_gllp_rate(
		self,
		decoy_estimates: DecoyEstimates,
		stats_map: Dict[str, TallyCounts]
		) -> KeyCalculationResult:
		signal_stats = stats_map.get("signal")
		if not signal_stats or signal_stats.sent == 0:
		return KeyCalculationResult(0, 0.0, 0.0, 0.5,
		error_codes=[ErrorCode.INSUFFICIENT_STATISTICS])
		
		mu = self.p.source.get_pulse_config_by_name("signal").mean_photon_number
		q_factor = 0.5
		
		Q_mu = signal_stats.sifted / signal_stats.sent
		E_mu = signal_stats.errors_sifted / signal_stats.sifted \
		if signal_stats.sifted > 0 else 0.5
		f_ec = self._calculate_f_ec(E_mu)
		
		Y1_L = decoy_estimates.yield_1_lower_bound
		e1_U = decoy_estimates.error_rate_1_upper_bound
		
		prob_1_mu = self._get_prob_n(mu, 1)
		Q_1 = Y1_L * prob_1_mu
	\end{lstlisting}
\end{LTR}

در خطوط ابتدایی، آمار سیگنال از \lr{\texttt{stats\_map}} استخراج می‌شود. اگر آمار موجود نباشد یا تعداد پالس‌های ارسالی صفر باشد، یک نتیجه با طول کلید صفر و کد خطای \lr{\texttt{INSUFFICIENT\_STATISTICS}} بازگردانده می‌شود. سپس شدت سیگنال $\mu$ از پیکرببری خوانده می‌شود و \lr{$q$-factor} برابر \lr{0.5} (برای \lr{BB84}) تنظیم می‌شود. مقادیر قابل مشاهده $Q_\mu$ و $E_\mu$ از آمار محاسبه می‌شوند. در صورتی که \lr{sifted = 0} (هیچ آشکارسازی در پایه درست نباشد)، $E_\mu = 0.5$ فرض می‌شود که بدترین حالت است. در نهایت، $Q_1 = Y_1 \cdot P(1 \mid \mu)$ محاسبه می‌شود که \lr{Gain} تک‌فوتونی است.

در ادامه، کد بررسی می‌کند که آیا حالت \lr{Weak GLLP} فعال است یا خیر:

\begin{LTR}
	\begin{lstlisting}[language=Python]
		weak_rate = self._calculate_weak_gllp_rate(decoy_estimates, stats_map)
		delta_bound = decoy_estimates.diagnostics.numeric_diagnostics.get(
		"Delta_tagged_bound_Eq36", None)
		
		if getattr(self.p, 'use_weak_gllp', False):
		key_len = max(0, int(weak_rate * signal_stats.sent))
		return KeyCalculationResult(
		secure_key_length=key_len,
		privacy_amplification_term=0.0,
		error_correction_leakage=Q_mu * f_ec *
		binary_entropy(E_mu) * signal_stats.sent,
		phase_error_rate_upper_bound=E_mu / (1.0 - delta_bound)
		if delta_bound and delta_bound < 1.0 else 0.5,
		diagnostics={"mode": "WEAK_GLLP_EQ_43",
			"rate_per_pulse": weak_rate, "Delta": delta_bound},
		weak_gllp_rate=weak_rate,
		tagged_fraction_bound=delta_bound
		)
	\end{lstlisting}
\end{LTR}

در این بخش، ابتدا نرخ \lr{Weak GLLP} محاسبه می‌شود (حتی اگر استفاده نشود، برای گزارش‌گیری \lr{diagnostics} مفید است). سپس اگر پرچم \lr{\texttt{use\_weak\_gllp}} فعال باشد، نتیجه بر اساس فرمول \lr{Weak GLLP} (معادله 43) ساخته می‌شود. طول کلید به صورت $N_{\text{sent}} \cdot R_{\text{weak}}$ محاسبه می‌شود و کران بالای خطای فاز به صورت $E_\mu / (1 - \Delta)$ تنظیم می‌شود. این تناظر با فرمول \lr{Weak GLLP} که در بخش تئوری توضیح داده شد، دقیق است.

در غیر این صورت، مسیر استاندارد \lr{GLLP} اجرا می‌شود:

\begin{LTR}
	\begin{lstlisting}[language=Python]
		h2_E_mu = binary_entropy(E_mu)
		term_ec_cost = Q_mu * f_ec * h2_E_mu
		
		h2_e1 = binary_entropy(e1_U)
		term_pa_gain = Q_1 * (1.0 - h2_e1)
		
		rate = q_factor * (term_pa_gain - term_ec_cost)
		key_len = max(0, int(rate * signal_stats.sent))
		
		diag = {
			"Q_mu": Q_mu, "E_mu": E_mu, "Q_1": Q_1,
			"Y1_L": Y1_L, "e1_U": e1_U, "rate_per_pulse": rate,
			"f_ec_used": f_ec, "weak_gllp_rate": weak_rate,
			"asym_retry_used": False,  # always False after paper-faithful removal
		}
		if decoy_estimates.diagnostics.numeric_diagnostics:
		diag.update(decoy_estimates.diagnostics.numeric_diagnostics)
		
		return KeyCalculationResult(
		secure_key_length=key_len,
		privacy_amplification_term=term_pa_gain * signal_stats.sent,
		error_correction_leakage=term_ec_cost * signal_stats.sent,
		phase_error_rate_upper_bound=e1_U,
		diagnostics=diag,
		weak_gllp_rate=weak_rate,
		tagged_fraction_bound=delta_bound
		)
	\end{lstlisting}
\end{LTR}

در این مسیر، ابتدا آنتروپی باینری $H_2(E_\mu)$ محاسبه می‌شود و ترم نشت تصحیح خطا به صورت $Q_\mu \cdot f_{\text{EC}} \cdot H_2(E_\mu)$ به دست می‌آید. سپس آنتروپی باینری $H_2(e_1^U)$ محاسبه می‌شود و ترم تقویت حریم خصوصی به صورت $Q_1 \cdot (1 - H_2(e_1^U))$ محاسبه می‌شود. نرخ نهایی برابر $q \cdot (\text{PA} - \text{EC})$ است و طول کلید به صورت $\max(0, \lfloor R \cdot N_{\text{sent}} \rfloor)$ محاسبه می‌شود. تابع \lr{\texttt{max}} تضمین می‌کند که طول کلید هرگز منفی نشود (که در صورت $\text{PA} < \text{EC}$ ممکن است رخ دهد). در نهایت، یک دیکشنری \lr{diagnostics} غنی شامل تمام مقادیر میانی ساخته می‌شود و در شیء نتیجه قرار می‌گیرد. یک فیلد مهم در این دیکشنری \lr{\texttt{asym\_retry\_used}} است که پس از حذف منطق \lr{asymptotic retry} غیروفادارانه به مقاله، همواره مقدار \lr{False} دارد؛ این فیلد برای حفظ سازگاری با کدهای گزارش‌گیری موجود و برای انطباق آینده در صورت فعال‌سازی مجدد نگه داشته شده است.

\subsection{پیاده‌سازی پروتکل \lr{Vacuum + Weak}}

کلاس \lr{\texttt{Ma2005VacuumWeakProof}} پروتکل کلاسیک \lr{Vacuum + Weak} را با بهبود مقاومت در برابر نوسان شدت پیاده‌سازی می‌کند. در ادامه، بخش کلیدی متد \lr{\texttt{estimate\_yields\_and\_errors}} این کلاس را بررسی می‌کنیم:

\begin{LTR}
	\begin{lstlisting}[language=Python]
		def estimate_yields_and_errors(self, stats_map: Dict[str, TallyCounts]) -> DecoyEstimates:
		stats_map = self._normalize_stats_map(stats_map, self._get_pulse_names())
		if self.is_thermal:
		raise ConfigurationError("Ma2005VacuumWeakProof supports only Poisson statistics.")
		
		try:
		cfg_sig = self.p.source.get_pulse_config_by_name("signal")
		cfg_decoy = self.p.source.get_pulse_config_by_name("decoy")
		cfg_vac = self.p.source.get_pulse_config_by_name("vacuum")
		except Exception as e:
		return DecoyEstimates(0.0, 0.5, False, 1.0,
		SolverDiagnostics("Ma2005", False, str(e)))
		
		mu_base = cfg_sig.mean_photon_number
		nu_base = cfg_decoy.mean_photon_number
		
		delta = self.p.source.intensity_jitter
		n_total = sum(s.sent for s in stats_map.values())
		if n_total > 0:
		stat_fluc = 1.0 / math.sqrt(n_total)
		delta += stat_fluc
		
		mu_min, mu_max = mu_base * (1 - delta), mu_base * (1 + delta)
		nu_min, nu_max = nu_base * (1 - delta), nu_base * (1 + delta)
	\end{lstlisting}
\end{LTR}

در خطوط ابتدایی، ابتدا بررسی می‌شود که منبع از نوع \lr{Poisson} باشد (این کلاس از \lr{Thermal} پشتیبانی نمی‌کند). سپس سه پیکرببری پالس (سیگنال، \lr{decoy} و خلأ) از پیکرببری منبع استخراج می‌شوند. اگر هر کدام موجود نباشند، یک \lr{DecoyEstimates} با طول صفر و \lr{diagnostics} خطا بازگردانده می‌شود. سپس شدت پایه $\mu$ و $\nu$ خوانده می‌شوند و کران خطای منبع $\delta$ از پیکرببری استخراج می‌شود. یک نکته کلیدی در این کد، افزودن نوسان آماری به $\delta$ است: $\delta_{\text{total}} = \delta_{\text{jitter}} + 1/\sqrt{N_{\text{total}}}$ که در آن $N_{\text{total}}$ کل پالس‌های ارسالی است. این نوسان آماری مدل می‌کند که برای یک $N$ متناهی، توزیع فوتون‌های مشاهده‌شده دقیقاً با توزیع تئوریک همخوانی ندارد. در نهایت، کران‌های بالا و پایین شدت به صورت $\mu_{\min} = \mu(1 - \delta)$ و $\mu_{\max} = \mu(1 + \delta)$ محاسبه می‌شوند.

سپس کران‌های آماری برای $Q_\mu$، $Q_\nu$ و $Q_{\text{vac}}$ محاسبه می‌شوند:

\begin{LTR}
	\begin{lstlisting}[language=Python]
		alpha = self.eps_alloc.eps_pe
		
		Q_mu_L, Q_mu_U, E_mu_L, E_mu_U = self._get_bounds_stats(
		stats_map, "signal", alpha)
		Q_nu_L, Q_nu_U, E_nu_L, E_nu_U = self._get_bounds_stats(
		stats_map, "decoy", alpha)
		Q_vac_L, Q_vac_U, E_vac_L, E_vac_U = self._get_bounds_stats(
		stats_map, "vacuum", alpha)
		
		Y0_U = Q_vac_U
		Y0_L = Q_vac_L
		e0 = 0.5
	\end{lstlisting}
\end{LTR}

در این بخش، $\alpha = \varepsilon_{\text{PE}}$ به عنوان سطح اطمینان استفاده می‌شود و کران‌های دوطرفه برای $Q$ و $E$ در هر سه شدت محاسبه می‌شوند. بازده خلأ $Y_0$ به صورت مستقیم از آمار \lr{vacuum} استخراج می‌شود: $Y_0^U = Q_{\text{vac}}^U$ و $Y_0^L = Q_{\text{vac}}^L$. این یک ویژگی متمایزکننده پروتکل \lr{Vacuum + Weak} است که در آن یک \lr{decoy} خلأ به طور مستقیم $Y_0$ را اندازه می‌گیرد. $e_0 = 0.5$ ثابت فرض می‌شود زیرا \lr{dark count} کاملاً تصادفی است.

سپس فرمول کران پایین $Y_1$ با کران‌های بدترین حالت اعمال می‌شود:

\begin{LTR}
	\begin{lstlisting}[language=Python]
		# [Paper Eq 34] Lower Bound on Y1
		term_Q_nu = Q_nu_L * math.exp(nu_min)
		term_Q_mu = Q_mu_U * math.exp(mu_max) * (nu_max**2 / mu_min**2)
		term_Y0 = Y0_U * (mu_max**2 - nu_min**2) / (mu_min**2)
		
		denom = mu_max * nu_min - nu_max**2
		if denom <= 0:
		prefactor = 0.0
		else:
		prefactor = mu_min / denom
		
		Y1_L = prefactor * (term_Q_nu - term_Q_mu - term_Y0)
		Y1_L = max(0.0, Y1_L)
		
		# [Paper Eq 37] Upper Bound on e1
		if Y1_L * nu_min <= 1e-20:
		e1_U = 0.5
		else:
		numerator_e1 = (E_nu_U * Q_nu_U * math.exp(nu_max)) - (e0 * Y0_L)
		e1_U = numerator_e1 / (Y1_L * nu_min)
		e1_U = min(0.5, max(0.0, e1_U))
		
		# [Paper Eq 36] Tagged Fraction Delta Upper Bound
		if Q_nu_L > 0:
		term1 = (nu_max / (mu_min - nu_max))
		term2 = (nu_max * math.exp(-nu_min) * Q_mu_U) / \\
		(mu_min * math.exp(-mu_max) * Q_nu_L) - 1.0
		term3 = (nu_max * math.exp(-nu_min) * Y0_U) / (mu_min * Q_nu_L)
		delta_bound_strict = term1 * term2 + term3
		else:
		delta_bound_strict = 1.0
		
		# Asymptotic Reference
		eta = self.p.channel.transmittance * self.p.detector.det_eff_d0
		y0_val = 2.0 * self.p.detector.dark_rate   # Gap [B] fix: two-detector BB84
		e_det = self.p.detector.qber_intrinsic
		
		y1_asymptotic = y0_val + eta
		e1_asymptotic = (0.5 * y0_val + e_det * eta) / y1_asymptotic \\
		if y1_asymptotic > 0 else 0.5
		
		beta_y1 = (y1_asymptotic - Y1_L) / y1_asymptotic \\
		if y1_asymptotic > 0 else 0.0
		beta_e1 = (e1_U - e1_asymptotic) / e1_asymptotic \\
		if e1_asymptotic > 0 else 0.0
		
		diag_data = {
			"Y0_L": Y0_L, "Y0_U": Y0_U,
			"Delta_tagged_bound_Eq36": delta_bound_strict,
			"beta_Y1_deviation": beta_y1,
			"beta_e1_deviation": beta_e1,
			"used_asymptotic_fallback": False,  # always False after paper-faithful removal
			"Y1_asymptotic": y1_asymptotic,
			"e1_asymptotic": e1_asymptotic,
		}
		
		return DecoyEstimates(
		yield_1_lower_bound=Y1_L,
		error_rate_1_upper_bound=e1_U,
		is_feasible=True,
		failure_prob_used=0.0,
		diagnostics=SolverDiagnostics("Ma2005_VacWeak_Robust", True, "OK",
		numeric_diagnostics=diag_data)
		)
	\end{lstlisting}
\end{LTR}

این پیاده‌سازی معادلات (34)، (37) و (36) مقاله را با کران‌های بدترین حالت اعمال می‌کند. در ادامه، بخش‌های کلیدی این کد را به تفصیل توضیح می‌دهیم:

\textbf{کران پایین $Y_1$ (معادله 34):} هر شدت در هر جایگاه فرمول با کران مناسب (بالا یا پایین) جایگزین شده است تا کران نهایی تا حد ممکن سخت‌گیرانه باشد. به طور خاص، $Q_\nu$ با کران پایین $Q_\nu^L$ و $Q_\mu$ با کران بالا $Q_\mu^U$ جایگزین می‌شوند و $Y_0$ با کران بالا $Y_0^U$. اگر مخرج $\mu_{\max} \nu_{\min} - \nu_{\max}^2 \le 0$ باشد (که در پیکرببری‌های غیرفیزیکی ممکن است رخ دهد)، ضریب پیشانی به صفر تنظیم می‌شود و در نتیجه $Y_1 = 0$ می‌شود. در نهایت، $\max(0, \cdot)$ تضمین می‌کند که $Y_1$ هرگز منفی نشود.

\textbf{کران بالای $e_1$ (معادله 37):} این فرمول کران بالای خطای تک‌فوتونی را بر اساس مشاهدات \lr{decoy} ضعیف محاسبه می‌کند. صورت کسر شامل $E_\nu^U \cdot Q_\nu^U \cdot e^{\nu_{\max}}$ (با کران‌های بالا برای محافظه‌کاری) منهای $e_0 \cdot Y_0^L$ است. مخرج شامل $Y_1^L \cdot \nu_{\min}$ است که با کران‌های پایین محافظه‌کارانه انتخاب شده. اگر $Y_1^L \cdot \nu_{\min} \le 10^{-20}$ باشد (تقریباً صفر)، $e_1 = 0.5$ تنظیم می‌شود که معادل بدترین حالت (کامل تصادفی) است. در نهایت، $e_1$ با $\min(0.5, \max(0, \cdot))$ در بازه $[0, 0.5]$ محدود می‌شود.

\textbf{کران بالای کسر برچسب‌گذاری‌شده $\Delta$ (معادله 36):} این کران نشان می‌دهد چه کسری از پالس‌های سیگنال می‌توانند چندفوتونی باشند. اگر $Q_\nu^L > 0$ باشد، فرمول کامل اعمال می‌شود؛ در غیر این صورت، $\Delta = 1.0$ تنظیم می‌شود که به معنای عدم امکان تخمین است. این مقدار $\Delta$ در دیکشنری \lr{diagnostics} با کلید \lr{\texttt{Delta\_tagged\_bound\_Eq36}} ذخیره می‌شود و توسط متد \lr{\texttt{\_calculate\_weak\_gllp\_rate}} برای محاسبه نرخ \lr{Weak GLLP} استفاده می‌شود.

\textbf{مرجع حد \lr{Asymptotic}:} پس از محاسبه کران‌های متناهی، کد یک مرجع \lr{asymptotic} محاسبه می‌کند تا انحراف کران متناهی از حد ایده‌آل قابل اندازه‌گیری باشد. در این محاسبات، $\eta$ از پارامترهای کانال و آشکارساز خوانده می‌شود ($\eta = \eta_{\text{channel}} \cdot \eta_{\text{det}}$) و $y_0$ به صورت $2 \cdot \text{dark\_rate}$ تنظیم می‌شود که \lr{Gap [B] fix} برای پشتیبانی از \lr{BB84} دو آشکارساز است. در \lr{BB84} استاندارد، آلیس در دو پایه $Z$ و $X$ ارسال می‌کند و باب نیز در دو پایه اندازه‌گیری می‌کند، بنابراین نرخ کلی \lr{dark count} مؤثر دو برابر نرخ تک‌آشکارساز است. سپس $Y_1^{\text{asym}} = y_0 + \eta$ و $e_1^{\text{asym}} = (0.5 \cdot y_0 + e_{\text{det}} \cdot \eta) / Y_1^{\text{asym}}$ محاسبه می‌شوند.

\textbf{ضریب انحراف $\beta$:} دو کمیت تشخیصی $\beta_{Y_1}$ و $\beta_{e_1}$ به صورت زیر محاسبه می‌شوند:
\begin{equation}
	\beta_{Y_1} = \frac{Y_1^{\text{asym}} - Y_1^L}{Y_1^{\text{asym}}}, \qquad
	\beta_{e_1} = \frac{e_1^U - e_1^{\text{asym}}}{e_1^{\text{asym}}}.
\end{equation}
این ضرایب نشان می‌دهند که کران متناهی چقدر از حد ایده‌آل فاصله دارد. اگر $\beta_{Y_1}$ نزدیک به صفر باشد، یعنی کران متناهی تقریباً به حد ایده‌آل رسیده است و اندازه بلوک کافی است. اگر $\beta_{Y_1}$ بزرگ باشد، یعنی افزایش اندازه بلوک می‌تواند کران را به طور قابل توجهی بهبود بخشد. این اطلاعات برای بهینه‌سازی طراحی آزمایش حیاتی است. فیلد \lr{\texttt{used\_asymptotic\_fallback}} همواره \lr{False} است زیرا پس از حذف منطق \lr{fallback} غیروفادارانه به مقاله، کد همواره مسیر تخمین اصلی را طی می‌کند و فقط مقادیر \lr{asymptotic} را برای مقایسه در \lr{diagnostics} گزارش می‌کند.

\subsection{پیاده‌سازی پروتکل \lr{General Two-Decoy}}

کلاس \lr{\texttt{Ma2005GeneralTwoDecoyProof}} پروتکل دو \lr{decoy} تعمیم‌یافته را پیاده‌سازی می‌کند که بر خلاف \lr{Vacuum + Weak}، فرض خلأ بودن یکی از \lr{decoy}‌ها را ندارد. این کلاس نیز مانند سایر کلاس‌های \lr{Ma 2005}، در ابتدای متد \lr{\texttt{estimate\_yields\_and\_errors}} فراخوانی \lr{\texttt{\_normalize\_stats\_map}} را انجام می‌دهد تا کلیدهای دیکشنری ورودی به فرمت استاندارد نام پالس تبدیل شوند. سپس بررسی می‌کند که منبع از نوع \lr{Poisson} باشد (این کلاس از \lr{Thermal} پشتیبانی نمی‌کند) و در صورت عدم تطابق، استثنای \lr{\texttt{ConfigurationError}} \lr{raise} می‌شود. در ادامه، دو پیکرببری \lr{decoy} با کوچکترین و بزرگترین شدت انتخاب می‌شوند ($\nu_2 < \nu_1$). این کلاس از معادلات (18)، (21) و (25) مقاله استفاده می‌کند:

\begin{LTR}
	\begin{lstlisting}[language=Python]
		# Eq 18: Y0 lower bound
		term_num = nu1 * Q_nu2_L * math.exp(nu2) - nu2 * Q_nu1_U * math.exp(nu1)
		Y0_L = max(0.0, term_num / (nu1 - nu2))
		
		# Eq 21: Y1 lower bound
		denom_Y1 = (mu * nu1) - (mu * nu2) - (nu1**2) + (nu2**2)
		term_Q_diff = Q_nu1_L * math.exp(nu1) - Q_nu2_U * math.exp(nu2)
		term_mu_sub = (Q_mu_U * math.exp(mu) - Y0_L)
		factor_sq = (nu1**2 - nu2**2) / (mu**2)
		
		Y1_L = max(0.0, (mu / denom_Y1) * (term_Q_diff - (factor_sq * term_mu_sub)))
		
		# Eq 25: e1 upper bound
		if Y1_L * (nu1 - nu2) <= 1e-20:
		e1_U = 0.5
		else:
		num_e1 = (E_nu1_U * Q_nu1_U * math.exp(nu1)) - (E_nu2_L * Q_nu2_L * math.exp(nu2))
		e1_U = min(0.5, max(0.0, num_e1 / ((nu1 - nu2) * Y1_L)))
	\end{lstlisting}
\end{LTR}

این پیاده‌سازی چند ویژگی مهم دارد: نخست، $Y_0$ به جای اندازه‌گیری مستقیم (که در \lr{Vacuum + Weak} با \lr{decoy} خلأ انجام می‌شد)، از ترکیب دو \lr{decoy} ضعیف تخمین زده می‌شود (معادله 18). این تخمین با استفاده از کران پایین $Q_{\nu_2}$ و کران بالا $Q_{\nu_1}$ انجام می‌شود تا کران پایین $Y_0$ محافظه‌کارانه باشد. دوم، در معادله (21)، $Y_0^L$ به جای $Y_0$ استفاده می‌شود تا جمله کم‌شونده بزرگ‌تر و کران $Y_1$ سخت‌گیرانه‌تر شود. سوم، در معادله (25)، $E_{\nu_1} Q_{\nu_1}$ با کران بالا و $E_{\nu_2} Q_{\nu_2}$ با کران پایین جایگزین می‌شوند تا صورت کسر بزرگ‌تر و کران $e_1$ سخت‌گیرانه‌تر شود. این انتخاب‌های جهت‌دار در سراسر کد، تضمین می‌کنند که تمام کران‌ها در بدترین حالت نوسان آماری معتبر باشند.

\subsection{پیاده‌سازی پروتکل \lr{One-Decoy}}

کلاس \lr{\texttt{Ma2005OneDecoyProof}} پروتکل تک‌‌‌decoy را با دو حالت (\lr{zero-background} یا \lr{standard}) پیاده‌سازی می‌کند. مانند سایر کلاس‌های \lr{Ma 2005}، این کلاس نیز در ابتدای متد \lr{\texttt{estimate\_yields\_and\_errors}} فراخوانی \lr{\texttt{\_normalize\_stats\_map}} را انجام می‌دهد تا کلیدهای دیکشنری ورودی به فرمت استاندارد نام پالس تبدیل شوند. سپس بررسی می‌کند که منبع از نوع \lr{Poisson} باشد (این کلاس از \lr{Thermal} پشتیبانی نمی‌کند) و در صورت عدم تطابق، استثنای \lr{\texttt{ConfigurationError}} \lr{raise} می‌شود. در حالت \lr{zero-background}، فرض می‌شود $Y_0 = 0$ که یک تقریب محافظه‌کارانه برای کانال‌های با \lr{dark count} بسیار پایین است:

\begin{LTR}
	\begin{lstlisting}[language=Python]
		if self.assume_zero_background:
		# Eq. 41: Tighter Bound (Y0=0 implicitly)
		prefactor = mu / (mu * nu - nu**2)
		term_Q_nu = Q_nu_L * math.exp(nu)
		term_Q_mu = Q_mu_U * math.exp(mu) * (nu**2 / mu**2)
		Y1_L = max(0.0, prefactor * (term_Q_nu - term_Q_mu))
		
		if Y1_L * nu <= 1e-20:
		e1_U = 0.5
		else:
		e1_U = min(0.5, max(0.0,
		(E_nu_U * Q_nu_U * math.exp(nu)) / (Y1_L * nu)))
		diag_msg = "Ma2005_OneDecoy_ZeroBackground"
	\end{lstlisting}
\end{LTR}

در حالت \lr{zero-background}، فرمول کران $Y_1$ ساده‌تر می‌شود زیرا جمله $Y_0$ حذف می‌شود. این فرمول معادله (41) مقاله را پیاده‌سازی می‌کند. در حالت استاندارد (با $Y_0$ غیر صفر)، فرمول کامل معادلات (38) تا (40) استفاده می‌شود که در آن $Y_0$ از روی $E_\mu Q_\mu$ تخمین زده می‌شود:

\begin{LTR}
	\begin{lstlisting}[language=Python]
		else:
		# Standard One-Decoy Bound (Eq. 38-40)
		Y0_U = (E_mu_U * Q_mu_U * math.exp(mu)) / e0
		
		term1 = Q_nu_L * math.exp(nu)
		term2 = Q_mu_U * math.exp(mu) * (nu**2 / mu**2)
		term3 = (Y0_U) * ((mu**2 - nu**2) / mu**2)
		Y1_L = max(0.0, (mu / (mu * nu - nu**2)) * (term1 - term2 - term3))
		
		if Y1_L * mu <= 1e-20:
		e1_U = 0.5
		else:
		e1_U = min(0.5, max(0.0,
		(E_mu_U * Q_mu_U * math.exp(mu)) / (Y1_L * mu)))
		diag_msg = "Ma2005_OneDecoy_Simple"
	\end{lstlisting}
\end{LTR}

در حالت استاندارد، $Y_0^U$ از روی خطای سیگنال تخمین زده می‌شود: $Y_0^U = E_\mu^U Q_\mu^U e^\mu / e_0$. این تخمین بر اساس این فرض است که تمام خطاهای مشاهده‌شده در سیگنال از \lr{dark count} ناشی می‌شوند (که یک فرض محافظه‌کارانه است). سپس این $Y_0^U$ در فرمول کران $Y_1$ استفاده می‌شود.

\subsection{پیاده‌سازی پروتکل چند \lr{Decoy} با \lr{LP}}

کلاس \lr{\texttt{Ma2005GeneralMultiDecoyProof}} برای حالت چند \lr{decoy} (بیش از دو) از برنامه‌ریزی خطی (\lr{LP}) برای یافتن کران‌های بهینه $Y_1$ و $e_1$ استفاده می‌کند. این تنها کلاس در ماژول است که از توزیع \lr{Thermal} پشتیبانی می‌کند. مانند سایر کلاس‌های \lr{Ma 2005}، این کلاس نیز در ابتدای متد \lr{\texttt{estimate\_yields\_and\_errors}} فراخوانی \lr{\texttt{\_normalize\_stats\_map}} را انجام می‌دهد. سپس ماتریس \lr{LP} با استفاده از تمام پیکرببری‌های پالس موجود ساخته می‌شود. مسأله \lr{LP} به صورت زیر فرمول‌بندی می‌شود:

\begin{equation}
	\min_{\{Y_n\}} Y_1 \quad \text{s.t.} \quad Q_k^L \le \sum_{n=0}^{N_{\text{cap}}} Y_n\,P(n \mid \mu_k) \le Q_k^U, \quad \forall k.
\end{equation}

در این فرمول‌بندی، $\{Y_n\}_{n=0}^{N_{\text{cap}}}$ متغیرهای تصمیم هستند، $Q_k^L$ و $Q_k^U$ کران‌های آماری \lr{Gain} در شدت $k$، و $P(n \mid \mu_k)$ احتمال ارسال $n$ فوتون با شدت $\mu_k$ است. این مسأله با استفاده از تابع \lr{\texttt{solve\_lp}} از ماژول \lr{\texttt{qkd.proofs.utils\_lp}} حل می‌شود. در ادامه، کد کلیدی این کلاس را ارائه می‌دهیم:

\begin{LTR}
	\begin{lstlisting}[language=Python]
		c_y1 = np.zeros(n_cap + 1)
		c_y1[1] = 1.0  # Minimize Y_1
		
		A_eq = np.zeros((k_intensities, n_cap + 1))
		b_eq = np.zeros(k_intensities)
		b_eq_L = np.zeros(k_intensities)
		b_eq_U = np.zeros(k_intensities)
		
		valid_indices = []
		for i, pc in enumerate(pulse_configs):
		if pc.name not in stats_map or stats_map[pc.name].sent == 0:
		continue
		valid_indices.append(i)
		mu = pc.mean_photon_number
		
		Q_L, Q_U, _, _ = self._get_bounds_stats(stats_map, pc.name, alpha)
		b_eq_L[i] = Q_L
		b_eq_U[i] = Q_U
		
		for n in range(n_cap + 1):
		A_eq[i, n] = self._get_prob_n(mu, n)
	\end{lstlisting}
\end{LTR}

در این کد، بردار هزینه $c$ به صورتی تنظیم می‌شود که $Y_1$ مینیمم شود (یعنی $c = [0, 1, 0, \ldots, 0]$). سپس ماتریس محدودیت $A_{\text{eq}}$ ساخته می‌شود که در آن سطر $i$ معادله $\sum_n Y_n P(n \mid \mu_i) = Q_i$ را نشان می‌دهد. کران‌های بالا و پایین $Q_i$ به صورت دو مجموعه از محدودیت‌های نامساوی اعمال می‌شوند: $\sum \le Q^U$ و $-\sum \le -Q^L$. این تبدیل استاندارد، مسأله با محدودیت‌های دوطرفه را به مسأله با محدودیت‌های یکطرفه تبدیل می‌کند که توسط \lr{solver}‌های استاندارد \lr{LP} قابل حل است.

سپس مسأله \lr{LP} حل می‌شود و برای $e_1$ نیز یک مسأله مشابه حل می‌شود:

\begin{LTR}
	\begin{lstlisting}[language=Python]
		A_ub = np.vstack([A_eq, -A_eq])
		b_ub = np.hstack([b_eq_U, -b_eq_L])
		
		sol_y1, diag_y1 = solve_lp(c_y1, A_ub, b_ub, n_cap + 1,
		self.p.lp_solver_method)
		Y1_L = max(0.0, sol_y1[1])
		
		# For e1: minimize -e1*Y1 (which maximizes e1*Y1)
		c_ey1 = np.zeros(n_cap + 1)
		c_ey1[1] = -1.0
		
		# ... build A_ub_err, b_ub_err similar to above ...
		
		sol_ey1, diag_ey1 = solve_lp(c_ey1, A_ub_err, b_ub_err, n_cap + 1,
		self.p.lp_solver_method)
		e1Y1_U = max(0.0, sol_ey1[1])
		
		if Y1_L <= 1e-20:
		e1_U = 0.5
		else:
		e1_U = min(0.5, e1Y1_U / Y1_L)
	\end{lstlisting}
\end{LTR}

در این بخش، مسأله \lr{LP} برای $Y_1$ حل می‌شود و سپس یک مسأله \lr{LP} دیگر برای $e_1 Y_1$ حل می‌شود (با مینیمم‌کردن $-e_1 Y_1$ که معادل ماکزیمم‌کردن $e_1 Y_1$ است). سپس $e_1 = (e_1 Y_1) / Y_1$ محاسبه می‌شود. این رویکرد دو مرحله‌ای ضروری است زیرا $e_1$ به صورت نسبتی تعریف می‌شود و نمی‌تواند مستقیماً در \lr{LP} مینیمم شود. اگر $Y_1 \le 10^{-20}$ باشد، $e_1 = 0.5$ تنظیم می‌شود که معادل بدترین حالت است. این طراحی اجازه می‌دهد کاربر هر تعداد \lr{decoy} را بدون محدودیت فرمول‌های تحلیلی استفاده کند.

\subsection{پیاده‌سازی حد \lr{Asymptotic}}

کلاس \lr{\texttt{Ma2005AsymptoticProof}} حد نظری \lr{Infinite Decoy State} را با معادلات (27) و (28) پیاده‌سازی می‌کند. این کلاس به عنوان مرجع بالای عملکرد عمل می‌کند:

\begin{LTR}
	\begin{lstlisting}[language=Python]
		# 1. Estimate Q_mu
		Q_mu = signal_stats.sifted / signal_stats.sent
		
		# 2. Derive eta from Q_mu
		Y0 = 2.0 * self.p.detector.dark_rate   # Gap [B] fix: two-detector BB84
		e0 = 0.5
		e_det = self.p.detector.qber_intrinsic
		
		if "vacuum" in stats_map and stats_map["vacuum"].sent > 0:
		vac_stats = stats_map["vacuum"]
		Y0 = vac_stats.sifted / vac_stats.sent
		
		mu = self.p.source.get_pulse_config_by_name("signal").mean_photon_number
		
		# Solve Q_mu = Y0 + 1 - exp(-eta * mu) for eta
		arg = 1.0 + Y0 - Q_mu
		if arg <= 0:
		eta = self.p.channel.transmittance * self.p.detector.det_eff_d0
		else:
		eta_effective = -math.log(arg) / mu
		eta = max(0.0, min(1.0, eta_effective))
		
		# 3. Calculate Asymptotic Limits (Eq 27 and 28)
		Y1_asymptotic = Y0 + eta
		
		if Y1_asymptotic > 0:
		numerator = (e0 * Y0) + (e_det * eta)
		e1_asymptotic = numerator / Y1_asymptotic
		else:
		e1_asymptotic = 0.5
		
		e1_asymptotic = min(0.5, max(0.0, e1_asymptotic))
	\end{lstlisting}
\end{LTR}

در این کد، ابتدا $Q_\mu$ از آمار محاسبه می‌شود. سپس $Y_0 = 2 \cdot \text{dark\_rate}$ تنظیم می‌شود که \lr{Gap [B] fix} برای پشتیبانی از \lr{BB84} دو آشکارساز است. در \lr{BB84} استاندارد، باب از دو آشکارساز (یکی برای پایه $Z$ و دیگری برای پایه $X$) استفاده می‌کند و هر دو آشکارساز در صورت عدم تطابق پایه، \lr{dark count} تولید می‌کنند. بنابراین نرخ مؤثر \lr{dark count} که در فرمول‌های تخمین ظاهر می‌شود، دو برابر نرخ تک‌آشکارساز است. عدم اعمال این ضریب 2 در نسخه‌های قبلی کد منجر به تخمین خوش‌بینانه $Y_1$ و در نتیجه طول کلید ناامن می‌شد. سپس $e_{\text{det}}$ از پیکرببری خوانده می‌شود. اگر آمار خلأ موجود باشد، $Y_0$ از آن استخراج می‌شود (که به واقعیت نزدیک‌تر است و مستقل از مقدار تئوریک \lr{dark\_rate} است). سپس $\eta$ به صورت غیرمستقیم از روی $Q_\mu$ با معکوس‌کردن رابطه $Q_\mu = Y_0 + 1 - e^{-\eta \mu}$ استخراج می‌شود. اگر $1 + Y_0 - Q_\mu \le 0$ (که ناهنجاری فیزیکی است و معمولاً ناشی از نویز بالا یا خطای کالیبراسیون است)، از مقدار تئوریک $\eta = \eta_{\text{channel}} \cdot \eta_{\text{det}}$ از پیکرببری به عنوان \lr{fallback} استفاده می‌شود. در نهایت، $Y_1$ و $e_1$ با معادلات (27) و (28) محاسبه می‌شوند. این مقادیر به عنوان کران‌های بالای تئوریک عمل می‌کنند و اجازه می‌دهند کاربر ببیند در شرایط ایده‌آل (تعداد نامتناهی \lr{decoy} و اندازه بلوک بی‌نهایت)، طول کلید چقدر می‌تواند باشد.

\subsection{پیاده‌سازی چارچوب \lr{Wang 2008}}

کلاس \lr{\texttt{Wang2005Proof}} چارچوب \lr{Wang 2008} را با مدیریت خطای منبع پیاده‌سازی می‌کند. در ادامه، بخش کلیدی متد \lr{\texttt{estimate\_yields\_and\_errors}} این کلاس را بررسی می‌کنیم:

\begin{LTR}
	\begin{lstlisting}[language=Python]
		delta = getattr(self.p.source, 'intensity_jitter', 0.0)
		source_fidelity = getattr(self.p.source, 'source_fidelity', 1.0)
		
		if source_fidelity < 1.0:
		delta = max(delta, 1.0 - source_fidelity)
		
		# Bounds on intensities
		mu_U = mu * (1 + delta)
		mu_L = mu * (1 - delta)
		nu_U = nu * (1 + delta)
		nu_L = nu * (1 - delta)
	\end{lstlisting}
\end{LTR}

در این بخش، $\delta$ از دو منبع ممکن است به دست آید: یا مستقیماً از \lr{\texttt{intensity\_jitter}} خوانده می‌شود، یا از طریق \lr{source\_fidelity} به صورت غیرمستقیم محاسبه می‌شود. اگر \lr{fidelity} $F < 1$ باشد، $\delta = \max(\delta_{\text{jitter}}, 1 - F)$ تنظیم می‌شود. این \lr{mapping} محافظه‌کارانه فرض می‌کند که انحراف از حالت \lr{coherent} ایده‌آل، حداقل به اندازه $(1 - F)$ در شدت مؤثر اثر دارد. سپس کران‌های بالا و پایین شدت به صورت $\mu^U = \mu(1 + \delta)$ و $\mu^L = \mu(1 - \delta)$ محاسبه می‌شوند. این کران‌ها در سراسر فرمول‌های تخمین $Y_1$ و $e_1$ استفاده می‌شوند تا در بدترین حالت نوسان منبع، امنیت حفظ شود.

سپس فرمول کران پایین $Y_1$ با کران‌های بدترین حالت اعمال می‌شود:

\begin{LTR}
	\begin{lstlisting}[language=Python]
		term1 = Q_nu * math.exp(nu_L)
		term2 = Q_mu * math.exp(mu_U) * (nu_U**2) / (mu_L**2)
		term3 = (mu_U**2 - nu_L**2)/(mu_L**2) * Q_vac  # Y0 approx Q_vac
		
		numerator = term1 - term2 - term3
		denominator = mu_U * nu_L - nu_U**2
		
		if denominator <= 0:
		Y1_L = 0.0
		else:
		Y1_L = (mu_L / denominator) * numerator
		
		Y1_L = max(0.0, Y1_L)
	\end{lstlisting}
\end{LTR}

این پیاده‌سازی معادله (56) مقاله \lr{Wang 2008} را با کران‌های بدترین حالت اعمال می‌کند. تفاوت کلیدی با فرمول \lr{Ma 2005} این است که تمام شدت‌ها با کران‌های بالا یا پایین جایگزین شده‌اند. اگر مخرج $\mu^U \nu^L - (\nu^U)^2 \le 0$ باشد، $Y_1 = 0$ تنظیم می‌شود. در نهایت، $\max(0, \cdot)$ تضمین می‌کند که $Y_1$ هرگز منفی نشود.

سپس کران بالای $e_1$ محاسبه می‌شود:

\begin{LTR}
	\begin{lstlisting}[language=Python]
		term_err = E_nu * Q_nu * math.exp(nu_U)
		term_vac = E_vac * Q_vac  # e0*Y0
		
		if Y1_L * nu_L <= 0:
		e1_U = 0.5
		else:
		e1_U = (term_err - term_vac) / (Y1_L * nu_L)
		
		e1_U = min(0.5, max(0.0, e1_U))
	\end{lstlisting}
\end{LTR}

در این فرمول، $e^{\nu}$ با $e^{\nu^U}$ (کران بالا) جایگزین شده تا صورت کسر بزرگ‌تر و کران $e_1$ سخت‌گیرانه‌تر شود. همچنین $Y_1$ و $\nu$ در مخرج با کران پایین $Y_1^L$ و $\nu^L$ جایگزین شده‌اند. اگر $Y_1^L \cdot \nu^L \le 0$ باشد، $e_1 = 0.5$ تنظیم می‌شود که معادل بدترین حالت است.

در نهایت، نرخ کلید با فرمول استاندارد \lr{GLLP} محاسبه می‌شود:

\begin{LTR}
	\begin{lstlisting}[language=Python]
		mu = self.p.source.get_pulse_config_by_name("signal").mean_photon_number
		p1 = mu * math.exp(-mu)
		
		Y1_L = decoy_estimates.yield_1_lower_bound
		e1_U = decoy_estimates.error_rate_1_upper_bound
		
		Q_mu = stats.sifted / n_total
		E_mu = stats.errors_sifted / stats.sifted if stats.sifted > 0 else 0.5
		
		gain_1 = Y1_L * p1
		
		from ..utils.math import binary_entropy
		pa_factor = 1.0 - binary_entropy(e1_U)
		secure_gain = gain_1 * pa_factor
		
		leakage = Q_mu * self.p.f_error_correction * binary_entropy(E_mu)
		
		rate = 0.5 * (secure_gain - leakage)  # q=0.5 for BB84
		
		key_len = max(0, int(rate * n_total))
	\end{lstlisting}
\end{LTR}

در این بخش، ابتدا $P(1 \mid \mu) = \mu e^{-\mu}$ محاسبه می‌شود (با فرض \lr{Poisson}). سپس \lr{Gain} تک‌فوتونی $Q_1 = Y_1^L \cdot P(1 \mid \mu)$ محاسبه می‌شود. ترم تقویت حریم خصوصی به صورت $Q_1 \cdot (1 - H_2(e_1^U))$ و ترم نشت تصحیح خطا به صورت $Q_\mu \cdot f_{\text{EC}} \cdot H_2(E_\mu)$ محاسبه می‌شوند. نرخ نهایی $R = 0.5 \cdot (\text{PA} - \text{EC})$ (با $q = 0.5$ برای \lr{BB84}) و طول کلید به صورت $\max(0, \lfloor R \cdot N \rfloor)$ محاسبه می‌شود.

\section{ارسال و دریافت با سایر ماژول‌ها}

\subsection{ورودی‌های ماژول \lr{ma2005.py}}

ماژول \lr{ma2005.py} برای انجام محاسبات خود به چندین ورودی از سایر ماژول‌های فریم‌ورک نیاز دارد. این ورودی‌ها به دو دسته پیکرببری و آمار زمان اجرا تقسیم می‌شوند. ورودی‌های پیکرببری در زمان ساخت شیء اثبات امنیتی خوانده می‌شوند و در طول اجرا ثابت می‌مانند، در حالی که ورودی‌های آمار در هر فراخوانی \lr{\texttt{calculate\_key\_length}} به روز می‌شوند.

\begin{itemize}
	\item \lr{\texttt{QKDParams}}: این شیء از ماژول \lr{\texttt{qkd.params}} وارد می‌شود و شامل تمام پارامترهای پیکرببری سیستم است. زیرشئه‌های کلیدی که این ماژول استفاده می‌کند عبارتند از: \lr{\texttt{p.source}} (شامل \lr{\texttt{pulse\_configs}}، \lr{\texttt{statistics\_type}}، \lr{\texttt{intensity\_jitter}} و \lr{\texttt{source\_fidelity}})، \lr{\texttt{p.detector}} (شامل \lr{\texttt{dark\_rate}}، \lr{\texttt{det\_eff\_d0}} و \lr{\texttt{qber\_intrinsic}})، \lr{\texttt{p.channel}} (شامل \lr{\texttt{transmittance}})، \lr{\texttt{p.eps\_sec}}، \lr{\texttt{p.eps\_pe}}، \lr{\texttt{p.eps\_smooth}}، \lr{\texttt{p.eps\_cor}}، \lr{\texttt{p.f\_error\_correction}}، \lr{\texttt{p.f\_ec\_dynamic\_config}}، \lr{\texttt{p.photon\_number\_cap}}، \lr{\texttt{p.lp\_solver\_method}} و \lr{\texttt{p.ci\_method}}.
	\item \lr{\texttt{stats\_map}}: این یک دیکشنری از \lr{\texttt{str}} به \lr{\texttt{TallyCounts}} است که از لایه شبیه‌ساز آشکارساز به اثبات امنیتی داده می‌شود. کلیدهای این دیکشنری نام شدت‌ها (\lr{``signal''}، \lr{``decoy''}، \lr{``vacuum''} و در حالت چند \lr{decoy}، نام‌های اختصاصی) هستند و مقادیر، شیء \lr{\texttt{TallyCounts}} با فیلدهای \lr{\texttt{sent}} (تعداد پالس‌های ارسالی)، \lr{\texttt{sifted}} (تعداد آشکارسازی‌های پایه درست) و \lr{\texttt{errors\_sifted}} (تعداد خطاها در پایه درست) هستند.
	\item \lr{\texttt{EpsilonAllocation}}: این شیء از ماژول \lr{\texttt{qkd.datatypes}} وارد می‌شود و شامل فیلدهای $\varepsilon_{\text{sec}}$، $\varepsilon_{\text{cor}}$، $\varepsilon_{\text{pe}}$، $\varepsilon_{\text{smooth}}$، $\varepsilon_{\text{pa}}$ و $\varepsilon_{\text{phase\_est}}$ است. این شیء از متد \lr{\texttt{allocate\_epsilons}} تولید می‌شود و در محاسبه کران‌های آماری استفاده می‌شود.
\end{itemize}

\subsection{ورودی‌های ماژول \lr{wang2005.py}}

ماژول \lr{wang2005.py} ورودی‌های مشابهی با \lr{ma2005.py} دارد اما به دو فیلد اضافی در پیکرببری منبع وابسته است:

\begin{itemize}
	\item \lr{\texttt{p.source.intensity\_jitter}}: این فیلد کران نوسان شدت منبع را مشخص می‌کند. در ماژول \lr{wang2005.py}، این فیلد به صورت \lr{$\delta$} در فرمول‌های کران شدت ($\mu^U = \mu(1 + \delta)$ و $\mu^L = \mu(1 - \delta)$) استفاده می‌شود. اگر این فیلد تنظیم نشده باشد، مقدار پیش‌فرض 0.0 استفاده می‌شود که معادل منبع ایده‌آل است.
	\item \lr{\texttt{p.source.source\_fidelity}}: این فیلد نشان می‌دهد که حالت واقعی منبع چقدر به یک \lr{coherent state} ایده‌آل نزدیک است. مقدار \lr{1.0} به معنای منبع ایده‌آل و مقدار کمتر از \lr{1.0} به معنای انحراف از حالت ایده‌آل است. اگر این فیلد کمتر از \lr{1.0} باشد، ماژول یک \lr{mapping} محافظه‌کارانه به صورت $\delta = \max(\delta_{\text{jitter}}, 1 - F)$ انجام می‌دهد.
\end{itemize}

این دو فیلد با استفاده از تابع \lr{\texttt{getattr}} با مقدار پیش‌فرض خوانده می‌شوند تا در صورت عدم وجود در پیکرببری، ماژول با خطا مواجه نشود. این طراحی انعطاف‌پذیر اجازه می‌دهد ماژول با نسخه‌های مختلف پیکرببری سازگار باشد.

\subsection{خروجی‌های ماژول}

خروجی اصلی هر دو ماژول یک شیء \lr{\texttt{KeyCalculationResult}} است که از کلاس پایه \lr{\texttt{FiniteKeyProof}} به ارث رسیده. این شیء شامل فیلدهای زیر است:

\begin{enumerate}
	\item \lr{\texttt{secure\_key\_length}}: طول کلید امن به صورت یک عدد صحیح غیرمنفی. این فیلد خروجی نهایی محاسبه است و توسط لایه بالایی برای تولید کلید واقعی استفاده می‌شود.
	\item \lr{\texttt{privacy\_amplification\_term}}: ترم تقویت حریم خصوصی به صورت $Q_1 (1 - H_2(e_1)) \cdot N_{\text{sent}}$. این فیلد برای تحلیل و گزارش‌گیری مفید است.
	\item \lr{\texttt{error\_correction\_leakage}}: نشت تصحیح خطا به صورت $Q_\mu f_{\text{EC}} H_2(E_\mu) \cdot N_{\text{sent}}$. این فیلد نشان می‌دهد چه مقدار اطلاعات در مرحله تصحیح خطا به مهاجم نشت کرده است.
	\item \lr{\texttt{phase\_error\_rate\_upper\_bound}}: کران بالای خطای فاز $e_1^U$. این فیلد برای تحلیل امنیتی و مقایسه با آستانه امنیت \lr{BB84} (حدود 11\%) استفاده می‌شود.
	\item \lr{\texttt{diagnostics}}: یک دیکشنری شامل تمام مقادیر میانی مانند $Q_\mu$، $E_\mu$، $Q_1$، $Y_1^L$، $e_1^U$، نرخ هر پالس، $f_{\text{EC}}$ استفاده‌شده، نرخ \lr{Weak GLLP}، $\Delta$ و اطلاعات \lr{solver}. این فیلد برای \lr{debugging} و بازبینی امنیتی ضروری است.
	\item \lr{\texttt{error\_codes}}: لیستی از کدهای خطا (مانند \lr{\texttt{INSUFFICIENT\_STATISTICS}}) که در صورت بروز مشکل، علت آن را نشان می‌دهند.
	\item \lr{\texttt{weak\_gllp\_rate}} و \lr{\texttt{tagged\_fraction\_bound}}: فیلدهای اضافی که در صورت استفاده از حالت \lr{Weak GLLP}، اطلاعات بیشتری ارائه می‌دهند.
\end{enumerate}

علاوه بر \lr{\texttt{KeyCalculationResult}}، متد \lr{\texttt{estimate\_yields\_and\_errors}} یک شیء \lr{\texttt{DecoyEstimates}} (یا \lr{\texttt{DecoyEstimatesWang2005}} در ماژول \lr{wang2005.py}) را برمی‌گرداند که شامل $Y_1^L$، $e_1^U$، \lr{\texttt{is\_feasible}}، \lr{\texttt{failure\_prob\_used}} و \lr{\texttt{diagnostics}} است. این شیء برای کاربرانی که می‌خواهند تخمین‌های میانی را بدون محاسبه کامل طول کلید بررسی کنند، مفید است.

\subsection{وابستگی‌های داخلی و خارجی}

هر دو ماژول وابستگی‌های داخلی و خارجی متعددی دارند. وابستگی‌های داخلی شامل موارد زیر است:

\begin{itemize}
	\item \lr{\texttt{qkd.proofs.base}}: کلاس پایه \lr{\texttt{FiniteKeyProof}}، \lr{\texttt{DecoyEstimates}}، \lr{\texttt{KeyCalculationResult}}، \lr{\texttt{SolverDiagnostics}} و \lr{\texttt{ErrorCode}} از این ماژول وارد می‌شوند. این کلاس پایه رابط عمومی و خدمات مشترک مانند آنتروپی باینری، \lr{clamp} عددی و گزارش‌گیری \lr{audit} را فراهم می‌سازد.
	\item \lr{\texttt{qkd.proofs.utils\_lp}}: تابع \lr{\texttt{solve\_lp}} برای حل مسأله برنامه‌ریزی خطی در کلاس \lr{\texttt{Ma2005GeneralMultiDecoyProof}} وارد می‌شود. این تابع از \lr{scipy.optimize.linprog} یا \lr{cvxopt} به عنوان \lr{solver} زیرین استفاده می‌کند.
	\item \lr{\texttt{qkd.datatypes}}: \lr{\texttt{TallyCounts}}، \lr{\texttt{EpsilonAllocation}}، \lr{\texttt{SourceStatisticsType}} و \lr{\texttt{ConfidenceBoundMethod}} از این ماژول وارد می‌شوند. \lr{\texttt{SourceStatisticsType}} یک \lr{enum} است که مقادیر \lr{POISSON} و \lr{THERMAL} را دارد و \lr{\texttt{ConfidenceBoundMethod}} یک \lr{enum} است که روش‌های فاصله اطمینان را مشخص می‌کند.
	\item \lr{\texttt{qkd.exceptions}}: استثناهای سفارشی \lr{\texttt{ConfigurationError}}، \lr{\texttt{LPFailureError}} و \lr{\texttt{ParameterValidationError}} برای گزارش خطاهای ساختاریافته.
	\item \lr{\texttt{qkd.utils.math}}: تابع \lr{\texttt{binary\_entropy}} برای محاسبه آنتروپی باینری $H_2(p) = -p \log_2 p - (1 - p) \log_2 (1 - p)$.
	\item \lr{\texttt{qkd.params.QKDParams}}: کلاس پارامترهای QKD که شامل تمام پارامترهای پیکرببری سیستم است.
\end{itemize}

وابستگی‌های خارجی شامل \lr{\texttt{numpy}} برای محاسبات برداری در مسیر \lr{LP}، \lr{\texttt{math}} برای توابع ریاضی پایه (مانند \lr{\texttt{exp}}، \lr{\texttt{log}}، \lr{\texttt{factorial}}) و \lr{\texttt{logging}} برای ثبت رویدادها و هشدارها هستند.

\subsection{جریان داده در سراسر فریم‌ورک}

جریان داده کلی به این شکل است:
\begin{enumerate}
	\item لایه شبیه‌ساز پالس‌های نوری را با شدت‌های $\mu$ (سیگنال)، $\nu$ (\lr{decoy} ضعیف) و 0 (خلأ) تولید و از کانال کوانتومی عبور می‌دهد.
	\item آشکارسازها در سمت باب، آشکارسازی‌ها و خطاها را برای هر شدت ثبت کرده و در \lr{\texttt{TallyCounts}} ذخیره می‌کنند.
	\item این شمارش‌ها به صورت یک دیکشنری \lr{\texttt{stats\_map}} با کلیدهای نام شدت به اثبات امنیتی داده می‌شوند.
	\item متد \lr{\texttt{allocate\_epsilons}} در زمان ساخت شیء فراخوانی می‌شود و بودجه $\varepsilon$ را میان اجزای مختلف تقسیم می‌کند.
	\item متد \lr{\texttt{estimate\_yields\_and\_errors}} با دریافت \lr{\texttt{stats\_map}}، کران‌های $Y_1^L$ و $e_1^U$ را با استفاده از فرمول‌های تحلیلی یا \lr{LP} محاسبه می‌کند.
	\item متد \lr{\texttt{calculate\_key\_length}} با دریافت تخمین‌ها و \lr{\texttt{stats\_map}}، نرخ کلید را با فرمول \lr{GLLP} محاسبه می‌کند و یک شیء \lr{\texttt{KeyCalculationResult}} برمی‌گرداند.
	\item لایه بالایی می‌تواند با متدهای \lr{\texttt{explain\_key\_decision}} و \lr{\texttt{audit\_dump}} (از کلاس پایه) علت نتیجه را به صورت متنی دریافت کند و یک گزارش کامل برای بازبینی امنیتی تولید کند.
\end{enumerate}

این جداسازی مراحل به گونه‌ای طراحی شده که هر مرحله قابل آزمایش و بازبینی مستقل باشد. به عنوان مثال، می‌توان \lr{\texttt{DecoyEstimates}} را به صورت دستی ساخت و رفتار مرحله نهایی محاسبه کلید را در شرایط خاص بررسی کرد. این ماژول‌بندی از اصل \lr{Single Responsibility} پیروی می‌کند و قابلیت نگهداری و توسعه کد را به طور قابل توجهی افزایش می‌دهد. علاوه بر این، شفافیت کامل در \lr{diagnostics} به کاربر اجازه می‌دهد هر مرحله از محاسبه را ردیابی کرده و در صورت نیاز، فرضیات تئوریک را برای کاربرد خاص خود بازبینی کند.

\subsection{انتخاب میان \lr{Ma 2005} و \lr{Wang 2008}}

یکی از تصمیمات مهم در استفاده از این فریم‌ورک، انتخاب میان ماژول \lr{ma2005.py} و \lr{wang2005.py} است. این انتخاب بر اساس شرایط واقعی منبع نوری انجام می‌شود:

\begin{itemize}
	\item اگر منبع ایده‌آل باشد (یعنی \lr{coherent state} با شدت دقیقاً مشخص و بدون نوسان)، ماژول \lr{ma2005.py} با یکی از پنج کلاس (مثلاً \lr{\texttt{Ma2005VacuumWeakProof}}) کافی است و نرخ کلید بالاتری تولید می‌کند زیرا فرض‌های سخت‌گیرانه‌تری دارد.
	\item اگر منبع دارای نوسان شدت یا انحراف از حالت \lr{coherent} باشد، ماژول \lr{wang2005.py} باید استفاده شود زیرا این نوسانات را در فرمول‌های کران لحاظ می‌کند. استفاده از \lr{ma2005.py} در این حالت می‌تواند منجر به کران‌های خوش‌بینانه (ناامن) شود.
	\item اگر تعداد \lr{decoy}‌ها بیش از دو باشد، تنها \lr{\texttt{Ma2005GeneralMultiDecoyProof}} از ماژول \lr{ma2005.py} قابل استفاده است (زیرا \lr{wang2005.py} فقط پروتکل \lr{Vacuum + Weak} را پشتیبانی می‌کند).
	\item اگر منبع از نوع \lr{Thermal} باشد، تنها \lr{\texttt{Ma2005GeneralMultiDecoyProof}} قابل استفاده است (سایر کلاس‌ها در صورت \lr{Thermal} بودن، استثنای \lr{\texttt{ConfigurationError}} \lr{raise} می‌کنند).
	\item اگر کاربر بخواهد حد نظری \lr{Infinite Decoy} را به عنوان مرجع بالای عملکرد محاسبه کند، \lr{\texttt{Ma2005AsymptoticProof}} مناسب است.
\end{itemize}

این انعطاف‌پذیری در انتخاب، به کاربر اجازه می‌دهد تحلیل امنیتی متناسب با شرایط خاص سیستم خود را انتخاب کند و در عین حال، از کران‌های قابل اثبات ریاضی بهره‌مند شود. طراحی ماژولار همچنین اجازه می‌دهد در آینده، چارچوب‌های جدید (مانند \lr{Lim 2014} یا اثبات‌های مبتنی بر \lr{MDI-QKD}) به سادگی با ارث‌بری از \lr{\texttt{FiniteKeyProof}} اضافه شوند بدون آن‌که کد موجود تغییر کند.

	
	\chapter{اثبات امنیت فشرده}
	\label{ch:tight}
	\sloppy
\emergencystretch=5em
\hbadness=10000
\vbadness=10000
\tolerance=2000
\providecolor{codebg}{RGB}{248,248,248}
\providecolor{codeframe}{RGB}{200,200,200}
\providecolor{keyword}{RGB}{0,0,180}
\providecolor{string}{RGB}{163,21,21}
\providecolor{comment}{RGB}{0,128,0}
\lstset{
	basicstyle=\ttfamily\scriptsize,
	breaklines=true,
	breakatwhitespace=false,
	columns=fullflexible,
	keepspaces=true,
	frame=single,
	framesep=2pt,
	showstringspaces=false,
	backgroundcolor=\color{codebg},
	rulecolor=\color{codeframe},
	keywordstyle=\color{keyword}\bfseries,
	stringstyle=\color{string},
	commentstyle=\color{comment}\itshape,
	numberstyle=\tiny\color{gray},
	numbers=left,
	numbersep=6pt,
	xleftmargin=14pt,
	framexleftmargin=14pt,
	literate=
	{ą}{{\v{a}}}1 {ę}{{\k{e}}}1 {ż}{{\.z}}1 {ś}{{\'s}}1
	{ł}{{\l}}1 {ć}{{\'c}}1 {ń}{{\'n}}1 {ó}{{\'o}}1 {ź}{{\'z}}1
	{→}{{$\rightarrow$}}1 {←}{{$\leftarrow$}}1 {≠}{{$\neq$}}1
}

\section{نمای کلی و هدف ماژول}

فایل \lr{\texttt{qkd/proofs/tight.py}} پیاده‌سازی مرجع و تولید-آمادۀ (\lr{production-grade}) اثبات امنیتی کلید متناهی \lr{BB84-Tight} است که از چارچوب \lr{Lim et al. 2014} مشتق شده و به طور گسترده برای رفع استانداردهای تولید در زمینه‌های صحت‌سنجی، یکپارچگی \lr{API}، پایداری عددی، تشخیص‌گری (\lr{diagnostics}) و ردپایی (\lr{provenance}) بازطراحی شده است. این ماژول بر خلاف پیاده‌سازی‌های اولیه‌ی آکادمیک که صرفاً به دنبال بازتولید فرمول‌های مقاله بودند، با هدف قرار گرفتن در محیط‌های عملیاتی واقعی توسعه یافته است؛ محیط‌هایی که در آن‌ها خطاهای عددی، نوسانات ورودی و رفتار غیرعادی سیستم نباید منجر به فروپاشی خط لوله شبیه‌سازی شوند. نسخه‌ی فعلی با شناسه \lr{\texttt{3.7.6-fixed}} علامت‌گذاری شده است که نشان‌دهنده‌ی اصلاح یک باگ منطقی حیاتی در تخصیص بودجۀ امنیتی است؛ این باگ پیش از آن، در ریاضیات ممیز شناور ۶۴-بیت، می‌توانست منجر به نقض نامرئی نامساوی امنیتی $\varepsilon_{\text{sec}}$ شود.

هدف اصلی این ماژول، تبدیل شمارش‌های آماری شبکۀ آشکارسازی (در قالب شیء \lr{\texttt{TallyCounts}}) و تخمین‌های \lr{Decoy-State} (در قالب شیء \lr{\texttt{DecoyEstimatesInput}} یا هر شیء دارای فیلدهای \lr{\texttt{yield\_1\_lower\_bound}} و \lr{\texttt{error\_rate\_1\_upper\_bound}}) به یک عدد صحیح مثبت به عنوان طول کلید امن نهایی است. این تبدیل با رعایت سخت‌گیرانۀ تمام محافظه‌کاری‌های ریاضی، آماری و فیزیکی انجام می‌شود تا هیچ طول کلیدی که از نظر تئوری قابل اثبات نیست، تولید نشود. فلسفۀ طراحی این ماژول بر پایۀ اصل \lr{``fail-safe, not fail-fast''} استوار است؛ به این معنا که در صورت بروز هرگونه ناسازگاری پیکربندی، آمار ناکافی، یا خطای عددی، خروجی به جای \lr{raise} کردن استثنا و متوقف کردن خط لوله، یک شیء \lr{\texttt{KeyCalculationResult}} با طول کلید صفر و کد خطای مناسب تولید می‌کند. این رفتار برای کاربردهای تولیدی که در آن‌ها قطع شبیه‌سازی هزینه‌بر است، حیاتی است و اجازه می‌دهد خط لوله ادامه یابد در حالی که علت شکست در گزارش تشخیصی ثبت می‌گردد.

این ماژول از کلاس پایه \lr{\texttt{Lim2014Proof}} (که در فایل \lr{\texttt{qkd/proofs/lim2014.py}} تعریف شده) ارث می‌برد و متدهای \lr{\texttt{allocate\_epsilons}} و \lr{\texttt{calculate\_key\_length}} را بازنویسی می‌کند. ارث‌بری از \lr{\texttt{Lim2014Proof}} به این ماژول اجازه می‌دهد تا تمام زیرساخت‌های مشترک مانند محاسبۀ آنتروپی باینری، تخصیص \lr{$\varepsilon$} پایه، و ابزارهای \lr{clamp} عددی را بدون تکرار کد به کار گیرد و در عین حال، منطق محاسبۀ کلید را با رویکردی متفاوت و سخت‌گیرانه‌تر بازنویسی کند. تفاوت کلیدی میان \lr{\texttt{BB84TightProof}} و کلاس پایۀ \lr{\texttt{Lim2014Proof}} در این است که نسخۀ \lr{Tight} از یک رویکرد \lr{``entropy-budget''} استفاده می‌کند که در آن، آنتروپی موجود از پایۀ تولید‌کنندۀ کلید ($X$) محاسبه شده و سپس نشت تصحیح خطا و جریمه‌های تقویت حریم خصوصی از آن کم می‌شوند، در حالی که نسخۀ پایه از فرمول مستقیم نرخ کلید استفاده می‌کند. این رویکرد آنتروپی-محور، تطابق دقیق‌تری با ادبیات مدرن امنیتی ترکیبی (\lr{composable security}) دارد.

یکی از ویژگی‌های متمایزکنندۀ این ماژول، پشتیبانی از هر دو حالت \lr{Decoy-State} و \lr{non-Decoy} است. در حالت \lr{Decoy-State}، ماژول تخمین‌های $Y_1$ و $e_1$ را از موتور تخمین \lr{Decoy} دریافت کرده و از آن‌ها برای محاسبۀ شمارش تک‌فوتونی در هر دو پایۀ $X$ و $Z$ استفاده می‌کند. در حالت \lr{non-Decoy}، ماژول به صورت محافظه‌کارانه فرض می‌کند که تمام بیت‌های \lr{sifted} از تک‌فوتون‌ها نشأت گرفته‌اند و از \lr{QBER} پایۀ $Z$ به عنوان تخمین خطای فاز استفاده می‌کند. این انعطاف‌پذیری اجازه می‌دهد ماژول هم برای شبیه‌سازی‌های پیشرفته با \lr{Decoy-State} و هم برای شبیه‌سازی‌های ساده \lr{BB84} پایه به کار رود. علاوه بر این، ماژول از یک طرح غنی گزارش‌گیری \lr{diagnostics} استفاده می‌کند که در آن تمام مقادیر میانی محاسبه (مانند $n_X$، $m_X$، $n_Z$، $m_Z$، $s_{X,1}^L$، $s_{Z,1}^L$، \lr{QBER}، نشت تصحیح خطا، خطای فاز، جریمه‌های \lr{PA} و آنتروپی موجود) در یک دیکشنری ساختاریافته ذخیره می‌شوند و برای تحلیل و بازبینی امنیتی توسط شخص ثالث قابل دسترس هستند.

\section{مبانی تئوری و فیزیکی}

\subsection{امنیت ترکیبی و بودجه‌ی $\varepsilon$}

در چارچوب مدرن امنیت ترکیبی (\lr{composable security}) که توسط \lr{Canetti} (2001) و \lr{Ben-Or, Mayers} (2004) توسعه یافته، امنیت یک پروتکل \lr{QKD} به صورت یک پارامتر کوچک $\varepsilon_{\text{sec}}$ تعریف می‌شود که نشان‌دهندۀ حداکثر احتمال شکست امنیتی است. یک پروتکل $\varepsilon_{\text{sec}}$-امن است اگر خروجی آن از یک کلید ایده‌آل (توزیع یکنواخت و کاملاً مخفی) با فاصله‌ی $\varepsilon_{\text{sec}}$ در معیار \lr{trace distance} فاصله داشته باشد. این تعریف اجازه می‌دهد امنیت \lr{QKD} به صورت قابل ترکیب با سایر پروتکل‌های رمزنگاری تحلیل شود. کل بودجه‌ی $\varepsilon_{\text{sec}}$ باید میان چندین رویداد شکست مستقل تقسیم شود:
\begin{equation}
	\varepsilon_{\text{sec}} \;\ge\; \varepsilon_{\text{PE}} + 2\,\varepsilon_{\text{smooth}} + \varepsilon_{\text{PA}} + \varepsilon_{\text{cor}} + \varepsilon_{\text{phase\_est}}.
\end{equation}
در این نامساوی، $\varepsilon_{\text{PE}}$ بودجه‌ی تخمین پارامتر، $\varepsilon_{\text{smooth}}$ بودجه‌ی هموارسازی (\lr{smoothing}) که دو بار استفاده می‌شود (یک بار برای آنتروپی حداقل هموار و یک بار برای آنتروپی حداقل شرطی هموار)، $\varepsilon_{\text{PA}}$ بودجه‌ی تقویت حریم خصوصی (\lr{Privacy Amplification})، $\varepsilon_{\text{cor}}$ بودجه‌ی تصحیح خطا (\lr{Error Correction}) و $\varepsilon_{\text{phase\_est}}$ بودجه‌ی تخمین خطای فاز است. در کد ماژول \lr{\texttt{tight.py}}، این تقسیم در متد \lr{\texttt{allocate\_epsilons}} پیاده‌سازی شده و با فراخوانی \lr{\texttt{allocation.validate()}} اعتبارسنجی می‌شود. یک نکتۀ ظریف و حیاتی این است که در ریاضیات ممیز شناور ۶۴-بیت، مجموع \lr{$\varepsilon_{\text{cor}} + (\varepsilon_{\text{sec}} - \varepsilon_{\text{cor}})$} ممکن است به دلیل خطای گردکردن، کمی بزرگتر از $\varepsilon_{\text{sec}}$ شود و اعتبارسنجی شکست بخورد. برای رفع این مشکل، ماژول یک حاشیۀ امنیتی مطلق به اندازه‌ی $10^{-20}$ از بودجه‌ی باقی‌مانده کم می‌کند تا تضمین شود نامساوی به صورت عددی برقرار می‌ماند.

\subsection{خطای فاز در مقابل خطای بیت}

در پروتکل \lr{BB84}، دو نوع خطا تعریف می‌شود: خطای بیت (\lr{bit error}) و خطای فاز (\lr{phase error}). خطای بیت، خطایی است که در پایۀ اندازه‌گیری قابل مشاهده است و با مقایسۀ بیت‌های آلیس و باب در همان پایه اندازه‌گیری می‌شود. در مقابل، خطای فاز خطایی است که اگر آلیس و باب در پایۀ مزدوج (\lr{conjugate basis}) اندازه‌گیری می‌کردند، قابل مشاهده می‌بود. این دو خطا در یک کانال بدون نویز برابرند، اما در حضور حمله‌های کوانتومی مانند \lr{PNS} یا حمله‌های \lr{coherent} می‌توانند متفاوت باشند. رابطۀ ریاضی میان خطای فاز و خطای بیت با استفاده از نامساوی \lr{Hoeffding} یا \lr{Azuma} برقرار می‌شود. در این ماژول، خطای فاز از پایۀ $Z$ تخمین زده می‌شود و خطای بیت از پایۀ $X$ (که پایۀ تولیدکنندۀ کلید است).

نامساوی \lr{Hoeffding} برای تخمین فاصلۀ میان خطای فاز و خطای بیت به صورت زیر اعمال می‌شود:
\begin{equation}
	e_{\text{phase}} \;\le\; e_{\text{bit}} + \sqrt{\frac{-\ln(\varepsilon_{\text{phase\_est}})}{2\,s_{Z,1}^L}},
\end{equation}
که در آن $e_{\text{bit}} = e_1^U$ کران بالای خطای بیت تک‌فوتونی، $\varepsilon_{\text{phase\_est}}$ بودجه‌ی تخمین خطای فاز و $s_{Z,1}^L$ کران پایین شمارش تک‌فوتونی در پایۀ $Z$ است. این نامساوی نشان می‌دهد که هرچه شمارش تک‌فوتونی در پایۀ $Z$ بیشتر باشد، عدم قطعیت در تخمین خطای فاز کمتر می‌شود. در کد، این فرمول در متد \lr{\texttt{\_compute\_phase\_error\_ub}} پیاده‌سازی شده و در صورت بزرگتر بودن کران نهایی از 0.5 (حداکثر آنتروپی باینری)، به 0.5 بریده می‌شود. علاوه بر این، اگر پرچم \lr{\texttt{assume\_phase\_equals\_bit\_error}} در پیکربندی فعال باشد، ماژول به جای محاسبۀ نامساوی \lr{Hoeffding}، فرض می‌کند $e_{\text{phase}} = e_{\text{bit}}$ که یک تقریب ساده‌تر اما محافظه‌کارانه‌تر در شرایطی است که شمارش پایۀ $Z$ ناکافی است.

\subsection{آنتروپی باینری و برش عددی}

آنتروپی باینری $H_2(p)$ یک تابع کلیدی در تمام محاسبات امنیتی \lr{QKD} است و به صورت زیر تعریف می‌شود:
\begin{equation}
	H_2(p) = -p\,\log_2(p) - (1 - p)\,\log_2(1 - p).
\end{equation}
این تابع در $p = 0.5$ به حداکثر مقدار ۱ می‌رسد و در $p = 0$ یا $p = 1$ به صفر میل می‌کند. در محاسبۀ طول کلید، آنتروپی باینری در دو جایگاه ظاهر می‌شود: نخست، در محاسبۀ نشت تصحیح خطا (\lr{EC leakage}) به صورت $f_{\text{EC}} \cdot H_2(E_{\text{bit}}) \cdot n_{\text{key}}$ که نشان‌دهندۀ تعداد بیت‌هایی است که در فرآیند تصحیح خطا به باب نشت می‌کنند و در نتیجه توسط مهاجم قابل دسترس هستند؛ دوم، در محاسبۀ آنتروپی تک‌فوتونی به صورت $s_{X,1}^L \cdot (1 - H_2(e_{\text{phase}}))$ که نشان‌دهندۀ آنتروپی قابل استخراج از تک‌فوتون‌های پایۀ تولیدکنندۀ کلید است.

یک مشکل عددی ظریف این است که آنتروپی باینری در $p \to 0$ یا $p \to 1$ به سمت صفر میل می‌کند، اما عبارت $\log_2(p)$ به سمت $-\infty$ میل می‌کند. در ریاضیات ممیز شناور، این می‌تواند منجر به ضرب یک عدد بسیار کوچک در یک عدد بسیار بزرگ شود و نتیجه‌ی ناپایدار یا حتی \lr{NaN} تولید کند. برای رفع این مشکل، ماژول از یک تابع \lr{\texttt{\_binary\_entropy\_safe}} استفاده می‌کند که مقدار $p$ را با استفاده از ثابت \lr{\texttt{ENTROPY\_PROB\_CLAMP}} (که معمولاً $10^{-12}$ است) به بازۀ $[\varepsilon, 1 - \varepsilon]$ بریده می‌کند:
\begin{equation}
	p_{\text{clamped}} = \max(\varepsilon_{\text{clamp}}, \min(p, 1 - \varepsilon_{\text{clamp}})).
\end{equation}
این برش عددی، در حالی که به طور ظاهری کوچک است، تضمین می‌کند که تمام محاسبات آنتروپی پایدار و قابل بازتولید باشند. این برش همچنین از تقسیم بر صفر در عبارت $\log_2(p)$ جلوگیری می‌کند. در ادبیات امنیتی، این برش به عنوان یک ``بهای عددی'' شناخته می‌شود که در صورتی که $\varepsilon_{\text{clamp}} \ll \varepsilon_{\text{sec}}$ باشد، از نظر تئوری ناچیز است.

\subsection{آنتروپی موجود و معادلۀ کلید نهایی}

فرمول اصلی محاسبۀ طول کلید امن در چارچوب \lr{BB84-Tight} به صورت زیر است:
\begin{equation}
	\ell \;=\; \underbrace{s_{X,1}^L \cdot \left(1 - H_2(e_{\text{phase}})\right)}_{\text{available entropy}} - \underbrace{f_{\text{EC}} \cdot H_2(e_{\text{bit}}) \cdot n_{\text{key}}}_{\text{EC leakage}} - \underbrace{\left(\log_2\frac{1}{2\varepsilon_{\text{smooth}}} + \log_2\frac{1}{\varepsilon_{\text{PA}}} + \log_2\frac{2}{\varepsilon_{\text{cor}}}\right)}_{\text{PA \& correction penalties}}.
\end{equation}
در این معادله، $\ell$ طول کلید امن نهایی، $s_{X,1}^L$ کران پایین شمارش تک‌فوتونی در پایۀ $X$ (پایۀ تولیدکنندۀ کلید)، $e_{\text{phase}}$ کران بالای خطای فاز، $e_{\text{bit}}$ کران بالای خطای بیت، $n_{\text{key}}$ تعداد بیت‌های \lr{sifted} در پایۀ $X$، $f_{\text{EC}}$ بازدهی تصحیح خطا و $\varepsilon_{\text{smooth}}$، $\varepsilon_{\text{PA}}$ و $\varepsilon_{\text{cor}}$ بودجه‌های امنیتی مربوطه هستند. ترم آنتروپی موجود نشان‌دهندۀ حداکثر اطلاعات قابل استخراج از تک‌فوتون‌های امن است. ترم نشت تصحیح خطا نشان‌دهندۀ اطلاعاتی است که در فرآیند تصحیح خطا به باب نشت می‌کند و توسط مهاجم قابل رهگیری است. در نهایت، ترم جریمه‌های \lr{PA} و تصحیح، بیت‌های ثابتی هستند که در فرآیند تقویت حریم خصوصی و تأیید تصحیح خطا از دست می‌روند.

این فرمول در متد \lr{\texttt{calculate\_key\_length}} پیاده‌سازی شده است. یک نکتۀ کلیدی این است که اگر $\ell < 0$ شود (یعنی نشت و جریمه‌ها از آنتروپی موجود بیشتر شوند)، طول کلید به صفر برده می‌شود. این رفتار با استفاده از متد \lr{\texttt{\_clamp\_key\_length}} از کلاس پایه تضمین می‌شود. علاوه بر این، هر مقدار میانی غیرمتناهی (\lr{NaN} یا \lr{Inf}) باعث ثبت کد خطای \lr{\texttt{NON\_FINITE\_VALUE}} می‌شود. این بررسی‌ها در متد \lr{\texttt{\_finite\_or\_error}} پیاده‌سازی شده‌اند.

\subsection{تخمین شمارش تک‌فوتونی در دو پایه}

یکی از چالش‌های کلیدی در محاسبۀ طول کلید، تخمین شمارش تک‌فوتونی در دو پایۀ $X$ و $Z$ است. در حالت \lr{Decoy-State}، این شمارش‌ها به صورت زیر تخمین زده می‌شوند:
\begin{align}
	s_{X,1}^L &= N_{\text{sent},X} \cdot P(1 \mid \mu) \cdot Y_1^L \cdot p_X^A \cdot p_X^B, \\
	s_{Z,1}^L &= N_{\text{sent},Z} \cdot P(1 \mid \mu) \cdot Y_1^L \cdot p_Z^A \cdot p_Z^B,
\end{align}
که در آن $N_{\text{sent},X}$ و $N_{\text{sent},Z}$ تعداد پالس‌های ارسالی در پایۀ $X$ و $Z$، $P(1 \mid \mu) = \mu\,e^{-\mu}$ احتمال ارسال تک‌فوتون با شدت $\mu$، $Y_1^L$ کران پایین بازده تک‌فوتونی (که فرض می‌شود مستقل از پایه است)، $p_X^A$ و $p_X^B$ احتمال انتخاب پایۀ $X$ توسط آلیس و باب، و به طور متقابل $p_Z^A$ و $p_Z^B$ احتمال انتخاب پایۀ $Z$ است. این فرمول‌ها در متد \lr{\texttt{\_compute\_base\_counts\_decoy}} پیاده‌سازی شده‌اند. یک نکتۀ محافظه‌کارانه این است که اگر $s_{X,1}^L > n_X$ شود (که از نظر فیزیکی غیرممکن است زیرا تعداد تک‌فوتون‌ها نمی‌تواند از کل آشکارسازی‌ها بیشتر باشد)، مقدار به $n_X$ بریده می‌شود. این برش با استفاده از ثابت \lr{\texttt{NUMERIC\_ABS\_TOL}} انجام می‌شود تا خطاهای کوچک گردکردن نادیده گرفته شوند.

در حالت \lr{non-Decoy}، ماژول به صورت محافظه‌کارانه فرض می‌کند که تمام بیت‌های \lr{sifted} از تک‌فوتون‌ها نشأت گرفته‌اند، یعنی $s_{X,1}^L = n_X$ و $s_{Z,1}^L = n_Z$. این فرض، در حالی که در حضور پالس‌های چندفوتونی ناامن است، برای شبیه‌سازی‌های \lr{BB84} ایده‌آل (با $\mu \to 0$) قابل قبول است. در این حالت، خطای فاز به طور مستقیم از \lr{QBER} پایۀ $Z$ تخمین زده می‌شود: $e_{\text{bit}} = m_Z / n_Z$.

\subsection{جریمه‌های \lr{PA} و تصحیح}

جریمه‌های \lr{PA} و تصحیح، بیت‌های ثابتی هستند که در فرآیند تقویت حریم خصوصی و تأیید تصحیح خطا از دست می‌روند. این جریمه‌ها به صورت زیر محاسبه می‌شوند:
\begin{equation}
	\text{PA penalty} = \log_2\!\left(\frac{1}{2\,\varepsilon_{\text{smooth}}}\right) + \log_2\!\left(\frac{1}{\varepsilon_{\text{PA}}}\right), \qquad \text{Correction penalty} = \log_2\!\left(\frac{2}{\varepsilon_{\text{cor}}}\right).
\end{equation}
این جریمه‌ها در متد \lr{\texttt{\_compute\_pa\_terms}} پیاده‌سازی شده‌اند. در این متد، ابتدا هر مقدار $\varepsilon$ با استفاده از متد \lr{\texttt{\_log\_inv\_clamp}} برش داده می‌شود تا از $\log_2(0)$ یا مقادیر بسیار ناپایدار جلوگیری شود. این متد، اگر مقدار ورودی کمتر از \lr{\texttt{\_TIGHT\_PROOF\_EPS\_MIN}} (که برابر $10^{-300}$ است) باشد، آن را به این حداقل برش می‌دهد. این برش اطمینان می‌دهد که حتی در صورت تنظیم $\varepsilon = 0$ به اشتباه، محاسبات منجر به \lr{$-\infty$} نشوند. در عمل، این جریمه‌ها معمولاً در حدود چند ده بیت هستند و در شبیه‌سازی‌های با کلیدهای طولانی (میلیون‌ها بیت) ناچیز هستند، اما در کلیدهای کوتاه می‌توانند تأثیر قابل توجهی بر طول کلید نهایی داشته باشند.

\section{تحلیل معماری و ساختار داده‌ها}

\subsection{سلسلۀ کلاس‌ها و وارثت}

ماژول \lr{\texttt{tight.py}} از یک طراحی سلسله‌مراتبی استفاده می‌کند که در آن کلاس اصلی \lr{\texttt{BB84TightProof}} از کلاس \lr{\texttt{Lim2014Proof}} (که خود از \lr{\texttt{FiniteKeyProof}} انتزاعی ارث می‌برد) ارث می‌برد. این طراحی چند لایه به ماژول اجازه می‌دهد تا از زیرساخت‌های مشترک مانند محاسبۀ آنتروپی باینری، تخصیص پایۀ \lr{$\varepsilon$} و ابزارهای \lr{clamp} عددی بهره ببرد و در عین حال، منطق محاسبۀ کلید را با رویکردی متفاوت بازنویسی کند. در ادامه، ساختار اصلی کلاس را به صورت خلاصه نمایش می‌دهیم:

\begin{LTR}
	\begin{lstlisting}[language=Python]
		class BB84TightProof(Lim2014ProofBase):
		__implementation_version__ = "3.7.6-fixed"
		
		def __repr__(self) -> str:
		return f"BB84TightProof(impl_version='{self.__implementation_version__}')"
	\end{lstlisting}
\end{LTR}

استفاده از ویژگی کلاس \lr{\texttt{\_\_implementation\_version\_\_}} به عنوان یک شناسۀ نسخه، اجازه می‌دهد تا در زمان اجرا، نسخۀ دقیق پیاده‌سازی مورد استفاده را شناسایی کنیم. این ویژگی برای ردیابی پروونانس (\lr{provenance}) در شبیه‌سازی‌های تولیدی حیاتی است، زیرا به کاربر اجازه می‌دهد بداند کدام نسخۀ اثبات امنیتی نتایج یک شبیه‌سازی خاص را تولید کرده است. متد \lr{\texttt{\_\_repr\_\_}} نیز یک نمایش متنی خوانا از کلاس ارائه می‌دهد که در \lr{debugging} و گزارش‌گیری مفید است. عبارت \lr{``3.7.6-fixed''} نشان‌دهندۀ نسخۀ 3.7.6 با اصلاح یک باگ منطقی حیاتی است.

\subsection{شمارۀ خطاهای ساختاریافته}

ماژول یک \lr{enum} به نام \lr{\texttt{ErrorCode}} تعریف می‌کند که تمام کدهای خطای ممکن را شامل می‌شود:

\begin{LTR}
	\begin{lstlisting}[language=Python]
		class ErrorCode(Enum):
		INVALID_PARAMS = "INVALID_PARAMS"
		INVALID_DECOY_ESTIMATES_TYPE = "INVALID_DECOY_ESTIMATES_TYPE"
		MISSING_DECOY_FIELDS = "MISSING_DECOY_FIELDS"
		INVALID_STATS_MAP = "INVALID_STATS_MAP"
		MISSING_SIGNAL_CONFIG = "MISSING_SIGNAL_CONFIG"
		INVALID_PULSE_CONFIG = "INVALID_PULSE_CONFIG"
		NO_SIGNAL_PULSES_SENT = "NO_SIGNAL_PULSES_SENT"
		INSUFFICIENT_STATISTICS = "INSUFFICIENT_STATISTICS"
		INCONSISTENT_MEASUREMENT = "INCONSISTENT_MEASUREMENT"
		NON_FINITE_VALUE = "NON_FINITE_VALUE"
		PHASE_ERROR_CLAMPED = "PHASE_ERROR_CLAMPED"
		INVALID_YIELD_BOUND = "INVALID_YIELD_BOUND"
		INVALID_ERROR_RATE_BOUND = "INVALID_ERROR_RATE_BOUND"
	\end{lstlisting}
\end{LTR}

استفاده از \lr{enum} به جای رشته‌های خام، چندین مزیت دارد: نخست، جلوگیری از خطاهای املایی در نام کدها؛ دوم، تکمیل خودکار در ویرایشگرهای کد؛ سوم، امکان جستجو و بازبینی سیستماتیک. هر کد خطا معنای مشخصی دارد. به عنوان مثال، \lr{\texttt{INVALID\_DECOY\_ESTIMATES\_TYPE}} نشان می‌دهد که شیء ورودی \lr{decoy\_estimates} نه از نوع \lr{\texttt{DecoyEstimatesInput}} است و نه دارای فیلدهای مورد نیاز است. \lr{\texttt{MISSING\_DECOY\_FIELDS}} نشان می‌دهد که فیلدهای \lr{\texttt{yield\_1\_lower\_bound}} یا \lr{\texttt{error\_rate\_1\_upper\_bound}} در شیء ورودی موجود نیستند. \lr{\texttt{NON\_FINITE\_VALUE}} نشان می‌دهد که یک مقدار میانی در محاسبات به \lr{NaN} یا \lr{Inf} تبدیل شده است. این کدها در شیء \lr{\texttt{KeyCalculationResult}} در فیلد \lr{\texttt{error\_codes}} ذخیره می‌شوند و در متد \lr{\texttt{\_finalize\_result}} به رشته تبدیل می‌شوند تا در خروجی \lr{JSON} قابل سریال‌سازی باشند.

\subsection{ساختار \lr{dataclass} برای ورودی و خروجی}

ماژول دو \lr{dataclass} اصلی تعریف می‌کند. نخست، \lr{\texttt{DecoyEstimatesInput}} که ورودی نرمال‌شدۀ تخمین‌های \lr{Decoy} را نگه می‌دارد:

\begin{LTR}
	\begin{lstlisting}[language=Python]
		@dataclass
		class DecoyEstimatesInput:
		"""
		Normalized decoy estimates required for the key calculation.
		This proof expects the bit error rate (e1_bit_ub) from the Z-basis
		and the yield (Y1_L_rate) which is assumed basis-independent.
		"""
		Y1_L_rate: float
		e1_bit_ub: float  # This is e1_z_ub (error rate in Z-basis)
	\end{lstlisting}
\end{LTR}

این \lr{dataclass} دو فیلد دارد: \lr{\texttt{Y1\_L\_rate}} که کران پایین بازده تک‌فوتونی است (به صورت یک نسبت، مستقل از پایه) و \lr{\texttt{e1\_bit\_ub}} که کران بالای خطای بیت تک‌فوتونی در پایۀ $Z$ است. در این طراحی، فرض می‌شود که $Y_1$ مستقل از پایه است (یعنی $Y_{1,X} = Y_{1,Z} = Y_1$) که یک فرض استاندارد در پروتکل‌های \lr{Decoy-State BB84} است. این \lr{dataclass} در متد \lr{\texttt{\_normalize\_decoy\_estimates}} از ورودی‌های مختلف (شیء \lr{\texttt{DecoyEstimates}} از \lr{base.py}، دیکشنری، یا خود \lr{\texttt{DecoyEstimatesInput}}) ساخته می‌شود تا یک رابط یکنواخت برای ادامه‌ی محاسبات فراهم شود.

دوم، \lr{\texttt{KeyCalculationResult}} که خروجی نهایی ماژول است:

\begin{LTR}
	\begin{lstlisting}[language=Python]
		@dataclass
		class KeyCalculationResult:
		secure_key_length: int = 0
		diagnostics: Dict[str, Any] = field(default_factory=dict)
		error_codes: List[str] = field(default_factory=list)
		metadata: Dict[str, Any] = field(default_factory=dict)
	\end{lstlisting}
\end{LTR}

این \lr{dataclass} چهار فیلد دارد: \lr{\texttt{secure\_key\_length}} که طول کلید امن نهایی به صورت یک عدد صحیح غیرمنفی است؛ \lr{\texttt{diagnostics}} که یک دیکشنری شامل تمام مقادیر میانی محاسبه (مانند $n_X$، $m_X$، $n_Z$، $m_Z$، \lr{QBER}، نشت تصحیح خطا، خطای فاز، جریمه‌های \lr{PA} و آنتروپی موجود) است؛ \lr{\texttt{error\_codes}} که لیستی از کدهای خطای رخ داده است؛ و \lr{\texttt{metadata}} که شامل اطلاعاتی مانند زمان‌شناسه (\lr{timestamp}) و مدت زمان محاسبه به میلی‌ثانیه است. استفاده از \lr{\texttt{field(default\_factory=dict)}} به جای مقدار پیش‌فرض \lr{\texttt{\{\}}} یک الگوی استاندارد \lr{Python} است که از اشتراک‌گذاری \lr{state} تصادفی میان نمونه‌ها جلوگیری می‌کند.

\subsection{استثناهای سفارشی}

ماژول دو استثنای سفارشی تعریف می‌کند:

\begin{LTR}
	\begin{lstlisting}[language=Python]
		class ParameterValidationError(ValueError): pass
		class QKDSimulationError(RuntimeError): pass
	\end{lstlisting}
\end{LTR}

\lr{\texttt{ParameterValidationError}} از \lr{\texttt{ValueError}} ارث می‌برد و برای خطاهای مربوط به پیکربندی پارامترها (مانند $\varepsilon_{\text{sec}} \le \varepsilon_{\text{cor}}$) استفاده می‌شود. \lr{\texttt{QKDSimulationError}} از \lr{\texttt{RuntimeError}} ارث می‌برد و برای خطاهای زمان اجرا (مانند ناسازگاری اندازه‌گیری) استفاده می‌شود. این جداسازی استثناها به کاربر اجازه می‌دهد به صورت انتخابی آن‌ها را بگیرد (\lr{catch}) و رفتار متفاوتی برای هر نوع خطا نشان دهد. در متد \lr{\texttt{calculate\_key\_length}}، این استثناها در یک بلوک \lr{\texttt{try/except}} گرفته می‌شوند و به جای \lr{raise} مجدد، در فیلد \lr{\texttt{error\_codes}} ثبت می‌شوند تا رفتار \lr{fail-safe} حفظ شود.

\subsection{ثابت‌های عددی و برش‌ها}

ماژول چندین ثابت عددی برای پایداری محاسبات تعریف می‌کند:

\begin{LTR}
	\begin{lstlisting}[language=Python]
		_TIGHT_PROOF_EPS_MIN = 1e-300
		DEFAULT_ENTROPY_PROB_CLAMP = 1e-12
		DEFAULT_NUMERIC_TOL = 1e-9
	\end{lstlisting}
\end{LTR}

\lr{\texttt{\_TIGHT\_PROOF\_EPS\_MIN}} کمترین مقدار مجاز برای $\varepsilon$ در محاسبات لگاریتمی است و از \lr{$\log_2(0) = -\infty$} جلوگیری می‌کند. \lr{\texttt{DEFAULT\_ENTROPY\_PROB\_CLAMP}} حد برش برای احتمالات در محاسبۀ آنتروپی باینری است و از \lr{$\log_2(0)$} در $p = 0$ یا $p = 1$ جلوگیری می‌کند. \lr{\texttt{DEFAULT\_NUMERIC\_TOL}} تلورانس عددی عمومی برای مقایسه‌ها و برش‌ها است. این ثابت‌ها از ماژول \lr{\texttt{qkd.constants}} نیز وارد می‌شوند (\lr{\texttt{NUMERIC\_ABS\_TOL}} و \lr{\texttt{ENTROPY\_PROB\_CLAMP}}) تا با سایر بخش‌های فریم‌ورک یکنواخت باشند. این طراحی، مقادیر برش را به صورت متمرکز قابل تنظیم می‌کند و از پراکندگی مقادیر در کد جلوگیری می‌کند.

\section{پیاده‌سازی ریاضی و الگوریتمی}

\subsection{تخصیص بودجۀ $\varepsilon$ با حاشیۀ عددی}

متد \lr{\texttt{allocate\_epsilons}} یکی از حیاتی‌ترین بخش‌های ماژول است که بودجۀ امنیتی $\varepsilon_{\text{sec}}$ را میان اجزای مختلف تقسیم می‌کند. در ادامه، نسخۀ کامل این متد را با توضیح خط به خط ارائه می‌دهیم:

\begin{LTR}
	\begin{lstlisting}[language=Python]
		def allocate_epsilons(self) -> EpsilonAllocation:
		eps_sec = self.p.eps_sec
		eps_cor = self.p.eps_cor
		remaining_eps = eps_sec - eps_cor
		if remaining_eps <= 0:
		raise ParameterValidationError(
		f"eps_sec ({eps_sec:.2e}) must be greater than "
		f"eps_cor ({eps_cor:.2e}).")
		
		# Use an absolute safety margin to handle 64-bit float rounding
		absolute_margin = 1e-20
		budget_to_alloc = remaining_eps - absolute_margin
		
		if budget_to_alloc <= 0:
		raise ParameterValidationError(
		f"Remaining epsilon budget ({remaining_eps:.2e}) "
		f"is too small after applying margin.")
		
		# Allocate 4 parts: 1 for pe, 2 for smooth, 1 for pa
		eps_pe = budget_to_alloc / 4.0
		eps_smooth = budget_to_alloc / 4.0
		eps_pa = budget_to_alloc - eps_pe - (2.0 * eps_smooth)
		
		num_pulse_configs = len(getattr(self.p.source, 'pulse_configs', []))
		denominator = max(1, 4 * num_pulse_configs + 1)
		eps_phase_est = eps_pe / denominator
		
		allocation = EpsilonAllocation(
		eps_sec=eps_sec, eps_cor=eps_cor, eps_pe=eps_pe,
		eps_smooth=eps_smooth, eps_pa=eps_pa, eps_phase_est=eps_phase_est
		)
		allocation.validate()
		return allocation
	\end{lstlisting}
\end{LTR}

در خطوط ابتدایی، $\varepsilon_{\text{sec}}$ و $\varepsilon_{\text{cor}}$ از پیکربندی خوانده می‌شوند و بودجۀ باقی‌مانده به صورت $\varepsilon_{\text{sec}} - \varepsilon_{\text{cor}}$ محاسبه می‌شود. اگر این مقدار غیرمثبت باشد (یعنی $\varepsilon_{\text{cor}} \ge \varepsilon_{\text{sec}}$)، استثنای \lr{\texttt{ParameterValidationError}} \lr{raise} می‌شود زیرا این حالت از نظر تئوری غیرممکن است. سپس یک حاشیۀ امنیتی مطلق به اندازۀ $10^{-20}$ از بودجۀ باقی‌مانده کم می‌شود. این حاشیه برای رفع یک باگ ظریف در ریاضیات ممیز شناور ۶۴-بیت ضروری است: در برخی حالت‌ها، مجموع $\varepsilon_{\text{cor}} + (\varepsilon_{\text{sec}} - \varepsilon_{\text{cor}})$ به دلیل خطای گردکردن، کمی بزرگتر از $\varepsilon_{\text{sec}}$ می‌شود و اعتبارسنجی \lr{\texttt{allocation.validate()}} شکست می‌خورد. با کم کردن این حاشیۀ کوچک، تضمین می‌شود که نامساوی به صورت عددی برقرار می‌ماند.

سپس بودجه میان چهار جزء تقسیم می‌شود: $\varepsilon_{\text{PE}}$ یک چهارم، $\varepsilon_{\text{smooth}}$ یک چهارم (که در فرمول دو بار استفاده می‌شود، یعنی مجموع $2 \cdot \varepsilon_{\text{smooth}}$ نیمی از بودجه می‌شود) و $\varepsilon_{\text{PA}}$ به عنوان باقی‌مانده. این تقسیم متقارن تضمین می‌کند که نامساوی $\varepsilon_{\text{PE}} + 2\,\varepsilon_{\text{smooth}} + \varepsilon_{\text{PA}} \le \text{budget}$ به طور دقیق برقرار باشد. در نهایت، $\varepsilon_{\text{phase\_est}}$ به صورت $\varepsilon_{\text{PE}} / (4 N_{\text{configs}} + 1)$ محاسبه می‌شود که در آن $N_{\text{configs}}$ تعداد پیکربندی‌های پالس است. این فرمول، بودجۀ تخمین خطای فاز را به تعداد رویدادهای تخمین تقسیم می‌کند و هر رویداد یک سهم مساوی دریافت می‌کند. در پایان، اعتبارسنجی نهایی با \lr{\texttt{allocation.validate()}} انجام می‌شود که در صورت بروز هرگونه ناسازگاری، استثنا \lr{raise} می‌کند.

\subsection{متد اصلی محاسبۀ طول کلید}

متد \lr{\texttt{calculate\_key\_length}} هستۀ اصلی ماژول است که تمام مراحل محاسبۀ طول کلید را هماهنگ می‌کند. در ادامه، نسخۀ کامل این متد را با توضیح بخش‌به‌بخش ارائه می‌دهیم:

\begin{LTR}
	\begin{lstlisting}[language=Python]
		def calculate_key_length(
		self,
		stats_map: Dict[str, TallyCounts],
		decoy_estimates: Optional[Union[Dict[str, Any], object]] = None
		) -> KeyCalculationResult:
		run_id = getattr(self.p, 'run_id', None)
		self.logger.info("event=key_calc_start run_id=%s", run_id)
		start_time = time.perf_counter()
		result = KeyCalculationResult()
		
		try:
		if decoy_estimates is None:
		return super().calculate_key_length(stats_map)
		
		is_decoy_protocol = bool(decoy_estimates)
		
		if is_decoy_protocol:
		signal_stats, p_sig_cfg = self._validate_stats_map_decoy(
		stats_map, result)
		if result.error_codes:
		return self._finalize_result(result, start_time)
		
		decoy_norm = self._normalize_decoy_estimates(
		decoy_estimates, result)
		if result.error_codes:
		return self._finalize_result(result, start_time)
		
		s_x_1_L_count, s_z_1_L_count, n_x, _, n_z = \
		self._compute_base_counts_decoy(
		signal_stats, p_sig_cfg, decoy_norm, result)
		e1_bit_ub = decoy_norm.e1_bit_ub
	\end{lstlisting}
\end{LTR}

در خطوط ابتدایی، یک شناسۀ اجرا (\lr{run\_id}) از پیکرببری خوانده می‌شود (اگر موجود باشد) و یک رویداد شروع در لاگ ثبت می‌شود. زمان شروع با \lr{\texttt{time.perf\_counter()}} (که دارای دقت میکروثانیه است) ثبت می‌شود تا در پایان، مدت زمان محاسبه گزارش شود. سپس یک شیء \lr{\texttt{KeyCalculationResult}} خالی ساخته می‌شود. اگر \lr{\texttt{decoy\_estimates}} برابر \lr{None} باشد، متد کلاس پایه فراخوانی می‌شود تا از فروپاشی جلوگیری شود. در غیر این صورت، بررسی می‌شود که آیا پروتکل از نوع \lr{Decoy} است یا خیر (با بررسی \lr{\texttt{bool(decoy\_estimates)}} که در صورت دیکشنری خالی یا شیء فاقد محتوا، \lr{False} برمی‌گرداند).

در مسیر \lr{Decoy-State}، ابتدا \lr{\texttt{stats\_map}} با متد \lr{\texttt{\_validate\_stats\_map\_decoy}} اعتبارسنجی می‌شود. این متد بررسی می‌کند که کلید \lr{``signal''} در \lr{\texttt{stats\_map}} موجود باشد و مقدار آن از نوع \lr{\texttt{TallyCounts}} باشد. همچنین پیکرببری پالس سیگنال را از پیکرببری منبع استخراج می‌کند و بررسی می‌کند که فیلد \lr{\texttt{mean\_photon\_number}} موجود باشد. در صورت بروز هرگونه خطا، کد خطای مناسب ثبت می‌شود و متد با فراخوانی \lr{\texttt{\_finalize\_result}} به طور زودهنگام بازمی‌گردد. سپس تخمین‌های \lr{Decoy} با متد \lr{\texttt{\_normalize\_decoy\_estimates}} نرمال‌سازی می‌شوند. این متد، ورودی‌های مختلف (شیء \lr{\texttt{DecoyEstimates}} از \lr{base.py}، دیکشنری، یا خود \lr{\texttt{DecoyEstimatesInput}}) را به یک شیء یکنواخت \lr{\texttt{DecoyEstimatesInput}} تبدیل می‌کند. سپس شمارش‌های پایۀ \lr{Decoy} با متد \lr{\texttt{\_compute\_base\_counts\_decoy}} محاسبه می‌شوند که در ادامه به تفصیل توضیح داده خواهد شد.

در مسیر \lr{non-Decoy}، کد به صورت متفاوت عمل می‌کند:

\begin{LTR}
	\begin{lstlisting}[language=Python]
		else:
		result.diagnostics['info'] = (
		"Running in non-decoy mode; using conservative estimates.")
		signal_stats = self._get_total_stats_non_decoy(
		stats_map, result)
		if result.error_codes:
		return self._finalize_result(result, start_time)
		
		# Assume key is from X-basis, phase error from Z-basis
		n_x = float(signal_stats.sifted_x)
		m_x = float(signal_stats.errors_sifted_x)
		n_z = float(signal_stats.sifted_z)
		m_z = float(signal_stats.errors_sifted_z)
		
		# Conservative: all sifted bits are single-photon
		s_x_1_L_count = n_x
		s_z_1_L_count = n_z
		
		e1_bit_ub = self._safe_divide(
		m_z, n_z, name="qber_z",
		error_codes=result.error_codes)
		result.diagnostics['intermediate'] = {
			"n_x": n_x, "m_x": m_x, "n_z": n_z, "m_z": m_z,
			"qber_z": e1_bit_ub
		}
	\end{lstlisting}
\end{LTR}

در مسیر \lr{non-Decoy}، ابتدا یک پیام اطلاعاتی در \lr{diagnostics} ثبت می‌شود تا کاربر بداند سیستم در حالت \lr{non-Decoy} اجرا می‌شود. سپس تمام آمار از \lr{\texttt{stats\_map}} با متد \lr{\texttt{\_get\_total\_stats\_non\_decoy}} ادغام می‌شود. این متد تمام شیءهای \lr{\texttt{TallyCounts}} در \lr{\texttt{stats\_map}} را با متد \lr{\texttt{merged}} ترکیب می‌کند تا یک شیء \lr{\texttt{TallyCounts}} واحد به دست آید. سپس شمارش‌های پایۀ $X$ و $Z$ استخراج می‌شوند. در این حالت، فرض محافظه‌کارانه این است که تمام بیت‌های \lr{sifted} از تک‌فوتون‌ها نشأت گرفته‌اند، یعنی $s_{X,1}^L = n_X$ و $s_{Z,1}^L = n_Z$. این فرض در حضور پالس‌های چندفوتونی ناامن است، اما برای شبیه‌سازی‌های \lr{BB84} ایده‌آل قابل قبول است. خطای بیت به طور مستقیم از \lr{QBER} پایۀ $Z$ تخمین زده می‌شود: $e_{1,\text{bit}} = m_Z / n_Z$.

سپس، در هر دو مسیر، مراحل نهایی محاسبه انجام می‌شود:

\begin{LTR}
	\begin{lstlisting}[language=Python]
		min_s_z_1_L = getattr(self.p, 'S_Z_1_L_MIN_FOR_PHASE_EST', 100)
		
		if n_z <= 0 or s_z_1_L_count < min_s_z_1_L:
		result.error_codes.append(ErrorCode.INSUFFICIENT_STATISTICS)
		return self._finalize_result(result, start_time)
		
		if is_decoy_protocol:
		n_key, m_key = signal_stats.sifted_x, signal_stats.errors_sifted_x
		else:
		n_key, m_key = n_x, m_x
		
		leak_ec = self._compute_leak_ec(n_key, m_key, result)
		e1_phase_ub = self._compute_phase_error_ub(
		s_z_1_L_count, e1_bit_ub, result)
		pa_corr_bits = self._compute_pa_terms(result)
		
		available_entropy = s_x_1_L_count * (
		1.0 - self._binary_entropy_safe(e1_phase_ub))
		
		key_len_float = available_entropy - leak_ec - pa_corr_bits
		
		self._finite_or_error("final_key_length", key_len_float, result)
		result.secure_key_length = self._clamp_key_length(key_len_float)
		self._populate_final_diagnostics(
		result, available_entropy, key_len_float, pa_corr_bits)
		
		except (ParameterValidationError, QKDSimulationError) as e:
		self.logger.error(
		"event=key_calc_error run_id=%s error='%s'", run_id, e)
		result.error_codes.append(ErrorCode.INVALID_PARAMS)
		
		return self._finalize_result(result, start_time)
	\end{lstlisting}
\end{LTR}

در این بخش، ابتدا حداقل شمارش تک‌فوتونی در پایۀ $Z$ برای تخمین خطای فاز از پیکرببری خوانده می‌شود (با مقدار پیش‌فرض ۱۰۰). اگر $n_Z \le 0$ یا $s_{Z,1}^L$ کمتر از این حداقل باشد، کد خطای \lr{\texttt{INSUFFICIENT\_STATISTICS}} ثبت می‌شود و متد به طور زودهنگام بازمی‌گردد. این بررسی تضمین می‌کند که تخمین خطای فاز به اندازۀ کافی قابل اعتماد است. سپس شمارش‌های پایۀ تولیدکنندۀ کلید ($X$) استخراج می‌شوند. نشت تصحیح خطا با متد \lr{\texttt{\_compute\_leak\_ec}}، کران بالای خطای فاز با متد \lr{\texttt{\_compute\_phase\_error\_ub}} و جریمه‌های \lr{PA} و تصحیح با متد \lr{\texttt{\_compute\_pa\_terms}} محاسبه می‌شوند. آنگاه آنتروپی موجود به صورت $s_{X,1}^L \cdot (1 - H_2(e_{\text{phase}}))$ محاسبه می‌شود و طول کلید به صورت $\ell = \text{entropy} - \text{leak} - \text{penalties}$ به دست می‌آید. در نهایت، بررسی \lr{\texttt{\_finite\_or\_error}} تضمین می‌کند که مقدار نهایی متناهی است و متد \lr{\texttt{\_clamp\_key\_length}} از کلاس پایه، مقدار را به یک عدد صحیح غیرمنفی تبدیل می‌کند. در صورت بروز هرگونه استثنای \lr{\texttt{ParameterValidationError}} یا \lr{\texttt{QKDSimulationError}}، استثنا گرفته می‌شود، در لاگ ثبت می‌گردد و کد خطای \lr{\texttt{INVALID\_PARAMS}} به \lr{result} اضافه می‌شود.

\subsection{محاسبۀ شمارش تک‌فوتونی در حالت \lr{Decoy-State}}

متد \lr{\texttt{\_compute\_base\_counts\_decoy}} شمارش‌های تک‌فوتونی را در هر دو پایۀ $X$ و $Z$ محاسبه می‌کند:

\begin{LTR}
	\begin{lstlisting}[language=Python]
		def _compute_base_counts_decoy(self, signal_stats, p_sig_cfg,
		decoy_norm, result):
		mu_s = p_sig_cfg.mean_photon_number
		p1_s = self._single_photon_prob(mu_s)
		
		# Key is from X-basis
		alice_x_prob = 1.0 - getattr(
		self.p.protocol, "alice_z_basis_prob", 0.5)
		bob_x_prob = 1.0 - getattr(
		self.p.protocol, "bob_z_basis_prob", 0.5)
		s_x_1_L_count = float(signal_stats.sent_x) * p1_s * \
		decoy_norm.Y1_L_rate * alice_x_prob * bob_x_prob
		n_x = float(signal_stats.sifted_x)
		m_x = float(signal_stats.errors_sifted_x)
		if s_x_1_L_count > n_x + NUMERIC_ABS_TOL:
		s_x_1_L_count = n_x
		
		# Phase error is from Z-basis
		alice_z_prob = getattr(self.p.protocol, "alice_z_basis_prob", 0.5)
		bob_z_prob = getattr(self.p.protocol, "bob_z_basis_prob", 0.5)
		s_z_1_L_count = float(signal_stats.sent_z) * p1_s * \
		decoy_norm.Y1_L_rate * alice_z_prob * bob_z_prob
		n_z = float(signal_stats.sifted_z)
		if s_z_1_L_count > n_z + NUMERIC_ABS_TOL:
		s_z_1_L_count = n_z
		
		result.diagnostics['intermediate'] = {
			"mu_s": mu_s, "p1_s": p1_s,
			"s_x_1_L_count": s_x_1_L_count, "n_x": n_x, "m_x": m_x,
			"s_z_1_L_count": s_z_1_L_count, "n_z": n_z
		}
		return s_x_1_L_count, s_z_1_L_count, n_x, m_x, n_z
	\end{lstlisting}
\end{LTR}

در این متد، ابتدا شدت سیگنال $\mu_s$ از پیکرببری خوانده می‌شود و احتمال ارسال تک‌فوتون $P(1 \mid \mu_s) = \mu_s \cdot e^{-\mu_s}$ با متد \lr{\texttt{\_single\_photon\_prob}} محاسبه می‌شود. سپس احتمال انتخاب پایۀ $X$ توسط آلیس و باب از پیکرببری خوانده می‌شود (با مقدار پیش‌فرض ۰.۵ در صورت عدم تنظیم). شمارش تک‌فوتونی در پایۀ $X$ به صورت $N_{\text{sent},X} \cdot P(1 \mid \mu_s) \cdot Y_1^L \cdot p_X^A \cdot p_X^B$ محاسبه می‌شود. سپس شمارش‌های \lr{sifted} و خطاهای پایۀ $X$ استخراج می‌شوند. یک بررسی محافظه‌کارانه این است که اگر $s_{X,1}^L > n_X$ شود (که از نظر فیزیکی غیرممکن است)، مقدار به $n_X$ بریده می‌شود. این برش با تلورانس \lr{\texttt{NUMERIC\_ABS\_TOL}} انجام می‌شود تا خطاهای کوچک گردکردن نادیده گرفته شوند. به طور متقابل، شمارش تک‌فوتونی در پایۀ $Z$ به صورت مشابه محاسبه می‌شود. در پایان، تمام مقادیر میانی در \lr{diagnostics} ذخیره می‌شوند تا برای تحلیل و \lr{debugging} قابل دسترس باشند.

\subsection{محاسبۀ نشت تصحیح خطا}

متد \lr{\texttt{\_compute\_leak\_ec}} نشت تصحیح خطا را محاسبه می‌کند:

\begin{LTR}
	\begin{lstlisting}[language=Python]
		def _compute_leak_ec(self, n_key_basis: float, m_key_basis: float,
		result: KeyCalculationResult) -> float:
		qber_key_basis = self._safe_divide(
		m_key_basis, n_key_basis, name="qber_key_basis",
		error_codes=result.error_codes)
		f_ec = getattr(self.p, 'f_error_correction', 1.1)
		leak_ec = f_ec * self._binary_entropy_safe(qber_key_basis) * n_key_basis
		
		di = result.diagnostics.setdefault('intermediate', {})
		di['qber_key_basis'] = qber_key_basis
		di['leak_ec_bits'] = leak_ec
		return leak_ec
	\end{lstlisting}
\end{LTR}

در این متد، ابتدا \lr{QBER} پایۀ تولیدکنندۀ کلید به صورت $E_{\text{bit}} = m_{\text{key}} / n_{\text{key}}$ محاسبه می‌شود. تابع \lr{\texttt{\_safe\_divide}} (از کلاس پایه) در صورت تقسیم بر صفر، مقدار ۰.۵ (بدترین حالت) را برمی‌گرداند و کد خطای مناسب را ثبت می‌کند. سپس بازدهی تصحیح خطا $f_{\text{EC}}$ از پیکرببری خوانده می‌شود (با مقدار پیش‌فرض ۱.۱). نشت تصحیح خطا به صورت $f_{\text{EC}} \cdot H_2(E_{\text{bit}}) \cdot n_{\text{key}}$ محاسبه می‌شود که در آن $H_2$ آنتروپی باینری است. این فرمول نشان می‌دهد که نشت متناسب با تعداد بیت‌های \lr{sifted} و آنتروپی خطا است. در پایان، مقادیر میانی در \lr{diagnostics} ذخیره می‌شوند. یک نکتۀ مهم این است که نام متغیرها در این متد به صورت \lr{\texttt{n\_key\_basis}} و \lr{\texttt{m\_key\_basis}} انتخاب شده تا به وضوح نشان دهد که این مقادیر مربوط به پایۀ تولیدکنندۀ کلید (نه لزوماً $Z$) هستند. این تغییر نام، یک اصلاح \lr{semantic} نسبت به نسخه‌های قبلی است که به اشتباه از \lr{\texttt{n\_z}} و \lr{\texttt{m\_z}} استفاده می‌کردند.

\subsection{محاسبۀ کران بالای خطای فاز}

متد \lr{\texttt{\_compute\_phase\_error\_ub}} کران بالای خطای فاز را با استفاده از نامساوی \lr{Hoeffding} محاسبه می‌کند:

\begin{LTR}
	\begin{lstlisting}[language=Python]
		def _compute_phase_error_ub(self, s_z_1_L_count: float,
		e1_bit_ub: float,
		result: KeyCalculationResult) -> float:
		di = result.diagnostics.setdefault('intermediate', {})
		if self.p.assume_phase_equals_bit_error:
		di['e1_phase_ub'] = e1_bit_ub
		return e1_bit_ub
		
		eps_ph_est_safe, _ = self._log_inv_clamp(self.eps_alloc.eps_phase_est)
		log_term = -math.log(eps_ph_est_safe)
		
		if s_z_1_L_count <= 0:
		di['e1_phase_ub'] = 0.5
		return 0.5
		
		delta = self._safe_sqrt(log_term / (2 * s_z_1_L_count))
		e1_phase_ub = e1_bit_ub + delta
		
		if e1_phase_ub > 0.5:
		e1_phase_ub = 0.5
		
		di['e1_phase_ub'] = e1_phase_ub
		di['phase_error_delta'] = delta
		return e1_phase_ub
	\end{lstlisting}
\end{LTR}

در این متد، ابتدا بررسی می‌شود که آیا پرچم \lr{\texttt{assume\_phase\_equals\_bit\_error}} در پیکرببری فعال است یا خیر. اگر فعال باشد، فرض می‌شود $e_{\text{phase}} = e_{\text{bit}}$ که یک تقریب ساده اما محافظه‌کارانه در شرایطی است که شمارش پایۀ $Z$ ناکافی است. در غیر این صورت، نامساوی \lr{Hoeffding} اعمال می‌شود. ابتدا $\varepsilon_{\text{phase\_est}}$ با متد \lr{\texttt{\_log\_inv\_clamp}} برش داده می‌شود تا از $\log(0)$ جلوگیری شود. سپس $\delta = \sqrt{-\ln(\varepsilon_{\text{phase\_est}}) / (2 \cdot s_{Z,1}^L)}$ محاسبه می‌شود. تابع \lr{\texttt{\_safe\_sqrt}} در صورت ورودی منفی، صفر برمی‌گرداند. کران بالای خطای فاز به صورت $e_{\text{phase}} = e_{\text{bit}} + \delta$ محاسبه می‌شود. اگر این مقدار از ۰.۵ بیشتر شود، به ۰.۵ بریده می‌شود زیرا آنتروپی باینری در ۰.۵ بیشینه است و مقدار بیشتر از ۰.۵ از نظر فیزیکی معنایی ندارد. در پایان، مقادیر میانی در \lr{diagnostics} ذخیره می‌شوند. این متد به طور ویژه از شمارش پایۀ $Z$ (یعنی \lr{\texttt{s\_z\_1\_L\_count}}) استفاده می‌کند زیرا خطای فاز از پایۀ $Z$ تخمین زده می‌شود؛ این یک اصلاح حیاتی نسبت به نسخه‌های قبلی است که به اشتباه از شمارش پایۀ $X$ استفاده می‌کردند.

\subsection{محاسبۀ جریمه‌های \lr{PA} و تصحیح}

متد \lr{\texttt{\_compute\_pa\_terms}} جریمه‌های تقویت حریم خصوصی و تصحیح را محاسبه می‌کند:

\begin{LTR}
	\begin{lstlisting}[language=Python]
		def _compute_pa_terms(self, result):
		eps_smooth_safe, _ = self._log_inv_clamp(self.eps_alloc.eps_smooth)
		eps_pa_safe, _ = self._log_inv_clamp(self.eps_alloc.eps_pa)
		eps_cor_safe, _ = self._log_inv_clamp(self.eps_alloc.eps_cor)
		pa_term_bits = _safe_log2(1.0 / (2.0 * eps_smooth_safe)) + \
		_safe_log2(1.0 / eps_pa_safe)
		corr_term_bits = _safe_log2(2.0 / eps_cor_safe)
		
		di = result.diagnostics.setdefault('intermediate', {})
		di['pa_corr_bits'] = pa_term_bits + corr_term_bits
		di['pa_bits'] = pa_term_bits
		di['corr_bits'] = corr_term_bits
		
		return pa_term_bits + corr_term_bits
	\end{lstlisting}
\end{LTR}

در این متد، ابتدا هر یک از $\varepsilon_{\text{smooth}}$، $\varepsilon_{\text{PA}}$ و $\varepsilon_{\text{cor}}$ با متد \lr{\texttt{\_log\_inv\_clamp}} برش داده می‌شوند. این برش تضمین می‌کند که در صورت تنظیم صفر به اشتباه، محاسبات منجر به \lr{$-\infty$} نشوند. سپس جریمۀ \lr{PA} به صورت $\log_2(1 / (2 \varepsilon_{\text{smooth}})) + \log_2(1 / \varepsilon_{\text{PA}})$ محاسبه می‌شود و جریمۀ تصحیح به صورت $\log_2(2 / \varepsilon_{\text{cor}})$. تابع \lr{\texttt{\_safe\_log2}} (در سطح ماژول) در صورت ورودی غیرمثبت، صفر برمی‌گرداند. در پایان، مجموع جریمه‌ها برگردانده می‌شود و مقادیر میانی در \lr{diagnostics} ذخیره می‌شوند. این جریمه‌ها معمولاً در حدود چند ده بیت هستند و در شبیه‌سازی‌های با کلیدهای طولانی ناچیز هستند، اما در کلیدهای کوتاه می‌توانند تأثیر قابل توجهی بر طول کلید نهایی داشته باشند.

\subsection{توابع کمکی پایداری عددی}

چندین تابع کمکی برای پایداری عددی در ماژول تعریف شده‌اند:

\begin{LTR}
	\begin{lstlisting}[language=Python]
		def _log_inv_clamp(self, val):
		return (_TIGHT_PROOF_EPS_MIN, True) if val < _TIGHT_PROOF_EPS_MIN \
		else (val, False)
		
		def _finite_or_error(self, name, val, result, non_negative=False):
		if not math.isfinite(val) or (non_negative and val < 0):
		result.error_codes.append(ErrorCode.NON_FINITE_VALUE)
		
		def _binary_entropy_safe(self, p: float) -> float:
		p_clamped = max(ENTROPY_PROB_CLAMP, min(p, 1 - ENTROPY_PROB_CLAMP))
		if p_clamped <= 0 or p_clamped >= 1: return 0.0
		return -p_clamped * math.log2(p_clamped) - \
		(1 - p_clamped) * math.log2(1 - p_clamped)
		
		def _single_photon_prob(self, mu: float) -> float:
		if mu < 0: raise ParameterValidationError(f"mu cannot be negative: {mu}")
		if mu == 0: return 0.0
		return mu * math.exp(-mu)
		
		def _safe_sqrt(self, x: float) -> float:
		if x < 0: return 0.0
		return math.sqrt(x)
		
		def _safe_log2(val: float) -> float:
		return 0.0 if val <= 0 else math.log2(val)
	\end{lstlisting}
\end{LTR}

متد \lr{\texttt{\_log\_inv\_clamp}} مقدار ورودی را به حداقل \lr{\texttt{\_TIGHT\_PROOF\_EPS\_MIN}} برش می‌دهد و یک تاپل شامل مقدار برش‌داده‌شده و یک پرچم (که نشان می‌دهد آیا برش انجام شده یا خیر) برمی‌گرداند. متد \lr{\texttt{\_finite\_or\_error}} بررسی می‌کند که آیا مقدار ورودی متناهی است (نه \lr{NaN} یا \lr{Inf}) و در صورت عدم تناهی، کد خطای \lr{\texttt{NON\_FINITE\_VALUE}} ثبت می‌کند. متد \lr{\texttt{\_binary\_entropy\_safe}} آنتروپی باینری را با برش $p$ به بازۀ $[\varepsilon, 1 - \varepsilon]$ محاسبه می‌کند تا از \lr{$\log_2(0)$} جلوگیری شود. متد \lr{\texttt{\_single\_photon\_prob}} احتمال ارسال تک‌فوتون $P(1 \mid \mu) = \mu \cdot e^{-\mu}$ را محاسبه می‌کند و در صورت $\mu < 0$، استثنا \lr{raise} می‌کند. متد \lr{\texttt{\_safe\_sqrt}} در صورت ورودی منفی، صفر برمی‌گرداند تا از خطای \lr{math domain} جلوگیری شود. تابع سطح ماژول \lr{\texttt{\_safe\_log2}} نیز در صورت ورودی غیرمثبت، صفر برمی‌گرداند. این توابع کمکی در سراسر ماژول برای تضمین پایداری عددی استفاده می‌شوند و از رفتار غیرقابل پیش‌بینی در شرایط مرزی جلوگیری می‌کنند.

\subsection{نرمال‌سازی تخمین‌های \lr{Decoy}}

متد \lr{\texttt{\_normalize\_decoy\_estimates}} ورودی‌های مختلف تخمین \lr{Decoy} را به یک شیء یکنواخت تبدیل می‌کند:

\begin{LTR}
	\begin{lstlisting}[language=Python]
		def _normalize_decoy_estimates(self, decoy_estimates, result):
		# If the input is already the correct type, return it immediately
		if isinstance(decoy_estimates, DecoyEstimatesInput):
		return decoy_estimates
		
		# Attempt to get attributes from the DecoyEstimates object (from base.py)
		Y1_L = getattr(decoy_estimates, 'yield_1_lower_bound', None)
		e1_U = getattr(decoy_estimates, 'error_rate_1_upper_bound', None)
		
		if Y1_L is None or e1_U is None:
		# Fallback for dictionary inputs
		if isinstance(decoy_estimates, dict):
		Y1_L = decoy_estimates.get('Y1_L')
		e1_U = decoy_estimates.get('e1_U')
		
		if Y1_L is None or e1_U is None:
		result.error_codes.append(ErrorCode.MISSING_DECOY_FIELDS)
		return None
		
		return DecoyEstimatesInput(float(Y1_L), float(e1_U))
	\end{lstlisting}
\end{LTR}

این متد یک طراحی انعطاف‌پذیر برای پذیرش ورودی‌های مختلف دارد. نخست، اگر ورودی از قبل از نوع \lr{\texttt{DecoyEstimatesInput}} باشد، مستقیماً برگردانده می‌شود. در غیر این صورت، تلاش می‌شود فیلدهای \lr{\texttt{yield\_1\_lower\_bound}} و \lr{\texttt{error\_rate\_1\_upper\_bound}} با \lr{\texttt{getattr}} استخراج شوند (این فیلدها در شیء \lr{\texttt{DecoyEstimates}} از \lr{base.py} موجود هستند). اگر این فیلدها موجود نباشند و ورودی یک دیکشنری باشد، تلاش می‌شود کلیدهای \lr{``Y1\_L''} و \lr{``e1\_U''} استخراج شوند (این کلیدها در برخی پیاده‌سازی‌های قدیمی استفاده می‌شدند). اگر هیچ‌کدام موجود نباشند، کد خطای \lr{\texttt{MISSING\_DECOY\_FIELDS}} ثبت می‌شود و \lr{None} برگردانده می‌شود. در غیر این صورت، یک شیء جدید \lr{\texttt{DecoyEstimatesInput}} با مقادیر تبدیل‌شده به \lr{float} ساخته می‌شود. این طراحی، ماژول را در برابر تغییرات در ساختار ورودی مقاوم می‌کند و اجازه می‌دهد با انواع مختلف شیءهای تخمین \lr{Decoy} از بخش‌های مختلف فریم‌ورک کار کند.

\subsection{نهایی‌سازی نتیجه و ثبت پروونانس}

متد \lr{\texttt{\_finalize\_result}} نتیجه نهایی را نهایی می‌کند و پروونانس را ثبت می‌کند:

\begin{LTR}
	\begin{lstlisting}[language=Python]
		def _finalize_result(self, result, start_time):
		end_time = time.perf_counter()
		result.metadata['timestamp_utc'] = datetime.now(timezone.utc).isoformat()
		result.metadata['calculation_time_ms'] = (end_time - start_time) * 1000
		result.error_codes = [
		code.value if isinstance(code, Enum) else code
		for code in result.error_codes
		]
		run_id = getattr(self.p, 'run_id', None)
		if result.error_codes:
		self.logger.warning(
		"event=key_calc_abort run_id=%s errors=%s",
		run_id, result.error_codes)
		else:
		self.logger.info(
		"event=key_calc_finish run_id=%s key_length=%d",
		run_id, result.secure_key_length)
		return result
	\end{lstlisting}
\end{LTR}

در این متد، زمان پایان با \lr{\texttt{time.perf\_counter()}} ثبت می‌شود و مدت زمان محاسبه به میلی‌ثانیه محاسبه می‌شود. یک زمان‌شناسه‌ی \lr{UTC} با فرمت \lr{ISO 8601} در \lr{metadata} ذخیره می‌شود تا برای ردیابی پروونانس قابل استفاده باشد. سپس کدهای خطا از \lr{enum} به رشته تبدیل می‌شوند تا در خروجی \lr{JSON} قابل سریال‌سازی باشند. در پایان، بر اساس وجود یا عدم وجود خطا، یک رویداد لاگ مناسب ثبت می‌شود: \lr{``key\_calc\_abort''} در صورت خطا و \lr{``key\_calc\_finish''} در صورت موفقیت. این لاگ‌ها برای مانیتورینگ و \lr{debugging} در محیط‌های تولیدی حیاتی هستند. طراحی این متد به گونه‌ای است که همیشه یک شیء \lr{\texttt{KeyCalculationResult}} کامل با پروونانس و \lr{diagnostics} برمی‌گرداند، حتی در صورت بروز خطا.

\section{ارسال و دریافت با سایر ماژول‌ها}

\subsection{ورودی‌های ماژول}

ماژول \lr{\texttt{tight.py}} برای انجام محاسبات خود به چندین ورودی از سایر ماژول‌های فریم‌ورک نیاز دارد. این ورودی‌ها به دو دستۀ \lr{پیکرببری} و \lr{آمار زمان اجرا} تقسیم می‌شوند. ورودی‌های پیکرببری در زمان ساخت شیء اثبات امنیتی خوانده می‌شوند و در طول اجرا ثابت می‌مانند، در حالی که ورودی‌های آمار در هر فراخوانی \lr{\texttt{calculate\_key\_length}} به‌روز می‌شوند.

\begin{itemize}
	\item \lr{\texttt{QKDParams}}: این شیء از ماژول \lr{\texttt{qkd.params}} وارد می‌شود و شامل تمام پارامترهای پیکرببری سیستم است. زیرشئه‌های کلیدی که این ماژول استفاده می‌کند عبارتند از: \lr{\texttt{p.eps\_sec}}، \lr{\texttt{p.eps\_cor}}، \lr{\texttt{p.f\_error\_correction}}، \lr{\texttt{p.assume\_phase\_equals\_bit\_error}}، \lr{\texttt{p.S\_Z\_1\_L\_MIN\_FOR\_PHASE\_EST}} (حداقل شمارش تک‌فوتونی پایۀ $Z$ برای تخمین خطای فاز، با مقدار پیش‌فرض ۱۰۰)، \lr{\texttt{p.run\_id}} (شناسۀ اجرا برای ردیابی پروونانس)، \lr{\texttt{p.source.pulse\_configs}} (لیست پیکرببری‌های پالس) و \lr{\texttt{p.protocol.alice\_z\_basis\_prob}} و \lr{\texttt{p.protocol.bob\_z\_basis\_prob}} (احتمال انتخاب پایۀ $Z$ توسط آلیس و باب).
	\item \lr{\texttt{stats\_map}}: این یک دیکشنری از \lr{\texttt{str}} به \lr{\texttt{TallyCounts}} است که از لایۀ شبیه‌ساز آشکارساز به اثبات امنیتی داده می‌شود. کلید \lr{``signal''} باید موجود باشد و مقدار آن یک شیء \lr{\texttt{TallyCounts}} با فیلدهای \lr{\texttt{sent\_x}}، \lr{\texttt{sent\_z}}، \lr{\texttt{sifted\_x}}، \lr{\texttt{sifted\_z}}، \lr{\texttt{errors\_sifted\_x}} و \lr{\texttt{errors\_sifted\_z}} باشد. در حالت \lr{non-Decoy}، تمام شیءهای \lr{\texttt{TallyCounts}} در \lr{\texttt{stats\_map}} با متد \lr{\texttt{merged}} ادغام می‌شوند.
	\item \lr{\texttt{decoy\_estimates}}: این پارامتر اختیاری است و می‌تواند یکی از موارد زیر باشد: یک شیء \lr{\texttt{DecoyEstimatesInput}} (با فیلدهای \lr{\texttt{Y1\_L\_rate}} و \lr{\texttt{e1\_bit\_ub}})، یک شیء \lr{\texttt{DecoyEstimates}} از \lr{base.py} (با فیلدهای \lr{\texttt{yield\_1\_lower\_bound}} و \lr{\texttt{error\_rate\_1\_upper\_bound}})، یا یک دیکشنری با کلیدهای \lr{``Y1\_L''} و \lr{``e1\_U''}. این طراحی انعطاف‌پذیر اجازه می‌دهد ماژول با انواع مختلف شیءهای تخمین \lr{Decoy} از بخش‌های مختلف فریم‌ورک کار کند. اگر این پارامتر \lr{None} باشد، متد کلاس پایه فراخوانی می‌شود.
	\item \lr{\texttt{EpsilonAllocation}}: این شیء از متد \lr{\texttt{allocate\_epsilons}} تولید می‌شود و شامل فیلدهای $\varepsilon_{\text{sec}}$، $\varepsilon_{\text{cor}}$، $\varepsilon_{\text{pe}}$، $\varepsilon_{\text{smooth}}$، $\varepsilon_{\text{pa}}$ و $\varepsilon_{\text{phase\_est}}$ است. این شیء در محاسبۀ جریمه‌های \lr{PA} و خطای فاز استفاده می‌شود.
\end{itemize}

\subsection{خروجی‌های ماژول}

خروجی اصلی ماژول یک شیء \lr{\texttt{KeyCalculationResult}} است که شامل چهار فیلد است:

\begin{enumerate}
	\item \lr{\texttt{secure\_key\_length}}: طول کلید امن به صورت یک عدد صحیح غیرمنفی. این فیلد خروجی نهایی محاسبه است و توسط لایۀ بالایی برای تولید کلید واقعی استفاده می‌شود. در صورت بروز هرگونه خطا، این مقدار صفر است.
	\item \lr{\texttt{diagnostics}}: یک دیکشنری شامل تمام مقادیر میانی محاسبه. این دیکشنری شامل یک کلید \lr{``intermediate''} است که خود یک دیکشنری با مقادیری مانند $\mu_s$، $P(1 \mid \mu_s)$، $s_{X,1}^L$، $n_X$، $m_X$، $s_{Z,1}^L$، $n_Z$، \lr{QBER}، نشت تصحیح خطا، خطای فاز، $\delta$ خطای فاز، جریمه‌های \lr{PA} و تصحیح، آنتروپی موجود و طول کلید اعشاری است. این فیلد برای تحلیل و بازبینی امنیتی ضروری است.
	\item \lr{\texttt{error\_codes}}: لیستی از کدهای خطا (به صورت رشته) که در صورت بروز مشکل، علت آن را نشان می‌دهند. اگر این لیست خالی باشد، محاسبه با موفقیت انجام شده است.
	\item \lr{\texttt{metadata}}: یک دیکشنری شامل اطلاعاتی مانند زمان‌شناسه‌ی \lr{UTC} (با فرمت \lr{ISO 8601}) و مدت زمان محاسبه به میلی‌ثانیه. این فیلد برای ردیابی پروونانس و مانیتورینگ عملکرد استفاده می‌شود.
\end{enumerate}

\subsection{وابستگی‌های داخلی و خارجی}

ماژول وابستگی‌های داخلی و خارجی متعددی دارد. وابستگی‌های داخلی شامل موارد زیر است:

\begin{itemize}
	\item \lr{\texttt{qkd.proofs.lim2014.Lim2014Proof}}: کلاس پایه که \lr{\texttt{BB84TightProof}} از آن ارث می‌برد. این کلاس زیرساخت‌های مشترک مانند محاسبۀ آنتروپی باینری، تخصیص پایۀ \lr{$\varepsilon$}، ابزارهای \lr{clamp} عددی و متد \lr{\texttt{\_safe\_divide}} را فراهم می‌سازد.
	\item \lr{\texttt{qkd.datatypes.TallyCounts}} و \lr{\texttt{qkd.datatypes.EpsilonAllocation}}: این دو \lr{dataclass} به ترتیب برای شمارش‌های آماری و تخصیص بودجۀ \lr{$\varepsilon$} استفاده می‌شوند. \lr{\texttt{TallyCounts}} شامل فیلدهای \lr{\texttt{sent\_x}}، \lr{\texttt{sent\_z}}، \lr{\texttt{sifted\_x}}، \lr{\texttt{sifted\_z}}، \lr{\texttt{errors\_sifted\_x}} و \lr{\texttt{errors\_sifted\_z}} است و متد \lr{\texttt{merged}} برای ادغام چندین شیء \lr{\texttt{TallyCounts}} را فراهم می‌سازد.
	\item \lr{\texttt{qkd.constants.NUMERIC\_ABS\_TOL}} و \lr{\texttt{qkd.constants.ENTROPY\_PROB\_CLAMP}}: این دو ثابت برای تلورانس عددی و برش احتمال در محاسبۀ آنتروپی باینری استفاده می‌شوند.
\end{itemize}

وابستگی‌های خارجی شامل \lr{\texttt{math}} برای توابع ریاضی پایه (مانند \lr{\texttt{log}}، \lr{\texttt{log2}}، \lr{\texttt{exp}}، \lr{\texttt{sqrt}}، \lr{\texttt{isfinite}})، \lr{\texttt{logging}} برای ثبت رویدادها، \lr{\texttt{time}} برای اندازه‌گیری مدت محاسبه، و \lr{\texttt{datetime}} برای تولید زمان‌شناسه‌ی \lr{UTC} هستند. همچنین از ماژول‌های استاندارد \lr{Python} مانند \lr{\texttt{typing}}، \lr{\texttt{dataclasses}} و \lr{\texttt{enum}} برای ساختاردهی کد استفاده می‌شود.

\subsection{جریان داده در سراسر فریم‌ورک}

جریان داده کلی به این شکل است:
\begin{enumerate}
	\item لایۀ شبیه‌ساز پالس‌های نوری را با شدت‌های $\mu$ (سیگنال)، $\nu$ (\lr{decoy} ضعیف) و 0 (خلأ) تولید و از کانال کوانتومی عبور می‌دهد.
	\item آشکارسازها در سمت باب، آشکارسازی‌ها و خطاها را در دو پایۀ $X$ و $Z$ و برای هر شدت ثبت کرده و در \lr{\texttt{TallyCounts}} ذخیره می‌کنند.
	\item موتور تخمین \lr{Decoy} (مانند \lr{\texttt{Ma2005VacuumWeakProof}} یا \lr{\texttt{Lim2014Proof}}) با دریافت این شمارش‌ها، تخمین‌های $Y_1$ و $e_1$ را به صورت یک شیء \lr{\texttt{DecoyEstimates}} تولید می‌کند.
	\item متد \lr{\texttt{calculate\_key\_length}} با دریافت \lr{\texttt{stats\_map}} و \lr{\texttt{decoy\_estimates}} فراخوانی می‌شود. این متد ابتدا \lr{\texttt{stats\_map}} را اعتبارسنجی می‌کند و سپس تخمین‌های \lr{Decoy} را با متد \lr{\texttt{\_normalize\_decoy\_estimates}} نرمال‌سازی می‌کند.
	\item سپس شمارش‌های تک‌فوتونی در هر دو پایه با متد \lr{\texttt{\_compute\_base\_counts\_decoy}} محاسبه می‌شوند.
	\item نشت تصحیح خطا با متد \lr{\texttt{\_compute\_leak\_ec}}، کران بالای خطای فاز با متد \lr{\texttt{\_compute\_phase\_error\_ub}} و جریمه‌های \lr{PA} و تصحیح با متد \lr{\texttt{\_compute\_pa\_terms}} محاسبه می‌شوند.
	\item آنتروپی موجود به صورت $s_{X,1}^L \cdot (1 - H_2(e_{\text{phase}}))$ محاسبه می‌شود و طول کلید به صورت $\ell = \text{entropy} - \text{leak} - \text{penalties}$ به دست می‌آید.
	\item در نهایت، متد \lr{\texttt{\_finalize\_result}} نتیجه را نهایی می‌کند و پروونانس (زمان‌شناسه و مدت محاسبه) را ثبت می‌کند.
	\item لایۀ بالایی می‌تواند با بررسی \lr{\texttt{error\_codes}} وضعیت موفقیت یا شکست را تشخیص دهد و در صورت موفقیت، از \lr{\texttt{secure\_key\_length}} برای تولید کلید واقعی استفاده کند.
\end{enumerate}

این جداسازی مراحل به گونه‌ای طراحی شده که هر مرحله قابل آزمایش و بازبینی مستقل باشد. به عنوان مثال، می‌توان \lr{\texttt{DecoyEstimatesInput}} را به صورت دستی ساخت و رفتار مرحلۀ نهایی محاسبۀ کلید را در شرایط خاص بررسی کرد. این ماژول‌بندی از اصل \lr{Single Responsibility} پیروی می‌کند و قابلیت نگهداری و توسعه کد را به طور قابل توجهی افزایش می‌دهد. علاوه بر این، شفافیت کامل در \lr{diagnostics} به کاربر اجازه می‌دهد هر مرحله از محاسبه را ردیابی کرده و در صورت نیاز، فرضیات تئوریک را برای کاربرد خاص خود بازبینی کند. طراحی \lr{fail-safe} این ماژول، آن را به انتخابی ایده‌آل برای محیط‌های تولیدی تبدیل می‌کند که در آن‌ها پایداری و قابل اعتمادبودن سیستم حیاتی است.

	
	\chapter{بهینه‌سازی پارامترها با استفاده از برنامه‌ریزی خطی (\lr{LP})}
	\label{ch:optimization}
	\sloppy
\emergencystretch=5em
\hbadness=10000
\vbadness=10000
\tolerance=2000
\providecolor{codebg}{RGB}{248,248,248}
\providecolor{codeframe}{RGB}{200,200,200}
\providecolor{keyword}{RGB}{0,0,180}
\providecolor{string}{RGB}{163,21,21}
\providecolor{comment}{RGB}{0,128,0}
\lstset{
	basicstyle=\ttfamily\scriptsize,
	breaklines=true,
	breakatwhitespace=false,
	columns=fullflexible,
	keepspaces=true,
	frame=single,
	framesep=2pt,
	showstringspaces=false,
	backgroundcolor=\color{codebg},
	rulecolor=\color{codeframe},
	keywordstyle=\color{keyword}\bfseries,
	stringstyle=\color{string},
	commentstyle=\color{comment}\itshape,
	numberstyle=\tiny\color{gray},
	numbers=left,
	numbersep=6pt,
	xleftmargin=14pt,
	framexleftmargin=14pt,
	literate=
	{ą}{{\v{a}}}1 {ę}{{\k{e}}}1 {ż}{{\.z}}1 {ś}{{\'s}}1
	{ł}{{\l}}1 {ć}{{\'c}}1 {ń}{{\'n}}1 {ó}{{\'o}}1 {ź}{{\'z}}1
	{→}{{$\rightarrow$}}1 {←}{{$\leftarrow$}}1 {≠}{{$\neq$}}1
}

\section{نمای کلی و هدف ماژول}

فایل \lr{\texttt{qkd/proofs/optimization.py}} مجموعه‌ای از ابزارهای بهینه‌سازی پارامتر برای پروتکل‌های مختلف توزیع کلید کوانتومی (\lr{QKD}) را پیاده‌سازی می‌کند. این ماژول بر اساس دو چارچوب مرجع توسعه یافته است: چارچوب \lr{Ma et al. 2005} برای پروتکل‌های \lr{Decoy-State BB84} و چارچوب \lr{Chapman et al. 2018} برای پروتکل‌های \lr{state-dependent} مبتنی بر \lr{Amplitude Damping Channel (ADC)}. هدف اصلی این ماژول، پیدا کردن مقادیر بهینۀ پارامترهای قابل کنترل آلیس (مانند شدت سیگنال $\mu$، شدت \lr{decoy} $\nu$، احتمال انتخاب سیگنال $P_s$، احتمال انتخاب \lr{decoy} $P_d$، و زاویۀ چرخش کدگذاری $\theta$) به گونه‌ای است که نرخ کلید امن نهایی به حداکثر برسد. این بهینه‌سازی در عمل برای استقرار سیستم‌های \lr{QKD} حیاتی است، زیرا انتخاب پارامترهای غیربهینه می‌تواند نرخ کلید را تا چندین براکم کند یا حتی آن را به صفر برساند.

اهمیت این ماژول از این واقعیت ناشی می‌شود که نرخ کلید امن $R$ در پروتکل‌های \lr{QKD} یک تابع غیرخطی و معمولاً غیرمحدب (\lr{non-convex}) از پارامترهای سیستم است. به عنوان مثال، در پروتکل \lr{Decoy-State BB84}، نرخ کلید به صورت $R = q\{Q_1[1 - H_2(e_1)] - Q_\mu f(E_\mu) H_2(E_\mu)\}$ محاسبه می‌شود که در آن $Q_1 = Y_1 \mu e^{-\mu}$ و $Q_\mu = \sum_n Y_n P(n|\mu)$ هر دو به صورت غیرخطی به $\mu$ وابسته‌اند. افزایش $\mu$ باعث افزایش $Q_1$ (و در نتیجه افزایش ترم تقویت حریم خصوصی) می‌شود، اما همزمان باعث افزایش $Q_\mu$ (و در نتیجه افزایش نشت تصحیح خطا) و افزایش آسیب‌پذیری به حمله \lr{PNS} نیز می‌شود. بنابراین یک نقطۀ بهینۀ $\mu^*$ وجود دارد که در آن توازن میان این دو ترم بهینه است. پیدا کردن این نقطه به صورت تحلیلی به دلیل پیچیدگی فرمول‌ها ممکن نیست و باید از روش‌های عددی استفاده شود.

این ماژول ده تابع بهینه‌سازی مستقل ارائه می‌کند که هر کدام برای یک سناریوی خاص طراحی شده‌اند. این توابع عبارتند از: \lr{\texttt{optimize\_decoy\_intensities}} (بهینه‌سازی اکتشافی شدت سیگنال و \lr{decoy} بر اساس مدل نرخ ساده‌شده)، \lr{\texttt{optimize\_rotation\_angle}} (بهینه‌سازی زاویۀ چرخش کدگذاری برای یک کانال عمومی)، \lr{\texttt{calculate\_optimal\_mu\_ma2005}} (محاسبۀ تحلیلی-عددی $\mu$ بهینه بر اساس معادلۀ 12 مقالۀ \lr{Ma 2005})، \lr{\texttt{optimize\_nu\_ma2005}} (بهینه‌سازی $\nu$ برای $N$ متناهی بر اساس بخش 4.3 مقاله)، \lr{\texttt{optimize\_pulse\_distribution\_ma2005}} (بهینه‌سازی احتمالات توزیع پالس با روش \lr{SLSQP})، \lr{\texttt{optimize\_pulse\_probabilities}} (پوشش‌دهنده برای تابع قبلی)، \lr{\texttt{optimize\_simultaneous\_ma2005}} (بهینه‌سازی همزمان $P_s$، $P_d$ و $\nu$ با استفاده از کلاس \lr{\texttt{Ma2005GlobalOptimizer}} و گرادیان عددی)، \lr{\texttt{optimize\_global\_ma2005}} (پوشش‌دهنده برای تابع قبلی)، \lr{\texttt{optimize\_lim2014\_finite\_key}} (بهینه‌سازی پنج پارامتر چارچوب \lr{Lim 2014} برای کلید متناهی با روش \lr{Nelder-Mead})، و \lr{\texttt{optimize\_rotation\_angle\_chapman}} (بهینه‌سازی زاویۀ چرخش برای \lr{Coherent Scheme} در \lr{Chapman 2018} با محاسبۀ فاصلۀ \lr{trace distance} میان حالت‌های خروجی کانال).

طراحی این ماژول بر اساس اصل جداسازی دغدغه‌ها (\lr{Separation of Concerns}) استوار است؛ به این ترتیب که توابع بهینه‌سازی صرفاً مسئول پیدا کردن پارامترهای بهینه هستند و محاسبۀ نرخ کلید (یا هر کمیت هدف دیگر) به یک فراخوان (\lr{callback}) تابعی سپرده می‌شود که توسط کاربر ارائه می‌گردد. این طراحی چندین مزیت دارد: 

نخست، ماژول بهینه‌سازی از جزئیات پیاده‌سازی اثبات امنیتی مستقل است و می‌تواند با هر چارچوب اثبات (مانند \lr{Ma 2005}، \lr{Lim 2014} یا \lr{Tight}) کار کند؛ دوم، کاربر می‌تواند تابع هدف خود را به سادگی جایگزین کند (مثلاً به جای نرخ کلید، \lr{QBER} را کمینه کند یا طول کلید متناهی را بیشینه سازد) بدون آنکه نیازی به تغییر کد بهینه‌سازی باشد؛ سوم، فرآیند تست‌کردن و ارزیابی آسان‌تر می‌شود، زیرا می‌توان از \lr{callback}‌های مصنوعی با رفتار مشخص و شناخته‌شده برای اعتبارسنجی توابع بهینه‌سازی استفاده کرد. این الگو که به تزریق وابستگی (\lr{«dependency injection»}) نیز معروف است، یک اصل کلیدی در مهندسی نرم‌افزار به شمار می‌رود که قابلیت نگهداری و توسعهٔ کد را به طور قابل‌توجهی افزایش می‌دهد.


در صورت شکست بهینه‌سازی (به هر دلیل: عدم همگرایی، ناحیه غیرممکن، یا خطای عددی)، توابع این ماژول به جای \lr{raise} کردن استثنا، مقادیر پیش‌فرض امن و محافظه‌کارانه را برمی‌گردانند. این رفتار \lr{fail-safe} تضمین می‌کند که خط لوله شبیه‌سازی متوقف نشود و کاربر بتواند با مقادیر پیش‌فرض ادامه دهد. به عنوان مثال، اگر بهینه‌سازی $\mu$ شکست بخورد، مقدار پیش‌فرض 0.48 (برای $f \le 1.05$) یا 0.54 (برای $f > 1.05$) برگردانده می‌شود که مقادیر تجربی خوبی برای پروتکل‌های \lr{Decoy-State BB84} هستند. این مقادیر پیش‌فرض بر اساس بخش 3.1 مقالۀ \lr{Ma 2005} انتخاب شده‌اند.

\section{مبانی تئوری و فیزیکی}

\subsection{نرخ کلید \lr{GLLP} و وابستگی غیرخطی به $\mu$}

فرمول استاندارد نرخ کلید در چارچوب \lr{Gottesman-Lo-Lütkenhaus-Preskill (GLLP)} به صورت زیر است:
\begin{equation}
	R = q\left\{Q_1\left[1 - H_2(e_1)\right] - Q_\mu\,f(E_\mu)\,H_2(E_\mu)\right\},
\end{equation}
که در آن $q$ فاکتور پایه‌گذاری (برای \lr{BB84} برابر \lr{0.5})، $Q_1 = Y_1 \mu e^{-\mu}$ \lr{Gain} تک‌فوتونی، $Q_\mu = \sum_n Y_n P(n|\mu)$ \lr{Gain} کلی سیگنال، $e_1$ خطای تک‌فوتونی، $E_\mu$ \lr{QBER} سیگنال، $f(E_\mu)$ بازدهی تصحیح خطا و $H_2(\cdot)$ آنتروپی باینری است. ترم اول $Q_1[1 - H_2(e_1)]$ نشان‌دهندۀ آنتروفی قابل استخراج از تک‌فوتون‌های امن است و ترم دوم $Q_\mu f(E_\mu) H_2(E_\mu)$ نشان‌دهندۀ نشت تصحیح خطا است. 

برای پیدا کردن $\mu$ بهینه، باید $\partial R / \partial \mu = 0$ را حل کنیم. با فرض ساده‌کنندۀ $e_1 \approx E_\mu$ و $Y_1 \approx \eta$ (در حد نویز کم)، این مشتق‌گیری به یک معادلۀ تحلیلی-عددی منجر می‌شود که در معادلۀ (12) مقالۀ \lr{Ma 2005} ارائه شده است. این معادله به صورت زیر فرمول‌بندی می‌شود:
\begin{equation}
	(1 - \mu)\,e^{-\mu} = C, \qquad C = \frac{f_{\text{EC}}\,H_2(e_{\text{det}})}{1 - H_2(e_{\text{det}})},
\end{equation}
که در آن $e_{\text{det}}$ خطای ذاتی آشکارساز است. تابع $(1-\mu)e^{-\mu}$ در $\mu = 0$ برابر 1 است و به طور یکنواخت تا $\mu \to \infty$ کاهش می‌یابد. بنابراین برای هر $C \in (0, 1)$ یک جواب یکتا در بازۀ $(0, 1)$ وجود دارد. اگر $C \ge 1$ (یعنی خطای آشکارساز خیلی بزرگ باشد)، جوابی وجود ندارد و نرخ کلید منفی می‌شود؛ در این حالت، ماژول به مقدار پیش‌فرض $\mu = 0.48$ برمی‌گردد. این محاسبۀ تحلیلی-عددی در تابع \lr{\texttt{calculate\_optimal\_mu\_ma2005}} با استفاده از روش \lr{Brent's method} (\lr{\texttt{brentq}} از \lr{SciPy}) پیاده‌سازی شده است.

\subsection{بهینه‌سازی $\nu$ برای کلید متناهی}

در حالت کلید متناهی ($N < \infty$)، نرخ کلید به اندازۀ نمونۀ آماری \lr{decoy} حساس می‌شود. اگر $\nu$ خیلی کوچک باشد، تعداد آشکارسازی‌های \lr{decoy} $N_d = N \cdot P_d \cdot Q_\nu$ کم می‌شود و کران‌های آماری برای $Y_1$ و $e_1$ سخت‌گیرانه‌تر (یعنی بدتر) می‌شوند. اگر $\nu$ خیلی بزرگ باشد، شباهت آن به $\mu$ زیاد می‌شود و قابلیت تفکیک آماری بین سیگنال و \lr{decoy} کاهش می‌یابد. بنابراین یک $\nu^*$ بهینه وجود دارد. بخش 4.3 مقالۀ \lr{Ma 2005} این بهینه‌سازی را به صورت یک مینیمم‌سازی اسکالر روی $\nu$ با کران‌های $[\nu_{\min}, \mu \cdot 0.9]$ فرمول‌بندی می‌کند. در کد، این بهینه‌سازی در تابع \lr{\texttt{optimize\_nu\_ma2005}} با استفاده از روش \lr{Brent's bounded method} (\lr{\texttt{minimize\_scalar}} با \lr{\texttt{method='bounded'}}) پیاده‌سازی شده است. این روش برای توابع یک‌بعدی هموار بهینه است و همگرایی تضمین‌شده‌ای در بازۀ مشخص‌شده دارد.

\subsection{بهینه‌سازی توامان $P_s$، $P_d$ و $\nu$}

بهینه‌سازی توامان سه پارامتر $P_s$، $P_d$ و $\nu$ یک مسألۀ بهینه‌سازی چندبعدی است که نیاز به روش‌های \lr{gradient-based} دارد. در کد، این مسأله با استفاده از روش \lr{SLSQP (Sequential Least SQuares Programming)} از \lr{SciPy} حل می‌شود. این روش برای مسائل بهینه‌سازی مقید با محدودیت‌های نامساوی مناسب است و از گرادیان تابع هدف برای هدایت جستجو استفاده می‌کند. کلاس \lr{\texttt{Ma2005GlobalOptimizer}} برای این منظور طراحی شده و یک متد \lr{\texttt{gradient}} برای محاسبۀ گرادیان به روش تفاضل محدود (\lr{finite difference}) ارائه می‌کند. 

گرادیان به روش تفاضل محدود مرکزی به صورت زیر محاسبه می‌شود:
\begin{equation}
	\frac{\partial f}{\partial x_i} \approx \frac{f(x + \epsilon\,e_i) - f(x - \epsilon\,e_i)}{2\,\epsilon},
\end{equation}
که در آن $e_i$ بردار واحد در جهت $i$-ام و $\epsilon$ گام تفاضل محدود است (در کد $\epsilon = 10^{-8}$). این روش برای $P_s$ و $P_d$ به صورت تفاضل پیشرو (\lr{forward difference}) و برای $\nu$ به صورت تفاضل مرکزی پیاده‌سازی شده است. این انتخاب به این دلیل است که $P_s$ و $P_d$ در محدودیت‌های سخت ($P_s + P_d \le 1$) شرکت می‌کنند و تفاضل پیشرو در نزدیک مرز محدودیت پایدارتر است، در حالی که $\nu$ یک متغیر آزاد است و تفاضل مرکزی دقت بالاتری دارد.

\subsection{کانال \lr{Amplitude Damping} و \lr{Chapman 2018}}

در پروتکل‌های \lr{state-dependent QKD} که توسط \lr{Chapman et al. 2018} پیشنهاد شده‌اند، آلیس اطلاعات را در حالت‌های کوانتومی $\rho_0$ و $\rho_1$ کدگذاری می‌کند و این حالت‌ها از یک کانال \lr{Amplitude Damping Channel (ADC)} با پارامتر $\gamma$ عبور می‌کنند. \lr{ADC} کانالی است که افت انرژی (مانند واهلش آرام اتم‌ها) را مدل می‌کند و عملگرهای \lr{Kraus} آن به صورت زیر تعریف می‌شوند:
\begin{align}
	E_0 &= \begin{pmatrix} 1 & 0 \\ 0 & \sqrt{1 - \gamma} \end{pmatrix}, \\
	E_1 &= \begin{pmatrix} 0 & \sqrt{\gamma} \\ 0 & 0 \end{pmatrix}.
\end{align}
حالت خروجی کانال از رابطۀ $\rho_{\text{out}} = E_0 \rho_{\text{in}} E_0^\dagger + E_1 \rho_{\text{in}} E_1^\dagger$ محاسبه می‌شود. در \lr{Coherent Scheme}، آلیس حالت‌های ورودی $|0\rangle\langle 0|$ و $|1\rangle\langle 1|$ را با یک چرخش \lr{Y} به زاویۀ $\theta$ تبدیل می‌کند:
\begin{equation}
	\rho_i^{\text{in}}(\theta) = R_y(\theta)\,|i\rangle\langle i|\,R_y(\theta)^\dagger, \qquad R_y(\theta) = \begin{pmatrix} \cos(\theta/2) & -\sin(\theta/2) \\ \sin(\theta/2) & \cos(\theta/2) \end{pmatrix}.
\end{equation}
احتمال موفقیت تشخیص حالت توسط باب متناسب با فاصلۀ \lr{trace distance} میان $\rho_0^{\text{out}}$ و $\rho_1^{\text{out}}$ است:
\begin{equation}
	P_{\text{succ}} = \frac{1}{2}\left(1 + D(\rho_0^{\text{out}}, \rho_1^{\text{out}})\right), \qquad D(\rho_0, \rho_1) = \frac{1}{2}\,\text{Tr}\left|\rho_0 - \rho_1\right|,
\end{equation}
که در آن $|A| = \sqrt{A^\dagger A}$ و \lr{Tr} عملگر \lr{trace} است. برای ماتریس‌های \lr{Hermitian}، $D$ را می‌توان به صورت $D = \frac{1}{2}\sum_i |\lambda_i|$ محاسبه کرد که در آن $\lambda_i$ مقادیر ویژۀ $\rho_0 - \rho_1$ هستند. بهینه‌سازی زاویۀ $\theta$ به گونه‌ای که $D$ بیشینه شود، در تابع \lr{\texttt{optimize\_rotation\_angle\_chapman}} پیاده‌سازی شده است.

\subsection{بهینه‌سازی کلید متناهی \lr{Lim 2014}}

چارچوب \lr{Lim et al. 2014} پنج پارامتر آزاد دارد: $q_x$ (احتمال پایۀ $X$)، $P_{\mu_1}$ (احتمال سیگنال)، $P_{\mu_2}$ (احتمال \lr{decoy} ضعیف)، $\mu_1$ (شدت سیگنال) و $\mu_2$ (شدت \lr{decoy}). بهینه‌سازی نرخ کلید $R = \ell / N$ روی این پنج پارامتر یک مسألۀ پیچیده است زیرا چشم‌انداز (\lr{landscape}) تابع هدف معمولاً ناهموار (\lr{non-smooth}) است و دارای چندین محلی است. در کد، این مسأله با روش \lr{Nelder-Mead} (یک روش بدون گرادیان) حل می‌شود که برای توابع ناهموار مقاوم‌تر از روش‌های مبتنی بر گرادیان است. این روش از یک \lr{simplex} (چندوجهی محدب) برای کاوش فضای پارامتر استفاده می‌کند و با منبسط‌کردن، منقبض‌کردن و انعکاس دادن \lr{simplex}، به سمت نقطۀ بهینه حرکت می‌کند. حداکثر تعداد تکرارها در کد 200 تنظیم شده و تلورانس $10^{-3}$ برای توقف استفاده می‌شود.

\subsection{محدودیت‌های فیزیکی و منطق جریمه}

در تمام توابع بهینه‌سازی این ماژول، محدودیت‌های فیزیکی به صورت جریمه‌های (\lr{penalty}) بزرگ (مانند $10^9$) پیاده‌سازی شده‌اند. این محدودیت‌ها عبارتند از: $\mu > \nu$ (شدت سیگنال باید بزرگتر از \lr{decoy} باشد)، $\nu > 0$ (شدت باید مثبت باشد)، $\mu \le 1$ (برای جلوگیری از پالس‌های بسیار چندفوتونی)، $P_s + P_d < 1$ (احتمالات باید نرمالیزه شوند و فضایی برای $P_v$ باقی بماند)، و $\nu < 0.9 \mu$ (برای حفظ قابلیت تفکیک آماری). استفاده از جریمه به جای محدودیت‌های سخت به روش \lr{Nelder-Mead} اجازه می‌دهد در ناحیه‌های غیرممکن هم کاوش کند (با جریمه بالا) و در عین حال به روش \lr{SLSQP} (که از محدودیت‌های صریح پشتیبانی می‌کند) انعطاف‌پذیری بدهد. این طراحی ترکیبی به ماژول اجازه می‌دهد با انواع روش‌های بهینه‌سازی کار کند بدون آنکه نیاز به بازنویسی تابع هدف باشد.

\section{تحلیل معماری و ساختار داده‌ها}

\subsection{ساختار کلی ماژول}

ماژول \lr{\texttt{optimization.py}} بر اساس یک طراحی تابعی (\lr{functional design}) استوار است که در آن نه توابع مستقل بهینه‌سازی ارائه می‌شوند و یک کلاس کمکی \lr{\texttt{Ma2005GlobalOptimizer}} برای بهینه‌سازی توامان استفاده می‌شود. این طراحی بر اساس اصل \lr{``functions first, classes when needed''} است که در آن توابع ساده برای کاربردهای یک‌بار مصرف و کلاس‌ها برای حالت‌هایی که نیاز به نگهداری \lr{state} (مانند \lr{callback} و گام تفاضل محدود) هستند، استفاده می‌شوند. تمام توابع در لیست \lr{\texttt{\_\_all\_\_}} صریحاً اعلام شده‌اند تا به کاربر یک رابط واضح از API عمومی ماژول ارائه دهند.

\subsection{کلاس \lr{Ma2005GlobalOptimizer}}

کلاس \lr{\texttt{Ma2005GlobalOptimizer}} برای بهینه‌سازی توامان $P_s$، $P_d$ و $\nu$ طراحی شده است. در ادامه، ساختار این کلاس را به صورت خلاصه نمایش می‌دهیم:

\begin{LTR}
	\begin{lstlisting}[language=Python]
		class Ma2005GlobalOptimizer:
		"""
		Implements the global optimization for Vacuum+Weak decoy protocol
		using analytical derivatives derived from Eqs. (34) and (37) w.r.t nu
		and Ps, Pd.
		"""
		def __init__(self, mu_target, n_total, cost_callback):
		self.mu = mu_target
		self.N = n_total
		self.callback = cost_callback
		self.epsilon = 1e-8  # small finite diff step
	\end{lstlisting}
\end{LTR}

سازندۀ این کلاس چهار پارامتر دریافت می‌کند: \lr{\texttt{mu\_target}} (شدت هدف سیگنال $\mu$)، \lr{\texttt{n\_total}} (تعداد کل پالس‌های $N$)، \lr{\texttt{cost\_callback}} (تابع هدف که نرخ منفی کلید را محاسبه می‌کند) و \lr{\texttt{epsilon}} (گام تفاضل محدود با مقدار پیش‌فرض $10^{-8}$). این کلاس به صورت یک \lr{stateful wrapper} عمل می‌کند که \lr{callback} و پارامترهای سیستم را نگه می‌دارد و متدهای \lr{\texttt{objective}} و \lr{\texttt{gradient}} را برای استفاده توسط \lr{solver} \lr{SciPy} ارائه می‌دهد. این طراحی از الگوی \lr{``Strategy''} پیروی می‌کند که در آن الگوریتم بهینه‌سازی (\lr{SLSQP}) از تابع هدف (که توسط \lr{callback} تعریف می‌شود) مستقل است.

\subsection{متد \lr{\texttt{objective}} و جریمه‌ها}

متد \lr{\texttt{objective}} تابع هدف را برای \lr{solver} فراهم می‌کند:

\begin{LTR}
	\begin{lstlisting}[language=Python]
		def objective(self, x):
		"""Objective function for minimization (-Rate)."""
		p_s, p_d, nu = x
		if p_s <= 0.01 or p_d <= 0.01 or (p_s + p_d) >= 1.0:
		return 1e9
		if nu <= 0.001 or nu >= self.mu * 0.99:
		return 1e9
		# Call the provided cost function (which computes -Rate)
		return self.callback(p_s, p_d, nu)
	\end{lstlisting}
\end{LTR}

در این متد، ابتدا محدودیت‌های فیزیکی بررسی می‌شوند: $P_s > 0.01$، $P_d > 0.01$، $P_s + P_d < 1$، $\nu > 0.001$ و $\nu < 0.99 \mu$. اگر هر کدام از این محدودیت‌ها نقض شود، یک جریمۀ بزرگ ($10^9$) برگردانده می‌شود. این جریمه به \lr{solver} می‌گوید که این ناحیه از فضای پارامتر نامطلوب است و باید از آن دوری کند. انتخاب کران‌های 0.01 و 0.001 به این دلیل است که مقادیر خیلی کوچک از نظر آماری معنایی ندارند (تعداد پالس‌های \lr{decoy} خیلی کم می‌شود) و از نظر عددی نیز ناپایدار هستند. اگر محدودیت‌ها نقض نشوند، \lr{callback} فراخوانی می‌شود و مقدار تابع هدف (نرخ منفی کلید) برگردانده می‌شود. این طراحی به کاربر اجازه می‌دهد تابع هدف خود را بدون نگرانی در مورد محدودیت‌های فیزیکی پیاده‌سازی کند.

\subsection{متد \lr{\texttt{gradient}} با تفاضل محدود}

متد \lr{\texttt{gradient}} گرادیان تابع هدف را به روش تفاضل محدود محاسبه می‌کند:

\begin{LTR}
	\begin{lstlisting}[language=Python]
		def gradient(self, x):
		"""Calculates the gradient [d(-R)/dPs, d(-R)/dPd, d(-R)/dnu]."""
		p_s, p_d, nu = x
		
		# 1. Calculate current rate
		val_center = self.objective(x)
		if val_center >= 1e8:
		return np.zeros(3)  # Infeasible region
		
		# 2. Finite Difference for Probabilities (Ps, Pd)
		grad = np.zeros(3)
		
		# d/dPs (forward difference)
		x_ps = x.copy(); x_ps[0] += self.epsilon
		val_ps = self.objective(x_ps)
		grad[0] = (val_ps - val_center) / self.epsilon
		
		# d/dPd (forward difference)
		x_pd = x.copy(); x_pd[1] += self.epsilon
		val_pd = self.objective(x_pd)
		grad[1] = (val_pd - val_center) / self.epsilon
		
		# 3. Centered difference for nu
		x_nu_p = x.copy(); x_nu_p[2] += self.epsilon
		x_nu_m = x.copy(); x_nu_m[2] -= self.epsilon
		grad[2] = (self.objective(x_nu_p) - self.objective(x_nu_m)) / (2 * self.epsilon)
		
		return grad
	\end{lstlisting}
\end{LTR}

این متد گرادیان را به صورت یک بردار سه‌بعدی $[\partial(-R)/\partial P_s, \partial(-R)/\partial P_d, \partial(-R)/\partial \nu]$ محاسبه می‌کند. ابتدا مقدار تابع هدف در نقطۀ فعلی محاسبه می‌شود. اگر این مقدار در ناحیۀ غیرممکن باشد (بزرگتر از $10^8$)، گرادیان صفر برگردانده می‌شود تا \lr{solver} در همان نقطه بماند. سپس برای $P_s$ و $P_d$ از تفاضل پیشرو استفاده می‌شود: $\partial f / \partial x_i \approx [f(x + \epsilon e_i) - f(x)] / \epsilon$. این روش به یک ارزیابی تابع اضافی برای هر پارامتر نیاز دارد (در مجموع دو ارزیابی برای $P_s$ و $P_d$). برای $\nu$ از تفاضل مرکزی استفاده می‌شود: $\partial f / \partial \nu \approx [f(\nu + \epsilon) - f(\nu - \epsilon)] / (2\epsilon)$. این روش به دو ارزیابی اضافی نیاز دارد اما دقت مرتبۀ دوم (\lr{second-order accuracy}) دارد (خطای $O(\epsilon^2)$ در مقابل $O(\epsilon)$ برای تفاضل پیشرو). انتخاب ترکیبی تفاضل پیشرو برای $P_s$ و $P_d$ و تفاضل مرکزی برای $\nu$ به این دلیل است که $P_s$ و $P_d$ در محدودیت‌های سخت شرکت می‌کنند و تفاضل پیشرو در نزدیک مرز پایدارتر است، در حالی که $\nu$ یک متغیر آزاد است و تفاضل مرکزی دقت بالاتری ارائه می‌دهد.

\subsection{الگوی \lr{callback} و تزریق وابستگی}

یکی از ویژگی‌های کلیدی این ماژول، استفاده گسترده از الگوی \lr{callback} است. توابع بهینه‌سازی به جای آنکه خود نرخ کلید را محاسبه کنند، یک تابع \lr{callback} از کاربر دریافت می‌کنند که این محاسبه را انجام می‌دهد. به عنوان مثال، تابع \lr{\texttt{optimize\_simultaneous\_ma2005}} یک \lr{callback} با امضای \lr{\texttt{Callable[[float, float, float], float]}} دریافت می‌کند که به این معناست تابعی است که سه عدد اعشاری ($P_s$، $P_d$، $\nu$) دریافت می‌کند و یک عدد اعشاری (نرخ منفی کلید) برمی‌گرداند. این طراحی چندین مزیت دارد:

\begin{itemize}
	\item جداسازی لایه‌ها: لایۀ بهینه‌سازی از لایۀ اثبات امنیتی مستقل است.
	\item انعطاف‌پذیری: کاربر می‌تواند تابع هدف خود را به‌سادگی جایگزین کند (مثلاً به‌جای نرخ کلید، \lr{QBER} را کمینه کند).
	\item تست‌پذیری: می‌توان \lr{callback}‌های مصنوعی با رفتار شناخته‌شده برای اعتبارسنجی استفاده کرد.
	\item قابلیت توسعه: اضافه‌کردن چارچوب‌های اثبات جدید نیازی به تغییر کد بهینه‌سازی ندارد.
\end{itemize}

این الگو که به \lr{``dependency injection''} نیز شناخته می‌شود، یک اصل کلیدی مهندسی نرم‌افزار است که قابلیت نگهداری و توسعه کد را به طور قابل توجهی افزایش می‌دهد.

\subsection{Type Hint ها و ایمنی نوع}

ماژول به طور گسترده از \lr{Type Hint} های \lr{Python} استفاده می‌کند تا ایمنی نوع و قابلیت خواندن کد را افزایش دهد. به عنوان مثال:

\begin{LTR}
	\begin{lstlisting}[language=Python]
		def optimize_simultaneous_ma2005(
		mu_target: float,
		total_pulses_N: int,
		cost_function_callback: Callable[[float, float, float], float]
		) -> Tuple[float, float, float, float]:
	\end{lstlisting}
\end{LTR}

در این امضا، \lr{\texttt{mu\_target: float}} مشخص می‌کند که \lr{\texttt{mu\_target}} باید یک عدد اعشاری باشد، \lr{\texttt{total\_pulses\_N: int}} مشخص می‌کند که \lr{\texttt{total\_pulses\_N}} باید یک عدد صحیح باشد، و \lr{\texttt{cost\_function\_callback: Callable[[float, float, float], float]}} یک \lr{type hint} پیچیده‌تر است که مشخص می‌کند \lr{callback} یک تابع است که سه عدد اعشاری دریافت می‌کند و یک عدد اعشاری برمی‌گرداند. خروجی تابع نیز به صورت \lr{\texttt{Tuple[float, float, float, float]}} مشخص شده که نشان می‌دهد یک تاپل چهار عنصری از اعداد اعشاری برگردانده می‌شود ($P_s^*$، $P_d^*$، $P_v^*$، $\nu^*$). این \lr{Type Hint} ها به ویرایشگرهای کد اجازه می‌دهند خطاهای نوع را در زمان توسعه شناسایی کنند و به کاربر در فهم رابط تابع کمک می‌کنند.

\section{پیاده‌سازی ریاضی و الگوریتمی}

\subsection{محاسبۀ $\mu$ بهینه بر اساس معادلۀ 12 مقالۀ \lr{Ma 2005}}

تابع \lr{\texttt{calculate\_optimal\_mu\_ma2005}} محاسبۀ تحلیلی-عددی $\mu$ بهینه را پیاده‌سازی می‌کند. در ادامه، نسخۀ کامل این تابع را با توضیح خط به خط ارائه می‌دهیم:

\begin{LTR}
	\begin{lstlisting}[language=Python]
		def calculate_optimal_mu_ma2005(
		distance_km: float,
		fiber_loss_db_km: float,
		det_eff: float,
		dark_rate: float,
		error_correction_f: float = 1.22,
		e_detector: float = 0.033
		) -> float:
		"""
		Calculates optimal signal intensity (mu) maximizing the GLLP rate.
		Implements the numerical solution to Eq. (12) in Ma et al. 2005.
		"""
		# 1. Calculate RHS constant C based on detector error
		h2_e = binary_entropy(e_detector)
		numerator = error_correction_f * h2_e
		denominator = 1.0 - h2_e
		
		if denominator <= 0:
		logger.warning(
		"Denominator in optimal mu calc is <= 0 "
		"(Error rate too high). Defaulting to 0.48.")
		return 0.48
		
		C = numerator / denominator
	\end{lstlisting}
\end{LTR}

در خطوط ابتدایی، ابتدا آنتروپی باینری خطای آشکارساز $H_2(e_{\text{det}})$ با تابع \lr{\texttt{binary\_entropy}} از ماژول \lr{\texttt{qkd.utils.math}} محاسبه می‌شود. سپس صورت و مخرج ثابت $C = f_{\text{EC}} H_2(e_{\text{det}}) / (1 - H_2(e_{\text{det}}))$ محاسبه می‌شوند. اگر مخرج غیرمثبت باشد (یعنی $H_2(e_{\text{det}}) \ge 1$ که فقط در $e_{\text{det}} = 0.5$ رخ می‌دهد و نشان‌دهندۀ خطای بسیار بالای آشکارساز است)، یک هشدار در لاگ ثبت می‌شود و مقدار پیش‌فرض \lr{0.48} برگردانده می‌شود. این رفتار \lr{fail-safe} تضمین می‌کند که تابع حتی در شرایط غیرفیزیکی هم یک خروجی معتبر برمی‌گرداند. سپس ثابت $C$ محاسبه می‌شود که در معادلۀ $(1 - \mu) e^{-\mu} = C$ استفاده می‌شود.

سپس تابع هدف و جواب آن به صورت زیر تعریف می‌شود:

\begin{LTR}
	\begin{lstlisting}[language=Python]
		# 2. Define function to find root: (1-mu)e^-mu - C = 0
		def objective(mu):
		return (1.0 - mu) * math.exp(-mu) - C
		
		# Define dynamic fallback based on Paper Section 3.1
		fallback_mu = 0.54 if error_correction_f <= 1.05 else 0.48
		
		# 3. Solve for mu in (0, 1]
		try:
		val_low = objective(0.001)
		val_high = objective(0.999)
		
		if val_low * val_high < 0:
		mu_opt = brentq(objective, 0.001, 0.999)
		return float(mu_opt)
		else:
		logger.warning(
		f"Optimal mu equation has no solution in [0,1]. "
		f"Falling back to {fallback_mu}.")
		return fallback_mu
		except Exception as e:
		logger.warning(
		f"Numerical solver failed for optimal mu: {e}. "
		f"Defaulting to {fallback_mu}.")
		return fallback_mu
	\end{lstlisting}
\end{LTR}

در این بخش، ابتدا تابع هدف به صورت $f(\mu) = (1 - \mu) e^{-\mu} - C$ تعریف می‌شود. این تابع در $\mu = 0$ برابر $1 - C$ (مثبت برای $C < 1$) و در $\mu \to \infty$ به $-C$ (منفی) میل می‌کند. بنابراین یک ریشه در بازۀ $(0, \infty)$ وجود دارد. برای اطمینان از وجود ریشه در بازۀ $[0.001, 0.999]$، مقدار تابع در دو انتها بررسی می‌شود: اگر $f(0.001) \cdot f(0.999) < 0$ (یعنی علامت‌های متضاد)، طبق \lr{Intermediate Value Theorem} یک ریشه در این بازه وجود دارد و \lr{\texttt{brentq}} می‌تواند آن را پیدا کند. در غیر این صورت، یک هشدار ثبت می‌شود و مقدار پیش‌فرض \lr{\texttt{fallback\_mu}} (که بر اساس $f_{\text{EC}}$ انتخاب می‌شود) برگردانده می‌شود. این \lr{fallback} بر اساس بخش 3.1 مقالۀ \lr{Ma 2005} انتخاب شده است: برای $f_{\text{EC}} \le 1.05$ مقدار \lr{0.54} و برای $f_{\text{EC}} > 1.05$ مقدار \lr{0.48} پیشنهاد شده است. در صورت بروز هرگونه استثنا در \lr{\texttt{brentq}}، بلوک \lr{\texttt{try/except}} آن را می‌گیرد و مقدار پیش‌فرض برمی‌گرداند.

\subsection{بهینه‌سازی $\nu$ با روش \lr{Brent}}

تابع \lr{\texttt{optimize\_nu\_ma2005}} بهینه‌سازی $\nu$ را برای $N$ متناهی پیاده‌سازی می‌کند:

\begin{LTR}
	\begin{lstlisting}[language=Python]
		def optimize_nu_ma2005(
		mu: float,
		total_pulses_N: float,
		cost_function_callback: Callable[[float], float],
		distance_km: float = 0.0
		) -> float:
		"""
		Optimizes the weak decoy intensity (nu) for a fixed data size N.
		Based on Ma et al. 2005 Section 4.3.
		"""
		lower_bound = 0.001
		upper_bound = mu * 0.9
		
		# Robust 1D scalar minimization (Brent's method)
		res = minimize_scalar(
		cost_function_callback,
		bounds=(lower_bound, upper_bound),
		method='bounded',
		options={'xatol': 1e-4}
		)
		
		if res.success:
		return float(res.x)
		
		# Fallback if optimization fails
		return 0.1
	\end{lstlisting}
\end{LTR}

در این تابع، ابتدا کران‌های جستجو به صورت $[\nu_{\min}, \nu_{\max}] = [0.001, 0.9 \mu]$ تعریف می‌شوند. کران پایین \lr{0.001} برای جلوگیری از مقادیر خیلی کوچک (که از نظر آماری معنایی ندارند) و کران بالا $0.9 \mu$ برای حفظ قابلیت تفکیک آماری بین سیگنال و \lr{decoy} انتخاب شده‌اند. سپس \lr{\texttt{minimize\_scalar}} با روش \lr{\texttt{'bounded'}} فراخوانی می‌شود. این روش بر اساس الگوریتم \lr{Brent's bounded method} است که برای توابع یک‌بعدی هموار بهینه است و همگرایی تضمین‌شده‌ای در بازۀ مشخص‌شده دارد. گزینه \lr{\texttt{xatol=1e-4}} تلورانس توقف را روی $10^{-4}$ تنظیم می‌کند که برای کاربردهای عملی کافی است. اگر بهینه‌سازی موفق شود، مقدار بهینۀ $\nu^*$ برگردانده می‌شود. در غیر این صورت، مقدار پیش‌فرض \lr{0.1} (یک مقدار تجربی خوب برای \lr{decoy} ضعیف) برگردانده می‌شود.

\subsection{بهینه‌سازی احتمالات توزیع پالس}

تابع \lr{\texttt{optimize\_pulse\_distribution\_ma2005}} احتمالات توزیع پالس $P_s$ و $P_d$ را بهینه‌سازی می‌کند:

\begin{LTR}
	\begin{lstlisting}[language=Python]
		def optimize_pulse_distribution_ma2005(
		mu: float,
		nu: float,
		total_pulses_N: int,
		cost_function_callback: Callable[[float, float], float]
		) -> Tuple[float, float, float]:
		"""
		Optimizes the pulse probabilities (Signal, Decoy, Vacuum) for a fixed data size N.
		"""
		def objective(x):
		p_s, p_d = x
		if p_s <= 0 or p_d <= 0 or (p_s + p_d) >= 1.0:
		return 1e9  # Penalty
		return cost_function_callback(p_s, p_d)
		
		# Initial guess
		x0 = [0.66, 0.29]
		bounds = [(0.01, 0.98), (0.01, 0.98)]
		constraints = [{'type': 'ineq', 'fun': lambda x: 1.0 - (x[0] + x[1]) - 0.01}]
		
		res = minimize(objective, x0, bounds=bounds,
		constraints=constraints, method='SLSQP')
		
		if res.success:
		p_s, p_d = res.x
		p_v = 1.0 - p_s - p_d
		return p_s, p_d, p_v
		
		logger.warning("Pulse distribution optimization failed. Using defaults.")
		return 0.8, 0.1, 0.1
	\end{lstlisting}
\end{LTR}

در این تابع، ابتدا تابع هدف داخلی تعریف می‌شود که محدودیت‌های فیزیکی را بررسی می‌کند: $P_s > 0$، $P_d > 0$ و $P_s + P_d < 1$. اگر هر کدام نقض شود، جریمۀ $10^9$ برگردانده می‌شود. در غیر این صورت، \lr{callback} فراخوانی می‌شود. سپس نقطۀ شروع $x_0 = [0.66, 0.29]$ (یک توزیع متعادل با $P_v \approx 0.05$) و کران‌های $[0.01, 0.98]$ برای هر دو پارامتر تعریف می‌شوند. یک محدودیت نامساوی نیز اضافه می‌شود: $P_s + P_d \le 0.99$ که فضایی برای $P_v \ge 0.01$ باقی می‌گذارد. سپس \lr{\texttt{minimize}} با روش \lr{\texttt{SLSQP}} فراخوانی می‌شود. این روش برای مسائل بهینه‌سازی مقید مناسب است و از محدودیت‌های نامساوی پشتیبانی می‌کند. اگر بهینه‌سازی موفق شود، $P_s^*$، $P_d^*$ و $P_v^* = 1 - P_s^* - P_d^*$ برگردانده می‌شوند. در غیر این صورت، مقادیر پیش‌فرض $(0.8, 0.1, 0.1)$ برگردانده می‌شوند که یک توزیع متعادل برای پروتکل \lr{Vacuum + Weak} است.

\subsection{بهینه‌سازی توامان سه پارامتر}

تابع \lr{\texttt{optimize\_simultaneous\_ma2005}} بهینه‌سازی توامان $P_s$، $P_d$ و $\nu$ را با استفاده از کلاس \lr{\texttt{Ma2005GlobalOptimizer}} پیاده‌سازی می‌کند:

\begin{LTR}
	\begin{lstlisting}[language=Python]
		def optimize_simultaneous_ma2005(
		mu_target: float,
		total_pulses_N: int,
		cost_function_callback: Callable[[float, float, float], float]
		) -> Tuple[float, float, float, float]:
		"""
		Simultaneously optimizes pulse probabilities (P_s, P_d) and decoy intensity (nu).
		"""
		optimizer = Ma2005GlobalOptimizer(
		mu_target, total_pulses_N, cost_function_callback)
		
		# Initial guess: Start with standard values
		x0 = [0.66, 0.29, 0.127]
		
		# Bounds
		bounds = [
		(0.1, 0.95), (0.01, 0.8), (0.001, mu_target * 0.99)
		]
		
		# Relax vacuum constraint to allow P_v -> 0
		constraints = [{'type': 'ineq',
			'fun': lambda x: 1.0 - (x[0] + x[1]) - 1e-9}]
		
		# Use SLSQP with the Jacobian method implemented in the optimizer class
		res = minimize(
		optimizer.objective,
		x0,
		jac=optimizer.gradient,
		bounds=bounds,
		constraints=constraints,
		method='SLSQP',
		tol=1e-5
		)
		
		if res.success:
		p_s_opt, p_d_opt, nu_opt = res.x
		p_v_opt = 1.0 - p_s_opt - p_d_opt
		return p_s_opt, p_d_opt, p_v_opt, nu_opt
		else:
		# If optimization fails, try fallback or return initial guess
		return 0.66, 0.29, 0.05, 0.127
	\end{lstlisting}
\end{LTR}

در این تابع، ابتدا یک نمونۀ \lr{\texttt{Ma2005GlobalOptimizer}} با \lr{\texttt{mu\_target}}، \lr{\texttt{total\_pulses\_N}} و \lr{\texttt{cost\_function\_callback}} ساخته می‌شود. سپس نقطۀ شروع $x_0 = [0.66, 0.29, 0.127]$ (مقادیر استاندارد برای پروتکل \lr{Vacuum + Weak}) تعریف می‌شود. کران‌ها به صورت $[0.1, 0.95]$ برای $P_s$، $[0.01, 0.8]$ برای $P_d$ و $[0.001, 0.99 \mu]$ برای $\nu$ تعریف می‌شوند. یک محدودیت نامساوی نیز اضافه می‌شود: $P_s + P_d \le 1 - 10^{-9}$ که اجازه می‌دهد $P_v$ به صفر میل کند (برای حالتی که \lr{decoy} خلأ لازم نیست). سپس \lr{\texttt{minimize}} با روش \lr{\texttt{SLSQP}} فراخوانی می‌شود و هم تابع هدف (\lr{\texttt{optimizer.objective}}) و هم گرادیان (\lr{\texttt{optimizer.gradient}}) به آن پاس داده می‌شوند. تلورانس توقف $10^{-5}$ تنظیم می‌شود. اگر بهینه‌سازی موفق شود، $P_s^*$، $P_d^*$، $P_v^*$ و $\nu^*$ برگردانده می‌شوند. در غیر این صورت، مقادیر پیش‌فرض $(0.66, 0.29, 0.05, 0.127)$ برگردانده می‌شوند.

\subsection{بهینه‌سازی کلید متناهی \lr{Lim 2014}}

تابع \lr{\texttt{optimize\_lim2014\_finite\_key}} بهینه‌سازی پنج پارامتر چارچوب \lr{Lim 2014} را پیاده‌سازی می‌کند:

\begin{LTR}
	\begin{lstlisting}[language=Python]
		def optimize_lim2014_finite_key(
		params: QKDParams,
		dist_km: float,
		cost_callback: Callable[[float, float, float, float, float], float]
		) -> Tuple[float, float, float, float, float]:
		"""
		Numerically optimizes the secret key rate R = l/N over the free parameters
		{q_x, p_mu1, p_mu2, mu1, mu2} as specified in Lim et al. (2014) Section IV.
		"""
		x0 = [0.9, 0.6, 0.2, 0.5, 0.1]
		
		bounds = [
		(0.5, 0.99),  # q_x
		(0.1, 0.8),   # p_s
		(0.05, 0.5),  # p_d
		(0.2, 0.8),   # mu
		(0.01, 0.2)   # nu
		]
		
		def objective(x):
		q_x, p_s, p_d, mu, nu = x
		if p_s + p_d >= 0.99: return 1e9
		if nu >= mu * 0.9: return 1e9
		if mu <= 0 or nu <= 0: return 1e9
		return cost_callback(q_x, p_s, p_d, mu, nu)
		
		# Run optimization (Nelder-Mead is robust for non-smooth finite-key landscapes)
		res = minimize(
		objective,
		x0,
		bounds=bounds,
		method='Nelder-Mead',
		options={'maxiter': 200, 'xatol': 1e-3}
		)
		
		if res.success:
		return tuple(res.x)
		
		# Fallback to initial guess if optimization fails
		logger.warning("Lim2014 parameter optimization failed. Using defaults.")
		return tuple(x0)
	\end{lstlisting}
\end{LTR}

در این تابع، نقطۀ شروع $x_0 = [0.9, 0.6, 0.2, 0.5, 0.1]$ برای پنج پارامتر $q_x$، $P_s$، $P_d$، $\mu$ و $\nu$ تعریف می‌شود. کران‌ها به صورت $[0.5, 0.99]$ برای $q_x$ (احتمال پایۀ $X$، باید بیشتر از \lr{0.5} باشد زیرا پایۀ تولیدکنندۀ کلید است)، $[0.1, 0.8]$ برای $P_s$، $[0.05, 0.5]$ برای $P_d$، $[0.2, 0.8]$ برای $\mu$ و $[0.01, 0.2]$ برای $\nu$ تعریف می‌شوند. تابع هدف داخلی محدودیت‌های فیزیکی را بررسی می‌کند: $P_s + P_d < 0.99$، $\nu < 0.9 \mu$ و $\mu, \nu > 0$. اگر هر کدام نقض شود، جریمۀ $10^9$ برگردانده می‌شود. سپس \lr{\texttt{minimize}} با روش \lr{\texttt{Nelder-Mead}} فراخوانی می‌شود. این روش بدون گرادیان است و برای توابع ناهموار مقاوم‌تر از روش‌های مبتنی بر گرادیان است. حداکثر تعداد تکرارها 200 و تلورانس $10^{-3}$ تنظیم می‌شود. اگر بهینه‌سازی موفق شود، تاپل پنج عنصری $(q_x^*, P_s^*, P_d^*, \mu^*, \nu^*)$ برگردانده می‌شود. در غیر این صورت، مقادیر پیش‌فرض (همان $x_0$) برگردانده می‌شوند.

\subsection{بهینه‌سازی زاویۀ چرخش در \lr{Chapman 2018}}

تابع \lr{\texttt{optimize\_rotation\_angle\_chapman}} بهینه‌سازی زاویۀ چرخش $\theta$ برای \lr{Coherent Scheme} در \lr{Chapman 2018} را پیاده‌سازی می‌کند:

\begin{LTR}
	\begin{lstlisting}[language=Python]
		def optimize_rotation_angle_chapman(gamma: float, M: int = 1) -> float:
		"""
		Finds the optimal encoding rotation angle theta_gamma for the Coherent Scheme
		in Chapman et al. (2018).
		
		We maximize P_succ = 0.5 * (1 + D(rho_0, rho_1)) where rho_i are outputs of ADC.
		"""
		# Kraus ops for ADC(gamma)
		E0 = np.array([[1, 0], [0, math.sqrt(1-gamma)]])
		E1 = np.array([[0, math.sqrt(gamma)], [0, 0]])
		
		# Rotation Y
		Ry = lambda theta: np.array([
		[math.cos(theta/2), -math.sin(theta/2)],
		[math.sin(theta/2),  math.cos(theta/2)]])
		
		def get_rho_out(state_in):
		return E0 @ state_in @ E0.conj().T + E1 @ state_in @ E1.conj().T
		
		def objective(theta):
		# Rotate |0> and |1> input states by theta
		U = Ry(theta)
		state0 = U @ np.array([[1,0],[0,0]]) @ U.conj().T
		state1 = U @ np.array([[0,0],[0,1]]) @ U.conj().T
		
		rho0_out = get_rho_out(state0)
		rho1_out = get_rho_out(state1)
		
		# Maximize Trace Distance for M=1
		diff = rho0_out - rho1_out
		evals = np.linalg.eigvalsh(diff)
		dist = 0.5 * np.sum(np.abs(evals))
		return -dist
		
		res = minimize_scalar(objective, bounds=(0, math.pi), method='bounded')
		return float(res.x)
	\end{lstlisting}
\end{LTR}

در این تابع، ابتدا عملگرهای \lr{Kraus} برای \lr{ADC} با پارامتر $\gamma$ ساخته می‌شوند: $E_0$ و $E_1$ به صورت ماتریس‌های \lr{NumPy} تعریف می‌شوند. سپس یک \lr{lambda} برای ماتریس چرخش \lr{Y} به صورت $R_y(\theta) = \begin{pmatrix} \cos(\theta/2) & -\sin(\theta/2) \\ \sin(\theta/2) & \cos(\theta/2) \end{pmatrix}$ تعریف می‌شود. تابع \lr{\texttt{get\_rho\_out}} حالت خروجی کانال را از رابطۀ $\rho_{\text{out}} = E_0 \rho_{\text{in}} E_0^\dagger + E_1 \rho_{\text{in}} E_1^\dagger$ محاسبه می‌کند. در تابع هدف، ابتدا ماتریس چرخش $U = R_y(\theta)$ ساخته می‌شود و حالت‌های ورودی $|0\rangle\langle 0|$ و $|1\rangle\langle 1|$ با این چرخش تبدیل می‌شوند: $\rho_i^{\text{in}} = U |i\rangle\langle i| U^\dagger$. سپس حالت‌های خروجی محاسبه می‌شوند و فاصلۀ \lr{trace} میان آن‌ها به صورت $D = \frac{1}{2} \sum_i |\lambda_i|$ محاسبه می‌شود که در آن $\lambda_i$ مقادیر ویژۀ $\rho_0^{\text{out}} - \rho_1^{\text{out}}$ هستند. تابع \lr{\texttt{eigvalsh}} از \lr{NumPy} برای محاسبۀ مقادیر ویژۀ یک ماتریس \lr{Hermitian} استفاده می‌شود (که سریع‌تر و پایدارتر از \lr{\texttt{eigvals}} عمومی است). چون بهینه‌سازی به دنبال بیشینه‌کردن $D$ است و \lr{\texttt{minimize\_scalar}} به دنبال کمینه‌کردن است، علامت منفی $-D$ برگردانده می‌شود. در نهایت، \lr{\texttt{minimize\_scalar}} با روش \lr{\texttt{'bounded'}} در بازۀ $[0, \pi]$ فراخوانی می‌شود و زاویۀ بهینۀ $\theta^*$ برگردانده می‌شود.

\subsection{بهینه‌سازی اکتشافی شدت‌ها}

تابع \lr{\texttt{optimize\_decoy\_intensities}} یک بهینه‌سازی اکتشافی برای شدت سیگنال و \lr{decoy} ارائه می‌کند:

\begin{LTR}
	\begin{lstlisting}[language=Python]
		def optimize_decoy_intensities(
		distance_km: float,
		fiber_loss_db_km: float,
		det_eff: float,
		dark_rate: float,
		error_correction_f: float
		) -> Tuple[float, float]:
		"""
		Heuristic optimization for signal and decoy intensities.
		Returns: (mu_opt, nu_opt)
		"""
		transmittance = 10**(-fiber_loss_db_km * distance_km / 10.0)
		
		def neg_key_rate(x):
		mu, nu = x[0], x[1]
		if mu <= nu or nu <= 0 or mu > 1.0:
		return 1.0
		
		# Heuristic simplified rate model
		Q_mu = 1 - math.exp(-transmittance * mu) + 2 * dark_rate
		E_mu = (0.5 * 2 * dark_rate) / Q_mu
		Y1 = transmittance * det_eff
		Q1 = mu * math.exp(-mu) * Y1
		
		rate = Q1 - error_correction_f * Q_mu * binary_entropy(E_mu)
		return -rate
		
		res = minimize(neg_key_rate, [0.5, 0.1],
		bounds=[(0.1, 0.9), (0.01, 0.3)])
		
		if res.success:
		return res.x[0], res.x[1]
		return 0.48, 0.05
	\end{lstlisting}
\end{LTR}

در این تابع، ابتدا عبور کانال به صورت $\eta = 10^{-\alpha L / 10}$ محاسبه می‌شود که در آن $\alpha$ ضریب افت فیبر (به \lr{dB/km}) و $L$ فاصله (به \lr{km}) است. این فرمول از قانون \lr{Beer-Lambert} ناشی می‌شود. سپس تابع هدف به صورت نرخ منفی کلید با یک مدل ساده‌شده تعریف می‌شود: $Q_\mu = 1 - e^{-\eta \mu} + 2 Y_0$ (احتمال آشکارسازی سیگنال به همراه \lr{dark count})، $E_\mu = Y_0 / Q_\mu$ (نرخ خطای سیگنال که عمدتاً از \lr{dark count} ناشی می‌شود)، $Y_1 = \eta \cdot \text{det\_eff}$ (بازده تک‌فوتونی) و $Q_1 = \mu e^{-\mu} Y_1$ (\lr{Gain} تک‌فوتونی). نرخ کلید به صورت $R = Q_1 - f_{\text{EC}} Q_\mu H_2(E_\mu)$ محاسبه می‌شود. این یک مدل اکتشافی است که فرض می‌کند $e_1 \approx 0$ و $q = 1$ (یعنی پایه‌گذاری بی‌اتلاف). سپس \lr{\texttt{minimize}} با نقطۀ شروع $[0.5, 0.1]$ و کران‌های $[0.1, 0.9]$ برای $\mu$ و $[0.01, 0.3]$ برای $\nu$ فراخوانی می‌شود. اگر بهینه‌سازی موفق شود، $(\mu^*, \nu^*)$ برگردانده می‌شود. در غیر این صورت، مقادیر پیش‌فرض $(0.48, 0.05)$ برگردانده می‌شوند.

\section{ارسال و دریافت با سایر ماژول‌ها}

\subsection{ورودی‌های ماژول}

ماژول \lr{\texttt{optimization.py}} برای انجام بهینه‌سازی‌های خود به چندین نوع ورودی نیاز دارد. این ورودی‌ها به دو دستۀ پارامترهای فیزیکی و \lr{callback}‌های تابعی تقسیم می‌شوند.

\begin{itemize}
	\item پارامترهای فیزیکی: این پارامترها مشخصات سیستم \lr{QKD} را توصیف می‌کنند و شامل موارد زیر هستند: \lr{\texttt{distance\_km}} (فاصله به کیلومتر)، \lr{\texttt{fiber\_loss\_db\_km}} (ضریب افت فیبر به \lr{dB/km}، معمولاً 0.2 برای فیبر سیلیکایی در طول موج 1550 \lr{nm})، \lr{\texttt{det\_eff}} (راندمان آشکارساز، معمولاً 0.1 تا 0.3 برای \lr{SPAD})، \lr{\texttt{dark\_rate}} (نرخ شمارش تاریک آشکارساز)، \lr{\texttt{error\_correction\_f}} (بازدهی تصحیح خطا، معمولاً 1.1 تا 1.2) و \lr{\texttt{e\_detector}} (خطای ذاتی آشکارساز). این پارامترها یا به صورت آرگومان‌های مستقیم به توابع پاس داده می‌شوند یا از شیء \lr{\texttt{QKDParams}} استخراج می‌شوند.
	\item فراخوان‌های تابعی (\lr{callback}ها): این \lr{callback}‌ها تابع هدف را برای بهینه‌سازی تعریف می‌کنند و توسط کاربر ارائه می‌شوند. امضای این \lr{callback}‌ها بسته به تابع بهینه‌سازی متفاوت است: \lr{\texttt{Callable[[float], float]}} برای \lr{\texttt{optimize\_nu\_ma2005}}، \lr{\texttt{Callable[[float, float], float]}} برای \lr{\texttt{optimize\_pulse\_distribution\_ma2005}}، \lr{\texttt{Callable[[float, float, float], float]}} برای \lr{\texttt{optimize\_simultaneous\_ma2005}} و \lr{\texttt{Callable[[float, float, float, float, float], float]}} برای \lr{\texttt{optimize\_lim2014\_finite\_key}}. این \lr{callback}‌ها معمولاً نرخ منفی کلید را محاسبه می‌کنند (زیرا حل‌گرها (\lr{solver}ها) به دنبال کمینه‌کردن هستند) اما می‌توانند هر کمیت هدف دیگری را نیز محاسبه کنند.
	\item \lr{\texttt{QKDParams}}: این شیء از ماژول \lr{\texttt{qkd.params}} وارد می‌شود و در تابع \lr{\texttt{optimize\_lim2014\_finite\_key}} استفاده می‌شود. این شیء شامل تمام پارامترهای پیکربندی سیستم است و به \lr{callback} اجازه می‌دهد به پارامترهای اثبات امنیتی (مانند $\varepsilon_{\text{sec}}$، $\varepsilon_{\text{cor}}$ و $\varepsilon_{\text{pe}}$) دسترسی داشته باشد.
	\item \lr{\texttt{gamma}} (پارامتر کانال): این پارامتر در توابع \lr{\texttt{optimize\_rotation\_angle}} و \lr{\texttt{optimize\_rotation\_angle\_chapman}} استفاده می‌شود و نشان‌دهندۀ پارامتر کانال \lr{ADC} (در \lr{Chapman 2018}) یا یک کانال عمومی است.
\end{itemize}

\subsection{خروجی‌های ماژول}

خروجی‌های ماژول بسته به تابع بهینه‌سازی متفاوت هستند:

\begin{enumerate}
	\item خروجی‌های تک‌مقداری: توابع \lr{\texttt{calculate\_optimal\_mu\_ma2005}}، \lr{\texttt{optimize\_nu\_ma2005}}، \lr{\texttt{optimize\_rotation\_angle}} و \lr{\texttt{optimize\_rotation\_angle\_chapman}} یک مقدار اعشاری واحد را برمی‌گردانند (به ترتیب $\mu^*$، $\nu^*$، $\theta^*$ عمومی و $\theta^*$ برای \lr{Chapman}).
	\item خروجی‌های دومقداری: تابع \lr{\texttt{optimize\_decoy\_intensities}} یک تاپل دو عنصری شامل $(\mu^*, \nu^*)$ را برمی‌گرداند.
	\item خروجی‌های سه‌مقداری: تابع \lr{\texttt{optimize\_pulse\_distribution\_ma2005}} (و پوشش‌دهندۀ آن \lr{\texttt{optimize\_pulse\_probabilities}}) یک تاپل سه عنصری شامل $(P_s^*, P_d^*, P_v^*)$ را برمی‌گرداند.
	\item خروجی‌های چهارمقداری: تابع \lr{\texttt{optimize\_simultaneous\_ma2005}} (و پوشش‌دهندۀ آن \lr{\texttt{optimize\_global\_ma2005}}) یک تاپل چهار عنصری شامل $(P_s^*, P_d^*, P_v^*, \nu^*)$ را برمی‌گرداند.
	\item خروجی‌های پنج‌مقداری: تابع \lr{\texttt{optimize\_lim2014\_finite\_key}} یک تاپل پنج عنصری شامل $(q_x^*, P_s^*, P_d^*, \mu^*, \nu^*)$ را برمی‌گرداند.
\end{enumerate}

در تمام موارد، اگر بهینه‌سازی شکست بخورد، مقادیر پیش‌فرض محافظه‌کارانه برگردانده می‌شوند. این مقادیر پیش‌فرض بر اساس تجربۀ عملی و مقالات مرجع انتخاب شده‌اند و برای پروتکل‌های \lr{Decoy-State BB84} معمولاً خوب عمل می‌کنند.

\subsection{وابستگی‌های داخلی و خارجی}

ماژول وابستگی‌های داخلی و خارجی متعددی دارد. وابستگی‌های داخلی شامل موارد زیر است:

\begin{itemize}
	\item \lr{\texttt{qkd.params.QKDParams}}: کلاس پارامترهای \lr{QKD} که در تابع \lr{\texttt{optimize\_lim2014\_finite\_key}} استفاده می‌شود.
	\item \lr{\texttt{qkd.utils.math.binary\_entropy}}: تابع محاسبۀ آنتروپی باینری $H_2(p) = -p \log_2 p - (1-p) \log_2 (1-p)$ که در توابع \lr{\texttt{calculate\_optimal\_mu\_ma2005}} و \lr{\texttt{optimize\_decoy\_intensities}} استفاده می‌شود.
	\item \lr{\texttt{qkd.constants.is\_valid\_probability}}: تابع اعتبارسنجی احتمال که بررسی می‌کند یک مقدار در بازۀ $[0, 1]$ قرار دارد.
	\item \lr{\texttt{qkd.datatypes.TallyCounts}}: کلاس دادهٔ (\lr{dataclass}) شمارش‌های آماری که در برخی \lr{callback}‌ها برای محاسبۀ نرخ کلید استفاده می‌شود.
\end{itemize}

وابستگی‌های خارجی شامل موارد زیر است:

\begin{itemize}
	\item \lr{\texttt{scipy.optimize.minimize}}: تابع بهینه‌سازی چندبعدی \lr{SciPy} که در توابع \lr{\texttt{optimize\_decoy\_intensities}}، \lr{\texttt{optimize\_pulse\_distribution\_ma2005}}، \lr{\texttt{optimize\_simultaneous\_ma2005}} و \lr{\texttt{optimize\_lim2014\_finite\_key}} استفاده می‌شود. این تابع از روش‌های مختلفی (مانند \lr{SLSQP} و \lr{Nelder-Mead}) پشتیبانی می‌کند.
	\item \lr{\texttt{scipy.optimize.minimize\_scalar}}: تابع بهینه‌سازی یک‌بعدی \lr{SciPy} که در توابع \lr{\texttt{optimize\_rotation\_angle}}، \lr{\texttt{optimize\_nu\_ma2005}} و \lr{\texttt{optimize\_rotation\_angle\_chapman}} استفاده می‌شود. این تابع از روش \lr{Brent's bounded} پشتیبانی می‌کند.
	\item \lr{\texttt{scipy.optimize.brentq}}: تابع یافتن ریشه \lr{SciPy} که در تابع \lr{\texttt{calculate\_optimal\_mu\_ma2005}} برای حل معادلۀ $(1-\mu)e^{-\mu} = C$ استفاده می‌شود.
	\item \lr{\texttt{numpy}}: کتابخانۀ محاسبات برداری که در تابع \lr{\texttt{optimize\_rotation\_angle\_chapman}} برای محاسبۀ ماتریس‌ها و مقادیر ویژه استفاده می‌شود.
	\item \lr{\texttt{math}}: ماژول ریاضی استاندارد \lr{Python} برای توابع پایه مانند \lr{\texttt{exp}}، \lr{\texttt{log}}، \lr{\texttt{cos}}، \lr{\texttt{sin}} و \lr{\texttt{sqrt}}.
	\item \lr{\texttt{logging}}: ماژول لاگ‌گذاری استاندارد \lr{Python} برای ثبت هشدارها و اطلاعات.
	\item \lr{\texttt{typing}}: ماژول \lr{Type Hint} برای تعریف امضای توابع.
\end{itemize}

\subsection{جریان داده در سراسر فریم‌ورک}

جریان داده کلی به این شکل است:
\begin{enumerate}
	\item لایۀ بالایی (مانند اسکریپت پیکربندی یا شبیه‌ساز) پارامترهای فیزیکی سیستم (مانند فاصله، افت فیبر، راندمان آشکارساز) را مشخص می‌کند.
	\item یک \lr{callback} تابعی تعریف می‌شود که نرخ کلید (یا هر کمیت هدف دیگر) را برای مجموعه‌ای از پارامترهای داده‌شده محاسبه می‌کند. این \lr{callback} معمولاً از ماژول‌های اثبات امنیتی (مانند \lr{\texttt{lim2014.py}}، \lr{\texttt{tight.py}} یا \lr{\texttt{ma2005.py}}) استفاده می‌کند.
	\item تابع بهینه‌سازی مناسب (مانند \lr{\texttt{optimize\_lim2014\_finite\_key}}) با پارامترهای فیزیکی و \lr{callback} فراخوانی می‌شود.
	\item تابع بهینه‌سازی، \lr{callback} را در نقاط مختلف فضای پارامتر فراخوانی می‌کند و با استفاده از روش‌های عددی (مانند \lr{SLSQP} یا \lr{Nelder-Mead}) به سمت نقطۀ بهینه حرکت می‌کند.
	\item در هر فراخوانی \lr{callback}، تابع هدف ابتدا محدودیت‌های فیزیکی را بررسی می‌کند و در صورت نقض، جریمۀ بزرگی برگردانده می‌شود.
	\item پس از همگرایی (یا رسیدن به حداکثر تعداد تکرار)، تابع بهینه‌سازی پارامترهای بهینه را برمی‌گرداند.
	\item لایۀ بالایی از این پارامترهای بهینه برای پیکربندی شبیه‌سازی نهایی استفاده می‌کند.
\end{enumerate}

این جداسازی مراحل به گونه‌ای طراحی شده که هر مرحله به صورت مستقل قابل آزمایش و بازبینی باشد. به عنوان مثال، می‌توان از \lr{callback}‌های مصنوعی با رفتار مشخص برای اعتبارسنجی توابع بهینه‌سازی استفاده کرد. این ماژول‌بندی از اصل تک‌مسئولیتی (\lr{Single Responsibility}) پیروی می‌کند و قابلیت نگهداری و توسعه کد را به طور قابل توجهی افزایش می‌دهد. علاوه بر این، الگوی \lr{callback} به کاربر اجازه می‌دهد تابع هدف خود را بدون تغییر کد بهینه‌سازی جایگزین کند. این انعطاف‌پذیری برای تحقیقات \lr{QKD} که در آن‌ها معیارهای مختلف بهینه‌سازی (نرخ کلید، \lr{QBER}، طول کلید متناهی و غیره) مورد بررسی قرار می‌گیرند، حیاتی است. طراحی مطمئن از خرابی (\lr{fail-safe}) این ماژول، آن را به انتخابی ایده‌آل برای محیط‌های عملیاتی تبدیل می‌کند که در آن‌ها پایداری و قابلیت اطمینان سیستم اهمیت زیادی دارد.

	
	% =========================== بخش سوم: اجرا و شبکه ===========================
	\part{خط لوله اجرا و لایه شبکه}
	
	\chapter{اسکریپت اصلی اجرا}
	\label{ch:main}
	\sloppy
\emergencystretch=5em
\hbadness=10000
\vbadness=10000
\tolerance=2000

\providecolor{codebg}{RGB}{248,248,248}
\providecolor{codeframe}{RGB}{200,200,200}
\providecolor{keyword}{RGB}{0,0,180}
\providecolor{string}{RGB}{163,21,21}
\providecolor{comment}{RGB}{0,128,0}

\lstset{
	backgroundcolor=\color{codebg},
	frame=single,
	rulecolor=\color{codeframe},
	language=Python,
	basicstyle=\ttfamily\scriptsize,
	keywordstyle=\color{keyword}\bfseries,
	stringstyle=\color{string},
	commentstyle=\color{comment}\itshape,
	showstringspaces=false,
	breaklines=true,
	breakatwhitespace=true,
	numbers=left,
	numberstyle=\tiny\color{gray},
	numbersep=8pt,
	tabsize=4,
	literate=
	{ف}{{{\bfseries\itshape ف}}}1
	{ی}{{{\bfseries\itshape ی}}}1
}

\section{نمای کلی و هدف ماژول}

ماژول \lr{\texttt{main\_optimized.py}} هسته‌ی مرکزی فریم‌ورک شبیه‌سازی \lr{QKD} است و نقش \emph{هماهنگ‌کننده‌ی سرتاسری} (\lr{end-to-end orchestrator}) را در معماری کلی ایفا می‌کند. این ماژول در نسخه‌ی فعلی (نسخه‌ی بازنویسی‌شده با نشانگرهای \lr{F-01} تا \lr{F-29}) از یک اسکریپت ساده‌ی پیمایش \lr{Monte-Carlo} به یک \lr{pipeline} بالغ با قابلیت‌های زیر ارتقا یافته است: پشتیبانی کامل از حالت \lr{DWDM} (\lr{Dense Wavelength Division Multiplexing}) با تزریق نویز \lr{Raman}، بهینه‌سازی خودکار پارامترهای پالس به ازای هر فاصله با \lr{SLSQP}، الگوی استاندارد \lr{SimParams} برای جداسازی داده‌های اعتبارسنجی‌شده از پیکربندی خام، و یک لایه‌ی خطا‌‌نگاری یکپارچه بر پایه‌ی سلسله‌مراتب استثناهای \lr{QKDException}.

هدف اصلی این ماژول، پاسخ به پرسش بنیادین زیر است: \emph{چگونه می‌توان طیف وسیعی از پارامترهای فیزیکی (فاصله، شدت پالس، بازده آشکارساز، نویز دارک کانت، نویز \lr{Raman} در \lr{DWDM}، و...) را به‌صورت سیستماتیک در یک شبیه‌سازی موازی پویش (\lr{sweep}) نمود و نتایج را به‌صورت یک فایل \lr{CSV} ساختارمند با ستون‌های ثابت و قابل ممیزی ثبت کرد؟} این پرسش به سه دلیل اهمیت دارد. نخست، در شبیه‌سازی‌های علمی \lr{QKD}، نیاز به پیمایش صدها یا هزاران ترکیب پارامتر برای رسم منحنی‌های نرخ کلید-فاصله یا مقایسه‌ی پروتکل‌ها، امری روزمره است. دوم، موازی‌سازی با \lr{ProcessPoolExecutor} نیازمند مدیریت دقیق \lr{state} هر کارگر، بذرگیری مستقل آماری، و الگوی \lr{lazy rebuild} برای جلوگیری از بازسازی اشیای سنگین فیزیکی در هر تکرار است. سوم، خروجی \lr{CSV} باید ساختار یکسانی برای همه‌ی ردیف‌ها (چه موفق و چه خطا) داشته باشد تا ابزارهای تحلیل بعدی (\lr{pandas}، \lr{R}) بتوانند آن را به‌درستی پردازش کنند.

ساختار کلی این ماژول در نسخه‌ی جدید شامل چندین لایه‌ی منطقی است. لایه‌ی پیکربندی شامل \lr{\texttt{DEFAULT\_CONFIG}} (دیکشنری تودرتوی پیش‌فرض) و \lr{\texttt{DEFAULT\_SWEEPS}} (تعریف پیمایش‌های رایج، اکثراً \lr{comment} شده) است. لایه‌ی \lr{adapter} توابع \lr{\texttt{build\_*}} را در بر می‌گیرد که از پیکربندی خام، اشیای \lr{dataclass} اعتبارسنجی‌شده می‌سازند. لایه‌ی \lr{SimParams} شامل \lr{\texttt{build\_channel\_sim\_params}}، \lr{\texttt{build\_source\_sim\_params}}، \lr{\texttt{build\_detector\_sim\_params}} است که داده‌های اعتبارسنجی‌شده‌ی \lr{frozen} را برای حلقه‌ی شبیه‌سازی فراهم می‌سازند. لایه‌ی \lr{worker} شامل کلاس \lr{\texttt{\_WorkerState}} و توابع \lr{\texttt{init\_worker}} و \lr{\texttt{run\_single\_simulation}} است. لایه‌ی \lr{orchestrator} تابع \lr{\texttt{run\_and\_save\_csv}} را در بر می‌گیرد که کارهای شبیه‌سازی را تولید کرده و آن‌ها را به \lr{\texttt{dispatch\_work\_items}} (در ماژول جداگانه‌ی \lr{\texttt{main\_optimized\_dispatch}}) ارسال می‌نماید.

ویژگی‌های برجسته‌ی نسخه‌ی جدید عبارت‌اند از: (۱) الگوی \lr{SimParams} که خواندن مستقیم از \lr{config dict} در حلقه‌ی شبیه‌سازی را حذف کرده و داده‌های اعتبارسنجی‌شده را از اشیای \lr{frozen} می‌خواند؛ (۲) پشتیبانی کامل از \lr{DWDM} با تزریق نویز \lr{Raman} به دو مسیر مستقل (مسیر \lr{Monte-Carlo} آشکارساز و مسیر اثبات امنیتی \lr{Lim2014})؛ (۳) بهینه‌سازی خودکار پارامترهای پالس با کش per-(\lr{combo}, \lr{distance}) برای اشتراک پارامترهای بهینه بین ردیف‌های \lr{noise=True} و \lr{noise=False}؛ (۴) الگوی استاندارد خطا‌‌نگاری با \lr{\texttt{\_as\_qkd\_exception}} و \lr{\texttt{\_log\_qkd\_exception}} که همه‌ی \lr{catch-block}ها را به یک الگوی واحد تبدیل می‌کند؛ (۵) الگوی \lr{from\_config\_dict} برای آشکارساز که تمام پارامترها را در یک گام اعتبارسنجی می‌کند و نیاز به \lr{setattr} پس از ساخت را حذف می‌نماید؛ (۶) افزودن ستون \lr{analytic\_rogers\_sbr} به \lr{CSV} برای مقایسه‌ی مستقیم \lr{Monte-Carlo} با فرمول بسته‌ی \lr{Rogers et al. (2007)}؛ (۷) ستون‌های \lr{params\_*} برای ممیزی \lr{DWDM}؛ (۸) ستون‌های \lr{det\_override\_*} برای ممیزی \lr{override}های آشکارساز.

در سطح عملیاتی، این ماژول به‌عنوان \lr{entry point} اصلی فریم‌ورک عمل می‌کند و از طریق \lr{argparse} یک \lr{CLI} غنی با بیش از ۳۰ گزینه ارائه می‌نماید. کاربر می‌تواند با خط فرمان، پروتکل (\lr{BB84Decoy}، \lr{B92}، \lr{MDI}، \lr{RedundantTransmission})، اثبات امنیتی (\lr{Lim2014}، \lr{Ma2005}، \lr{Wang2005}، \lr{Tight})، پارامترهای فیزیکی (فاصله، بازده، نویز)، پارامترهای امنیتی (\lr{epsilons})، حالت \lr{DWDM} (با تعداد کانال، توان لانچ، ضریب \lr{Raman})، و بهینه‌سازی خودکار پالس را پیکربندی نماید. خروجی در فایل \lr{\texttt{qkd\_results\_optimized\_sweep.csv}} ذخیره می‌شود و شامل ستون‌های ثابت برای مسیر موفق و مسیر خطا است (الگوی \lr{F-07}).

\section{مبانی تئوری و فیزیکی}

\subsection{ساختار پیکربندی سلسله‌مراتبی و اصل Single Source of Truth}

پیکربندی فریم‌ورک در یک دیکشنری تودرتوی چهار سطحی سازماندهی شده است: \lr{\texttt{protocol\_params}} (نوع پروتکل و آشکارساز)، \lr{\texttt{protocol\_runtime}} (کلاس پروتکل، احتمال پایه، خط‌مشی کلیک دوگانه، پرچم‌های انکودر/دیکودر درهم‌تنیده)، \lr{\texttt{channel}} (کاهیدگی فیبر، پاشش)، \lr{\texttt{detector}} (بازده، نرخ دارک کانت، \lr{QBER} ذاتی، نامساوی، \lr{dead time}، \lr{jitter}، \lr{afterpulse}، دما، ولتاژ بایاس)، \lr{\texttt{source}} (کلاس منبع، توزیع آماری، مدل خطا، شدت پالس‌ها، \lr{MZM}، نویز الکتریکی، \lr{security\_metadata})، \lr{\texttt{protocol}} (تخصیص \lr{epsilon})، و \lr{\texttt{simulation}} (بذر، \lr{log-range} پالس، بازه فاصله). این سلسله‌مراتب امکان پیکربندی دقیق هر لایه‌ی فیزیکی را فراهم می‌سازد.

با این حال، در نسخه‌ی جدید، اصل \emph{Single Source of Truth} به‌طور جدی اعمال شده است. به‌جای خواندن مستقیم از \lr{config dict} در حلقه‌ی شبیه‌سازی (که ممکن است ناسازگار با اشیای اعتبارسنجی‌شده باشد)، توابع \lr{\texttt{build\_*}} اشیای \lr{dataclass} می‌سازند و سپس توابع \lr{\texttt{build\_*\_sim\_params}} اشیای \lr{SimParams} \lr{frozen} را از آن‌ها استخراج می‌کنند. حلقه‌ی شبیه‌سازی فقط از \lr{SimParams} می‌خواند. این الگو تضمین می‌کند که اگر کاربر پیکربندی را با \lr{sweep} تغییر دهد، تمام مقادیر اعتبارسنجی شده و سپس به حلقه ارسال شوند.

\subsection{تخصیص بودجه امنیتی $\varepsilon$ و Composable Security}

در رمزنگاری کوانتومی مدرن، امنیت به‌صورت \lr{composable} تعریف می‌شود؛ یعنی مجموع احتمال شکست امنیتی باید کم‌تر از یک بودجه‌ی کل $\varepsilon_{\text{total}}$ باشد. این بودجه بین چندین مرحله تقسیم می‌شود:

\begin{equation}
	\varepsilon_{\text{total}} = \varepsilon_{\text{sec}} + \varepsilon_{\text{cor}} + \varepsilon_{\text{pe}} + \varepsilon_{\text{smooth}} + \varepsilon_{\text{pa}} + \varepsilon_{\text{phase\_est}}
\end{equation}

در این رابطه، $\varepsilon_{\text{sec}}$ امنیت کلی، $\varepsilon_{\text{cor}}$ تصحیح خطا، $\varepsilon_{\text{pe}}$ تخمین پارامتر، $\varepsilon_{\text{smooth}}$ هموارسازی، $\varepsilon_{\text{pa}}$ تقویت حریم خصوصی، و $\varepsilon_{\text{phase\_est}}$ تخمین خطای فاز است. تابع \lr{\texttt{build\_epsilon\_allocation}} این مقادیر را از پیکربندی می‌خواند و اعتبارسنجی می‌نماید که همگی مثبت باشند و مجموعشان کم‌تر از $1$ باشد. در صورت نقض، \lr{\texttt{ParameterValidationError}} پرتاب می‌شود.

\subsection{فیزیک دارک کانت و فرض Appendix B مقالۀ Lim 2014}

در مقاله‌ی \lr{Lim et al. 2014} (\lr{PhysRevA.89.022307})، \lr{Appendix B} به‌صراحت فرض می‌کند که \emph{«مشارکت‌های واکوئم هیچ اطلاعاتی درباره‌ی مقادیر بیت انتخابی ندارند و بیت‌ها به‌طور یکنواخت توزیع شده‌اند.»} این فرض، توجیه می‌کند که رویدادهای واکوئم در فرمول طول کلید (معادله‌ی ۱) اعتبار آنتروپی کامل دریافت می‌کنند بدون آنکه نیاز به کسر $h(\cdot)$ باشد.

از نظر فیزیکی، آشکارسازی واکوئم در واقع دارک کانت است: آشکارساز به‌صورت خود به خودی کلیک می‌کند و مقدار بیت که \lr{Bob} تخصیص می‌دهد، به‌طور یکنواخت تصادفی در $\{0, 1\}$ است. بنابراین، \lr{QBER} \emph{اندازه‌گیری‌شده} روی پالس‌های واکوئم باید حدود $50\%$ باشد، نه \lr{QBER} ذاتی آشکارساز ($\sim 1\%$). در کد، دو بلاک توضیحی (در خطوط ۱۲۵۳–۱۲۹۶ و ۱۲۹۸–۱۳۲۶) این موضوع را به‌تفصیل شرح می‌دهند. باگ در \lr{qkd.detectors} باعث می‌شود \lr{QBER} واکوئم $\sim 1\%$ باشد که فرض \lr{Appendix B} را نقض می‌کند. در نسخه‌ی فعلی، اصلاح در سطح مرز \lr{stats-map} به‌صورت یک \lr{workaround} اعمال شده و توابع \lr{\texttt{correct\_vacuum\_qber\_for\_proof}} و \lr{\texttt{correct\_signal\_decoy\_qber\_for\_proof}} (که در نسخه‌ی قبلی وجود داشتند) به‌صورت comment یا placeholder باقی مانده‌اند. این موضوع در بخش \ref{sec:csv_row} به‌تفصیل بررسی می‌شود.

\subsection{مدل فیزیکی QBER سیگنال و دکوی}

\lr{QBER} (\lr{Quantum Bit Error Rate}) سیگنال و دکوی از ترکیب خطای ذاتی سیستم و نویز دارک کانت به دست می‌آید. در حالت ایده‌آل، \lr{QBER} متوسط برابر است با:

\begin{equation}
	E_\mu = \frac{e_d \cdot P_{\text{real}} + 0.5 \cdot P_{\text{dark}}}{P_{\text{real}} + P_{\text{dark}}}
\end{equation}

در این رابطه، $e_d$ خطای ذاتی (مجموع \lr{qber\_intrinsic} و \lr{misalignment})، $P_{\text{real}} = 1 - \exp(-\mu \cdot \eta_{\text{det}} \cdot \eta_{\text{channel}})$ احتمال کلیک فوتون واقعی (مدل \lr{Poisson} + بازده)، و $P_{\text{dark}} = 1 - (1 - Y_0 \cdot \Delta t)^2$ احتمال کلیک دارک کانت (دو آشکارساز) است. در کد، بلاک توضیحی در خطوط ۱۲۹۸–۱۳۲۶ این مدل را شرح می‌دهد و نشان می‌دهد که باگ \lr{qkd.detectors} باعث می‌شود \lr{QBER} سیگنال در فاصله‌ی صفر به‌جای $\sim 0.015$ به $\sim 0.235$ برسد.

\subsection{تئوری Rogers et al. (2007) برای نرخ بیت sifted}

\lr{Rogers, Bienfang, Nakassis, Xu \& Clark (2007)} فرمول بسته‌ای برای نرخ بیت \lr{sifted} در سیستم \lr{BB84} تحت اثر \lr{dead time} استخراج کرده‌اند. معادله‌ی کلیدی (معادله‌ی ۱۵ مقاله) به‌صورت زیر است:

\begin{equation}
	\text{SBR} = \rho_{\text{TX}} \cdot 8 \cdot p \cdot P_{00}(p, k) \cdot S(p, k)
\end{equation}

در این رابطه، $\rho_{\text{TX}}$ نرخ ارسال (هرتز)، $p = L/8$ احتمال کلیک به ازای هر آشکارساز در هر چرخه (که در آن $L$ احتمال آشکارسازی یک فوتون از \lr{Alice} در \lr{Bob} است)، $k = \rho_{\text{TX}} \cdot \tau$ تعداد دوره‌های ارسال در هر \lr{dead time} (که در آن $\tau$ زمان مرده است)، $P_{00}$ احتمال حالت پایدار که هر دو آشکارساز در یک پایه زنده باشند (معادله‌ی ۸)، و $S(p, k)$ احتمال \lr{sifting} (معادله‌ی ۹) است.

تابع \lr{\texttt{\_compute\_rogers\_analytic\_sbr}} این فرمول را پیاده‌سازی می‌کند و از \lr{\texttt{link\_efficiency}} از \lr{qkd.channel} برای محاسبه‌ی $L = \eta_{\text{ch}} \cdot \eta_{\text{det}} \cdot \eta_{\text{coup}}$ استفاده می‌نماید. خروجی این تابع به‌عنوان ستون \lr{\texttt{analytic\_rogers\_sbr}} در \lr{CSV} ثبت می‌شود تا کاربر بتواند مستقیماً \lr{Monte-Carlo} را با منحنی بسته‌ی مقاله (شکل‌های ۳–۵) مقایسه کند.

\subsection{تئوری Dead Time آشکارساز: Non-Paralyzable}

مدل \lr{Non-Paralyzable} برای \lr{dead time} آشکارساز \lr{SPAD} به این معناست که پس از هر کلیک، آشکارساز برای مدت $\tau$ غیرفعال می‌شود و کلیک‌های در این بازه به‌طور کامل نادیده گرفته می‌شوند (نه این که زمان مرده را تمدید کنند). نرخ کلیک اندازه‌گیری‌شده $R_{\text{meas}}$ با نرخ واقعی $R_{\text{real}}$ رابطه‌ی زیر را دارد:

\begin{equation}
	R_{\text{meas}} = \frac{R_{\text{real}}}{1 + R_{\text{real}} \cdot \tau}
\end{equation}

در نسخه‌ی جدید، \lr{DeadTimeModel.PARALYZABLE} به‌عنوان \lr{deprecated} علامت‌گذاری شده (طبق \lr{Rogers et al.}) و \lr{ROGERS\_2007\_POLICY} به‌عنوان \lr{sentinel string} برای فعال‌سازی خط‌مشی \lr{sequence-collapsing sifting} معرفی شده است. این خط‌مشی در سطح پایه (نه آشکارساز) عمل می‌کند و حداکثر یک بیت \lr{sifted} به ازای هر دنباله‌ی آشکارسازی که \lr{dead time} را پوشش می‌دهد، تولید می‌کند.

\subsection{مدل Afterpulse آشکارساز}

\lr{Afterpulse} پدیده‌ای است که در آن یک کلیک واقعی باعث تولید کلیک‌های کاذب بعدی در آشکارساز می‌شود. دو مدل اصلی وجود دارد: \lr{Exponential} (که در آن احتمال \lr{afterpulse} با زمان به‌صورت نمایی کاهش می‌یابد) و \lr{Geometric} (که در آن احتمال در هر دوره‌ی ارسال ثابت است). در پیکربندی پیش‌فرض، مدل \lr{Exponential} با احتمال $0.001$ و طول عمر $10$ نانوثانیه به‌کار رفته است.

\subsection{نویز Raman در DWDM و مدل فیزیکی}

در حالت \lr{DWDM}، چندین کانال کلاسیک در کنار کانال کوانتومی در یک فیبر منتقل می‌شوند. پراکندگی \lr{Raman} باعث می‌شود فوتون‌های کانال‌های کلاسیک به طول‌موج کانال کوانتومی پخش شوند و نویز پس‌زمینه ایجاد نمایند. طبق \lr{Lim2014 Section IV}، نرخ فوتون‌های \lr{Raman} به‌صورت زیر مدل می‌شود:

\begin{equation}
	P_{\text{Raman}} = P_{\text{launch}} \cdot (N - 1) \cdot \gamma_{\text{Raman}} \cdot L_{\text{eff}}
\end{equation}

در این رابطه، $P_{\text{launch}} = 10^{P_{\text{dBm}}/10} \times 10^{-3}$ توان لانچ هر کانال (وات)، $N$ تعداد کانال‌های کلاسیک، $\gamma_{\text{Raman}}$ ضریب پراکندگی \lr{Raman}، و $L_{\text{eff}} = (1 - e^{-\alpha L}) / \alpha$ طول مؤثر غیرخطی است (که در آن $\alpha = \alpha_{\text{dB/km}} \cdot 0.23026$ نپر بر کیلومتر). نرخ فوتون به ازای هرتز با تقسیم بر انرژی فوتون $E_{\text{ph}} = hc/\lambda$ به دست می‌آید:

\begin{equation}
	R_{\text{Raman}} = \frac{P_{\text{Raman}}}{E_{\text{photon}}}
\end{equation}

تابع \lr{\texttt{\_compute\_raman\_dark\_rate\_hz}} این محاسبه را پیاده‌سازی می‌کند و خروجی را به‌صورت نرخ دارک کانت معادل (هرتز) برمی‌گرداند. این مقدار سپس به \lr{\texttt{detector\_run.dark\_rate}} و \lr{\texttt{params.detector.dark\_rate}} تزریق می‌شود.

\subsection{بهینه‌سازی پارامترهای پالس با GLLP}

هنگامی که پرچم \lr{\texttt{-{}-auto-optimize-pulses}} فعال است، به ازای هر (\lr{combo}, \lr{distance})، پارامترهای $\mu$ (شدت سیگنال)، $\nu$ (شدت دکوی)، $p_s$ (احتمال سیگنال)، $p_d$ (احتمال دکوی)، و $p_v$ (احتمال واکوئم) با \lr{SLSQP} بهینه می‌شوند. تابع هدف، نرخ کلید مجانبی \lr{GLLP} با جریمه‌ی متناهی است:

\begin{equation}
	R = q_{\text{sift}} \cdot \left[ Q_1 \cdot (1 - H_2(e_1)) - Q_\mu \cdot f_{\text{EC}} \cdot H_2(E_\mu) \right] - \text{finite-size penalty}
\end{equation}

این بهینه‌سازی در ماژول \lr{\texttt{qkd.proofs.optimization}} پیاده‌سازی شده و \lr{\texttt{main\_optimized}} فقط یک \lr{thin wrapper} به نام \lr{\texttt{\_optimize\_pulses\_for\_distance}} ارائه می‌دهد. پارامترهای بهینه در \lr{\texttt{\_ws.optimized\_params}} کش می‌شوند تا ردیف‌های \lr{noise=True} و \lr{noise=False} در همان (\lr{combo}, \lr{distance}) از پارامترهای یکسان استفاده کنند.

\section{تحلیل معماری و ساختار داده‌ها}

\subsection{DEFAULT\_CONFIG و DEFAULT\_SWEEPS}

\lr{\texttt{DEFAULT\_CONFIG}} یک دیکشنری تودرتوی با هفت بخش است که در خطوط ۱۱۸–۲۰۳ تعریف شده است. هر بخش مجموعه‌ای از پارامترهای پیش‌فرض را در بر می‌گیرد. در نسخه‌ی جدید، چندین تغییر کلیدی اعمال شده: \lr{\texttt{error\_correction\_efficiency}} از $1.1$ به $1.16$ افزایش یافته (مطابق مقادیر عملی \lr{LDPC})، \lr{\texttt{dark\_rate}} به $10^{-6}$ تنظیم شده، \lr{\texttt{dark\_rate\_d1}} به $2 \times 10^{-6}$ (نامتقارن)، \lr{\texttt{dead\_time\_ns}} به $10.0$ نانوثانیه، \lr{\texttt{jitter\_fwhm\_ns}} به $0.050$ نانوثانیه، و \lr{\texttt{afterpulse\_prob}} به $0.001$.

\lr{\texttt{DEFAULT\_SWEEPS}} در خطوط ۲۰۹–۲۴۵ تعریف شده و اکثراً به‌صورت \lr{comment} است. این طراحی عمدی است: کاربر می‌تواند با \lr{uncomment} کردن خطوط مورد نظر، پیمایش‌های رایج را فعال کند. یک کامنت مهم در خطوط ۲۳۰–۲۳۶ توضیح می‌دهد که \lr{DeadTimeModel.PARALYZABLE} به‌عنوان \lr{deprecated} علامت‌گذاری شده و به‌جای آن، \lr{ROGERS\_2007\_POLICY} به sweep اضافه شده است.

\subsection{الگوی Adapter: توابع build\_*}

توابع \lr{\texttt{build\_*}} به‌عنوان \lr{adapter} بین پیکربندی خام (دیکشنری) و اشیای \lr{dataclass} اعتبارسنجی‌شده عمل می‌کنند. این الگو چندین مزیت دارد: نخست، جداسازی \lr{concern}ها (پیکربندی در برابر اعتبارسنجی در برابر ساخت شیء)؛ دوم، امکان \lr{test} مستقل هر \lr{adapter}؛ سوم، مهاجرت تدریجی از \lr{dict} به \lr{typed object} (طبق \lr{TODO F-25}). توابع کلیدی عبارت‌اند از:

\lr{\texttt{build\_optical\_source\_config}}: از \lr{config["source"]} یک \lr{\texttt{OpticalSourceConfig}} می‌سازد. در نسخه‌ی جدید، \lr{\texttt{source\_rate}} از \lr{\texttt{pulse\_period\_ns}} استنتاج می‌شود (یا برعکس) چون \lr{\texttt{OpticalSourceConfig}} \lr{\texttt{pulse\_period\_ns}} را به‌عنوان \lr{@property} از \lr{\texttt{source\_rate}} محاسبه می‌کند. این استنتاج در خطوط ۵۲۲–۵۳۱ پیاده‌سازی شده و در صورت فقدان هر دو، \lr{\texttt{ConfigurationError}} پرتاب می‌شود.

\lr{\texttt{build\_intensity\_config}}: یک \lr{adapter} قدیمی است که فرض می‌کند دقیقاً ۳ نوع پالس (\lr{signal}، \lr{decoy}، \lr{vacuum}) با نام‌های ثابت وجود دارد. در نسخه‌ی جدید، تابع \lr{\texttt{build\_intensity\_config\_from\_source}} به‌عنوان مسیر ترجیحی معرفی شده که از \lr{\texttt{source.intensity\_config}} استفاده می‌کند و از جستجوی مبتنی بر نقش (یافتن \lr{"signal"} با نام، نه ایندکس) پشتیبانی می‌نماید.

\lr{\texttt{build\_detector}}: در نسخه‌ی جدید کاملاً بازنویسی شده است. رویکرد قدیمی (حذف شده) شامل دو مرحله بود: (۱) \lr{\texttt{build\_detection\_config}} که فقط ۳ فیلد داشت و (۲) \lr{\texttt{from\_config}} + ۱۰ \lr{setattr} برای پارامترهای از دست رفته. رویکرد جدید از \lr{\texttt{\_build\_detector\_config\_dict}} برای ساخت یک دیکشنری تخت با تمام پارامترها استفاده می‌کند و سپس \lr{\texttt{SinglePhotonDetector.from\_config\_dict}} را فراخوانی می‌نماید که در یک گام اعتبارسنجی کامل انجام می‌دهد.

\subsection{تابع parse\_enum و یکپارچه‌سازی Enum}

تابع \lr{\texttt{parse\_enum}} در خطوط ۳۳۲–۳۵۵ پیاده‌سازی شده و یک \lr{adapter} برای تبدیل رشته‌ها به اعضای \lr{Enum} است. این تابع ابتدا بررسی می‌کند که آیا مقدار از قبل عضو \lr{Enum} است (در این صورت بدون تغییر برمی‌گردد)، سپس رشته را با \lr{strip} و \lr{upper}/\lr{lower} نرمال‌سازی می‌کند و با نام و مقدار اعضای \lr{Enum} مقایسه می‌نماید. در صورت عدم تطابق، \lr{\texttt{ParameterValidationError}} با فهرست مقادیر معتبر پرتاب می‌شود.

یک مورد خاص در \lr{\texttt{normalize\_detector\_config}} (خط ۳۵۷) وجود دارد: \lr{\texttt{ROGERS\_2007\_POLICY}} به‌عنوان \lr{sentinel string} (نه عضو \lr{Enum}) تعریف شده و در صورت مشاهده، بدون تغییر عبور می‌کند (\lr{pass through}). این طراحی به این دلیل است که \lr{DoubleClickPolicy} در \lr{qkd.datatypes} تعریف شده و ما عمداً آن را تغییر نمی‌دهیم.

\subsection{کلاس \_WorkerState و الگوی State Encapsulation}

کلاس \lr{\texttt{\_WorkerState}} در خطوط ۱۳۶۶–۱۴۰۵ تعریف شده و \lr{state} mutable هر کارگر را در یک شیء واحد کپسوله می‌کند. این الگو (طبق \lr{F-03}) جایگزین متغیرهای سراسری ماژول شده است. هر فرآیند کارگر نسخه‌ی خود را از طریق \lr{\texttt{init\_worker}} دریافت می‌کند و این شیء بین فرآیندها به اشتراک گذاشته نمی‌شود.

\lr{\texttt{\_WorkerState}} از \lr{\texttt{\_\_slots\_\_}} استفاده می‌کند که در پایتون، حافظه‌ی کم‌تری مصرف می‌کند و دسترسی به ویژگی‌ها را سریع‌تر می‌سازد. اسلات‌های کلیدی عبارت‌اند از: \lr{\texttt{base\_config}}، \lr{\texttt{combo\_map}}، \lr{\texttt{last\_combo\_idx}} (برای \lr{lazy rebuild})، \lr{\texttt{source}}، \lr{\texttt{detector}}، \lr{\texttt{protocol}} (اشیای فیزیکی ساخته‌شده)، \lr{\texttt{source\_meta\_cache}} (کش metadata منبع برای جلوگیری از سریالی‌سازی مکرر، طبق \lr{F-14})، \lr{\texttt{detector\_overrides}} (ممیزی \lr{override}های آشکارساز)، \lr{\texttt{auto\_optimize\_pulses}}، \lr{\texttt{last\_optimized\_dist}}، \lr{\texttt{optimized\_params}} (کش بهینه‌ساز per-(\lr{combo}, \lr{distance}))، \lr{\texttt{dwdm}}، \lr{\texttt{wdm\_channel\_count}}، \lr{\texttt{wdm\_channel\_power\_dbm}}، \lr{\texttt{wdm\_raman\_coefficient}} (پارامترهای \lr{DWDM})، و \lr{\texttt{params}} (کش \lr{QKDParams} برای مسیر خطا).

\subsection{ساختار داده‌ی TallyCounts}

\lr{\texttt{TallyCounts}} یک \lr{dataclass} است که آمار شمارش پالس‌ها را در بر می‌گیرد: \lr{\texttt{sent}} (کل ارسال‌شده)، \lr{\texttt{sifted}} (کل \lr{sifted})، \lr{\texttt{errors\_sifted}} (خطاهای \lr{sifted})، \lr{\texttt{double\_clicks\_discarded}}، \lr{\texttt{sent\_z}}، \lr{\texttt{sent\_x}}، \lr{\texttt{sifted\_z}}، \lr{\texttt{sifted\_x}}، \lr{\texttt{errors\_sifted\_z}}، \lr{\texttt{errors\_sifted\_x}}. این ساختار امکان تفکیک پایه‌های $Z$ (مستقیم) و $X$ (قطری) را فراهم می‌سازد که در پروتکل \lr{BB84} برای تخمین خطای فاز ضروری است.

تابع \lr{\texttt{tally\_counts\_from\_sifting\_results}} (خط ۱۲۰۴) این ساختار را از نتایج \lr{sifting} می‌سازد. این تابع با استفاده از \lr{\texttt{\_basis\_is\_z\_mask}} ماسک‌های پایه‌ی $Z$ را استخراج می‌کند و سپس برای هر نوع پالس (\lr{pulse\_type\_index})، شمارش‌های جداگانه برای پایه‌های $Z$ و $X$ را محاسبه می‌نماید. در صورت عدم دسترسی به ماسک پایه، از تقسیم برابر $50/50$ به‌عنوان fallback استفاده می‌کند.

\subsection{SimulationResults به عنوان خروجی جامع}

\lr{\texttt{SimulationResults}} یک \lr{dataclass} جامع است که خروجی کامل یک نقطه‌ی شبیه‌سازی را در بر می‌گیرد: \lr{\texttt{params}} (\lr{ProtocolParameters})، \lr{\texttt{metadata}}، \lr{\texttt{security\_certificate}}، \lr{\texttt{decoy\_estimates}}، \lr{\texttt{secure\_key\_length}}، \lr{\texttt{raw\_sifted\_key\_length}}، \lr{\texttt{simulation\_time\_seconds}}، \lr{\texttt{status}}، \lr{\texttt{tally\_stats}}، \lr{\texttt{weak\_gllp\_rate}}، \lr{\texttt{tagged\_fraction\_bound}}، \lr{\texttt{beta\_y1\_deviation}}، \lr{\texttt{beta\_e1\_deviation}}، و \lr{\texttt{success\_probability}}. تابع \lr{\texttt{build\_simulation\_results}} (خط ۱۳۳۹) این شیء را می‌سازد.

\subsection{DETECTOR\_METADATA\_KEYS، DETECTOR\_OVERRIDE\_KEYS و PARAMS\_METADATA\_KEYS}

در نسخه‌ی جدید، سه گروه ستون اضافی به \lr{CSV} اضافه شده‌اند. \lr{\texttt{DETECTOR\_METADATA\_KEYS}} (خط ۱۴۱۶) از \lr{\texttt{DetectionDiagnostics.metadata\_keys()}} استخراج می‌شود تا ستون‌های \lr{CSV} همیشه با \lr{dataclass} هماهنگ باشد. \lr{\texttt{DETECTOR\_OVERRIDE\_KEYS}} (خط ۱۴۲۳) شامل ۹ ستون برای ممیزی \lr{override}های آشکارساز است (مانند \lr{\texttt{det\_override\_det\_eff\_d0}}). \lr{\texttt{PARAMS\_METADATA\_KEYS}} (خط ۱۴۳۹) شامل ۷ ستون برای ممیزی \lr{DWDM} است (مانند \lr{\texttt{params\_wdm\_channel\_count}}، \lr{\texttt{params\_raman\_dark\_rate\_hz}}).

\subsection{الگوی SimParams: ChannelSimParams، SourceSimParams، DetectorSimParams}

در نسخه‌ی جدید، الگوی \lr{SimParams} معرفی شده است. سه تابع \lr{\texttt{build\_channel\_sim\_params}}، \lr{\texttt{build\_source\_sim\_params}}، و \lr{\texttt{build\_detector\_sim\_params}} اشیای \lr{frozen} می‌سازند که تمام مقادیر فیزیکی اعتبارسنجی‌شده را در بر دارند. حلقه‌ی شبیه‌سازی به‌جای خواندن از \lr{config dict}، از این اشیای \lr{frozen} می‌خواند. این الگو تضمین می‌کند که اگر کاربر پیکربندی را با \lr{sweep} تغییر دهد، تمام مقادیر اعتبارسنجی شده و سپس به حلقه ارسال شوند.

\section{پیاده‌سازی ریاضی و الگوریتمی}

\subsection{تابع set\_nested\_value: اعمال overrides روی پیکربندی تودرتو}

تابع \lr{\texttt{set\_nested\_value}} (خط ۲۴۸) یک مقدار را در دیکشنری تودرتو با مسیر نقطه‌گذاری (\lr{dot-separated}) تنظیم می‌کند. این تابع از پارامتر \lr{\texttt{strict}} پشتیبانی می‌کند که در حالت \lr{True}، در صورت عدم وجود کلید میانی، \lr{\texttt{ConfigurationError}} پرتاب می‌کند (طبق \lr{F-10}).

\begin{LTR}
	\begin{lstlisting}[language=Python]
		def set_nested_value(d: Dict, key_path: str, value: Any, *, strict: bool = False):
		if not key_path or not isinstance(key_path, str):
		raise ConfigurationError(
		"Sweep override key_path must be a non-empty dotted string.",
		context={"key_path": key_path},
		)
		keys = key_path.split('.')
		current = d
		for key in keys[:-1]:
		if not isinstance(current, dict):
		raise ConfigurationError(
		"Cannot apply sweep override through a non-dictionary config node.",
		context={"key_path": key_path, "blocked_at": key},
		)
		if strict and key not in current:
		raise ConfigurationError(
		"Sweep override key_path references a missing intermediate key.",
		context={"key_path": key_path, "missing_key": key},
		)
		current = current.setdefault(key, {})
		current[keys[-1]] = value
	\end{lstlisting}
\end{LTR}

در خطوط ابتدایی، اعتبارسنجی \lr{key\_path} انجام می‌شود. سپس با \lr{\texttt{split('.')}} مسیر به فهرست کلیدها تجزیه می‌شود. حلقه‌ی \lr{\texttt{for}} تا کلید یکی‌مانده‌به‌آخر می‌چرخد و در هر مرحله، بررسی می‌کند که \lr{current} از نوع \lr{dict} باشد. در حالت \lr{strict}، عدم وجود کلید میانی باعث خطا می‌شود. در پایان، کلید آخر با مقدار مورد نظر تنظیم می‌گردد.

\subsection{تابع build\_epsilon\_allocation: اعتبارسنجی بودجۀ امنیتی}

تابع \lr{\texttt{build\_epsilon\_allocation}} (خط ۴۰۳) تخصیص \lr{epsilon} را از پیکربندی می‌خواند و اعتبارسنجی می‌نماید.

\begin{LTR}
	\begin{lstlisting}[language=Python]
		def build_epsilon_allocation(config: Dict[str, Any]) -> EpsilonAllocation:
		eps = config["protocol"]["epsilons"]
		allocation = EpsilonAllocation(
		eps_sec=eps["eps_sec"], eps_cor=eps["eps_cor"], eps_pe=eps["eps_pe"],
		eps_smooth=eps["eps_smooth"], eps_pa=eps["eps_pa"],
		eps_phase_est=eps["eps_phase_est"],
		)
		total_eps = (
		allocation.eps_sec + allocation.eps_cor + allocation.eps_pe +
		allocation.eps_smooth + allocation.eps_pa + allocation.eps_phase_est
		)
		if total_eps >= 1.0:
		raise ParameterValidationError(
		f"Total epsilon allocation ({total_eps:.2e}) must be less than 1.",
		param_name="protocol.epsilons",
		param_value=total_eps,
		context={"total_epsilon": total_eps},
		)
		for field_name in ("eps_sec", "eps_cor", "eps_pe", "eps_smooth", "eps_pa", "eps_phase_est"):
		val = getattr(allocation, field_name)
		if val <= 0:
		raise ParameterValidationError(
		f"Epsilon value {field_name}={val} must be positive.",
		param_name=f"protocol.epsilons.{field_name}",
		param_value=val,
		)
		return allocation
	\end{lstlisting}
\end{LTR}

ابتدا \lr{EpsilonAllocation} از پیکربندی ساخته می‌شود. سپس مجموع \lr{epsilon} محاسبه می‌شود و در صورت بزرگ‌تر یا مساوی $1$، خطا پرتاب می‌گردد. در پایان، هر مقدار بررسی می‌شود که مثبت باشد. این اعتبارسنجی (طبق \lr{F-17}) تضمین می‌کند که بودجه‌ی امنیتی معتبر باشد.

\subsection{تابع build\_optical\_source\_config: استنتاج source\_rate}

تابع \lr{\texttt{build\_optical\_source\_config}} (خط ۵۱۳) یک \lr{\texttt{OpticalSourceConfig}} از پیکربندی می‌سازد. نکته‌ی کلیدی، استنتاج \lr{\texttt{source\_rate}} از \lr{\texttt{pulse\_period\_ns}} است:

\begin{LTR}
	\begin{lstlisting}[language=Python]
		if "source_rate" in source_cfg and source_cfg["source_rate"] is not None:
		source_rate = float(source_cfg["source_rate"])
		elif "pulse_period_ns" in source_cfg and source_cfg["pulse_period_ns"] is not None:
		pulse_period_ns = float(source_cfg["pulse_period_ns"])
		source_rate = 1e9 / pulse_period_ns if pulse_period_ns > 0 else 1e8
		else:
		raise ConfigurationError(
		"source config must contain either 'source_rate' or 'pulse_period_ns'.",
		context={"section": "source", "available_keys": sorted(source_cfg.keys())},
		)
	\end{lstlisting}
\end{LTR}

در این بلاک، اولویت با \lr{\texttt{source\_rate}} مستقیم است. در صورت عدم وجود، \lr{\texttt{pulse\_period\_ns}} به \lr{\texttt{source\_rate}} تبدیل می‌شود. در صورت فقدان هر دو، خطا پرتاب می‌شود. این الگو به این دلیل ضروری است که \lr{\texttt{OpticalSourceConfig}} \lr{\texttt{pulse\_period\_ns}} را به‌عنوان \lr{@property} از \lr{\texttt{source\_rate}} محاسبه می‌کند و در نتیجه، فقط یکی از این دو باید در پیکربندی ارائه شود.

\subsection{تابع build\_detector: الگوی from\_config\_dict}

تابع \lr{\texttt{build\_detector}} (خط ۹۰۷) در نسخه‌ی جدید کاملاً بازنویسی شده است. رویکرد قدیمی (حذف شده) شامل دو مرحله بود: \lr{\texttt{build\_detection\_config}} (با ۳ فیلد) + \lr{\texttt{from\_config}} + ۱۰ \lr{setattr}. رویکرد جدید از \lr{\texttt{\_build\_detector\_config\_dict}} برای ساخت یک دیکشنری تخت با تمام پارامترها استفاده می‌کند:

\begin{LTR}
	\begin{lstlisting}[language=Python]
		def build_detector(config: Dict[str, Any]) -> Tuple[SinglePhotonDetector, Dict[str, Any]]:
		det_config_dict = _build_detector_config_dict(config)
		detector = SinglePhotonDetector.from_config_dict(det_config_dict)
		actual_config = detector.to_config_dict(include_experimental=True)
		input_config = det_config_dict
		_overrides_applied: Dict[str, Any] = {}
		override_keys = (
		"det_eff_d0", "det_eff_d1", "dark_rate", "dark_rate_d1",
		"bias_voltage", "breakdown_voltage", "temperature_k",
		"ref_temperature_k", "ref_bias_voltage",
		)
		for key in override_keys:
		input_val = input_config.get(key)
		actual_val = actual_config.get(key)
		if actual_val is not None:
		_overrides_applied[key] = actual_val
		if input_val != actual_val:
		logger.debug(
		"Detector parameter %s: input=%s, actual=%s (default/fallback applied).",
		key, input_val, actual_val,
		)
		return detector, _overrides_applied
	\end{lstlisting}
\end{LTR}

ابتدا \lr{\texttt{\_build\_detector\_config\_dict}} دیکشنری تخت را با نرمال‌سازی \lr{Enum} و تبدیل نوع‌های عددی می‌سازد. سپس \lr{\texttt{from\_config\_dict}} آشکارساز را در یک گام اعتبارسنجی می‌کند. در پایان، با مقایسه‌ی \lr{\texttt{to\_config\_dict()}} با ورودی، \lr{overrides} واقعی استخراج می‌شود. این الگو تضمین می‌کند که \lr{CSV} همیشه مقادیر واقعی آشکارساز را نشان دهد، حتی اگر \lr{fallback} یا نرمال‌سازی داخلی اعمال شده باشد.

\subsection{تابع tally\_counts\_from\_sifting\_results}

تابع \lr{\texttt{tally\_counts\_from\_sifting\_results}} (خط ۱۲۰۴) آمار شمارش را از نتایج \lr{sifting} می‌سازد. این تابع از \lr{\texttt{\_basis\_is\_z\_mask}} برای استخراج ماسک پایه‌ی $Z$ استفاده می‌کند:

\begin{LTR}
	\begin{lstlisting}[language=Python]
		def _basis_is_z_mask(values) -> np.ndarray:
		arr = np.asarray(values)
		if arr.dtype == np.bool_:
		return arr.astype(bool)
		if np.issubdtype(arr.dtype, np.integer):
		return arr == 0
		out = np.zeros(arr.shape, dtype=bool)
		flat = arr.ravel()
		for i, v in enumerate(flat):
		if isinstance(v, str):
		s = v.strip().upper()
		else:
		name = getattr(v, "name", None)
		s = str(name).strip().upper() if name is not None else str(v).strip().upper()
		out.ravel()[i] = s in {"Z", "Z_BASIS", "RECTILINEAR", "STANDARD", "0"}
		return out
	\end{lstlisting}
\end{LTR}

این تابع با آرایه‌های بولی، صحیح و رشته‌ای کار می‌کند. در حالت رشته‌ای، نام پایه را با مجموعه‌ای از نام‌های متداول (\lr{Z}، \lr{Z\_BASIS}، \lr{RECTILINEAR}، \lr{STANDARD}، \lr{0}) مقایسه می‌نماید. این طراحی انعطاف‌پذیر، امکان پذیرش فرمت‌های مختلف ورودی را فراهم می‌سازد.

\subsection{تابع \_compute\_raman\_dark\_rate\_hz: نویز Raman}

تابع \lr{\texttt{\_compute\_raman\_dark\_rate\_hz}} (خط ۱۵۴۴) نرخ دارک کانت معادل \lr{Raman} را محاسبه می‌کند. این تابع در دو مسیر استفاده می‌شود: مسیر \lr{Monte-Carlo} (تزریق به \lr{\texttt{detector\_run.dark\_rate}}) و مسیر اثبات (تزریق به \lr{\texttt{params.detector.dark\_rate}}):

\begin{LTR}
	\begin{lstlisting}[language=Python]
		def _compute_raman_dark_rate_hz(params) -> float:
		if params is None:
		return 0.0
		wdm_count = getattr(params, "wdm_channel_count", 1)
		if wdm_count is None or wdm_count <= 1:
		return 0.0
		wdm_power_dbm = float(getattr(params, "wdm_channel_power_dbm", 0.0))
		p_launch = (10.0 ** (wdm_power_dbm / 10.0)) * 1e-3  # Watts
		ch = getattr(params, "channel", None)
		if ch is None:
		return 0.0
		fiber_loss_db_km = float(getattr(params, "fiber_loss_db_km",
		getattr(ch, "fiber_loss_db_km", 0.2)))
		distance_km = float(getattr(ch, "distance_km", 0.0))
		alpha = fiber_loss_db_km * 0.23026  # Np/km
		if alpha > 1e-12:
		l_eff = (1.0 - math.exp(-alpha * distance_km)) / alpha
		else:
		l_eff = distance_km
		wdm_raman_coeff = float(getattr(params, "wdm_raman_coefficient", 0.0))
		p_raman = p_launch * (wdm_count - 1) * wdm_raman_coeff * l_eff
		lam_nm = float(getattr(params, "carrier_wavelength_nm", 1550.0))
		photon_energy = (_const.h * _const.c) / (lam_nm * 1e-9)
		rate_hz = p_raman / photon_energy
		det = getattr(params, "detector", None)
		det_eff_d0 = float(getattr(det, "det_eff_d0", 1.0)) if det is not None else 1.0
		return rate_hz * det_eff_d0
	\end{lstlisting}
\end{LTR}

در ابتدا، در صورت \lr{wdm\_count <= 1}، مقدار صفر بازگردانده می‌شود (حالت تک‌کانال). سپس $P_{\text{launch}}$ از \lr{dBm} به وات تبدیل می‌شود. $\alpha$ از \lr{dB/km} به \lr{Np/km} تبدیل می‌شود (با ضریب $0.23026 = \ln(10)/10$). $L_{\text{eff}}$ محاسبه می‌شود. سپس $P_{\text{Raman}}$ و $R_{\text{Raman}}$ محاسبه می‌گردند. در پایان، ضرب در \lr{\texttt{det\_eff\_d0}} باعث می‌شود که نرخ نهایی، نرخ کلیک پس از بازده آشکارساز باشد (مطابق با قرارداد \lr{\texttt{dark\_rate}} در \lr{qkd.detectors}).

نکته‌ی مهم (طبق کامنت \lr{Option J}) این است که این تابع \lr{\texttt{params.fiber\_loss\_db\_km}} را به‌جای \lr{\texttt{channel.fiber\_loss\_km}} می‌خواند. در حالت \lr{DWDM}، کانال از \lr{\texttt{FiberChannel.from\_total\_loss}} بازسازی می‌شود که \lr{alpha} را از کل تلفات (فیبر + فیلتر \lr{AWG}) استنتاج می‌کند. استفاده از \lr{alpha} آلوده باعث می‌شود \lr{Raman} ۹–۱۹\% کم‌برآورد شود.

\subsection{تابع \_compute\_rogers\_analytic\_sbr}

تابع \lr{\texttt{\_compute\_rogers\_analytic\_sbr}} (خط ۱۷۴۲) نرخ بیت \lr{sifted} تحلیلی را طبق معادله‌ی ۱۵ مقاله‌ی \lr{Rogers et al. (2007)} محاسبه می‌کند:

\begin{LTR}
	\begin{lstlisting}[language=Python]
		def _compute_rogers_analytic_sbr(
		*, pulse_period_ns, dead_time_ns, channel_transmittance,
		det_eff_d0, coupling_efficiency=1.0, eps=0.0,
		) -> Optional[float]:
		if dead_time_ns <= 0.0 or pulse_period_ns <= 0.0:
		return None
		L = link_efficiency(channel_transmittance, det_eff_d0, coupling_efficiency)
		if L <= 0.0:
		return None
		L = min(L, 1.0)
		rho_tx = 1e9 / float(pulse_period_ns)
		tau_s = float(dead_time_ns) * 1e-9
		try:
		return float(rogers_sifted_bit_rate(
		rho_tx=rho_tx, L=L, tau=tau_s, eps=eps,
		))
		except (ValueError, OverflowError, ArithmeticError) as exc:
		logger.debug("Rogers analytic SBR computation failed: %s.", exc)
		return None
	\end{lstlisting}
\end{LTR}

در ابتدا، اعتبارسنجی \lr{dead\_time\_ns} و \lr{pulse\_period\_ns} انجام می‌شود. سپس $L$ با \lr{\texttt{link\_efficiency}} محاسبه می‌شود (این تابع از \lr{qkd.channel} وارد شده و $L = \eta_{\text{ch}} \cdot \eta_{\text{det}} \cdot \eta_{\text{coup}}$ را محاسبه می‌نماید). در پایان، \lr{\texttt{rogers\_sifted\_bit\_rate}} (که از \lr{qkd.detectors} وارد شده) فراخوانی می‌شود. در صورت بروز خطای محاسباتی، \lr{None} بازگردانده می‌شود تا در \lr{CSV} به‌صورت سلول خالی ثبت شود.

\subsection{تابع \_balance\_decoy\_probabilities}

تابع \lr{\texttt{\_balance\_decoy\_probabilities}} (خط ۱۷۹۰) احتمال‌های پالس را به‌صورت خودکار متعادل می‌کند. هنگامی که یک \lr{sweep} فقط یکی از احتمال‌ها (\lr{signal} یا \lr{decoy}) را override می‌کند، این تابع دیگری را محاسبه می‌نماید:

\begin{LTR}
	\begin{lstlisting}[language=Python]
		def _balance_decoy_probabilities(overrides, config):
		sig_p = overrides.get("source.pulses.signal.prob")
		dec_p = overrides.get("source.pulses.decoy.prob")
		vac_p = config["source"]["pulses"]["vacuum"]["prob"]
		if sig_p is not None and dec_p is None:
		balanced_dec = round(1.0 - sig_p - vac_p, 9)
		if balanced_dec < 0:
		raise ConfigurationError(
		f"Probability balancing produced negative decoy probability ({balanced_dec}).",
		context={"signal_prob": sig_p, "vacuum_prob": vac_p, "decoy_prob": balanced_dec},
		)
		overrides["source.pulses.decoy.prob"] = balanced_dec
		elif dec_p is not None and sig_p is None:
		balanced_sig = round(1.0 - dec_p - vac_p, 9)
		if balanced_sig < 0:
		raise ConfigurationError(
		f"Probability balancing produced negative signal probability ({balanced_sig}).",
		context={"decoy_prob": dec_p, "vacuum_prob": vac_p, "signal_prob": balanced_sig},
		)
		overrides["source.pulses.signal.prob"] = balanced_sig
	\end{lstlisting}
\end{LTR}

این تابع (طبق \lr{F-11}) اعتبارسنجی می‌کند که احتمال‌های متعادل‌شده در بازه‌ی $[0, 1]$ باشند. در صورت منفی شدن، خطا پرتاب می‌شود. این الگو امکان \lr{sweep} فقط یک احتمال را فراهم می‌سازد و بقیه به‌صورت خودکار محاسبه می‌شوند.

\subsection{تابع \_optimize\_pulses\_for\_distance: بهینه‌سازی پالس}

تابع \lr{\texttt{\_optimize\_pulses\_for\_distance}} (خط ۱۸۴۵) یک \lr{thin wrapper} به \lr{\texttt{optimize\_pulses\_for\_distance}} در \lr{\texttt{qkd.proofs.optimization}} است:

\begin{LTR}
	\begin{lstlisting}[language=Python]
		def _optimize_pulses_for_distance(config, dist_km, total_pulses, proof_name):
		return optimize_pulses_for_distance(config, dist_km, total_pulses, proof_name)
	\end{lstlisting}
\end{LTR}

این تابع پارامترهای بهینه $(\mu, \nu, p_s, p_d, p_v)$ را برمی‌گرداند. در \lr{\texttt{run\_single\_simulation}}، این پارامترها در \lr{\texttt{\_ws.optimized\_params}} کش می‌شوند و کلید کش \lr{(combo\_idx, dist)} است. این طراحی تضمین می‌کند که ردیف‌های \lr{noise=True} و \lr{noise=False} در همان (\lr{combo}, \lr{distance}) از پارامترهای یکسان استفاده کنند که برای مقایسه‌ی منصفانه ضروری است.

\subsection{تابع init\_worker و الگوی initializer}

تابع \lr{\texttt{init\_worker}} (خط ۱۴۴۹) در نسخه‌ی جدید امضای گسترده‌ای دارد:

\begin{LTR}
	\begin{lstlisting}[language=Python]
		def init_worker(base_config, combo_map, auto_optimize_pulses=False, dwdm=False,
		wdm_channel_count=4, wdm_channel_power_dbm=None,
		wdm_raman_coefficient=None):
		logging.basicConfig(
		level=logging.INFO,
		format="%(levelname)s [%(name)s] %(message)s",
		force=True,
		)
		_ws.base_config = base_config
		_ws.combo_map = combo_map
		_ws.last_combo_idx = -1
		_ws.source_meta_cache = {}
		_ws.detector_overrides = {}
		_ws.auto_optimize_pulses = bool(auto_optimize_pulses)
		_ws.last_optimized_dist = None
		_ws.optimized_params = None
		_ws.params = None
		_ws.dwdm = bool(dwdm)
		_ws.wdm_channel_count = int(wdm_channel_count) if wdm_channel_count is not None else None
		_ws.wdm_channel_power_dbm = float(wdm_channel_power_dbm) if wdm_channel_power_dbm is not None else None
		_ws.wdm_raman_coefficient = float(wdm_raman_coefficient) if wdm_raman_coefficient is not None else None
	\end{lstlisting}
\end{LTR}

در ابتدا، \lr{\texttt{logging.basicConfig}} با \lr{force=True} فراخوانی می‌شود تا \lr{ERROR}-level در فرآیندهای کارگر قابل مشاهده باشد (بدون این، کارگران با \lr{WARNING}-level پیش‌فرض شروع می‌شوند و \lr{ERROR}-ها خاموش می‌شوند). سپس تمام فیلدهای \lr{\texttt{\_ws}} مقداردهی می‌شوند. نکته‌ی مهم این است که \lr{wdm\_channel\_count}، \lr{wdm\_channel\_power\_dbm} و \lr{wdm\_raman\_coefficient} با تبدیل نوع (\lr{int}/\lr{float}) و بررسی \lr{None} ذخیره می‌شوند (طبق \lr{FIX-WS-WDM-V2}).

\subsection{تابع run\_single\_simulation: هستۀ اصلی}

تابع \lr{\texttt{run\_single\_simulation}} (خط ۱۸۵۷) هسته‌ی اصلی شبیه‌سازی است. این تابع یک \lr{args\_tuple} شامل (\lr{dist}، \lr{total\_pulses}، \lr{apply\_noise}، \lr{worker\_seed}، \lr{combo\_idx}) می‌گیرد و یک دیکشنری \lr{row} برمی‌گرداند. جریان اصلی به چند مرحله تقسیم می‌شود:

\textbf{مرحله‌ی ۱: Lazy rebuild.} در صورت تغییر \lr{combo\_idx}، پیکربندی با \lr{deepcopy} کپی شده، \lr{overrides} اعمال می‌شود، و اشیای \lr{source}، \lr{detector}، \lr{protocol} بازسازی می‌شوند. سپس \lr{metadata} در \lr{\texttt{\_ws.source\_meta\_cache}} کش می‌شود.

\textbf{مرحله‌ی ۲: Auto-optimize pulses.} در صورت فعال بودن \lr{\texttt{\_ws.auto\_optimize\_pulses}} و تغییر (\lr{combo}, \lr{dist})، \lr{\texttt{\_optimize\_pulses\_for\_distance}} فراخوانی می‌شود، پارامترهای بهینه در \lr{config} نوشته می‌شوند، و \lr{source} و \lr{protocol} بازسازی می‌گردند.

\textbf{مرحله‌ی ۳: بذرگیری.} با \lr{\texttt{SeedSequence.spawn}} (طبق \lr{F-05})، ۵ جریان \lr{RNG} مستقل از \lr{worker\_seed} استخراج می‌شوند:

\begin{LTR}
	\begin{lstlisting}[language=Python]
		seed_seq = np.random.SeedSequence(worker_seed)
		child_seqs = seed_seq.spawn(5)
		rng_prepare = np.random.default_rng(child_seqs[0])
		rng_photons = np.random.default_rng(child_seqs[1])
		rng_detect = np.random.default_rng(child_seqs[2])
		rng_sift = np.random.default_rng(child_seqs[3])
		rng_pe = np.random.default_rng(child_seqs[4])
	\end{lstlisting}
\end{LTR}

این الگو (به‌جای \lr{seed+1}، \lr{seed+2} و...) تضمین می‌کند که جریان‌های \lr{RNG} از نظر آماری مستقل باشند.

\textbf{مرحله‌ی ۴: بارگذاری QKDParams.} در نسخه‌ی جدید (طبق \lr{FIX-PARAMS-ORDER})، \lr{\texttt{load\_lim2014\_*\_params}} \emph{قبل} از بلاک \lr{apply\_noise} فراخوانی می‌شود تا تزریق \lr{Raman} از \lr{params} صحیح (مربوط به همین فاصله) استفاده کند. در حالت \lr{DWDM}، \lr{\texttt{load\_lim2014\_dwdm\_params}} با پارامترهای \lr{WDM} فراخوانی می‌شود.

\textbf{مرحله‌ی ۵: تزریق Raman به params.} در حالت \lr{DWDM} و \lr{apply\_noise=True}، نرخ \lr{Raman} محاسبه شده و به \lr{\texttt{params.detector.dark\_rate}} تزریق می‌شود (با \lr{\texttt{dataclasses.replace}} چون \lr{params} \lr{frozen} است).

\textbf{مرحله‌ی ۶: ساخت کانال.} \lr{\texttt{build\_channel\_from\_distance}} کانال را می‌سازد. در حالت \lr{DWDM}، تلفات فیلتر \lr{AWG} به کانال اضافه می‌شود.

\textbf{مرحله‌ی ۷: ساخت SimParams.} \lr{\texttt{build\_channel\_sim\_params}}، \lr{\texttt{build\_source\_sim\_params}} و (بعداً) \lr{\texttt{build\_detector\_sim\_params}} اشیای \lr{frozen} می‌سازند. اعتبارسنجی \lr{pulse\_period\_ns} انجام می‌شود.

\textbf{مرحله‌ی ۸: آماده‌سازی حالت‌ها و فوتون‌ها.} \lr{\texttt{protocol.prepare\_states}}، \lr{\texttt{generate\_photons\_for\_prepared\_states}}، و \lr{\texttt{infer\_ideal\_outcomes\_d0}} فراخوانی می‌شوند.

\textbf{مرحله‌ی ۹: شبیه‌سازی آشکارساز.} در حالت \lr{apply\_noise=True}، \lr{\texttt{detector.clone()}} ساخته می‌شود و نویز \lr{Raman} به \lr{\texttt{detector\_run.dark\_rate}} تزریق می‌گردد. در حالت \lr{noise=False}، \lr{\texttt{detector.noiseless\_copy()}} استفاده می‌شود. سپس \lr{\texttt{simulate\_detection}} فراخوانی می‌شود.

\textbf{مرحله‌ی ۱۰: Sifting و tally.} \lr{\texttt{detector\_result\_to\_protocol\_detection\_results}}، \lr{\texttt{protocol.sift\_results}}، \lr{\texttt{sample\_for\_parameter\_estimation}}، و \lr{\texttt{tally\_counts\_from\_sifting\_results}} فراخوانی می‌شوند.

\textbf{مرحله‌ی ۱۱: محاسبه کلید امن.} این مرحله در بخش بعد به‌تفصیل بررسی می‌شود.

\textbf{مرحله‌ی ۱۲: ساخت ردیف CSV.} با محاسبه‌ی \lr{QBER} اصلاح‌شده و فاصله‌ی اطمینان، ردیف نهایی ساخته می‌شود.

\subsection{مدیریت API متفاوت کلاس‌های اثبات}

در نسخه‌ی جدید، مدیریت اثبات‌های مختلف در خطوط ۲۲۹۳–۲۳۸۰ پیاده‌سازی شده است:

\begin{LTR}
	\begin{lstlisting}[language=Python]
		proof_name = config["source"].get("intended_proof", "LIM_2014").upper()
		if proof_name != "LIM_2014":
		import dataclasses
		updates = {}
		if hasattr(params, 'auto_optimize_lim2014'):
		updates['auto_optimize_lim2014'] = False
		if not hasattr(params, 'eps_pa'):
		updates['eps_pa'] = 1e-10
		if not hasattr(params, 'eps_sif'):
		updates['eps_sif'] = 1e-10
		if updates:
		params = dataclasses.replace(params, **updates)
		
		if proof_name == "MA_2005":
		from qkd.proofs.ma2005 import Ma2005VacuumWeakProof
		proof = Ma2005VacuumWeakProof(params)
		elif proof_name == "WANG_2005":
		from qkd.proofs.wang2005 import Wang2005Proof
		proof = Wang2005Proof(params)
		elif proof_name == "TIGHT":
		from qkd.proofs.tight import BB84TightProof
		proof = BB84TightProof(params)
		else:
		proof = Lim2014Proof(params)
	\end{lstlisting}
\end{LTR}

ابتدا نام اثبات از پیکربندی خوانده می‌شود. در صورت عدم استفاده از \lr{Lim2014}، پارامترهای \lr{auto\_optimize\_lim2014} به \lr{False} و \lr{eps\_pa}/\lr{eps\_sif} (در صورت فقدان) به مقادیر پیش‌فرض تنظیم می‌شوند. سپس کلاس اثبات مناسب \lr{import} و نمونه‌سازی می‌شود. \lr{import} به‌صورت \lr{lazy} (داخل \lr{if}) انجام می‌شود تا ماژول‌های غیرضروری در \lr{startup} بارگذاری نشوند.

سپس فراخوانی اثبات به دو مسیر تقسیم می‌شود:

\begin{LTR}
	\begin{lstlisting}[language=Python]
		if proof_name == "WANG_2005":
		mapped_stats = {
			"signal": stats.get("0", TallyCounts()),
			"decoy": stats.get("1", TallyCounts()),
			"vacuum": stats.get("2", TallyCounts())
		}
		decoy_estimates = proof.estimate_yields_and_errors(stats_map=mapped_stats)
		key_result = proof.calculate_key_length(decoy_estimates, stats_map=mapped_stats)
		elif proof_name == "MA_2005":
		decoy_estimates = proof.estimate_yields_and_errors(stats_map=stats)
		key_result = proof.calculate_key_length(decoy_estimates, stats_map=stats)
		else:
		key_result = proof.calculate_key_length(stats_map=stats)
	\end{lstlisting}
\end{LTR}

\lr{Wang2005} به نقشه‌ی \lr{stats} با کلیدهای رشته‌ای (\lr{"signal"}، \lr{"decoy"}، \lr{"vacuum"}) نیاز دارد، در حالی که \lr{stats} از \lr{pulse\_type\_index} صحیح (\lr{"0"}، \lr{"1"}، \lr{"2"}) استفاده می‌کند. بنابراین، یک نگاشت انجام می‌شود. \lr{Ma2005} و \lr{Lim2014/Tight} مستقیماً از \lr{stats} استفاده می‌کنند. در پایان، \lr{\texttt{secure\_key\_length}} از \lr{key\_result} استخراج می‌شود.

\subsection{مدیریت خطا و ساخت ردیف CSV خطا}

در صورت بروز خطا در محاسبه‌ی کلید امن یا در جریان اصلی، الگوی خطای یکپارچه (طبق \lr{F-23}) اعمال می‌شود:

\begin{LTR}
	\begin{lstlisting}[language=Python]
		except (LPFailureError, ParameterValidationError, ConfigurationError) as exc:
		qkd_exc = exc
		_log_qkd_exception(qkd_exc, level=logging.WARNING)
		status = f"{qkd_exc.code.value}: {qkd_exc.message[:80]}"
		qkd_exc = QKDSimulationError(
		f"{proof_name} secure-key calculation failed for this simulation point.",
		context={"distance_km": dist, "total_pulses": total_pulses, "combo_idx": combo_idx},
		cause=exc,
		)
		_log_qkd_exception(qkd_exc, level=logging.WARNING)
		status = f"{qkd_exc.code.value}: {qkd_exc.message[:80]}"
	\end{lstlisting}
\end{LTR}

در ابتدا، استثناهای \lr{LPFailureError}، \lr{ParameterValidationError}، \lr{ConfigurationError} گرفته می‌شوند و با \lr{\texttt{\_log\_qkd\_exception}} ثبت می‌گردند. سپس به \lr{\texttt{QKDSimulationError}} تبدیل می‌شوند تا در سطح بالاتر مدیریت شوند. \lr{status} برای ثبت در \lr{CSV} ساخته می‌شود.

در بلاک \lr{except} بیرونی، خطاهای دیگر (\lr{ValueError}، \lr{TypeError} و...) با \lr{\texttt{\_as\_qkd\_exception}} به \lr{QKDException} مناسب تبدیل می‌شوند:

\begin{LTR}
	\begin{lstlisting}[language=Python]
		except (QKDException, ValueError, TypeError, KeyError, RuntimeError, ArithmeticError) as exc:
		orig_exc_type = type(exc).__name__
		orig_exc_msg = str(exc)
		qkd_exc = _as_qkd_exception(
		exc,
		f"Simulation point failed ({orig_exc_type}: {orig_exc_msg[:200]})",
		context={"distance_km": dist, "total_pulses": total_pulses,
			"combo_idx": combo_idx, "original_exception_type": orig_exc_type},
		)
		_log_qkd_exception(qkd_exc, level=logging.ERROR)
	\end{lstlisting}
\end{LTR}

تابع \lr{\texttt{\_as\_qkd\_exception}} (خط ۳۸۶) یک \lr{helper} استاندارد است که استثناهای خام را به \lr{QKDException} مناسب تبدیل می‌کند: \lr{TypeError}/\lr{ValueError}/\lr{KeyError} به \lr{ConfigurationError}، و بقیه به \lr{QKDSimulationError}. این الگو تضمین می‌کند که همه‌ی \lr{catch-block}ها از یک الگوی واحد پیروی کنند.

سپس ردیف خطای \lr{CSV} ساخته می‌شود که همان ستون‌های ردیف موفق را دارد (با مقادیر صفر یا \lr{None}). این الگو (طبق \lr{F-07}) تضمین می‌کند که \lr{CSV} ساختار یکسانی برای همه‌ی ردیف‌ها داشته باشد.

\subsection{تابع run\_and\_save\_csv: هماهنگ‌کننده اصلی}

تابع \lr{\texttt{run\_and\_save\_csv}} (خط ۲۵۸۳) هماهنگ‌کننده‌ی اصلی است. این تابع در نسخه‌ی جدید امضای گسترده‌ای دارد:

\begin{LTR}
	\begin{lstlisting}[language=Python]
		def run_and_save_csv(config, num_workers=None, sweeps=None,
		auto_optimize_pulses=False, dwdm=False,
		wdm_channel_count=4, wdm_channel_power_dbm=None,
		wdm_raman_coefficient=None):
	\end{lstlisting}
\end{LTR}

پارامترهای جدید شامل \lr{auto\_optimize\_pulses}، \lr{dwdm}، \lr{wdm\_channel\_count}، \lr{wdm\_channel\_power\_dbm}، و \lr{wdm\_raman\_coefficient} هستند. این پارامترها به \lr{\texttt{init\_worker}} ارسال می‌شوند.

جریان اصلی شامل این مراحل است: (۱) پیکربندی \lr{logging}؛ (۲) اعتبارسنجی \lr{num\_workers} و بررسی \lr{RAM} (با \lr{psutil}، طبق \lr{F-27})؛ (۳) اعتبارسنجی پیکربندی شبیه‌سازی؛ (۴) ساخت \lr{source\_tmp} و \lr{protocol\_name}؛ (۵) اعتبارسنجی \lr{sweeps}؛ (۶) تولید ترکیب‌های \lr{Cartesian} با \lr{\texttt{itertools.product}} و هشدار در صورت تجاوز از $10000$ ترکیب (طبق \lr{F-15})؛ (۷) تولید \lr{work\_items} با پیمایش (\lr{combo}، \lr{noise}، \lr{pulses}، \lr{dist})؛ (۸) ساخت \lr{fieldnames} \lr{CSV} با ترکیب ستون‌های ثابت و ستون‌های پویا (\lr{sweep\_*})؛ (۹) فراخوانی \lr{\texttt{dispatch\_work\_items}} با \lr{on\_row} \lr{callback}.

\subsection{محاسبۀ QBER اصلاح‌شده و فاصلۀ اطمینان}
\label{sec:csv_row}

در نسخه‌ی جدید، \lr{QBER} اصلاح‌شده و فاصله‌ی اطمینان از \lr{stats} (که ممکن است اصلاحات واکوئم/سیگنال/دکوی اعمال شده باشد) محاسبه می‌شود:

\begin{LTR}
	\begin{lstlisting}[language=Python]
		corrected_total_sifted = sum(tc.sifted for tc in stats.values())
		corrected_total_errors = sum(tc.errors_sifted for tc in stats.values())
		if corrected_total_sifted > 0:
		corrected_qber = corrected_total_errors / corrected_total_sifted
		else:
		corrected_qber = 0.0
		_raw_qber = float(summary.get("qber", 0.0) or 0.0)
		_raw_ci_low = float(summary.get("qber_ci_low", 0.0) or 0.0)
		_raw_ci_high = float(summary.get("qber_ci_high", 0.0) or 0.0)
		_ci_half_width = max(_raw_qber - _raw_ci_low, _raw_ci_high - _raw_qber)
		corrected_ci_low = max(0.0, corrected_qber - _ci_half_width)
		corrected_ci_high = min(1.0, corrected_qber + _ci_half_width)
	\end{lstlisting}
\end{LTR}

در این بلاک، \lr{QBER} از مجموع خطاها تقسیم بر مجموع \lr{sifted} محاسبه می‌شود. نیم‌عرض فاصله‌ی اطمینان از \lr{summary} پروتکل گرفته می‌شود (چون \lr{summary} از \lr{stats} خام محاسبه شده و ممکن است نادرست باشد، اما نیم‌عرض به‌عنوان تقریب عدم‌قطعیت آماری قابل قبول است). سپس فاصله‌ی اطمینان اصلاح‌شده با اعمال این نیم‌عرض روی \lr{QBER} اصلاح‌شده ساخته می‌شود.

ردیف نهایی \lr{CSV} شامل تمام ستون‌های ثابت، \lr{metadata} آشکارساز، \lr{metadata} منبع، ستون‌های \lr{override} آشکارساز، ستون‌های \lr{DWDM}، و ستون‌های \lr{sweep} است. در پایان، \lr{overrides} به‌صورت ستون‌های \lr{sweep\_*} اضافه می‌شوند.

\subsection{استفاده از dispatch\_work\_items}

تابع \lr{\texttt{dispatch\_work\_items}} از ماژول \lr{\texttt{main\_optimized\_dispatch}} وارد شده است. این تابع مسیر تک‌کارگره (حلقه‌ی \lr{for} ترتیبی) و مسیر چندکارگره (\lr{mp.Pool.imap\_unordered}) را مدیریت می‌کند. \lr{on\_row} یک \lr{callback} است که هر ردیف را در \lr{CSV} می‌نویسد و پیشرفت را گزارش می‌دهد. این جداسازی (طبق کامنت در خط ۱۹) امکان \lr{test} \lr{synchronization} و \lr{ordering contract} با \lr{stub workers} را فراهم می‌سازد.

\subsection{تابع build\_pn\_lookup\_table}

تابع \lr{\texttt{build\_pn\_lookup\_table}} (خط ۱۶۷۵) یک جدول \lr{lookup} برای $P(n)$ می‌سازد:

\begin{LTR}
	\begin{lstlisting}[language=Python]
		def build_pn_lookup_table(source, max_n=PN_TRUNCATION_DIM):
		pulse_names = source.pulse_names()
		n_pulses = len(pulse_names)
		table = np.zeros((n_pulses, max_n), dtype=np.float64)
		for i, name in enumerate(pulse_names):
		mu = source.get_mean_photon_number_by_name(name)
		table[i, :] = photon_number_distribution(
		mu, source.statistics_type, max_n=max_n
		)
		return table
	\end{lstlisting}
\end{LTR}

این تابع یک آرایه‌ی \lr{NumPy} به شکل \lr{(n\_pulse\_types, max\_n)} می‌سازد که در آن \lr{table[i, n]} مقدار $P(n)$ برای نوع پالس $i$ را نگه می‌دارد. این جدول می‌تواند توسط حلقه‌های \lr{Numba} یا محاسبات \lr{vectorized} بدون \lr{overhead} شیء \lr{Python} مصرف شود. این طراحی مستقیماً به معیار بازآوری «افشایش آرایه‌های خام \lr{NumPy}» پاسخ می‌دهد.

\subsection{توابع log\_source\_details و log\_detector\_details}

این دو تابع (در خطوط ۷۴۴ و ۹۸۱) برای ثبت اطلاعات تشخیصی در سطح \lr{DEBUG} طراحی شده‌اند. تابع \lr{\texttt{log\_source\_details}} همهٔ ویژگی‌های منبع را با استفاده از متدهای پرس‌وجو (\lr{query})، مانند \lr{\texttt{source.pulse\_names()}} و \lr{\texttt{source.base\_mean\_photon\_numbers()}}، در لاگ ثبت می‌کند. تابع \lr{\texttt{log\_detector\_details}} نیز به‌طور مشابه برای آشکارساز عمل می‌کند و با استفاده از \lr{\texttt{to\_config\_dict()}}، یک نمایهٔ کامل از پیکربندی آن را ثبت می‌نماید. این الگو، مطابق با \lr{F-26}، تضمین می‌کند که مقادیر ثبت‌شده در لاگ با حالت داخلیِ اعتبارسنجی‌شدهٔ سامانه سازگار باشند.


\subsection{توابع کمکی pipeline}

چندین تابع کمکی برای \lr{pipeline} شبیه‌سازی وجود دارند:

\lr{\texttt{build\_channel\_from\_distance}} (خط ۱۰۵۱): یک \lr{\texttt{FiberChannel}} اعتبارسنجی‌شده می‌سازد و در صورت \lr{lossless} بودن، لاگ می‌کند.

\lr{\texttt{build\_protocol}} (خط ۱۰۷۴): یک شیء \lr{Protocol} می‌سازد. این تابع از \lr{dispatch} بر اساس \lr{protocol\_class} استفاده می‌کند و در صورت عدم تطابق، \lr{ConfigurationError} با فهرست کلاس‌های معتبر پرتاب می‌نماید.

\lr{\texttt{get\_pulse\_indices\_from\_prepared\_states}}، \lr{\texttt{generate\_photons\_for\_prepared\_states}}، \lr{\texttt{infer\_ideal\_outcomes\_d0}}، \lr{\texttt{extract\_detection\_arrays}}، \lr{\texttt{detector\_result\_to\_protocol\_detection\_results}}: این توابع برای استخراج و تبدیل داده‌ها بین مراحل مختلف \lr{pipeline} طراحی شده‌اند.

\lr{\texttt{get\_source\_metadata\_row}} (خط ۱۴۸۸) و \lr{\texttt{get\_params\_metadata\_row}} (خط ۱۶۲۰): این توابع برای استخراج \lr{metadata} از اشیای \lr{source} و \lr{params} به ستون‌های \lr{CSV} طراحی شده‌اند. \lr{\texttt{get\_source\_metadata\_row}} از متدهای \lr{query} منبع استفاده می‌کند (نه \lr{to\_dict()}) تا تضمین کند \lr{metadata} با حالت داخلی هماهنگ باشد. \lr{\texttt{get\_params\_metadata\_row}} فیلدهای \lr{DWDM} را استخراج می‌کند و \lr{\texttt{\_compute\_raman\_dark\_rate\_hz}} را برای محاسبه‌ی نرخ \lr{Raman} فراخوانی می‌نماید.

\subsection{الگوی \_\_main\_\_ و CLI}

نقطه‌ی ورود \lr{\texttt{\_\_main\_\_}} در خطوط ۲۷۳۲–۳۰۰۳ پیاده‌سازی شده است. این بلاک از \lr{\texttt{argparse}} با بیش از ۳۰ گزینه استفاده می‌کند. پارامترهای کلیدی عبارت‌اند از: \lr{\texttt{-{}-protocol-class}}، \lr{\texttt{-{}-fiber-loss}}، \lr{\texttt{-{}-min/max-pulses-log}}، \lr{\texttt{-{}-rng-seed}}، \lr{\texttt{-{}-workers}}، پارامترهای آشکارساز (\lr{\texttt{-{}-det-eff-d0}}، \lr{\texttt{-{}-dark-rate}}، \lr{\texttt{-{}-qber-intrinsic}}، ...)، پارامترهای منبع (\lr{\texttt{-{}-statistics-type}}، \lr{\texttt{-{}-error-model}}، \lr{\texttt{-{}-intended-proof}}، ...)، پارامترهای \lr{DWDM} (\lr{\texttt{-{}-dwdm}}، \lr{\texttt{-{}-wdm-channel-count}}، \lr{\texttt{-{}-wdm-channel-power-dbm}}، \lr{\texttt{-{}-wdm-raman-coefficient}})، پارامترهای امنیتی (\lr{\texttt{-{}-eps-sec}}، \lr{\texttt{-{}-eps-cor}}، \lr{\texttt{-{}-eps-pe}}، \lr{\texttt{-{}-eps-smooth}}، \lr{\texttt{-{}-f-error-correction}})، و \lr{\texttt{-{}-auto-optimize-pulses}}.

در صورت \lr{KeyboardInterrupt}، \lr{\texttt{SimulationInterruptedError}} ساخته می‌شود و با کد خروج $130$ (کد استاندارد \lr{POSIX} برای \lr{SIGINT}) خارج می‌شود. در صورت \lr{QKDException}، خطا با کد خروج $1$ خارج می‌شود. در صورت خطاهای دیگر، \lr{\texttt{\_as\_qkd\_exception}} اعمال می‌شود و سپس با کد خروج $1$ خارج می‌گردد.

\section{ارسال و دریافت با سایر ماژول‌ها}

\subsection{ورودی‌های اصلی ماژول}

ورودی‌های اصلی این ماژول عبارت‌اند از: (۱) \lr{\texttt{DEFAULT\_CONFIG}} به‌عنوان پیکربندی پایه؛ (۲) آرگومان‌های \lr{CLI} که روی پیکربندی پایه اعمال می‌شوند؛ (۳) \lr{\texttt{-{}-sweep-json}} به‌عنوان \lr{JSON} تعریف پیمایش؛ (۴) پارامترهای \lr{DWDM} و \lr{Raman} از \lr{CLI}.

\subsection{خروجی‌های اصلی ماژول}

خروجی‌های اصلی عبارت‌اند از: (۱) فایل \lr{\texttt{qkd\_results\_optimized\_sweep.csv}} با ساختار ستون‌های ثابت؛ (۲) لاگ‌های \lr{stdout}/\lr{stderr} برای پیشرفت و خطاها؛ (۳) کد خروج \lr{POSIX} برای \lr{CI/CD}.

\subsection{تعامل با ماژول qkd.sources}

این ماژول از \lr{\texttt{qkd.sources}} کلاس‌های \lr{OpticalSource}، \lr{PoissonSource}، \lr{DensityMatrixSource} و توابع \lr{\texttt{photon\_number\_distribution}}، \lr{\texttt{poisson\_pn\_array}}، \lr{\texttt{thermal\_pn\_array}}، \lr{\texttt{PN\_TRUNCATION\_DIM}}، \lr{\texttt{SourceSimParams}}، \lr{\texttt{build\_source\_sim\_params}} را وارد می‌کند. تابع \lr{\texttt{build\_source}} از \lr{\texttt{OpticalSource.create}} استفاده می‌کند که \lr{dispatch} بین \lr{PoissonSource} و \lr{DensityMatrixSource} را مدیریت می‌نماید و اعتبارسنجی ماتریس چگالی و \lr{security metadata} را داخلی انجام می‌دهد.

\subsection{تعامل با ماژول qkd.detectors}

این ماژول از \lr{\texttt{qkd.detectors}} کلاس‌های \lr{SinglePhotonDetector}، \lr{AfterpulseModel}، \lr{DeadTimeModel}، \lr{EntanglingDecoder}، \lr{DetectionDiagnostics}، \lr{DetectionResult}، \lr{DetectorSimParams}، \lr{\texttt{build\_detector\_sim\_params}} و \lr{ROGERS\_2007\_POLICY} و توابع تحلیلی \lr{Rogers} را وارد می‌کند. تابع \lr{\texttt{build\_detector}} از \lr{\texttt{from\_config\_dict}} استفاده می‌کند که تمام پارامترها را در یک گام اعتبارسنجی می‌نماید.

\subsection{تعامل با ماژول qkd.channel}

این ماژول از \lr{\texttt{qkd.channel}} کلاس \lr{FiberChannel}، توابع \lr{\texttt{build\_attenuation\_config\_from\_sim\_config}}، \lr{\texttt{transmittance\_scalar}}، \lr{\texttt{transmittance\_array}}، \lr{\texttt{link\_efficiency}}، \lr{\texttt{ChannelSimParams}}، \lr{\texttt{build\_channel\_sim\_params}} و ثابت‌های \lr{MAX\_DISTANCE\_KM}، \lr{MAX\_FIBER\_LOSS\_DB\_KM} را وارد می‌کند. تابع \lr{\texttt{build\_attenuation\_config}} به‌عنوان \lr{thin wrapper} به \lr{\texttt{build\_attenuation\_config\_from\_sim\_config}} تفویض می‌کند تا تمام اعتبارسنجی در یک ماژول متمرکز باشد.

\subsection{تعامل با ماژول qkd.protocols}

این ماژول از \lr{\texttt{qkd.protocols}} کلاس‌های \lr{Protocol}، \lr{BB84DecoyProtocol}، \lr{B92Protocol}، \lr{MDIQKDProtocol}، \lr{RedundantTransmissionProtocol}، \lr{DetectionResults}، \lr{SiftingResults}، \lr{BB84PreparedStates}، \lr{B92PreparedStates}، \lr{MDIPreparedStates} و توابع \lr{\texttt{make\_worker\_rngs}}، \lr{\texttt{sample\_for\_parameter\_estimation}} را وارد می‌کند. تابع \lr{\texttt{build\_protocol}} بر اساس \lr{protocol\_class} کلاس مناسب را نمونه‌سازی می‌کند.

\subsection{تعامل با ماژول‌های اثبات امنیتی}

این ماژول با چهار ماژول اثبات در تعامل است: \lr{\texttt{qkd.proofs.lim2014}} (\lr{Lim2014Proof})، \lr{\texttt{qkd.proofs.ma2005}} (\lr{Ma2005VacuumWeakProof})، \lr{\texttt{qkd.proofs.wang2005}} (\lr{Wang2005Proof})، \lr{\texttt{qkd.proofs.tight}} (\lr{BB84TightProof}) و \lr{\texttt{qkd.proofs.optimization}} (\lr{optimize\_pulses\_for\_distance}). \lr{import} به‌صورت \lr{lazy} (داخل \lr{if}) انجام می‌شود. مدیریت \lr{API} متفاوت این کلاس‌ها در بخش پیاده‌سازی بررسی شد.

\subsection{تعامل با ماژول main\_optimized\_dispatch}

این ماژول تابع \lr{\texttt{dispatch\_work\_items}} را از \lr{\texttt{main\_optimized\_dispatch}} وارد می‌کند. این جداسازی امکان \lr{test} \lr{synchronization} و \lr{ordering contract} را با \lr{stub workers} فراهم می‌سازد، بدون نیاز به بارگذاری \lr{qkd.*} در زمان \lr{import}.

\subsection{تعامل با qkd.params، qkd.io، qkd.exceptions و qkd.datatypes}

این ماژول از \lr{\texttt{qkd.params}} توابع \lr{\texttt{load\_lim2014\_dedicated\_params}} و \lr{\texttt{load\_lim2014\_dwdm\_params}} را وارد می‌کند. از \lr{\texttt{qkd.io}} توابع \lr{\texttt{safe\_json\_dumps}}، \lr{\texttt{parse\_json\_strict}}، \lr{\texttt{open\_atomic\_text}} را وارد می‌نماید. از \lr{\texttt{qkd.exceptions}} سلسله‌مراتب \lr{QKDException} را وارد می‌کند. از \lr{\texttt{qkd.datatypes}} تمام \lr{dataclass}ها و \lr{Enum}ها را وارد می‌نماید.

\subsection{مدیریت خطا و تبدیل استثنا}

الگوی یکپارچه‌ی خطا با \lr{\texttt{\_as\_qkd\_exception}} و \lr{\texttt{\_log\_qkd\_exception}} (طبق \lr{F-23}) در تمام \lr{catch-block}ها اعمال می‌شود. این الگو تضمین می‌کند که خطاها به‌صورت یکپارچه ثبت شوند و \lr{CSV} ساختار یکسانی برای همه‌ی ردیف‌ها داشته باشد.

\subsection{مدیریت KeyboardInterrupt}

در صورت \lr{KeyboardInterrupt}، \lr{\texttt{SimulationInterruptedError}} ساخته می‌شود و با کد خروج $130$ خارج می‌شود. این کد، استاندارد \lr{POSIX} برای \lr{SIGINT} است و به \lr{shell} اجازه می‌دهد تشخیص دهد که کاربر اجرا را لغو کرده است.

\subsection{نقش ماژول در اکوسیستم کلی}

در اکوسیستم کلی فریم‌ورک، \lr{\texttt{main\_optimized.py}} به‌عنوان \lr{entry point} اصلی عمل می‌کند و تمام ماژول‌های تخصصی (\lr{sources}، \lr{detectors}، \lr{channel}، \lr{protocols}، \lr{proofs}، \lr{params}، \lr{io}، \lr{exceptions}، \lr{datatypes}) را در یک \lr{pipeline} یکپارچه هماهنگ می‌نماید. خروجی \lr{CSV} این ماژول، ورودی ابزارهای تحلیل بعدی (\lr{pandas}، \lr{R}، \lr{matplotlib}) است. این معماری، جداسازی واضحی بین لایه‌ی هماهنگ‌کننده و لایه‌های تخصصی برقرار می‌سازد و امکان توسعه‌ی مستقل هر ماژول را فراهم می‌نماید.

	
	\chapter{لایه شبکه QKD و مسیریابی}
	\label{ch:network}
	\sloppy
\emergencystretch=5em
\hbadness=10000
\vbadness=10000
\tolerance=2000

\providecolor{codebg}{RGB}{248,248,248}
\providecolor{codeframe}{RGB}{200,200,200}
\providecolor{keyword}{RGB}{0,0,180}
\providecolor{string}{RGB}{163,21,21}
\providecolor{comment}{RGB}{0,128,0}

\lstset{
	backgroundcolor=\color{codebg},
	frame=single,
	rulecolor=\color{codeframe},
	language=Python,
	basicstyle=\ttfamily\scriptsize,
	keywordstyle=\color{keyword}\bfseries,
	stringstyle=\color{string},
	commentstyle=\color{comment}\itshape,
	showstringspaces=false,
	breaklines=true,
	breakatwhitespace=true,
	numbers=left,
	numberstyle=\tiny\color{gray},
	numbersep=8pt,
	tabsize=4,
	literate=
	{ف}{{{\bfseries\itshape ف}}}1
	{ی}{{{\bfseries\itshape ی}}}1
}

\section{نمای کلی و هدف ماژول}

ماژول \lr{\texttt{qkd\_network\_7nodes.py}} پیاده‌سازی کامل یک شبکه توزیع کلید کوانتومی (\lr{Quantum Key Distribution}) با هفت گره است که در آن چهار گره ارسال‌کننده \lr{D}، \lr{E}، \lr{F} و \lr{G} می‌توانند از طریق دو گره میانی \lr{B} و \lr{C} که نقش \lr{Trusted Relay} را ایفا می‌کنند، با گره گیرنده مرکزی \lr{A} کلید مشترک تولید کنند. این ماژول به‌عنوان بالاترین لایه انتزاعی در فریم‌ورک شبیه‌سازی \lr{QKD} عمل می‌کند؛ یعنی به جای آنکه صرفاً یک نشست منفرد توزیع کلید را روی یک لینک نقطه‌به‌نقطه اجرا کند، یک \textbf{شبکه کامل} متشکل از ده لینک، بیست استخر کلید و سه سوییچ نوری را مدیریت می‌کند و نشست‌های متعدد \lr{QKD} را به‌صورت هماهنگ اجرا، انبار و توزیع می‌نماید. هدف اصلی این ماژول، ارائه یک محیط مرجع برای مطالعه مسائل مربوط به مدیریت کلید در سطح شبکه است؛ مسائلی نظیر انتخاب مسیر بهینه، \lr{\texttt{failover}}، \lr{replenishment} استخر کلید، و هماهنگ‌سازی نشست‌ها میان گره‌های مختلف.

ساختار توپولوژی شبکه به‌صورت یک درخت دوبخشی است که در رأس آن گره \lr{A} به‌عنوان \lr{Receiver} قرار دارد و از طریق دو لینک مستقیم \lr{A-B} و \lr{A-C} به دو رله قابل‌اعتماد متصل می‌شود. هر رله سپس از طریق چهار لینک جداگانه به ارسال‌کننده‌های \lr{D}، \lr{E}، \lr{F} و \lr{G} متصل است؛ بنابراین مجموع لینک‌ها برابر با $2 + 2 \times 4 = 10$ لینک خواهد بود. برای هر لینک، دو استخر کلید جداگانه در نظر گرفته می‌شود — یکی در هر انتهای لینک — که در مجموع $10 \times 2 = 20$ استخر کلید را تشکیل می‌دهد. دلیل این طراحی آن است که هر لینک \lr{QKD} کلید را به‌صورت دوطرفه تولید می‌کند ولی هر جهت نیازمند استخر مستقل است تا مصرف کلید در یک جهت، فشار استخر جهت مقابل را تحت تأثیر قرار ندهد. این جداسازی همچنین امکان \lr{back-pressure} مستقل را در الگوریتم \lr{SPF+} فراهم می‌آورد.

یکی از ویژگی‌های برجستهٔ این ماژول، پشتیبانی از دو لایهٔ شبیه‌سازیِ متمایز است که انتخاب آن‌ها به‌صورت شفاف و در زمان اجرا انجام می‌شود. لایهٔ نخست، \lr{\texttt{PipelineQKDLinkSimulator}}، موتور فیزیکیِ کاملِ \lr{\texttt{main\_optimized.run\_single\_simulation()}} را از طریق بارگذاری پویای ماژول \lr{\texttt{main\_optimized.py}} به کار می‌گیرد؛ این موتور، خط لولهٔ کاملِ \lr{\texttt{Proof}} $\to$ \lr{\texttt{Sifting}} $\to$ \lr{\texttt{Detector}} $\to$ \lr{\texttt{Protocol}} $\to$ \lr{\texttt{Channel}} $\to$ \lr{\texttt{Source}} را به‌همراه همهٔ اصلاحات \lr{\texttt{F-01}} تا \lr{\texttt{F-22}} اجرا می‌کند. لایهٔ دوم، \lr{\texttt{StatisticalQKDLinkSimulator}}، یک مدل تحلیلیِ فاصله‌محور است که در صورت عدم دسترسی به موتور فیزیکی، به‌عنوان \lr{\texttt{fallback}} به کار می‌رود. این معماریِ دولایه تضمین می‌کند که ماژول شبکه حتی در محیط‌هایی که بستهٔ \lr{\texttt{qkd}} در آن‌ها نصب نشده است نیز قابل اجرا باشد؛ در عین حال، در محیط توسعهٔ اصلی، بیشترین دقت فیزیکیِ در دسترس را فراهم می‌کند.



علاوه بر خود شبیه‌سازی لینک، این ماژول چهار مؤلفه الگوریتمی کلیدی را پیاده‌سازی می‌کند: \lr{Key Management System (KMS)} که به‌عنوان ارکستراتور مرکزی عمل می‌کند، الگوریتم \lr{SPF+} که نسخه توسعه‌یافته‌ای از الگوریتم \lr{Dijkstra} برای انتخاب مسیر است و علاوه بر فاصله، پارامترهای فیزیکی نظیر \lr{QBER}، تلفات نوری، نرخ کلید و فشار استخر را در نظر می‌گیرد، مکانیزم \lr{Trusted Relay} مبتنی بر \lr{XOR} که امکان انتقال کلید از سراسر شبکه را بدون نیاز به رمزگشایی در نقاط میانی فراهم می‌سازد، و سامانه مدیریت استخر کلید با مدل \lr{watermark} که به‌صورت پویا استخرها را پر و تخلیه می‌کند. این چهار مؤلفه در هم تنیده شده‌اند تا یک سیستم توزیع کلید واقع‌گرایانه را شبیه‌سازی کنند که در آن محدودیت‌های فیزیکی و عملیاتی همگی در انتخاب مسیر نهایی تأثیرگذارند.

در سطح کاربردی، ماژول با فراخوانی تابع \lr{\texttt{demo()}} در انتهای فایل اجرا می‌شود. این تابع ابتدا شبکه را مقداردهی اولیه می‌کند، سپس با اجرای نشست‌های متعدد \lr{QKD} روی هر لینک، استخرهای کلید را تا سطح \lr{high watermark} پر می‌کند (\lr{priming})، پس از آن چهار درخواست کلید خدمت (\lr{service key}) از طرف ارسال‌کننده‌ها به گیرنده \lr{A} را پردازش می‌کند، یک خرابی لینک را شبیه‌سازی می‌کند و در نهایت فرآیند \lr{replenishment} را برای استخرهای تخلیه‌شده اجرا می‌نماید. خروجی این \lr{demo} یک گزارش متنی کامل از وضعیت توپولوژی، معیارهای هر لینک، سطح هر استخر، گام‌های رمزگذاری/رمزگشایی \lr{XOR}، و مسیرهای جایگزین پس از خرابی است که برای مستندسازی و اعتبارسنجی رفتار شبکه بسیار ارزشمند است.

\section{مبانی تئوری و فیزیکی}

\subsection{مدل کانال فیبر نوری و انتقال}

در شبیه‌سازی لینک‌های \lr{QKD} در این ماژول، پایه‌ای‌ترین کمیت فیزیکی، ضریب انتقال کانال (\lr{channel transmittance}) است که نسبت تعداد فوتون‌های دریافتی به فوتون‌های ارسالی را نشان می‌دهد. برای یک فیبر نوری با طول $d$ کیلومتر و ضریب تلفت $\alpha$ دسی‌بل بر کیلومتر، تلفات کلی کانال برابر با $\alpha \cdot d$ دسی‌بل است. چون مقیاس دسی‌بل لگاریتمی است، تبدیل آن به ضریب انتقال خطی از رابطه زیر پیروی می‌کند:

\begin{equation}
	\eta_{\text{channel}}(d) = 10^{-\frac{\alpha \cdot d}{10}}
\end{equation}

در این رابطه، $d$ بر حسب کیلومتر و $\alpha$ بر حسب \lr{dB/km} است. به‌طور معمول، فیبر تک‌حالت استاندارد در طول‌موج $1550$ نانومتر ضریب تلفت $\alpha \approx 0.2$ \lr{dB/km} دارد که در این ماژول به‌عنوان مقدار پیش‌فرض به کار رفته است. ضریب انتقال کلی کانال شامل بازده آشکارساز $\eta$ نیز می‌شود؛ یعنی:

\begin{equation}
	\eta_{\text{total}}(d) = \eta_{\text{det}} \cdot \eta_{\text{channel}}(d) = \eta_{\text{det}} \cdot 10^{-\frac{\alpha \cdot d}{10}}
\end{equation}

این کمیت نقشی محوری در محاسبه نرخ کلید دارد چرا که احتمال شناسایی هر تک‌فوتون ارسالی مستقیماً با $\eta_{\text{total}}$ متناسب است. در فواصل کوتاه متروپولیتن ($30$ تا $75$ کیلومتر)، $\eta_{\text{channel}}$ از $10^{-0.6} \approx 0.25$ تا $10^{-1.5} \approx 0.032$ متغیر است که نشان‌دهنده افت شدید در فواصل دورتر است.

\subsection{مدل خطای کوانتومی (\lr{QBER})}

نرخ خطای کوانتومی بیت (\lr{Quantum Bit Error Rate}) که با نماد $Q$ نشان داده می‌شود، کسری از بیت‌های \lr{sifted} است که بین دو طرف ناهماهنگ هستند. این کمیت معیار اصلی کیفیت لینک \lr{QKD} است و اگر از یک آستانه بحرانی عبور کند، تولید کلید امن از نظر تئوری غیرممکن می‌شود. در پروتکل \lr{BB84} با منبع تک‌فوتونی آرمانی، این آستانه بحرانی برابر با تقریباً $11\%$ است که در ماژول به‌صورت ثابت \lr{\texttt{QBER\_THRESHOLD = 0.11}} تعریف شده است. مدل \lr{QBER} به‌کار رفته در \lr{StatisticalQKDLinkSimulator} چهار مؤلفه مستقل را با هم جمع می‌کند:

\begin{equation}
	Q = Q_{\text{intrinsic}} + Q_{\text{polarization}}(d) + Q_{\text{dark}}(d) + Q_{\text{afterpulse}}
\end{equation}

که در آن هر جزء به شرح زیر تعریف می‌شود. $Q_{\text{intrinsic}}$ خطای پایه ناشی از نقص‌های اپتیکی نظیر \lr{misalignment} تداخل‌سنج است که در ماژول برابر با $0.015$ در نظر گرفته شده است؛ این مقدار شامل $1\%$ خطای ذاتی و $0.5\%$ \lr{misalignment} است که با مقادیر پیش‌فرض \lr{DEFAULT\_CONFIG} در \lr{main\_optimized.py} هم‌خوانی دارد. $Q_{\text{polarization}}(d)$ رشد خطی خطای ناشی از رانش قطبش به دلیل پدیده \lr{Polarization Mode Dispersion} و دوپسیسم بیرفرینژانس فیبر است:

\begin{equation}
	Q_{\text{polarization}}(d) = \kappa_{\text{pol}} \cdot d, \quad \kappa_{\text{pol}} = 10^{-4} \,\text{km}^{-1}
\end{equation}

ضریب $\kappa_{\text{pol}} = 10^{-4}$ به‌ازای هر کیلومتر، متناظر با یک سیستم \lr{QKD} متروپولیتن با تثبیت فعال قطبش است؛ مقدار قبلی $5 \times 10^{-4}$ به‌ازای هر کیلومتر به‌شدت تهاجمی بود و برای لینک‌های $30$ تا $40$ کیلومتری، \lr{QBER} نهایی را در محدوده $2$ تا $2.5\%$ تولید می‌کرد که ناتوان از اشتباه‌شدنی بودن از مقادیر گزارش‌شده در ادبیات تجربی است. مؤلفه $Q_{\text{dark}}$ ناشی از آشکارسازی‌های کاذب (\lr{dark counts}) است و به‌صورت نسبت نرخ دارک به \lr{gain} کلی محاسبه می‌شود:

\begin{equation}
	Q_{\text{dark}}(d) = \frac{R_{\text{dark}}}{g(d)}, \quad g(d) = \max\left(\mu \cdot \eta_{\text{total}}(d) + R_{\text{dark}},\; 10^{-30}\right)
\end{equation}

این تعریف نشان می‌دهد که با افزایش فاصله و کاهش $\eta_{\text{total}}(d)$، سهم دارک کانت در کل شناسایی‌ها رشد می‌کند و در نهایت \lr{QBER} را به سمت آستانه بحرانی هل می‌دهد. مقدار پیش‌فرض $R_{\text{dark}} = 10^{-6}$ به‌ازای هر پالس، متناظر با آشکارسازهای \lr{InGaAs SPAD} با نرخ دارک $10$ هرتز و نرخ پالس $10$ مگاهرتز است. در نهایت، $Q_{\text{afterpulse}} = 0.001$ سهم ثابت پالس‌های پسینی (\lr{afterpulses}) است که از به دام افتادن حامل‌ها در آشکارساز ناشی می‌شود.

\subsection{آنتروپی باینری و کسر امن کلید}

نرخ کلید امن در پروتکل \lr{decoy-state BB84} به‌صورت مجانبی از رابطه بنیادین زیر پیروی می‌کند که آنتروپی باینری \lr{Shannon} را به کسر اطلاعات مشترک امن تبدیل می‌کند:

\begin{equation}
	R = q \cdot \eta_{\text{total}}(d) \cdot \mu \cdot f_{\text{secret}}(Q)
\end{equation}

\begin{equation}
	f_{\text{secret}}(Q) = \max\left(0,\; 1 - 2 \, h_2(Q)\right)
\end{equation}

در این رابطه، $q = 0.5$ ضریب \lr{sifting} است که نشان می‌دهد در \lr{BB84} تنها نیمی از پالس‌ها پس از مقایسه پایه‌ها \lr{sifted} می‌شوند. $\mu$ میانگین تعداد فوتون در پالس سیگنال است که مقدار پیش‌فرض $\mu = 0.5$ دارد؛ این مقدار بهینه تقریبی برای \lr{decoy-state} در رژیم مجانبی است. تابع $h_2(Q)$ آنتروپی باینری است:

\begin{equation}
	h_2(x) = -x \log_2(x) - (1-x) \log_2(1-x)
\end{equation}

این تابع در $x = 0.5$ به بیشینه مقدار $1$ می‌رسد و در $x = 0$ و $x = 1$ به صفر میل می‌کند. عامل $1 - 2 h_2(Q)$ نشان‌دهنده کسر امن کلید پس از کسر اطلاعات لو رفته به حامل به دلیل وجود \lr{QBER} است؛ ضریب $2$ به این دلیل است که در \lr{BB84}، هر خطای بیت هم به نشت اطلاعات به ایو و هم به نیاز تصحیح خطا مربوط می‌شود. وقتی $Q = 0.11$، مقدار $h_2(0.11) \approx 0.5$ می‌شود و در نتیجه $f_{\text{secret}} = 0$؛ این همان آستانه‌ای است که در آن تولید کلید امن متوقف می‌گردد.

\subsection{مدل استخر کلید و فشار (\lr{Pool Pressure})}

یکی از نوآوری‌های کلیدی ماژول، معرفی مدل \lr{piecewise-linear} برای محاسبه فشار استخر کلید است. این مدل امکان تمایز میان سطوح مختلف کمبود کلید را فراهم می‌کند و به الگوریتم \lr{SPF+} اجازه می‌دهد مسیرهایی را برگزیند که استخر کلید سالم‌تری دارند. تعریف ریاضی فشار استخر به‌صورت زیر است:

\begin{equation}
	p_{\text{pool}}(a) = \begin{cases} 0 & \text{if } a \geq h \\ 0.5 \left(1 - \dfrac{a - \ell}{h - \ell}\right) & \text{if } \ell \leq a < h \\ 0.5 + 0.5 \left(1 - \dfrac{a}{\ell}\right) & \text{if } 0 \leq a < \ell \end{cases}
\end{equation}

در این رابطه، $a$ تعداد بیت‌های موجود در استخر است، $\ell$ (\lr{low watermark}) آستانه پایین و $h$ (\lr{high watermark}) آستانه بالای استخر است. مدل به سه رژیم تقسیم می‌شود: رژیم سالم ($a \geq h$) با فشار صفر، رژیم هشدار ($\ell \leq a < h$) که فشار به‌صورت خطی از صفر تا $0.5$ رشد می‌کند، و رژیم بحرانی ($a < \ell$) که فشار از $0.5$ تا $1$ افزایش می‌یابد. این طراحی دو-بخشی باعث می‌شود که الگوریتم مسیریابی حساسیت متفاوتی نسبت به استخرهای نزدیک به آستانه بحرانی در مقایسه با استخرهای نزدیک به سالم داشته باشد. ارزش عملی این مدل در آن است که اگر دو مسیر هم‌ارز از نظر فیزیکی وجود داشته باشند، مسیری که استخر کلید سالم‌تری دارد به‌طور خودکار انتخاب می‌شود.

\subsection{تابع هزینه \lr{SPF+} و الگوریتم \lr{Dijkstra}}

الگوریتم \lr{SPF+} (\lr{Shortest Path First Plus}) توسعه‌ای از الگوریتم کلاسیک \lr{Dijkstra} است که به جای استفاده از یک معیار منفرد (مانند فاصله یا تعداد جهش)، یک تابع هزینه ترکیبی استفاده می‌کند که پنج مؤلفه مختلف را با هم وزن‌دهی می‌کند. تابع هزینه برای هر یال در گراف شبکه به‌صورت زیر تعریف می‌شود:

\begin{equation}
	C_{\text{link}} = \alpha \cdot \frac{L}{L_{\max}} + \beta \cdot \frac{Q}{Q_{\text{th}}} + \gamma \cdot \frac{d}{d_{\max}} + \delta \cdot p_{\text{pool}} - \epsilon \cdot \frac{R}{R_{\max}}
\end{equation}

در این رابطه، $L$ تلفات نوری لینک بر حسب دسی‌بل، $Q$ مقدار \lr{QBER} فعلی، $d$ فاصله فیزیکی، $p_{\text{pool}}$ فشار استخر کلید، و $R$ نرخ تولید کلید امن است. هر یک از این کمیت‌ها به مقدار بیشینه خود نرمال‌سازی می‌شود تا در محدوده $[0, 1]$ قرار گیرند و سپس با ضرایب وزن $\alpha$، $\beta$، $\gamma$، $\delta$ و $\epsilon$ ترکیب می‌شوند. مقادیر پیش‌فرض این ضرایب در ماژول به‌صورت $\alpha = 1.0$، $\beta = 5.0$، $\gamma = 0.5$، $\delta = 2.0$ و $\epsilon = 5.0$ تعریف شده‌اند.

یک ویژگی مهم تابع هزینه بالا، عبارت منفی $-\epsilon \cdot (R/R_{\max})$ است که به‌عنوان یک \lr{bonus} عمل می‌کند؛ یعنی لینک‌هایی که نرخ کلید بالاتری دارند، هزینه کمتری خواهند داشت و در نتیجه به احتمال بیشتری انتخاب می‌شوند. این طراحی با این واقعیت فیزیکی هم‌خوان است که لینک‌های پربازده‌تر معمولاً استخر کلید پرتری دارند و می‌توانند در آینده نیز بهتر پاسخگو باشند. اما اگر \lr{QBER} از آستانه بحرانی $Q_{\text{th}}$ عبور کند، یک جریمه سنگین اما متناهی به هزینه اضافه می‌شود:

\begin{equation}
	C_{\text{link}} \mathrel{+}= H_{\text{penalty}} \cdot \frac{Q}{Q_{\text{th}}}, \quad \text{if } Q > Q_{\text{th}}
\end{equation}

که در آن $H_{\text{penalty}} = 10^6$ است. دلیل استفاده از جریمه متناهی به جای هزینه بی‌نهایت آن است که اگر همه لینک‌ها به دلیل پیکربندی نامناسب \lr{QBER} بالا داشته باشند، الگوریتم همچنان قادر به یافتن مسیر \emph{کم‌ترین ضرر} باشد. این رفتار به‌صورت عملی در حین توسعه مشاهده شد که در آن پیکربندی نامناسب موتور فیزیکی باعث \lr{QBER} بالا در همه لینک‌ها شد و الگوریتم با هزینه بی‌نهایت قادر به یافتن هیچ مسیری نبود.

الگوریتم \lr{Dijkstra} با این تابع هزینه از طریق صف اولویت \lr{min-heap} پیاده‌سازی می‌شود. پیچیدگی زمانی آن با استفاده از \lr{heapq} برابر با $\mathcal{O}((V + E) \log V)$ است که در آن $V$ تعداد گره‌ها و $E$ تعداد یال‌هاست. در شبکه هفت‌گره‌ای با ده یال، این پیچیدگی ناچیز است و امکان اجرای مکرر \lr{SPF+} در هر درخواست کلید خدمت را بدون تأثیر قابل توجه بر کارایی فراهم می‌کند.

\subsection{پروتکل \lr{Trusted Relay} مبتنی بر \lr{XOR}}

یکی از مفاهیم بنیادین در شبکه‌های \lr{QKD}، انتقال کلید به سراسر شبکه از طریق رله‌های قابل‌اعتماد است. در این ماژول، این انتقال با استفاده از \lr{One-Time Pad} مبتنی بر \lr{XOR} انجام می‌شود که از نظر تئوری امنیت اطلاعات کامل را به شرط آنکه کلیدهای هر جهش کاملاً مخفی و یک‌بار مصرف باشند، فراهم می‌سازد. برای یک مسیر سه‌گرهی مانند $D \to B \to A$، فرآیند انتقال به‌صورت زیر است:

\begin{equation}
	C_1 = K_{\text{service}} \oplus K_{DB}
\end{equation}

\begin{equation}
	C_2 = C_1 \oplus K_{DB} \oplus K_{BA} = K_{\text{service}} \oplus K_{BA}
\end{equation}

\begin{equation}
	K_{\text{service}} = C_2 \oplus K_{BA}
\end{equation}

در معادله اول، فرستنده $D$ کلید خدمت را با کلید جهش $K_{DB}$ (که با رله $B$ مشترک است) \lr{XOR} می‌کند و رمزنگاشت $C_1$ را می‌سازد. در گام دوم، رله $B$ ابتدا $C_1$ را با $K_{DB}$ رمزگشایی می‌کند تا به کلید خدمت برسد، سپس آن را با $K_{BA}$ (کلید مشترک با گیرنده $A$) رمزگذاری می‌کند و $C_2$ را تولید می‌نماید. معادله دوم نشان می‌دهد که از نظر ریاضی، $C_2 = K_{\text{service}} \oplus K_{BA}$ چون دو \lr{XOR} با $K_{DB}$ یکدیگر را خنثی می‌کنند. در گام سوم، گیرنده $A$ با \lr{XOR} کردن $C_2$ با $K_{BA}$ به کلید خدمت اصلی می‌رسد. این پروتکل از نظر تئوری امن است زیرا هر ناظر می‌تواند تنها $C_1$ یا $C_2$ را ببیند، ولی بدون دانستن $K_{DB}$ یا $K_{BA}$ نمی‌تواند اطلاعات کمی از $K_{\text{service}}$ به دست آورد.

شرط حیاتی این پروتکل آن است که رله‌ها \emph{قابل‌اعتماد} باشند؛ یعنی کلید جهش را فاش نکنند. به همین دلیل، این رله‌ها در ماژول به‌صورت کلاس \lr{TrustedRelayNode} تعریف شده‌اند که نقش آن‌ها به‌صورت صریح از گره‌های ارسال‌کننده و گیرنده متمایز است. در شبکه‌های عملیاتی، این رله‌ها معمولاً در سایت‌های قابل‌اعتماد مانند دیتاسنترهای امن قرار دارند و کنترل فیزیکی و منطقی روی آن‌ها اعمال می‌شود. در صورتی که رله قابل‌اعتماد نباشد، باید از پروتکل‌های \lr{Measurement-Device-Independent QKD} یا \lr{Entanglement-Based QKD} استفاده کرد که در آن‌ها اعتماد به رله حذف می‌شود.

\subsection{سوییچ‌های نوری و پیکربندی پویا}

در این ماژول، سه سوییچ نوری مدل شده‌اند که امکان پیکربندی پویای توپولوژی شبکه را فراهم می‌سازند. سوییچ \lr{1x2} در گره \lr{A} به این گره اجازه می‌دهد به‌صورت انتخابی به یکی از دو رله \lr{B} یا \lr{C} متصل شود؛ این سوییچ از نوع \lr{MEMS} یا \lr{electro-optic} است که زمان سوییچ در محدوده میکروثانیه تا میلی‌ثانیه دارد. دو سوییچ \lr{2x4} در گره‌های \lr{B} و \lr{C} امکان اتصال هر رله به یکی از چهار ارسال‌کننده \lr{D}، \lr{E}، \lr{F} یا \lr{G} را فراهم می‌سازند. در عمل، این سوییچ‌ها در زمان اجرای \lr{demo} به‌صورت خودکار پیش از اجرای هر درخواست کلید خدمت پیکربندی می‌شوند تا مسیر انتخاب‌شده توسط \lr{SPF+} در دسترس باشد.

\section{تحلیل معماری و ساختار داده‌ها}

\subsection{شمارش‌ها (\lr{Enums}) و ثابت‌های سراسری}

ماژول با تعریف چهار \lr{Enum} پایه آغاز می‌شود که وضعیت‌های ممکن برای مؤلفه‌های مختلف شبکه را مشخص می‌کنند. این شمارش‌ها به‌جای استفاده از رشته‌های خام، نوع‌داده ایمن و قابل اعتبارسنجی فراهم می‌سازند و از خطاهای املایی جلوگیری می‌کنند.

\begin{LTR}
	\begin{lstlisting}[language=Python]
		class KeyBlockState(Enum):
		READY = "READY"
		RESERVED = "RESERVED"
		USED = "USED"
		EXPIRED = "EXPIRED"
		REVOKED = "REVOKED"
		
		class KeyPoolStatus(Enum):
		ACTIVE = "ACTIVE"
		LOW = "LOW"
		EMPTY = "EMPTY"
		DISABLED = "DISABLED"
		ERROR = "ERROR"
		REPLENISHING = "REPLENISHING"
		
		class LinkStatus(Enum):
		UP = "UP"
		DOWN = "DOWN"
		DEGRADED = "DEGRADED"
		
		class NodeRole(Enum):
		RECEIVER = "RECEIVER"
		RELAY = "RELAY"
		SENDER = "SENDER"
	\end{lstlisting}
\end{LTR}

شمارش \lr{KeyBlockState} چرخه حیات یک بلوک کلید را مشخص می‌کند: \lr{READY} به معنای آماده برای مصرف، \lr{RESERVED} به معنای رزرو شده برای یک درخواست در حال انجام، \lr{USED} به معنای مصرف‌شده، \lr{EXPIRED} به معنای منقضی‌شده پس از یک ساعت، و \lr{REVOKED} به معنای ابطال‌شده به دلیل خرابی. شمارش \lr{KeyPoolStatus} وضعیت سلامت استخر را نشان می‌دهد و شش حالت را پشتیبانی می‌کند؛ دو حالت \lr{DISABLED} و \lr{ERROR} به‌صورت \emph{چسبناک} (\lr{sticky}) عمل می‌کنند به این معنا که تنها مداخله خارجی می‌تواند آن‌ها را تغییر دهد. شمارش \lr{LinkStatus} وضعیت فیزیکی لینک را نشان می‌دهد و شامل حالت \lr{DEGRADED} نیز می‌شود که برای لینک‌هایی با کیفیت پایین اما قابل استفاده به کار می‌رود. شمارش \lr{NodeRole} نقش هر گره را در توپولوژی شبکه مشخص می‌کند و در الگوریتم‌های مسیریابی و رله استفاده می‌شود.

در بخش ثابت‌های سراسری، مهم‌ترین کمیت‌ها عبارت‌اند از: \lr{\texttt{INFINITY\_COST = 1e18}} که به‌عنوان هزینه بی‌نهایت در الگوریتم \lr{Dijkstra} به کار می‌رود و از مقدار عددی بزرگ اما قابل نمایش در \lr{float64} استفاده می‌کند؛ \lr{\texttt{QBER\_THRESHOLD = 0.11}} که آستانه بحرانی پروتکل \lr{BB84} است؛ و \lr{\texttt{HIGH\_PENALTY = 1e6}} که جریمه‌ای متناهی اما سنگین برای لینک‌های با \lr{QBER} بالای آستانه است. علاوه بر این، دو مجموعه ثابت \lr{watermark} برای دو نوع شبیه‌ساز تعریف شده است: مجموعه \lr{STAT} برای شبیه‌ساز آماری با مقادیر $1$ مگابیت، $10$ مگابیت و $20$ مگابیت، و مجموعه \lr{PIPELINE} برای شبیه‌ساز خط لوله با مقادیر $10$ کیلوبیت، $100$ کیلوبیت و $200$ کیلوبیت. این تفکیک لازم است چرا که خروجی دو شبیه‌ساز در مقیاس‌های متفاوتی قرار دارد؛ شبیه‌ساز آماری در حدود $0.01$ تا $0.04$ بیت بر پالس تولید می‌کند در حالی که خط لوله واقعی در حدود $10^{-4}$ تا $3 \times 10^{-4}$ بیت بر پالس تولید می‌نماید.

\subsection{داده‌کلاس \lr{KeyBlock}}

کلاس \lr{KeyBlock} با استفاده از \lr{\texttt{@dataclass}} پیاده‌سازی شده و واحد بنیادی ذخیره کلید در استخر را تشکیل می‌دهد. هر بلوک شامل یک شناسه یکتا، شناسه لینک تولیدکننده، گره محلی، گره همتا، ماده کلید به‌صورت بایت‌های خام، طول به بیت، وضعیت، زمان ایجاد، زمان انقضا، شناسه نشست \lr{QKD} تولیدکننده، \lr{QBER} حین تولید و نرخ کلید امن است. استفاده از \lr{dataclass} مزایای متعددی دارد: تولید خودکار سازنده، \lr{\_\_repr\_\_} و \lr{\_\_eq\_\_}، پشتیبانی از \lr{type hint}، و امکان \lr{frozen=True} برای نسخه‌های تغییرناپذیر.

\begin{LTR}
	\begin{lstlisting}[language=Python]
		@dataclass
		class KeyBlock:
		key_id: str
		link_id: str
		local_node: str
		peer_node: str
		key_material: bytes
		length_bits: int
		state: KeyBlockState
		created_at: float
		expires_at: Optional[float]
		qkd_session_id: str
		qber: float
		secret_key_rate: float
	\end{lstlisting}
\end{LTR}

طراحی این داده‌کلاس بر اساس اصل \lr{provenance tracking} است؛ یعنی هر بلوک کلید حامل تمام متادیتای لازم برای ردیابی تاریخچه تولید خود است. این متادیتا در موارد زیر کاربرد دارد: در زمان انتخاب مسیر \lr{SPF+}، اگر \lr{QBER} حین تولید بالاتر از حد باشد، بلوک می‌تواند اولویت پایین‌تری داشته باشد؛ در گزارش‌گیری و حسابرسی، شناسه نشست \lr{QKD} امکان پیگیری کلید تا نشست اصلی را فراهم می‌سازد؛ و در زمان انقضا، زمان ایجاد به‌عنوان مبنا استفاده می‌شود. طول پیش‌فرض هر بلوک $256$ بیت است که برای کلیدهای \lr{AES-256} یا \lr{HMAC-256} مناسب است؛ این مقدار به‌صورت ثابت در تابع \lr{run\_qkd\_session} تعریف شده است.

\subsection{کلاس \lr{KeyPool} و مدیریت چرخه حیات}

کلاس \lr{KeyPool} مسئول مدیریت مجموعه‌ای از بلوک‌های کلید برای یک جهت مشخص از یک لینک \lr{QKD} است. این کلاس به‌جای ذخیره ساده بلوک‌ها، شمارش‌گرهای جداگانه برای بیت‌های موجود (\lr{available})، رزروشده (\lr{reserved}) و مصرف‌شده (\lr{used}) را نگه می‌دارد و یک ماشین وضعیت (\lr{state machine}) برای انتقال میان وضعیت‌های استخر را پیاده‌سازی می‌کند. این طراحی باعث می‌شود که پرس‌وجوهای متداول مانند \lr{available\_key\_bits} در زمان $\mathcal{O}(1)$ اجرا شوند، در حالی که عملیات پیچیده‌تر مانند \lr{reserve\_keys} به $\mathcal{O}(n \log n)$ نیاز دارد که در آن $n$ تعداد بلوک‌ها است.

\begin{LTR}
	\begin{lstlisting}[language=Python]
		class KeyPool:
		def __init__(
		self,
		pool_id: str,
		local_node: str,
		peer_node: str,
		link_id: str,
		max_capacity_bits: int = DEFAULT_MAX_CAPACITY_BITS,
		low_watermark_bits: int = DEFAULT_LOW_WATERMARK_BITS,
		high_watermark_bits: int = DEFAULT_HIGH_WATERMARK_BITS,
		) -> None:
		self.pool_id: str = pool_id
		self.local_node: str = local_node
		self.peer_node: str = peer_node
		self.link_id: str = link_id
		self.key_blocks: Dict[str, KeyBlock] = {}
		self.available_key_bits: int = 0
		self.reserved_key_bits: int = 0
		self.used_key_bits: int = 0
		self.max_capacity_bits: int = max_capacity_bits
		self.low_watermark_bits: int = low_watermark_bits
		self.high_watermark_bits: int = high_watermark_bits
		self.status: KeyPoolStatus = KeyPoolStatus.EMPTY
		self.last_updated: float = time.time()
	\end{lstlisting}
\end{LTR}

روش \lr{add\_key\_block} دارای یک محافظ دفاعی (\lr{defensive cap}) است که از سرریز استخر جلوگیری می‌کند. اگر افزودن یک بلوک جدید باعث شود مجموع بیت‌های موجود و رزرو شده از ظرفیت بیشینه تجاوز کند، بلوک به‌صورت خاموش رد می‌شود. این محافظ در زمان توسعه به‌دلیل یک باگ مشاهده شد که در آن شبیه‌ساز آماری با بازده آشکارساز $\eta = 0.65$ در برابر استخرهای با مقیاس خط لوله $200$ کیلوبیتی اجرا می‌شد و میلیون‌ها بیت در هر نشست تولید می‌کرد که استخر را به‌سرعت پر می‌کرد و حافظه را تخلیه می‌نمود.

یکی از پیچیده‌ترین روش‌های این کلاس، \lr{reserve\_keys} است که بلوک‌ها را به‌صورت نزولی بر اساس اندازه مرتب می‌کند تا تعداد بلوک‌های مصرف‌شده کمینه شود و قطعه‌بندی (\lr{fragmentation}) کاهش یابد. اگر بلوکی بزرگ‌تر از مقدار مورد نیاز باشد، روش \lr{\_split\_block\_for\_reservation} آن را به دو بخش تقسیم می‌کند: یک بلوک رزروشده با اندازه دقیق مورد نیاز و یک بلوک آماده با باقی‌مانده. این رفتار تضمین می‌کند که \lr{available\_key\_bits} دقیقاً به اندازه \lr{length\_bits} کاهش یابد.

\begin{LTR}
	\begin{lstlisting}[language=Python]
		ready_blocks = sorted(
		[b for b in self.key_blocks.values() if b.state == KeyBlockState.READY],
		key=lambda b: -b.length_bits,
		)
		
		for block in ready_blocks:
		if accumulated >= length_bits:
		break
		needed = length_bits - accumulated
		if block.length_bits <= needed:
		block.state = KeyBlockState.RESERVED
		self.available_key_bits -= block.length_bits
		self.reserved_key_bits += block.length_bits
		accumulated += block.length_bits
		reserved.append(block)
		else:
		self._split_block_for_reservation(block, needed, reserved)
		accumulated += needed
	\end{lstlisting}
\end{LTR}

دلیل اهمیت مرتب‌سازی نزولی در این روش آن است که \lr{dict} در پایتون به ترتیب درج را حفظ می‌کند. اگر دو نشست آماده‌سازی هر کدام یک بلوک کوچک \lr{remaining\_bits} (مثلاً $32$ بیت) باقی بگذارند و این بلوک کوچک قبل از بلوک $256$ بیتی بعدی درج شود، الگوریتم اصلی (با ترتیب درج) ابتدا بلوک $32$ بیتی را برمی‌داشت (انباشته $= 32$)، سپس بلوک $256$ بیتی بعدی را (انباشته $= 288$)؛ این یعنی برای یک درخواست $256$ بیتی، $32$ بیت اضافه رزرو می‌شد. در چهار درخواست کلید خدمت، این رفتار به $1062$ بیت مصرف به جای $1024$ بیت مورد انتظار منجر می‌شد. با مرتب‌سازی نزولی، الگوریتم ترجیحاً یک بلوک $256$ بیتی منفرد را برمی‌دارد (تطبیق دقیق) به جای ترکیب بلوک کوچک با بلوک بزرگ.

روش \lr{\_split\_block\_for\_reservation} پیاده‌سازی دقیق تقسیم را با \lr{byte alignment} انجام می‌دهد. \lr{needed\_bits} به نزدیک‌ترین مضرب $8$ گرد می‌شود تا برش بایت‌ها دقیق باشد؛ این کار ممکن است در موارد نادر تا $7$ بیت \lr{over-reservation} ایجاد کند. در موارد لبه‌ای که گرد کردن از اندازه بلوک تجاوز کند، مقدار به اندازه بلوک محدود می‌شود.

\begin{LTR}
	\begin{lstlisting}[language=Python]
		needed_aligned = (needed_bits + 7) // 8 * 8
		if needed_aligned > block.length_bits:
		needed_aligned = block.length_bits
		remainder_bits = block.length_bits - needed_aligned
		needed_bytes = needed_aligned // 8
		
		del self.key_blocks[block.key_id]
		
		taken_id = f"{block.key_id}#R-{uuid.uuid4().hex[:6]}"
		taken_block = KeyBlock(
		key_id=taken_id,
		link_id=block.link_id,
		local_node=block.local_node,
		peer_node=block.peer_node,
		key_material=block.key_material[:needed_bytes],
		length_bits=needed_aligned,
		state=KeyBlockState.RESERVED,
		...
		)
	\end{lstlisting}
\end{LTR}

روش \lr{calculate\_pressure} مدل \lr{piecewise-linear} توضیح‌داده‌شده در بخش مبانی تئوری را پیاده‌سازی می‌کند. این روش در هر بار فراخوانی \lr{SPF+} برای هر دو استخر هر یال اجرا می‌شود و حداکثر آن دو به‌عنوان فشار یال در نظر گرفته می‌شود. این به این معناست که اگر یک استخر کاملاً پر و دیگری خالی باشد، فشار کل برابر $1$ خواهد بود — رفتاری منطقی چرا که جهت برگشت نیز به کلید نیاز دارد.

روش \lr{check\_status} ماشین وضعیت استخر را مدیریت می‌کند. قواعد انتقال آن به‌صورت زیر است: وضعیت‌های \lr{DISABLED} و \lr{ERROR} چسبناک هستند و تنها مداخله خارجی می‌تواند آن‌ها را تغییر دهد؛ \lr{REPLENISHING} به \lr{ACTIVE} تبدیل می‌شود وقتی استخر به \lr{low watermark} برسد (بدون نیاز به \lr{high watermark}، چرا که در حالت خط لوله یک نشست ممکن است تا \lr{high watermark} پر نکند)؛ \lr{EMPTY} به \lr{LOW} تبدیل می‌شود با افزودن هر بیت؛ و \lr{LOW} به \lr{ACTIVE} با رسیدن به \lr{low watermark}.

\subsection{کلاس \lr{QKDLink} و تابع هزینه}

کلاس \lr{QKDLink} یک لینک فیزیکی بین دو گره را مدل می‌کند و شامل پارامترهای فیزیکی (فاصله، تلفات فیبر) و معیارهای عملیاتی (\lr{QBER}، نرخ کلید، تلفات کلی) است. تابع \lr{get\_cost} این کلاس تابع هزینه \lr{SPF+} را پیاده‌سازی می‌کند که در بخش تئوری به‌صورت کامل توضیح داده شد.

\begin{LTR}
	\begin{lstlisting}[language=Python]
		class QKDLink:
		def __init__(
		self,
		link_id: str,
		node_a: str,
		node_b: str,
		distance_km: float,
		fiber_loss_db_km: float = 0.2,
		) -> None:
		self.link_id: str = link_id
		self.node_a: str = node_a
		self.node_b: str = node_b
		self.distance_km: float = distance_km
		self.fiber_loss_db_km: float = fiber_loss_db_km
		self.status: LinkStatus = LinkStatus.UP
		self.qber: float = 0.0
		self.secret_key_rate: float = 0.0
		self.optical_loss_db: float = fiber_loss_db_km * distance_km
		self.last_qkd_session_id: str = ""
	\end{lstlisting}
\end{LTR}

در سازنده این کلاس، تلفات کلی لینک به‌صورت $L = \alpha \cdot d$ محاسبه و در ویژگی \lr{optical\_loss\_db} ذخیره می‌شود. این مقدار در تابع \lr{get\_cost} به‌عنوان یکی از پنج مؤلفه هزینه استفاده می‌شود. تابع \lr{get\_cost} ابتدا بررسی می‌کند که آیا لینک از نظر فیزیکی قطع است؛ اگر چنین باشد، هزینه \lr{INFINITY\_COST} برمی‌گرداند که در الگوریتم \lr{Dijkstra} به‌معنای عدم استفاده از این یال است. در غیر این صورت، پنج مؤلفه هزینه به‌صورت نرمال‌شده محاسبه و با ضرایب وزن ترکیب می‌شوند.

\begin{LTR}
	\begin{lstlisting}[language=Python]
		def get_cost(self, spf_params: Dict[str, float]) -> float:
		if self.status == LinkStatus.DOWN:
		return INFINITY_COST
		
		alpha = spf_params.get("alpha", DEFAULT_ALPHA)
		beta = spf_params.get("beta", DEFAULT_BETA)
		gamma = spf_params.get("gamma", DEFAULT_GAMMA)
		delta = spf_params.get("delta", DEFAULT_DELTA)
		epsilon = spf_params.get("epsilon", DEFAULT_EPSILON)
		pool_pressure = spf_params.get("pool_pressure", 0.0)
		key_rate = spf_params.get("key_rate", 0.0)
		
		norm_loss = self.optical_loss_db / SPF_MAX_LOSS_DB
		norm_qber = self.qber / QBER_THRESHOLD
		norm_distance = self.distance_km / SPF_MAX_DISTANCE_KM
		norm_key_rate = key_rate / SPF_MAX_KEY_RATE
		
		cost = (
		alpha * norm_loss
		+ beta * norm_qber
		+ gamma * norm_distance
		+ delta * pool_pressure
		- epsilon * norm_key_rate
		)
		
		if self.qber > QBER_THRESHOLD:
		qber_excess_ratio = self.qber / QBER_THRESHOLD
		cost += HIGH_PENALTY * qber_excess_ratio
		
		return max(cost, 0.01)
	\end{lstlisting}
\end{LTR}

خط نهایی \lr{return max(cost, 0.01)} یک کف مثبت برای هزینه تعیین می‌کند تا از بروز هزینه صفر یا منفی جلوگیری شود. این کف ضروری است چرا که عبارت منفی $-\epsilon \cdot (R/R_{\max})$ ممکن است در لینک‌های با نرخ کلید بالا، هزینه کل را منفی کند که در الگوریتم \lr{Dijkstra} باعث رفتار غیرتعریف‌شده می‌شود.

\subsection{کلاس \lr{OpticalSwitch}}

کلاس \lr{OpticalSwitch} یک سوییچ نوری را مدل می‌کند که امکان پیکربندی پویای توپولوژی شبکه را فراهم می‌سازد. این کلاس بسیار ساده است و تنها شامل یک نگاشت از پورت ورودی به پورت خروجی و یک پرچم \lr{available} است.

\begin{LTR}
	\begin{lstlisting}[language=Python]
		class OpticalSwitch:
		def __init__(self, switch_id: str, switch_type: str) -> None:
		self.switch_id: str = switch_id
		self.switch_type: str = switch_type
		self.configuration: Dict[str, str] = {}
		self.available: bool = True
		
		def configure(self, mapping: Dict[str, str]) -> None:
		self.configuration = dict(mapping)
		
		def get_connection(self, input_port: str) -> str:
		return self.configuration.get(input_port, "")
	\end{lstlisting}
\end{LTR}

در پیاده‌سازی واقعی، سه سوییچ ایجاد می‌شود: \lr{SW-A-1x2} در گره \lr{A} با نوع \lr{1x2} که به این گره اجازه می‌دهد به یکی از دو رله \lr{B} یا \lr{C} متصل شود، \lr{SW-B-2x4} در گره \lr{B} با نوع \lr{2x4} برای اتصال به چهار ارسال‌کننده، و \lr{SW-C-2x4} در گره \lr{C} با همان نوع. این سوییچ‌ها در زمان اجرای هر درخواست کلید خدمت توسط روش \lr{configure\_switches\_for\_path} در \lr{KMS} پیکربندی می‌شوند تا مسیر انتخاب‌شده توسط \lr{SPF+} در دسترس باشد.

\subsection{سلسله‌مراتب کلاس‌های گره}

گروه کلاس‌های گره از یک کلاس پایه \lr{QKDNode} مشتق می‌شوند که سه زیرکلاس تخصصی دارد: \lr{ReceiverNode}، \lr{TrustedRelayNode} و \lr{SenderNode}. این طراحی از الگوی \lr{Template Method} و \lr{Role-Based Specialization} پیروی می‌کند که در آن هر نقش رفتار خاص خود را دارد ولی زیرساخت مشترک (مانند مدیریت استخر کلید و فهرست همسایه‌ها) در کلاس پایه قرار می‌گیرد.

\begin{LTR}
	\begin{lstlisting}[language=Python]
		class QKDNode:
		def __init__(self, node_id: str, role: NodeRole) -> None:
		self.node_id: str = node_id
		self.role: NodeRole = role
		self.key_pools: Dict[str, KeyPool] = {}
		self.neighbors: List[str] = []
		
		def get_key_pool(self, peer_id: str) -> Optional[KeyPool]:
		return self.key_pools.get(f"{self.node_id}-{peer_id}")
		
		def add_neighbor(self, peer_id: str) -> None:
		if peer_id not in self.neighbors:
		self.neighbors.append(peer_id)
		
		
		class TrustedRelayNode(QKDNode):
		def __init__(self, node_id: str) -> None:
		super().__init__(node_id, NodeRole.RELAY)
		
		def relay_key(
		self,
		incoming_key_material: bytes,
		from_node: str,
		to_node: str,
		hop_keys: Dict[str, bytes],
		) -> bytes:
		key_from = hop_keys[f"{from_node}-{self.node_id}"]
		decrypted = bytes(a ^ b for a, b in zip(incoming_key_material, key_from))
		
		key_to = hop_keys[f"{self.node_id}-{to_node}"]
		encrypted = bytes(a ^ b for a, b in zip(decrypted, key_to))
		
		return encrypted
	\end{lstlisting}
\end{LTR}

روش \lr{relay\_key} در \lr{TrustedRelayNode} پیاده‌سازی پروتکل \lr{XOR} را به‌صورت دو مرحله‌ای انجام می‌دهد: ابتدا رمزگشایی با کلید جهش ورودی، سپس رمزگذاری با کلید جهش خروجی. استفاده از \lr{\texttt{zip}} با \lr{generator expression} به‌صورت \lr{\texttt{bytes(a \^{} b for a, b in zip(...))}} بسیار کارآمد است و از لیست میانی اجتناب می‌کند. در مواردی که طول کلیدها نابرابر باشد، \lr{zip} در کوتاه‌ترین طول متوقف می‌شود که یک رفتار امن است (هرگز سرریز نمی‌کند).

\subsection{شبیه‌سازهای \lr{QKD}}

دو کلاس شبیه‌ساز در این ماژول پیاده‌سازی شده‌اند که هر دو رابط یکسانی (\lr{simulate\_link}) را ارائه می‌دهند اما با پیاده‌سازی متفاوت. \lr{StatisticalQKDLinkSimulator} یک مدل تحلیلی فاصله‌محور است که فرمول‌های استاندارد \lr{decoy-state BB84} را در رژیم مجانبی استفاده می‌کند. \lr{PipelineQKDLinkSimulator} به موتور فیزیکی کامل \lr{main\_optimized.py} وابسته است و خط لوله کامل را با تمام اصلاحات \lr{F-01} تا \lr{F-22} اجرا می‌کند.

\begin{LTR}
	\begin{lstlisting}[language=Python]
		class StatisticalQKDLinkSimulator:
		def __init__(
		self,
		det_eff: float = 0.15,
		dark_rate: float = 1e-6,
		fiber_loss: float = 0.2,
		intrinsic_qber: float = 0.015,
		mu_signal: float = 0.5,
		) -> None:
		self.det_eff = det_eff
		self.dark_rate = dark_rate
		self.fiber_loss = fiber_loss
		self.intrinsic_qber = intrinsic_qber
		self.mu_signal = mu_signal
		
		@staticmethod
		def _binary_entropy(x: float) -> float:
		if x <= 0.0 or x >= 1.0:
		return 0.0
		return -x * math.log2(x) - (1.0 - x) * math.log2(1.0 - x)
	\end{lstlisting}
\end{LTR}

مقادیر پیش‌فرض سازنده \lr{StatisticalQKDLinkSimulator} به‌دقت با \lr{DEFAULT\_CONFIG} در \lr{main\_optimized.py} هم‌خوان است: \lr{det\_eff=0.15} متناظر با بازده آشکارساز \lr{D0} و \lr{D1} در خط لوله، \lr{dark\_rate=1e-6} متناظر با نرخ دارک کانت، \lr{intrinsic\_qber=0.015} شامل $1\%$ خطای ذاتی و $0.5\%$ خطای \lr{misalignment}، و \lr{mu\_signal=0.5} متناظر با شدت پالس سیگنال. این هم‌خوانی تضمین می‌کند که در صورت استفاده از \lr{fallback} آماری، نرخ کلید تولیدشده در همان مرتبه بزرگی با خروجی خط لوله باشد.

روش \lr{\_binary\_entropy} یک متد \lr{static} است که به‌صورت مستقل از نمونه کلاس عمل می‌کند. این روش برای محاسبه آنتروپی باینری استفاده می‌شود و دو شرط مرزی برای جلوگیری از \lr{log(0)} دارد: اگر $x \leq 0$ یا $x \geq 1$، مقدار $0$ برمی‌گردد. این رفتار با تعریف حد آنتروپی باینری هم‌خوان است: $\lim_{x \to 0} x \log_2(x) = 0$.

کلاس \lr{PipelineQKDLinkSimulator} بسیار پیچیده‌تر است و شامل منطق بارگذاری پویای ماژول، پیکربندی موتور فیزیکی، و اعتبارسنجی خروجی است. در سازنده آن، ابتدا یک شبیه‌ساز آماری به‌عنوان \lr{fallback} ایجاد می‌شود با همان بازده آشکارساز خط لوله؛ این تضمین می‌کند که اگر خط لوله در یک لینک به‌خصوص شکست بخورد، \lr{fallback} نرخ کلیدی در همان مرتبه بزرگی تولید کند (در غیر این صورت، \lr{fallback} با \lr{det\_eff=0.65} نرخ‌هایی $20$ برابر بزرگ‌تر از خط لوله تولید می‌کرد).

\subsection{ارکستراتور \lr{KMS} و کلاس \lr{QKDNetwork7Nodes}}

کلاس \lr{KMS} (\lr{Key Management System}) ارکستراتور مرکزی ماژول است و مسئول ساخت توپولوژی، اجرای نشست‌های \lr{QKD}، انتخاب مسیر \lr{SPF+}، اجرای رله قابل‌اعتماد، و مدیریت چرخه حیات استخرها است. این کلاس وضعیت کلی شبکه را در سه دیکشنری \lr{nodes}، \lr{links} و \lr{switches} نگه می‌دارد و شامل روش‌هایی برای هر عملیات کلیدی است. کلاس \lr{QKDNetwork7Nodes} یک لایه بالاتر است که چرخه حیات کامل \lr{demo} را مدیریت می‌کند و خروجی قالب‌بندی‌شده برای کاربر تولید می‌نماید.

\section{پیاده‌سازی ریاضی و الگوریتمی}

\subsection{بارگذاری پویای موتور فیزیکی}

یکی از جنبه‌های مهم پیاده‌سازی، بارگذاری پویای ماژول \lr{main\_optimized.py} است که با استفاده از \lr{importlib.util.spec\_from\_file\_location} انجام می‌شود. این روش به ماژول شبکه اجازه می‌دهد بدون توجه به مسیر کاری فعلی، نسخه محلی موتور فیزیکی را بارگذاری کند. این طراحی در برابر تغییر مسیر کاری (\lr{cwd}) مقاوم است و از بارگذاری اشتباه نسخه‌ای از ماژول از مسیر \lr{site-packages} یا دایرکتوری والد جلوگیری می‌کند.

\begin{LTR}
	\begin{lstlisting}[language=Python]
		_PIPELINE_MODULE_NAME = "_qkd_network_main_optimized"
		
		
		def _load_local_pipeline_module():
		framework_dir = Path(__file__).resolve().parent
		module_path = framework_dir / "main_optimized.py"
		if not module_path.is_file():
		raise ImportError(f"Pipeline engine not found at {module_path}")
		
		framework_dir_str = str(framework_dir)
		if not sys.path or sys.path[0] != framework_dir_str:
		sys.path.insert(0, framework_dir_str)
		
		cached_module = sys.modules.get(_PIPELINE_MODULE_NAME)
		if cached_module is not None:
		cached_path = getattr(cached_module, "__file__", None)
		if cached_path and Path(cached_path).resolve() == module_path:
		return cached_module
		sys.modules.pop(_PIPELINE_MODULE_NAME, None)
		
		spec = importlib.util.spec_from_file_location(_PIPELINE_MODULE_NAME, module_path)
		if spec is None or spec.loader is None:
		raise ImportError(f"Could not create an import spec for {module_path}")
		
		module = importlib.util.module_from_spec(spec)
		sys.modules[_PIPELINE_MODULE_NAME] = module
		try:
		spec.loader.exec_module(module)
		except Exception:
		sys.modules.pop(_PIPELINE_MODULE_NAME, None)
		raise
		
		loaded_path = Path(module.__file__).resolve()
		if loaded_path != module_path:
		raise ImportError(
		f"Loaded pipeline engine from {loaded_path}, expected {module_path}"
		)
		return module
	\end{lstlisting}
\end{LTR}

این تابع از چندین تکنیک دفاعی استفاده می‌کند: ابتدا مسیر ماژول را از \lr{\_\_file\_\_} استخراج می‌کند که به مسیر کاری فعلی وابسته نیست؛ سپس مسیر دایرکتوری را به ابتدای \lr{sys.path} اضافه می‌کند تا وابستگی‌های موتور فیزیکی (مانند بسته \lr{qkd}) نیز قابل بارگذاری باشند؛ در ادامه، اگر ماژول قبلاً بارگذاری شده باشد و مسیر آن تغییر نکرده باشد، از نسخه \lr{cached} استفاده می‌کند؛ در غیر این صورت، نسخه قبلی را از \lr{sys.modules} حذف و نسخه جدید را بارگذاری می‌نماید. در نهایت، یک بررسی امنیتی انجام می‌شود که مطمئن می‌شود ماژول بارگذاری‌شده واقعاً از مسیر مورد انتظار است؛ این بررسی از حملات \lr{path traversal} و بارگذاری اشتباه جلوگیری می‌کند.

\subsection{محاسبه نرخ کلید در شبیه‌ساز آماری}

روش \lr{simulate\_link} در \lr{StatisticalQKDLinkSimulator} مدل کامل \lr{QBER} و نرخ کلید را برای یک لینک پیاده‌سازی می‌کند. این روش با دریافت شیء \lr{QKDLink} و تعداد پالس‌ها، سه خروجی تولید می‌کند: تعداد بیت‌های کلید امن، مقدار \lr{QBER} و نرخ کلید بر پالس.

\begin{LTR}
	\begin{lstlisting}[language=Python]
		def simulate_link(
		self, link: QKDLink, num_pulses: int = 10 ** 7
		) -> Tuple[int, float, float]:
		d = link.distance_km
		alpha = self.fiber_loss
		eta = self.det_eff
		mu = self.mu_signal
		q = 0.5  # sifting factor
		
		# Channel transmittance
		transmittance = eta * 10.0 ** (-alpha * d / 10.0)
		
		# Gain (detection probability per pulse) -- linear approximation
		gain_approx = mu * transmittance + self.dark_rate
		gain = max(gain_approx, 1e-30)
		
		# Distance-dependent QBER model
		qber_intrinsic = self.intrinsic_qber
		qber_polarization = 0.0001 * d
		qber_dark = self.dark_rate / gain
		qber_afterpulse = 0.001
		
		qber = qber_intrinsic + qber_polarization + qber_dark + qber_afterpulse
		qber = min(qber, 1.0)
		
		# If QBER exceeds threshold, no secure key is possible
		if qber > QBER_THRESHOLD:
		return 0, qber, 0.0
		
		# Secret fraction (using asymptotic decoy-state analysis)
		h2_qber = self._binary_entropy(qber)
		secret_fraction = max(0.0, 1.0 - 2.0 * h2_qber)
		
		# Key rate per pulse
		key_rate = q * transmittance * mu * secret_fraction
		
		# Total secure key bits
		secure_key_bits = int(key_rate * num_pulses)
		
		return secure_key_bits, qber, key_rate
	\end{lstlisting}
\end{LTR}

تحلیل خط‌به‌خط این روش به شرح زیر است: ابتدا متغیرهای فیزیکی از شیء لینک و پیکربندی شبیه‌ساز استخراج می‌شوند. سپس ضریب انتقال کلی $\eta_{\text{total}} = \eta \cdot 10^{-\alpha d/10}$ محاسبه می‌شود. در گام بعدی، \lr{gain} (احتمال شناسایی به ازای هر پالس) با تقریب خطی $g \approx \mu \cdot \eta_{\text{total}} + R_{\text{dark}}$ محاسبه می‌شود؛ این تقریب برای $\mu \ll 1$ و $\eta_{\text{total}} \ll 1$ بسیار دقیق است. برای جلوگیری از تقسیم بر صفر، \lr{gain} با کف $10^{-30}$ محدود می‌شود.

چهار مؤلفه \lr{QBER} سپس محاسبه و جمع می‌شوند. اگر \lr{QBER} کل از آستانه $0.11$ تجاوز کند، روش به‌سرعت با خروجی صفر برمی‌گردد چرا که تولید کلید امن از نظر تئوری غیرممکن است. در غیر این صورت، آنتروپی باینری $h_2(Q)$ محاسبه و کسر امن $f_{\text{secret}} = \max(0, 1 - 2 h_2(Q))$ به دست می‌آید. نرخ کلید بر پالس از $R = q \cdot \eta_{\text{total}} \cdot \mu \cdot f_{\text{secret}}$ محاسبه و تعداد کل بیت‌های امن با ضرب در تعداد پالس‌ها به دست می‌آید.

\subsection{اجراهای موتور خط لوله}

روش \lr{\_run\_pipeline} در \lr{PipelineQKDLinkSimulator} موتور فیزیکی کامل را فراخوانی می‌کند. این روش با ساخت پیکربندی، مقداردهی اولیه کارگر، تولید بذر تصادفی، و فراخوانی \lr{run\_single\_simulation} کار می‌کند.

\begin{LTR}
	\begin{lstlisting}[language=Python]
		def _run_pipeline(
		self, link: QKDLink, num_pulses: int
		) -> Tuple[int, float, float]:
		m = self._pipeline_modules
		np = m["np"]
		dist = link.distance_km
		
		if type(num_pulses) is not int or num_pulses <= 0:
		raise ValueError("num_pulses must be a positive integer")
		
		cfg = self._make_pipeline_config()
		
		# Initialise the worker state for this config.
		# combo_map {0: {}} means no sweep overrides.
		m["init_worker"](cfg, {0: {}})
		
		# Derive a high-entropy seed for this session.
		# run_single_simulation uses SeedSequence.spawn(5) internally (F-05).
		rng_seed = int(np.random.default_rng().integers(0, 2**63))
		
		# Run the complete pipeline via main_optimized.
		# args_tuple: (distance_km, total_pulses, apply_noise, seed, combo_idx)
		result = m["run_single_simulation"](
		(dist, num_pulses, self.apply_noise, rng_seed, 0)
		)
		
		secure_key = result["secure_key_bits"]
		qber = result["qber"]
		key_rate = result["secure_key_rate"]
		
		# Store full result for callers that want richer diagnostics
		self._last_result = result
		
		return secure_key, qber, key_rate
	\end{lstlisting}
\end{LTR}

در این روش، ابتدا پیکربندی با استفاده از \lr{\_make\_pipeline\_config} ساخته می‌شود که یک کپی عمیق از \lr{DEFAULT\_CONFIG} گرفته و بازده آشکارساز \lr{D0} و \lr{D1} را با مقدار نمونه تنظیم می‌کند. سپس \lr{init\_worker} با یک \lr{combo\_map} خالی فراخوانی می‌شود که به‌معنای عدم اعمال \lr{sweep} است. بذر تصادفی با استفاده از \lr{numpy.random.default\_rng().integers(0, 2**63)} تولید می‌شود که یک بذر $63$ بیتی با آنتروپی بالا فراهم می‌سازد؛ این بذر سپس به \lr{run\_single\_simulation} ارجاع داده می‌شود که در داخل خود از \lr{SeedSequence.spawn(5)} (اصلاح \lr{F-05}) برای تولید جریان‌های مستقل \lr{RNG} استفاده می‌کند. خروجی شامل سه مقدار کلیدی است که از دیکشنری نتیجه استخراج می‌شوند؛ همچنین کل دیکشنری برای اهداف تشخیصی در \lr{\_last\_result} ذخیره می‌شود.

روش \lr{simulate\_link} در \lr{PipelineQKDLinkSimulator} دارای منطق اعتبارسنجی و \lr{fallback} پیچیده‌ای است. اگر موتور خط لوله \lr{QBER} بالاتر از آستانه یا نرخ کلید صفر برگرداند، به‌صورت خودکار به شبیه‌ساز آماری \lr{fallback} می‌کند. این رفتار به این دلیل ضروری است که گاهی موتور فیزیکی به دلیل پیکربندی نامناسب یا شرایط لبه‌ای عددی، نتایج غیرواقعی تولید می‌کند که در آن صورت، مدل آماری اعتبارسنجی‌شده نتایج قابل اتکاتر ارائه می‌دهد.

\begin{LTR}
	\begin{lstlisting}[language=Python]
		if qber > QBER_THRESHOLD or (key_rate == 0.0 and secure_bits == 0):
		if not hasattr(self, '_fallback_warned'):
		self._fallback_warned: set = set()
		if link.link_id not in self._fallback_warned:
		self._fallback_warned.add(link.link_id)
		import logging
		logging.getLogger(__name__).warning(
		"Pipeline returned unrealistic QBER=%.4f / key_rate=%.2e for "
		"link %s (%.0f km). Falling back to statistical model.",
		qber, key_rate, link.link_id, link.distance_km,
		)
		return self.fallback.simulate_link(link, num_pulses)
	\end{lstlisting}
\end{LTR}

برای جلوگیری از هجوم هشدارهای تکراری، یک مجموعه \lr{\_fallback\_warned} نگه داشته می‌شود که شناسه لینک‌های قبلاً هشدار داده‌شده را در خود ذخیره می‌کند. این الگو (\lr{lazy attribute}) به این صورت کار می‌کند که ویژگی تنها در اولین فراخوانی ایجاد می‌شود و در فراخوانی‌های بعدی مجدداً استفاده می‌گردد.

\subsection{الگوریتم \lr{SPF+} و پیاده‌سازی \lr{Dijkstra}}

روش \lr{spf\_plus} در \lr{KMS} الگوریتم \lr{Dijkstra} را با تابع هزینه ترکیبی \lr{SPF+} پیاده‌سازی می‌کند. این روش از \lr{heapq} به‌عنوان صف اولویت استفاده می‌کند و در هر تکرار، گره با کمترین هزینه را از صف خارج و همسایه‌های آن را بررسی می‌کند.

\begin{LTR}
	\begin{lstlisting}[language=Python]
		def spf_plus(self, source: str, destination: str) -> List[str]:
		dist: Dict[str, float] = {nid: INFINITY_COST for nid in self.nodes}
		prev: Dict[str, Optional[str]] = {nid: None for nid in self.nodes}
		dist[source] = 0.0
		visited: set = set()
		heap: List[Tuple[float, str]] = [(0.0, source)]
		
		while heap:
		current_dist, u = heapq.heappop(heap)
		if u in visited:
		continue
		visited.add(u)
		if u == destination:
		break
		
		for v in self.nodes[u].neighbors:
		if v in visited:
		continue
		
		link_id = self._get_link_id(u, v)
		link = self.links.get(link_id)
		if link is None:
		continue
		
		pool_u_v = self.nodes[u].key_pools.get(f"{u}-{v}")
		pool_v_u = self.nodes[v].key_pools.get(f"{v}-{u}")
		
		pressures = []
		for pool in (pool_u_v, pool_v_u):
		if pool is not None:
		pressures.append(pool.calculate_pressure())
		else:
		pressures.append(1.0)
		pool_pressure = max(pressures)
		
		if (pool_u_v and pool_u_v.status in (KeyPoolStatus.ERROR, KeyPoolStatus.DISABLED)) or \
		(pool_v_u and pool_v_u.status in (KeyPoolStatus.ERROR, KeyPoolStatus.DISABLED)):
		edge_cost = INFINITY_COST
		else:
		cost_params = {
			"alpha": self.spf_weights["alpha"],
			"beta": self.spf_weights["beta"],
			"gamma": self.spf_weights["gamma"],
			"delta": self.spf_weights["delta"],
			"epsilon": self.spf_weights["epsilon"],
			"pool_pressure": pool_pressure,
			"key_rate": link.secret_key_rate,
		}
		edge_cost = link.get_cost(cost_params)
		new_dist = current_dist + edge_cost
		
		if new_dist < dist[v]:
		dist[v] = new_dist
		prev[v] = u
		heapq.heappush(heap, (new_dist, v))
		
		if dist[destination] >= INFINITY_COST:
		return []
		
		path: List[str] = []
		current: Optional[str] = destination
		while current is not None:
		path.append(current)
		current = prev[current]
		path.reverse()
		return path
	\end{lstlisting}
\end{LTR}

تحلیل این الگوریتم به شرح زیر است: ابتدا دو دیکشنری \lr{dist} و \lr{prev} با مقادیر اولیه $\infty$ و \lr{None} مقداردهی می‌شوند؛ تنها فاصله گره مبدأ به صفر تنظیم می‌گردد. در حلقه اصلی، گره با کمترین فاصله از \lr{heap} خارج می‌شود؛ اگر قبلاً بازدید شده باشد، نادیده گرفته می‌شود (این ممکن است رخ دهد چون \lr{heap} ممکن است چندین ورودی برای یک گره داشته باشد). پس از بازدید، اگر گره مقصد باشد، حلقه پایان می‌یابد. برای هر همسایه $v$ از گره فعلی $u$، استخر کلید هر دو جهت بررسی می‌شود و فشار کل به‌عنوان حداکثر فشار دو استخر محاسبه می‌شود. اگر هر یک از دو استخر در وضعیت \lr{ERROR} یا \lr{DISABLED} باشد، هزینه یال به $\infty$ تنظیم می‌شود که در عمل معادل حذف آن یال است. در غیر این صورت، تابع \lr{get\_cost} با پارامترهای هزینه فراخوانی می‌شود. اگر فاصله جدید کمتر از فاصله ثبت‌شده باشد، به‌روزرسانی انجام و ورودی جدید به \lr{heap} افزوده می‌شود.

در پایان، بازسازی مسیر از گره مقصد به مبدأ با پیروی از زنجیره \lr{prev} انجام می‌شود و سپس مسیر معکوس می‌گردد تا از مبدأ به مقصد باشد. اگر فاصله نهایی مقصد همچنان $\infty$ باشد، یعنی هیچ مسیری وجود ندارد و لیست خالی برگردانده می‌شود.

\subsection{درخواست کلید خدمت و رله قابل‌اعتماد}

روش \lr{request\_service\_key} در \lr{KMS} فرآیند کامل تولید و توزیع کلید خدمت را مدیریت می‌کند. این روش شامل شش گام است: یافتن مسیر با \lr{SPF+}، رزرو کلیدهای جهش در هر هاپ، تولید ماده کلید خدمت، اجرای رله قابل‌اعتماد، علامت‌گذاری کلیدهای مصرف‌شده، و ثبت کلید خدمت در رجیستری.

\begin{LTR}
	\begin{lstlisting}[language=Python]
		def request_service_key(
		self, sender: str, receiver: str, key_length_bits: int = 256
		) -> dict:
		# Step 1: Find path
		path = self.spf_plus(sender, receiver)
		if not path:
		return {
			"status": "error",
			"message": f"No available path from {sender} to {receiver}",
			...
		}
		
		# Step 2: Reserve hop keys
		hop_reservations: List[Dict[str, Any]] = []
		key_bytes = key_length_bits // 8
		
		for i in range(len(path) - 1):
		u, v = path[i], path[i + 1]
		pool_u = self.nodes[u].key_pools.get(f"{u}-{v}")
		pool_v = self.nodes[v].key_pools.get(f"{v}-{u}")
		
		if pool_u is None or pool_v is None:
		self._rollback_reservations(hop_reservations)
		return {
			"status": "error",
			"message": f"Missing key pool for hop {u}-{v}",
			...
		}
		
		blocks_u = pool_u.reserve_keys(key_length_bits)
		blocks_v = pool_v.reserve_keys(key_length_bits)
		
		if not blocks_u or not blocks_v:
		self._rollback_reservations(hop_reservations)
		for b in blocks_u:
		b.state = KeyBlockState.READY
		pool_u.available_key_bits += b.length_bits
		pool_u.reserved_key_bits -= b.length_bits
		for b in blocks_v:
		b.state = KeyBlockState.READY
		pool_v.available_key_bits += b.length_bits
		pool_v.reserved_key_bits -= b.length_bits
		return {
			"status": "error",
			"message": f"Insufficient key material in pool for hop {u}-{v}",
			...
		}
		
		hop_reservations.append({
			"u": u, "v": v,
			"blocks_u": blocks_u, "blocks_v": blocks_v,
			"pool_u": pool_u, "pool_v": pool_v,
		})
		
		# Step 3: Generate service key material
		service_key_material = os.urandom(key_bytes)
		
		# Step 4: Perform trusted relay
		relay_success = self.perform_trusted_relay(path, service_key_material)
		
		# Step 5: Mark hop keys as USED
		all_used_key_ids: List[str] = []
		for hop in hop_reservations:
		ids_u = [b.key_id for b in hop["blocks_u"]]
		ids_v = [b.key_id for b in hop["blocks_v"]]
		hop["pool_u"].mark_used(ids_u)
		hop["pool_v"].mark_used(ids_v)
		all_used_key_ids.extend(ids_u)
		all_used_key_ids.extend(ids_v)
		
		# Step 6: Register service key
		self._svc_key_counter += 1
		service_key_id = f"SVC-{sender}{receiver}-{self._svc_key_counter:06d}"
		self.service_key_registry[service_key_id] = {
			"service_key_id": service_key_id,
			"sender": sender,
			"receiver": receiver,
			"key_length_bits": key_length_bits,
			"path": path,
			"hop_keys_used": all_used_key_ids,
			"created_at": time.time(),
			"relay_success": relay_success,
		}
		return {"status": "ok", ...}
	\end{lstlisting}
\end{LTR}

نکته کلیدی در این روش، پیاده‌سازی صحیح \lr{rollback} در صورت شکست است. اگر در هر مرحله از رزرو، یک هاپ به اندازه کافی کلید نداشته باشد، تمام رزروهای قبلی باید به حالت \lr{READY} بازگردانده شوند تا استخرها در وضعیت سازگار باقی بمانند. این الگو به‌عنوان \lr{compensating transaction} در سیستم‌های توزیع‌شده شناخته می‌شود و تضمین می‌کند که حتی در صورت شکست، سیستم در وضعیت سازگار باقی بماند.

روش \lr{perform\_trusted\_relay} پروتکل \lr{XOR} را در طول مسیر اجرا می‌کند. این روش با رمزگذاری اولیه در فرستنده آغاز می‌شود، سپس در هر رله میانی، رمزگشایی و رمزگذاری مجدد انجام می‌شود، و در نهایت در گیرنده، رمزگشایی نهایی انجام می‌گردد. یک ویژگی جالب این روش آن است که در پایان، صحت رله را با مقایسه ماده نهایی با ماده اولیه اعتبارسنجی می‌کند:

\begin{LTR}
	\begin{lstlisting}[language=Python]
		def perform_trusted_relay(
		self, path: List[str], service_key_material: bytes
		) -> bool:
		if len(path) < 2:
		return False
		
		current_material = service_key_material
		
		# Sender encrypts with first hop key
		first_hop_u, first_hop_v = path[0], path[1]
		pool_u = self.nodes[first_hop_u].key_pools.get(f"{first_hop_u}-{first_hop_v}")
		if pool_u is None:
		return False
		
		hop_key_send = self._get_reserved_key_material(pool_u)
		if hop_key_send is None:
		return False
		
		# Step 1: Sender encrypts
		current_material = bytes(a ^ b for a, b in zip(current_material, hop_key_send))
		
		# Intermediate relay nodes
		for i in range(1, len(path) - 1):
		relay_node = self.nodes[path[i]]
		from_node = path[i - 1]
		to_node = path[i + 1]
		
		if not isinstance(relay_node, TrustedRelayNode):
		return False
		
		pool_in = relay_node.key_pools.get(f"{relay_node.node_id}-{from_node}")
		if pool_in is None:
		pool_in = relay_node.key_pools.get(f"{from_node}-{relay_node.node_id}")
		pool_out = relay_node.key_pools.get(f"{relay_node.node_id}-{to_node}")
		
		if pool_in is None or pool_out is None:
		return False
		
		key_in = self._get_reserved_key_material(pool_in)
		key_out = self._get_reserved_key_material(pool_out)
		
		if key_in is None or key_out is None:
		return False
		
		hop_keys = {
			f"{from_node}-{relay_node.node_id}": key_in,
			f"{relay_node.node_id}-{to_node}": key_out,
		}
		
		current_material = relay_node.relay_key(
		current_material, from_node, to_node, hop_keys
		)
		
		# Final receiver decrypts with last hop key
		last_hop_u, last_hop_v = path[-2], path[-1]
		pool_recv = self.nodes[last_hop_v].key_pools.get(f"{last_hop_v}-{last_hop_u}")
		if pool_recv is None:
		return False
		
		hop_key_recv = self._get_reserved_key_material(pool_recv)
		if hop_key_recv is None:
		return False
		
		final_material = bytes(a ^ b for a, b in zip(current_material, hop_key_recv))
		
		# Verify relay correctness (final material should match original)
		return final_material == service_key_material
	\end{lstlisting}
\end{LTR}

اعتبارسنجی نهایی \lr{return final\_material == service\_key\_material} یک تضمین سلامت (\lr{integrity check}) است؛ اگر به هر دلیلی (مثلاً عدم تطابق طول کلیدها یا خطای منطقی در رله) ماده نهایی با ماده اولیه تطابق نداشته باشد، روش مقدار \lr{False} برمی‌گرداند و در نتیجه در رجیستری کلید خدمت، \lr{relay\_success=False} ثبت می‌شود. این رفتار به کاربر اجازه می‌دهد شکست‌ها را تشخیص دهد.

\subsection{مدیریت نشست \lr{QKD} و پر کردن استخر}

روش \lr{run\_qkd\_session} در \lr{KMS} یک نشست \lr{QKD} را روی یک لینک اجرا می‌کند و کلیدهای تولیدشده را در استخرهای هر دو گره انتهای لینک توزیع می‌کند. این روش با تقسیم کلید تولیدشده به بلوک‌های $256$ بیتی، بلوک‌های متناظر در هر دو گره با \emph{ماده کلیدی یکسان} می‌سازد.

\begin{LTR}
	\begin{lstlisting}[language=Python]
		def run_qkd_session(
		self, link_id: str, num_pulses: int = 10 ** 7
		) -> Tuple[int, float, float]:
		link = self.links.get(link_id)
		if link is None:
		raise ValueError(f"Unknown link: {link_id}")
		if link.status == LinkStatus.DOWN:
		raise RuntimeError(f"Link {link_id} is DOWN")
		
		secure_bits, qber, key_rate = self.simulator.simulate_link(link, num_pulses)
		
		link.qber = qber
		link.secret_key_rate = key_rate
		session_id = f"QKD-{link_id}-{uuid.uuid4().hex[:8]}"
		link.last_qkd_session_id = session_id
		
		if secure_bits <= 0:
		return secure_bits, qber, key_rate
		
		block_size_bits = 256
		num_blocks = secure_bits // block_size_bits
		remaining_bits = secure_bits % block_size_bits
		
		node_a, node_b = link.node_a, link.node_b
		
		for i in range(num_blocks):
		self._key_counter += 1
		kid = f"{node_a}{node_b}-{self._key_counter:06d}"
		
		block_a = KeyBlock(
		key_id=kid,
		link_id=link_id,
		local_node=node_a,
		peer_node=node_b,
		key_material=os.urandom(block_size_bits // 8),
		length_bits=block_size_bits,
		state=KeyBlockState.READY,
		created_at=time.time(),
		expires_at=time.time() + 3600.0,
		qkd_session_id=session_id,
		qber=qber,
		secret_key_rate=key_rate,
		)
		
		self._key_counter += 1
		kid_b = f"{node_b}{node_a}-{self._key_counter:06d}"
		
		block_b = KeyBlock(
		key_id=kid_b,
		...
		key_material=block_a.key_material,  # Same key at both ends
		...
		)
		
		pool_a = self.nodes[node_a].key_pools.get(f"{node_a}-{node_b}")
		pool_b = self.nodes[node_b].key_pools.get(f"{node_b}-{node_a}")
		if pool_a:
		pool_a.add_key_block(block_a)
		if pool_b:
		pool_b.add_key_block(block_b)
		
		return secure_bits, qber, key_rate
	\end{lstlisting}
\end{LTR}

نکته مهم در این پیاده‌سازی آن است که ماده کلیدی بلوک \lr{block\_b} با \lr{block\_a.key\_material} مقداردهی می‌شود، نه با یک فراخوانی جدید \lr{os.urandom}. این تضمین می‌کند که کلید در هر دو گره یکسان است که برای رمزنگاری \lr{XOR} ضروری است. در عمل، در یک سیستم \lr{QKD} واقعی، این کلید از طریق پروتکل \lr{QKD} به‌صورت مستقل در هر دو گره تولید و سپس از طریق کانال احراز هویت عمومی تطبیق داده می‌شود؛ در این شبیه‌سازی، به‌جای شبیه‌سازی کامل پروتکل تطبیق، مستقیماً کلید یکسان در هر دو گره قرار داده می‌شود.

روش \lr{replenish\_key\_pools} استخرها را اسکن می‌کند و برای استخرهای در وضعیت \lr{LOW} یا \lr{EMPTY}، نشست \lr{QKD} جدیدی اجرا می‌کند. نکته طراحی مهم در این روش آن است که هر دو استخر متناظر در یک لینک به‌صورت همزمان به وضعیت \lr{REPLENISHING} تغییر می‌یابند تا وضعیت همگام حفظ شود.

\begin{LTR}
	\begin{lstlisting}[language=Python]
		def replenish_key_pools(self) -> List[str]:
		replenished: List[str] = []
		visited_links: set = set()
		
		for node_id, node in self.nodes.items():
		for pool_key, pool in node.key_pools.items():
		if pool.status in (KeyPoolStatus.LOW, KeyPoolStatus.EMPTY):
		link_id = pool.link_id
		
		if link_id in visited_links:
		continue
		visited_links.add(link_id)
		
		link = self.links.get(link_id)
		if link is None:
		continue
		pool_a = self.nodes[link.node_a].key_pools.get(
		f"{link.node_a}-{link.node_b}"
		)
		pool_b = self.nodes[link.node_b].key_pools.get(
		f"{link.node_b}-{link.node_a}"
		)
		if pool_a:
		pool_a.status = KeyPoolStatus.REPLENISHING
		if pool_b:
		pool_b.status = KeyPoolStatus.REPLENISHING
		
		try:
		replenish_pulses = 10 ** 7
		self.run_qkd_session(link_id, num_pulses=replenish_pulses)
		replenished.append(link_id)
		except Exception as e:
		if pool_a:
		pool_a.status = KeyPoolStatus.ERROR
		if pool_b:
		pool_b.status = KeyPoolStatus.ERROR
		print(f"  [ERROR] Replenishment failed for {link_id}: {e}")
		
		return replenished
	\end{lstlisting}
\end{LTR}

استفاده از \lr{visited\_links} برای جلوگیری از \lr{double-replenishment} همان لینک است؛ چون هر لینک دو استخر دارد، اسکن بدون این محافظ، همان لینک را دو بار شناسایی و پردازش می‌کرد. در صورت شکست نشست \lr{QKD}، هر دو استخر به وضعیت \lr{ERROR} تغییر می‌یابند که چسبناک است و تنها با مداخله خارجی قابل پاک‌سازی است.

\subsection{کالیبراسیون برای نوع شبیه‌ساز}

روش \lr{calibrate\_for\_simulator} در \lr{KMS} یکی از مهم‌ترین روش‌های تنظیم پویا است. این روش با دریافت یک پرچم بولی \lr{is\_pipeline}، \lr{watermark} استخرها و ثابت نرمال‌سازی \lr{SPF+} را برای نوع شبیه‌ساز تنظیم می‌کند. این کالیبراسیون ضروری است چرا که خروجی دو شبیه‌ساز در مقیاس‌های کاملاً متفاوتی قرار دارد.

\begin{LTR}
	\begin{lstlisting}[language=Python]
		def calibrate_for_simulator(self, is_pipeline: bool) -> None:
		self._is_pipeline_mode = is_pipeline
		
		if is_pipeline:
		new_low = PIPELINE_LOW_WATERMARK_BITS
		new_high = PIPELINE_HIGH_WATERMARK_BITS
		new_max = PIPELINE_MAX_CAPACITY_BITS
		new_max_rate = PIPELINE_SPF_MAX_KEY_RATE
		else:
		new_low = STAT_LOW_WATERMARK_BITS
		new_high = STAT_HIGH_WATERMARK_BITS
		new_max = STAT_MAX_CAPACITY_BITS
		new_max_rate = STAT_SPF_MAX_KEY_RATE
		
		# Update all pools
		for node_id, node in self.nodes.items():
		for pool_key, pool in node.key_pools.items():
		pool.low_watermark_bits = new_low
		pool.high_watermark_bits = new_high
		pool.max_capacity_bits = new_max
		pool.status = pool.check_status()
		
		# Update SPF+ key rate normalization
		global SPF_MAX_KEY_RATE
		SPF_MAX_KEY_RATE = new_max_rate
	\end{lstlisting}
\end{LTR}

اگر این کالیبراسیون انجام نشود، در حالت خط لوله همه استخرها به‌طور دائمی در وضعیت \lr{LOW} یا \lr{REPLENISHING} گیر می‌کنند و فشار استخر برای همه هاپ‌ها به $\sim 1.0$ می‌رسد که باعث می‌شود \lr{SPF+} نتواند میان مسیرها تمایز قائل شود. این رفتار در حین توسعه مشاهده شد و به‌عنوان یک باگ بحرانی شناسایی و اصلاح گردید. تغییر متغیر سراسری \lr{SPF\_MAX\_KEY\_RATE} با کلیدواژه \lr{global} انجام می‌شود که در پایتون برای اصلاح متغیرهای ماژول سطح بالا ضروری است.

\section{ارسال و دریافت با سایر ماژول‌ها}

\subsection{وابستگی به موتور فیزیکی \lr{main\_optimized.py}}

مهم‌ترین وابستگی خارجی این ماژول، به فایل \lr{main\_optimized.py} است که در همان دایرکتوری قرار دارد. این وابستگی از طریق تابع \lr{\_load\_local\_pipeline\_module} برقرار می‌شود که موتور فیزیکی را به‌صورت پویا بارگذاری می‌کند. از موتور فیزیکی، پنج مؤلفه کلیدی استخراج می‌شود: \lr{numpy} به‌عنوان کتابخانه آرایه، \lr{DEFAULT\_CONFIG} به‌عنوان پیکربندی پایه، \lr{init\_worker} برای مقداردهی اولیه کارگر، \lr{run\_single\_simulation} به‌عنوان نقطه ورود اصلی شبیه‌سازی، و \lr{set\_nested\_value} برای اصلاح پیکربندی.

این وابستگی به‌صورت اختیاری (\lr{soft dependency}) طراحی شده است؛ یعنی اگر موتور فیزیکی در دسترس نباشد، ماژول شبکه همچنان با شبیه‌ساز آماری قابل اجرا است. این طراحی \lr{graceful degradation} نامیده می‌شود و تضمین می‌کند که ماژول در محیط‌های مختلف قابل اجرا باشد. در زمان آماده‌سازی شبکه، روش \lr{initialize} در \lr{QKDNetwork7Nodes} ابتدا تلاش می‌کند موتور خط لوله را بارگذاری کند و در صورت شکست، به‌صورت شفاف به شبیه‌ساز آماری روی می‌آورد و پیام مناسبی به کاربر نمایش می‌دهد.

\subsection{ورودی‌ها به ماژول}

ورودی اصلی ماژول از طریق فراخوانی تابع \lr{demo()} در انتهای فایل است که نیازی به پارامتر ندارد. این تابع یک نمونه از \lr{QKDNetwork7Nodes} می‌سازد و چرخه حیات کامل را اجرا می‌کند. پارامترهای قابل تنظیم در طول اجرا عبارت‌اند از: \lr{num\_sessions\_per\_link} در روش \lr{prime\_pools} که تعداد نشست‌های \lr{QKD} برای آماده‌سازی هر لینک را تعیین می‌کند (پیش‌فرض $50$)، و \lr{pulses\_per\_session} که تعداد پالس‌ها در هر نشست است (پیش‌فرض $10^7$). در مورد \lr{pulses\_per\_session}، مقدار قبلی $10^8$ به دلیل فشار حافظه کاهش یافت؛ در $10^8$ پالس، خط لوله چندین آرایه \lr{int64} به حجم حدود $800$ مگابایت (مانند \lr{prepared\_states}، \lr{photons}، \lr{click0}، \lr{click1}) تخصیص می‌داد که باعث فشار حافظه و موارد لبه‌ای عددی در کد \lr{sifting}/\lr{parameter-estimation} بسته \lr{qkd} می‌شد.

علاوه بر این، کلاس \lr{StatisticalQKDLinkSimulator} پارامترهای فیزیکی قابل تنظیمی دارد: \lr{det\_eff} (بازده آشکارساز)، \lr{dark\_rate} (نرخ دارک کانت)، \lr{fiber\_loss} (تلفات فیبر بر کیلومتر)، \lr{intrinsic\_qber} (\lr{QBER} ذاتی) و \lr{mu\_signal} (شدت پالس سیگنال). این پارامترها در سازنده پیش‌فرض‌گذاری شده‌اند تا با \lr{DEFAULT\_CONFIG} در \lr{main\_optimized.py} هم‌خوان باشند. کلاس \lr{PipelineQKDLinkSimulator} نیز دو پارامتر دارد: \lr{detector\_efficiency} که بازده آشکارساز \lr{D0} و \lr{D1} را تنظیم می‌کند، و \lr{apply\_noise} که مشخص می‌کند آیا مسیر کامل آشکارساز نویزدار استفاده شود.

توپولوژی شبکه به‌صورت سخت‌کد در روش \lr{build\_network} تعریف شده است. هفت گره با نقش مشخص، ده لینک با فواصل تعریف‌شده در دیکشنری \lr{NODE\_DISTANCES}، و سه سوییچ نوری با پیکربندی پیش‌فرض ایجاد می‌شوند. فواصل لینک‌ها به‌صورت متروپولیتن در نظر گرفته شده‌اند: لینک \lr{A-C} با $50$ کیلومتر کوتاه‌ترین مسیر مستقیم به رله است، در حالی که \lr{A-B} با $75$ کیلومتر طولانی‌ترین است. لینک‌های رله به ارسال‌کننده از $30$ کیلومتر (\lr{C-D}) تا $70$ کیلومتر (\lr{B-G}) متغیر هستند.

\subsection{خروجی‌ها و گزارش‌گیری}

خروجی اصلی ماژول، نمایش متنی قالب‌بندی‌شده در \lr{stdout} است که شامل چندین بخش است: جدول توپولوژی گره‌ها با نقش و همسایه‌ها، جدول معیارهای لینک شامل فاصله، وضعیت، \lr{QBER}، نرخ کلید و تلفات، جدول سطح استخر با بیت‌های موجود، رزرو شده و مصرف‌شده، پیکربندی سوییچ‌های نوری، جزئیات هر درخواست کلید خدمت شامل مقایسه مسیرهای کاندید، گام‌های \lr{XOR relay}، و شناسه کلید خدمت صادر شده.

علاوه بر خروجی متنی، روش \lr{get\_network\_status} در \lr{KMS} یک دیکشنری ساختار یافته برمی‌گرداند که شامل وضعیت کامل شبکه است. این دیکشنری برای یکپارچه‌سازی با سامانه‌های مانیتورینگ خارجی یا تولید گزارش خودکار مناسب است. ساختار آن به این صورت است:

\begin{LTR}
	\begin{lstlisting}[language=Python]
		{
			"nodes": {
				"A": {
					"role": "RECEIVER",
					"neighbors": ["B", "C"],
					"pools": {
						"A-B": {
							"pool_id": "POOL-A-B",
							"status": "ACTIVE",
							"available_bits": 12345678,
							"reserved_bits": 0,
							"used_bits": 1024,
							"num_blocks": 42
						},
						...
					}
				},
				...
			},
			"links": {
				"A-B": {
					"distance_km": 75.0,
					"status": "UP",
					"qber": 0.023456,
					"secret_key_rate": 0.0123,
					"optical_loss_db": 15.0,
					"last_session": "QKD-A-B-abc12345"
				},
				...
			},
			"switches": {
				"SW-A-1x2": {
					"type": "1x2",
					"available": true,
					"configuration": {"A": "B"}
				},
				...
			},
			"service_keys_issued": 4
		}
	\end{lstlisting}
\end{LTR}

این ساختار تودرتو امکان سریع پیمایش وضعیت شبکه را فراهم می‌سازد و برای ابزارهای گزارش‌گیری خودکار یا داشبورد مانیتورینگ مناسب است.

\subsection{یکپارچه‌سازی با اصلاحات \lr{F-01} تا \lr{F-22}}

هنگامی که موتور خط لوله فعال است، ماژول شبکه به‌صورت خودکار از تمام اصلاحات اعمال‌شده در \lr{main\_optimized.py} بهره می‌برد. این اصلاحات شامل موارد زیر است: \lr{F-01} (تثبیت ولتاژ بایاس آشکارساز و ممیزی \lr{override})، \lr{F-05} (استفاده از \lr{SeedSequence.spawn} برای جریان‌های مستقل \lr{RNG})، \lr{F-11} (اعتبارسنجی تعادل احتمال)، \lr{F-14} (ذخیره‌سازی متادیتای منبع) و \lr{F-18} (حذف کد مرده \lr{debug pulse override}). این اصلاحات به‌صورت شفاف از طریق فراخوانی \lr{run\_single\_simulation} فعال می‌شوند و نیازی به پیکربندی اضافی در ماژول شبکه ندارند.

در مورد \lr{F-05}، نکته مهم آن است که بذر تصادفی در سطح ماژول شبکه با \lr{numpy.random.default\_rng().integers(0, 2**63)} تولید می‌شود و سپس به \lr{run\_single\_simulation} ارسال می‌گردد. در داخل موتور فیزیکی، این بذر با \lr{SeedSequence.spawn(5)} به پنج جریان مستقل تقسیم می‌شود که هر یک برای یک بخش مختلف خط لوله (منبع، کانال، پروتکل، آشکارساز، \lr{sifting}) استفاده می‌شود. این تضمین می‌کند که نشست‌های متوالی روی همان لینک، نتایج مستقل از نظر آماری داشته باشند و همبستگی‌های ناخواسته میان اجزای مختلف خط لوله ایجاد نشود.

\subsection{رفتار در حالت خرابی و بازیابی}

ماژول شبکه شامل منطق پیچیده‌ای برای مدیریت خرابی و بازیابی است. روش \lr{handle\_failure} در \lr{KMS} با دریافت شناسه یک لینک خراب، آن لینک را به وضعیت \lr{DOWN} تبدیل می‌کند، استخرهای متناظر را به وضعیت \lr{ERROR} (چسبناک) تغییر می‌دهد، و سپس برای همه جفت‌های ارسال‌کننده–گیرنده تحت تأثیر، مسیرهای جایگزین را با \lr{SPF+} جستجو می‌کند. این روش مسیرهای جایگزین را به‌صورت لیست رشته‌ای برمی‌گرداند که در خروجی \lr{demo} نمایش داده می‌شوند.

\begin{LTR}
	\begin{lstlisting}[language=Python]
		def handle_failure(self, failed_link_id: str) -> List[str]:
		link = self.links.get(failed_link_id)
		if link is None:
		return []
		
		link.status = LinkStatus.DOWN
		
		node_a, node_b = link.node_a, link.node_b
		pool_a = self.nodes[node_a].key_pools.get(f"{node_a}-{node_b}")
		pool_b = self.nodes[node_b].key_pools.get(f"{node_b}-{node_a}")
		if pool_a:
		pool_a.status = KeyPoolStatus.ERROR
		if pool_b:
		pool_b.status = KeyPoolStatus.ERROR
		
		affected_pairs: List[Tuple[str, str]] = []
		senders = ["D", "E", "F", "G"]
		for s in senders:
		affected_pairs.append((s, "A"))
		
		alternative_paths: List[str] = []
		for sender, receiver in affected_pairs:
		alt_path = self.spf_plus(sender, receiver)
		if alt_path:
		path_str = " -> ".join(alt_path)
		alternative_paths.append(path_str)
		
		return alternative_paths
	\end{lstlisting}
\end{LTR}

در \lr{demo}، پس از شبیه‌سازی خرابی لینک \lr{B-D}، درخواست کلید خدمت \lr{D-A} مجدداً اجرا می‌شود. در این حالت، \lr{SPF+} به‌جای مسیر کوتاه \lr{D-B-A}، مسیر جایگزین \lr{D-C-A} را انتخاب می‌کند چون لینک \lr{B-D} با هزینه $\infty$ عملاً از گراف حذف شده است. این رفتار به‌صورت خودکار توسط الگوریتم \lr{Dijkstra} مدیریت می‌شود و نیازی به منطق خاصی در سطح بالاتر ندارد.

پس از خرابی، روش \lr{run\_demo} لینک \lr{B-D} را به‌صورت دستی به وضعیت \lr{UP} بازمی‌گرداند و استخرهای آن را به وضعیت \lr{LOW} تنظیم می‌کند تا فرآیند \lr{replenishment} آن‌ها را شناسایی و پر کند. این الگو (\lr{manual state injection}) برای شبیه‌سازی سناریوهای بازیابی استفاده می‌شود. در یک سیستم تولیدی واقعی، این تغییرات وضعیت به‌صورت خودکار توسط سیستم‌های مانیتورینگ که فاصله، \lr{QBER} یا سایر معیارها را پایش می‌کنند، اعمال می‌شوند.

\subsection{رفتار حافظه و کارایی}

از نظر حافظه، ماژول شبکه به‌صورت نسبتاً سبک طراحی شده است. هر استخر کلید به‌صورت یک دیکشنری از بلوک‌های کلید پیاده‌سازی شده که هر بلوک شامل یک شیء بایت با طول $32$ بایت (برای $256$ بیت کلید) است. در حالت آماری با \lr{watermark} پیش‌فرض $20$ مگابیت، هر استخر می‌تواند تا حدود $78125$ بلوک $256$ بیتی نگه دارد که در مجموع حدود $2.5$ مگابایت حافظه برای هر استخر و $50$ مگابایت برای کل شبکه نیاز دارد. در حالت خط لوله با \lr{watermark} $200$ کیلوبیت، هر استخر تنها حدود $25$ کیلوبایت حافظه نیاز دارد.

از نظر کارایی، عملیات \lr{SPF+} با پیچیدگی $\mathcal{O}((V+E) \log V) \approx \mathcal{O}(17 \log 7) \approx 50$ عملیات برای هر درخواست، بسیار سریع است. عملیات \lr{reserve\_keys} با پیچیدگی $\mathcal{O}(n \log n)$ که در آن $n$ تعداد بلوک‌ها است، در بدترین حالت با $n \approx 78000$ به حدود $1.3$ میلیون عملیات می‌رسد که در عملی روی سخت‌افزار مدرن در کمتر از $100$ میلی‌ثانیه اجرا می‌شود. عملیات \lr{simulate\_link} در حالت آماری تقریباً بی‌درنگ است، در حالی که در حالت خط لوله با $10^7$ پالس، حدود چند ثانیه تا چند دقیقه زمان نیاز دارد بسته به پیکربندی.

به‌طور کلی، این ماژول به‌عنوان لایه بالاترین انتزاع در فریم‌ورک شبیه‌سازی \lr{QKD} عمل می‌کند و در عین حال از موتور فیزیکی دقیق بهره می‌برد. طراحی ماژول‌ار آن امکان توسعه آسان را فراهم می‌سازد؛ برای مثال، افزودن گره‌های جدید تنها نیازمند اصلاح \lr{build\_network} و \lr{NODE\_DISTANCES} است، و افزودن پروتکل‌های رله جدید تنها با ایجاد زیرکلاس جدید از \lr{QKDNode} امکان‌پذیر است. این انعطاف‌پذیری، ماژول را به ابزاری مناسب برای مطالعه شبکه‌های \lr{QKD} در مقیاس متروپولیتن تبدیل می‌کند.

	
	% =========================== بخش چهارم: اعتبارسنجی و تست‌ها ===========================
	\part{اعتبارسنجی و آزمون}
	
	\chapter{اعتبارسنجی لایه فیزیکی}
	\label{ch:val_physics}
	\sloppy
\emergencystretch=5em
\hbadness=10000
\vbadness=10000
\tolerance=2000

\providecolor{codebg}{RGB}{248,248,248}
\providecolor{codeframe}{RGB}{200,200,200}
\providecolor{keyword}{RGB}{0,0,180}
\providecolor{string}{RGB}{163,21,21}
\providecolor{comment}{RGB}{0,128,0}

\lstset{
	backgroundcolor=\color{codebg},
	frame=single,
	rulecolor=\color{codeframe},
	language=Python,
	basicstyle=\ttfamily\scriptsize,
	keywordstyle=\color{keyword}\bfseries,
	stringstyle=\color{string},
	commentstyle=\color{comment}\itshape,
	showstringspaces=false,
	breaklines=true,
	breakatwhitespace=true,
	numbers=left,
	numberstyle=\tiny\color{gray},
	numbersep=8pt,
	tabsize=4,
	literate=
	{ف}{{{\bfseries\itshape ف}}}1
	{ی}{{{\bfseries\itshape ی}}}1
}

\section{نمای کلی و هدف ماژول}

ماژول \lr{\texttt{validate\_physics.py}} یک اسکریپت اعتبارسنجی فیزیکی است که نقش \emph{بین‌اللایه‌ای حیاتی} در فریم‌ورک شبیه‌سازی \lr{QKD} ایفا می‌کند. این ماژول به‌عنوان یک \lr{golden reference} عمل می‌کند که خروجی موتور شبیه‌سازی کامل \lr{\texttt{main\_optimized.py}} را در برابر معادلات تحلیلی بسته و مستقل از پیاده‌سازی، مقایسه می‌نماید و با استفاده از آزمون آماری \lr{Z-test} دودویی (\lr{binomial Z-test}) در سطح معنی‌داری $\alpha = 10^{-4}$، انطباق آماری نتایج را بررسی می‌سازد. وجود چنین اسکریپت اعتبارسنجی جداگانه‌ای در یک فریم‌ورک علمی، نشانه بلوغ مهندسی نرم‌افزار است؛ زیرا هرگونه تغییر در پیاده‌سازی موتور فیزیکی، چه عمدی (مانند اصلاح یک باگ) و چه ناخواسته (مانند \lr{regression} پس از \lr{refactor})، بلافاصله توسط این آزمون‌ها شناسایی می‌شود و پیش از انتشار نتایج، به اطلاع توسعه‌دهنده می‌رسد.

هدف اصلی این ماژول، پاسخ به پرسش بنیادین زیر است:
\emph{آیا خروجی شبیه‌سازی مونت‌کارلو که از خط لولهٔ کامل
	\lr{\texttt{Source -> Channel -> Protocol -> Detector -> Sifting -> Proof}}
	به دست می‌آید، از نظر آماری با پیش‌بینی تئوریکِ معادلات تحلیلیِ بسته هم‌خوان است؟}
این پرسش به دو دلیل اهمیت دارد. نخست، در شبیه‌سازی‌های فیزیکی که در آن‌ها فرایندهای تصادفی متعددی (نظیر توزیع \lr{Poisson} فوتون‌ها، شمارش تاریکِ آشکارساز (\lr{dark count})، غربال‌سازیِ پایه‌ها (\lr{sifting}) و غیره) با یکدیگر ترکیب می‌شوند، خطاهای ظریف در پیاده‌سازی ممکن است به نتایجی منجر شود که از نظر فیزیکی قابل قبول به نظر می‌رسند، اما از نظر آماری \emph{biased} هستند. دوم، معادلات تحلیلیِ بسته که در ادبیات \lr{QKD} به‌طور گسترده پذیرفته شده‌اند (مانند معادلات \lr{Gain} و \lr{QBER} در مدل \lr{Gottesman--Lo--Lütkenhaus--Preskill}) به‌عنوان مرجع علمی عمل می‌کنند و انطباق با آن‌ها، اعتبار فیزیکی پیاده‌سازی را تضمین می‌نماید.


ساختار کلی ماژول شامل چهار تابع محاسباتی تحلیلی، یک تابع آزمون آماری و یک تابع اجرای شبیه‌سازی ایزوله است. تابع \lr{\texttt{channel\_transmittance}} ضریب انتقال کانال فیبر نوری را بر اساس قانون \lr{Beer-Lambert} محاسبه می‌کند. تابع \lr{\texttt{Q\_mu\_poisson}} \lr{Gain} (نرخ کلیک به ازای هر پالس) را برای منبع \lr{Weak Coherent Pulse} با توزیع \lr{Poisson} محاسبه می‌کند. تابع \lr{\texttt{Q\_mu\_thermal}} همین کمیت را برای منبع \lr{Thermal} (تک‌حالت یا چندحالت) محاسبه می‌نماید. تابع \lr{\texttt{E\_mu\_poisson}} \lr{QBER} متوسط را برای منبع \lr{Poisson} به دست می‌آورد. تابع \lr{\texttt{dead\_time\_correction}} تصحیح \lr{dead time} آشکارساز را بر اساس مدل \lr{Non-Paralyzable} اعمال می‌کند. تابع \lr{\texttt{validate\_statistic}} آزمون \lr{Z-test} دودویی را روی یک مشاهده انجام می‌دهد و نتیجه \lr{PASS}/\lr{FAIL} را برمی‌گرداند. در نهایت، تابع \lr{\texttt{run\_isolated\_simulation}} یک نشست شبیه‌سازی مستقل روی موتور \lr{\texttt{main\_optimized.py}} اجرا می‌کند و خروجی را به فرمت مناسب برای آزمون آماری بازمی‌گرداند.

سه گروه آزمون در این ماژول تعریف شده است که هر یک به جنبه‌ای متفاوت از فیزیک شبیه‌سازی می‌پردازد. \lr{TEST 1} منبع \lr{Weak Coherent Pulse} با توزیع \lr{Poisson} را در سه فاصله $0$، $50$ و $100$ کیلومتر اعتبارسنجی می‌کند؛ این آزمون هم \lr{Gain} و هم \lr{QBER} را در رژیم مجانبی (یعنی بدون \lr{decoy} و بدون \lr{dead time}) بررسی می‌سازد. \lr{TEST 2} منبع \lr{Thermal} تک‌حالت را در فاصله صفر آزمون می‌کند تا توزیع آماری متفاوت منبع را راستی‌آزمایی نماید. \lr{TEST 3} فیزیک \lr{dead time} آشکارساز \lr{SPAD} را در حضور دو آشکارساز مستقل بررسی می‌کند و این موضوع مهم را که هر \lr{SPAD} به‌طور مستقل زمان مرده خود را تجربه می‌کند، به آزمون می‌گذارد.

یکی از ویژگی‌های برجسته این ماژول، انتخاب سطح معنی‌داری $\alpha = 10^{-4}$ است که به \lr{Z-crit} $\approx 3.8906$ متناظر می‌شود. این سطح به‌مراتب سخت‌گیرانه‌تر از سطح متداول $\alpha = 0.05$ (\lr{Z-crit} $\approx 1.96$) یا حتی $\alpha = 0.001$ (\lr{Z-crit} $\approx 3.29$) است. دلیل این انتخاب سخت‌گیرانه آن است که در شبیه‌سازی \lr{QKD} با $N = 100000$ پالس، حتی انحرافات کوچک از تئوری می‌توانند به‌سرعت شناسایی شوند و توسعه‌دهنده را از پذیرش نتایج \emph{تقریباً درست} که در عمل \emph{biased} هستند، بازدارند. با این حال، این سطح به‌قدری سخت‌گیرانه نیست که نویز آماری ذاتی شبیه‌سازی مونت‌کارلو باعث \lr{false positive} مکرر شود؛ تعادل میان حساسیت و \lr{false positive rate} با دقت انتخاب شده است.

در سطح عملیاتی، این ماژول به‌عنوان یک \lr{CI/CD gate} عمل می‌کند؛ یعنی پیش از \lr{merge} هر تغییر در شاخه اصلی، این اسکریپت اجرا می‌شود و در صورت شکست هر یک از آزمون‌ها، \lr{pipeline} متوقف می‌گردد. این رویکرد \lr{shift-left testing} تضمین می‌کند که باگ‌های فیزیکی در همان مراحل اولیه توسعه شناسایی شوند، نه پس از آنکه نتایج شبیه‌سازی در مقالات یا گزارش‌های فنی منتشر شده باشند. در ادبیات مهندسی نرم‌افزار علمی، این الگو به‌عنوان \lr{regression testing for computational physics} شناخته می‌شود و در پروژه‌های بالغی مانند \lr{Quantum ESPRESSO}، \lr{LAMMPS} و \lr{Qiskit} به‌صورت استاندارد به کار می‌رود.

\section{مبانی تئوری و فیزیکی}

\subsection{کانال فیبر نوری و قانون Beer-Lambert}

انتقال فوتون در فیبر نوری تک‌حالت با قانون \lr{Beer-Lambert} مدل می‌شود که رابطه بنیادین بین تلفات فیبر، فاصله و ضریب انتقال را بیان می‌دارد. برای یک فیبر به طول $L$ کیلومتر با ضریب تلفت خطی $\alpha_{\text{dB/km}}$ (بر حسب دسی‌بل بر کیلومتر)، تلفات کلی کانال بر حسب دسی‌بل برابر با $\alpha_{\text{dB/km}} \cdot L$ است. چون مقیاس دسی‌بل لگاریتمی است، تبدیل آن به ضریب انتقال خطی (\lr{power transmittance}) از رابطه زیر پیروی می‌کند:

\begin{equation}
	\eta_{\text{channel}}(L, \alpha_{\text{dB/km}}) = 10^{-\frac{\alpha_{\text{dB/km}} \cdot L}{10}}
\end{equation}

در این رابطه، $L$ بر حسب کیلومتر و $\alpha_{\text{dB/km}}$ بر حسب \lr{dB/km} است. در طول‌موج مخابراتی استاندارد $1550$ نانومتر، فیبر تک‌حالت استاندارد ضریب تلفت $\alpha_{\text{dB/km}} \approx 0.2$ دارد؛ بنابراین در فاصله $50$ کیلومتر، ضریب انتقال کانال برابر با $10^{-1} = 0.1$ و در فاصله $100$ کیلومتر، برابر با $10^{-2} = 0.01$ می‌شود. این کاهش نمایی ضریب انتقال با فاصله، عامل اصلی محدودیت برد \lr{QKD} در فیبر نوری است و دلیل آن است که در فواصل بزرگ‌تر از $200$ کیلومتر، حتی با استفاده از \lr{decoy-state}، نرخ کلید امن به صفر میل می‌کند.

ضریب انتقال کلی سیستم شامل بازده آشکارساز $\eta_{\text{det}}$ نیز می‌شود که در این ماژول به‌صورت یک کمیت مستقل در نظر گرفته شده است:

\begin{equation}
	\eta_{\text{total}}(L) = \eta_{\text{det}} \cdot \eta_{\text{channel}}(L, \alpha_{\text{dB/km}})
\end{equation}

در پیاده‌سازی پیش‌فرض این ماژول، $\eta_{\text{det}} = 0.15$ است که متناظر با آشکارساز \lr{InGaAs SPAD} در دمای $-50$ درجه سانتی‌گراد با بازده کوانتومی متداول است. در عمل، این مقدار در بازه $0.1$ تا $0.25$ متغیر است و به دمای آشکارساز، ولتاژ بایاس و طول‌موج فوتون بستگی دارد.

\subsection{منبع Weak Coherent Pulse و توزیع Poisson}

در اکثر پیاده‌سازی‌های عملی \lr{QKD}، به‌جای منبع تک‌فوتونی آرمانی (که هنوز در دسترس تجاری نیست)، از منبع \lr{Weak Coherent Pulse} (\lr{WCP}) استفاده می‌شود. در این منبع، لیزر با شدت متوسط $\mu$ فوتون در هر پالس تحریک می‌شود و تعداد فوتون‌های هر پالس از توزیع \lr{Poisson} پیروی می‌کند:

\begin{equation}
	P(n | \mu) = \frac{\mu^n e^{-\mu}}{n!}, \quad n = 0, 1, 2, \ldots
\end{equation}

با این توزیع، احتمال اینکه یک پالس حاوی $n$ فوتون باشد به‌صورت بالا داده می‌شود. پارامتر $\mu$ میانگین تعداد فوتون در هر پالس است و در \lr{decoy-state BB84} معمولاً در بازه $0.1$ تا $0.5$ انتخاب می‌شود. مقدار $\mu = 0.5$ که در این ماژول به‌کار رفته، یک انتخاب متعادل میان نرخ کلید (که با $\mu$ رشد می‌کند) و امنیت (که با $\mu$ کاهش می‌یابد، چون پالس‌های چندفوتونی به حمله \lr{Photon-Number Splitting} حساس هستند) است.

\lr{Gain} کل سیستم (\lr{detection yield}) که با $Q_\mu$ نشان داده می‌شود، احتمال آن است که یک پالس منجر به ثبت کلیک در آشکارساز شود. برای منبع \lr{WCP} با توزیع \lr{Poisson}، در حضور دارک کانت با نرخ $Y_0$ و بازده کلی $\eta$، این کمیت از رابطه بسته زیر به دست می‌آید:

\begin{equation}
	Q_\mu^{\text{Poisson}}(\mu, \eta, Y_0) = 1 - (1 - Y_0) \cdot e^{-\mu \cdot \eta}
\end{equation}

این رابطه از جمع احتمالات تمام حالات ممکن به دست می‌آید. به‌طور شهودی، عبارت $e^{-\mu \cdot \eta}$ احتمال آن است که هیچ فوتونی شناسایی نشود (یعنی یا هیچ فوتونی ارسال نشده، یا فوتون‌های ارسالی از دست رفته‌اند) و عامل $(1 - Y_0)$ احتمال آن است که دارک کانت نیز رخ ندهد. بنابراین، $1 - (1-Y_0) e^{-\mu \eta}$ احتمال متمم است: یعنی احتمال آنکه حداقل یک کلیک (چه از فوتون واقعی و چه از دارک کانت) ثبت شود. این فرمول در ادبیات \lr{QKD} به‌ویژه در کار کلاسیک \lr{Lo-Ma-Chen (2005)} به‌تفصیل استخراج شده است.

\subsection{منبع Thermal و توزیع Bose-Einstein}

در برخی سناریوهای \lr{QKD}، به‌جای لیزر \lr{WCP} از منبع \lr{Thermal} استفاده می‌شود. منبع حرارتی فوتون‌ها با توزیع \lr{Bose-Einstein} تولید می‌کند که برای یک منبع تک‌حالتی به‌صورت زیر تعریف می‌شود:

\begin{equation}
	P(n | \mu) = \frac{\mu^n}{(1 + \mu)^{n+1}}, \quad n = 0, 1, 2, \ldots
\end{equation}

برای یک منبع چندحالتی با $M$ حالت مستقل، توزیع مرجع به توزیع \lr{Negative Binomial} میل می‌کند. \lr{Gain} کل سیستم برای منبع حرارتی با $M$ حالت، از رابطه زیر به دست می‌آید:

\begin{equation}
	Q_\mu^{\text{Thermal}}(\mu, \eta, Y_0, M) = 1 - \frac{1 - Y_0}{\left(1 + \frac{\mu \cdot \eta}{M}\right)^M}
\end{equation}

این رابطه تعمیم طبیعی فرمول \lr{Poisson} است. در حد $M \to \infty$، توزیع \lr{Thermal} به \lr{Poisson} میل می‌کند و فرمول بالا به فرمول \lr{Poisson} تقلیل می‌یابد. برای $M = 1$ (تک‌حالت)، این رابطه به $Q_\mu^{\text{Thermal}} = 1 - (1-Y_0)/(1+\mu \eta)$ ساده می‌شود که فرمول کلاسیک منبع حرارتی تک‌حالت است. تفاوت فیزیکی میان \lr{Poisson} و \lr{Thermal} در آن است که در \lr{Thermal}، فوتون‌ها تمایل به \emph{bunching} دارند (یعنی احتمال پالس‌های چندفوتونی بیشتر است) که این امر امنیت را در برابر حمله \lr{PNS} کاهش می‌دهد. با این حال، مطالعه \lr{Thermal} از منظر آکادمیک اهمیت دارد چرا که برخی منابع نور تک‌فوتونی (مانند \lr{quantum dot}) در رژیم خاصی رفتار حرارتی نشان می‌دهند.

\subsection{مدل QBER برای منبع Poisson}

نرخ خطای کوانتومی بیت (\lr{Quantum Bit Error Rate}) که با $E_\mu$ نشان داده می‌شود، کسری از بیت‌های \lr{sifted} است که بین دو طرف ناهماهنگ هستند. در مدل تحلیلی، این کمیت از نسبت خطاهای کلیک به کل کلیک‌ها به دست می‌آید:

\begin{equation}
	E_\mu^{\text{Poisson}}(\mu, \eta, Y_0, e_d) = \frac{0.5 \cdot Y_0 + e_d \cdot (Q_\mu - Y_0)}{Q_\mu}
\end{equation}

در این رابطه، $e_d$ خطای نوری ذاتی (\lr{misalignment error}) است که از نقص‌های تداخل‌سنج ناشی می‌شود و $Y_0$ نرخ دارک کانت است. صورت این کسر نشان می‌دهد که خطا از دو منبع می‌آید: دارک کانت که کاملاً تصادفی است و نرخ خطای $0.5$ دارد (یعنی نیمی از دارک کانت‌ها با بیت ارسالی هم‌خوان و نیمی ناهم‌خوان هستند)، و فوتون‌های واقعی که با احتمال $e_d$ دچار خطا می‌شوند. مخرج $Q_\mu$ کل کلیک‌ها است که مجموع دارک و سیگنال است. در نتیجه، $E_\mu$ کسر وزنی خطاها است.

این فرمول نشان می‌دهد که در حضور دارک کانت، حتی اگر $e_d = 0$ باشد (یعنی هیچ خطای نوری ذاتی وجود نداشته باشد)، \lr{QBER} همچنان بزرگ‌تر از صفر خواهد بود؛ این پدیده به‌ویژه در فواصل طولانی که سیگنال ضعیف می‌شود، اهمیت پیدا می‌کند و دلیل اصلی آن است که در فواصل بزرگ، \lr{QBER} به $0.5$ میل می‌کند (زیرا بیشتر کلیک‌ها دارک هستند). در حد $\eta \to 0$ (یعنی کانال با تلفات بسیار بالا)، $Q_\mu \to Y_0$ و در نتیجه $E_\mu \to 0.5$ که نشان‌دهنده از دست رفتن کامل اطلاعات است.

در صورت کسر، عبارت $0.5 \cdot Y_0$ نشان‌دهنده آن است که دارک کانت‌ها به‌صورت کاملاً تصادفی رخ می‌دهند و بنابراین نرخ خطای آن‌ها برابر $0.5$ است. عبارت $e_d \cdot (Q_\mu - Y_0)$ نشان‌دهنده خطای ناشی از فوتون‌های واقعی است؛ $(Q_\mu - Y_0)$ بخشی از کلیک‌ها را نشان می‌دهد که از فوتون‌های واقعی می‌آیند (نه از دارک کانت) و این بخش با احتمال $e_d$ دچار خطا می‌شود. بنابراین، صورت کل کسر برابر است با کل خطاهای ثبت‌شده در واحد پالس.

\subsection{مدل Dead Time آشکارساز SPAD}

آشکارسازهای \lr{Single-Photon Avalanche Diode} (\lr{SPAD}) پس از هر کلیک، به مدتی برای بازیابی نیاز دارند که به آن \lr{dead time} می‌گویند. در این بازه زمانی $\tau_d$، آشکارساز قادر به ثبت کلیک جدیدی نیست. دو مدل متداول برای این پدیده وجود دارد: \lr{Non-Paralyzable} و \lr{Paralyzable}. در این ماژول، مدل \lr{Non-Paralyzable} به‌کار رفته است که در آن، کلیک‌های رخ‌داده در زمان مرده آشکارساز را برای مدت بیشتر \emph{بازنشانی نمی‌کنند} بلکه صرفاً نادیده گرفته می‌شوند.

برای یک آشکارساز \lr{Non-Paralyzable} با احتمال کلیک ایده‌آل $Q_{\text{ideal}}$ به ازای هر پالس و بازه زمانی میان پالس‌ها $T_{\text{pulse}}$، احتمال کلیک مشاهده‌شده $Q_{\text{obs}}$ از رابطه زیر به دست می‌آید:

\begin{equation}
	Q_{\text{obs}} = \frac{Q_{\text{ideal}}}{1 + Q_{\text{ideal}} \cdot \frac{\tau_d}{T_{\text{pulse}}}}
\end{equation}

این رابطه از حل معادله تابع ضمنی $Q_{\text{obs}} = Q_{\text{ideal}} \cdot (1 - Q_{\text{obs}} \cdot \tau_d / T_{\text{pulse}})$ به دست می‌آید که نشان می‌دهد کسر زمانی که آشکارساز در دسترس است، برابر با $1 - Q_{\text{obs}} \cdot \tau_d / T_{\text{pulse}}$ است. در حد $\tau_d \to 0$، $Q_{\text{obs}} \to Q_{\text{ideal}}$ (بدون کاهش) و در حد $\tau_d \to \infty$ یا $Q_{\text{ideal}} \to 1$، $Q_{\text{obs}} \to T_{\text{pulse}} / \tau_d$ (اشباع).

یکی از نکات ظریف در پیاده‌سازی این ماژول، آن است که در سیستم \lr{BB84} دو آشکارساز مستقل \lr{D0} و \lr{D1} وجود دارند که هر کدام زمان مرده \emph{مستقل} خود را تجربه می‌کنند. بنابراین، احتمال کلیک \lr{Gain} ایده‌آل ابتدا باید به دو نیمه تقسیم شود (هر آشکارساز نیمی از کلیک‌ها را ثبت می‌کند)، سپس تصحیح \lr{dead time} به‌صورت مستقل روی هر آشکارساز اعمال شود:

\begin{equation}
	Q_{\text{per-spad}} = \frac{Q_{\text{ideal}}^{\text{raw}}}{2}, \quad Q_{\text{obs}}^{\text{per-spad}} = \frac{Q_{\text{per-spad}}}{1 + Q_{\text{per-spad}} \cdot \tau_d / T_{\text{pulse}}}
\end{equation}

\begin{equation}
	Q_{\text{obs}}^{\text{total}} = 2 \cdot Q_{\text{obs}}^{\text{per-spad}}
\end{equation}

در نهایت، چون \lr{sifting} در \lr{BB84} نیمی از پالس‌ها را حذف می‌کند (پالس‌هایی که پایه گیرنده با فرستنده هم‌خوان نداشته)، \lr{Gain} نهایی که در خروجی شبیه‌سازی مشاهده می‌شود، برابر با $0.5 \cdot Q_{\text{obs}}^{\text{total}}$ است:

\begin{equation}
	Q_{\text{final}} = Q_{\text{obs}}^{\text{per-spad}} \cdot 2 \cdot 0.5 = Q_{\text{obs}}^{\text{per-spad}}
\end{equation}

این رابطه ظریف نشان می‌دهد که در نهایت، \lr{Gain} مشاهده‌شده با \lr{Gain} هر آشکارساز برابر است؛ چیزی که در نگاه اول غیرمنتظره به نظر می‌رسد اما با بررسی دقیق عوامل $2$ و $0.5$ که یکدیگر را خنثی می‌کنند، آشکار می‌گردد.

\subsection{آزمون آماری Z-test دودویی}

برای اعتبارسنجی آماری نتایج شبیه‌سازی، این ماژول از \lr{Z-test} دودویی استفاده می‌کند که یکی از ساده‌ترین و در عین حال قدرتمندترین آزمون‌های آماری برای متغیرهای دودویی است. فرض کنید $N$ آزمایش مستقل انجام شده و $k$ مورد موفقیت مشاهده شده است؛ اگر احتمال موفقیت تئوریک $p$ باشد، تعداد موفقیت‌های مشاهده‌شده از توزیع دودویی \lr{Binomial(N, p)} پیروی می‌کند. در حد $N \to \infty$، این توزیع به توزیع نرمال $\mathcal{N}(Np, Np(1-p))$ میل می‌کند (\lr{De Moivre-Laplace theorem}). در نتیجه، متغیر استانداردشده $Z$ از توزیع نرمال استاندارد پیروی می‌کند:

\begin{equation}
	Z = \frac{|k - Np|}{\sqrt{Np(1-p)}}
\end{equation}

در این رابطه، $k - Np$ انحراف مشاهده از امید ریاضی است و $\sqrt{Np(1-p)}$ انحراف معیار توزیع دودویی است. اگر $Z$ از یک آستانه بحرانی $Z_{\text{crit}}$ تجاوز کند، فرض صفر (اینکه مشاهده با تئوری هم‌خوان است) رد می‌شود. در سطح معنی‌داری $\alpha = 10^{-4}$، آستانه بحرانی (دوسویه) برابر با $Z_{\text{crit}} \approx 3.8906$ است. این مقدار از تابع توزیع نرمال استاندارد به دست می‌آید:

\begin{equation}
	\alpha = 2 \cdot \Phi(-Z_{\text{crit}}), \quad Z_{\text{crit}} = -\Phi^{-1}(\alpha/2)
\end{equation}

که در آن $\Phi$ تابع توزیع تجمعی نرمال استاندارد است. انتخاب $\alpha = 10^{-4}$ به‌جای $\alpha = 0.05$ (که متداول در علوم اجتماعی است) دو دلیل دارد. نخست، در شبیه‌سازی‌های فیزیکی، ما کنترل کاملی روی منبع تصادفی داریم و می‌توانیم $N$ را به اندازه کافی بزرگ انتخاب کنیم تا نویز آماری کوچک باشد؛ بنابراین آستانه سخت‌گیرانه قابل قبول است. دوم، در \lr{QKD} حتی انحرافات کوچک از تئوری می‌توانند نشانگر باگ‌های ظریف در پیاده‌سازی باشند که در نهایت به نتایج ناامن منجر می‌شوند؛ بنابراین حساسیت بالا ضروری است.

\section{تحلیل معماری و ساختار داده‌ها}

\subsection{ساختار کلی ماژول}

ساختار این ماژول به‌صورت یک \lr{stateless utility module} طراحی شده است؛ یعنی هیچ کلاسی تعریف نشده و تمام توابع به‌صورت مستقل (\lr{pure functions} یا \lr{near-pure functions}) عمل می‌کنند. این طراحی به این دلیل انتخاب شده که اعتبارسنجی فیزیکی باید کاملاً قطعی و قابل تکرار باشد و نباید به وضعیت داخلی هیچ شیء وابسته باشد. هر تابع با دریافت ورودی‌های صریح، خروجی مشخصی تولید می‌کند و عوارض جانبی (\lr{side effects}) آن به چاپ نتیجه آزمون محدود است.

واردسازی (\lr{imports}) در ابتدای ماژول شامل چهار کتابخانه استاندارد پایتون است: \lr{\texttt{math}} برای توابع ریاضی پایه (نظیر \lr{sqrt}، \lr{exp})، \lr{\texttt{sys}} برای دسترسی به \lr{argv} (اگرچه در نسخه فعلی استفاده نمی‌شود، برای توسعه‌های آینده نگه داشته شده)، \lr{\texttt{copy}} برای کپی عمیق پیکربندی، و \lr{\texttt{csv}} برای خواندن فایل خروجی شبیه‌سازی. این حداقل وابستگی تضمین می‌کند که ماژول در هر محیطی قابل اجرا باشد و نیازی به کتابخانه‌های خارجی ندارد.

\begin{LTR}
	\begin{lstlisting}[language=Python]
		import math
		import sys
		import copy
		import csv
	\end{lstlisting}
\end{LTR}

تنها وابستگی خارجی ماژول، به \lr{\texttt{main\_optimized.py}} است که به‌صورت محلی (\lr{local import}) در داخل تابع \lr{run\_isolated\_simulation} بارگذاری می‌شود. این الگو (\lr{lazy import}) دو مزیت دارد: نخست، وابستگی چرخه‌ای (\lr{circular dependency}) را که ممکن است در زمان بارگذاری ماژول رخ دهد، اجتناب می‌کند؛ دوم، اجازه می‌دهد ماژول حتی در صورتی که \lr{\texttt{main\_optimized.py}} در دسترس نباشد، حداقل برای توابع تحلیلی (\lr{channel\_transmittance}، \lr{Q\_mu\_poisson} و غیره) قابل استفاده باشد.

\subsection{توابع تحلیلی channel\_transmittance}

این تابع ساده‌ترین واحد محاسباتی ماژول است و ضریب انتقال کانال را بر اساس قانون \lr{Beer-Lambert} محاسبه می‌کند. با وجود سادگی، این تابع به‌عنوان \lr{building block} پایه برای تمام محاسبات بعدی عمل می‌کند و هر خطا در آن به خطاهای فاحش در تمام آزمون‌ها منجر می‌شود.

\begin{LTR}
	\begin{lstlisting}[language=Python]
		def channel_transmittance(L, alpha_dB_km):
		return 10 ** (-alpha_dB_km * L / 10)
	\end{lstlisting}
\end{LTR}

پارامترهای این تابع عبارت‌اند از: $L$ فاصله به کیلومتر و $\alpha_{\text{dB/km}}$ ضریب تلفت فیبر به \lr{dB/km}. خروجی، ضریب انتقال خطی (در محدوده $[0, 1]$) است. این تابع به‌صورت \lr{stateless} عمل می‌کند و هیچ متغیر سراسری را تغییر نمی‌دهد؛ این ویژگی آن را به یک \lr{pure function} تبدیل می‌کند که آزمون واحد (\lr{unit test}) برای آن بسیار ساده است.

\subsection{توابع تحلیلی Q\_mu\_poisson و Q\_mu\_thermal}

این دو تابع، \lr{Gain} (نرخ کلیک) را برای دو نوع منبع نوری متفاوت محاسبه می‌کنند. تابع \lr{Q\_mu\_poisson} برای منبع \lr{WCP} با توزیع \lr{Poisson} استفاده می‌شود که رایج‌ترین منبع در سیستم‌های \lr{QKD} عملی است.

\begin{LTR}
	\begin{lstlisting}[language=Python]
		def Q_mu_poisson(mu, eta, Y0):
		return 1 - (1 - Y0) * math.exp(-mu * eta)
		
		def Q_mu_thermal(mu, eta, Y0, M=1):
		return 1 - (1 - Y0) / ((1 + mu * eta / M) ** M)
	\end{lstlisting}
\end{LTR}

پارامترهای مشترک این دو تابع عبارت‌اند از: $\mu$ میانگین تعداد فوتون در هر پالس، $\eta$ ضریب انتقال کلی (شامل بازده آشکارساز) و $Y_0$ نرخ دارک کانت. تابع \lr{thermal} پارامتر اضافی $M$ را نیز می‌پذیرد که تعداد حالت‌های مستقل منبع حرارتی را نشان می‌دهد و مقدار پیش‌فرض آن $1$ (تک‌حالت) است.

طراحی این دو تابع به‌گونه‌ای است که هر دو \lr{signature} مشابه دارند (با تفاوت تنها در پارامتر $M$)؛ این امکان را فراهم می‌سازد که در کد فراخوانی، به‌سادگی بین دو مدل منبع جابه‌جا شوند. این الگو (\lr{Strategy pattern}) در مهندسی نرم‌افزار برای پیاده‌سازی الگوریتم‌های قابل تعویض به کار می‌رود.

نکته مهم فیزیکی: در حد $M \to \infty$، تابع \lr{thermal} به \lr{Poisson} میل می‌کند. این را می‌توان با استفاده از حد معروف $\lim_{M \to \infty} (1 + x/M)^M = e^x$ نشان داد:

\begin{equation}
	\lim_{M \to \infty} \frac{1 - Y_0}{\left(1 + \frac{\mu \eta}{M}\right)^M} = (1 - Y_0) \cdot e^{-\mu \eta}
\end{equation}

بنابراین، در حد $M \to \infty$، $Q_\mu^{\text{Thermal}} \to Q_\mu^{\text{Poisson}}$. این ویژگی ریاضی تضمین می‌کند که دو تابع در حد فیزیکی مناسب همگرا می‌شوند و پیاده‌سازی سازگار است.

\subsection{تابع تحلیلی E\_mu\_poisson}

این تابع \lr{QBER} متوسط را برای منبع \lr{WCP} محاسبه می‌کند و بر اساس مدل کلاسیک \lr{GLLP} پیاده‌سازی شده است.

\begin{LTR}
	\begin{lstlisting}[language=Python]
		def E_mu_poisson(mu, eta, Y0, e_d):
		Q = Q_mu_poisson(mu, eta, Y0)
		if Q <= 0: return 0.5
		return (0.5 * Y0 + e_d * (Q - Y0)) / Q
	\end{lstlisting}
\end{LTR}

این تابع ابتدا \lr{Gain} $Q$ را با فراخوانی \lr{Q\_mu\_poisson} محاسبه می‌کند، سپس فرمول \lr{QBER} را اعمال می‌نماید. خط \lr{if Q <= 0: return 0.5} یک محافظ عددی است که در صورتی که $Q$ صفر یا منفی شود (که فیزیکی نیست اما ممکن است در موارد لبه‌ای عددی رخ دهد)، مقدار $0.5$ را برمی‌گرداند. این مقدار $0.5$ نشان‌دهنده آن است که در غیاب هرگونه کلیک، هیچ اطلاعاتی در دست نیست و در نتیجه \lr{QBER} به‌صورت \emph{maximum uncertainty} یعنی $0.5$ تنظیم می‌شود.

پارامتر $e_d$ در این تابع، خطای نوری ذاتی (\lr{misalignment error}) است که از نقص‌های تداخل‌سنج ناشی می‌شود. در عمل، این مقدار معمولاً در بازه $0.005$ تا $0.03$ است و در این ماژول مقدار پیش‌فرض $e_d = 0.01$ به کار رفته است.

\subsection{تابع تحلیلی dead\_time\_correction}

این تابع تصحیح \lr{dead time} را بر اساس مدل \lr{Non-Paralyzable} اعمال می‌کند. این تابع به‌عنوان \lr{building block} پایه برای محاسبه \lr{Gain} در حضور \lr{dead time} عمل می‌کند.

\begin{LTR}
	\begin{lstlisting}[language=Python]
		def dead_time_correction(Q_ideal, tau_d_ns, T_pulse_ns):
		return Q_ideal / (1 + Q_ideal * (tau_d_ns / T_pulse_ns))
	\end{lstlisting}
\end{LTR}

پارامترهای این تابع عبارت‌اند از: $Q_{\text{ideal}}$ احتمال کلیک ایده‌آل (بدون \lr{dead time})، $\tau_d$ زمان مرده به نانوثانیه و $T_{\text{pulse}}$ بازه زمانی میان پالس‌ها به نانوثانیه. خروجی، احتمال کلیک مشاهده‌شده پس از اعمال اثر \lr{dead time} است.

نکته طراحی مهم: این تابع \emph{به‌صورت مستقل} روی یک آشکارساز عمل می‌کند، نه روی کل سیستم. این به این معناست که اگر چندین آشکارساز مستقل وجود داشته باشد (مانند \lr{D0} و \lr{D1} در \lr{BB84})، هر کدام باید به‌صورت جداگانه تصحیح شوند. این موضوع در بخش \lr{TEST 3} به‌صورت صریح پیاده‌سازی شده است.

\subsection{تابع آزمون آماری validate\_statistic}

این تابع هسته آماری ماژول است و آزمون \lr{Z-test} دودویی را روی یک مشاهده انجام می‌دهد. این تابع با دریافت تعداد مشاهده‌شده موفقیت‌ها، تعداد کل آزمایش‌ها، احتمال تئوریک و نام آزمون، نتیجه \lr{PASS}/\lr{FAIL} را برمی‌گرداند.

\begin{LTR}
	\begin{lstlisting}[language=Python]
		def validate_statistic(observed_count, total_trials, expected_prob, test_name):
		expected_count = total_trials * expected_prob
		variance = total_trials * expected_prob * (1 - expected_prob)
		if variance <= 0: variance = 1e-10
		std_dev = math.sqrt(variance)
		z_score = abs(observed_count - expected_count) / std_dev
		
		passed = z_score <= 3.8906
		print(f"{test_name:25} | Exp_p: {expected_prob:.6e} | "
		f"Sim_p: {observed_count/total_trials:.6e} | "
		f"Z: {z_score:6.2f} | {'PASS' if passed else 'FAIL'}")
		return passed
	\end{lstlisting}
\end{LTR}

تحلیل خط‌به‌خط این تابع به شرح زیر است. در خط اول، تعداد مورد انتظار موفقیت‌ها از ضرب تعداد کل آزمایش‌ها در احتمال تئوریک محاسبه می‌شود: $k_{\text{exp}} = N \cdot p$. در خط دوم، واریانس توزیع دودویی از فرمول $\text{Var} = N p (1-p)$ محاسبه می‌شود. در خط سوم، یک محافظ عددی برای جلوگیری از تقسیم بر صفر اعمال می‌شود: اگر واریانس صفر یا منفی باشد (که ممکن است در موارد لبه‌ای مانند $p = 0$ یا $p = 1$ رخ دهد)، مقدار بسیار کوچک $10^{-10}$ جایگزین می‌شود. در خط چهارم، انحراف معیار از جذر واریانس محاسبه می‌شود. در خط پنجم، \lr{Z-score} از نسبت قدرمطلق انحراف به انحراف معیار به دست می‌آید. در خط هفتم، نتیجه آزمون با مقایسه \lr{Z-score} با آستانه بحرانی $Z_{\text{crit}} = 3.8906$ (متناظر با $\alpha = 10^{-4}$ دوسویه) تعیین می‌شود.

قابل توجه است که این تابع در خط چاپ، علاوه بر نتیجه \lr{PASS}/\lr{FAIL}، مقادیر احتمال تئوریک، احتمال شبیه‌سازی‌شده و \lr{Z-score} را نیز نمایش می‌دهد. این اطلاعات تشخیصی (\lr{diagnostic information}) برای تحلیل شکست آزمون‌ها بسیار ارزشمند است؛ زیرا به توسعه‌دهنده اجازه می‌دهد ببیند انحراف از تئوری چقدر بزرگ است و در چه جهتی (مثبت یا منفی) قرار دارد. این الگو (\lr{verbose assertion}) در توسعه نرم‌افزار علمی به‌عنوان بهترین روش شناخته می‌شود.

\subsection{تابع اجرای شبیه‌سازی run\_isolated\_simulation}

این تابع پیچیده‌ترین بخش ماژول است و یک نشست شبیه‌سازی مستقل روی موتور \lr{\texttt{main\_optimized.py}} اجرا می‌کند. این تابع با دریافت یک دیکشنری \lr{config\_overrides}، پیکربندی پیش‌فرض را با مقادیر دلخواه بازنویسی می‌کند و شبیه‌سازی را با پارامترهای ایزوله اجرا می‌نماید.

\begin{LTR}
	\begin{lstlisting}[language=Python]
		def run_isolated_simulation(config_overrides):
		import main_optimized as m
		cfg = copy.deepcopy(m.DEFAULT_CONFIG)
		
		N_pulses = 100000
		dist = config_overrides.get("dist", 0)
		
		cfg["simulation"]["apply_statistical_noise"] = True
		cfg["simulation"]["distance_start_km"] = dist
		cfg["simulation"]["distance_stop_km"] = dist
		cfg["simulation"]["distance_points"] = 1
		cfg["simulation"]["min_pulses_log"] = 5
		cfg["simulation"]["max_pulses_log"] = 5
		
		mu = config_overrides.get("mu", 0.5)
		cfg["source"]["statistics_type"] = config_overrides.get("stat_type", "POISSON")
		cfg["source"]["intensity_jitter"] = 0.0
		
		cfg["source"]["pulses"] = {
			"signal": {"mu": mu, "prob": 1.0},
			"decoy": {"mu": 0.1, "prob": 0.0},
			"vacuum": {"mu": 0.0, "prob": 0.0}
		}
		
		eta_det = config_overrides.get("det_eff", 0.15)
		Y0 = config_overrides.get("Y0", 3e-6)
		ed = config_overrides.get("ed", 0.01)
		
		cfg["detector"]["det_eff_d0"] = eta_det
		cfg["detector"]["det_eff_d1"] = eta_det
		cfg["detector"]["dark_rate"] = Y0
		cfg["detector"]["dark_rate_d1"] = Y0
		cfg["detector"]["qber_intrinsic"] = ed
		cfg["detector"]["misalignment"] = 0.0
		
		cfg["detector"]["dead_time_ns"] = config_overrides.get("dt", 0.0)
		cfg["detector"]["dead_time_model"] = "NON_PARALYZABLE"
		
		m.run_and_save_csv(cfg, num_workers=1, sweeps={})
		
		with open("qkd_results_optimized_sweep.csv", "r") as f:
		rows = list(csv.DictReader(f))
		latest = rows[-1]
		
		y_sig = float(latest.get("detection_yield_signal", 0))
		q_sig = float(latest.get("qber_signal", 0))
		t_pulse = float(cfg["source"]["pulse_period_ns"])
		
		return {
			"clicks": int(y_sig * N_pulses),
			"errors": int(q_sig * y_sig * N_pulses),
			"pulse_period": t_pulse
		}
	\end{lstlisting}
\end{LTR}

تحلیل این تابع به شرح زیر است. در خط اول، ماژول \lr{\texttt{main\_optimized}} به‌صورت محلی (\lr{lazy import}) بارگذاری می‌شود؛ این الگو امکان اجرای ماژول \lr{validate\_physics} را حتی در صورتی که \lr{\texttt{main\_optimized}} در دسترس نباشد، برای توابع تحلیلی فراهم می‌سازد. در خط دوم، یک کپی عمیق از پیکربندی پیش‌فرض \lr{DEFAULT\_CONFIG} ساخته می‌شود. استفاده از \lr{deepcopy} ضروری است زیرا پیکربندی یک دیکشنری تودرتو است و کپی سطحی باعث به اشتراک‌گذاری زیردیکشنری‌ها می‌شود که می‌تواند به تغییرات ناخواسته در پیکربندی پیش‌فرض منجر گردد.

در خطوط بعدی، پیکربندی شبیه‌سازی به‌صورت ایزوله تنظیم می‌شود: \lr{apply\_statistical\_noise = True} فعال‌سازی نویز آماری، \lr{distance\_start\_km = distance\_stop\_km = dist} اجرای شبیه‌سازی در یک فاصله منفرد (بدون \lr{sweep})، \lr{distance\_points = 1} تولید تنها یک نقطه داده، و \lr{min\_pulses\_log = max\_pulses\_log = 5} تنظیم تعداد پالس‌ها به $10^5 = 100000$.

یکی از نکات کلیدی این تابع، تنظیم \lr{prob = 1.0} برای سیگنال و \lr{prob = 0.0} برای \lr{decoy} و \lr{vacuum} است. این تنظیم به این معناست که \emph{تمام} پالس‌ها به‌عنوان سیگنال ارسال می‌شوند و هیچ پالس \lr{decoy} یا \lr{vacuum} وجود ندارد. این ایزوله‌سازی ضروری است زیرا معادلات تحلیلی مرجع (مانند \lr{Q\_mu\_poisson}) برای پالس سیگنال خالص به‌دست آمده‌اند و حضور \lr{decoy} یا \lr{vacuum} نیاز به تجمیع متفاوت دارد.

در بخش آشکارساز، چهار پارامتر کلیدی تنظیم می‌شوند: \lr{det\_eff\_d0} و \lr{det\_eff\_d1} بازده آشکارساز \lr{D0} و \lr{D1}، \lr{dark\_rate} و \lr{dark\_rate\_d1} نرخ دارک کانت هر آشکارساز، \lr{qber\_intrinsic} خطای ذاتی، و \lr{misalignment = 0.0} حذف خطای \lr{misalignment} (زیرا این کمیت به‌صورت جداگانه در پارامتر $e_d$ کنترل می‌شود).

در خط فراخوانی \lr{run\_and\_save\_csv}، شبیه‌سازی با پیکربندی نهایی اجرا می‌شود و خروجی در فایل \lr{qkd\_results\_optimized\_swap.csv} ذخیره می‌گردد. سپس این فایل خوانده شده و آخرین ردیف (که نتیجه شبیه‌سازی است) استخراج می‌شود. در نهایت، سه مقدار کلیدی برگردانده می‌شود: تعداد کلیک‌ها (از ضرب \lr{detection\_yield\_signal} در تعداد پالس‌ها)، تعداد خطاها (از ضرب \lr{QBER} در کلیک‌ها)، و بازه زمانی پالس (برای محاسبه \lr{dead time}).

\section{پیاده‌سازی ریاضی و الگوریتمی}

\subsection{الگوریتم TEST 1: اعتبارسنجی منبع Poisson}

این آزمون، منبع \lr{WCP} با توزیع \lr{Poisson} را در سه فاصله $0$، $50$ و $100$ کیلومتر اعتبارسنجی می‌کند. برای هر فاصله، دو کمیت مستقل آزمون می‌شود: \lr{Gain} (نرخ کلیک) و \lr{QBER} (نرخ خطای بیت).

\begin{LTR}
	\begin{lstlisting}[language=Python]
		print("\n[TEST 1] Poisson WCP Asymptotic")
		for d in [0, 50, 100]:
		res = run_isolated_simulation({
			"dist": d, "mu": 0.5, "stat_type": "POISSON",
			"det_eff": eta_det, "Y0": Y0, "ed": ed, "dt": 0.0
		})
		eta = channel_transmittance(d, 0.2) * eta_det
		Q_th = Q_mu_poisson(0.5, eta, Y0) * 0.5
		E_th = E_mu_poisson(0.5, eta, Y0, ed)
		validate_statistic(res["clicks"], N, Q_th, f"WCP Gain L={d}km")
		validate_statistic(res["errors"], res["clicks"], E_th, f"WCP QBER L={d}km")
	\end{lstlisting}
\end{LTR}

تحلیل ریاضی این آزمون به شرح زیر است. در هر تکرار، ابتدا شبیه‌سازی با پارامترهای ایزوله اجرا می‌شود و تعداد کلیک‌ها و خطاها به دست می‌آید. سپس ضریب انتقال تئوریک از ضرب ضریب انتقال کانال (با $\alpha_{\text{dB/km}} = 0.2$) در بازده آشکارساز محاسبه می‌شود:

\begin{equation}
	\eta_{\text{total}}(d) = \eta_{\text{det}} \cdot 10^{-\frac{0.2 \cdot d}{10}}
\end{equation}

برای $d = 0$ کیلومتر: $\eta_{\text{total}} = 0.15$، برای $d = 50$ کیلومتر: $\eta_{\text{total}} = 0.015$، و برای $d = 100$ کیلومتر: $\eta_{\text{total}} = 0.0015$.

سپس \lr{Gain} تئوریک از فرمول \lr{Q\_mu\_poisson} محاسبه می‌شود و در $0.5$ ضرب می‌گردد تا اثر \lr{sifting} اعمال شود. این ضرب در $0.5$ به این دلیل است که در \lr{BB84}، گیرنده به‌صورت تصادفی یکی از دو پایه را انتخاب می‌کند و تنها در نیمی از موارد، پایه گیرنده با فرستنده هم‌خوان است (\lr{sifting}). در نتیجه، نرخ کلیک مشاهده‌شده نیمی از \lr{Gain} خام است.

\lr{QBER} تئوریک از فرمول \lr{E\_mu\_poisson} به دست می‌آید. نکته مهم آن است که در این محاسبه، \lr{sifting} اعمال نمی‌شود زیرا \lr{QBER} نسبت خطا به کل کلیک‌ها است و این نسبت تحت \lr{sifting} تغییر نمی‌کند (هم خطاها و هم کل کلیک‌ها به همان نسبت $0.5$ کاهش می‌یابند).

دو آزمون آماری جداگانه انجام می‌شود: یکی برای \lr{Gain} (با $N = 100000$ آزمایش و احتمال تئوریک $Q_{\text{th}}$) و دیگری برای \lr{QBER} (با $N = $ تعداد کلیک‌ها و احتمال تئوریک $E_{\text{th}}$). در آزمون \lr{QBER}، تعداد آزمایش‌ها به تعداد کلیک‌ها تغییر می‌کند زیرا \lr{QBER} بر روی کلیک‌ها اعمال می‌شود، نه بر روی کل پالس‌ها. این تمایز ظریف برای صحت آماری آزمون ضروری است.

\subsection{الگوریتم TEST 2: اعتبارسنجی منبع Thermal}

این آزمون، منبع حرارتی تک‌حالت را در فاصله صفر اعتبارسنجی می‌کند. این آزمون کوتاه‌تر از \lr{TEST 1} است زیرا فقط \lr{Gain} بررسی می‌شود و فقط در یک فاصله.

\begin{LTR}
	\begin{lstlisting}[language=Python]
		print("\n[TEST 2] Thermal Single-Mode")
		res = run_isolated_simulation({
			"dist": 0, "mu": 0.5, "stat_type": "THERMAL",
			"det_eff": eta_det, "Y0": Y0, "dt": 0.0
		})
		Q_th = Q_mu_thermal(0.5, eta_det, Y0, M=1) * 0.5
		validate_statistic(res["clicks"], N, Q_th, "Thermal Gain L=0km")
	\end{lstlisting}
\end{LTR}

در این آزمون، پارامتر \lr{stat\_type = "THERMAL"} به موتور شبیه‌سازی دستور می‌دهد که به‌جای توزیع \lr{Poisson}، از توزیع \lr{Bose-Einstein} تک‌حالت برای تولید فوتون استفاده کند. سپس \lr{Gain} تئوریک از فرمول \lr{Q\_mu\_thermal} با $M = 1$ محاسبه می‌شود:

\begin{equation}
	Q_\mu^{\text{Thermal}}(\mu = 0.5, \eta = 0.15, Y_0 = 3 \times 10^{-6}, M = 1) = 1 - \frac{1 - 3 \times 10^{-6}}{1 + 0.5 \times 0.15} \approx 0.075
\end{equation}

این مقدار در $0.5$ ضرب می‌شود تا \lr{sifting} اعمال شود و در نتیجه $Q_{\text{th}} \approx 0.0375$ به دست می‌آید. این یعنی از $100000$ پالس، انتظار می‌رود حدود $3750$ کلیک ثبت شود.

نکته فیزیکی قابل توجه: برای همان $\mu$ و $\eta$، \lr{Gain} منبع حرارتی بزرگ‌تر از \lr{Gain} منبع \lr{Poisson} است. این به دلیل پدیده \lr{bunching} در نور حرارتی است که احتمال پالس‌های چندفوتونی را افزایش می‌دهد. در مقابل، نور \lr{Poisson} دارای آمار \lr{coherent} است که در آن فوتون‌ها مستقل توزیع می‌شوند. این تفاوت آماری یکی از مفاهیم بنیادین اپتیک کوانتومی است که در آزمون ما به‌صورت کمی نشان داده می‌شود.

\subsection{الگوریتم TEST 3: اعتبارسنجی Dead Time مستقل SPAD}

این آزمون پیچیده‌ترین آزمون ماژول است و فیزیک \lr{dead time} در حضور دو آشکارساز مستقل را بررسی می‌کند. این آزمون به‌طور خاص این موضوع مهم را راستی‌آزمایی می‌کند که هر \lr{SPAD} به‌طور مستقل زمان مرده خود را تجربه می‌کند، نه اینکه زمان مرده روی کل سیستم اعمال شود.

\begin{LTR}
	\begin{lstlisting}[language=Python]
		print("\n[TEST 3] Dead-Time Physics (Independent SPADs)")
		dt_ns = 50.0
		res = run_isolated_simulation({
			"dist": 0, "mu": 0.5, "stat_type": "POISSON",
			"det_eff": eta_det, "Y0": Y0, "dt": dt_ns
		})
		
		# فیزیک صحیح: محاسبه زمان مرده روی هر SPAD به طور مستقل
		Q_ideal_raw = Q_mu_poisson(0.5, eta_det, Y0)
		Q_per_spad = Q_ideal_raw / 2.0
		T_pulse = res["pulse_period"]
		
		Q_dead_per_spad = dead_time_correction(Q_per_spad, dt_ns, T_pulse)
		Q_dt_exp = Q_dead_per_spad * 2.0 * 0.5  # 2 SPADs * Sifting probability
		
		validate_statistic(res["clicks"], N, Q_dt_exp, f"Dead-time tau={dt_ns}ns")
	\end{lstlisting}
\end{LTR}

تحلیل ریاضی این آزمون به شرح زیر است. ابتدا \lr{Gain} خام ایده‌آل (بدون \lr{dead time}) از فرمول \lr{Q\_mu\_poisson} محاسبه می‌شود. سپس این مقدار به دو نیمه تقسیم می‌شود زیرا در سیستم \lr{BB84}، هر کلیک به‌صورت تصادفی به یکی از دو آشکارساز \lr{D0} یا \lr{D1} می‌رود (با احتمال برابر $0.5$). در نتیجه، \lr{Gain} هر آشکارساز برابر با $Q_{\text{per-spad}} = Q_{\text{ideal}}^{\text{raw}} / 2$ است.

سپس تصحیح \lr{dead time} به‌صورت مستقل روی هر آشکارساز اعمال می‌شود:

\begin{equation}
	Q_{\text{obs}}^{\text{per-spad}} = \frac{Q_{\text{per-spad}}}{1 + Q_{\text{per-spad}} \cdot \frac{\tau_d}{T_{\text{pulse}}}}
\end{equation}

در نهایت، برای محاسبه \lr{Gain} کل سیستم، مقدار هر آشکارساز در $2$ ضرب می‌شود (چون دو آشکارساز مستقل وجود دارد) و سپس در $0.5$ ضرب می‌شود تا اثر \lr{sifting} اعمال گردد:

\begin{equation}
	Q_{\text{final}} = Q_{\text{obs}}^{\text{per-spad}} \cdot 2 \cdot 0.5 = Q_{\text{obs}}^{\text{per-spad}}
\end{equation}

این رابطه جالب نشان می‌دهد که در نهایت، \lr{Gain} مشاهده‌شده با \lr{Gain} هر آشکارساز (پس از تصحیح \lr{dead time}) برابر است. این ساده‌سازی ریاضی به این دلیل رخ می‌دهد که عوامل $2$ (برای دو آشکارساز) و $0.5$ (برای \lr{sifting}) یکدیگر را خنثی می‌کنند. با این حال، \emph{ترتیب اعمال} این عوامل مهم است: ابتدا باید \lr{dead time} روی هر آشکارساز به‌صورت مستقل اعمال شود و سپس جمع انجام گردد. اگر به‌اشتباه \lr{dead time} روی \lr{Gain} کل اعمال شود، نتیجه متفاوت و نادرستی به دست می‌آید.

برای درک کمی، فرض کنید $Q_{\text{ideal}}^{\text{raw}} = 0.075$ (مقدار تقریبی برای $\mu = 0.5$، $\eta = 0.15$، $Y_0 = 3 \times 10^{-6}$ در فاصله صفر). آنگاه $Q_{\text{per-spad}} = 0.0375$. اگر $\tau_d = 50$ نانوثانیه و $T_{\text{pulse}} = 1000$ نانوثانیه (نرخ پالس $1$ مگاهرتز) باشد:

\begin{equation}
	Q_{\text{obs}}^{\text{per-spad}} = \frac{0.0375}{1 + 0.0375 \times \frac{50}{1000}} = \frac{0.0375}{1.001875} \approx 0.03743
\end{equation}

این نشان می‌دهد که در این شرایط، \lr{dead time} باعث کاهش $\approx 0.18\%$ در \lr{Gain} می‌شود. هرچند این کاهش کوچک به نظر می‌رسد، اما در آزمون آماری با $N = 100000$ قابل شناسایی است زیرا انحراف معیار تئوریک برابر با $\sqrt{100000 \times 0.0375 \times 0.9625} \approx 60$ است و انحراف $\approx 0.18\%$ از $0.0375$ معادل $\approx 6.75$ کلیک است که بسیار کوچک‌تر از $3.89 \times 60 \approx 233$ (\lr{Z-crit}) است. اما اگر \lr{dead time} به‌اشتباه روی کل \lr{Gain} اعمال شود، خطا به‌مراتب بزرگ‌تر خواهد بود و آزمون شکست می‌خورد.

\subsection{محاسبه Q-score و آستانه بحرانی}

انتخاب آستانه بحرانی $Z_{\text{crit}} = 3.8906$ بر اساس سطح معنی‌داری $\alpha = 10^{-4}$ دوسویه است. این مقدار از معادله زیر به دست می‌آید:

\begin{equation}
	\alpha = P(|Z| > Z_{\text{crit}} | H_0) = 2 \cdot \Phi(-Z_{\text{crit}})
\end{equation}

که در آن $\Phi$ تابع توزیع تجمعی نرمال استاندارد است. با حل معادله $\Phi(-Z_{\text{crit}}) = 5 \times 10^{-5}$، مقدار $Z_{\text{crit}} \approx 3.8906$ به دست می‌آید. این مقدار در کد به‌صورت یک ثابت سخت‌کد شده (\lr{hard-coded}) تعریف شده است.

در عمل، انتخاب این سطح سخت‌گیرانه دو دلیل دارد. نخست، در شبیه‌سازی \lr{QKD} با $N = 100000$ پالس، انحراف معیار نسبی برابر با $\sqrt{p(1-p)/N} \approx 0.0016$ برای $p \approx 0.04$ است. این یعنی نویز آماری ذاتی شبیه‌سازی در حد $0.16\%$ است و انحرافات بزرگ‌تر از این مقدار باید شناسایی شوند. با $Z_{\text{crit}} = 3.89$، آستانه شناسایی برابر با $3.89 \times 0.16\% \approx 0.62\%$ است که به‌طور معناداری بزرگ‌تر از نویز است. دوم، این سطح به‌قدری سخت‌گیرانه است که \lr{false positive} بسیار نادر است (احتمال $\alpha = 10^{-4}$ یعنی یک شکست در هر $10000$ اجرای مستقل) اما هنوز به اندازه کافی حساس است که باگ‌های واقعی را شناسایی کند.

\subsection{تحلیل پیچیدگی و کارایی}

از نظر پیچیدگی زمانی، تمام توابع تحلیلی (\lr{channel\_transmittance}، \lr{Q\_mu\_poisson}، \lr{Q\_mu\_thermal}، \lr{E\_mu\_poisson}، \lr{dead\_time\_correction}) در زمان $\mathcal{O}(1)$ اجرا می‌شوند زیرا تنها شامل چند عملیات حسابی ساده هستند. تابع \lr{validate\_statistic} نیز در زمان $\mathcal{O}(1)$ اجرا می‌شود.

تابع \lr{run\_isolated\_simulation} پیچیده‌تر است: زمان اجرای آن به زمان اجرای \lr{run\_and\_save\_csv} در \lr{\texttt{main\_optimized.py}} بستگی دارد که خود به تعداد پالس‌ها $N$ و پیچیدگی خط لوله شبیه‌سازی وابسته است. برای $N = 100000$ پالس، این زمان معمولاً در حد چند ثانیه تا چند ده ثانیه است. خواندن فایل \lr{CSV} در زمان $\mathcal{O}(R)$ اجرا می‌شود که در آن $R$ تعداد ردیف‌های فایل است (که در این حالت تنها $1$ ردیف است زیرا \lr{sweep} خالی است).

از نظر پیچیدگی حافظه، ماژول بسیار سبک است؛ تنها ساختار داده‌های قابل توجه دیکشنری \lr{cfg} (با اندازه ثابت) و لیست \lr{rows} (با اندازه $1$) هستند. در نتیجه، کل ماژول می‌تواند به‌راحتی در محیط‌های با حافظه محدود اجرا شود.

\section{ارسال و دریافت با سایر ماژول‌ها}

\subsection{وابستگی به main\_optimized.py}

مهم‌ترین وابستگی خارجی این ماژول، به فایل \lr{\texttt{main\_optimized.py}} است که به‌عنوان موتور اصلی شبیه‌سازی عمل می‌کند. این وابستگی از طریق سه مؤلفه کلیدی برقرار می‌شود: \lr{DEFAULT\_CONFIG} به‌عنوان پیکربندی پایه، \lr{run\_and\_save\_csv} به‌عنوان نقطه ورود شبیه‌سازی، و فرمت فایل \lr{CSV} خروجی به‌عنوان پروتکل ارتباطی.

\begin{LTR}
	\begin{lstlisting}[language=Python]
		import main_optimized as m
		cfg = copy.deepcopy(m.DEFAULT_CONFIG)
		# ... (modify cfg) ...
		m.run_and_save_csv(cfg, num_workers=1, sweeps={})
	\end{lstlisting}
\end{LTR}

این الگوی وابستگی به این معناست که هرگونه تغییر در \lr{DEFAULT\_CONFIG} (مانند افزودن فیلدهای جدید یا تغییر نام فیلدها) یا در فرمت خروجی \lr{CSV} (مانند تغییر نام ستون‌ها) می‌تواند به شکست ماژول \lr{validate\_physics} منجر شود. برای کاهش این ریسک، ماژول از کلیدهای رشته‌ای (\lr{string keys}) با مقادیر پیش‌فرض استفاده می‌کند (مانند \lr{\texttt{latest.get("detection\_yield\_signal", 0)}}) که در صورت نبود یک کلید، مقدار پیش‌فرض برمی‌گردد و از خطای \lr{KeyError} جلوگیری می‌کند.

\subsection{ورودی‌های ماژول}

ورودی اصلی ماژول از طریق دیکشنری \lr{config\_overrides} در تابع \lr{run\_isolated\_simulation} ارائه می‌شود. این دیکشنری شامل هفت کلید اختیاری است: \lr{dist} (فاصله به کیلومتر، پیش‌فرض $0$)، \lr{mu} (میانگین فوتون، پیش‌فرض $0.5$)، \lr{stat\_type} (نوع توزیع آماری، پیش‌فرض \lr{"POISSON"})، \lr{det\_eff} (بازده آشکارساز، پیش‌فرض $0.15$)، \lr{Y0} (نرخ دارک کانت، پیش‌فرض $3 \times 10^{-6}$)، \lr{ed} (خطای ذاتی، پیش‌فرض $0.01$) و \lr{dt} (زمان مرده به نانوثانیه، پیش‌فرض $0$).

این الگو (\lr{configuration override pattern}) مزایای متعددی دارد. نخست، امکان تنظیم دقیق پارامترهای شبیه‌سازی برای هر آزمون به‌صورت مستقل فراهم می‌شود. دوم، مقادیر پیش‌فرض معنی‌دار هستند و در صورت عدم ارائه یک کلید، مقدار معقولی استفاده می‌شود. سوم، افزودن پارامترهای جدید در آینده بدون شکستن کدهای فراخوان موجود امکان‌پذیر است.

در سطح بالاتر، ماژول سه گروه پارامتر ثابت در بدنه اصلی (\lr{\_\_main\_\_}) تعریف می‌کند: $N = 100000$ (تعداد پالس‌ها)، $\eta_{\text{det}} = 0.15$ (بازده آشکارساز)، $Y_0 = 3 \times 10^{-6}$ (نرخ دارک کانت) و $e_d = 0.01$ (خطای ذاتی). این مقادیر به‌عنوان \lr{golden parameters} عمل می‌کنند و در تمام آزمون‌ها به‌صورت ثابت استفاده می‌شوند. ثابت نگه‌داشتن این پارامترها ضروری است زیرا تغییر آن‌ها می‌تواند به نتایج متفاوتی منجر شود که مقایسه با مقادیر تئوریک را دشوار سازد.

\subsection{خروجی‌های ماژول}

خروجی اصلی ماژول به دو صورت است: خروجی متنی در \lr{stdout} و مقدار بازگشتی بولی از تابع \lr{validate\_statistic}. خروجی متنی شامل یک جدول قالب‌بندی‌شده است که برای هر آزمون، نام، احتمال تئوریک، احتمال شبیه‌سازی‌شده، \lr{Z-score} و نتیجه \lr{PASS}/\lr{FAIL} را نمایش می‌دهد. این خروجی برای تحلیل بصری توسط توسعه‌دهنده طراحی شده است.

علاوه بر خروجی متنی، تابع \lr{validate\_statistic} مقدار بولی \lr{passed} را برمی‌گرداند که می‌تواند برای اتوماسیون در \lr{CI/CD pipeline} استفاده شود. در یکپارچه‌سازی مستمر، اسکریپت می‌تواند تمام مقادیر بازگشتی را جمع کرده و در صورت وجود حتی یک \lr{FAIL}، کد خطا برگرداند تا \lr{pipeline} متوقف شود. این الگو (\lr{exit code based gating}) در ابزارهایی مانند \lr{GitHub Actions} و \lr{GitLab CI} به‌صورت استاندارد پشتیبانی می‌شود.

تابع \lr{run\_isolated\_simulation} نیز یک دیکشنری با سه کلید برمی‌گرداند: \lr{clicks} (تعداد کلیک‌های ثبت‌شده)، \lr{errors} (تعداد خطاهای ثبت‌شده) و \lr{pulse\_period} (بازه زمانی پالس به نانوثانیه). این مقادیر به‌عنوان ورودی به توابع آزمون آماری استفاده می‌شوند و امکان محاسبه مقادیر تئوریک متناظر را فراهم می‌سازند.

\subsection{ارتباط با فایل CSV خروجی}

یکی از جنبه‌های مهم ارتباط با موتور شبیه‌سازی، فرمت فایل \lr{CSV} خروجی است. پس از اجرای \lr{run\_and\_save\_csv}، نتایج در فایل \lr{qkd\_results\_optimized\_sweep.csv} ذخیره می‌شوند و ماژول \lr{validate\_physics} این فایل را می‌خواند.

\begin{LTR}
	\begin{lstlisting}[language=Python]
		with open("qkd_results_optimized_sweep.csv", "r") as f:
		rows = list(csv.DictReader(f))
		latest = rows[-1]
		
		y_sig = float(latest.get("detection_yield_signal", 0))
		q_sig = float(latest.get("qber_signal", 0))
		t_pulse = float(cfg["source"]["pulse_period_ns"])
	\end{lstlisting}
\end{LTR}

استفاده از \lr{csv.DictReader} به این معناست که فایل \lr{CSV} باید دارای سطر سرستون باشد و نام ستون‌ها باید دقیقاً با کلیدهای مورد انتظار (\lr{detection\_yield\_signal} و \lr{qber\_signal}) مطابقت داشته باشد. این یک \emph{قرارداد ارتباطی} (\lr{interface contract}) بین دو ماژول است که باید در صورت تغییر فرمت \lr{CSV}، به‌روزرسانی شود.

نکته طراحی مهم: استفاده از \lr{rows[-1]} برای گرفتن آخرین ردیف به این دلیل است که در حالت \lr{sweep} با چندین نقطه فاصله، چندین ردیف تولید می‌شود. در این ماژول، چون \lr{sweep} خالی است و فقط یک فاصله مشخص شده، تنها یک ردیف تولید می‌شود و \lr{rows[-1]} همان ردیف اول است. اما استفاده از \lr{rows[-1]} به‌جای \lr{rows[0]} یک الگوی دفاعی است که در صورتی که به‌اشتباه چندین ردیف تولید شود، آخرین ردیف (که معمولاً متناظر با آخرین فاصله است) را برمی‌گرداند.

نکته دیگر آن است که \lr{pulse\_period} از \lr{cfg} خوانده می‌شود، نه از فایل \lr{CSV}. این به این دلیل است که \lr{pulse\_period} یک پارامتر پیکربندی است که در طول شبیه‌سازی ثابت می‌ماند و نیازی به ذخیره در خروجی ندارد. با این حال، این الگو به یکپارچگی داده‌ها آسیب می‌زند زیرا اگر \lr{cfg} پس از اجرای شبیه‌سازی تغییر کند، مقدار خوانده‌شده ممکن است با مقدار واقعی استفاده‌شده در شبیه‌سازی متفاوت باشد. در یک طراحی بهتر، این مقدار نیز در فایل \lr{CSV} ذخیره می‌شد.

\subsection{یکپارچه‌سازی با سیستم CI/CD}

این ماژول به‌عنوان یک \lr{test suite} مستقل عمل می‌کند و می‌تواند به‌راحتی در سیستم‌های \lr{CI/CD} یکپارچه شود. الگوی متداول استفاده به این صورت است:

\begin{LTR}
	\begin{lstlisting}[language=Python]
		# In CI script:
		python validate_physics.py
		if [ $? -ne 0 ]; then
		echo "Physics validation failed!"
		exit 1
		fi
	\end{lstlisting}
\end{LTR}

با این حال، در نسخه فعلی ماژول، کد خروج (\lr{exit code}) به‌صورت صریح تنظیم نمی‌شود و حتی در صورت شکست آزمون‌ها، اسکریپت با کد $0$ (موفقیت) خارج می‌شود. این یک محدودیت است که در نسخه‌های آینده باید اصلاح شود؛ به‌عنوان مثال، می‌توان در انتهای بدنه اصلی، تعداد شکست‌ها را شمارش کرده و در صورت بزرگ‌تر بودن از صفر، با \lr{sys.exit(1)} خارج شد.

\subsection{محدودیت‌ها و توسعه‌های آینده}

این ماژول دارای چند محدودیت است که در توسعه‌های آینده قابل بهبود است. نخست، تعداد پالس‌ها $N = 100000$ ثابت است و قابل تنظیم نیست؛ برای آزمون‌های حساس‌تر، می‌توان $N$ را به $10^6$ یا $10^7$ افزایش داد، هرچند این به زمان اجرای طولانی‌تر نیاز دارد. دوم، تنها سه آزمون تعریف شده است و سناریوهای پیشرفته‌تری مانند \lr{decoy-state} کامل، \lr{afterpulse}، و \lr{polarization drift} پوشش داده نمی‌شوند. سوم، خروجی به‌صورت متنی است و برای تحلیل خودکار نیاز به تجزیه دارد؛ می‌توان خروجی را به فرمت \lr{JSON} یا \lr{JUnit XML} تغییر داد تا با ابزارهای \lr{CI} بهتر یکپارچه شود. چهارم، آزمون‌ها به‌صورت قطعی (\lr{deterministic}) هستند زیرا از بذر تصادفی ثابت استفاده نمی‌کنند؛ در نتیجه، در اجراهای مختلف ممکن است نتایج کمی متفاوت باشند. برای اصلاح این، می‌توان از \lr{SeedSequence} با بذر ثابت استفاده کرد.

علاوه بر این، ماژول می‌تواند با معرفی \lr{parametrized tests} گسترش یابد؛ به‌عنوان مثال، می‌توان به‌جای سه فاصله ثابت در \lr{TEST 1}، یک \lr{sweep} روی بازه $[0, 100]$ کیلومتر با گام $10$ کیلومتر انجام داد و آزمون را برای هر نقطه تکرار کرد. این الگو (\lr{parameterized testing}) در فریم‌ورک‌هایی مانند \lr{pytest} به‌صورت استاندارد پشتیبانی می‌شود و می‌تواند اعتبار فیزیکی ماژول را در طیف وسیع‌تری از پارامترها تضمین نماید.

در نهایت، این ماژول به‌عنوان یک \lr{golden reference} برای فیزیک \lr{QKD} در فریم‌ورک عمل می‌کند و هرگونه تغییر در پیاده‌سازی موتور فیزیکی باید پیش از \lr{merge} با این آزمون‌ها اعتبارسنجی شود. این رویکرد \lr{shift-left testing} تضمین می‌کند که باگ‌های فیزیکی در همان مراحل اولیه توسعه شناسایی شوند و به انتشار نتایج نادرست در مقالات یا گزارش‌های فنی منجر نشوند. وجود چنین ماژول اعتبارسنجی مستقل، نشانه بلوغ مهندسی نرم‌افزار در پروژه‌های علمی است و در پروژه‌های تحقیقاتی جدی به‌عنوان یک ضرورت در نظر گرفته می‌شود.

	
	\chapter{اعتبارسنجی رمزنگاری و تحلیلی}
	\label{ch:val_crypto}
	\sloppy
\emergencystretch=5em
\hbadness=10000
\vbadness=10000
\tolerance=2000

\providecolor{codebg}{RGB}{248,248,248}
\providecolor{codeframe}{RGB}{200,200,200}
\providecolor{keyword}{RGB}{0,0,180}
\providecolor{string}{RGB}{163,21,21}
\providecolor{comment}{RGB}{0,128,0}

\lstset{
	backgroundcolor=\color{codebg},
	frame=single,
	rulecolor=\color{codeframe},
	language=Python,
	basicstyle=\ttfamily\scriptsize,
	keywordstyle=\color{keyword}\bfseries,
	stringstyle=\color{string},
	commentstyle=\color{comment}\itshape,
	showstringspaces=false,
	breaklines=true,
	breakatwhitespace=true,
	numbers=left,
	numberstyle=\tiny\color{gray},
	numbersep=8pt,
	tabsize=4,
	literate=
	{ف}{{{\bfseries\itshape ف}}}1
	{ی}{{{\bfseries\itshape ی}}}1
}

\section{نمای کلی و هدف ماژول}

ماژول \lr{\texttt{absolute\_crypto\_validation.py}} یک \emph{سوییت اعتبارسنجی مطلق کریپتوگرافیک} است که نقشی متمایز و مکمل نسبت به سایر اسکریپت‌های آزمون در فریم‌ورک شبیه‌سازی \lr{QKD} ایفا می‌کند. در حالی که ماژول \lr{\texttt{validate\_physics.py}} بر آزمون آماری توزیع‌های منبع (\lr{Poisson} و \lr{Thermal}) متمرکز است و از آزمون \lr{Z-test} دودویی برای راستی‌آزمایی \lr{Gain} و \lr{QBER} استفاده می‌نماید، این ماژول به \emph{صفت کریپتوگرافیک هسته‌ای} پروتکل می‌پردازد: آیا کمیت‌های بنیادین امنیت کلید نظیر \lr{single-photon yield} ($Y_1$)، \lr{single-photon QBER} ($e_1$) و \lr{secure key rate} ($R$) که موتور کامل شبیه‌سازی محاسبه می‌کند، با معادلات تحلیلی بسته‌ای که در ادبیات کلاسیک \lr{decoy-state QKD} (به‌ویژه کار مرجع \lr{Lo-Ma-Chen 2005}) پذیرفته شده‌اند، انطباق عددی دارند؟

هدف اصلی این ماژول، پاسخ به پرسش زیر است: \emph{آیا پیاده‌سازی کامل موتور شبیه‌سازی، از منبع \lr{WCP} و کانال فیبر و آشکارساز \lr{SPAD} و \lr{sifting} پایه‌ها و پروتکل \lr{BB84-Decoy} و اثبات امنیتی \lr{Lim-2014} تا استخراج کلید نهایی، خروجی عددی تولید می‌کند که از نظر کریپتوگرافیک معتبر باشد؟} این پرسش از آنجا اهمیت می‌یابد که در شبیه‌سازی‌های عددی پیچیده، خطاهای ظریف در پیاده‌سازی — مانند اشتباه در فرمول \lr{Bayes} برای استخراج $Y_1$، یا نادیده گرفتن کسر \lr{sifting} $q_{\text{sift}}$ در نرخ کلید، یا اشتباه در نرمال‌سازی \lr{QBER} — می‌توانند منجر به نتایجی شوند که از نظر \lr{QBER} ظاهراً قابل‌قبول‌اند اما از نظر امنیت کریپتوگرافیک \emph{غیرقابل دفاع} هستند. به‌عنوان مثال، اگر استخراج $Y_1$ به‌جای کران پایین \lr{Lo2005} از کران بالا استفاده کند، شبیه‌سازی نرخ کلیدی گزارش می‌دهد که در عمل قابل دستیابی نیست و هر مقاله‌ای که بر این پایه منتشر شود، از نظر رمزنگاری قابل نقض خواهد بود.

ساختار کلی این ماژول شامل سه تابع تحلیلی و یک تابع اجرای آزمون است. تابع \lr{\texttt{H2}} آنتروپی باینری شانون را با برش عددی (\lr{numerical clipping}) پیاده‌سازی می‌کند. تابع \lr{\texttt{channel\_transmittance}} ضریب انتقال کانال فیبر را بر اساس قانون \lr{Beer-Lambert} محاسبه می‌نماید. تابع \lr{\texttt{get\_analytics}} مجموعه‌ای از پنج کمیت تحلیلی کلیدی — $Q_\mu$، $E_\mu$، $Y_1$، $e_1$ و $R_{\text{asymp}}$ — را بر اساس معادلات بسته \lr{decoy-state} محاسبه می‌کند و آن‌ها را در قالب یک دیکشنری برمی‌گرداند. تابع \lr{\texttt{run\_absolute\_test}} آزمون اصلی است که با مقداردهی پیکربندی پیش‌فرض \lr{\texttt{DEFAULT\_CONFIG}} از ماژول \lr{\texttt{main\_optimized}}، اجرای شبیه‌سازی کامل، خواندن خروجی \lr{CSV} و مقایسه با مقادیر تحلیلی، در پنج پارامتر با تلورانس‌های متفاوت، نتیجه نهایی \lr{PASS}/\lr{FAIL} را اعلام می‌کند.

ویژگی برجسته این ماژول، \emph{تفکیک تلورانس‌ها بر اساس حساسیت کریپتوگرافیک} است. تلورانس $25\%$ برای $Q_\mu$ به‌کار رفته است، زیرا این کمیت در شبیه‌سازی مونت‌کارلو با $10^7$ پالس، نویز آماری قابل‌توجهی دارد و انحراف از مقدار تحلیلی طبیعی است. در مقابل، تلورانس $1\%$ برای $Y_1$، $e_1$ و $R_{\text{asymp}}$ به‌کار رفته، چون این کمیت‌ها از ترکیب خطی چندین کمیت به دست می‌آیند و خطای سیستماتیک در هر یک از آن‌ها می‌تواند به خطای بزرگ در نرخ کلید منجر شود. این تلورانس‌های سخت‌گیرانه، تضمین می‌کنند که هرگونه \lr{regression} در پیاده‌سازی موتور — مثلاً پس از \lr{refactor} کلاس اثبات یا تغییر مدل آشکارساز — بلافاصله شناسایی شود.

در سطح عملیاتی، این ماژول به‌عنوان یک \lr{cryptographic gate} عمل می‌کند؛ یعنی پیش از \lr{merge} هر تغییر در شاخه اصلی و پیش از انتشار هر نتیجه شبیه‌سازی در مقالات یا گزارش‌های فنی، این اسکریپت باید با موفقیت اجرا شود. خروجی کد خروج (\lr{exit code}) این اسکریپت نیز به‌گونه‌ای طراحی شده که در \lr{CI/CD pipeline} قابل استفاده باشد: در صورت موفقیت همه آزمون‌ها، کد خروج $0$ و در غیر این صورت، کد خروج $1$ بازگردانده می‌شود. این الگو در مهندسی نرم‌افزار علمی به‌عنوان \lr{cryptographic regression testing} شناخته می‌شود و در پروژه‌های بالغی مانند \lr{Open Quantum Safe} و \lr{Qiskit} به‌صورت استاندارد به کار می‌رود.

نکته مهم دیگر آن است که این ماژول \emph{بدون اعمال هرگونه \lr{patch}} روی موتور اصلی، آن را آزمون می‌کند. به عبارت دیگر، برخلاف برخی اسکریپت‌های آزمون که با \lr{monkey-patching} یا \lr{mock} کردن بخش‌هایی از کد، آن را برای آزمون قابل‌کنترل می‌کنند، این ماژول از همان مسیر اجرای واقعی \lr{\texttt{main\_optimized.run\_and\_save\_csv}} استفاده می‌کند و فقط پیکربندی را override می‌نماید. این رویکرد تضمین می‌کند که آزمون، رفتار واقعی موتور را در شرایط عملیاتی منعکس می‌کند و نه رفتار آزمایشگاهی مصنوعی.

\section{مبانی تئوری و فیزیکی}

\subsection{آنتروپی باینری شانون و نقش آن در امنیت QKD}

آنتروپی باینری شانون، یکی از مفاهیم بنیادین نظریه اطلاعات است که در رمزنگاری کوانتومی نقشی مرکزی ایفا می‌کند. این کمیت که با $H_2(p)$ نشان داده می‌شود، میزان عدم‌قطعیت (\lr{uncertainty}) متغیر تصادفی باینری با توزیع \lr{Bernoulli} را اندازه می‌گیرد و به‌صورت زیر تعریف می‌شود:

\begin{equation}
	H_2(p) = -p \log_2(p) - (1 - p) \log_2(1 - p)
\end{equation}

در این رابطه، $p \in [0, 1]$ احتمال وقوع یکی از دو حالت متغیر باینری است. تابع $H_2(p)$ تقارن نسبت به $p = 0.5$ دارد، در این نقطه به مقدار بیشینه $1$ می‌رسد و در نقاط $p = 0$ و $p = 1$ به صفر میل می‌کند. این رفتار بازتاب‌دهنده این واقعیت استنباطی است که در حالت $p = 0.5$، متغیر تصادفی کاملاً غیرقابل پیش‌بینی است و در حالت‌های $p = 0$ یا $p = 1$ کاملاً قطعی است.

در رمزنگاری کوانتومی، $H_2(p)$ در دو نقطه بحرانی ظاهر می‌شود. نخست، در محاسبه \lr{privacy amplification} که حجم اطلاعات نشت‌کرده به حملهور (\lr{eavesdropper}) را تعیین می‌کند؛ اگر \lr{QBER} سیگنال $E_\mu$ باشد، آنتروپی شرطی شانون $H_2(E_\mu)$ حجم اطلاعات بیت-خطای حملهور را تقریب می‌زند. دوم، در محاسبه \lr{single-photon error rate} $e_1$ که آنتروپی $H_2(e_1)$ حجم اطلاعات حملهور درباره فوتون‌های تک‌فوتونی را تعیین می‌کند و مستقیماً در نرخ کلید امن وارد می‌شود.

پیاده‌سازی تابع $H_2$ در کد، با برش عددی (\lr{numerical clipping}) در نقاط $p = 0$ و $p = 1$ همراه است، زیرا در این نقاط، عبارت $\log_2(p)$ به $-\infty$ میل می‌کند و عبارت $p \log_2(p)$ به حالت تعریف‌نشده $\lim_{p \to 0^+} p \log_2 p = 0$ می‌رسد. در محاسبات عددی کامپیوتری، این حالت می‌تواند منجر به \lr{NaN} یا \lr{exception} شود؛ بنابراین پیاده‌سازی پیش از فراخوانی \lr{\texttt{math.log2}}، شرایط $p \leq 0$ یا $p \geq 1$ را بررسی کرده و در این موارد، مقدار $0$ را برمی‌گرداند. این برش با حد گسسته \lr{L'Hôpital} هم‌خوان است و از نظر ریاضی صحیح است.

\subsection{کاهیدگی فیبر و قانون Beer-Lambert}

انتقال فوتون در فیبر نوری تک‌حالت در طول‌موج مخابراتی $1550$ نانومتر با قانون \lr{Beer-Lambert} مدل می‌شود. این قانون که اصلی‌ترین مدل کاهیدگی در فیبرهای مخابراتی است، رابطه‌ای نمایی بین طول فیبر و ضریب انتقال برقرار می‌سازد. برای یک فیبر به طول $L$ کیلومتر با ضریب کاهیدگی خطی $\alpha_{\text{dB/km}}$، کاهیدگی کل بر حسب دسی‌بل برابر با $\alpha_{\text{dB/km}} \cdot L$ است و تبدیل آن به ضریب انتقال خطی (\lr{power transmittance}) از رابطه زیر پیروی می‌کند:

\begin{equation}
	\eta_{\text{channel}}(L, \alpha_{\text{dB/km}}) = 10^{-\frac{\alpha_{\text{dB/km}} \cdot L}{10}}
\end{equation}

در فیبر تک‌حالت استاندارد در $1550$ نانومتر، $\alpha_{\text{dB/km}} \approx 0.2$ است؛ بنابراین در فاصله $L = 50$ کیلومتر (که نقطه آزمون پیش‌فرض این ماژول است)، ضریب انتقال کانال برابر با $10^{-1} = 0.1$ می‌شود. این کاهش نمایی، عامل اصلی محدودیت برد \lr{QKD} در فیبر است و دلیل آن است که در فواصل بزرگ‌تر از $200$ کیلومتر، حتی با \lr{decoy-state}، نرخ کلید امن به صفر میل می‌کند.

ضریب انتقال کلی سیستم شامل بازده آشکارساز $\eta_{\text{det}}$ نیز می‌شود:

\begin{equation}
	\eta_{\text{total}}(L) = \eta_{\text{det}} \cdot \eta_{\text{channel}}(L, \alpha_{\text{dB/km}})
\end{equation}

در این پیاده‌سازی، $\eta_{\text{det}} = 0.15$ که متناظر با \lr{InGaAs SPAD} در دمای $-50$ درجه سانتی‌گراد است. با این مقادیر، در $L = 50$ کیلومتر، $\eta_{\text{total}} = 0.015$ می‌شود که نشان‌دهنده شرایط عملیاتی واقعی برای \lr{QKD} در فواصل مخابراتی متوسط است.

\subsection{تئوری Decoy-State و معادلات Lo2005}

ایده اصلی \lr{decoy-state QKD} که توسط \lr{Hwang (2003)} پیشنهاد و توسط \lr{Lo-Ma-Chen (2005)} به فرم مدرن آن بسط داده شد، آن است که فرستنده با تغییر تصادفی شدت پالس (\lr{intensity}) در میان چندین سطح (\lr{signal}، \lr{decoy}، \lr{vacuum})، امکان تخمین آماری ویژگی‌های کانال را بدون افشای اینکه کدام پالس حامل کلید است، فراهم می‌سازد. این تکنیک، حمله \lr{Photon-Number Splitting} (\lr{PNS}) را که علیه منبع \lr{WCP} با توزیع \lr{Poisson} قابل اعمال است، خنثی می‌کند و امکان دستیابی به نرخ کلیدی نزدیک به حد امن تک‌فوتونی آرمانی را فراهم می‌سازد.

در ادامه، معادلات بنیادین این تئوری را که همگی در تابع \lr{\texttt{get\_analytics}} پیاده‌سازی شده‌اند، معرفی می‌نماییم.

\subsubsection{Gain سیگنال و QBER سیگنال}

برای منبع \lr{WCP} با توزیع \lr{Poisson} با شدت میانگین $\mu$، \lr{Gain} سیگنال (یعنی احتمال آنکه یک پالس منجر به ثبت کلیک در آشکارساز شود) برابر است با:

\begin{equation}
	Q_\mu^{\text{raw}}(\mu, \eta, Y_0) = 1 - (1 - Y_0) \cdot e^{-\eta \mu}
\end{equation}

در این رابطه، $Y_0$ نرخ \lr{dark count} به ازای هر پالس (احتمال کلیک اشتباه در نبود فوتون واقعی)، $\eta$ ضریب انتقال کلی سیستم و $\mu$ شدت میانگین سیگنال است. عبارت $e^{-\eta \mu}$ احتمال آنکه هیچ فوتونی شناسایی نشود (یعنی یا هیچ فوتونی ارسال نشده، یا فوتون‌های ارسالی از دست رفته‌اند) و عامل $(1 - Y_0)$ احتمال آنکه دارک کانت نیز رخ ندهد را نشان می‌دهد. بنابراین، $1 - (1-Y_0) e^{-\eta \mu}$ احتمال متمم، یعنی احتمال آنکه حداقل یک کلیک ثبت شود، است.

در پروتکل \lr{BB84}، فرستنده و گیرنده هر یک به‌طور تصادفی یکی از دو پایه (\lr{rectilinear} یا \lr{diagonal}) را انتخاب می‌کنند و فقط در حدود نیمی از پالس‌ها (یعنی زمانی که پایه‌ها هم‌خوان باشند) برای استخراج کلید مفید هستند. این کسر به \lr{sifting factor} موسوم است:

\begin{equation}
	q_{\text{sift}} = \frac{1}{2}
\end{equation}

\lr{Gain} سیگنالِ \emph{sifted} شده برابر است با:

\begin{equation}
	Q_\mu^{\text{sifted}} = q_{\text{sift}} \cdot Q_\mu^{\text{raw}}
\end{equation}

\lr{QBER} سیگنال، یعنی کسر کلیک‌های اشتباه در میان کلیک‌های ثبت‌شده، از رابطه زیر به دست می‌آید:

\begin{equation}
	E_\mu = \frac{0.5 \cdot Y_0 + e_d \cdot (1 - e^{-\eta \mu})}{Q_\mu^{\text{raw}}}
\end{equation}

در این رابطه، $e_d$ خطای ذاتی سیستم (\lr{misalignment error}) است که ناشی از عدم‌تقارن پایه‌ها، ناپایداری قطعات اپتیکی و نویز الکترونیکی است. در این پیاده‌سازی، $e_d = 0.01$ به‌کار رفته که در عمل برای سیستم‌های متعادل قابل دستیابی است. صورت این کسر از دو بخش تشکیل شده است: $0.5 \cdot Y_0$ که نشان‌دهنده کلیک‌های اشتباه ناشی از دارک کانت است (چون دارک کانت کاملاً تصادفی است، نصف آن در پایه درست و نصف در پایه اشتباه رخ می‌دهد) و $e_d \cdot (1 - e^{-\eta \mu})$ که نشان‌دهنده کلیک‌های اشتباه ناشی از فوتون‌های واقعی با احتمال خطای ذاتی $e_d$ است.

\subsubsection{Gain دکوی ضعیف}

در پروتکل \lr{decoy-state}، علاوه بر سیگنال با شدت $\mu$، از یک \lr{weak decoy} با شدت $\nu < \mu$ نیز استفاده می‌شود. \lr{Gain} این دکوی ضعیف به‌صورت مشابه محاسبه می‌شود:

\begin{equation}
	Q_\nu(\nu, \eta, Y_0) = 1 - (1 - Y_0) \cdot e^{-\eta \nu}
\end{equation}

این کمیت برای استخراج $Y_1$ (بازده تک‌فوتونی) به‌کار می‌رود. ایده اصلی آن است که با مقایسه $Q_\mu$ و $Q_\nu$ که از دو شدت متفاوت به دست آمده‌اند، می‌توان ویژگی‌های کانال برای فوتون‌های تک‌عددی را استخراج کرد، چون وابستگی $Q_\mu$ به $\mu$ از ترکیب خطی توزیع \lr{Poisson} و ویژگی‌های کانال به دست می‌آید.

\subsubsection{استخراج کران پایین بازده تک‌فوتونی Y1}

مرحله کلیدی در \lr{decoy-state QKD}، استخراج کران پایین برای $Y_1$ (احتمال آنکه یک پالس تک‌فوتونی منجر به کلیک شود) از مشاهدات $Q_\mu$ و $Q_\nu$ است. این کران، که در کار مرجع \lr{Lo-Ma-Chen (2005)} به‌تفصیل استخراج شده، به‌صورت زیر بیان می‌شود:

\begin{equation}
	Y_1^{L} = \frac{\mu^2 e^{-\mu}}{\mu \nu - \mu^2} \left( Q_\nu e^{\nu} - \frac{\nu^2}{\mu^2} Q_\mu^{\text{raw}} e^{\mu} \right)
\end{equation}

این فرمول از ایده زیر نشأت می‌گیرد: $Q_\mu$ را می‌توان به‌صورت جمع وزن‌دار بر روی تعداد فوتون $n$ نوشت:

\begin{equation}
	Q_\mu = \sum_{n=0}^{\infty} Y_n \cdot P(n | \mu) = \sum_{n=0}^{\infty} Y_n \cdot \frac{\mu^n e^{-\mu}}{n!}
\end{equation}

در این رابطه، $Y_n$ بازده کلیک برای پالس $n$-فوتونی است که فرض می‌شود به شدت پالس وابسته نیست (فرض کلیدی \lr{decoy-state}). این فرض بر آن است که حمله حاملهور نمی‌تواند شدت پالس را تشخیص دهد و بنابراین، $Y_n$ برای سیگنال و دکوی یکسان است. با تفریق معادلات $Q_\mu$ و $Q_\nu$ و حذف چندجمله‌ای‌های درجه بالا، می‌توان کران پایین برای $Y_1$ را به دست آورد.

ضریب $\frac{\mu^2 e^{-\mu}}{\mu \nu - \mu^2}$ در فرمول بالا، نرمال‌سازی ریاضی است که از حل سیستم معادلات خطی به دست می‌آید. صورت کران، یعنی $Q_\nu e^{\nu} - \frac{\nu^2}{\mu^2} Q_\mu^{\text{raw}} e^{\mu}$، تفاوت وزن‌دار بین دو مشاهده است که بازده تک‌فوتونی را از نویز چندفوتونی جدا می‌کند. در صورتی که $\mu \approx \nu$ (یعنی شدت سیگنال و دکوی نزدیک به هم باشند)، مخرج $\mu \nu - \mu^2 = \mu(\nu - \mu) \to 0$ میل می‌کند و فرمول به حالت تکین (\lr{singular}) می‌رسد. در این حالت، استخراج $Y_1$ از این روش غیرممکن است و در پیاده‌سازی، مقدار $Y_1^L = 0$ گرفته می‌شود.

\subsubsection{نرخ خطای تک‌فوتونی e1}

نرخ خطای تک‌فوتونی $e_1$، یعنی کسر کلیک‌های اشتباه در میان پالس‌های تک‌فوتونی، از رابطه زیر به دست می‌آید:

\begin{equation}
	e_1 = \frac{0.5 \cdot Y_0 + e_d \cdot \eta}{Y_1}
\end{equation}

این رابطه مشابه $E_\mu$ است، با این تفاوت که به‌جای احتمال کلیک فوتون واقعی $1 - e^{-\eta \mu}$، احتمال کلیک تک‌فوتونی $\eta$ قرار دارد. این جایگزینی برآمده از آن است که برای یک پالس تک‌فوتونی، احتمال شناسایی فوتون دقیقاً برابر با $\eta$ (ضریب انتقال کلی) است، در حالی که برای پالس چندفوتونی، این احتمال ترکیبی از احتمالات کلیک برای هر تعداد فوتون است. در پیاده‌سازی، شرط $Y_1 > 0$ بررسی می‌شود تا از تقسیم بر صفر جلوگیری شود؛ در صورت $Y_1 = 0$، مقدار $e_1 = 0$ گرفته می‌شود.

\subsubsection{بازده تک‌فوتونی و نرخ کلید مجانبی}

\lr{Gain} تک‌فوتونی، یعنی کسری از کلیک‌های کل که از پالس‌های تک‌فوتونی ناشی می‌شوند، از رابطه زیر به دست می‌آید:

\begin{equation}
	Q_1 = \mu \cdot e^{-\mu} \cdot Y_1
\end{equation}

این رابطه از ضرب احتمال آنکه پالس تک‌فوتونی باشد ($P(1 | \mu) = \mu e^{-\mu}$) در بازده تک‌فوتونی $Y_1$ به دست می‌آید. در نهایت، نرخ کلید امن مجانبی (\lr{asymptotic secure key rate}) از فرمول کلاسیک \lr{GLLP} به دست می‌آید:

\begin{equation}
	R_{\text{asymp}} = q_{\text{sift}} \cdot \left[ Q_1 \cdot (1 - H_2(e_1)) - Q_\mu^{\text{raw}} \cdot f_{\text{EC}} \cdot H_2(E_\mu) \right]
\end{equation}

این فرمول، که در ادبیات \lr{QKD} به‌عنوان \lr{GLLP key rate formula} شناخته می‌شود، از دو بخش تشکیل شده است. بخش اول $Q_1 \cdot (1 - H_2(e_1))$ نشان‌دهنده حجم اطلاعات مشترک بین فرستنده و گیرنده از پالس‌های تک‌فوتونی پس از \lr{privacy amplification} است؛ عبارت $1 - H_2(e_1)$ ضریب \lr{privacy amplification} است که حجم اطلاعات نشت‌کرده به حاملهور را کم می‌کند. بخش دوم $Q_\mu^{\text{raw}} \cdot f_{\text{EC}} \cdot H_2(E_\mu)$ نشان‌دهنده حجم اطلاعات فاش‌شده در مرحله \lr{error correction} است؛ $f_{\text{EC}} \geq 1$ بازدهی \lr{error correction} است که در عمل $f_{\text{EC}} \approx 1.1$ برای کد \lr{LDPC} بهینه است. در این پیاده‌سازی، $f_{\text{EC}} = 1.1$ به‌کار رفته است.

تفاوت این فرمول با حالت \lr{finite-key} (مثل اثبات \lr{Lim-2014}) آن است که در اینجا از حد مجانبی $N \to \infty$ استفاده می‌شود و جریمه‌های آماری نظیر \lr{Hoeffding bound} و \lr{composable security} اعمال نمی‌گردند. به همین دلیل، $R_{\text{asymp}}$ همیشه بزرگ‌تر یا مساوی با نرخ کلید متناهی است و در عمل به‌عنوان \emph{کران بالای} نرخ کلید قابل دستیابی عمل می‌کند. با این حال، تطابق موتور شبیه‌سازی با این کران بالا، شرط لازم — هرچند کافی نیست — برای صحت پیاده‌سازی است.

\subsection{تلورانس اعتبارسنجی و تفسیر کریپتوگرافیک}

انتخاب تلورانس‌های متفاوت برای پنج پارامتر، بازتاب‌دهنده درک عمیق از فیزیک و آمار شبیه‌سازی است. تلورانس $25\%$ برای $Q_\mu$ به این دلیل انتخاب شده که در شبیه‌سازی مونت‌کارلو با $N = 10^7$ پالس، انحراف معیار $Q_\mu$ تقریباً برابر با $\sqrt{Q_\mu (1 - Q_\mu) / N} \approx 10^{-4}$ است؛ با این حال، چون دارک کانت $Y_0 = 3 \times 10^{-6}$ بسیار کوچک است، نویز نسبی می‌تواند تا چندین درصد باشد. تلورانس $5\%$ برای $E_\mu$ سخت‌گیرانه‌تر است، چون $E_\mu$ از نسبت دو کمیت به دست می‌آید و واریانس نسبی آن کمتر است.

تلورانس $1\%$ برای $Y_1$، $e_1$ و $R_{\text{asymp}}$ بسیار سخت‌گیرانه است و دلیل آن آن است که این کمیت‌ها در فرمول کلید نهایی به‌صورت ضرب و تفریق وارد می‌شوند و خطای کوچک در هر یک می‌تواند به خطای بزرگ در $R$ منجر شود. به‌ویژه، $R_{\text{asymp}}$ از تفریق دو کمیت نزدیک به هم به دست می‌آید و حساس به خطای سیستماتیک است. در صورتی که تلورانس $1\%$ برای $R_{\text{asymp}}$ رعایت شود، می‌توان با اطمینان گفت که پیاده‌سازی موتور، از نظر کریپتوگرافیک معتبر است.

\section{تحلیل معماری و ساختار داده‌ها}

\subsection{ساختار کلی ماژول و الگوی Module-Level Functions}

این ماژول از الگوی \lr{module-level functions} پیروی می‌کند، یعنی به‌جای تعریف کلاس‌ها، از توابع مستقل در سطح ماژول استفاده می‌نماید. این انتخاب طراحی، بازتاب‌دهنده ماهیت \lr{procedural} این ماژول است: توابع تحلیلی (\lr{\texttt{H2}}، \lr{\texttt{channel\_transmittance}}، \lr{\texttt{get\_analytics}}) توابع ریاضی خالص (\lr{pure functions}) هستند که ورودی‌هایشان را به خروجی‌هایشان نگاشت می‌کنند و هیچ \lr{state} داخلی ندارند. این الگو مزایای زیر را به همراه دارد: تست‌پذیری آسان (بدون نیاز به \lr{setup} و \lr{teardown}، هر تابع مستقل قابل آزمون است)، قابلیت ترکیب (خروجی یک تابع مستقیماً به ورودی تابع بعدی می‌رود) و قابلیت خواندن بالا (جریان کنترل از بالا به پایین واضح است).

تنها تابع دارای \lr{side-effect}، \lr{\texttt{run\_absolute\_test}} است که فایل \lr{CSV} می‌نویسد و گزارش چاپ می‌کند. این جداسازی بین توابع خالص و توابع دارای \lr{side-effect}، یکی از اصول مهندسی نرم‌افزار بالغ است و در الگوهای طراحی نظیر \lr{functional core, imperative shell} تبیین می‌شود. در این الگو، هسته محاسباتی (\lr{functional core}) از توابع خالص تشکیل می‌شود و لایه بیرونی (\lr{imperative shell}) مسئول \lr{I/O} و مدیریت \lr{side-effect} است.

\subsection{تحلیل تابع H2}

تابع \lr{\texttt{H2(p)}} یک تابع ریاضی خالص است که آنتروپی باینری شانون را محاسبه می‌نماید. این تابع با بررسی شرایط $p \leq 0$ یا $p \geq 1$، از خطای \lr{math domain error} در فراخوانی \lr{\texttt{math.log2}} جلوگیری می‌کند. در صورتی که $p$ خارج از بازه $(0, 1)$ باشد، مقدار $0$ بازگردانده می‌شود که با حد ریاضی هم‌خوان است. این پیاده‌سازی، الگوی \lr{defensive programming} را به نمایش می‌گذارد: ورودی‌ها پیش از استفاده در محاسبات، اعتبارسنجی می‌شوند و در صورت غیرمعتبر بودن، مقدار پیش‌فرض امن بازگردانده می‌شود.

\begin{LTR}
	\begin{lstlisting}[language=Python]
		def H2(p):
		if p <= 0 or p >= 1: return 0.0
		return -p * math.log2(p) - (1 - p) * math.log2(1 - p)
	\end{lstlisting}
\end{LTR}

در خط اول، شرط $p \leq 0$ یا $p \geq 1$ بررسی می‌شود. این شرط شامل نقاط انتهایی $p = 0$ و $p = 1$ نیز می‌شود که در آن‌ها آنتروپی به‌صورت تحلیلی صفر است. در خط دوم، فرمول اصلی $-p \log_2 p - (1-p) \log_2(1-p)$ پیاده‌سازی می‌شود. چون این خط فقط در صورتی اجرا می‌شود که $p \in (0, 1)$، فراخوانی \lr{\texttt{math.log2}} همیشه با ورودی مثبت انجام می‌شود و خطای \lr{math domain error} رخ نمی‌دهد.

\subsection{تحلیل تابع channel\_transmittance}

این تابع، ضریب انتقال کانال فیبر را بر اساس قانون \lr{Beer-Lambert} محاسبه می‌کند. این تابع نیز یک تابع خالص است که دو ورودی $L$ (طول فیبر به کیلومتر) و $\alpha_{\text{dB/km}}$ (ضریب کاهیدگی به \lr{dB/km}) را می‌گیرد و ضریب انتقال خطی را برمی‌گرداند.

\begin{LTR}
	\begin{lstlisting}[language=Python]
		def channel_transmittance(L, alpha_dB_km):
		return 10 ** (-alpha_dB_km * L / 10)
	\end{lstlisting}
\end{LTR}

این پیاده‌سازی ساده، مستقیماً فرمول $\eta_{\text{channel}} = 10^{-\alpha_{\text{dB/km}} \cdot L / 10}$ را پیاده می‌سازد. انتخاب نمادگذاری $10^{**}$ به‌جای \lr{\texttt{math.pow}} یا \lr{\texttt{exp}}، عمدی است: در پایتون، عملگر \lr{**} برای توان‌دهی عددی به‌کار می‌رود و خوانایی بالاتری دارد. در این فرمول، علامت منفی در نماد $-\alpha_{\text{dB/km}} \cdot L / 10$ بازتاب‌دهنده آن است که کاهیدگی باعث کاهش انتقال می‌شود و بنابراین، ضریب انتقال با افزایش طول فیبر، به‌صورت نمایی کاهش می‌یابد.

\subsection{تحلیل تابع get\_analytics و ساختار خروجی}

تابع \lr{\texttt{get\_analytics}} اصلی‌ترین تابع تحلیلی این ماژول است که پنج کمیت کلیدی کریپتوگرافیک را محاسبه می‌کند. این تابع هفت پارامتر ورودی دارد: $\mu$ (شدت سیگنال)، $\nu$ (شدت دکوی ضعیف)، $\eta$ (ضریب انتقال کلی)، $Y_0$ (نرخ دارک کانت)، $e_d$ (خطای ذاتی)، $f_{\text{EC}}$ (بازدهی \lr{error correction}) و $q_{\text{sift}}$ (کسر \lr{sifting}) که مقدار پیش‌فرض $0.5$ دارد. خروجی این تابع یک دیکشنری است که کلیدهای \lr{\texttt{Q\_mu}}، \lr{\texttt{E\_mu}}، \lr{\texttt{Y1}}، \lr{\texttt{e1}} و \lr{\texttt{R\_asymp}} را دارد.

\begin{LTR}
	\begin{lstlisting}[language=Python]
		def get_analytics(mu, nu, eta, Y0, ed, f_ec, q_sift=0.5):
		# 1. پارامترهای سیگنال (mu)
		Q_mu_raw = 1 - (1 - Y0) * math.exp(-eta * mu)
		Q_mu_sifted = Q_mu_raw * q_sift
		E_mu = (0.5 * Y0 + ed * (1 - math.exp(-eta * mu))) / Q_mu_raw
		
		# 2. پارامترهای فریب ضعیف (nu)
		Q_nu = 1 - (1 - Y0) * math.exp(-eta * nu)
	\end{lstlisting}
\end{LTR}

در خطوط ابتدایی، پارامترهای سیگنال محاسبه می‌شوند. $Q_\mu^{\text{raw}}$ از فرمول بسته محاسبه می‌شود؛ $Q_\mu^{\text{sifted}}$ با ضرب در $q_{\text{sift}}$ به دست می‌آید؛ $E_\mu$ از نسبت خطای کل به \lr{Gain} کل محاسبه می‌شود. در ادامه، $Q_\nu$ برای دکوی ضعیف به‌صورت مشابه محاسبه می‌گردد. این مقادیر، پایه محاسبات بعدی هستند.

\begin{LTR}
	\begin{lstlisting}[language=Python]
		# 3. استخراج کران پایین تک‌فوتونیتی (Lo2005)
		if abs(mu - nu) < 1e-9:
		Y1_L = 0
		else:
		term1 = Q_nu * math.exp(nu)
		term2 = (nu**2 / mu**2) * Q_mu_raw * math.exp(mu)
		factor = (mu**2 * math.exp(-mu)) / (mu * nu - mu**2)
		Y1_L = factor * (term1 - term2)
		
		Y1 = max(Y1_L, 0)
	\end{lstlisting}
\end{LTR}

در این بخش، کران پایین $Y_1^L$ محاسبه می‌شود. شرط $\left| \mu - \nu \right| < 10^{-9}$ حالت تکین را شناسایی می‌کند؛ در این حالت، فرمول به‌صورت تحلیلی نامحدود می‌شود و مقدار $Y_1^L = 0$ گرفته می‌شود. در غیر این صورت، فرمول اصلی \lr{Lo2005} پیاده می‌شود: \lr{\texttt{term1}} و \lr{\texttt{term2}} دو جمله اصلی صورت کسر هستند و \lr{\texttt{factor}} ضریب نرمال‌سازی است. در نهایت، \lr{\texttt{max(Y1\_L, 0)}} تضمین می‌کند که $Y_1$ منفی نشود؛ این برش عددی ضروری است، چون در برخی موارد (به‌ویژه وقتی $Q_\nu$ به‌دلیل نویز آماری کوچک‌تر از حد تئوریک باشد)، $Y_1^L$ می‌تواند منفی شود که فیزیکاً بی‌معناست.

\begin{LTR}
	\begin{lstlisting}[language=Python]
		# 4. محاسبه e1
		e1 = (0.5 * Y0 + ed * eta) / Y1 if Y1 > 0 else 0.0
		
		# 5. محاسبه نرخ کلید مجانبی
		Q1 = mu * math.exp(-mu) * Y1
		R_asymp = q_sift * (Q1 * (1 - H2(e1)) - Q_mu_raw * f_ec * H2(E_mu))
		R_asymp = max(R_asymp, 0.0)
	\end{lstlisting}
\end{LTR}

در این بخش، $e_1$، $Q_1$ و $R_{\text{asymp}}$ محاسبه می‌شوند. شرط $Y_1 > 0$ از تقسیم بر صفر جلوگیری می‌کند؛ در صورت $Y_1 = 0$، $e_1 = 0$ گرفته می‌شود. $Q_1$ از ضرب $\mu e^{-\mu}$ (احتمال پالس تک‌فوتونی) در $Y_1$ (بازده تک‌فوتونی) به دست می‌آید. $R_{\text{asymp}}$ از فرمول \lr{GLLP} محاسبه می‌شود و در نهایت، \lr{\texttt{max(R\_asymp, 0.0)}} تضمین می‌کند که نرخ کلید منفی نشود؛ این برش بازتاب‌دهنده آن است که در رژیم‌های با \lr{QBER} بالا، نرخ کلید به‌صورت تحلیلی منفی می‌شود که معنای فیزیکی آن این است که هیچ کلید امنی قابل استخراج نیست و در عمل، مقدار صفر گزارش می‌شود.

\begin{LTR}
	\begin{lstlisting}[language=Python]
		return {
			"Q_mu": Q_mu_sifted, "E_mu": E_mu, "Y1": Y1, "e1": e1, "R_asymp": R_asymp
		}
	\end{lstlisting}
\end{LTR}

خروجی تابع به‌صورت یک دیکشنری با پنج کلید است. این انتخاب طراحی، به دلیل خوانایی بالا و انعطاف‌پذیری در افزودن کمیت‌های جدید در آینده است. در صورتی که در نسخه‌های بعدی، کمیت‌های جدیدی نظیر $Y_0^L$ (کران پایین دارک کانت) یا $e_1^U$ (کران بالا خطای تک‌فوتونی) اضافه شوند، فقط یک کلید جدید به دیکشنری اضافه می‌شود و کدهای فراخوان نیاز به تغییر ندارند.

\subsection{تحلیل تابع run\_absolute\_test و معماری پیکربندی}

تابع \lr{\texttt{run\_absolute\_test}} اصلی‌ترین تابع آزمون است که کل جریان کار را هماهنگ می‌کند. این تابع چهار پارامتر ورودی دارد: $L$ (فاصله به کیلومتر)، $\mu$ (شدت سیگنال)، $\nu$ (شدت دکوی) و $f_{\text{EC}}$ (بازدهی \lr{error correction}). در ابتدا، ماژول \lr{\texttt{main\_optimized}} به‌صورت \lr{lazy} وارد می‌شود:

\begin{LTR}
	\begin{lstlisting}[language=Python]
		def run_absolute_test(dist_km, mu, nu, f_ec):
		try:
		import main_optimized as m
		except ImportError:
		print("Error: main_optimized not found.")
		return False
	\end{lstlisting}
\end{LTR}

این الگوی \lr{lazy import} دو مزیت دارد: نخست، در صورتی که ماژول به‌عنوان کتابخانه وارد شود (نه به‌عنوان اسکریپت اجرا شود)، وارد کردن \lr{\texttt{main\_optimized}} به تأخیر می‌افتد و زمان بارگذاری ماژول کاهش می‌یابد. دوم، در صورتی که \lr{\texttt{main\_optimized}} در دسترس نباشد (مثلاً در محیط \lr{CI} با وابستگی‌های ناقص)، خطا به‌صورت کنترل‌شده مدیریت می‌شود و تابع مقدار \lr{False} برمی‌گرداند.

سپس پیکربندی پیش‌فرض با \lr{\texttt{copy.deepcopy}} کپی می‌شود:

\begin{LTR}
	\begin{lstlisting}[language=Python]
		cfg = copy.deepcopy(m.DEFAULT_CONFIG)
		
		cfg["simulation"]["apply_statistical_noise"] = True
		cfg["simulation"]["distance_start_km"] = dist_km
		cfg["simulation"]["distance_stop_km"] = dist_km
		cfg["simulation"]["distance_points"] = 1
		cfg["simulation"]["min_pulses_log"] = 7
		cfg["simulation"]["max_pulses_log"] = 7
	\end{lstlisting}
\end{LTR}

استفاده از \lr{\texttt{deepcopy}} به‌جای انتساب مستقیم، حیاتی است: \lr{\texttt{DEFAULT\_CONFIG}} یک دیکشنری تودرتوی است و در صورت انتساب مستقیم، تغییرات روی کپی اعمال می‌شود و \lr{\texttt{DEFAULT\_CONFIG}} اصلی نیز تغییر می‌کند. این می‌تواند باعث \lr{regression} در آزمون‌های بعدی شود. \lr{\texttt{deepcopy}} یک کپی عمیق و مستقل ایجاد می‌کند که تغییرات روی آن، روی شیء اصلی اثر نمی‌گذارد. در ادامه، پیکربندی شبیه‌سازی تنظیم می‌شود: فقط یک نقطه فاصله (یعنی $L = $ \lr{\texttt{dist\_km}}) آزمون می‌شود، $N = 10^7$ پالس استفاده می‌شود و نویز آماری فعال است.

\begin{LTR}
	\begin{lstlisting}[language=Python]
		cfg["source"]["intended_proof"] = "LIM_2014"
		cfg["protocol_params"]["protocol"] = "BB84_DECOY"
		cfg["protocol_runtime"]["protocol_class"] = "BB84DecoyProtocol"
		cfg["source"]["intensity_jitter"] = 0.0
		
		cfg["source"]["pulses"] = {
			"signal": {"mu": mu, "prob": 0.5},
			"decoy": {"mu": nu, "prob": 0.25},
			"vacuum": {"mu": 0.0, "prob": 0.25}
		}
	\end{lstlisting}
\end{LTR}

در این بخش، پروتکل و اثبات امنیتی تنظیم می‌شوند. \lr{\texttt{LIM\_2014}} به‌عنوان اثبات امنیتی انتخاب می‌شود، زیرا این اثبات، نسخه متناهی-\lr{key} \lr{decoy-state BB84} است که در فریم‌ورک پیاده‌سازی شده. با این حال، چون فرمول‌های تحلیلی در تابع \lr{\texttt{get\_analytics}} حد مجانبی را محاسبه می‌کنند، آزمون در واقع کران بالای نرخ کلید را راستی‌آزمایی می‌کند. این یک انتخاب طراحی هوشمندانه است: اگر موتور شبیه‌سازی نتواند به حد مجانبی نزدیک شود، قطعاً نمی‌تواند نرخ کلید متناهی صحیح را تولید کند. در ادامه، توزیع پالس‌ها تنظیم می‌شود: $50\%$ سیگنال، $25\%$ دکوی و $25\%$ واکوئم. این توزیع، توزیع کلاسیک \lr{decoy-state} است که در کار \lr{Lim-2014} نیز به‌کار رفته. \lr{\texttt{intensity\_jitter = 0.0}} تضمین می‌کند که شدت پالس‌ها دقیقاً برابر با مقدار تنظیم‌شده باشد و نویز شدت اعمال نشود؛ این برای آزمون ضروری است، چون فرمول‌های تحلیلی فرض می‌کنند که شدت پالس‌ها دقیقاً $\mu$ و $\nu$ هستند.

\begin{LTR}
	\begin{lstlisting}[language=Python]
		eta_det = 0.15
		Y0 = 3e-6
		ed = 0.01
		cfg["detector"]["det_eff_d0"] = eta_det
		cfg["detector"]["det_eff_d1"] = eta_det
		cfg["detector"]["dark_rate"] = Y0
		cfg["detector"]["dark_rate_d1"] = Y0
		cfg["detector"]["qber_intrinsic"] = ed
		cfg["detector"]["misalignment"] = 0.0
		cfg["detector"]["dead_time_ns"] = 0.0
	\end{lstlisting}
\end{LTR}

پارامترهای آشکارساز تنظیم می‌شوند. $\eta_{\text{det}} = 0.15$، $Y_0 = 3 \times 10^{-6}$ و $e_d = 0.01$ انتخاب شده‌اند. این مقادیر با مقادیر مرجع در ادبیات \lr{QKD} هم‌خوان هستند. \lr{\texttt{misalignment = 0.0}} و \lr{\texttt{dead\_time\_ns = 0.0}} تنظیم شده‌اند تا اثرات اضافی حذف شوند و آزمون فقط روی فیزیک اصلی متمرکز باشد. این رویکرد \lr{minimal configuration}، تضمین می‌کند که هرگونه انحراف از فرمول تحلیلی، ناشی از باگ در موتور باشد نه از اثرات فیزیکی اضافی.

\subsection{الگوی Comparison و Tolerance-Based Validation}

قلب تابع \lr{\texttt{run\_absolute\_test}}، تابع داخلی \lr{\texttt{check}} است که مقادیر شبیه‌سازی را با مقادیر تحلیلی مقایسه می‌کند:

\begin{LTR}
	\begin{lstlisting}[language=Python]
		def check(name, fw_val, th_val, tol):
		if th_val == 0: 
		err = abs(fw_val - th_val)
		passed = err < 1e-9
		else:
		err = abs(fw_val - th_val) / abs(th_val) * 100
		passed = err <= tol
		print(f"{name:<15} | {fw_val:<15.6e} | {th_val:<15.6e} | {err:<10.4f} | {tol:<5} | {'PASS' if passed else 'FAIL'}")
		return passed
	\end{lstlisting}
\end{LTR}

این تابع از الگوی \lr{relative tolerance} پیروی می‌کند: در صورتی که مقدار تحلیلی غیرصفر باشد، خطای نسبی به درصد محاسبه می‌شود و با تلورانس مقایسه می‌گردد. در صورتی که مقدار تحیلی صفر باشد، خطای مطلق با آستانه $10^{-9}$ مقایسه می‌شود. این دوگانه‌سازی ضروری است، چون در صورتی که مقدار تحلیلی صفر باشد، خطای نسبی تعریف‌نشده (\lr{division by zero}) است. آستانه $10^{-9}$ برای حالت صفر، بسیار سخت‌گیرانه است و تضمین می‌کند که در این حالت، فقط مقادیر عملاً صفر قابل قبول هستند.

خروجی این تابع، یک خط فرمت‌شده با ستون‌های نام پارامتر، مقدار شبیه‌سازی، مقدار تحلیلی، خطا، تلورانس و نتیجه \lr{PASS}/\lr{FAIL} است. این قالب، خوانایی بالایی دارد و امکان تحلیل سریع نتایج را فراهم می‌سازد. در نسخه اصلی کد، از ایموجی \lr{✅} و \lr{❌} برای نمایش نتیجه استفاده شده، اما در پیاده‌سازی مستند، برای سازگاری با خروجی کنسول استاندارد، می‌توان از \lr{PASS} و \lr{FAIL} استفاده کرد.

\subsection{ساختار پنج‌آزمونی و خروجی نهایی}

پنج آزمون به ترتیب با تلورانس‌های $25\%$، $5\%$، $1\%$، $1\%$ و $1\%$ اجرا می‌شوند:

\begin{LTR}
	\begin{lstlisting}[language=Python]
		p1 = check("Q_mu", fw_Q_mu, th["Q_mu"], 25.0)      
		p2 = check("E_mu", fw_E_mu, th["E_mu"], 5.0)
		p3 = check("Y1 (Extraction)", fw_Y1, th["Y1"], 1.0)  
		p4 = check("e1 (Extraction)", fw_e1, th["e1"], 1.0)
		p5 = check("SKR (Absolute)", fw_R, th["R_asymp"], 1.0) 
		
		all_passed = all([p1, p2, p3, p4, p5])
		if all_passed:
		print("100% CRYPTOGRAPHIC VALIDATION PASSED!")
		else:
		print("MATH/CRYPTO VALIDATION FAILED. CHECK FRAMEWORK ENGINE.")
		
		return all_passed
	\end{lstlisting}
\end{LTR}

در این بخش، پنج آزمون به ترتیب اجرا می‌شوند. تابع \lr{\texttt{all}} بررسی می‌کند که آیا همه نتایج \lr{True} هستند یا خیر. در صورت موفقیت همه، پیام \lr{100\% CRYPTOGRAPHIC VALIDATION PASSED!} چاپ می‌شود و در غیر این صورت، پیام خطا چاپ می‌گردد. خروجی تابع، یک مقدار بولین است که در \lr{\texttt{\_\_main\_\_}} برای تعیین کد خروج استفاده می‌شود. این الگو، تضمین می‌کند که در \lr{CI/CD pipeline}، کد خروج به‌درستی تنظیم شود و در صورت شکست، \lr{pipeline} متوقف گردد.

\subsection{تحلیل نقطه ورود \_\_main\_\_}

\begin{LTR}
	\begin{lstlisting}[language=Python]
		if __name__ == "__main__":
		print("=" * 80)
		print("ABSOLUTE CRYPTOGRAPHIC & MATHEMATICAL VALIDATION SUITE")
		print("Testing Native Framework Engine (No Patching) vs Analytical Ground Truth")
		print("=" * 80)
		
		is_valid = run_absolute_test(dist_km=50.0, mu=0.5, nu=0.1, f_ec=1.1)
		sys.exit(0 if is_valid else 1)
	\end{lstlisting}
\end{LTR}

نقطه ورود \lr{\texttt{\_\_main\_\_}} یک سرآیند تزئینی چاپ می‌کند و سپس تابع \lr{\texttt{run\_absolute\_test}} را با پارامترهای پیش‌فرض $L = 50$ کیلومتر، $\mu = 0.5$، $\nu = 0.1$ و $f_{\text{EC}} = 1.1$ فراخوانی می‌کند. این پارامترها، نقطه عملیاتی متعارف برای \lr{QKD} در فیبر نوری هستند و انتخاب آن‌ها بر اساس کار مرجع \lr{Lim-2014} است. در نهایت، \lr{\texttt{sys.exit}} با کد خروج $0$ (در صورت موفقیت) یا $1$ (در صورت شکست) فراخوانی می‌شود. این الگو، استاندارد \lr{POSIX} برای اسکریپت‌های \lr{CI/CD} است و تضمین می‌کند که \lr{shell} بتواند کد خروج را به‌درستی تفسیر کند.

\section{پیاده‌سازی ریاضی و الگوریتمی}

\subsection{پیاده‌سازی آنتروپی باینری با برش عددی}

پیاده‌سازی تابع $H_2$ در این ماژول، نمونه‌ای از \lr{defensive programming} در محاسبات عددی است. فرمول تحلیلی $-p \log_2 p - (1-p) \log_2(1-p)$ در نقاط $p = 0$ و $p = 1$ به‌صورت تحلیلی به $0$ میل می‌کند، اما در محاسبات عددی، فراخوانی \lr{\texttt{math.log2(0)}} منجر به \lr{ValueError} می‌شود. برای حل این مشکل، شرط $p \leq 0$ یا $p \geq 1$ پیش از فراخوانی \lr{\texttt{math.log2}} بررسی می‌شود و در این موارد، مقدار $0$ بازگردانده می‌شود.

این برش عددی با حد گسسته \lr{L'Hôpital} هم‌خوان است:

\begin{equation}
	\lim_{p \to 0^+} p \log_2 p = \lim_{p \to 0^+} \frac{\log_2 p}{1/p} = \lim_{p \to 0^+} \frac{1/(p \ln 2)}{-1/p^2} = \lim_{p \to 0^+} \frac{-p}{\ln 2} = 0
\end{equation}

بنابراین، پیاده‌سازی عددی با حد تحلیلی هم‌خوان است و از نظر ریاضی صحیح است. این الگو در بسیاری از توابع ریاضی در علوم کامپیوتر به‌کار می‌رود و به‌عنوان \lr{numerical stability via boundary clipping} شناخته می‌شود.

\subsection{پیاده‌سازی فرمول Lo2005 برای Y1}

پیاده‌سازی فرمول $Y_1^L$ در تابع \lr{\texttt{get\_analytics}} چندین نکته مهم دارد. نخست، شرط $\left| \mu - \nu \right| < 10^{-9}$ حالت تکین را شناسایی می‌کند. این شرط بازتاب‌دهنده آن است که در فرمول $Y_1^L$، مخرج $\mu \nu - \mu^2 = \mu(\nu - \mu)$ در صورتی که $\nu \to \mu$ باشد، به صفر میل می‌کند و فرمول به‌صورت تحلیلی نامحدود می‌شود. در عمل، این حالت رخ نمی‌دهد (چون $\nu$ و $\mu$ معمولاً بسیار متفاوت هستند)، اما برای پایداری عددی، شرط بررسی می‌شود.

دوم، ساختار محاسبه به سه مرحله تفکیک شده است: محاسبه \lr{\texttt{term1}} و \lr{\texttt{term2}} (دو جمله اصلی صورت) و محاسبه \lr{\texttt{factor}} (ضریب نرمال‌سازی). این تفکیک، خوانایی کد را بالا می‌برد و امکان \lr{debug} آسان‌تر را فراهم می‌سازد. در صورتی که مقدار $Y_1^L$ منفی شود (که می‌تواند در شرایط نویز آماری رخ دهد)، \lr{\texttt{max(Y1\_L, 0)}} تضمین می‌کند که $Y_1$ نهایی غیرمنفی باشد. این برش، بازتاب‌دهنده مفهوم فیزیکی $Y_1$ به‌عنوان احتمال است که نمی‌تواند منفی باشد.

\subsection{پیاده‌سازی فرمول GLLP برای نرخ کلید}

پیاده‌سازی فرمول \lr{GLLP} برای $R_{\text{asymp}}$ مستقیم است:

\begin{equation}
	R_{\text{asymp}} = q_{\text{sift}} \cdot \left[ Q_1 \cdot (1 - H_2(e_1)) - Q_\mu^{\text{raw}} \cdot f_{\text{EC}} \cdot H_2(E_\mu) \right]
\end{equation}

در کد، این فرمول به‌صورت زیر پیاده می‌شود:

\begin{LTR}
	\begin{lstlisting}[language=Python]
		Q1 = mu * math.exp(-mu) * Y1
		R_asymp = q_sift * (Q1 * (1 - H2(e1)) - Q_mu_raw * f_ec * H2(E_mu))
		R_asymp = max(R_asymp, 0.0)
	\end{lstlisting}
\end{LTR}

نکته مهم در این پیاده‌سازی، استفاده از \lr{\texttt{Q\_mu\_raw}} (نه \lr{\texttt{Q\_mu\_sifted}}) در فرمول است. دلیل آن آن است که در فرمول \lr{GLLP}، \lr{Gain} بدون \lr{sifting} به‌کار می‌رود، چون \lr{sifting} در سطح کل فرمول (یعنی در $q_{\text{sift}}$ بیرونی) اعمال می‌شود. این یک نکته ظریف است که در پیاده‌سازی‌های نادرست به‌کررات نادیده گرفته می‌شود و می‌تواند منجر به نرخ کلید اشتباه شود.

برش \lr{\texttt{max(R\_asymp, 0.0)}} تضمین می‌کند که نرخ کلید منفی نشود. این برش در رژیم‌های با \lr{QBER} بالا (مثلاً $E_\mu > 0.11$ که آستانه بحرانی \lr{BB84} است) ضروری است، چون در این رژیم، عبارت $Q_\mu^{\text{raw}} \cdot f_{\text{EC}} \cdot H_2(E_\mu)$ بزرگ‌تر از $Q_1 \cdot (1 - H_2(e_1))$ می‌شود و $R_{\text{asymp}}$ منفی می‌گردد. در عمل، نرخ کلید منفی به این معناست که هیچ کلید امنی قابل استخراج نیست و مقدار صفر گزارش می‌شود.

\subsection{الگوریتم Lazy Import و مدیریت خطا}

پیاده‌سازی \lr{lazy import} برای \lr{\texttt{main\_optimized}} الگوی مهمی در مهندسی نرم‌افزار است. این الگو به دو دلیل به‌کار می‌رود: نخست، زمان بارگذاری ماژول را کاهش می‌دهد (در صورتی که ماژول به‌عنوان کتابخانه وارد شود و تابع \lr{\texttt{run\_absolute\_test}} فراخوانی نشود، \lr{\texttt{main\_optimized}} بارگذاری نمی‌شود). دوم، در صورت عدم دسترسی به \lr{\texttt{main\_optimized}}، خطا به‌صورت کنترل‌شده مدیریت می‌شود و تابع مقدار \lr{False} برمی‌گرداند.

\begin{LTR}
	\begin{lstlisting}[language=Python]
		try:
		import main_optimized as m
		except ImportError:
		print("Error: main_optimized not found.")
		return False
	\end{lstlisting}
\end{LTR}

این الگو، نمونه‌ای از \lr{graceful degradation} است: در صورت عدم دسترسی به یک وابستگی، سیستم به‌جای \lr{crash} کردن، رفتار قابل پیش‌بینی نشان می‌دهد. این الگو در محیط‌های \lr{CI} که ممکن است وابستگی‌ها ناقص باشند، حیاتی است.

\subsection{الگوریتم Deep Copy و جلوگیری از Mutation}

استفاده از \lr{\texttt{copy.deepcopy}} برای کپی \lr{\texttt{DEFAULT\_CONFIG}} یک الگوی حیاتی در پایتون است. در پایتون، انتساب دیکشنری (\lr{\texttt{cfg = m.DEFAULT\_CONFIG}}) فقط یک مرجع (\lr{reference}) جدید ایجاد می‌کند و تغییرات روی \lr{\texttt{cfg}} روی \lr{\texttt{DEFAULT\_CONFIG}} اصلی نیز اعمال می‌شود. این می‌تواند منجر به \lr{regression} در آزمون‌های بعدی شود، چون \lr{\texttt{DEFAULT\_CONFIG}} دیگر \emph{پیش‌فرض} نیست.

\begin{LTR}
	\begin{lstlisting}[language=Python]
		cfg = copy.deepcopy(m.DEFAULT_CONFIG)
	\end{lstlisting}
\end{LTR}

\lr{\texttt{deepcopy}} یک کپی عمیق و مستقل ایجاد می‌کند: تمام اشیای تودرتوی (دیکشنری‌ها، لیست‌ها و غیره) به‌صورت بازگشتی کپی می‌شوند و شیء جدید، کاملاً مستقل از شیء اصلی است. این الگو در فریم‌ورک‌های آزمون که در آن‌ها پیکربندی به‌کررات تغییر می‌کند، استاندارد است. هزینه \lr{\texttt{deepcopy}} از مرتبه $O(n)$ است که در اینجا $n$ تعداد کلیدهای دیکشنری است و قابل صرف‌نظر است.

\subsection{الگوریتم Relative vs Absolute Tolerance}

پیاده‌سازی تابع \lr{\texttt{check}} از الگوریتم دوگانه \lr{relative}/\lr{absolute tolerance} استفاده می‌کند:

\begin{LTR}
	\begin{lstlisting}[language=Python]
		if th_val == 0: 
		err = abs(fw_val - th_val)
		passed = err < 1e-9
		else:
		err = abs(fw_val - th_val) / abs(th_val) * 100
		passed = err <= tol
	\end{lstlisting}
\end{LTR}

این الگوریتم دو حالت را تفکیک می‌کند: در صورتی که مقدار تحلیلی صفر باشد، خطای مطلق با آستانه $10^{-9}$ مقایسه می‌شود (چون خطای نسبی در این حالت تعریف‌نشده است). در غیر این صورت، خطای نسبی به درصد محاسبه می‌شود و با تلورانس مقایسه می‌گردد. این الگوریتم، الگوی استاندارد در آزمون‌های عددی است و در کتابخانه‌هایی نظیر \lr{NumPy} (\lr{\texttt{numpy.isclose}}) و \lr{Python math.isclose} نیز به‌کار می‌رود. تفاوت در این است که در این پیاده‌سازی، کاربر می‌تواند تلورانس را به‌صورت صریح تنظیم کند، در حالی که در \lr{\texttt{math.isclose}}، تلورانس‌های پیش‌فرض $10^{-9}$ برای نسبی و $10^{-15}$ برای مطلق هستند.

\subsection{الگوریتم Comparison Pipeline و ساختار CSV}

پس از اجرای شبیه‌سازی، خروجی از فایل \lr{CSV} خوانده می‌شود:

\begin{LTR}
	\begin{lstlisting}[language=Python]
		with open("qkd_results_optimized_sweep.csv", "r") as f:
		rows = list(csv.DictReader(f))
		latest = rows[-1]
		
		fw_Q_mu = float(latest.get("detection_yield_signal", 0))
		fw_E_mu = float(latest.get("qber_signal", 0))
		fw_Y1 = float(latest.get("single_photon_yield", 0))
		fw_e1 = float(latest.get("single_photon_qber", 0))
		fw_R = float(latest.get("secure_key_rate", 0))
	\end{lstlisting}
\end{LTR}

استفاده از \lr{\texttt{csv.DictReader}} به‌جای \lr{\texttt{csv.reader}}، خوانایی را بالا می‌برد: دسترسی به ستون‌ها با نام به‌جای ایندکس. \lr{\texttt{rows[-1]}} آخرین ردیف را برمی‌گرداند که در صورتی که فقط یک نقطه فاصله آزمون شده باشد، همان ردیف آزمون است. در صورتی که فایل دارای چندین ردیف باشد (مثلاً از اجرای قبلی)، این رویکرد تضمین می‌کند که آخرین نتیجه استفاده شود. استفاده از \lr{\texttt{.get(key, 0)}} با مقدار پیش‌فرض $0$، در صورتی که ستون وجود نداشته باشد، از خطا جلوگیری می‌کند. این الگوی \lr{defensive programming} در خواندن \lr{CSV} که ممکن است ساختار آن بین نسخه‌های فریم‌ورک متغیر باشد، حیاتی است.

\subsection{محاسبه ضریب انتقال برای مقایسه}

پس از خواندن خروجی شبیه‌سازی، ضریب انتقال تحلیلی محاسبه می‌شود:

\begin{LTR}
	\begin{lstlisting}[language=Python]
		eta = channel_transmittance(dist_km, 0.2) * eta_det
		th = get_analytics(mu, nu, eta, Y0, ed, f_ec)
	\end{lstlisting}
\end{LTR}

در این خط، $\eta_{\text{channel}}$ با $\alpha_{\text{dB/km}} = 0.2$ (مقدار استاندارد فیبر تک‌حالت در $1550$ نانومتر) محاسبه می‌شود و سپس در $\eta_{\text{det}} = 0.15$ ضرب می‌گردد تا $\eta_{\text{total}}$ به دست آید. سپس تابع \lr{\texttt{get\_analytics}} با این $\eta$ فراخوانی می‌شود و مقادیر تحلیلی محاسبه می‌گردند. این الگو، تضمین می‌کند که مقایسه بین شبیه‌سازی و تحلیلی، در همان شرایط فیزیکی انجام شود.

\section{ارسال و دریافت با سایر ماژول‌ها}

\subsection{وابستگی به main\_optimized}

اصلی‌ترین وابستگیِ این ماژول، وابستگی به \lr{\texttt{main\_optimized.py}} است. این وابستگی در سه بخش نمود می‌یابد: نخست، \lr{\texttt{DEFAULT\_CONFIG}} به‌عنوان نقطهٔ آغاز پیکربندی به کار می‌رود. دوم، تابع \lr{\texttt{run\_and\_save\_csv}} برای اجرای شبیه‌سازی فراخوانی می‌شود. سوم، فایل \lr{\texttt{CSV}} با نام \lr{\texttt{qkd\_results\_optimized\_sweep.csv}} که توسط ماژول \lr{\texttt{main\_optimized}} تولید می‌شود، خوانده می‌شود.

این وابستگی، یک وابستگی \lr{tight coupling} است: تغییر در ساختار \lr{\texttt{DEFAULT\_CONFIG}} یا نام ستون‌های \lr{CSV} می‌تواند منجر به شکست این ماژول شود. برای کاهش این \lr{coupling}، در پیاده‌سازی‌های آینده می‌توان از الگوهای زیر استفاده کرد: تعریف یک \lr{schema} صریح برای ستون‌های \lr{CSV} (به‌عنوان مثال، در قالب \lr{dataclass} یا \lr{TypedDict})، یا استفاده از \lr{API} تابعی به‌جای فایل \lr{CSV} (یعنی تابع \lr{\texttt{run\_and\_save\_csv}} خروجی را به‌صورت یک شیء پایتونی برگرداند).

\subsection{دیکشنری DEFAULT\_CONFIG و Override Pattern}

\lr{\texttt{DEFAULT\_CONFIG}} یک دیکشنری تودرتوی است که در \lr{\texttt{main\_optimized.py}} تعریف شده و شامل تمام پارامترهای شبیه‌سازی است: تنظیمات شبیه‌سازی (\lr{\texttt{simulation}})، تنظیمات منبع (\lr{\texttt{source}})، تنظیمات آشکارساز (\lr{\texttt{detector}})، تنظیمات پروتکل (\lr{\texttt{protocol\_params}}) و غیره. این ماژول با \lr{\texttt{deepcopy}} یک کپی مستقل ایجاد می‌کند و سپس مقادیر مورد نظر را \lr{override} می‌نماید.

این الگوی \lr{configuration override} در فریم‌ورک‌های شبیه‌سازی استاندارد است و امکان آزمون آسان در شرایط مختلف را فراهم می‌سازد. مزیت اصلی آن آن است که فقط پارامترهای مورد نیاز برای آزمون تغییر می‌کنند و سایر پارامترها در مقدار پیش‌فرض باقی می‌مانند؛ این تضمین می‌کند که آزمون در شرایط نزدیک به واقعی انجام شود و در عین حال، تکرارپذیر باشد.

\subsection{خروجی CSV و اسکیمای ستون‌ها}

فایل \lr{CSV} که توسط \lr{\texttt{main\_optimized}} تولید می‌شود، شامل ستون‌های متعددی است که نتایج شبیه‌سازی را در بر دارد. در این ماژول، پنج ستون کلیدی خوانده می‌شوند:

\begin{itemize}
	\item \lr{\texttt{detection\_yield\_signal}}: \lr{Gain} سیگنال sifted شده ($Q_\mu^{\text{sifted}}$)
	\item \lr{\texttt{qber\_signal}}: \lr{QBER} سیگنال ($E_\mu$)
	\item \lr{\texttt{single\_photon\_yield}}: بازده تک‌فوتونی ($Y_1$)
	\item \lr{\texttt{single\_photon\_qber}}: خطای تک‌فوتونی ($e_1$)
	\item \lr{\texttt{secure\_key\_rate}}: نرخ کلید امن ($R$)
\end{itemize}

این اسکیما، در فریم‌ورک به‌عنوان \lr{contract} بین موتور شبیه‌سازی و ابزارهای تحلیل عمل می‌کند. در صورتی که ساختار \lr{CSV} تغییر کند (مثلاً ستون‌ها اضافه یا حذف شوند)، این ماژول و سایر ابزارهای تحلیل باید به‌روز شوند. برای کاهش این وابستگی، می‌توان از \lr{schema versioning} استفاده کرد: یک ستون \lr{\texttt{schema\_version}} به \lr{CSV} اضافه شود و ابزارهای تحلیل بر اساس نسخه، رفتار متفاوت داشته باشند.

\subsection{ارتباط با validate\_physics.py}

این ماژول و \lr{\texttt{validate\_physics.py}} دو لایه مکمل اعتبارسنجی را تشکیل می‌دهند. \lr{\texttt{validate\_physics.py}} بر آزمون آماری توزیع‌های منبع متمرکز است و از \lr{Z-test} برای راستی‌آزمایی \lr{Gain} و \lr{QBER} استفاده می‌کند. این ماژول (\lr{\texttt{absolute\_crypto\_validation.py}}) بر صفت کریپتوگرافیک پارامترهای کلیدی متمرکز است و از تلورانس عددی برای راستی‌آزمایی $Y_1$، $e_1$ و $R$ استفاده می‌نماید. این دو لایه، نقش‌های مکمل دارند: لایه اول (\lr{validate\_physics}) تضمین می‌کند که فیزیک پایه (توزیع \lr{Poisson} فوتون، دارک کانت و غیره) به‌درستی پیاده شده، و لایه دوم (این ماژول) تضمین می‌کند که استخراج کریپتوگرافیک (\lr{decoy-state} estimation، \lr{privacy amplification} و غیره) به‌درستی انجام شده.

\subsection{ارتباط با ماژول‌های اثبات Lim-2014 و Ma-2005}

این ماژول به‌طور غیرمستقیم با ماژول‌های اثبات امنیتی \lr{\texttt{lim2014.py}} و \lr{\texttt{ma2005.py}} در ارتباط است. از طریق \lr{\texttt{main\_optimized}}، این ماژول اثبات \lr{Lim-2014} را برای اجرای شبیه‌سازی استفاده می‌کند. با این حال، فرمول‌های تحلیلی در تابع \lr{\texttt{get\_analytics}} بر اساس حد مجانبی هستند (یعنی $N \to \infty$) که با فرمول‌های متناهی-\lr{key} در \lr{Lim-2014} متفاوت است. این تفاوت، عمدی است: آزمون کران بالای نرخ کلید را راستی‌آزمایی می‌کند و در صورتی که موتور شبیه‌سازی نتواند به کران بالا نزدیک شود، قطعاً نمی‌تواند نرخ کلید متناهی صحیح را تولید کند. این یک الگوی \lr{necessary condition testing} است که در آزمون‌های عددی به‌کررات به‌کار می‌رود.

\subsection{خروجی نهایی و استفاده در CI/CD}

خروجی نهایی این ماژول به دو شکل است: نخست، خروجی چاپی به کنسول که شامل جدول مقایسه و نتیجه \lr{PASS}/\lr{FAIL} است. دوم، کد خروج \lr{POSIX} ($0$ برای موفقیت، $1$ برای شکست) که توسط \lr{CI/CD pipeline} قابل تفسیر است. این الگو، استاندارد \lr{shell scripting} است و در ابزارهایی نظیر \lr{GitHub Actions}، \lr{GitLab CI} و \lr{Jenkins} به‌طور گسترده پشتیبانی می‌شود.

در یک \lr{CI/CD pipeline} نمونه، این ماژول به‌صورت زیر به‌کار می‌رود:

\begin{LTR}
	\begin{lstlisting}[language=bash]
		- name: Cryptographic Validation
		run: python absolute_crypto_validation.py
	\end{lstlisting}
\end{LTR}

در صورت شکست، \lr{pipeline} متوقف می‌شود و توسعه‌دهنده باید قبل از \lr{merge}، خطا را اصلاح کند. این الگو، تضمین می‌کند که هرگز کدی با خطای کریپتوگرافیک به شاخه اصلی \lr{merge} نشود. در ادبیات مهندسی نرم‌افزار علمی، این الگو به‌عنوان \lr{cryptographic regression testing} شناخته می‌شود و در پروژه‌های بالغی مانند \lr{Open Quantum Safe} و \lr{Qiskit} به‌صورت استاندارد به کار می‌رود.

\subsection{جمع‌بندی ارتباطات}

به‌طور خلاصه، این ماژول در نقش یک \lr{validator} مستقل عمل می‌کند که از طریق سه کانال با سایر بخش‌های فریم‌ورک در ارتباط است: ورودی پیکربندی از \lr{\texttt{main\_optimized.DEFAULT\_CONFIG}}، اجرای شبیه‌سازی از طریق \lr{\texttt{main\_optimized.run\_and\_save\_csv}}، و خواندن نتایج از فایل \lr{\texttt{qkd\_results\_optimized\_sweep.csv}}. خروجی آن به‌صورت کد خروج \lr{POSIX} و گزارش چاپی، برای مصرف توسط \lr{CI/CD pipeline} و توسعه‌دهنده طراحی شده است. این معماری، جداسازی واضحی بین \lr{validator} و موتور شبیه‌سازی برقرار می‌سازد و امکان توسعه مستقل هر یک را فراهم می‌نماید.

	
	\chapter{اعتبارسنجی حل‌گر برنامه‌ریزی خطی}
	\label{ch:val_lp}
	\sloppy
\emergencystretch=5em
\hbadness=10000
\vbadness=10000
\tolerance=2000

\providecolor{codebg}{RGB}{248,248,248}
\providecolor{codeframe}{RGB}{200,200,200}
\providecolor{keyword}{RGB}{0,0,180}
\providecolor{string}{RGB}{163,21,21}
\providecolor{comment}{RGB}{0,128,0}

\lstset{
	backgroundcolor=\color{codebg},
	frame=single,
	rulecolor=\color{codeframe},
	language=Python,
	basicstyle=\ttfamily\scriptsize,
	keywordstyle=\color{keyword}\bfseries,
	stringstyle=\color{string},
	commentstyle=\color{comment}\itshape,
	showstringspaces=false,
	breaklines=true,
	breakatwhitespace=true,
	numbers=left,
	numberstyle=\tiny\color{gray},
	numbersep=8pt,
	tabsize=4,
	literate=
	{ف}{{{\bfseries\itshape ف}}}1
	{ی}{{{\bfseries\itshape ی}}}1
}

\section{نمای کلی و هدف ماژول}

ماژول \lr{\texttt{lp\_validation.py}} یک اسکریپت اعتبارسنجی \emph{بهینه‌سازی پارامترهای پروتکل} است که نقش یک \lr{validator} مستقل را در فریم‌ورک شبیه‌سازی \lr{QKD} ایفا می‌کند. این ماژول به‌طور خاص برای پاسخ به پرسش بنیادین زیر طراحی شده است: \emph{آیا حلّال بهینه‌سازی پارامترهای پروتکل \lr{Lim-2014} (\lr{LP Solver}) که در ماژول \lr{\texttt{qkd.proofs.optimization}} پیاده‌سازی شده، در یافتن پارامترهای بهینه (\lr{optimal}) به‌طور قابل‌اتکا عمل می‌کند و آیا خروجی آن حداقل به‌اندازه پارامترهای سخت‌کدشده (\lr{hardcoded}) و مرجع، نرخ کلید امن (\lr{Secure Key Rate} یا \lr{SKR}) تولید می‌نماید؟} این پرسش از آنجا اهمیت می‌یابد که در پروتکل‌های \lr{decoy-state BB84}، انتخاب پارامترهای شدت پالس (\lr{$\mu$} و \lr{$\nu$}) و احتمال‌های تخصیص (\lr{$p_s$} و \lr{$p_d$}) تأثیر مستقیمی بر نرخ کلید نهایی دارد و یک بهینه‌ساز ناقص می‌تواند منجر به انتخاب پارامترهای زیربهینه (\lr{suboptimal}) شود که نرخ کلیدی کم‌تر از حد قابل دستیابی تولید می‌کنند.

هدف اصلی این ماژول، \emph{اعتبارسنجی تطبیقی} (\lr{comparative validation}) حلّال بهینه‌سازی در برابر یک مرجع سخت‌کدشده است. به‌عبارت دیگر، این ماژول یک نقطه عملیاتی مرجع با پارامترهای ثابت \lr{$\mu = 0.5$}، \lr{$\nu = 0.1$}، \lr{$q_x = 0.5$}، \lr{$p_s = 0.5$} و \lr{$p_d = 0.25$} (که در ادبیات \lr{decoy-state QKD} به‌عنوان مقادیر کلاسیک و شناخته‌شده هستند) تعریف می‌کند، سپس حلّال بهینه‌سازی را فراخوانی می‌نماید و در صورتی که \lr{SKR} بهینه حداقل برابر با \lr{SKR} مرجع باشد، آزمون موفق تلقی می‌گردد. این رویکرد، الگوی \lr{necessary condition testing} است: موفقیت آزمون تضمین می‌کند که حلّال حداقل به اندازه مرجع عمل می‌کند؛ هرچند ممکن است پارامترهای بهتر نیز یافت کند که در آن صورت، آزمون با حاشیه مثبت موفق می‌شود.

ساختار کلی این ماژول شامل سه تابع است. تابع \lr{\texttt{H2}} آنتروپی باینری شانون را با برش عددی (\lr{numerical clipping}) پیاده‌سازی می‌کند تا از خطای \lr{math domain error} در نقاط \lr{$p = 0$} و \lr{$p = 1$} جلوگیری نماید. تابع \lr{\texttt{analytical\_skr}} تابع هدف تحلیلی را محاسبه می‌کند که از حلّال بهینه‌سازی به‌عنوان \lr{cost function} استفاده می‌گردد؛ این تابع، پنج پارامتر پروتکل را به‌عنوان ورودی می‌گیرد و \lr{SKR} مجانبی را بر اساس فرمول \lr{GLLP} (\lr{Gottesman-Lo-Lütkenhaus-Preskill}) بازمی‌گرداند. تابع \lr{\texttt{run\_lp\_test}} آزمون اصلی است که در یک فاصله مشخص، حلّال را فراخوانی کرده و نتایج را مقایسه می‌نماید.

ویژگی برجسته این ماژول، \emph{استفاده از تابع callback} برای تزریق تابع هدف است. به‌جای آنکه حلّال بهینه‌سازی به‌صورت سخت‌کدشده به فرمول \lr{SKR} وابسته باشد، تابع \lr{\texttt{cost\_callback}} به‌عنوان پارامتر به حلّال ارسال می‌شود. این الگو، الگوی \lr{dependency injection} است و امکان جایگزینی تابع هدف با توابع دیگر (مثلاً \lr{SKR} متناهی-\lr{key} یا توابع هزینه پیچیده‌تر) را بدون تغییر در حلّال فراهم می‌سازد. این طراحی، جداسازی واضحی بین منطق بهینه‌سازی و تابع هدف برقرار می‌کند و ماژولار بودن سیستم را افزایش می‌دهد.

در سطح عملیاتی، این ماژول سه آزمون در فواصل مختلف ($0$، $50$ و $100$ کیلومتر) اجرا می‌کند تا عملکرد حلّال در رژیم‌های مختلف کانال (بدون کاهیدگی، کاهیدگی متوسط و کاهیدگی شدید) را راستی‌آزمایی نماید. در فاصله $L = 0$ کیلومتر، ضریب انتقال کانال برابر با $1$ است و حلّال باید پارامترهایی با \lr{SKR} نزدیک به حد نظری پیدا کند؛ در فاصله $L = 50$ کیلومتر، ضریب انتقال به $\eta_{\text{channel}} \approx 0.1$ کاهش می‌یابد و حلّال باید با کاهش $\mu$ برای جبران کاهیدگی، نرخ کلید را بهینه نماید؛ در فاصله $L = 100$ کیلومتر، با ضریب انتقال $\eta_{\text{channel}} \approx 0.01$، حلّال در رژیم نزدیک به آستانه بحرانی \lr{BB84} (\lr{$Q \approx 0.11$}) عمل می‌کند و انتخاب پارامترها به‌طور حساس به نتیجه تأثیر می‌گذارد. این سه آزمون، سه رژیم فیزیکی متفاوت را پوشش می‌دهند و تضمین می‌کنند که حلّال در طیف وسیعی از شرایط عملیاتی به‌طور قابل‌اتکا عمل می‌کند.

نکته مهم دیگر آن است که این ماژول \emph{یک تابع هدف تحلیلی مستقل} برای اعتبارسنجی استفاده می‌کند، نه همان تابعی که حلّال در داخل خود استفاده می‌نماید. این جداسازی، الگوی \lr{independent oracle} است و از خطای \lr{circular validation} جلوگیری می‌کند؛ یعنی اگر حلّال از یک تابع هدف باگ‌دار استفاده کند، این باگ در اعتبارسنجی نیز منتقل نشود و شناسایی گردد. تابع \lr{\texttt{analytical\_skr}} در این ماژول، مستقل از حلّال پیاده‌سازی شده و به‌عنوان یک \lr{oracle} عمل می‌کند که صحت خروجی حلّال را راستی‌آزمایی می‌نماید.

\section{مبانی تئوری و فیزیکی}

\subsection{پروتکل Decoy-State و نقش پارامترهای شدت}

در پروتکل \lr{decoy-state BB84}، فرستنده (\lr{Alice}) به‌جای ارسال پالس با شدت ثابت، به‌طور تصادفی پالس‌ها را در سه دسته با شدت‌های مختلف توزیع می‌کند: سیگنال با شدت $\mu$، دکوی ضعیف با شدت $\nu$ (\lr{$\nu < \mu$}) و واکوئم با شدت صفر. هدف از این توزیع، آن است که حمله \lr{Photon-Number Splitting} (\lr{PNS}) که علیه منبع \lr{Weak Coherent Pulse} (\lr{WCP}) با توزیع \lr{Poisson} قابل اعمال است، خنثی شود. در این حمله، حاملهور (\lr{eavesdropper}) با تقسیم پالس‌های چندفوتونی، یک کپی از فوتون را نگه می‌دارد و بقیه را به گیرنده ارسال می‌کند؛ اما در حضور دکوی، حاملهور نمی‌تواند تشخیص دهد که کدام پالس سیگنال و کدام دکوی است و بنابراین، حمله \lr{PNS} قابل اعمال نیست.

پارامترهای کلیدی این پروتکل عبارت‌اند از: $q_x$ (\lr{sifting factor}، کسر پالس‌هایی که پایه‌های \lr{Alice} و \lr{Bob} هم‌خوان هستند)، $p_s$ (احتمال تخصیص پالس به سیگنال)، $p_d$ (احتمال تخصیص پالس به دکوی)، $p_v = 1 - p_s - p_d$ (احتمال تخصیص به واکوئم)، $\mu$ (شدت میانگین سیگنال) و $\nu$ (شدت میانگین دکوی). انتخاب این پارامترها، یک مسئله بهینه‌سازی غیرخطی است که هدف آن، حداکثرسازی نرخ کلید امن است. این بهینه‌سازی، به‌دلیل وابستگی پیچیده \lr{SKR} به این پارامترها (از طریق استخراج $Y_1$ و محاسبه $e_1$ و $E_\mu$)، به‌سادگی قابل حل نیست و نیاز به حلّال‌های عددی نظیر \lr{SLSQP} یا \lr{Nelder-Mead} دارد.

\subsection{آنتروپی باینری شانون و نقش آن در امنیت}

آنتروپی باینری شانون که با $H_2(p)$ نشان داده می‌شود، یکی از مفاهیم بنیادین نظریه اطلاعات است که در رمزنگاری کوانتومی نقشی مرکزی ایفا می‌کند. این کمیت، میزان عدم‌قطعیت متغیر تصادفی باینری با توزیع \lr{Bernoulli} را اندازه می‌گیرد:

\begin{equation}
	H_2(p) = -p \log_2(p) - (1 - p) \log_2(1 - p)
\end{equation}

در این رابطه، $p \in [0, 1]$ احتمال وقوع یکی از دو حالت است. تابع $H_2(p)$ در $p = 0.5$ به مقدار بیشینه $1$ می‌رسد و در نقاط $p = 0$ و $p = 1$ به صفر میل می‌کند. در رمزنگاری کوانتومی، $H_2(p)$ در دو نقطه بحرانی ظاهر می‌شود: نخست در محاسبه \lr{privacy amplification} که حجم اطلاعات نشت‌کرده به حاملهور را تعیین می‌کند (\lr{$H_2(E_\mu)$} و \lr{$H_2(e_1)$})، و دوم در محاسبه ضریب \lr{privacy amplification} در فرمول \lr{GLLP}.

پیاده‌سازی تابع $H_2$ در کد، با برش عددی در نقاط $p = 0$ و $p = 1$ همراه است، زیرا در این نقاط، عبارت $\log_2(p)$ به $-\infty$ میل می‌کند. در محاسبات عددی، این حالت می‌تواند منجر به \lr{NaN} یا \lr{ValueError} شود؛ بنابراین پیاده‌سازی پیش از فراخوانی \lr{\texttt{math.log2}}، شرایط $p \leq 0$ یا $p \geq 1$ را بررسی کرده و در این موارد، مقدار $0$ را برمی‌گرداند. این برش با حد گسسته \lr{L'Hôpital} هم‌خوان است:

\begin{equation}
	\lim_{p \to 0^+} p \log_2 p = \lim_{p \to 0^+} \frac{\log_2 p}{1/p} = \lim_{p \to 0^+} \frac{-p}{\ln 2} = 0
\end{equation}

\subsection{کاهیدگی فیبر و قانون Beer-Lambert}

انتقال فوتون در فیبر نوری تک‌حالت در طول‌موج مخابراتی $1550$ نانومتر با قانون \lr{Beer-Lambert} مدل می‌شود. برای یک فیبر به طول $L$ کیلومتر با ضریب کاهیدگی خطی $\alpha_{\text{dB/km}}$، کاهیدگی کل بر حسب دسی‌بل برابر با $\alpha_{\text{dB/km}} \cdot L$ است و تبدیل آن به ضریب انتقال خطی از رابطه زیر پیروی می‌کند:

\begin{equation}
	\eta_{\text{channel}}(L, \alpha_{\text{dB/km}}) = 10^{-\frac{\alpha_{\text{dB/km}} \cdot L}{10}}
\end{equation}

در این پیاده‌سازی، $\alpha_{\text{dB/km}} = 0.2$ (مقدار استاندارد فیبر تک‌حالت در $1550$ نانومتر) به‌کار رفته است. در فاصله $L = 50$ کیلومتر، $\eta_{\text{channel}} = 10^{-1} = 0.1$ و در فاصله $L = 100$ کیلومتر، $\eta_{\text{channel}} = 10^{-2} = 0.01$ می‌شود. ضریب انتقال کلی سیستم شامل بازده آشکارساز $\eta_{\text{det}}$ نیز می‌شود:

\begin{equation}
	\eta_{\text{total}}(L) = \eta_{\text{det}} \cdot \eta_{\text{channel}}(L, \alpha_{\text{dB/km}})
\end{equation}

در این پیاده‌سازی، $\eta_{\text{det}} = 0.15$ که متناظر با آشکارساز \lr{InGaAs SPAD} در دمای $-50$ درجه سانتی‌گراد است. بنابراین در $L = 50$ کیلومتر، $\eta_{\text{total}} = 0.015$ می‌شود.

\subsection{Gain سیگنال و QBER سیگنال}

برای منبع \lr{WCP} با توزیع \lr{Poisson} با شدت میانگین $\mu$، \lr{Gain} سیگنال (یعنی احتمال آنکه یک پالس منجر به ثبت کلیک در آشکارساز شود) برابر است با:

\begin{equation}
	Q_\mu(\mu, \eta, Y_0) = 1 - (1 - Y_0) \cdot e^{-\eta \mu}
\end{equation}

در این رابطه، $Y_0$ نرخ \lr{dark count} به ازای هر پالس، $\eta$ ضریب انتقال کلی سیستم و $\mu$ شدت میانگین سیگنال است. عبارت $e^{-\eta \mu}$ احتمال آنکه هیچ فوتونی شناسایی نشود و عامل $(1 - Y_0)$ احتمال آنکه دارک کانت نیز رخ ندهد را نشان می‌دهد. \lr{QBER} سیگنال، یعنی کسر کلیک‌های اشتباه در میان کلیک‌های ثبت‌شده، از رابطه زیر به دست می‌آید:

\begin{equation}
	E_\mu = \frac{0.5 \cdot Y_0 + e_d \cdot (1 - e^{-\eta \mu})}{Q_\mu}
\end{equation}

در این رابطه، $e_d$ خطای ذاتی سیستم (\lr{misalignment error}) است. صورت این کسر از دو بخش تشکیل شده است: $0.5 \cdot Y_0$ که نشان‌دهنده کلیک‌های اشتباه ناشی از دارک کانت است (چون دارک کانت کاملاً تصادفی است، نصف آن در پایه درست و نصف در پایه اشتباه رخ می‌دهد) و $e_d \cdot (1 - e^{-\eta \mu})$ که نشان‌دهنده کلیک‌های اشتباه ناشی از فوتون‌های واقعی با احتمال خطای ذاتی $e_d$ است.

\subsection{Gain دکوی و استخراج Y1 به روش Lo2005}

\lr{Gain} دکوی ضعیف به‌صورت مشابه محاسبه می‌شود:

\begin{equation}
	Q_\nu(\nu, \eta, Y_0) = 1 - (1 - Y_0) \cdot e^{-\eta \nu}
\end{equation}

مرحله کلیدی در \lr{decoy-state QKD}، استخراج کران پایین برای $Y_1$ (احتمال آنکه یک پالس تک‌فوتونی منجر به کلیک شود) از مشاهدات $Q_\mu$ و $Q_\nu$ است. این کران، که در کار مرجع \lr{Lo-Ma-Chen (2005)} به‌تفصیل استخراج شده، در این ماژول به فرم زیر پیاده‌سازی شده است:

\begin{equation}
	Y_1^{L} = \frac{\mu}{\mu \nu - \nu^2} \left( Q_\nu \cdot e^{\nu} - \frac{\nu^2}{\mu^2} \cdot Q_\mu \cdot e^{\mu} \right)
\end{equation}

این فرمول، یک نسخه بازآرایی‌شده از فرمول کلاسیک \lr{Lo2005} است. ضریب $\frac{\mu}{\mu \nu - \nu^2} = \frac{\mu}{\nu(\mu - \nu)}$ در صورتی که $\nu \to \mu$ باشد، به بی‌نهایت میل می‌کند؛ بنابراین در پیاده‌سازی، شرط $\left| \mu - \nu \right| < 10^{-9}$ بررسی می‌شود و در این حالت، تابع مقدار صفر را برمی‌گرداند. در این پیاده‌سازی، یک اصلاح فیزیکی مهم نیز اعمال شده است: مقدار نهایی $Y_1$ با \lr{\texttt{max(Y1\_L, Y0)}} محاسبه می‌شود تا تضمین گردد که $Y_1$ هرگز کم‌تر از نرخ دارک کانت $Y_0$ نشود. این برش، بازتاب‌دهنده این واقعیت فیزیکی است که حتی در غیاب فوتون واقعی، دارک کانت می‌تواند منجر به کلیک شود و بنابراین، بازده تک‌فوتونی نمی‌تواند کم‌تر از $Y_0$ باشد.

\subsection{نرخ خطای تک‌فوتونی و بازده تک‌فوتونی}

نرخ خطای تک‌فوتونی $e_1$، یعنی کسر کلیک‌های اشتباه در میان پالس‌های تک‌فوتونی، از رابطه زیر به دست می‌آید:

\begin{equation}
	e_1 = \frac{0.5 \cdot Y_0 + e_d \cdot \eta}{Y_1}
\end{equation}

این رابطه مشابه $E_\mu$ است، با این تفاوت که به‌جای احتمال کلیک فوتون واقعی $1 - e^{-\eta \mu}$، احتمال کلیک تک‌فوتونی $\eta$ قرار دارد. در پیاده‌سازی، شرط $Y_1 > 0$ بررسی می‌شود تا از تقسیم بر صفر جلوگیری شود؛ در صورت $Y_1 = 0$، مقدار $e_1 = 0.5$ گرفته می‌شود که نشان‌دهنده حالت کاملاً تصادفی است (یعنی در غیاب سیگنال واقعی، نتیجه اندازهگیری کاملاً تصادفی و با احتمال برابر برای هر بیت است).

\lr{Gain} تک‌فوتونی، یعنی کسری از کلیک‌های کل که از پالس‌های تک‌فوتونی ناشی می‌شوند، از رابطه زیر به دست می‌آید:

\begin{equation}
	Q_1 = \mu \cdot e^{-\mu} \cdot Y_1
\end{equation}

این رابطه از ضرب احتمال آنکه پالس تک‌فوتونی باشد (\lr{$P(1 | \mu) = \mu e^{-\mu}$}) در بازده تک‌فوتونی $Y_1$ به دست می‌آید.

\subsection{نرخ کلید امن مجانبی (GLLP)}

نرخ کلید امن مجانبی (\lr{asymptotic secure key rate}) از فرمول کلاسیک \lr{GLLP} به دست می‌آید:

\begin{equation}
	R = q_x \cdot \left[ Q_1 \cdot (1 - H_2(e_1)) - Q_\mu \cdot f_{\text{EC}} \cdot H_2(E_\mu) \right]
\end{equation}

این فرمول از دو بخش تشکیل شده است. بخش اول $Q_1 \cdot (1 - H_2(e_1))$ نشان‌دهنده حجم اطلاعات مشترک بین فرستنده و گیرنده از پالس‌های تک‌فوتونی پس از \lr{privacy amplification} است؛ عبارت $1 - H_2(e_1)$ ضریب \lr{privacy amplification} است که حجم اطلاعات نشت‌کرده به حاملهور را کم می‌کند. بخش دوم $Q_\mu \cdot f_{\text{EC}} \cdot H_2(E_\mu)$ نشان‌دهنده حجم اطلاعات فاش‌شده در مرحله \lr{error correction} است؛ $f_{\text{EC}} \geq 1$ بازدهی \lr{error correction} است که در عمل $f_{\text{EC}} \approx 1.16$ برای کد \lr{LDPC} بهینه است. در این پیاده‌سازی، $f_{\text{EC}} = 1.16$ به‌کار رفته است که با مقادیر عملی در پیاده‌سازی‌های \lr{QKD} هم‌خوان است.

تفاوت این فرمول با حالت \lr{finite-key} (مثل اثبات \lr{Lim-2014}) آن است که در اینجا از حد مجانبی $N \to \infty$ استفاده می‌شود و جریمه‌های آماری نظیر \lr{Hoeffding bound} و \lr{composable security} اعمال نمی‌گردند. به همین دلیل، $R$ در اینجا همیشه بزرگ‌تر یا مساوی با نرخ کلید متناهی است و در عمل به‌عنوان \emph{کران بالای} نرخ کلید قابل دستیابی عمل می‌کند. با این حال، تطابق حلّال بهینه‌سازی با این کران بالا، شرط لازم — هرچند کافی نیست — برای صحت پیاده‌سازی است.

\subsection{قیدهای فیزیکی و ریاضی}

در تابع \lr{\texttt{analytical\_skr}}، چندین قید فیزیکی و ریاضی اعمال شده است که در صورت نقض، تابع مقدار صفر را برمی‌گرداند. این قیدها عبارت‌اند از:

نخست، قید \lr{$p_s + p_d \geq 1$}: این قید تضمین می‌کند که مجموع احتمال‌های تخصیص سیگنال و دکوی از $1$ تجاوز نکند؛ در غیر این صورت، احتمال واکوئم $p_v = 1 - p_s - p_d$ منفی می‌شود که از لحاظ فیزیکی بی‌معناست. در صورت نقض، تابع مقدار صفر برمی‌گرداند.

دوم، قید \lr{$\nu \geq 0.9 \cdot \mu$}: این قید تضمین می‌کند که شدت دکوی به‌طور قابل‌توجهی کم‌تر از شدت سیگنال باشد؛ در غیر این صورت، استخراج $Y_1$ از روش \lr{Lo2005} به‌دلیل نزدیک بودن $\mu$ و $\nu$ به تکینگی (\lr{singularity}) می‌رسد و نتایج غیرقابل اعتماد تولید می‌کند. در عمل، $\nu$ معمولاً در بازه $0.1 \mu$ تا $0.5 \mu$ انتخاب می‌شود.

سوم، قید \lr{$\left| \mu - \nu \right| < 10^{-9}$}: این قید حالت تکین را شناسایی می‌کند و در صورت وقوع، تابع مقدار صفر برمی‌گرداند.

این قیدها، الگوی \lr{defensive programming} را به نمایش می‌گذارند: ورودی‌ها پیش از محاسبه اعتبارسنجی می‌شوند و در صورت غیرمعتبر بودن، مقدار پیش‌فرض امن (صفر) بازگردانده می‌شود. این الگو در توابع بهینه‌سازی حیاتی است، چون حلّال ممکن است در جریان جستجوی فضای پارامتر، به مقادیر غیرفیزیکی برسد و بدون این قیدها، تابع هدف می‌تواند مقادیر غیرمتعارف یا \lr{NaN} تولید کند.

\section{تحلیل معماری و ساختار داده‌ها}

\subsection{ساختار کلی ماژول و الگوی Module-Level Functions}

این ماژول از الگوی \lr{module-level functions} پیروی می‌کند، یعنی به‌جای تعریف کلاس‌ها، از توابع مستقل در سطح ماژول استفاده می‌نماید. این انتخاب طراحی، بازتاب‌دهنده ماهیت \lr{procedural} این ماژول است: توابع تحلیلی (\lr{\texttt{H2}} و \lr{\texttt{analytical\_skr}}) توابع ریاضی خالص (\lr{pure functions}) هستند که ورودی‌هایشان را به خروجی‌هایشان نگاشت می‌کنند و هیچ \lr{state} داخلی ندارند. تابع \lr{\texttt{run\_lp\_test}} تنها تابع دارای \lr{side-effect} است که خروجی چاپ می‌کند و حلّال را فراخوانی می‌نماید. این جداسازی بین توابع خالص و توابع دارای \lr{side-effect}، یکی از اصول مهندسی نرم‌افزار بالغ است و در الگوهای طراحی نظیر \lr{functional core, imperative shell} تبیین می‌شود.

\subsection{تنظیم مسیر و Path Bootstrap}

در ابتدای ماژول، الگوی \lr{path bootstrap} پیاده‌سازی شده است:

\begin{LTR}
	\begin{lstlisting}[language=Python]
		sys.path.insert(0, os.path.abspath(os.path.dirname(__file__)))
	\end{lstlisting}
\end{LTR}

این خط، مسیر دایرکتوری حاوی اسکریپت را به ابتدای \lr{\texttt{sys.path}} اضافه می‌کند. این کار تضمین می‌نماید که ماژول بدون توجه به محل اجرای آن (از \lr{project root} یا از مسیر دیگر)، بتواند ماژول‌های \lr{\texttt{qkd.proofs.optimization}} و \lr{\texttt{qkd.params}} را وارد نماید. این الگو در اسکریپت‌هایی که به‌صورت مستقل اجرا می‌شوند و نیاز به واردسازی ماژول‌های پروژه دارند، استاندارد است.

\subsection{واردسازی توابع از فریم‌ورک}

دو واردسازی کلیدی در این ماژول انجام می‌شود:

\begin{LTR}
	\begin{lstlisting}[language=Python]
		from qkd.proofs.optimization import optimize_lim2014_finite_key
		from qkd.params import load_lim2014_dedicated_params
	\end{lstlisting}
\end{LTR}

نخست، تابع \lr{\texttt{optimize\_lim2014\_finite\_key}} از ماژول \lr{\texttt{qkd.proofs.optimization}} وارد می‌شود. این تابع، حلّال بهینه‌سازی پارامترهای پروتکل \lr{Lim-2014} است که از الگوریتم \lr{Nelder-Mead} یا \lr{SLSQP} برای یافتن پارامترهای بهینه (\lr{$q_x$}، \lr{$p_s$}، \lr{$p_d$}، \lr{$\mu$}، \lr{$\nu$}) استفاده می‌کند. دوم، تابع \lr{\texttt{load\_lim2014\_dedicated\_params}} از ماژول \lr{\texttt{qkd.params}} وارد می‌شود. این تابع، پارامترهای اختصاصی پروتکل \lr{Lim-2014} (مانند \lr{$N$}، \lr{$\varepsilon_{\text{PE}}$}، \lr{$\varepsilon_{\text{EC}}$} و سایر پارامترهای امنیتی) را برای یک فاصله و تعداد بیت مشخص بارگذاری می‌نماید.

این وابستگی‌ها نشان‌دهنده آن است که ماژول اعتبارسنجی، مستقیماً به موتور بهینه‌سازی فریم‌ورک متصل است و همان مسیر اجرای واقعی را آزمون می‌کند. این الگو، تضمین می‌کند که آزمون، رفتار واقعی موتور را در شرایط عملیاتی منعکس می‌کند و نه رفتار آزمایشگاهی مصنوعی.

\subsection{تحلیل تابع H2}

تابع \lr{\texttt{H2(p)}} یک تابع ریاضی خالص است که آنتروپی باینری شانون را محاسبه می‌نماید:

\begin{LTR}
	\begin{lstlisting}[language=Python]
		def H2(p):
		if p <= 0 or p >= 1: return 0.0
		return -p * math.log2(p) - (1 - p) * math.log2(1 - p)
	\end{lstlisting}
\end{LTR}

در خط اول، شرط $p \leq 0$ یا $p \geq 1$ بررسی می‌شود. این شرط شامل نقاط انتهایی $p = 0$ و $p = 1$ نیز می‌شود که در آن‌ها آنتروپی به‌صورت تحلیلی صفر است. در خط دوم، فرمول اصلی $-p \log_2 p - (1-p) \log_2(1-p)$ پیاده‌سازی می‌شود. چون این خط فقط در صورتی اجرا می‌شود که $p \in (0, 1)$، فراخوانی \lr{\texttt{math.log2}} همیشه با ورودی مثبت انجام می‌شود و خطای \lr{math domain error} رخ نمی‌دهد. این پیاده‌سازی، الگوی \lr{defensive programming} را به نمایش می‌گذارد.

\subsection{تحلیل تابع analytical\_skr و ساختار پارامترها}

تابع \lr{\texttt{analytical\_skr}} اصلی‌ترین تابع تحلیلی این ماژول است که نرخ کلید امن را بر اساس پنج پارامتر پروتکل محاسبه می‌کند. این تابع نه پارامتر ورودی دارد: \lr{$q_x$} (\lr{sifting factor})، \lr{$p_s$} (احتمال سیگنال)، \lr{$p_d$} (احتمال دکوی)، \lr{$\mu$} (شدت سیگنال)، \lr{$\nu$} (شدت دکوی)، \lr{$\eta$} (ضریب انتقال کلی)، \lr{$Y_0$} (نرخ دارک کانت)، \lr{$e_d$} (خطای ذاتی) و \lr{$f_{\text{EC}}$} (بازدهی \lr{error correction}). خروجی تابع، یک عدد اعشاری است که نرخ کلید امن را نشان می‌دهد.

\begin{LTR}
	\begin{lstlisting}[language=Python]
		def analytical_skr(q_x, p_s, p_d, mu, nu, eta, Y0, ed, f_ec):
		"""تابع هزینه تحلیلی برای بررسی پارامترها (اصلاح شده)"""
		if p_s + p_d >= 1.0: return 0.0
		if nu >= mu * 0.9: return 0.0
		
		Q_mu = 1 - (1 - Y0) * math.exp(-eta * mu)
		E_mu = (0.5 * Y0 + ed * (1 - math.exp(-eta * mu))) / Q_mu
		Q_nu = 1 - (1 - Y0) * math.exp(-eta * nu)
	\end{lstlisting}
\end{LTR}

در خطوط ابتدایی، دو قید فیزیکی بررسی می‌شوند: قید \lr{$p_s + p_d \geq 1$} و قید \lr{$\nu \geq 0.9 \mu$}. در صورت نقض هر یک، تابع بلافاصله مقدار صفر را برمی‌گرداند. این الگوی \lr{early return} از پردازش غیرضروری جلوگیری می‌کند و خوانایی کد را بالا می‌برد. سپس سه کمیت کلیدی $Q_\mu$، $E_\mu$ و $Q_\nu$ محاسبه می‌شوند. این کمیت‌ها پایه محاسبات بعدی هستند.

\begin{LTR}
	\begin{lstlisting}[language=Python]
		if abs(mu - nu) < 1e-9: return 0.0
		
		# اصلاح فرمول Lo2005 برای Y1
		term1 = Q_nu * math.exp(nu)
		term2 = (nu**2 / mu**2) * Q_mu * math.exp(mu)
		factor = mu / (mu * nu - nu**2)
		Y1_L = factor * (term1 - term2)
		
		Y1 = max(Y1_L, Y0)  # Y1 نمی‌تواند کمتر از نویز تاریک باشد
		e1 = (0.5 * Y0 + ed * eta) / Y1 if Y1 > 0 else 0.5
	\end{lstlisting}
\end{LTR}

در این بخش، کران پایین $Y_1^L$ محاسبه می‌شود. شرط $\left| \mu - \nu \right| < 10^{-9}$ حالت تکین را شناسایی می‌کند؛ در این حالت، تابع مقدار صفر را برمی‌گرداند. در غیر این صورت، فرمول اصلی \lr{Lo2005} پیاده می‌شود: \lr{\texttt{term1}} و \lr{\texttt{term2}} دو جمله اصلی صورت کسر هستند و \lr{\texttt{factor}} ضریب نرمال‌سازی است. در نهایت، \lr{\texttt{max(Y1\_L, Y0)}} تضمین می‌کند که $Y_1$ نهایی کم‌تر از $Y_0$ نشود؛ این برش فیزیکی، بازتاب‌دهنده این واقعیت است که حتی در غیاب فوتون واقعی، دارک کانت می‌تواند منجر به کلیک شود.

محاسبه $e_1$ نیز با شرط $Y_1 > 0$ همراه است تا از تقسیم بر صفر جلوگیری شود؛ در صورت $Y_1 = 0$، مقدار $e_1 = 0.5$ گرفته می‌شود که نشان‌دهنده حالت کاملاً تصادفی است.

\begin{LTR}
	\begin{lstlisting}[language=Python]
		Q1 = mu * math.exp(-mu) * Y1
		R = q_x * (Q1 * (1 - H2(e1)) - Q_mu * f_ec * H2(E_mu))
		return max(R, 0.0)
	\end{lstlisting}
\end{LTR}

در بخش پایانی، $Q_1$ و $R$ محاسبه می‌شوند. $Q_1$ از ضرب $\mu e^{-\mu}$ (احتمال پالس تک‌فوتونی) در $Y_1$ (بازده تک‌فوتونی) به دست می‌آید. $R$ از فرمول \lr{GLLP} محاسبه می‌شود و در نهایت، \lr{\texttt{max(R, 0.0)}} تضمین می‌کند که نرخ کلید منفی نشود؛ این برش در رژیم‌های با \lr{QBER} بالا (مثلاً $E_\mu > 0.11$ که آستانه بحرانی \lr{BB84} است) ضروری است، چون در این رژیم، عبارت $Q_\mu \cdot f_{\text{EC}} \cdot H_2(E_\mu)$ بزرگ‌تر از $Q_1 \cdot (1 - H_2(e_1))$ می‌شود و $R$ منفی می‌گردد. در عمل، نرخ کلید منفی به این معناست که هیچ کلید امنی قابل استخراج نیست و مقدار صفر گزارش می‌شود.

\subsection{تحلیل تابع run\_lp\_test و الگوی Cost Callback}

تابع \lr{\texttt{run\_lp\_test}} اصلی‌ترین تابع آزمون است که در یک فاصله مشخص، حلّال بهینه‌سازی را فراخوانی کرده و نتایج را مقایسه می‌نماید. این تابع یک پارامتر ورودی \lr{\texttt{dist\_km}} دارد که فاصله به کیلومتر است.

\begin{LTR}
	\begin{lstlisting}[language=Python]
		def run_lp_test(dist_km):
		print(f"\n[TEST] LP Solver at L={dist_km} km")
		eta_det = 0.15
		Y0 = 3e-6
		ed = 0.01
		f_ec = 1.16
		alpha = 0.2
		eta = 10 ** (-alpha * dist_km / 10) * eta_det
	\end{lstlisting}
\end{LTR}

در ابتدا، پارامترهای فیزیکی ثابت تعریف می‌شوند: $\eta_{\text{det}} = 0.15$، $Y_0 = 3 \times 10^{-6}$، $e_d = 0.01$، $f_{\text{EC}} = 1.16$ و $\alpha = 0.2$ (\lr{dB/km}). سپس ضریب انتقال کلی $\eta$ بر اساس قانون \lr{Beer-Lambert} محاسبه می‌شود. این پارامترها با مقادیر مرجع در ادبیات \lr{QKD} هم‌خوان هستند و در شرایط عملیاتی واقعی برای \lr{QKD} در فواصل مخابراتی متوسط قابل دستیابی هستند.

\begin{LTR}
	\begin{lstlisting}[language=Python]
		params = load_lim2014_dedicated_params(distance_km=dist_km, num_bits=10**7)
		
		def cost_callback(q_x, p_s, p_d, mu, nu):
		# بهینه‌ساز به دنبال مینیمم کردن هزینه است، پس منفی SKR را برمی‌گردانیم
		skr = analytical_skr(q_x, p_s, p_d, mu, nu, eta, Y0, ed, f_ec)
		return -skr
	\end{lstlisting}
\end{LTR}

در این بخش، پارامترهای پروتکل \lr{Lim-2014} برای فاصله و تعداد بیت مشخص بارگذاری می‌شوند. سپس تابع \lr{\texttt{cost\_callback}} به‌عنوان یک \lr{closure} تعریف می‌شود. این تابع، پنج پارامتر پروتکل را می‌گیرد و \emph{منفی} \lr{SKR} را برمی‌گرداند. دلیل این negate کردن آن است که حلّال بهینه‌سازی به‌طور پیش‌فرض به دنبال مینیمم‌کردن تابع هزینه است؛ بنابراین برای حداکثرسازی \lr{SKR}، باید منفی آن را مینیمم کرد. این الگو در ادبیات بهینه‌سازی به‌عنوان \lr{minimization of negated objective} شناخته می‌شود.

این تابع \lr{closure} به متغیرهای \lr{$\eta$}، \lr{$Y_0$}، \lr{$e_d$} و \lr{$f_{\text{EC}}$} از محیط بیرونی دسترسی دارد (\lr{closure capture})؛ این الگو امکان تزریق پارامترهای محیطی به تابع هدف را بدون نیاز به ارسال آن‌ها به‌عنوان پارامتر فراهم می‌سازد. این الگو در توابع \lr{callback} بسیار رایج است و خوانایی کد را بالا می‌برد.

\subsection{الگوی Comparison و Hardcoded Reference}

در ادامه، \lr{SKR} با پارامترهای سخت‌کدشده محاسبه می‌شود و سپس حلّال فراخوانی می‌گردد:

\begin{LTR}
	\begin{lstlisting}[language=Python]
		# ۱. محاسبه SKR با پارامترهای ثابت و سخت‌کد شده
		R_hardcoded = analytical_skr(0.5, 0.5, 0.25, 0.5, 0.1, eta, Y0, ed, f_ec)
		
		# ۲. اجرای LP Solver فریم‌ورک
		try:
		q_x_opt, p_s_opt, p_d_opt, mu_opt, nu_opt = optimize_lim2014_finite_key(
		params, dist_km, cost_callback)
		R_optimized = analytical_skr(q_x_opt, p_s_opt, p_d_opt, mu_opt, nu_opt, 
		eta, Y0, ed, f_ec)
	\end{lstlisting}
\end{LTR}

در این بخش، \lr{$R_{\text{hardcoded}}$} با پارامترهای مرجع \lr{$q_x = 0.5$}، \lr{$p_s = 0.5$}، \lr{$p_d = 0.25$}، \lr{$\mu = 0.5$} و \lr{$\nu = 0.1$} محاسبه می‌شود. این پارامترها، مقادیر کلاسیک و شناخته‌شده در ادبیات \lr{decoy-state QKD} هستند و به‌عنوان نقطه مرجع عمل می‌کنند. سپس حلّال \lr{\texttt{optimize\_lim2014\_finite\_key}} فراخوانی می‌شود و پارامترهای بهینه را برمی‌گرداند. \lr{$R_{\text{optimized}}$} با این پارامترهای بهینه محاسبه می‌شود.

استفاده از بلاک \lr{\texttt{try/except}} برای مدیریت خطاهای احتمالی حلّال ضروری است؛ چون حلّال بهینه‌سازی ممکن است در شرایط دشوار (مثلاً در فواصل بسیار بزرگ که \lr{SKR} نزدیک به صفر است) به همگرایی نرسد یا خطا تولید کند. در صورت وقوع خطا، بلاک \lr{\texttt{except}} پیام خطا و \lr{traceback} کامل را چاپ می‌نماید و تابع مقدار \lr{False} برمی‌گرداند.

\subsection{الگوریتم Comparison و Verdict}

در نهایت، نتایج مقایسه شده و \lr{verdict} اعلام می‌گردد:

\begin{LTR}
	\begin{lstlisting}[language=Python]
		print(f"  Hardcoded (mu=0.5, nu=0.1)     -> SKR: {R_hardcoded:.6e}")
		print(f"  Optimized  (mu={mu_opt:.3f}, nu={nu_opt:.3f}) -> SKR: {R_optimized:.6e}")
		
		if R_optimized >= R_hardcoded - 1e-9:
		print("  PASS: LP Solver successfully found optimal or equal parameters.")
		return True
		else:
		print("  FAIL: LP Solver performed worse than hardcoded parameters!")
		return False
	\end{lstlisting}
\end{LTR}

در این الگوریتم، شرط \lr{$R_{\text{optimized}} \geq R_{\text{hardcoded}} - 10^{-9}$} بررسی می‌شود. در صورت برقراری، آزمون موفق تلقی می‌گردد؛ در غیر این صورت، شکست. استفاده از \lr{$10^{-9}$} به‌عنوان تلورانس، به این دلیل است که در محاسبات عددی، خطای گردکردن می‌تواند منجر به تفاوت‌های بسیار کوچک شود و این تلورانس تضمین می‌کند که چنین تفاوت‌های جزئی به‌عنوان شکست تلقی نشوند. این الگو، الگوی \lr{tolerance-based comparison} است که در آزمون‌های عددی استاندارد است.

\subsection{نقطه ورود \_\_main\_\_ و سه‌آزمون}

\begin{LTR}
	\begin{lstlisting}[language=Python]
		if __name__ == "__main__":
		print("=" * 60)
		print("LP SOLVER (PARAMETER OPTIMIZATION) VALIDATION")
		print("=" * 60)
		
		# تست در فواصل کوتاه، متوسط و طولانی
		res1 = run_lp_test(0.0)
		res2 = run_lp_test(50.0)
		res3 = run_lp_test(100.0)
		
		print("\n" + "=" * 60)
		if res1 and res2 and res3:
		print("100% LP SOLVER VALIDATION PASSED!")
		else:
		print("LP SOLVER VALIDATION FAILED.")
		print("=" * 60)
	\end{lstlisting}
\end{LTR}

نقطه ورود \lr{\texttt{\_\_main\_\_}} یک سرآیند تزئینی چاپ می‌کند و سپس سه آزمون در فواصل $0$، $50$ و $100$ کیلومتر را اجرا می‌نماید. این سه فاصله، سه رژیم فیزیکی متفاوت را پوشش می‌دهند: فاصله $0$ کیلومتر (بدون کاهیدگی)، فاصله $50$ کیلومتر (کاهیدگی متوسط با $\eta_{\text{channel}} \approx 0.1$) و فاصله $100$ کیلومتر (کاهیدگی شدید با $\eta_{\text{channel}} \approx 0.01$). در صورت موفقیت همه سه آزمون، پیام \lr{100\% LP SOLVER VALIDATION PASSED!} چاپ می‌شود و در غیر این صورت، پیام شکست. این الگو، تضمین می‌کند که حلّال در طیف وسیعی از شرایط عملیاتی به‌طور قابل‌اتکا عمل می‌کند.

\section{پیاده‌سازی ریاضی و الگوریتمی}

\subsection{پیاده‌سازی آنتروپی باینری با برش عددی}

پیاده‌سازی تابع $H_2$ در این ماژول، نمونه‌ای از \lr{defensive programming} در محاسبات عددی است. فرمول تحلیلی $-p \log_2 p - (1-p) \log_2(1-p)$ در نقاط $p = 0$ و $p = 1$ به‌صورت تحلیلی به $0$ میل می‌کند، اما در محاسبات عددی، فراخوانی \lr{\texttt{math.log2(0)}} منجر به \lr{ValueError} می‌شود. برای حل این مشکل، شرط $p \leq 0$ یا $p \geq 1$ پیش از فراخوانی \lr{\texttt{math.log2}} بررسی می‌شود و در این موارد، مقدار $0$ بازگردانده می‌شود. این برش عددی با حد گسسته \lr{L'Hôpital} هم‌خوان است و از نظر ریاضی صحیح است.

\subsection{پیاده‌سازی فرمول Lo2005 با اصلاحات فیزیکی}

پیاده‌سازی فرمول $Y_1^L$ در تابع \lr{\texttt{analytical\_skr}} چندین اصلاح مهم دارد نسبت به فرمول کلاسیک \lr{Lo2005}. نخست، فرمول به فرم زیر بازآرایی شده است:

\begin{equation}
	Y_1^{L} = \frac{\mu}{\mu \nu - \nu^2} \left( Q_\nu \cdot e^{\nu} - \frac{\nu^2}{\mu^2} \cdot Q_\mu \cdot e^{\mu} \right)
\end{equation}

این بازآرایی، معادل فرمول کلاسیک \lr{Lo2005} است، اما با ساختار متفاوت. در فرمول کلاسیک، ضریب نرمال‌سازی به فرم $\frac{\mu^2 e^{-\mu}}{\mu \nu - \mu^2}$ است؛ اما در این پیاده‌سازی، ضریب به فرم $\frac{\mu}{\mu \nu - \nu^2}$ بازنویسی شده است. این بازنویسی، با استفاده از رابطه $\mu^2 - \mu \nu = \mu(\mu - \nu) = -\nu(\mu - \nu) + (\mu^2 - \nu^2)$ و simplification جبری، قابل اثبات است.

دوم، اصلاح فیزیکی \lr{\texttt{max(Y1\_L, Y0)}} اعمال شده است. این اصلاح، تضمین می‌کند که $Y_1$ هرگز کم‌تر از $Y_0$ نشود. این برش، بازتاب‌دهنده این واقعیت فیزیکی است که حتی در غیاب فوتون واقعی، دارک کانت می‌تواند منجر به کلیک شود و بنابراین، بازده تک‌فوتونی نمی‌تواند کم‌تر از نرخ دارک کانت باشد. این اصلاح، در برخی نسخه‌های فرمول \lr{Lo2005} به‌صورت ضمنی اعمال می‌شود، اما در این پیاده‌سازی به‌صورت صریح و قابل مشاهده اعمال گردیده است.

سوم، در محاسبه $e_1$، در صورت $Y_1 = 0$، مقدار $e_1 = 0.5$ گرفته می‌شود. این انتخاب، بازتاب‌دهنده حالت کاملاً تصادفی است (یعنی در غیاب سیگنال واقعی، نتیجه اندازه‌گیری کاملاً تصادفی و با احتمال برابر برای هر بیت است). این انتخاب، در طراحی امنیتی \lr{QKD} محافظه‌کارانه است؛ چون فرض می‌کند که در صورت عدم اطمینان، حاملهور می‌تواند کل اطلاعات را بداند.

\subsection{پیاده‌سازی قیدهای فیزیکی و Early Return}

پیاده‌سازی قیدهای فیزیکی با الگوی \lr{early return} انجام می‌شود:

\begin{LTR}
	\begin{lstlisting}[language=Python]
		if p_s + p_d >= 1.0: return 0.0
		if nu >= mu * 0.9: return 0.0
		...
		if abs(mu - nu) < 1e-9: return 0.0
	\end{lstlisting}
\end{LTR}

این الگو، چندین مزیت دارد: نخست، خوانایی کد را بالا می‌برد، چون هر قید به‌صورت مستقل و در ابتدای تابع بررسی می‌شود. دوم، از پردازش غیرضروری جلوگیری می‌کند؛ در صورت نقض قید، بقیه محاسبات انجام نمی‌شود. سوم، در توابع بهینه‌سازی حیاتی است، چون حلّال ممکن است در جریان جستجوی فضای پارامتر، به مقادیر غیرفیزیکی برسد و بدون این قیدها، تابع هدف می‌تواند مقادیر غیرمتعارف یا \lr{NaN} تولید کند که حلّال را در همگرایی مختل می‌نماید.

قید \lr{$\nu \geq 0.9 \mu$} به‌طور خاص جالب توجه است؛ این قید، تلورانسی برای حالت تکین فرمول \lr{Lo2005} فراهم می‌سازد. در فرمول کلاسیک، تکینگی در $\nu = \mu$ رخ می‌دهد، اما در عمل، حتی نزدیک بودن $\nu$ به $\mu$ (مثلاً $\nu = 0.95 \mu$) می‌تواند منجر به نتایج غیرقابل اعتماد شود. این قید، تضمین می‌کند که $\nu$ به‌طور قابل‌توجهی کم‌تر از $\mu$ باشد.

\subsection{الگوریتم Cost Callback و Negate Trick}

الگوریتم \lr{cost callback} با \lr{negate trick} یک الگوی کلاسیک در بهینه‌سازی است:

\begin{LTR}
	\begin{lstlisting}[language=Python]
		def cost_callback(q_x, p_s, p_d, mu, nu):
		skr = analytical_skr(q_x, p_s, p_d, mu, nu, eta, Y0, ed, f_ec)
		return -skr
	\end{lstlisting}
\end{LTR}

در این الگوریتم، تابع هدف واقعی \lr{SKR} محاسبه می‌شود، اما به‌جای بازگرداندن آن، منفی آن بازگردانده می‌شود. دلیل این کار آن است که اکثر حلّال‌های بهینه‌سازی (مانند \lr{scipy.optimize.minimize}) به‌طور پیش‌فرض به دنبال مینیمم‌کردن تابع هستند؛ بنابراین برای حداکثرسازی \lr{SKR}، باید منفی آن را مینیمم کرد. این الگو در ادبیات بهینه‌سازی به‌عنوان \lr{minimization of negated objective} شناخته می‌شود.

این الگو، یک ویژگی مهم دارد: امکان استفاده از همان حلّال برای هر دو مسئله مینیمم‌سازی و حداکثرسازی. به‌عنوان مثال، اگر بخواهیم \lr{QBER} را مینیمم کنیم، تابع هدف \lr{QBER} را مستقیماً برمی‌گردانیم؛ و اگر بخواهیم \lr{SKR} را حداکثر کنیم، منفی آن را برمی‌گردانیم. این یکپارچگی، در طراحی ماژولار سیستم بهینه‌سازی ارزشمند است.

\subsection{الگوریتم Closure Capture برای پارامترهای محیطی}

الگوریتم \lr{closure capture} در تعریف \lr{\texttt{cost\_callback}} به‌کار رفته است:

\begin{LTR}
	\begin{lstlisting}[language=Python]
		eta = 10 ** (-alpha * dist_km / 10) * eta_det
		...
		def cost_callback(q_x, p_s, p_d, mu, nu):
		skr = analytical_skr(q_x, p_s, p_d, mu, nu, eta, Y0, ed, f_ec)
		return -skr
	\end{lstlisting}
\end{LTR}

در این الگوریتم، متغیرهای \lr{$\eta$}، \lr{$Y_0$}، \lr{$e_d$} و \lr{$f_{\text{EC}}$} در محیط بیرونی تابع \lr{\texttt{cost\_callback}} تعریف می‌شوند و تابع \lr{closure} به آن‌ها دسترسی دارد. این الگو، چندین مزیت دارد: نخست، امضای تابع \lr{callback} ساده می‌ماند (فقط پنج پارامتر پروتکل) و حلّال نیازی به دانستن پارامترهای محیطی ندارد. دوم، پارامترهای محیطی می‌توانند در زمان تعریف \lr{closure} محاسبه شوند و در تمام فراخوانی‌های بعدی استفاده گردند؛ این کار از محاسبه مکرر آن‌ها جلوگیری می‌کند. سوم، در صورت نیاز به تغییر پارامترهای محیطی (مثلاً برای آزمون در فاصله دیگر)، فقط کافی است \lr{closure} در یک محیط جدید تعریف شود.

این الگو در برنامه‌نویسی تابعی به‌کررات به‌کار می‌رود و در پایتون با کلمات کلیدی \lr{\texttt{lambda}} یا \lr{\texttt{def}} در داخل تابع دیگر، قابل پیاده‌سازی است.

\subsection{الگوریتم Hardcoded Reference Comparison}

الگوریتم مقایسه با مرجع سخت‌کدشده، الگوریتم ساده اما مؤثری است:

\begin{LTR}
	\begin{lstlisting}[language=Python]
		R_hardcoded = analytical_skr(0.5, 0.5, 0.25, 0.5, 0.1, eta, Y0, ed, f_ec)
		...
		R_optimized = analytical_skr(q_x_opt, p_s_opt, p_d_opt, mu_opt, nu_opt, 
		eta, Y0, ed, f_ec)
		...
		if R_optimized >= R_hardcoded - 1e-9:
		return True
		else:
		return False
	\end{lstlisting}
\end{LTR}

در این الگوریتم، \lr{$R_{\text{hardcoded}}$} با پارامترهای مرجع محاسبه می‌شود و سپس با \lr{$R_{\text{optimized}}$} مقایسه می‌گردد. شرط موفقیت، آن است که \lr{$R_{\text{optimized}}$} حداقل برابر با \lr{$R_{\text{hardcoded}}$} (با تلورانس \lr{$10^{-9}$}) باشد. این الگوریتم، چندین ویژگی دارد:

نخست، استفاده از پارامترهای مرجع \lr{$\mu = 0.5$} و \lr{$\nu = 0.1$} که در ادبیات \lr{decoy-state QKD} به‌عنوان مقادیر کلاسیک شناخته شده‌اند. این پارامترها، نقطه عملیاتی متعارف هستند و در بسیاری از پیاده‌سازی‌های آزمایشگاهی به‌کار رفته‌اند. انتخاب آن‌ها به‌عنوان مرجع، تضمین می‌کند که حلّال حداقل به اندازه یک راه‌حل شناخته‌شده عمل می‌کند.

دوم، استفاده از تلورانس \lr{$10^{-9}$} برای جلوگیری از خطای گردکردن. در محاسبات عددی، خطای گردکردن می‌تواند منجر به تفاوت‌های بسیار کوچک شود و این تلورانس تضمین می‌کند که چنین تفاوت‌های جزئی به‌عنوان شکست تلقی نشوند.

سوم، این الگوریتم، الگوی \lr{necessary condition testing} است: موفقیت آزمون تضمین می‌کند که حلّال حداقل به اندازه مرجع عمل می‌کند؛ هرچند ممکن است پارامترهای بهتر نیز یافت کند. در صورت شکست، قطعاً حلّال دارای باگ است و نیاز به بررسی دارد.

\subsection{الگوریتم Three-Range Testing}

الگوریتم سه‌آزمون در فواصل $0$، $50$ و $100$ کیلومتر، الگوریتم هوشمندانه‌ای است:

\begin{LTR}
	\begin{lstlisting}[language=Python]
		res1 = run_lp_test(0.0)
		res2 = run_lp_test(50.0)
		res3 = run_lp_test(100.0)
	\end{lstlisting}
\end{LTR}

این سه فاصله، سه رژیم فیزیکی متفاوت را پوشش می‌دهند:

نخست، فاصله $L = 0$ کیلومتر: در این حالت، $\eta_{\text{channel}} = 1$ و $\eta_{\text{total}} = \eta_{\text{det}} = 0.15$. در این رژیم، کاهیدگی فیبر وجود ندارد و حلّال باید پارامترهایی با \lr{SKR} نزدیک به حد نظری پیدا کند. در این رژیم، مقدار بهینه $\mu$ معمولاً بزرگ‌تر از $0.5$ است، چون کاهیدگی کم است و پالس‌های چندفوتونی کم‌تر مشکل‌ساز می‌شوند.

دوم، فاصله $L = 50$ کیلومتر: در این حالت، $\eta_{\text{channel}} \approx 0.1$ و $\eta_{\text{total}} \approx 0.015$. در این رژیم، کاهیدگی متوسط است و حلّال باید با کاهش $\mu$ برای جبران کاهیدگی، نرخ کلید را بهینه نماید. مقدار بهینه $\mu$ معمولاً در بازه $0.3$ تا $0.5$ است.

سوم، فاصله $L = 100$ کیلومتر: در این حالت، $\eta_{\text{channel}} \approx 0.01$ و $\eta_{\text{total}} \approx 0.0015$. در این رژیم، کاهیدگی شدید است و حلّال در رژیم نزدیک به آستانه بحرانی \lr{BB84} (\lr{$Q \approx 0.11$}) عمل می‌کند. مقدار بهینه $\mu$ معمولاً در بازه $0.1$ تا $0.3$ است و انتخاب پارامترها به‌طور حساس به نتیجه تأثیر می‌گذارد.

این سه آزمون، تضمین می‌کنند که حلّال در طیف وسیعی از شرایط عملیاتی به‌طور قابل‌اتکا عمل می‌کند و در رژیم‌های مختلف فیزیکی، پارامترهای بهینه را پیدا می‌نماید.

\section{ارسال و دریافت با سایر ماژول‌ها}

\subsection{وابستگی به qkd.proofs.optimization}

اصلی‌ترین وابستگی این ماژول، به \lr{\texttt{qkd.proofs.optimization}} است. این وابستگی از طریق تابع \lr{\texttt{optimize\_lim2014\_finite\_key}} برقرار می‌شود که حلّال بهینه‌سازی پارامترهای پروتکل \lr{Lim-2014} است. این تابع، سه پارامتر ورودی دارد: \lr{\texttt{params}} (پارامترهای پروتکل)، \lr{\texttt{dist\_km}} (فاصله) و \lr{\texttt{cost\_callback}} (تابع هزینه). خروجی آن، پنج پارامتر بهینه (\lr{$q_x$}، \lr{$p_s$}، \lr{$p_d$}، \lr{$\mu$}، \lr{$\nu$}) است.

این وابستگی، یک وابستگی \lr{loose coupling} از طریق \lr{callback} است: حلّال به‌جز امضای تابع \lr{callback}، هیچ چیز دیگری از تابع هدف نمی‌داند. این الگو، امکان جایگزینی تابع هدف با توابع دیگر را فراهم می‌سازد. به‌عنوان مثال، می‌توان به‌جای تابع \lr{analytical\_skr} از یک تابع هزینه پیچیده‌تر (مثلاً \lr{SKR} متناهی-\lr{key} یا تابع هزینه با جریمه‌های اضافی) استفاده کرد، بدون اینکه حلّال نیاز به تغییر داشته باشد.

\subsection{وابستگی به qkd.params}

دومین وابستگی این ماژول، به \lr{\texttt{qkd.params}} است. این وابستگی از طریق تابع \lr{\texttt{load\_lim2014\_dedicated\_params}} برقرار می‌شود که پارامترهای اختصاصی پروتکل \lr{Lim-2014} را بارگذاری می‌نماید. این تابع، دو پارامتر ورودی دارد: \lr{\texttt{distance\_km}} و \lr{\texttt{num\_bits}}. خروجی آن، یک شیء \lr{dict} یا \lr{dataclass} است که شامل پارامترهای امنیتی پروتکل (مانند \lr{$\varepsilon_{\text{PE}}$}، \lr{$\varepsilon_{\text{EC}}$}، \lr{$\varepsilon_{\text{s}}$} و سایر پارامترهای متناهی-\lr{key}) است.

این پارامترها، در حلّال بهینه‌سازی برای محاسبه جریمه‌های متناهی-\lr{key} استفاده می‌شوند. در این ماژول اعتبارسنجی، از این پارامترها فقط برای فراخوانی حلّال استفاده می‌شود و در تابع هدف تحلیلی (\lr{\texttt{analytical\_skr}}) مستقیماً به‌کار نمی‌روند؛ چون تابع هدف، حد مجانبی را محاسبه می‌کند و جریمه‌های متناهی-\lr{key} را در بر نمی‌گیرد. این جداسازی، الگوی \lr{independent oracle} است که در بخش بعد توضیح داده می‌شود.

\subsection{الگوی Independent Oracle و جلوگیری از Circular Validation}

یکی از ویژگی‌های برجسته این ماژول، استفاده از یک تابع هدف تحلیلی مستقل برای اعتبارسنجی است. تابع \lr{\texttt{analytical\_skr}} در این ماژول، مستقل از حلّال پیاده‌سازی شده و به‌عنوان یک \lr{oracle} عمل می‌کند که صحت خروجی حلّال را راستی‌آزمایی می‌نماید. این جداسازی، از خطای \lr{circular validation} جلوگیری می‌کند.

\lr{circular validation} زمانی رخ می‌دهد که اعتبارسنجی از همان تابعی استفاده کند که در پیاده‌سازی اصلی به‌کار رفته است. در این حالت، اگر تابع دارای باگ باشد، این باگ در اعتبارسنجی نیز منتقل می‌شود و شناسایی نمی‌گردد. به‌عنوان مثال، اگر حلّال از یک تابع هدف باگ‌دار (مثلاً با فرمول اشتباه برای $Y_1$) استفاده کند و اعتبارسنجی نیز از همان تابع استفاده نماید، هر دو خروجی یکسانی تولید می‌کنند و آزمون موفق تلقی می‌شود، در حالی که خروجی هر دو نادرست است.

برای جلوگیری از این خطا، تابع \lr{\texttt{analytical\_skr}} در این ماژول، مستقل از حلّال پیاده‌سازی شده است. این تابع، فرمول \lr{GLLP} را به‌صورت تحلیلی و با اصلاحات فیزیکی (مانند \lr{\texttt{max(Y1\_L, Y0)}}) پیاده می‌سازد و به‌عنوان یک مرجع مستقل عمل می‌نماید. در صورت بروز باگ در حلّال، این تابع مستقل، خطا را آشکار می‌سازد.

\subsection{خروجی به stdout و گزارش متمرکز}

خروجی این ماژول به‌صورت چاپ به \lr{stdout} است. این خروجی شامل چندین بخش است: سرآیند تزئینی، نتیجه هر آزمون (شامل پارامترهای سخت‌کدشده و بهینه و \lr{SKR} مربوطه)، و \lr{verdict} نهایی. این قالب خروجی، برای مصرف توسط توسعه‌دهنده در حین \lr{debug} طراحی شده و امکان تحلیل سریع نتایج را فراهم می‌سازد.

در صورت نیاز، می‌توان خروجی را به فایل نیز نوشت تا در \lr{CI/CD pipeline} قابل بازبینی باشد. در نسخه فعلی، این قابلیت پیاده‌سازی نشده است، اما با افزودن چند خط کد (مانند \lr{\texttt{with open(...) as f: f.write(...)}}) قابل افزودن است.

\subsection{ارتباط با absolute\_crypto\_validation.py}

این ماژول و \lr{\texttt{absolute\_crypto\_validation.py}} دو لایه مکمل اعتبارسنجی را تشکیل می‌دهند. \lr{\texttt{absolute\_crypto\_validation.py}} بر راستی‌آزمایی کمیت‌های پایه (\lr{$Q_\mu$}، \lr{$E_\mu$}، \lr{$Y_1$}، \lr{$e_1$}، \lr{$R$}) در برابر فرمول‌های تحلیلی متمرکز است، در حالی که این ماژول (\lr{\texttt{lp\_validation.py}}) بر راستی‌آزمایی حلّال بهینه‌سازی پارامترها متمرکز است. این دو لایه، نقش‌های مکمل دارند: لایه اول تضمین می‌کند که موتور محاسباتی به‌درستی پیاده شده، و لایه دوم تضمین می‌نماید که حلّال بهینه‌سازی به‌درستی عمل می‌کند. در صورت شکست هر یک از این لایه‌ها، مشکلات متفاوتی را آشکار می‌سازند: شکست لایه اول نشان‌دهنده باگ در محاسبات پایه است، در حالی که شکست لایه دوم نشان‌دهنده باگ در حلّال بهینه‌سازی (مثلاً عدم همگرایی یا یافتن کمینه محلی) است.

\subsection{ارتباط با ماژول‌های اثبات Lim-2014 و Ma-2005}

این ماژول به‌طور غیرمستقیم با ماژول‌های اثبات امنیتی \lr{\texttt{lim2014.py}} و \lr{\texttt{ma2005.py}} در ارتباط است. از طریق تابع \lr{\texttt{optimize\_lim2014\_finite\_key}}، این ماژول از پارامترهای پروتکل \lr{Lim-2014} استفاده می‌کند. با این حال، تابع هدف تحلیلی (\lr{\texttt{analytical\_skr}}) بر اساس حد مجانبی است و با فرمول‌های متناهی-\lr{key} در \lr{Lim-2014} متفاوت است. این تفاوت، عمدی است: آزمون کران بالای نرخ کلید را راستی‌آزمایی می‌کند و در صورتی که حلّال نتواند به کران بالا نزدیک شود، قطعاً نمی‌تواند نرخ کلید متناهی صحیح را تولید کند. این یک الگوی \lr{necessary condition testing} است.

\subsection{خروجی نهایی و استفاده در CI/CD}

خروجی نهایی این ماژول به دو شکل است: نخست، خروجی چاپی به کنسول که شامل نتایج هر آزمون و \lr{verdict} نهایی است. دوم، کد خروج \lr{POSIX} (به‌طور ضمنی از طریق \lr{\texttt{sys.exit}}) که در صورت شکست، غیرصفر است. در پیاده‌سازی فعلی، کد خروج به‌صورت صریح تنظیم نشده است، اما می‌توان با افزودن \lr{\texttt{sys.exit(0 if all([res1, res2, res3]) else 1)}} در پایان \lr{\texttt{\_\_main\_\_}}، این قابلیت را اضافه کرد.

در یک \lr{CI/CD pipeline} نمونه، این ماژول به‌صورت زیر به‌کار می‌رود:

\begin{LTR}
	\begin{lstlisting}[language=bash]
		- name: LP Solver Validation
		run: python lp_validation.py
	\end{lstlisting}
\end{LTR}

در صورت شکست، \lr{pipeline} متوقف می‌شود و توسعه‌دهنده باید قبل از \lr{merge}، خطا را اصلاح کند. این الگو، تضمین می‌کند که هرگز کدی با خطای بهینه‌سازی به شاخه اصلی \lr{merge} نشود.

\subsection{جمع‌بندی ارتباطات}

به‌طور خلاصه، این ماژول در نقش یک \lr{validator} مستقل عمل می‌کند که از طریق سه کانال با سایر بخش‌های فریم‌ورک در ارتباط است: ورودی حلّال بهینه‌سازی از \lr{\texttt{qkd.proofs.optimization}}، ورودی پارامترهای پروتکل از \lr{\texttt{qkd.params}}، و خروجی به \lr{stdout} برای مصرف توسعه‌دهنده. این معماری، جداسازی واضحی بین لایه اعتبارسنجی و موتور بهینه‌سازی برقرار می‌سازد و امکان توسعه مستقل هر یک را فراهم می‌نماید. خروجی این ماژول، نه‌تنها برای توسعه‌دهندگان (در فرآیند \lr{debug})، بلکه برای بازبینان کریپتوگرافیک (در فرآیند \lr{security audit}) نیز مفید است، زیرا یک \lr{audit trail} قابل ممیزی از عملکرد حلّال بهینه‌سازی در رژیم‌های مختلف فیزیکی فراهم می‌سازد.

	
	\chapter{آزمون‌های هم‌روندی و حفظ هویت}
	\label{ch:val_concurrency}
	\sloppy
\emergencystretch=5em
\hbadness=10000
\vbadness=10000
\tolerance=2000

\providecolor{codebg}{RGB}{248,248,248}
\providecolor{codeframe}{RGB}{200,200,200}
\providecolor{keyword}{RGB}{0,0,180}
\providecolor{string}{RGB}{163,21,21}
\providecolor{comment}{RGB}{0,128,0}

\lstset{
	backgroundcolor=\color{codebg},
	frame=single,
	rulecolor=\color{codeframe},
	language=Python,
	basicstyle=\ttfamily\scriptsize,
	keywordstyle=\color{keyword}\bfseries,
	stringstyle=\color{string},
	commentstyle=\color{comment}\itshape,
	showstringspaces=false,
	breaklines=true,
	breakatwhitespace=true,
	numbers=left,
	numberstyle=\tiny\color{gray},
	numbersep=8pt,
	tabsize=4,
	literate=
	{ف}{{{\bfseries\itshape ف}}}1
	{ی}{{{\bfseries\itshape ی}}}1
}

\section{نمای کلی و هدف ماژول}

ماژول \lr{\texttt{pulse\_tracking\_trace.py}} یک اسکریپت \emph{ممیزی و ردیابی سرتاسری} (\lr{end-to-end audit}) است که نقش یک \lr{forensic tracer} را در فریم‌ورک شبیه‌سازی \lr{QKD} ایفا می‌کند. این ماژول به‌طور خاص برای پاسخ به پرسش بنیادین زیر طراحی شده است: \emph{وقتی هزاران پالس نوری از طریق خط لوله چهارمرحله‌ای شبیه‌سازی (\lr{dispatch}، \lr{worker}، \lr{receipt}، \lr{row}) عبور می‌کنند، آیا هویت فیزیکی هر پالس — شامل بیت ارسالی \lr{Alice}، پایه انتخابی، پایه گیرنده \lr{Bob} و حالت کوانتومی — به‌طور سالم و بدون تغییر از مرحله‌ای به مرحله بعد منتقل می‌شود، یا این که در طول مسیر به‌دلیل \lr{reordering} موازی، \lr{race condition}، یا خطاهای پیاده‌سازی دچار فساد (\lr{corruption}) می‌گردد؟} این پرسش در شبیه‌سازی‌های موازی که در آن \lr{ProcessPoolExecutor} با چند کارگر روی هم‌زمان (\lr{concurrency}) اجرا می‌شود، حیاتی است؛ زیرا ترتیب خروجی (\lr{arrival order}) با ترتیب ورودی (\lr{input order}) لزوماً یکسان نیست و این ناهماهنگی می‌تواند به‌طور پنهانی منجر به جابجایی هویت پالس‌ها شود.

هدف اصلی این ماژول، تولید یک \emph{رد Tracks قابل ممیزی} (\lr{auditable trace}) است که برای ۱۰۰ پالس تصادفی — از جمعیت ۱۰۰۰ پالس — نشان می‌دهد که دقیقاً چه بر سر هر پالس در هر یک از چهار مرحله می‌آید. چهار مرحله عبارت‌اند از: مرحله \lr{S1} (\lr{Pre-Dispatch Snapshot}) که هویت اصلی پالس پیش از ارسال به کارگران را ثبت می‌کند؛ مرحله \lr{S2} (\lr{Receipt Snapshot}) که هویت پالس پس از دریافت در کارگر را نشان می‌دهد؛ مرحله \lr{S3} (\lr{Worker Output}) که نتیجه اندازه‌گیری کارگر را در بر دارد؛ و مرحله \lr{S4} (\lr{Final Row}) که ردیف نهایی \lr{CSV} را نمایش می‌دهد. اگر هر چهار مرحله بر سر تمام \lr{identity fields} توافق کنند، یعنی هویت پالس به‌طور کامل حفظ شده است؛ در غیر این صورت، خطای فساد رخ داده است.

ساختار کلی این ماژول شامل چندین بخش است. بخش اول، توابع قالب‌بندی (\lr{formatting helpers}) شامل \lr{\texttt{\_format\_table}}، \lr{\texttt{\_short\_identity}}، \lr{\texttt{\_identity\_match}} و \lr{\texttt{\_measurement\_match}} است که مسئول تولید خروجی متنی خوانا هستند. بخش دوم، تابع اصلی \lr{\texttt{run\_trace}} است که جریان کامل آزمون را هماهنگ می‌کند: ساخت جمعیت پالس، انتخاب پالس‌های ردیابی‌شده، اجرای \lr{dispatch}، محاسبه اثبات \lr{reordering}، تولید جدول ممیزی چهارمرحله‌ای، و محاسبه آمار سطح-مرحله‌ای. بخش سوم، تابع \lr{\texttt{main}} و \lr{CLI} است که از \lr{\texttt{argparse}} برای پیکربندی پارامترها (تعداد کارگران، اندازه جمعیت، تعداد پالس‌های ردیابی‌شده، بذر تصادفی، تعداد ردیف‌های نمایش‌داده‌شده، فایل خروجی) استفاده می‌کند.

ویژگی برجسته این ماژول، الگوی \lr{reuse-by-import} آن است؛ یعنی به‌جای بازنویسی ماشین‌آلات آزمون، مستقیماً از توابع کمکی در ماژول \lr{\texttt{tests/test\_main\_optimized\_pulse\_tracking.py}} استفاده می‌کند. این توابع شامل \lr{\texttt{\_make\_pulse\_population}} برای ساخت جمعیت، \lr{\texttt{\_select\_tracked\_pids}} برای انتخاب تصادفی پالس‌های ردیابی‌شده با بذر قابل بازتولید، \lr{\texttt{\_run\_tracking\_dispatch}} برای اجرای \lr{dispatch} با ردیابی، \lr{\texttt{\_expected\_measurement}} برای پیش‌بینی نتیجه اندازه‌گیری بر اساس هویت پالس، و \lr{\texttt{\_tracking\_slow\_worker}} که یک کارگر با تأخیر دوجمله‌ای (\lr{bimodal delay}) است. این الگو، اصل \lr{DRY} (\lr{Don't Repeat Yourself}) را رعایت می‌کند و تضمین می‌نماید که اسکریپت ممیزی و آزمون واحد، از همان ماشین‌آلات استفاده کنند و در نتیجه، نتایج آن‌ها همیشه هم‌خوان باشند.

در سطح عملیاتی، این ماژول به‌عنوان یک \lr{diagnostic tool} عمل می‌کند؛ یعنی در صورت بروز مشکلات مربوط به حفظ هویت پالس (مثلاً پس از \lr{refactor} موتور \lr{dispatch} یا تغییر مدل کارگر)، توسعه‌دهنده می‌تواند این اسکریپت را اجرا کند و به‌سرعت علت فساد را شناسایی کند. خروجی آن به دو شکل است: چاپ به \lr{stdout} برای مشاهده زنده، و ذخیره به فایل \lr{\texttt{pulse\_tracking\_trace.txt}} در مسیر \lr{\texttt{/home/z/my-project/download/}} برای بازبینی بعدی. این الگو در مهندسی نرم‌افزار علمی به‌عنوان \lr{shift-left forensics} شناخته می‌شود و به توسعه‌دهنده اجازه می‌دهد خطاهای ظریف را پیش از انتشار نتایج شناسایی کند.

نکته مهم دیگر آن است که این ماژول \emph{اثبات \lr{reordering}} (\lr{reordering proof}) را به‌صورت صریح و قابل مشاهده ارائه می‌دهد. در شبیه‌سازی موازی، خروجی کارگران به ترتیب ورودی نیست؛ بلکه ترتیب ورودی به ترتیب اتمام کارگران بستگی دارد. این ماژول با محاسبه تعداد \lr{pairwise inversions} (یعنی زوج‌های پالسی که ترتیب ورودی‌شان با ترتیب خروجی متفاوت است)، نشان می‌دهد که \lr{reordering} رخ داده اما هیچ هویتی فاسد نشده است. این تمایز بنیادینی است: \lr{reordering} یک ویژگی ذاتی و قابل‌قبول شبیه‌سازی موازی است، در حالی که \lr{corruption} یک باگ است که باید برطرف شود.

\section{مبانی تئوری و فیزیکی}

\subsection{پروتکل BB84 و مفهوم هویت پالس}

پروتکل \lr{BB84} که توسط \lr{Bennett} و \lr{Brassard} در سال ۱۹۸۴ پیشنهاد شد، ساده‌ترین و پرکاربردترین پروتکل \lr{QKD} است. در این پروتکل، فرستنده (\lr{Alice}) یک بیت کلاسیک را با انتخاب تصادفی یکی از دو پایه (\lr{basis}) — پایه مستطیلی (\lr{rectilinear}) $\{|0\rangle, |1\rangle\}$ یا پایه قطری (\lr{diagonal}) $\{|+\rangle, |-\rangle\}$ — به گیرنده (\lr{Bob}) ارسال می‌کند. \lr{Bob} نیز به‌طور تصادفی یکی از این دو پایه را برای اندازه‌گیری انتخاب می‌کند. در صورتی که پایه‌های \lr{Alice} و \lr{Bob} یکسان باشند، نتیجه اندازه‌گیری \lr{Bob} با بیت ارسالی \lr{Alice} هم‌خوان است (به جز خطاهای ناشی از نویز کانال و آشکارساز)؛ در غیر این صورت، نتیجه تصادفی است.

هویت هر پالس در شبیه‌سازی \lr{BB84} با چندین فیلد کلیدی مشخص می‌شود که در مجموع \emph{identity fields} نامیده می‌شوند. این فیلدها در ثابت \lr{\texttt{IDENTITY\_FIELDS}} در ماژول آزمون تعریف شده‌اند و شامل موارد زیر هستند: \lr{\texttt{alice\_bit}} (بیت ارسالی \lr{Alice}، صفر یا یک)، \lr{\texttt{alice\_basis}} (پایه انتخابی \lr{Alice}، مستطیلی یا قطری)، \lr{\texttt{bob\_basis}} (پایه انتخابی \lr{Bob}، مستطیلی یا قطری)، \lr{\texttt{state\_label}} (برچسب حالت کوانتومی مانند \lr{|0⟩}، \lr{|1⟩}، \lr{|+⟩}، \lr{|−⟩})، \lr{\texttt{mu}} (شدت میانگین پالس) و \lr{\texttt{intensity\_label}} (برچسب شدت مانند \lr{signal}، \lr{decoy} یا \lr{vacuum}). حفظ یکپارچگی این فیلدها در طول خط لوله شبیه‌سازی، شرط لازم برای صحت نتایج است.

\subsection{پیش‌بینی نتیجه اندازه‌گیری و قانون Born}

نتیجه اندازه‌گیری یک پالس \lr{BB84} توسط \lr{Bob} را می‌توان بر اساس هویت پالس پیش‌بینی کرد. این پیش‌بینی بر اساس قانون \lr{Born} و قواعد پروتکل \lr{BB84} انجام می‌شود. در صورتی که پایه‌های \lr{Alice} و \lr{Bob} یکسان باشند (\lr{basis\_match = True}) و پالس به‌طور موفقیت‌آمیزی شناسایی شود (\lr{detected = True})، بیت اندازه‌گیری‌شده با بیت ارسالی هم‌خوان است (\lr{measured\_bit = alice\_bit}). در صورتی که پایه‌ها متفاوت باشند، نتیجه اندازه‌گیری کاملاً تصادفی است و احتمال آن $0.5$ برای هر بیت است. در شبیه‌سازی قطعی با بذر مشخص، این تصادفی بودن از یک \lr{PRNG} قابل بازتولید به دست می‌آید.

پیش‌بینی نتیجه اندازه‌گیری در تابع \lr{\texttt{\_expected\_measurement}} پیاده‌سازی شده است. این تابع یک \lr{dict} شامل سه کلید \lr{\texttt{basis\_match}}، \lr{\texttt{measured\_bit}} و \lr{\texttt{detected}} برمی‌گرداند که مقادیر پیش‌بینی‌شده را در بر دارد. مقایسه این مقادیر با مقادیر گزارش‌شده توسط کارگر (\lr{S3})، امکان تشخیص خطاهای پیاده‌سازی را فراهم می‌سازد. به‌عنوان مثال، اگر کارگر \lr{\texttt{basis\_match = True}} و \lr{\texttt{detected = True}} را گزارش دهد اما \lr{\texttt{measured\_bit}} با \lr{\texttt{alice\_bit}} متفاوت باشد، این یک ناسازگاری است که نشان‌دهنده باگ در منطق اندازه‌گیری کارگر است.

\subsection{تأخیر دوجمله‌ای و شبیه‌سازی کارگر واقعی}

یکی از چالش‌های کلیدی در آزمون شبیه‌سازی موازی، ایجاد شرایط واقعی برای آشکار کردن \lr{race condition} و \lr{reordering} است. اگر همه کارگران با تأخیر یکسان اجرا شوند، خروجی به‌طور قطعی به ترتیب ورودی است و \lr{reordering} رخ نمی‌دهد. برای غلبه بر این مشکل، تابع \lr{\texttt{\_tracking\_slow\_worker}} با تأخیر دوجمله‌ای (\lr{bimodal delay}) پیاده‌سازی شده است:

\begin{equation}
	T_{\text{worker}} = \begin{cases} 100 \text{ ms} & \text{با احتمال } 0.10 \\ U(0, 5) \text{ ms} & \text{با احتمال } 0.90 \end{cases}
\end{equation}

در این رابطه، $T_{\text{worker}}$ زمان اجرای کارگر است و $U(0, 5)$ توزیع یکنواخت در بازه $[0, 5]$ میلی‌ثانیه است. این مدل، دو رژیم را شبیه‌سازی می‌کند: $10\%$ کارگران \emph{کند} هستند (تأخیر ۱۰۰ میلی‌ثانیه که شبیه‌سازی بار \lr{CPU} سنگین یا \lr{I/O} است) و $90\%$ کارگران \emph{سریع} هستند (تأخیر تصادفی ۰ تا ۵ میلی‌ثانیه). این توزیع غیریکنواخت، احتمال \lr{reordering} را به‌طور قابل‌توجهی افزایش می‌دهد و امکان آشکارسازی باگ‌های مربوط به حفظ ترتیب را فراهم می‌سازد.

\subsection{ Kendall-Tau و معیار ناهماهنگی}

برای اندازه‌گیری کمی میزان \lr{reordering}، از معیار \lr{Kendall-Tau inversion count} استفاده می‌شود. این معیار، تعداد زوج‌های (\lr{i, j}) با $i < j$ را می‌شمارد که در آن ترتیب خروجی برعکس ترتیب ورودی است:

\begin{equation}
	\text{inv}(\pi) = \left| \{(i, j) : i < j \land \pi(i) > \pi(j)\} \right|
\end{equation}

در این رابطه، $\pi$ جایگشت خروجی نسبت به ورودی است. حداکثر تعداد وارونگی‌ها برای $n$ پالس برابر با $\binom{n}{2} = \frac{n(n-1)}{2}$ است که در صورت جایگشت معکوس کامل رخ می‌دهد. درصد وارونگی به‌صورت زیر تعریف می‌شود:

\begin{equation}
	\text{reorder\%} = \frac{\text{inv}(\pi)}{\binom{n}{2}} \times 100
\end{equation}

این معیار، یک اندازه عاری از مقیاس (\lr{scale-free}) برای سنجش شدت \lr{reordering} فراهم می‌سازد. در شبیه‌سازی با ۱۰۰ پالس ردیابی‌شده و چند کارگر با تأخیر دوجمله‌ای، مقدار معمول این معیار در محدوده $30\%$ تا $50\%$ است که نشان‌دهنده \lr{reordering} قابل‌توجه اما نه کامل است.

\subsection{مفهوم Identity-Bound-to-ID و Position-Independent Integrity}

مفهوم کلیدی که این ماژول آن را به‌اثبات می‌رساند، \emph{یکپارچگی وابسته به شناسه، نه به موقعیت} (\lr{identity bound to pulse\_id, not to row position}) است. در یک طراحی ساده اما خطاکار، ممکن است هویت پالس بر اساس موقعیت ردیف در خروجی استنتاج شود؛ یعنی فرض شود که ردیف $i$ در خروجی، مربوط به پالس $i$ در ورودی است. این فرض در حضور \lr{reordering} نادرست است و منجر به تخصیص اشتباه هویت به ردیف‌ها می‌شود.

طراحی صحیح، هویت را به \lr{\texttt{pulse\_id}} (یک شناسه یکتا که در زمان ساخت پالس تخصیص می‌یابد) می‌بندد. این شناسه به‌صورت یکتا در کل جمعیت است و در طول خط لوله بدون تغییر منتقل می‌شود. هر مرحله از خط لوله، این شناسه را در خروجی خود شامل می‌شود و امکان بازسازی هویت پالس را فراهم می‌سازد. این الگو در ادبیات مهندسی نرم‌افزار به‌عنوان \lr{correlation ID} یا \lr{trace ID} شناخته می‌شود و در سیستم‌های توزیع‌شده (\lr{distributed systems}) برای ردیابی درخواست‌ها به‌کررات به‌کار می‌رود.

\section{تحلیل معماری و ساختار داده‌ها}

\subsection{ساختار کلی ماژول و الگوی Script-with-Helpers}

این ماژول از الگوی \lr{script-with-helpers} پیروی می‌کند؛ یعنی یک تابع اصلی \lr{\texttt{run\_trace}} که جریان اجرا را هماهنگ می‌کند، و چندین تابع کمکی برای قالب‌بندی و مقایسه. این الگو برای اسکریپت‌های \lr{CLI} مناسب است که نیاز به خوانایی بالا و قابلیت آزمون دارند. تابع \lr{\texttt{run\_trace}} خروجی را به‌صورت یک رشته متنی بزرگ برمی‌گرداند (نه چاپ مستقیم)، که امکان نوشتن آن در \lr{stdout} یا فایل را فراهم می‌سازد. این جداسازی بین تولید خروجی و چاپ، الگوی \lr{separation of concerns} را رعایت می‌کند.

\subsection{ثابت IDENTITY\_FIELDS و فیلدهای هویتی}

ثابت \lr{\texttt{IDENTITY\_FIELDS}} که از ماژول آزمون وارد می‌شود، فهرست فیلدهایی است که برای تعیین هم‌خوانی هویت در چهار مرحله بررسی می‌شوند. این فیلدها در ساختار \lr{dict} هر پالس وجود دارند و شامل \lr{\texttt{alice\_bit}}، \lr{\texttt{alice\_basis}}، \lr{\texttt{bob\_basis}}، \lr{\texttt{state\_label}} و سایر فیلدهای هویتی هستند. استفاده از یک ثابت متمرکز برای این فیلدها، تضمین می‌کند که تمام بخش‌های ماژول از همان فهرست استفاده کنند و در صورت افزودن فیلد جدید، فقط یک نقطه نیاز به به‌روزرسانی دارد.

\begin{LTR}
	\begin{lstlisting}[language=Python]
		from tests.test_main_optimized_pulse_tracking import (
		IDENTITY_FIELDS,
		_make_pulse_population,
		_select_tracked_pids,
		_run_tracking_dispatch,
		_expected_measurement,
		_tracking_slow_worker,
		)
		from main_optimized_dispatch import dispatch_work_items
	\end{lstlisting}
\end{LTR}

این الگوی \lr{import} چندین نکته دارد. نخست، \lr{\texttt{IDENTITY\_FIELDS}} به‌عنوان یک ثابت وارد می‌شود تا در تمام ماژول قابل استفاده باشد. دوم، توابع کمکی با پیشوند \lr{\_} (مانند \lr{\texttt{\_make\_pulse\_population}}) وارد می‌شوند که در پایتون نشان‌دهنده \lr{private API} است؛ اما چون در همان پروژه هستیم، این واردسازی مجاز است. سوم، \lr{\texttt{dispatch\_work\_items}} از ماژول \lr{\texttt{main\_optimized\_dispatch}} وارد می‌شود که این نشان‌دهنده وابستگی مستقیم به موتور \lr{dispatch} اصلی فریم‌ورک است.

\subsection{تابع \_format\_table و قالب‌بندی متنی}

تابع \lr{\texttt{\_format\_table}} یک ابزار قالب‌بندی متنی است که لیستی از سطرها را به جدولی با ستون‌های هم‌تراز تبدیل می‌کند. این تابع دو پارامتر اصلی دارد: \lr{\texttt{headers}} (فهرست عناوین ستون‌ها) و \lr{\texttt{rows}} (لیست سطرها، که هر سطر یک لیست از مقادیر است). پارامتر سوم \lr{\texttt{col\_widths}} اختیاری است و در صورت عدم ارائه، به‌صورت خودکار بر اساس حداکثر طول مقدار در هر ستون محاسبه می‌شود.

\begin{LTR}
	\begin{lstlisting}[language=Python]
		def _format_table(headers: List[str], rows: List[List[Any]],
		col_widths: List[int] = None) -> str:
		"""Format a list of rows as an aligned text table."""
		if not rows:
		return "(no rows)"
		if col_widths is None:
		col_widths = []
		for i, h in enumerate(headers):
		w = max(len(str(h)), *(len(str(r[i])) for r in rows))
		col_widths.append(w)
		lines = []
		lines.append("  ".join(str(h).ljust(w) for h, w in zip(headers, col_widths)))
		lines.append("  ".join("-" * w for w in col_widths))
		for r in rows:
		lines.append("  ".join(str(c).ljust(w) for c, w in zip(r, col_widths)))
		return "\n".join(lines)
	\end{lstlisting}
\end{LTR}

در خط اول، شرط \lr{\texttt{not rows}} بررسی می‌شود و در صورت خالی بودن، رشته \lr{"(no rows)"} بازگردانده می‌شود. این الگو، الگوی \lr{graceful handling of empty input} است. در حلقه بعدی، عرض هر ستون بر اساس حداکثر طول مقدار در آن ستون محاسبه می‌شود؛ این کار با ترکیب \lr{\texttt{max}} و \lr{\texttt{generator expression}} انجام می‌شود که یک الگوی پایتونیک برای محاسبه حداکثر در یک مجموعه است. سپس خط عنوان، خط جداکننده (با \lr{-}‌ها) و سطرهای داده به‌ترتیب به خطوط اضافه می‌شوند. استفاده از \lr{\texttt{ljust(w)}} تضمین می‌کند که تمام مقادیر در یک ستون، عرض یکسان داشته باشند و جدول هم‌تراز باشد.

\subsection{تابع \_short\_identity و نمایش فشرده هویت}

تابع \lr{\texttt{\_short\_identity}} یک \lr{dict} هویت پالس را به یک رشته فشرده به فرمت \lr{bit/basis/bob/state} تبدیل می‌کند. این نمایش، خوانایی جدول ممیزی را به‌طور قابل‌توجهی افزایش می‌دهد؛ زیرا یک رشته کوتاه به‌جای چهار فیلد جداگانه نمایش داده می‌شود.

\begin{LTR}
	\begin{lstlisting}[language=Python]
		def _short_identity(d: Dict[str, Any]) -> str:
		"""Compact representation of an identity dict: 'bit/basis/bob/state'."""
		return (f"{d['alice_bit']}/{d['alice_basis']}/"
		f"{d['bob_basis']}/{d['state_label']}")
	\end{lstlisting}
\end{LTR}

این تابع با استفاده از \lr{f-string} و \lr{/} به‌عنوان جداکننده، یک رشته به فرمت \lr{0/0/1/|+⟩} تولید می‌کند که در آن اولین مقدار \lr{alice\_bit}، دومین \lr{alice\_basis}، سومین \lr{bob\_basis} و چهارمین \lr{state\_label} است. این فرمت، مشابه نمایش \lr{CSV} است اما با جداکننده \lr{/} که در خروجی متنی خواناتر است.

\subsection{تابع \_identity\_match و منطق مقایسه چهارمرحله‌ای}

تابع \lr{\texttt{\_identity\_match}} چهار \lr{dict} هویتی (\lr{s1}، \lr{s2}، \lr{s3}، \lr{s4}) را می‌گیرد و در صورتی که همه فیلدهای \lr{IDENTITY\_FIELDS} در هر چهار مرحله برابر باشند، مقدار \lr{"OK"} و در غیر این صورت، \lr{"FAIL"} برمی‌گرداند.

\begin{LTR}
	\begin{lstlisting}[language=Python]
		def _identity_match(s1, s2, s3, s4) -> str:
		"""Return 'OK' if all four stage views agree on every identity field."""
		for f in IDENTITY_FIELDS:
		if not (s1[f] == s2[f] == s3[f] == s4[f]):
		return "FAIL"
		return "OK"
	\end{lstlisting}
\end{LTR}

در این پیاده‌سازی، عبارت \lr{\texttt{s1[f] == s2[f] == s3[f] == s4[f]}} یک الگوی پایتونیک برای مقایسه زنجیره‌ای است که در آن، همه چهار مقدار باید برابر باشند. این عبارت معادل \lr{\texttt{s1[f] == s2[f] and s2[f] == s3[f] and s3[f] == s4[f]}} است اما خواناتر و کوتاه‌تر است. در صورت عدم تطابق در هر فیلد، تابع بلافاصله \lr{"FAIL"} برمی‌گرداند که الگوی \lr{early exit} است و از پردازش غیرضروری جلوگیری می‌کند.

\subsection{تابع \_measurement\_match و اعتبارسنجی پیش‌بینی}

تابع \lr{\texttt{\_measurement\_match}} دو پارامتر \lr{s1} (هویت پیش از \lr{dispatch}) و \lr{s3} (خروجی کارگر) را می‌گیرد و در صورتی که نتیجه اندازه‌گیری گزارش‌شده توسط کارگر با پیش‌بینی بر اساس \lr{s1} هم‌خوان باشد، \lr{"OK"} برمی‌گرداند.

\begin{LTR}
	\begin{lstlisting}[language=Python]
		def _measurement_match(s1, s3) -> str:
		"""Return 'OK' if the worker's measurement matches the S1 prediction."""
		exp = _expected_measurement(s1)
		if (s3["basis_match"] == exp["basis_match"] and
		s3["measured_bit"] == exp["measured_bit"] and
		s3["detected"] == exp["detected"]):
		return "OK"
		return "FAIL"
	\end{lstlisting}
\end{LTR}

در این تابع، ابتدا \lr{\texttt{\_expected\_measurement(s1)}} فراخوانی می‌شود تا پیش‌بینی نتیجه اندازه‌گیری بر اساس هویت پالس به دست آید. سپس سه فیلد کلیدی — \lr{\texttt{basis\_match}}، \lr{\texttt{measured\_bit}} و \lr{\texttt{detected}} — بین گزارش کارگر و پیش‌بینی مقایسه می‌شوند. این اعتبارسنجی، خطاهای احتمالی در منطق اندازه‌گیری کارگر را آشکار می‌سازد. به‌عنوان مثال، اگر کارگر به‌اشتباه \lr{\texttt{basis\_match = True}} را در حالی که پایه‌ها متفاوت هستند گزارش دهد، این تابع خطا را شناسایی می‌کند.

\subsection{تابع run\_trace و الگوی Closure برای خروجی}

تابع \lr{\texttt{run\_trace}} اصلی‌ترین تابع ماژول است. این تابع یک \lr{closure} به نام \lr{\texttt{out}} تعریف می‌کند که رشته‌ها را به لیست \lr{\texttt{lines}} اضافه می‌کند:

\begin{LTR}
	\begin{lstlisting}[language=Python]
		def run_trace(args) -> str:
		"""Run the 100-pulse trace and return the full text output."""
		lines: List[str] = []
		def out(s: str = ""):
		lines.append(s)
	\end{lstlisting}
\end{LTR}

این الگو چندین مزیت دارد: نخست، تمام خروجی در یک لیست جمع‌آوری می‌شود و در پایان به یک رشته با \lr{\texttt{"\textbackslash n".join(lines)}} تبدیل می‌گردد. این کار، از هزینه \lr{string concatenation} مکرر (که در پایتون از مرتبه $O(n^2)$ است) جلوگیری می‌کند. دوم، تابع \lr{\texttt{out}} به‌عنوان یک \lr{closure} به \lr{\texttt{lines}} دسترسی دارد و نیازی به ارسال \lr{\texttt{lines}} به عنوان پارامتر نیست. سوم، این الگو امکان تغییر آینده به \lr{file-like object} یا \lr{stream} را فراهم می‌سازد.

\subsection{تنظیم مسیر و الگوی Path Bootstrap}

در ابتدای ماژول، الگوی \lr{path bootstrap} پیاده‌سازی شده است:

\begin{LTR}
	\begin{lstlisting}[language=Python]
		SCRIPT_DIR = os.path.dirname(os.path.abspath(__file__))
		PROJECT_ROOT = os.path.dirname(SCRIPT_DIR)
		for candidate in (
		PROJECT_ROOT,
		os.path.join(PROJECT_ROOT, "src"),
		os.path.join(PROJECT_ROOT, "tests"),
		):
		if candidate not in sys.path:
		sys.path.insert(0, candidate)
	\end{lstlisting}
\end{LTR}

این الگو، تضمین می‌کند که ماژول بدون توجه به محل اجرای آن (از \lr{project root} یا از مسیر دیگر)، بتواند ماژول‌های \lr{\texttt{tests}} و \lr{\texttt{src}} را وارد کند. \lr{\texttt{SCRIPT\_DIR}} مسیر دایرکتوری حاوی اسکریپت است و \lr{\texttt{PROJECT\_ROOT}} یک سطح بالاتر است. سپس سه مسیر (ریشه پروژه، \lr{src}، \lr{tests}) به \lr{\texttt{sys.path}} اضافه می‌شوند. این الگو در اسکریپت‌هایی که به‌صورت مستقل اجرا می‌شوند و نیاز به واردسازی ماژول‌های پروژه دارند، استاندارد است.

\section{پیاده‌سازی ریاضی و الگوریتمی}

\subsection{الگوریتم انتخاب تصادفی پالس‌های ردیابی‌شده}

انتخاب ۱۰۰ پالس از جمعیت ۱۰۰۰ پالس، با بذر تصادفی مشخص، توسط تابع \lr{\texttt{\_select\_tracked\_pids}} انجام می‌شود. این تابع از \lr{\texttt{random.Random(seed)}} برای تولید یک \lr{PRNG} قابل بازتولید استفاده می‌کند و سپس با \lr{\texttt{random.sample}}، $k$ شناسه از جمعیت $N$ انتخاب می‌نماید. این الگوریتم، الگوریتم \lr{Fisher-Yates partial shuffle} است که در آن، فقط $k$ جایگشت اول محاسبه می‌شود:

\begin{equation}
	P(\text{pid } i \text{ selected}) = \frac{k}{N}, \quad \forall i \in \{0, 1, \ldots, N-1\}
\end{equation}

این الگوریتم تضمین می‌نماید که هر پالس با احتمال یکسان $k/N$ انتخاب شود و انتخاب‌ها مستقل باشند. با بذر مشخص، انتخاب قابل بازتولید است که برای آزمون‌های \lr{regression} حیاتی است.

\subsection{الگوریتم محاسبه Pairwise Inversions}

محاسبه تعداد \lr{pairwise inversions} در لیست ورودی، با الگوریتم ساده $O(n^2)$ انجام می‌شود:

\begin{LTR}
	\begin{lstlisting}[language=Python]
		n = len(tracked_arrival_order)
		inversions = 0
		for i in range(n):
		ai = tracked_arrival_order[i]
		for j in range(i + 1, n):
		if ai > tracked_arrival_order[j]:
		inversions += 1
		max_inv = n * (n - 1) // 2
	\end{lstlisting}
\end{LTR}

در این الگوریتم، برای هر زوج (\lr{i, j}) با $i < j$، بررسی می‌شود که آیا \lr{ai > arrival[j]} است. در صورت مثبت بودن، یک وارونگی شمرده می‌شود. حداکثر تعداد وارونگی‌ها برابر با $\binom{n}{2} = \frac{n(n-1)}{2}$ است که با عملگر \lr{//} (تقسیم صحیح) محاسبه می‌شود. این الگوریتم برای $n = 100$ کافی است ($10^4$ عملیات)، اما برای $n$ بزرگ‌تر، الگوریتم \lr{merge-sort-based} با پیچیدگی $O(n \log n)$ ترجیح داده می‌شود.

\subsection{الگوریتم یافتن Concrete Swap Examples}

برای نمایش شهودی \lr{reordering}، الگوریتمی برای یافتن نمونه‌های ملموس جابجایی پیاده‌سازی شده است:

\begin{LTR}
	\begin{lstlisting}[language=Python]
		swap_examples = []
		for i in range(len(tracked_arrival_order) - 1):
		a = tracked_arrival_order[i]
		b = tracked_arrival_order[i + 1]
		if a > b:  # inversion: larger pid arrived before smaller
		swap_examples.append((i, i + 1, a, b))
		if len(swap_examples) >= 5:
		break
	\end{lstlisting}
\end{LTR}

این الگوریتم، زوج‌های مجاور (\lr{adjacent pairs}) را بررسی می‌کند و در صورتی که \lr{pid} بزرگ‌تر قبل از \lr{pid} کوچک‌تر آمده باشد، آن را به‌عنوان یک جابجایی ثبت می‌نماید. این الگوریتم، فقط \lr{adjacent inversions} را شناسایی می‌کند که زیرمجموعه‌ای از \lr{pairwise inversions} هستند. هدف از این الگوریتم، ارائه نمونه‌های ملموس و قابل درک برای کاربر است؛ زیرا مشاهده دو \lr{pid} که جابجا شده‌اند، شهودی‌تر از عدد انتزاعی \lr{inversion count} است. حداکثر ۵ نمونه نمایش داده می‌شود تا خروجی بیش از حد طولانی نشود.

\subsection{الگوریتم محاسبه Position Shifts}

برای تحلیل کمی جابجایی هر پالس، الگوریتم محاسبه \lr{position shifts} پیاده‌سازی شده است:

\begin{LTR}
	\begin{lstlisting}[language=Python]
		sorted_idx = {pid: i for i, pid in enumerate(tracked_pids)}
		arrived_idx = {pid: i for i, pid in enumerate(tracked_arrival_order)}
		shifts = [(pid, sorted_idx[pid], arrived_idx[pid], arrived_idx[pid] - sorted_idx[pid])
		for pid in tracked_pids]
		max_up = max(shifts, key=lambda x: x[3])
		max_down = min(shifts, key=lambda x: x[3])
	\end{lstlisting}
\end{LTR}

در این الگوریتم، ابتدا دو دیکشنری \lr{\texttt{sorted\_idx}} و \lr{\texttt{arrived\_idx}} ساخته می‌شوند که برای هر \lr{pid}، موقعیت آن در لیست مرتب (ورودی) و در لیست ورودی (خروجی) را نگه می‌دارند. سپس برای هر پالس، اختلاف این دو موقعیت (\lr{shift}) محاسبه می‌شود. \lr{shift} مثبت به این معناست که پالس در خروجی زودتر از ورودی آمده (حرکت به بالا) و \lr{shift} منفی به این معناست که پالس دیرتر آمده (حرکت به پایین). \lr{\texttt{max\_up}} و \lr{\texttt{max\_down}} به‌ترتیب بیشترین حرکت به بالا و بیشترین حرکت به پایین را نشان می‌دهند. این الگوریتم، الگوی \lr{dictionary comprehension} و \lr{list comprehension} را به‌خوبی به نمایش می‌گذارد.

\subsection{الگوریتم Stage-Level Mismatch Counting}

برای محاسبه آمار ناهماهنگی در هر مرحله، الگوریتم زیر پیاده‌سازی شده است:

\begin{LTR}
	\begin{lstlisting}[language=Python]
		for pid in tracked_pids:
		if pid not in receipts_by_pid:
		missing_receipts += 1
		continue
		if pid not in rows_by_pid:
		missing_rows += 1
		continue
		s1 = snapshot[pid]
		s2 = receipts_by_pid[pid]["received_item"]
		s3 = rows_by_pid[pid]
		s4 = rows_by_pid[pid]
		if any(s1[f] != s2[f] for f in IDENTITY_FIELDS):
		s1_s2_mismatch += 1
		if any(s2[f] != s3[f] for f in IDENTITY_FIELDS):
		s2_s3_mismatch += 1
		if any(s3[f] != s4[f] for f in IDENTITY_FIELDS):
		s3_s4_mismatch += 1
		exp = _expected_measurement(s1)
		if (s3["basis_match"] != exp["basis_match"] or
		s3["measured_bit"] != exp["measured_bit"] or
		s3["detected"] != exp["detected"]):
		measurement_mismatch += 1
	\end{lstlisting}
\end{LTR}

در این الگوریتم، برای هر \lr{pid} در \lr{tracked\_pids}، ابتدا وجود \lr{receipt} و \lr{row} بررسی می‌شود. در صورت عدم وجود، شمارنده \lr{missing} افزایش می‌یابد و با \lr{\texttt{continue}} به پالس بعدی می‌رود. در غیر این صورت، چهار \lr{dict} هویتی \lr{s1}، \lr{s2}، \lr{s3}، \lr{s4} استخراج می‌شوند و برای هر جفت متوالی (\lr{S1-S2}، \lr{S2-S3}، \lr{S3-S4})، تطابق فیلدهای هویتی بررسی می‌گردد. تابع \lr{\texttt{any}} با \lr{generator expression}، الگوی پایتونیک برای بررسی وجود عدم تطابق است. در پایان، تطابق نتیجه اندازه‌گیری نیز بررسی می‌شود. این الگوریتم، الگوی \lr{fail-fast} را رعایت می‌کند: در صورت فقدان داده‌های پایه‌ای (\lr{receipt} یا \lr{row})، بقیه بررسی‌ها برای آن پالس انجام نمی‌شود.

\subsection{الگوریتم Verdict و تفسیر نهایی}

در پایان، الگوریتم \lr{verdict} نتیجه نهایی را اعلام می‌کند:

\begin{LTR}
	\begin{lstlisting}[language=Python]
		total_mismatch = (missing_receipts + missing_rows + s1_s2_mismatch +
		s2_s3_mismatch + s3_s4_mismatch + measurement_mismatch)
		if total_mismatch == 0:
		out("  VERDICT:  All 100 tracked pulses preserved identity through all 4 stages.")
		out(f"           {inversions} pairwise inversions occurred (reordering), yet NOT A")
		out("           SINGLE pulse's identity was corrupted. Identity is bound to")
		out("           pulse_id, NOT to row position -- exactly as the contract requires.")
		else:
		out(f"  VERDICT:  {total_mismatch} corruption(s) detected across the 4 stages.")
		out("           See the per-pulse audit table above for details.")
	\end{lstlisting}
\end{LTR}

این الگوریتم، مجموع تمام ناهماهنگی‌ها را محاسبه می‌کند و در صورت صفر بودن، پیام موفقیت با تأکید بر تمایز بین \lr{reordering} (که قابل‌قبول است) و \lr{corruption} (که باگ است) چاپ می‌نماید. این تمایز، پیام کلیدی ماژول است: \lr{reordering} یک ویژگی ذاتی شبیه‌سازی موازی است و در صورتی که هویت به \lr{pulse\_id} بسته شده باشد، هیچ مشکلی ایجاد نمی‌کند.

\subsection{الگوریتم CLI با argparse}

تابع \lr{\texttt{main}} از \lr{\texttt{argparse}} برای پیکربندی پارامترها استفاده می‌کند:

\begin{LTR}
	\begin{lstlisting}[language=Python]
		def main():
		parser = argparse.ArgumentParser(
		description="100-random-pulse end-to-end state trace",
		formatter_class=argparse.RawDescriptionHelpFormatter,
		)
		parser.add_argument("--workers", type=int, default=4)
		parser.add_argument("--population", type=int, default=1000)
		parser.add_argument("--tracked", type=int, default=100)
		parser.add_argument("--seed", type=int, default=42)
		parser.add_argument("--first", type=int, default=100)
		parser.add_argument("--output", type=str, default=None)
		args = parser.parse_args()
		
		text = run_trace(args)
		print(text)
		
		default_path = os.path.join(PROJECT_ROOT, "download",
		"pulse_tracking_trace.txt")
		os.makedirs(os.path.dirname(default_path), exist_ok=True)
		with open(default_path, "w") as f:
		f.write(text)
		print(f"\n[trace also written to {default_path}]", file=sys.stderr)
	\end{lstlisting}
\end{LTR}

این الگوریتم، الگوی استاندارد \lr{CLI} در پایتون است. هر پارامتر با \lr{\texttt{add\_argument}} تعریف می‌شود و دارای مقدار پیش‌فرض است. \lr{\texttt{RawDescriptionHelpFormatter}} امکان قالب‌بندی متن راهنما را فراهم می‌سازد. پس از اجرای \lr{\texttt{run\_trace}}، خروجی به \lr{stdout} چاپ می‌شود و همچنین به فایل پیش‌فرض در \lr{\texttt{download/}} نوشته می‌گردد. در صورت ارائه \lr{\texttt{-{}-output}}، خروجی به فایل اضافی نیز نوشته می‌شود. این الگو، تضمین می‌کند که خروجی همیشه برای بازبینی بعدی ذخیره می‌شود، حتی اگر کاربر فایل خاصی مشخص نکند.

\section{ارسال و دریافت با سایر ماژول‌ها}

\subsection{وابستگی به tests/test\_main\_optimized\_pulse\_tracking}

اصلی‌ترین وابستگی این ماژول، به \lr{\texttt{tests/test\_main\_optimized\_pulse\_tracking.py}} است. این وابستگی شامل پنج مورد است: ثابت \lr{\texttt{IDENTITY\_FIELDS}} و چهار تابع \lr{\texttt{\_make\_pulse\_population}}، \lr{\texttt{\_select\_tracked\_pids}}، \lr{\texttt{\_run\_tracking\_dispatch}} و \lr{\texttt{\_expected\_measurement}} و \lr{\texttt{\_tracking\_slow\_worker}}. این الگوی \lr{reuse-by-import}، تضمین می‌کند که اسکریپت ممیزی و آزمون واحد، از همان ماشین‌آلات استفاده کنند و نتایج آن‌ها همیشه هم‌خوان باشند. در صورت تغییر در یکی از این توابع (مثلاً افزودن فیلد جدید به \lr{\texttt{IDENTITY\_FIELDS}})، هر دو به‌طور خودکار به‌روز می‌شوند.

این وابستگی، یک وابستگی \lr{tight coupling} است: تغییر در ساختار \lr{dict} پالس یا در امضای توابع، می‌تواند منجر به شکست این ماژول شود. برای کاهش این \lr{coupling}، می‌توان توابع کمکی را به یک ماژول جداگانه (مانند \lr{\texttt{pulse\_tracking\_utils.py}}) منتقل کرد که هم اسکریپت ممیزی و هم آزمون واحد از آن استفاده کنند. این الگو، الگوی \lr{shared utility module} است که در پروژه‌های بالغ به‌کررات به‌کار می‌رود.

\subsection{ورودی از main\_optimized\_dispatch}

این ماژول تابع \lr{\texttt{dispatch\_work\_items}} را از \lr{\texttt{main\_optimized\_dispatch}} وارد می‌کند. این تابع، موتور اصلی \lr{dispatch} فریم‌ورک است که کارهای شبیه‌سازی را بین چند کارگر توزیع می‌کند. وابستگی به این تابع، تضمین می‌کند که اسکریپت ممیزی، همان مسیر اجرای واقعی فریم‌ورک را آزمون می‌کند و نه یک مسیر آزمایشگاهی مصنوعی.

\subsection{خروجی به stdout و فایل}

خروجی این ماژول به دو شکل است: نخست، چاپ به \lr{stdout} با \lr{\texttt{print(text)}} که امکان مشاهده زنده را فراهم می‌سازد. دوم، نوشتن به فایل پیش‌فرض در مسیر \lr{\texttt{/home/z/my-project/download/pulse\_tracking\_trace.txt}} با \lr{\texttt{open(default\_path, "w")}}. در صورت ارائه \lr{\texttt{-{}-output}}، خروجی به فایل اضافی نیز نوشته می‌شود. این الگو، تضمین می‌کند که خروجی همیشه برای بازبینی بعدی ذخیره می‌شود، حتی اگر کاربر آن را به فایل خاصی هدایت نکرده باشد.

\subsection{ارتباط با فریم‌ورک آزمون}

این ماژول و ماژول آزمون \lr{\texttt{tests/test\_main\_optimized\_pulse\_tracking.py}} دو لایه مکمل اعتبارسنجی را تشکیل می‌دهند. لایه آزمون، اعتبارسنجی خودکار را با \lr{pytest} فراهم می‌سازد و در \lr{CI/CD pipeline} اجرا می‌شود. لایه اسکریپت ممیزی (این ماژول)، خروجی قابل مشاهده و قابل تحلیل توسط انسان را فراهم می‌کند و برای \lr{debug} و تحلیل دستی استفاده می‌شود. این دو لایه، نقش‌های مکمل دارند: لایه آزمون، \lr{gate} خودکار را برای \lr{merge} فراهم می‌سازد و لایه اسکریپت، امکان تحلیل عمیق در صورت بروز مشکل را فراهم می‌نماید.

\subsection{خروجی قابل تحلیل توسط ابزارهای خارجی}

خروجی متنی این ماژول، به‌گونه‌ای طراحی شده که توسط ابزارهای تحلیل متنی (مانند \lr{grep}، \lr{awk} یا اسکریپت‌های \lr{Python}) قابل پردازش باشد. جدول‌ها با \lr{spaces} هم‌تراز شده‌اند و جداکننده‌ها واضح هستند. در صورت نیاز، می‌توان خروجی را به فرمت \lr{JSON} یا \lr{CSV} نیز تبدیل کرد تا توسط ابزارهای تحلیل داده (مانند \lr{pandas}) پردازش شود. این انعطاف‌پذیری، در پروژه‌های تحقیقاتی که نیاز به تحلیل آماری نتایج شبیه‌سازی دارند، ارزشمند است.

\subsection{جمع‌بندی ارتباطات}

به‌طور خلاصه، این ماژول در نقش یک \lr{forensic tracer} مستقل عمل می‌کند که از طریق سه کانال با سایر بخش‌های فریم‌ورک در ارتباط است: ورودی توابع کمکی از \lr{\texttt{tests/test\_main\_optimized\_pulse\_tracking.py}}، ورودی موتور \lr{dispatch} از \lr{\texttt{main\_optimized\_dispatch}}، و خروجی به \lr{stdout} و فایل \lr{\texttt{pulse\_tracking\_trace.txt}}. این معماری، جداسازی واضحی بین لایه ممیزی و موتور شبیه‌سازی برقرار می‌سازد و امکان توسعه مستقل هر یک را فراهم می‌نماید. خروجی این ماژول، نه‌تنها برای توسعه‌دهندگان (در فرآیند \lr{debug})، بلکه برای بازبینان کریپتوگرافیک (در فرآیند \lr{security audit}) نیز مفید است، زیرا یک \lr{audit trail} قابل ممیزی از حفظ هویت پالس‌ها در طول خط لوله شبیه‌سازی فراهم می‌سازد.

	
	\part{نتایج شبیه‌سازی}
	\label{part:results_final}     % تغییر نام از part:results به part:results_final
	% chapters/18_simulation_results.tex
% بخش پنجم: نتایج شبیه‌سازی‌ها
% این فایل شامل تمام نمودارهای خروجی ۱۰ کمپین شبیه‌سازی است.
% پیش‌نیاز: بسته‌های زیر باید در preamble فایل اصلی اضافه شوند:
%   \usepackage{subcaption}   % برای subfigure ها
%   \usepackage{float}         % برای گزینه [H] در صورت نیاز
% مسیر تصاویر: فرض شده PNG ها در پوشه figures/ کنار فایل اصلی قرار دارند.

% --- helper: include image with graceful fallback if file is missing ---
% \detokenize renders the filename as plain catcode-12 characters so that
% underscores (and other special tokens) inside the path do NOT trigger
% "Missing $ inserted" / "Extra }" / "There's no line here to end" errors
% when the file is absent and the fallback box is typeset.
\providecommand{\safeimg}[2][]{%
\IfFileExists{#2}{\includegraphics[#1]{#2}}{%
	\fbox{\parbox[c][5cm][c]{0.8\textwidth}{\centering
			\ttfamily\detokenize{#2}\par\vspace{4pt}%
			\normalfont\small(فایل تصویری یافت نشد --- تصویر را در مسیر فوق قرار دهید)%
	}}%
}%
}

% =====================================================================
% PART V: نتایج شبیه‌سازی
% =====================================================================
\part{نتایج شبیه‌سازی}
\label{part:results}

% ---------------------------------------------------------------------
\chapter{تنظیمات آزمایش و کمپین‌های شبیه‌سازی}
\label{ch:results_overview}

\section{نمای کلی کمپین‌ها}
این بخش نتایج حاصل از اجرای ده کمپین شبیه‌سازی مستقل با اسکریپت \lr{\texttt{main\_optimized.py}} را گزارش می‌کند. هر کمپین با ترکیب متفاوتی از پروتکل، اثبات امنیت، نوع آشکارساز و پارامترهای فیزیکی اجرا شده تا رفتار چارچوب در شرایط عملیاتی گوناگون ارزیابی شود. خروجی هر اجرا در یک فایل \lr{CSV} به نام \lr{\texttt{qkd\_results\_optimized\_sweep.csv}} نوشته می‌شود و سپس برای تحلیل به نمودار \lr{PNG} تبدیل شده است. این نمودارها در ادامهٔ فصل به‌ترتیب موضوعی ارائه می‌شوند.

تنها خروجی ملموس هر شبیه‌سازی همین فایل \lr{CSV} است؛ اسکریپت هیچ نمودار، فایل \lr{JSON} خلاصه یا گزارش تجمیعی تولید نمی‌کند. بنابراین تمام تحلیل‌های آماری بعدی بر مبنای پردازش همین فایل با کتابخانهٔ \lr{\texttt{pandas}} انجام شده است. تعداد کل شبیه‌سازی‌های اجرا شده در این ده کمپین تقریباً برابر ۱۸۰۰ نقطه است که هر کدام شامل تولید پالس، انتقال در کانال فیبر، آشکارسازی، غربال‌سازی پایه‌ها (\lr{sifting})، تخمین پارامتر و محاسبهٔ کلید امن هستند.

\section{خلاصهٔ کمپین‌ها}
جدول~\ref{tab:campaigns_summary} ده کمپین اجرا شده را به‌همراه پارامترهای متمایزکننده و تعداد نقاط شبیه‌سازی خلاصه می‌کند. این جدول به‌عنوان مرجع سریع برای تفسیر نمودارهای فصل‌های بعدی استفاده می‌شود.

\begin{table}[htbp]
\centering
\caption{خلاصهٔ کمپین‌های شبیه‌سازی اجرا شده با \lr{\texttt{main\_optimized.py}}}
\label{tab:campaigns_summary}
\small
\begin{tabular}{@{}cllc@{}}
	\toprule
	\textbf{شماره} & \textbf{نام فایل} & \textbf{محور تغییر اصلی} & \textbf{تعداد نقاط} \\
	\midrule
	۱ & \lr{\texttt{01\_baseline\_BB84\_Lim2014}}        & خط پایه \lr{BB84} با \lr{Lim2014}            & ۹۶  \\
	۲ & \lr{\texttt{02\_finite\_key\_scaling}}            & مقیاس پالس $10^4$ تا $10^8$       & ۲۸۸ \\
	۳ & \lr{\texttt{03\_dwdm\_*}}                          & \lr{DWDM} با توان‌های $-20$، $-34$ و $-40$ \lr{dBm} & ۲۸۸ \\
	۴ & \lr{\texttt{04\_auto\_optimized\_pulses}}        & بهینه‌سازی \lr{SLSQP} پارامتر پالس       & ۹۶  \\
	۵ & \lr{\texttt{05\_protocol\_*}}                      & چهار پروتکل مختلف                   & ۳۸۴ \\
	۶ & \lr{\texttt{06\_proof\_*}}                         & چهار اثبات امنیت متفاوت             & ۳۸۴ \\
	۷ & \lr{\texttt{07\_detector\_*}}                      & سه نوع آشکارساز \lr{SPD}/\lr{SNSPD}/\lr{PNRD}     & ۲۸۸ \\
	۸ & \lr{\texttt{08\_confidence\_*}}                    & سه روش بازهٔ اطمینان                 & ۲۸۸ \\
	۹ & \lr{\texttt{09\_rogers\_validation\_dense}}      & اعتبارسنجی \lr{Rogers} با ۵۱ نقطه        & ۱۰۲ \\
	۱۰ & \lr{\texttt{10\_qber\_misalignment\_sweep}}      & حساسیت به \lr{QBER} و \lr{misalignment}       & ۱۹۸ \\
	\bottomrule
\end{tabular}
\end{table}

\section{پیکربندی مشترک}
تمام کمپین‌ها از مقادیر پیش‌فرض زیر برای اطمینان از قابلیت مقایسه استفاده می‌کنند: افت فیبر \lr{0.2 dB/km}، بازهٔ فاصلهٔ \lr{0} تا \lr{150 km} با ۱۶ نقطه، دانهٔ تصادفی (\lr{seed}) برابر \lr{12345}، و هشت پردازشگر موازی. تعداد پالس‌ها مگر در کمپین ۲، در بازهٔ $10^4$ تا $10^6$ نگه داشته شده است. هر نقطه در دو حالت \lr{\texttt{noise\_applied=True}} و \lr{\texttt{noise\_applied=False}} اجرا می‌شود تا تأثیر نویز واقعی آشکارساز (شامل شمارش تاریک، پس‌پالس، جیتر، زمان مرده و عدم انطباق پایه‌ها) به‌صراحت قابل ارزیابی باشد.

% ---------------------------------------------------------------------
\chapter{نتایج خط پایه و مقیاس کلید محدود}
\label{ch:results_baseline}

\section{خط پایه BB84 با اثبات Lim2014}
شکل~\ref{fig:01_baseline} منحنی نرخ کلید امن به‌عنوان تابعی از فاصلهٔ فیبر را برای پروتکل \lr{BB84-Decoy} با اثبات \lr{Lim2014} نشان می‌دهد. این منحنی به‌عنوان مرجع اصلی گزارش عمل می‌کند و همهٔ مقایسه‌های بعدی نسبت به آن تفسیر می‌شوند. در حالت بدون نویز، نرخ کلید در فواصل کوتاه به مقدار نظری نزدیک می‌شود و با افزایش فاصله به‌دلیل افت نمایی فیبر ($\eta \propto 10^{-\alpha L/10}$) کاهش می‌یابد. در حالت با نویز واقعی، افت اضافی ناشی از شمارش تاریک (\lr{dark count}) و عدم انطباق پایه‌ها (\lr{misalignment}) باعث می‌شود نرخ کلید در فواصل طولانی‌تر سریع‌تر به صفر برسد.

\begin{figure}[htbp]
\centering
\safeimg[width=0.85\textwidth]{figures/01_baseline_BB84_Lim2014.csv.png}
\caption{نرخ کلید امن به‌ازای پالس در برابر فاصلهٔ فیبر برای پروتکل \lr{BB84-Decoy} با اثبات \lr{Lim2014}. دو منحنی حالت بدون نویز و با نویز واقعی آشکارساز مقایسه شده‌اند.}
\label{fig:01_baseline}
\end{figure}

نکتهٔ کلیدی در این شکل، شکستگی منحنی در فاصله‌ای است که در آن نرخ کلید به صفر می‌رسد؛ این نقطه نشان‌دهندهٔ حداکثر برد امن پروتکل در شرایط داده‌شده است. مقادیر دقیق این برد برای ترکیب‌های مختلف پارامتر در فصل‌های بعدی مقایسه می‌شود. علاوه بر این، ستون \lr{\texttt{status}} در \lr{CSV} خروجی، نقاطی که به وضعیت \lr{\texttt{ZERO\_KEY}} ختم شده‌اند را علامت‌گذاری می‌کند که برای تعیین برد امن معتبر هستند.

\section{مقیاس کلید محدود (\lr{Finite-Key Scaling})}
شکل~\ref{fig:02_finite_key} اثر اندازهٔ کلید بر نرخ کلید امن را در پنج اندازهٔ پالس از $10^4$ تا $10^7$ نشان می‌دهد. در اندازه‌های کوچک، جریمهٔ کلید محدود (\lr{finite-size penalty}) نرخ امن را به‌طور قابل توجهی کاهش می‌دهد؛ زیرا تخمین پارامتر با تعداد محدودی نمونه انجام می‌شود و بازهٔ اطمینان \lr{QBER} گسترده‌تر است. با افزایش تعداد پالس به $10^7$، نرخ کلید به مقدار مجانبی نزدیک می‌شود.

\begin{figure}[htbp]
\centering
\safeimg[width=0.85\textwidth]{figures/02_finite_key_scaling.csv.png}
\caption{اثر تعداد پالس بر نرخ کلید امن. با افزایش $N$ از $10^4$ به $10^7$، جریمهٔ کلید محدود کاهش می‌یابد و نرخ به مقدار مجانبی همگرا می‌شود.}
\label{fig:02_finite_key}
\end{figure}

این نتیجه اهمیت انتخاب اندازهٔ بلوک مناسب در پیاده‌سازی عملی \lr{QKD} را تأیید می‌کند؛ برای کاربردهای بلندبرد، اندازهٔ بلوک حداقل $10^6$ پالس ضروری است تا جریمهٔ \lr{finite-size} چشم‌پوشی‌پذیر شود. در مقابل، برای کاربردهای کوتاه‌برد با نرخ بالای رویداد، حتی $10^5$ پالس ممکن است کافی باشد.

\section{اعتبارسنجی Rogers 2007}
شکل~\ref{fig:09_rogers} مقایسهٔ نرخ کلید مونت‌کارلو (\lr{Monte Carlo}) به‌دست آمده از شبیه‌سازی کامل را با فرمول تحلیلی \lr{Rogers et al. 2007} (رابطهٔ ۱۵) برای نرخ بیت غربال‌شده (\lr{sifted bit}) نشان می‌دهد. این مقایسه در ۵۱ نقطهٔ متراکم از ۰ تا ۱۰۰ کیلومتر انجام شده تا هم‌خوانی شبیه‌سازی با تئوری به‌خوبی بررسی شود. ستون \lr{\texttt{analytic\_rogers\_sbr}} در فایل \lr{CSV} به‌طور خودکار برای هر نقطه محاسبه می‌شود و این امکان را فراهم می‌کند که نمودار دقیقاً بازتولید اشکال ۳ تا ۵ مقالهٔ \lr{Rogers} باشد.

\begin{figure}[htbp]
\centering
\safeimg[width=0.85\textwidth]{figures/09_rogers_validation_dense.csv.png}
\caption{مقایسهٔ نرخ \lr{sifted bit} تحلیلی \lr{Rogers 2007} با نرخ کلید امن مونت‌کارلو. هم‌خوانی دو منحنی، اعتبارسنجی صحت پیاده‌سازی لایهٔ فیزیکی و آشکارسازی را تأیید می‌کند.}
\label{fig:09_rogers}
\end{figure}


% ---------------------------------------------------------------------
\chapter{سناریوهای DWDM و بهینه‌سازی پالس}
\label{ch:results_dwdm_opt}

\section{تأثیر هم‌زیستی با کانال کلاسیک (\lr{DWDM})}
شکل~\ref{fig:03_dwdm_power} سه نمودار کنار هم را برای توان‌های گیرندهٔ مختلف \lr{$-20$}، \lr{$-34$} و \lr{$-40$ dBm} نشان می‌دهد و اثر نویز رامان بر نرخ کلید امن را روشن می‌سازد. در توان بالاتر (\lr{$-20$ dBm})، نویز رامان نرخ تیرگی مؤثر آشکارساز را به‌طور چشمگیری افزایش می‌دهد و حداکثر برد امن را به‌میزان قابل توجهی کاهش می‌دهد. در توان پایین‌تر (\lr{$-40$ dBm})، رفتار به حالت فیبر اختصاصی نزدیک می‌شود و افت فیلتر \lr{AWG} (۳ دسی‌بل) تنها منبع انحراف باقی می‌ماند.


\begin{figure}[htbp]
	\centering
	\safeimg[width=0.85\textwidth]{figures/03_dwdm_4channel.csv.png}
	\caption{توان $-34$ dBm (پیش‌فرض Lim2014)}
	\label{fig:03a_dwdm_default}
\end{figure}
\begin{figure}[htbp]
	\centering
	\safeimg[width=0.85\textwidth]{figures/03b_dwdm_power_minus20.csv.png}
	\caption{توان $-20$ dBm (نویز شدید)}
	\label{fig:03b_dwdm_high}
\end{figure}
\begin{figure}[htbp]
	\centering
	\safeimg[width=0.85\textwidth]{figures/03c_dwdm_power_minus40.csv.png}
	\caption{توان $-40$ dBm (نویز کم)}
	\label{fig:03c_dwdm_low}
\end{figure}


این نتایج با بخش IV مقالهٔ \lr{Lim 2014} همخوانی دارد و نشان می‌دهد که طراحی سیستم \lr{DWDM} برای \lr{QKD} نیازمند موازنهٔ دقیق بین توان کانال کلاسیک (برای ارتباط با نرخ بالا) و نویز رامان تحمیل‌شده بر کانال کوانتومی است. مقدار پیش‌فرض \lr{$-34$ dBm} نقطهٔ تعادل پیشنهادی مقاله است که در این شبیه‌سازی بازتولید شده است.

\section{بهینه‌سازی خودکار پارامتر پالس}
شکل~\ref{fig:04_auto_opt} نتیجهٔ بهینه‌سازی \lr{SLSQP} برای پنج پارامتر پالس $(\mu, \nu, p_s, p_d, p_v)$ به‌ازای هر فاصله را نشان می‌دهد. بهینه‌ساز از نرخ \lr{GLLP} مجانبی با جریمهٔ کلید محدود به‌عنوان تابع هدف استفاده می‌کند و شدت سیگنال $\mu$ را در فواصل کوتاه بالا و در فواصل طولانی پایین نگه می‌دارد تا تعادل بین بازدهی آشکارسازی و افشای اطلاعات بهینه شود.

\begin{figure}[htbp]
\centering
\safeimg[width=0.85\textwidth]{figures/04_auto_optimized_pulses.csv.png}
\caption{نرخ کلید امن با بهینه‌سازی خودکار پارامتر پالس در برابر خط پایه با پارامترهای ثابت. بهینه‌ساز \lr{SLSQP} پنج پارامتر $(\mu, \nu, p_s, p_d, p_v)$ را به‌ازای هر فاصله تنظیم می‌کند.}
\label{fig:04_auto_opt}
\end{figure}

مقایسه با شکل~\ref{fig:01_baseline} نشان می‌دهد که بهینه‌سازی خودکار، نرخ کلید را به‌ویژه در فواصل میانی (۲۰ تا ۸۰ کیلومتر) بهبود می‌بخشد. در فواصل بسیار کوتاه، پارامترهای ثابت پیش‌فرض از قبل نزدیک به بهینه هستند و سود بهینه‌سازی ناچیز است. در فواصل بسیار طولانی، بهینه‌سازی می‌تواند برد امن را چند کیلومتر افزایش دهد اما نرخ نهایی همچنان به صفر میل می‌کند.

% ---------------------------------------------------------------------
\chapter{مقایسه‌های سیستماتیک}
\label{ch:results_comparisons}

\section{مقایسهٔ پروتکل‌ها}
شکل~\ref{fig:05_protocols} چهار پروتکل \lr{BB84-Decoy}، \lr{B92}، \lr{MDI-QKD} و \lr{Redundant Transmission} را در شرایط یکسان مقایسه می‌کند. پروتکل \lr{BB84-Decoy} در فواصل کوتاه تا متوسط بهترین عملکرد را دارد؛ \lr{MDI-QKD} به‌دلیل نیاز به آشکارسازی وسط مسیر، نرخ پایین‌تری دارد اما در برابر حملات آشکارساز مقاوم است و در فواصل طولانی نسبتاً پایدارتر عمل می‌کند. \lr{B92} ساده‌ترین پروتکل است اما نرخ پایین آن، کاربردش را در سناریوهای عملی محدود می‌سازد.

\begin{figure}[htbp]
\centering
\safeimg[width=0.85\textwidth]{figures/05_protocols_comparison.png}
\caption{مقایسهٔ نرخ کلید امن چهار پروتکل \lr{BB84-Decoy}، \lr{B92}، \lr{MDI-QKD} و \lr{Redundant Transmission} در شرایط فیبر یکسان.}
\label{fig:05_protocols}
\end{figure}

\textbf{یادداشت:} اگر نمودار فوق به‌صورت چهار فایل جداگانه ذخیره شده است، می‌توانید به‌جای شکل واحد بالا، چهار \lr{subfigure} کنار هم قرار دهید. در آن صورت، نام فایل‌ها به‌ترتیب \lr{\texttt{05\_protocol\_BB84DecoyProtocol.csv.png}}، \lr{\texttt{05\_protocol\_B92Protocol.csv.png}}، \lr{\texttt{05\_protocol\_MDIQKDProtocol.csv.png}} و \lr{\texttt{05\_protocol\_RedundantTransmissionProtocol.csv.png}} خواهند بود.

\section{مقایسهٔ اثبات‌های امنیتی}
شکل~\ref{fig:06_proofs} چهار اثبات امنیت \lr{Lim2014}، \lr{Ma2005}، \lr{Wang2005} و \lr{Tight} را روی همان پروتکل \lr{BB84-Decoy} مقایسه می‌کند. اثبات \lr{Lim2014} پیش‌فرض چارچوب است و تعادل خوبی بین سخت‌گیری و کارایی برقرار می‌کند. اثبات \lr{Tight} نرخ بالاتری تولید می‌کند اما فرضیات قوی‌تری دارد. اثبات‌های \lr{Ma2005} و \lr{Wang2005} برای سناریوهای خاص (مثل منابع نوری غیرایده‌آل) طراحی شده‌اند و در حالت استاندارد ممکن است نرخ پایین‌تری ارائه دهند.

\begin{figure}[htbp]
\centering
\safeimg[width=0.85\textwidth]{figures/06_proofs_comparison.png}
\caption{مقایسهٔ نرخ کلید امن بر اساس چهار اثبات امنیت مختلف برای پروتکل \lr{BB84-Decoy}.}
\label{fig:06_proofs}
\end{figure}

تفاوت بین این چهار منحنی، حاشیهٔ امنیت (\lr{security margin}) هر اثبات را کمّی می‌کند؛ هرچه اثبات سخت‌گیرانه‌تر باشد، حاشیه بزرگ‌تر و نرخ نهایی پایین‌تر است. انتخاب اثبات در عمل به مدل تهدید و میزان اطمینان به منبع نوری و آشکارساز بستگی دارد.

\section{مقایسهٔ آشکارسازها}
شکل~\ref{fig:07_detectors} سه نوع آشکارساز \lr{SPD}، \lr{SNSPD} و \lr{PNRD} را مقایسه می‌کند. آشکارساز \lr{SNSPD} به‌دلیل نرخ شمارش تاریک بسیار پایین ($\sim 10^{-8}$) و راندمان بالا ($\sim 0.5$)، بهترین عملکرد را در فواصل طولانی دارد. آشکارساز \lr{SPD} راندمان پایین‌تری دارد اما هزینه و پیچیدگی کمتری داشته و برای فواصل کوتاه مناسب است. آشکارساز \lr{PNRD} (آشکارساز شمارش فوتون متمایز) در سناریوهایی که نیاز به تشخیص دقیق تعداد فوتون‌ها است، کاربرد دارد.

\begin{figure}[htbp]
\centering
\safeimg[width=0.85\textwidth]{figures/07_detectors_comparison.png}
\caption{مقایسهٔ نرخ کلید امن با سه نوع آشکارساز \lr{SPD}، \lr{SNSPD} و \lr{PNRD}.}
\label{fig:07_detectors}
\end{figure}

این مقایسه نشان می‌دهد که ارتقای آشکارساز از \lr{SPD} به \lr{SNSPD} می‌تواند برد امن را تا دو برابر افزایش دهد، اما هزینه و نیاز به سرمایش کرایوژنیک ملاحظات عملی مهمی هستند که باید در طراحی سیستم لحاظ شوند.

\section{مقایسهٔ روش‌های بازهٔ اطمینان}
شکل~\ref{fig:08_confidence} سه روش \lr{Gaussian}، \lr{Clopper-Pearson} و \lr{Hoeffding} را برای محاسبهٔ بازهٔ اطمینان \lr{QBER} مقایسه می‌کند. روش \lr{Clopper-Pearson} سخت‌گیرانه‌ترین (محافظه‌کارانه‌ترین) روش است و بازه‌ای گسترده‌تر تولید می‌کند که منجر به کاهش نرخ کلید نهایی می‌شود. \lr{Gaussian} روش تقریبی است که برای نمونه‌های بزرگ بهینه است و نرخ بالاتری می‌دهد. \lr{Hoeffding} حد متوسطی ارائه می‌کند که بدون فرض توزیع خاص، همواره معتبر است.

\begin{figure}[htbp]
\centering
\safeimg[width=0.85\textwidth]{figures/08_confidence_comparison.png}
\caption{تأثیر روش محاسبهٔ بازهٔ اطمینان بر نرخ کلید امن. \lr{Clopper-Pearson} محافظه‌کارانه‌ترین و \lr{Gaussian} کارآمدترین روش است.}
\label{fig:08_confidence}
\end{figure}

انتخاب روش به تعداد پالس و سطح اطمینان مورد نظر بستگی دارد؛ برای $N < 10^5$، روش \lr{Clopper-Pearson} به‌دلیل دقت بهتر در نمونه‌های کوچک توصیه می‌شود و برای $N \geq 10^6$، تقریب \lr{Gaussian} کافی است و نرخ کارآمدتری به دست می‌دهد.

% ---------------------------------------------------------------------
\chapter{تحلیل حساسیت}
\label{ch:results_sensitivity}

\section{حساسیت به QBER ذاتی و \lr{Misalignment}}
شکل‌های~\ref{fig:10a_sensitivity} و~\ref{fig:10b_sensitivity} تأثیر هم‌زمان \lr{QBER} ذاتی آشکارساز و جابه‌جایی زاویه‌ای پایه‌ها (\lr{misalignment}) بر نرخ کلید امن را نشان می‌دهند. این دو پارامتر مهم‌ترین منابع خطا در سیستم‌های \lr{QKD} عملی هستند و افزایش هر کدام به‌طور مستقیم \lr{QBER} کلی را بالا برده و نرخ کلید امن را کاهش می‌دهد. در مقادیر بالا (مثلاً \lr{\texttt{qber\_intrinsic=0.02}} و \lr{\texttt{misalignment=0.01}})، حتی در فواصل کوتاه نیز نرخ کلید به‌طور قابل توجهی افت می‌کند.

\begin{figure}[htbp]
\centering
\safeimg[width=0.95\textwidth]{figures/10a_qber_misalignment_sweep.csv.png}
\caption{تحلیل حساسیت نرخ کلید امن به \lr{QBER} ذاتی آشکارساز.}
\label{fig:10a_sensitivity}
\end{figure}

\begin{figure}[htbp]
\centering
\safeimg[width=0.95\textwidth]{figures/10b_qber_misalignment_sweep.csv.png}
\caption{تحلیل حساسیت نرخ کلید امن به عدم انطباق پایه‌ها (\lr{misalignment}).}
\label{fig:10b_sensitivity}
\end{figure}

این تحلیل حساسیت دو پیامد عملی دارد: نخست، اهمیت کالیبراسیون دقیق پایه‌ها در زمان نصب سیستم \lr{QKD}؛ چرا که \lr{misalignment} حتی در حد چند درجه می‌تواند نرخ کلید را تا ۵۰ درصد کاهش دهد. دوم، انتخاب آشکارساز با \lr{QBER} ذاتی پایین (مانند \lr{SNSPD} در مقایسه با \lr{SPD}) تأثیر مستقیمی بر عملکرد نهایی سیستم دارد.


	
\end{document}
